%%%%%%%%%%%%%%%%%%%%%%%%%%%%%%%%%%%%%%%%%%%%%%%%%%%%%%%%%%%%%%%%%%%%%%%%%%%%%%%
% FIN571: Financial Institutions \& Monetary Policy - 통합본
% 자동 생성됨
%%%%%%%%%%%%%%%%%%%%%%%%%%%%%%%%%%%%%%%%%%%%%%%%%%%%%%%%%%%%%%%%%%%%%%%%%%%%%%%

\documentclass[11pt,a4paper]{book}

%========================================================================================
% 기본 패키지
%========================================================================================

\usepackage{kotex}
\usepackage[top=20mm, bottom=20mm, left=20mm, right=18mm]{geometry}
\usepackage{setspace}
\onehalfspacing
\setlength{\parskip}{0.5em}
\setlength{\parindent}{0pt}

% --- 표 관련 ---
\usepackage{booktabs}
\usepackage{tabularx}
\usepackage{array}
\usepackage{longtable}
\usepackage{adjustbox}
\usepackage{colortbl}
\renewcommand{\arraystretch}{1.1}

%========================================================================================
% 헤더 및 푸터
%========================================================================================

\usepackage{fancyhdr}
\pagestyle{fancy}
\fancyhf{}
\fancyhead[LE,RO]{\thepage}
\fancyhead[LO]{\leftmark}
\fancyhead[RE]{FIN571}
\renewcommand{\headrulewidth}{0.5pt}
\renewcommand{\footrulewidth}{0.3pt}

%========================================================================================
% 색상 정의
%========================================================================================

\usepackage[dvipsnames]{xcolor}

\definecolor{lightblue}{RGB}{220, 235, 255}
\definecolor{lightgreen}{RGB}{220, 255, 235}
\definecolor{lightyellow}{RGB}{255, 250, 220}
\definecolor{lightpurple}{RGB}{240, 230, 255}
\definecolor{lightgray}{gray}{0.95}
\definecolor{lightpink}{RGB}{255, 235, 245}
\definecolor{boxgray}{gray}{0.95}
\definecolor{boxblue}{rgb}{0.9, 0.95, 1.0}
\definecolor{boxred}{rgb}{1.0, 0.95, 0.95}

\definecolor{darkblue}{RGB}{50, 80, 150}
\definecolor{darkgreen}{RGB}{40, 120, 70}
\definecolor{darkorange}{RGB}{200, 100, 30}
\definecolor{darkpurple}{RGB}{100, 60, 150}

% MIT/Harvard 추가 색상
\definecolor{mainblue}{RGB}{50, 80, 150}
\definecolor{subgray}{RGB}{180, 180, 180}
\definecolor{codebg}{RGB}{245, 245, 245}
\definecolor{boxbgcolor}{RGB}{245, 245, 250}

% Stanford CS230 스타일 색상
\definecolor{myblue}{RGB}{50, 100, 180}
\definecolor{sblue}{RGB}{70, 130, 180}
\definecolor{wgray}{RGB}{245, 245, 245}
\definecolor{codegreen}{rgb}{0,0.6,0}
\definecolor{codegray}{rgb}{0.5,0.5,0.5}
\definecolor{codepurple}{rgb}{0.58,0,0.82}
\definecolor{backcolour}{rgb}{0.95,0.95,0.92}

%========================================================================================
% 박스 환경
%========================================================================================

\usepackage[most]{tcolorbox}
\tcbuselibrary{skins, breakable}

\newtcolorbox{overviewbox}[1][]{
    enhanced,
    colback=lightpurple,
    colframe=darkpurple,
    fonttitle=\bfseries\large,
    title=강의 개요,
    arc=3mm,
    boxrule=1pt,
    breakable,
    #1
}

\newtcolorbox{summarybox}[1][]{
    enhanced,
    colback=lightblue,
    colframe=darkblue,
    fonttitle=\bfseries,
    title=핵심 요약,
    arc=2mm,
    boxrule=0.7pt,
    breakable,
    #1
}

\newtcolorbox{infobox}[1][]{
    enhanced,
    colback=lightgreen,
    colframe=darkgreen,
    fonttitle=\bfseries,
    title=핵심 정보,
    arc=2mm,
    boxrule=0.7pt,
    breakable,
    #1
}

\newtcolorbox{warningbox}[1][]{
    enhanced,
    colback=lightyellow,
    colframe=darkorange,
    fonttitle=\bfseries,
    title=주의,
    arc=2mm,
    boxrule=0.7pt,
    breakable,
    #1
}

\newtcolorbox{examplebox}[1][]{
    enhanced,
    colback=lightgray,
    colframe=black!60,
    fonttitle=\bfseries,
    title=예제,
    arc=2mm,
    boxrule=0.7pt,
    breakable,
    #1
}

\newtcolorbox{definitionbox}[1][]{
    enhanced,
    colback=lightpink,
    colframe=purple!70!black,
    fonttitle=\bfseries,
    title=정의,
    arc=2mm,
    boxrule=0.7pt,
    breakable,
    #1
}

\newtcolorbox{importantbox}[1][]{
    enhanced,
    colback=boxred,
    colframe=red!70!black,
    fonttitle=\bfseries,
    title=매우 중요,
    arc=2mm,
    boxrule=0.7pt,
    breakable,
    #1
}

\newtcolorbox{conceptbox}[1][]{
    enhanced,
    colback=lightgreen,
    colframe=darkgreen,
    fonttitle=\bfseries,
    title=개념,
    arc=2mm,
    boxrule=0.7pt,
    breakable,
    #1
}

\newtcolorbox{analogybox}[1][]{
    enhanced,
    colback=green!5!white,
    colframe=green!60!black,
    fonttitle=\bfseries,
    title=비유,
    arc=2mm,
    boxrule=0.7pt,
    breakable,
    #1
}

\let\cautionbox\warningbox
\let\endcautionbox\endwarningbox

% tcbset 스타일 정의 (\begin{tcolorbox}[summarybox] 형식 사용 시 호환)
\tcbset{
    summarybox/.style={enhanced, colback=lightblue, colframe=darkblue, fonttitle=\bfseries, title=핵심 요약, arc=2mm, boxrule=0.7pt, breakable},
    warningbox/.style={enhanced, colback=lightyellow, colframe=darkorange, fonttitle=\bfseries, title=주의, arc=2mm, boxrule=0.7pt, breakable},
    examplebox/.style={enhanced, colback=lightgray, colframe=black!60, fonttitle=\bfseries, title=예제, arc=2mm, boxrule=0.7pt, breakable},
    definitionbox/.style={enhanced, colback=lightpink, colframe=purple!70!black, fonttitle=\bfseries, title=정의, arc=2mm, boxrule=0.7pt, breakable},
    importantbox/.style={enhanced, colback=boxred, colframe=red!70!black, fonttitle=\bfseries, title=매우 중요, arc=2mm, boxrule=0.7pt, breakable},
    infobox/.style={enhanced, colback=lightgreen, colframe=darkgreen, fonttitle=\bfseries, title=핵심 정보, arc=2mm, boxrule=0.7pt, breakable},
    conceptbox/.style={enhanced, colback=lightgreen, colframe=darkgreen, fonttitle=\bfseries, title=개념, arc=2mm, boxrule=0.7pt, breakable},
    analogybox/.style={enhanced, colback=green!5!white, colframe=green!60!black, fonttitle=\bfseries, title=비유, arc=2mm, boxrule=0.7pt, breakable},
    overviewbox/.style={enhanced, colback=lightpurple, colframe=darkpurple, fonttitle=\bfseries\large, title=강의 개요, arc=3mm, boxrule=1pt, breakable},
    cautionbox/.style={enhanced, colback=lightyellow, colframe=darkorange, fonttitle=\bfseries, title=주의, arc=2mm, boxrule=0.7pt, breakable},
    mybox/.style={enhanced, colback=lightblue, colframe=darkblue, fonttitle=\bfseries, title={#1}, arc=2mm, boxrule=0.7pt, breakable},
    definition/.style={enhanced, colback=lightpink, colframe=purple!70!black, fonttitle=\bfseries, title={#1}, arc=2mm, boxrule=0.7pt, breakable},
    example/.style={enhanced, colback=lightgray, colframe=black!60, fonttitle=\bfseries, title={#1}, arc=2mm, boxrule=0.7pt, breakable},
    summary/.style={enhanced, colback=lightblue, colframe=darkblue, fonttitle=\bfseries, title={#1}, arc=2mm, boxrule=0.7pt, breakable},
    warning/.style={enhanced, colback=lightyellow, colframe=darkorange, fonttitle=\bfseries, title={#1}, arc=2mm, boxrule=0.7pt, breakable},
}

% mybox 환경 (UIUC용)
\newtcolorbox{mybox}[1]{
    enhanced,
    colback=lightblue,
    colframe=darkblue,
    fonttitle=\bfseries,
    title={#1},
    arc=2mm,
    boxrule=0.7pt,
    breakable
}

% 추가 환경 (Gemini 생성 콘텐츠에서 사용)
\newtcolorbox{tipbox}[1][]{enhanced, colback=green!5!white, colframe=green!60!black, fonttitle=\bfseries, title=팁, arc=2mm, boxrule=0.7pt, breakable, #1}
\newtcolorbox{solutionbox}[1][]{enhanced, colback=lightblue, colframe=darkblue, fonttitle=\bfseries, title=풀이, arc=2mm, boxrule=0.7pt, breakable, #1}
\newtcolorbox{downloadbox}[1][]{enhanced, colback=boxgray, colframe=black!40, fonttitle=\bfseries, title=다운로드, arc=2mm, boxrule=0.7pt, breakable, #1}
\newtcolorbox{practicebox}[1][]{enhanced, colback=lightyellow, colframe=darkorange, fonttitle=\bfseries, title=연습, arc=2mm, boxrule=0.7pt, breakable, #1}
\newtcolorbox{keypointbox}[1][]{enhanced, colback=lightyellow, colframe=darkorange, fonttitle=\bfseries, title=핵심 포인트, arc=2mm, boxrule=0.7pt, breakable, #1}
\newtcolorbox{storybox}[1][]{enhanced, colback=lightyellow, colframe=darkorange!80!black, fonttitle=\bfseries, title=이야기, arc=2mm, boxrule=0.7pt, breakable, #1}
\newtcolorbox{mathstep}[1][]{enhanced, colback=lightblue, colframe=darkblue, fonttitle=\bfseries, title=수학 단계, arc=2mm, boxrule=0.7pt, breakable, #1}
\newtcolorbox{systembox}[1][]{enhanced, colback=lightgray, colframe=black!60, fonttitle=\bfseries, title=시스템, arc=2mm, boxrule=0.7pt, breakable, #1}
\newtcolorbox{toolbox}[1][]{enhanced, colback=boxgray, colframe=black!40, fonttitle=\bfseries, title=도구, arc=2mm, boxrule=0.7pt, breakable, #1}
\newtcolorbox{bayesbox}[1][]{enhanced, colback=lightpurple, colframe=darkpurple, fonttitle=\bfseries, title=베이즈, arc=2mm, boxrule=0.7pt, breakable, #1}
\newtcolorbox{formulabox}[1][]{enhanced, colback=lightblue, colframe=darkblue, fonttitle=\bfseries, title=공식, arc=2mm, boxrule=0.7pt, breakable, #1}
\newtcolorbox{mathbox}[1][]{enhanced, colback=lightblue, colframe=darkblue, fonttitle=\bfseries, title=수학, arc=2mm, boxrule=0.7pt, breakable, #1}
\newtcolorbox{strategybox}[1][]{enhanced, colback=lightgreen, colframe=darkgreen, fonttitle=\bfseries, title=전략, arc=2mm, boxrule=0.7pt, breakable, #1}
\newtcolorbox{diagnosisbox}[1][]{enhanced, colback=lightpink, colframe=purple!70!black, fonttitle=\bfseries, title=진단, arc=2mm, boxrule=0.7pt, breakable, #1}
\newtcolorbox{featurebox}[1][]{enhanced, colback=lightgreen, colframe=darkgreen, fonttitle=\bfseries, title=특징, arc=2mm, boxrule=0.7pt, breakable, #1}
\newtcolorbox{faqbox}[1][]{enhanced, colback=lightyellow, colframe=darkorange, fonttitle=\bfseries, title=FAQ, arc=2mm, boxrule=0.7pt, breakable, #1}
\newtcolorbox{notebox}[1][]{enhanced, colback=lightyellow, colframe=darkorange, fonttitle=\bfseries, title=노트, arc=2mm, boxrule=0.7pt, breakable, #1}
\newtcolorbox{codeexamplebox}[1][]{enhanced, colback=lightgray, colframe=black!60, fonttitle=\bfseries, title=코드 예제, arc=2mm, boxrule=0.7pt, breakable, #1}
\newtcolorbox{keyconceptbox}[1][]{enhanced, colback=lightgreen, colframe=darkgreen, fonttitle=\bfseries, title=핵심 개념, arc=2mm, boxrule=0.7pt, breakable, #1}
\newtcolorbox{termbox}[1][]{enhanced, colback=lightpurple, colframe=darkpurple, fonttitle=\bfseries, title=용어, arc=2mm, boxrule=0.7pt, breakable, #1}
\let\cautionbox\warningbox
\let\endcautionbox\endwarningbox

% 커스텀 명령어 (Gemini 생성 콘텐츠에서 사용)
\newcommand{\excel}[1]{\texttt{#1}}
\newcommand{\code}[1]{\texttt{#1}}
\newcommand{\concept}[1]{\textbf{#1}}
\newcommand{\formula}[1]{\textit{#1}}
\providecommand{\captionof}[2]{\textbf{#2}}
\newcommand{\chaptertitle}[2]{\section*{#1: #2}}
\providecommand{\newsection}{\section}

%========================================================================================
% 코드 블록
%========================================================================================

\usepackage{listings}

\lstdefinestyle{mystyle}{
    backgroundcolor=\color{backcolour},
    commentstyle=\color{codegreen},
    keywordstyle=\color{magenta},
    numberstyle=\tiny\color{codegray},
    stringstyle=\color{codepurple},
    basicstyle=\ttfamily\footnotesize,
    breakatwhitespace=false,
    breaklines=true,
    captionpos=b,
    keepspaces=true,
    numbers=left,
    numbersep=5pt,
    showspaces=false,
    showstringspaces=false,
    showtabs=false,
    tabsize=2
}

\lstset{
    style=mystyle,
    basicstyle=\ttfamily\small,
    backgroundcolor=\color{lightgray},
    keywordstyle=\color{darkblue}\bfseries,
    commentstyle=\color{darkgreen}\itshape,
    stringstyle=\color{purple!80!black},
    numberstyle=\tiny\color{black!60},
    numbers=left,
    numbersep=8pt,
    breaklines=true,
    frame=single,
    frameround=tttt,
    rulecolor=\color{black!30},
    showstringspaces=false,
    tabsize=2,
    xleftmargin=15pt,
    xrightmargin=5pt
}

%========================================================================================
% 수학
%========================================================================================

\usepackage{amsmath, amssymb, amsthm, amsfonts}

\theoremstyle{definition}
\newtheorem{theorem}{정리}[chapter]
\newtheorem{lemma}[theorem]{보조정리}
\newtheorem{proposition}[theorem]{명제}
\newtheorem{corollary}[theorem]{따름정리}
\newtheorem{definition}{정의}[chapter]
\newtheorem{example}{예제}[chapter]

%========================================================================================
% 하이퍼링크
%========================================================================================

\usepackage[
    colorlinks=true,
    linkcolor=blue!80!black,
    urlcolor=blue!80!black,
    bookmarks=true,
    bookmarksnumbered=true
]{hyperref}

%========================================================================================
% 기타
%========================================================================================

\usepackage{enumitem}
\setlist{nosep, leftmargin=*, itemsep=0.3em}

\usepackage{microtype}
\usepackage{graphicx}
\usepackage{url}
\urlstyle{same}

%========================================================================================
% 메타 정보 박스 명령어
%========================================================================================

\newcommand{\metainfo}[4]{
\begin{tcolorbox}[
    colback=lightpurple,
    colframe=darkpurple,
    boxrule=1pt,
    arc=2mm,
    left=10pt,
    right=10pt,
    top=8pt,
    bottom=8pt
]
\begin{tabular}{@{}rl@{}}
\textbf{강의명:} & #1 \\[0.3em]
\textbf{주차:} & #2 \\[0.3em]
\textbf{교수명:} & #3 \\[0.3em]
\textbf{목적:} & \begin{minipage}[t]{0.75\textwidth}#4\end{minipage}
\end{tabular}
\end{tcolorbox}
}

%========================================================================================
% 문서 시작
%========================================================================================

\title{\textbf{FIN571: Financial Institutions \& Monetary Policy}}
\author{통합 강의 노트}
\date{}

\begin{document}

\maketitle
\tableofcontents


%=======================================================================
% Module 1: Lecture 1
%=======================================================================
\chapter{Lecture 1}
\label{ch:module1}

\metainfo{Banking and Financial Institutions}{Module 1: Money and Finance}{Prof. Rustom Manouchehri Irani}{금융 시스템의 이해와 경제적 영향 분석}


\section{개요}
\addcontentsline{toc}{section}{개요}
이 문서는 'Banking and Financial Institutions Module 1: Money and Finance'의 핵심 내용을 다룹니다.
돈의 정의와 기능, 지불 시스템의 진화, 통화량과 인플레이션의 관계를 탐구합니다.
또한 금융 시스템의 역할, 금융 시장과 금융 기관의 차이, 그리고 다양한 비은행 금융 기관의 기능을 이해합니다.
궁극적으로 금융 발전이 실물 경제 활동에 미치는 긍정적 및 부정적 영향을 종합적으로 파악하는 것이 목표입니다.
이 문서는 시험 대비를 위한 완결된 학습 자료로, 비전공자도 쉽게 이해할 수 있도록 상세한 설명과 예시를 제공합니다.


\section{용어 정리}
\addcontentsline{toc}{section}{용어 정리}
다음은 본 모듈에서 다루는 주요 용어들을 비전공자도 이해하기 쉽게 정리한 표입니다.

\begin{adjustbox}{width=\textwidth,center}
\begin{tabular}{lllc}
\toprule
\textbf{용어 (Term)} & \textbf{쉬운 설명 (Easy Explanation)} & \textbf{원어 (Original)} & \textbf{비고 (Notes)} \\
\midrule
상품 화폐 & 그 자체로 가치가 있는 물건을 돈으로 사용하는 것 & Commodity Money & 조개껍데기, 소금, 금속 등 \\
법정 화폐 & 정부의 명령으로 가치가 부여된 돈 & Fiat Money & 오늘날 대부분의 지폐 \\
지불 시스템 & 돈을 주고받는 방법과 관련된 모든 제도와 인프라 & Payments System & 현금, 카드, 전자 결제 등 \\
M1 & 가장 좁은 의미의 통화량; 현금과 요구불 예금 & M1 (Money Supply 1) & 즉시 사용 가능한 돈 \\
M2 & M1에 저축성 예금 등을 더한 넓은 의미의 통화량 & M2 (Money Supply 2) & M1보다 유동성은 낮지만 쉽게 현금화 가능 \\
인플레이션 & 물가가 지속적으로 오르는 현상 & Inflation & 돈의 가치가 하락함을 의미 \\
금융 시스템 & 돈이 필요한 사람과 돈이 남는 사람을 연결하는 체계 & Financial System & 경제 성장의 핵심 인프라 \\
금융 상품 & 돈을 빌려주거나 투자할 때 주고받는 약속 증서 & Financial Instrument & 주식, 채권, 예금, 보험 등 \\
금융 시장 & 금융 상품이 거래되는 장소 & Financial Market & 주식 시장, 채권 시장 등 \\
금융 기관 & 돈을 모으고 빌려주는 역할을 하는 회사 & Financial Institution & 은행, 보험사, 연기금 등 \\
은행 & 예금을 받고 대출을 해주는 전통적인 금융 기관 & Bank & 상업 은행, 저축 기관, 신용 협동조합 \\
비은행 금융 기관 & 예금을 받지 않으면서 금융 서비스를 제공하는 기관 & Non-Bank Financial Institution & 보험사, 연기금, 증권사, 할부금융사 등 \\
유동성 & 자산을 현금으로 얼마나 쉽고 빠르게 바꿀 수 있는지 & Liquidity & 현금이 가장 유동성이 높음 \\
만기 불일치 & 돈을 빌려주는 기간과 빌리는 기간이 다른 문제 & Maturity Mismatch & 은행이 단기 예금으로 장기 대출을 해줌 \\
정보 비대칭 & 거래 당사자 중 한쪽이 더 많은 정보를 가진 상황 & Information Asymmetry & 대출자의 신용 정보를 은행이 더 잘 앎 \\
정부 후원 기업 & 정부가 설립했으나 민간 형태로 운영되는 기업 & Government Sponsored Enterprise (GSE) & Fannie Mae, Freddie Mac 등 \\
\bottomrule
\end{tabular}
\end{adjustbox}


\section{핵심 개념 및 원리}
\addcontentsline{toc}{section}{핵심 개념 및 원리}

\subsection{1. 돈과 지불 시스템 (Money and Payments)}
\addcontentsline{toc}{subsection}{1. 돈과 지불 시스템}
이 섹션에서는 돈의 본질과 기능, 그리고 현대 사회에서 돈이 어떻게 교환되는지에 대한 지불 시스템을 탐구합니다. 돈은 우리가 매일 사용하는 개념이지만, 그 정의와 형태는 시대에 따라 크게 변화해 왔습니다.

\subsubsection*{1.1. 돈의 정의와 기능 (Money and Its Functions)}
\addcontentsline{toc}{subsubsection}{1.1. 돈의 정의와 기능}

\begin{tcolorbox}[summarybox]
  \textbf{핵심 요약:} 돈은 재화와 서비스의 지불 수단으로 널리 받아들여지는 특별한 자산이며, 교환의 매개, 가치 척도, 가치 저장의 세 가지 핵심 기능을 수행합니다.
\end{tcolorbox}

우리는 일상생활에서 항상 돈을 사용하지만, '돈이 무엇인가?'라는 질문은 생각보다 복잡합니다. 돈의 형태는 역사적으로 많이 변해왔고, 오늘날에도 그 정의에 대한 논의가 계속되고 있습니다.

\textbf{돈의 진화: 상품 화폐에서 법정 화폐까지}
돈의 역사를 살펴보면, 돈으로 사용된 자산들이 어떻게 변화해왔는지 알 수 있습니다.

\begin{itemize}
    \item \textbf{상품 화폐 (Commodity Money):}
    \begin{itemize}
        \item \textbf{한 줄 핵심 요약:} 그 자체로도 가치가 있는 물건을 돈으로 사용하는 형태입니다.
        \item \textbf{직관적 예시/비유:} 8세기 일본의 화살촉, 중국의 조개껍데기, 베네치아의 소금, 금과 같은 금속들이 상품 화폐의 예시입니다. 소금은 음식을 보존하는 데 유용하고, 화살촉은 사냥이나 방어에 사용될 수 있습니다.
        \item \textbf{기술적 정확한 설명:} 상품 화폐는 두 가지 특징을 가집니다. 첫째, 물리적인 물품에 의해 가치가 뒷받침됩니다. 둘째, 그 물품 자체에 본질적인 가치(intrinsic value)가 있습니다. 즉, 지불 수단으로 사용되지 않더라도 그 자체로 유용하거나 가치가 있는 것입니다. 16세기 스웨덴에서 사용된 구리판도 상품 화폐의 일종이었으나, 당시 구리의 가치가 낮아 거래를 위해 많은 양을 들고 다녀야 하는 불편함이 있었습니다.
    \end{itemize}

    \item \textbf{종이 화폐의 등장 (Paper Money):}
    \begin{itemize}
        \item \textbf{한 줄 핵심 요약:} 금속 화폐를 은행에 예치하고 그 증서로 거래하는 방식에서 시작되었습니다.
        \item \textbf{직관적 예시/비유:} 1657년 스웨덴의 Palmstruch가 설립한 Stockholms Banco는 유럽 최초로 종이 화폐를 발행했습니다. 이 종이 화폐는 은행에 예치된 구리 금속으로 언제든지 교환(redeemable on-demand)할 수 있었습니다. 사람들이 은행에 충분한 구리가 있다고 믿는 한, 종이 화폐는 가치를 가졌습니다.
        \item \textbf{기술적 정확한 설명:} 종이 화폐는 처음에는 예치된 상품 화폐(주로 금이나 은)에 의해 가치가 보증되었습니다. 미국도 1930년대까지 금본위제(gold standard)를 통해 종이 화폐를 금으로 교환할 수 있도록 했습니다. 그러나 Palmstruch 은행의 사례처럼, 은행이 충분한 금속을 보유하고 있지 않다는 소문이 돌면 사람들은 신뢰를 잃고 은행은 문을 닫을 수 있었습니다.
    \end{itemize}

    \item \textbf{법정 화폐 (Fiat Money):}
    \begin{itemize}
        \item \textbf{한 줄 핵심 요약:} 어떤 물리적 상품에도 의해 뒷받침되지 않고, 정부의 법적 강제력에 의해 가치가 부여되는 돈입니다.
        \item \textbf{직관적 예시/비유:} 오늘날 우리가 사용하는 지폐가 바로 법정 화폐입니다. 100달러 지폐를 인쇄하는 데 14센트밖에 들지 않지만, 우리는 그 지폐가 100달러의 가치를 가진다고 믿습니다. 이는 정부가 "이 지폐는 법정 통화이며, 모든 사적 사업체는 이를 지불 수단으로 받아들여야 한다"고 선언하고, 세금 납부에도 사용할 수 있기 때문입니다.
        \item \textbf{기술적 정확한 설명:} 법정 화폐는 그 자체로는 본질적인 가치가 없습니다. 그 가치는 전적으로 정부에 대한 신뢰와 법적 강제력에 기반합니다. 정부가 안정적이고 신뢰할 수 있는 한, 사람들은 법정 화폐를 가치 있는 지불 수단으로 받아들입니다.
    \end{itemize}
\end{itemize}

\textbf{현대 지불 수단의 진화}
돈의 형태는 종이 화폐를 넘어 플라스틱 카드와 디지털 방식으로 발전했습니다.

\begin{itemize}
    \item \textbf{직불 카드 (Debit Card):}
    \begin{itemize}
        \item \textbf{설명:} 내 은행 계좌에 있는 돈을 판매자에게 보내도록 은행에 지시하는 수단입니다. 계좌에 돈이 없으면 작동하지 않습니다.
    \end{itemize}
    \item \textbf{신용 카드 (Credit Card):}
    \begin{itemize}
        \item \textbf{설명:} 은행으로부터 대출을 받는 것입니다. 카드를 사용할 때 은행이 대신 돈을 지불하고, 나는 나중에 은행에 갚습니다. 1958년 Bank of America가 최초의 현대식 플라스틱 신용카드를 도입했습니다.
    \end{itemize}
    \item \textbf{스마트폰 앱 및 온라인 결제 시스템 (Smartphone Applications and Online Payment Systems):}
    \begin{itemize}
        \item \textbf{설명:} Apple Pay, Venmo, WeChat, PayPal 등은 플라스틱 카드 없이도 빠르고 안전하게 돈을 이체하거나 결제할 수 있게 해주는 21세기 시스템입니다.
    \end{itemize}
\end{itemize}

\textbf{돈의 공식적인 정의}
돈은 특정 자산으로, 재화와 서비스의 지불 또는 부채 상환(세금 납부 포함)을 위해 일반적으로 받아들여지는 수단입니다. 주머니 속 현금과 은행 계좌 잔액이 가장 중요한 예시입니다.

\begin{warningbox}
  \textbf{돈과 부(Wealth)의 차이:} 돈은 부와 같지 않습니다. 부유한 사람은 돈 외에도 집, 요트, 주식, 채권 등 다양한 자산을 가지고 있지만, 이들은 돈이 아닙니다. 돈은 오직 '지불 수단'으로서의 특별한 역할을 합니다.
\end{warningbox}

\textbf{돈의 세 가지 주요 기능}
돈이 지불 수단으로 기능하기 위해서는 다음과 같은 세 가지 특성을 가져야 합니다.

\begin{enumerate}
    \item \textbf{교환의 매개 (Medium of Exchange):}
    \begin{itemize}
        \item \textbf{한 줄 핵심 요약:} 거래를 원활하게 하고 물물교환의 불편함을 해소합니다.
        \item \textbf{직관적 예시/비유:} 내가 사과를 팔고 싶고 신발을 사고 싶을 때, 신발을 파는 사람이 사과를 원하지 않으면 거래가 어렵습니다 (물물교환의 이중 우연의 일치 문제). 돈이 있으면 사과를 팔아 돈을 받고, 그 돈으로 신발을 살 수 있습니다.
        \item \textbf{기술적 정확한 설명:} 현대 경제에서는 거래가 빈번하고 광범위하게 이루어집니다. 돈은 거래를 즉시 완료할 수 있게 하여, 지불을 기다릴 필요 없이 상품과 서비스를 교환할 수 있도록 합니다. 돈의 가치는 널리 이해되고 있어, 거래 시 상대방의 돈 가치를 일일이 확인할 필요가 없습니다.
    \end{itemize}

    \item \textbf{가치 척도 (Unit of Account):}
    \begin{itemize}
        \item \textbf{한 줄 핵심 요약:} 재화와 서비스의 가치를 측정하고 비교하는 기준이 됩니다.
        \item \textbf{직관적 예시/비유:} 모든 상품의 가격이 '달러와 센트'로 표시되어 있다면, 사과 한 개가 1달러이고 바나나 한 개가 50센트라는 것을 쉽게 비교할 수 있습니다.
        \item \textbf{기술적 정확한 설명:} 돈은 가격을 매기고 부채를 기록하는 단위가 됩니다. 이를 통해 다양한 제품의 가치를 쉽게 비교할 수 있으며, 경제 활동의 효율성을 높입니다.
    \end{itemize}

    \item \textbf{가치 저장 수단 (Store of Value):}
    \begin{itemize}
        \item \textbf{한 줄 핵심 요약:} 구매력을 현재에서 미래로 이전할 수 있게 해줍니다.
        \item \textbf{직관적 예시/비유:} 오늘 번 돈을 당장 쓰지 않고 저축했다가 나중에 휴가를 갈 때 사용할 수 있습니다. 돈은 그 가치를 비교적 안정적으로 유지해야 합니다.
        \item \textbf{기술적 정확한 설명:} 돈은 시간이 지나도 그 가치를 비교적 안정적으로 유지해야 합니다. 만약 돈의 가치가 불안정하게 급락한다면 (예: 초인플레이션), 상인들은 돈을 받으려 하지 않을 것입니다. 이 때문에 일부 국가에서는 자국 통화와 함께 미국 달러와 같은 안정적인 통화를 사용하기도 합니다.
    \end{itemize}
\end{enumerate}


\subsubsection*{1.2. 지불 시스템 (The Payments System)}
\addcontentsline{toc}{subsubsection}{1.2. 지불 시스템}

\begin{tcolorbox}[summarybox]
  \textbf{핵심 요약:} 지불 시스템은 재화, 서비스, 자산의 교환을 가능하게 하는 인프라로, 현금 결제와 비현금 결제로 나뉘며, 비현금 결제는 전자화되는 추세입니다.
\end{tcolorbox}

돈과 지불 시스템은 밀접하게 연결되어 있습니다. 지불 시스템은 경제의 '금융 배관(financial plumbing)'이라고 불리며, 상품, 서비스, 자산의 교환을 가능하게 하는 일련의 체계입니다. 이 시스템은 우리가 인식하는 것보다 더 많은 인프라를 필요로 하며, 끊임없이 혁신되고 있습니다. 효율적인 지불 시스템은 경제에 필수적이므로 정부와 정책 입안자들은 이에 큰 관심을 기울입니다. 많은 경우, 연방준비제도(Federal Reserve System)와 같은 중앙은행은 지불을 촉진하기 위해 설립되었습니다.

\textbf{현금 결제와 비현금 결제}
모든 거래가 지폐 교환을 수반하는 것은 아닙니다. 우리는 비현금 결제(non-cash payments)를 많이 사용합니다. 모든 비현금 결제는 궁극적으로 돈의 교환을 수반하며, 주로 한 은행 계좌에서 다른 은행 계좌로 이루어집니다.

\textbf{주요 비현금 결제 유형}

\begin{enumerate}
    \item \textbf{수표 (Paper Checks):}
    \begin{itemize}
        \item \textbf{한 줄 핵심 요약:} 여전히 미국에서 중요한 지불 수단으로, 은행에 자금 이체를 지시하는 서면 명령입니다.
        \item \textbf{직관적 예시/비유:} 누군가에게 수표를 주면, 그 수표는 내 은행에 '내 계좌에서 판매자의 계좌로 돈을 옮겨라'는 지시를 내리는 것과 같습니다. 수표 자체는 돈이나 통화가 아닙니다.
        \item \textbf{기술적 정확한 설명:} 미국에서는 지역 연방준비은행이 종종 수표 청산(clearing)을 담당합니다. 이는 은행 계좌 간의 자금 이체를 처리하는 것을 의미합니다.
    \end{itemize}

    \item \textbf{자동 어음 교환 시스템 (Automated Clearing House, ACH) 거래:}
    \begin{itemize}
        \item \textbf{한 줄 핵심 요약:} 수표와 유사하게 은행 간 자금 이체를 지시하지만, 전자적으로 이루어져 더 빠르고 편리합니다.
        \item \textbf{직관적 예시/비유:} 직불카드나 신용카드로 결제할 때 주로 ACH 이체가 발생합니다. 직불카드는 내 계좌에서, 신용카드는 은행의 자금에서 이체가 이루어지며, 신용카드 사용 시에는 대출이 발생합니다.
        \item \textbf{기술적 정확한 설명:} ACH 이체는 전자적으로 처리되므로 수표보다 행정적 번거로움이 적고 빠릅니다. 선불 직불카드(prepaid debit cards)도 중요한 지불 수단이며, 일부 국가에서는 고객의 은행 계좌와 연결되지 않기도 합니다.
    \end{itemize}

    \item \textbf{디지털 지갑 및 송금 앱 (Digital Wallets and Money Transfer Apps):}
    \begin{itemize}
        \item \textbf{한 줄 핵심 요약:} 컴퓨터나 스마트폰을 통해 전자 직불/신용 이체를 시작할 수 있게 해주는 시스템입니다.
        \item \textbf{직관적 예시/비유:} Apple Pay는 시장 선두 주자이며, Venmo나 Zelle과 같은 앱은 친구나 가족 간의 송금을 가능하게 합니다. 이들은 주로 기존 직불/신용카드 정보를 저장하여 사용합니다.
        \item \textbf{기술적 정확한 설명:} 이러한 시스템은 결제 정보를 디지털 방식으로 저장하고, 사용자가 편리하게 전자 거래를 시작할 수 있도록 합니다.
    \end{itemize}
\end{enumerate}

\begin{tcolorbox}[summarybox]
  \textbf{핵심 요약:} 현금과 은행 계좌 잔액은 경제학자들이 돈의 핵심 척도로 간주합니다. 2018년 미국에서는 약 1,750억 건의 비현금 결제가 약 100조 달러 규모로 이루어졌으며, 전자 결제가 수표를 대체하는 추세입니다.
\end{tcolorbox}

\textbf{비현금 결제 방식별 가치 비중 (2012년 이후 추세)}
제공된 자료에 따르면, 비현금 결제 가치에서 수표가 차지하는 비중은 감소하고 전자 결제가 증가하는 추세입니다. 2012년에는 신용카드 결제액이 수표 결제액을 처음으로 넘어섰습니다. 그럼에도 불구하고 수표는 여전히 중요한 역할을 하고 있습니다.

\textbf{사례 연구: 비트코인은 돈인가? (Is Bitcoin Money?)}
암호화폐는 새로운 지불 시스템과 정부가 발행하지 않는 새로운 통화를 결합한 것입니다. 통화 자체는 사고팔 수 있는 디지털 토큰이며, 거래는 분산원장기술(distributed ledger technology), 즉 블록체인을 사용하여 기록됩니다. 블록체인은 정부 기관의 검증 없이 거래를 기록하고 검증하는 전자 데이터베이스입니다. 2008년부터 존재해 온 비트코인(Bitcoin)이 가장 유명한 예시입니다.

\textbf{비트코인의 잠재적 장점:}
\begin{itemize}
    \item \textbf{정부 통제 불가능:} 정부가 디지털 토큰의 공급을 통제하지 않으므로, 정부에 의해 가치가 훼손될 수 없습니다.
    \item \textbf{익명성:} 거래가 정부 기관의 검증을 필요로 하지 않으므로, 사용자는 익명으로 전자 결제를 할 수 있습니다.
\end{itemize}

\textbf{비트코인이 돈의 세 가지 특성을 충족하는가?}
\begin{enumerate}
    \item \textbf{유용한 지불 수단인가? (Means of Payment)}
    \begin{itemize}
        \item \textbf{평가:} 아직은 아닙니다. 상품 화폐처럼 본질적인 가치가 없고, 법정 화폐도 아니므로 기업이 비트코인을 의무적으로 받아들일 필요가 없습니다. 세금 납부에도 사용할 수 없습니다. 현재까지 비트코인 거래는 주로 마약 거래나 사이버 공격과 같은 범죄 활동과 연관되어 있습니다. 이러한 이유로 중국은 암호화폐 거래를 단속했습니다.
    \end{itemize}
    \item \textbf{안정적인 가치 저장 수단인가? (Stable Store of Value)}
    \begin{itemize}
        \item \textbf{평가:} 아닙니다. 비트코인의 가격은 크게 변동하여 통화라기보다는 투기성 자산에 가깝습니다. 미국 세법상 비트코인 및 기타 암호화폐는 '재산(property)'으로 분류되어 투자로 과세됩니다. 가치 변동이 심하기 때문에 택시 요금이나 테슬라 자동차 가격을 비트코인으로 책정한다면 매우 불안정할 것입니다.
    \end{itemize}
    \item \textbf{가치 척도인가? (Unit of Account)}
    \begin{itemize}
        \item \textbf{평가:} 위의 두 가지 이유로 인해, 비트코인은 안정적인 가치 척도로서의 역할을 제대로 수행하지 못합니다.
    \end{itemize}
\end{enumerate}
\begin{tcolorbox}[warningbox]
  \textbf{결론:} 현재로서는 비트코인이 미국 달러나 유로와 같은 신뢰할 수 있는 정부 발행 통화와는 거리가 멀다고 평가됩니다. 미래에는 바뀔 수도 있지만, 아직은 아닙니다.
\end{tcolorbox}


\subsubsection*{1.3. 돈과 인플레이션 (Money and Inflation)}
\addcontentsline{toc}{subsubsection}{1.3. 돈과 인플레이션}

\begin{tcolorbox}[summarybox]
  \textbf{핵심 요약:} 통화량의 측정은 경제 성장에 영향을 미치므로 중요하며, 과거에는 통화량 증가가 인플레이션과 밀접한 관계를 보였으나, 21세기에는 그 관계가 불분명해졌습니다.
\end{tcolorbox}

이 세션에서는 돈의 개념을 넘어 경제에 유통되는 총 통화량을 측정하는 방법을 알아보고, 통화량과 인플레이션의 관계를 살펴봅니다. 통화량의 많고 적음은 경제 성장, 이자율, 인플레이션에 영향을 미칠 수 있으므로, 중앙은행과 정책 입안자들은 통화 공급과 지불 시스템의 혁신에 큰 관심을 기울입니다.

\textbf{돈과 유동성 (Money and Liquidity)}
돈은 우리가 지불에 사용하는 특별한 자산이며, 우리의 부(wealth)의 일부입니다. 우리는 거래를 위해 돈을 사용하지만, 우리의 부는 다양한 형태로 보유됩니다.

\begin{itemize}
    \item \textbf{자산 유동성 (Asset Liquidity):}
    \begin{itemize}
        \item \textbf{한 줄 핵심 요약:} 자산을 돈으로 전환하는 데 드는 비용과 시간을 나타냅니다.
        \item \textbf{직관적 예시/비유:} 현금은 이미 돈이므로 가장 유동성이 높습니다. 집이나 보석은 가치가 있지만, 적절한 가격에 팔려면 시간이 걸리므로 유동성이 낮습니다. 주식이나 저축 계좌는 그 중간 정도의 유동성을 가집니다.
        \item \textbf{기술적 정확한 설명:} 우리는 은행 계좌에 자금을 보유하고, 주식과 채권에 투자하며, 집을 소유할 수 있습니다. 이러한 다른 자산들도 지불에 사용될 수 있지만, 먼저 돈으로 전환(청산, liquidating)해야 하므로 더 많은 비용이 들 수 있습니다.
    \end{itemize}
\end{itemize}

\textbf{통화량 측정 (Measuring Money Supply)}
돈의 정의는 시간이 지남에 따라 변해왔으며, 이상적인 측정 방법은 없습니다. 현금만 돈으로 포함하는 것은 너무 엄격하고, 주식과 채권을 포함하는 것은 너무 느슨합니다. 경제학자들은 주로 두 가지 기본 정의인 M1과 M2를 사용합니다.

\begin{enumerate}
    \item \textbf{M1 (협의 통화):}
    \begin{itemize}
        \item \textbf{한 줄 핵심 요약:} 가장 엄격한 통화량 정의로, 대중이 보유한 현금과 은행의 요구불 예금(checking accounts)을 포함합니다.
        \item \textbf{직관적 예시/비유:} 내 지갑에 있는 지폐와 언제든지 인출하거나 직불카드로 사용할 수 있는 은행 계좌 잔액이 M1에 해당합니다.
        \item \textbf{기술적 정확한 설명:} M1은 즉시 지불 수단으로 사용될 수 있는 가장 유동적인 자산들로 구성됩니다.
    \end{itemize}

    \item \textbf{M2 (광의 통화):}
    \begin{itemize}
        \item \textbf{한 줄 핵심 요약:} M1에 일부 저축성 예금(savings accounts)을 추가한 더 넓은 통화량 측정치입니다.
        \item \textbf{직관적 예시/비유:} M1에 더해, 은행의 저축 계좌나 머니마켓 펀드(money market fund)와 같이 M1보다는 유동성이 약간 낮지만 쉽게 현금화할 수 있는 자금들이 M2에 포함됩니다. 저축 계좌는 이자를 지급하지만, 월별 인출 횟수 제한이 있을 수 있습니다.
        \item \textbf{기술적 정확한 설명:} 가계가 이자를 지급하는 저축 계좌로 자금을 옮기면서 M2가 경제 활동과 더 밀접하게 움직이는 경향이 있어, 경제학자들은 M2에 더 주목하는 경향이 있습니다.
    \end{itemize}
\end{enumerate}

\begin{tcolorbox}[examplebox]
  \textbf{통화량 구성 예시 (2020년 4월 미국 연방준비제도 데이터):}
  \begin{itemize}
      \item M1은 통화(currency)와 요구불 예금(checking accounts)으로 구성됩니다.
      \item M2는 M1에 저축성 예금(savings accounts)을 더한 것입니다.
  \end{itemize}
  만약 현금만 돈으로 측정한다면 전체 그림의 큰 부분을 놓치게 됩니다. 그래서 경제학자들은 M2를 선호합니다.
\end{tcolorbox}

\textbf{인플레이션 (Inflation)}
인플레이션은 시간이 지남에 따라 물가가 상승하는 비율을 나타냅니다. 인플레이션이 높으면 같은 상품을 사기 위해 더 많은 돈을 지불해야 합니다. 즉, 돈의 가치가 하락합니다.

\begin{itemize}
    \item \textbf{소비자 물가 지수 (Consumer Price Index, CPI):}
    \begin{itemize}
        \item \textbf{한 줄 핵심 요약:} 인플레이션율을 측정하는 가장 일반적인 지표로, 특정 기간 동안 동일한 재화 및 서비스 바구니(식료품, 주택, 휴대폰 등) 가격의 변화를 측정합니다.
        \item \textbf{직관적 예시/비유:} 2021년 1월 미국 CPI가 1.4\%였다면, 2020년 1월에 100달러였던 상품 바구니가 2021년 1월에는 101.40달러가 되었다는 의미입니다.
        \item \textbf{기술적 정확한 설명:} CPI의 성장은 돈의 구매력에 대한 정보를 제공합니다.
    \end{itemize}
\end{itemize}

\textbf{초인플레이션 (Hyperinflation)}
과거에는 매우 높은 인플레이션(초인플레이션) 사례들이 있었습니다.
\begin{itemize}
    \item 1923년 독일: 일일 인플레이션율 21\%, 3일마다 물가 두 배 상승.
    \item 1946년 헝가리, 2008년 짐바브웨, 2018년 베네수엘라에서도 유사한 사례 발생.
\end{itemize}
\begin{tcolorbox}[warningbox]
  \textbf{초인플레이션의 원인:} 일반적으로 정부가 재정 적자를 메우기 위해 중앙은행에 돈을 인쇄하도록 강요하여 통화 공급을 급격히 늘릴 때 발생합니다. 고정된 수의 상품을 쫓는 돈이 점점 많아지면 돈의 가치는 하락하고, 이는 가계와 기업에 심각한 결과를 초래합니다.
\end{tcolorbox}

\textbf{통화량 증가와 인플레이션의 관계}

\begin{enumerate}
    \item \textbf{1960년~1980년: 강한 양의 상관관계}
    \begin{itemize}
        \item \textbf{설명:} 1970년대와 80년대의 '대인플레이션(The Great Inflation)' 시기에는 인플레이션율이 10\%를 초과했습니다. 이 시기(1960~1980년)의 데이터 분석 결과, M2 통화량 증가율(2년 시차 적용)과 CPI 인플레이션율 사이에 강한 양의 상관관계가 나타났습니다. 즉, 높은 통화량 증가는 2년 후 높은 인플레이션과 연관되었고, 낮은 통화량 증가는 낮은 인플레이션과 연관되었습니다. 이는 우리의 직관과 일치합니다.
    \end{itemize}
    \item \textbf{2000년 이후: 상관관계 약화}
    \begin{itemize}
        \item \textbf{설명:} 2000년 이후 기간에 초점을 맞추면, M2 증가율과 CPI 인플레이션율 사이의 상관관계가 사라집니다. M2는 21세기 인플레이션을 예측하는 데 더 이상 유용하지 않습니다.
        \item \textbf{원인 가설:}
        \begin{enumerate}
            \item \textbf{낮은 인플레이션 환경:} 2000년 이후 인플레이션율이 역사적 기준으로 비정상적으로 낮은 약 2\% 수준을 유지했기 때문에, 통화량 증가와 인플레이션 간의 관계가 명확하게 나타나지 않을 수 있습니다.
            \item \textbf{돈의 측정 방식 변화:} 21세기에는 금융 시장의 거래 비용이 낮아지고, 신용카드가 널리 보급되었으며, 새로운 디지털 결제 기술이 매일 등장하면서 돈을 사용하고 지불하는 방식이 크게 변했습니다. 따라서 M2가 더 이상 돈을 정확하게 측정하는 지표가 아닐 수 있습니다. 새로운 통화량 측정 방식이 필요할 수도 있습니다.
        \end{enumerate}
    \end{itemize}
\end{enumerate}

\begin{tcolorbox}[summarybox]
  \textbf{결론:} 돈의 개념에서 통화량 측정으로 나아가, 통화량 증가와 인플레이션의 관계를 살펴보았습니다. 1970년대와 80년대의 대인플레이션 시기에는 통화량 증가가 높았지만, 21세기에는 이 관계가 불분명해졌습니다. 이는 미래에 정부와 중앙은행이 인플레이션을 통제하는 과거의 접근 방식이 효과적일지에 대한 의문을 제기합니다.
\end{tcolorbox}


\subsection{2. 금융 시스템 (The Financial System)}
\addcontentsline{toc}{subsection}{2. 금융 시스템}
이 섹션에서는 금융 시스템이 경제에서 어떤 역할을 하는지, 그리고 금융 시장과 금융 기관이 어떻게 자금을 중개하는지 상세히 살펴봅니다.

\subsubsection*{2.1. 금융 시스템의 역할 (What Does the Financial System Do?)}
\addcontentsline{toc}{subsubsection}{2.1. 금융 시스템의 역할}

\begin{tcolorbox}[summarybox]
  \textbf{핵심 요약:} 금융 시스템은 자금 여유가 있는 '저축자'와 자금 수요가 있는 '차입자'를 연결하여 자금을 가장 효율적인 곳으로 흐르게 하고, 그 대가로 '금융 상품'을 교환함으로써 경제 번영과 성장을 촉진합니다.
\end{tcolorbox}

금융 시스템의 출발점은 자금 수요가 있는 경제 주체입니다. 이는 수익성 있는 투자 기회를 가진 기업일 수도 있고, 주택 구매 자금이 필요한 가계일 수도 있습니다.

\textbf{저축자와 차입자 (Savers and Borrowers)}

\begin{itemize}
    \item \textbf{차입자 (Borrowers):}
    \begin{itemize}
        \item \textbf{한 줄 핵심 요약:} 자금이 필요한 개인, 기업, 정부를 포함합니다.
        \item \textbf{직관적 예시/비유:} 주택을 사거나 물건을 구매하기 위한 가계, 단기 운전자금이나 장기 사업 확장을 위한 기업, 공공재를 위한 정부, 트랙터를 사려는 농부 등이 차입자가 될 수 있습니다.
        \item \textbf{기술적 정확한 설명:} 기업은 자체 현금 보유액이나 자산 매각으로 투자 자금을 조달할 수 있지만, 내부 자금이 충분하지 않을 경우 외부에서 자금을 조달해야 합니다. 이때 저축자로부터 자금을 빌리는 차입자가 됩니다.
    \end{itemize}

    \item \textbf{저축자 (Savers):}
    \begin{itemize}
        \item \textbf{한 줄 핵심 요약:} 현재 지출보다 더 많은 소득을 가진 가계가 대부분입니다.
        \item \textbf{직관적 예시/비유:} 은퇴를 위해 연기금에 돈을 넣거나, 은행에 저축하거나, 비상금을 위해 보험에 가입하는 가계가 저축자입니다.
        \item \textbf{기술적 정확한 설명:} 저축은 미래를 위한 준비이며, 이는 연기금이나 뮤추얼 펀드를 통해 이루어지거나, 은행 예금으로 현금 관리에 활용되거나, 보험 가입을 통해 위험을 분산하는 형태로 나타납니다.
    \end{itemize}
\end{itemize}

\textbf{금융 시스템의 역할}
금융 시스템은 저축자와 차입자를 연결하여 저축된 자금을 차입자의 자금 수요에 맞게 전달합니다. 이 과정의 궁극적인 목표는 경제적 번영입니다. 금융 시스템은 자원을 가장 효율적인 용도로 배분하여 경제 성장을 이끌어냅니다. 저축자들은 자금을 제공하는 대가로 '금융 상품(financial instruments)'을 받습니다.

\textbf{금융 상품 (Financial Instruments)}
금융 상품은 기본적으로 '상환에 대한 공식적인 약속'입니다. 이는 법적 구속력이 있는 의무이며, 문제가 발생할 경우 법원과 정부가 강제할 수 있습니다. 특정 상황에서 참여자들 간에 자금이 이전되도록 요구합니다.

\begin{itemize}
    \item \textbf{저축자를 위한 금융 상품:}
    \begin{itemize}
        \item \textbf{은행 예금 계좌 (Bank Deposit Account):}
        \begin{itemize}
            \item \textbf{설명:} 저축자는 은행에 현금을 예치하고, 이 자금은 요구불(on-demand)로 인출 가능합니다. 안전하고 편리하며, 수표, 직불카드, 모바일 뱅킹을 통해 결제에 사용될 수 있습니다.
        \end{itemize}
        \item \textbf{뮤추얼 펀드 주식 (Mutual Fund Share):}
        \begin{itemize}
            \item \textbf{설명:} 돈을 투자 포트폴리오의 지분과 교환하여 부를 저장하고 성장시키는 데 도움을 줍니다.
        \end{itemize}
        \item \textbf{보험 (Insurance):}
        \begin{itemize}
            \item \textbf{설명:} 사고와 같은 불리한 사건에 대비하여 현재 지출을 포기하고 보험료를 지불합니다. 사고 발생 시 보험금(payout)을 받아 손실을 보전합니다. 보험은 저축자가 위험을 이전할 수 있게 해주는 금융 상품입니다.
        \end{itemize}
    \end{itemize}

    \item \textbf{차입자를 위한 금융 상품:}
    \begin{itemize}
        \item \textbf{주식 (Equity):}
        \begin{itemize}
            \item \textbf{설명:} 기업 소유주가 주식을 팔아 자금을 조달합니다. 주식 구매자는 기업의 지분을 소유하게 되며, 이는 미래 기업 이익에 대한 권리를 부여합니다.
        \end{itemize}
        \item \textbf{채무 (Debt):}
        \begin{itemize}
            \item \textbf{설명:} 채권(bonds), 은행 대출(bank loans), 모기지(mortgages) 등이 있습니다. 차입자는 오늘 일시불로 자금을 받고, 미래에 정해진 날짜에 정해진 금액을 상환하기로 약속합니다. 이 때문에 '고정 수입(fixed income)' 상품이라고도 불립니다.
            \item \textbf{예시:} 정부 채권은 과거에 쿠폰이 붙어 있어 은행에서 현금으로 교환할 수 있었습니다. 주택 담보 대출(residential mortgage)의 경우, 차입자는 토지와 주택을 담보(collateral)로 제공합니다.
        \end{itemize}
    \end{itemize}
\end{itemize}

\begin{tcolorbox}[warningbox]
  \textbf{왜 이렇게 많은 금융 상품이 존재할까?}
  다양한 금융 상품은 저축자와 차입자의 다양한 요구를 충족시키기 위해 맞춤 제작됩니다. 저축자는 편리한 부의 저장 및 성장 방법을 원하고, 위험을 분산하기 위해 보험을 원할 수 있습니다. 차입자는 장기 자금을 원하며, 부채 발행이나 주식 매각을 통해 자금을 조달할 수 있습니다.
\end{tcolorbox}

\begin{tcolorbox}[summarybox]
  \textbf{결론:} 금융 시스템은 자금 여유가 있는 저축자와 투자 수요가 있는 차입자 사이에 존재합니다. 이 시스템은 저축자로부터 차입자에게 자금을 전달하고, 그 대가로 저축자와 차입자의 요구를 충족시키도록 설계된 금융 상품을 교환합니다.
\end{tcolorbox}


\subsubsection*{2.2. 금융 시장 vs. 금융 기관 (Financial Markets Versus Financial Institutions)}
\addcontentsline{toc}{subsubsection}{2.2. 금융 시장 vs. 금융 기관}

\begin{tcolorbox}[summarybox]
  \textbf{핵심 요약:} 금융 시스템은 자금을 저축자에서 차입자로 전달하는 두 가지 주요 경로를 가집니다. 하나는 '금융 시장'을 통한 직접적인 경로이고, 다른 하나는 '금융 기관'을 통한 간접적인 경로입니다.
\end{tcolorbox}

금융 시스템은 저축자로부터 차입자에게 자금을 전달하고, 그 대가로 금융 상품을 교환합니다. 이 과정은 크게 두 가지 방식으로 이루어집니다.

\textbf{1. 금융 시장을 통한 직접 자금 흐름 (Direct Flow via Financial Markets)}
저축자가 금융 시장에 직접 참여하는 방식입니다.

\begin{itemize}
    \item \textbf{한 줄 핵심 요약:} 저축자가 주식이나 채권과 같은 금융 상품을 다른 투자자로부터 직접 구매하는 방식입니다.
    \item \textbf{직관적 예시/비유:} 내가 Apple 주식을 사고 싶다면, 보통 Apple 회사로부터 직접 사는 것이 아니라 뉴욕 증권 거래소(New York Stock Exchange)에서 다른 투자자로부터 주식을 구매합니다.
    \item \textbf{기술적 정확한 설명:} 이러한 '직접 교환(direct exchange)'은 보통 브로커(broker)나 딜러(dealer)와 같은 금융 회사의 도움을 받아 이루어집니다.
    \begin{itemize}
        \item \textbf{브로커 (Broker):} 수수료를 받고 거래 파트너를 찾아줍니다.
        \item \textbf{딜러 (Dealer):} 수수료를 받고 자신의 재고(inventory)에 있는 금융 상품을 직접 판매하여 거래 파트너 역할을 합니다.
    \end{itemize}
    오늘날 대부분의 이러한 상호작용은 E*Trade와 같은 온라인 거래 플랫폼을 통해 전자적으로 이루어집니다.
\end{itemize}

\textbf{금융 시장의 종류}
금융 시장은 거래되는 금융 상품의 종류에 따라 분류됩니다.

\begin{itemize}
    \item \textbf{주식 시장 (Stock or Equity Markets):} 기업의 주식이 거래되는 곳입니다.
    \item \textbf{채무 시장 (Debt Markets):} 대출, 모기지, 회사채 등이 거래되는 곳입니다.
    \item \textbf{단기 금융 시장 (Money Markets):} 만기가 1년 미만인 단기 채권(예: 기업어음, commercial paper)이 거래되는 곳입니다.
\end{itemize}
\begin{tcolorbox}[examplebox]
  \textbf{Apple의 자금 조달 예시:} Apple이 운영 자금을 조달할 때, 주식, 장기 채권, 또는 단기 채권(기업어음)을 발행할 수 있습니다. 발행 후, 주식은 주식 시장에서, 장기 채권은 채권 시장에서, 기업어음은 단기 금융 시장에서 거래됩니다.
\end{tcolorbox}

\textbf{잘 작동하는 금융 시장의 특징}
잘 작동하는 금융 시장은 현대 경제에서 핵심적인 역할을 합니다. 크고 활발한 시장을 가진 국가들은 다음 세 가지 특징을 보입니다.

\begin{enumerate}
    \item \textbf{낮은 거래 비용 (Low Transaction Costs):} 차입자가 증권을 발행하는 비용과 투자자들이 서로 거래하는 비용이 저렴해야 합니다.
    \item \textbf{정보의 효율적 흐름 (Efficient Information Flow):} 시장 참여자들이 차입자에 대한 정보를 수집, 처리, 전달할 수 있어야 합니다. 이는 시장 가격이 잠재적 위험과 수익을 정확하게 반영하도록 보장하며, 자원이 최적의 용도로 흐르도록 합니다.
    \item \textbf{투자자 보호 (Investor Protection):} 투자자들이 투자 손실로부터 보호받는다는 의미가 아니라, 차입자가 투자자를 속이지 않을 것이라는 신뢰를 가져야 한다는 의미입니다. 금융 상품은 법적 구속력이 있는 상환 약속이므로, 현대 금융 시장은 규제 기관과 법원이 계약을 강제할 것을 요구합니다. 투자자들이 남용으로부터 보호받을 때 기꺼이 투자할 것입니다.
\end{enumerate}

\textbf{2. 금융 기관을 통한 간접 자금 흐름 (Indirect Flow via Financial Institutions)}
금융 기관은 저축자로부터 자금을 조달하고, 그 자금을 차입자에게 투자하는 방식으로 자금을 중개합니다.

\begin{itemize}
    \item \textbf{한 줄 핵심 요약:} 금융 기관은 저축자로부터 자금을 모아 단기 부채(예: 예금)를 제공하고, 이 자금을 차입자에게 장기 투자(예: 대출)하는 방식으로 자금을 중개합니다.
    \item \textbf{기술적 정확한 설명:} 금융 기관은 먼저 저축자로부터 자금을 조달하고, 그 대가로 보통 단기 부채를 제공합니다. 그런 다음, 이 자금을 차입자가 발행한 금융 상품에 투자합니다. 이 과정에서 저축자로부터 차입자로 자금이 전달되지만, 중간에 금융 기관이라는 추가 단계가 있습니다. 브로커나 딜러와 달리, 금융 기관은 대차대조표(balance sheets)를 가지고 있으며 종종 장기 투자를 합니다.
\end{itemize}

\textbf{은행의 예시를 통한 금융 기관의 역할}
가장 단순화된 은행의 예를 통해 이 과정을 이해해 봅시다.

\begin{enumerate}
    \item \textbf{저축자의 예금:} 내가 저축자로서 1,000달러를 은행에 예치합니다. 은행은 나에게 '예금 계약(deposit contract)'을 제공하며, 이는 은행의 부채(liability)가 됩니다 (은행은 내가 원할 때 1,000달러를 돌려줄 의무가 있습니다). 은행은 이러한 예금을 여러 번 반복하여 총 100만 달러의 저축을 모읍니다.
    \item \textbf{은행의 대출:} 은행은 신용도 높은 사업체를 찾아 100만 달러를 장기 사업 대출로 빌려줍니다. 이 사업 대출은 은행의 자산(asset)이 됩니다 (사업체는 은행에 100만 달러를 갚을 의무가 있습니다).
    \item \textbf{자금 중개 및 대차대조표 형성:} 결과적으로 은행은 저축자로부터 차입자에게 자금을 전달했습니다. 이 과정에서 은행은 100만 달러의 대출(자산)과 100만 달러의 예금(부채)으로 구성된 대차대조표를 생성합니다.
\end{enumerate}

\textbf{왜 금융 기관이 필요한가? (금융 기관의 가치 창출)}
금융 시장과 은행이 공존하는 것을 보면, 둘 다 존재할 만한 좋은 이유가 있을 것입니다. 경제학자들은 금융 기관이 다음 세 가지 문제를 해결하는 데 도움을 준다고 봅니다.

\begin{enumerate}
    \item \textbf{만기 불일치 (Maturity Mismatch):}
    \begin{itemize}
        \item \textbf{한 줄 핵심 요약:} 차입자는 장기 자금을 원하고, 저축자는 단기 유동성을 원하는데, 은행이 이 둘을 연결해줍니다.
        \item \textbf{직관적 예시/비유:} 주택 구매자는 15년 또는 30년 만기의 모기지를 원하고, 기업은 장기 회수 기간을 가진 프로젝트를 위한 자금을 원합니다. 반면 저축자들은 자신의 자금에 쉽게 접근하고 싶어 하며, 모든 돈이 장기간 묶이는 것을 원치 않습니다.
        \item \textbf{기술적 정확한 설명:} 은행은 요구불 예금과 같은 단기 부채를 발행하고, 이를 장기 대출로 전환함으로써 이러한 만기 불일치 문제를 해결합니다.
    \end{itemize}

    \item \textbf{가격 위험 (Price Risk):}
    \begin{itemize}
        \item \textbf{한 줄 핵심 요약:} 주식이나 채권 가격의 변동성으로 인한 위험을 은행 예금은 줄여줍니다.
        \item \textbf{직관적 예시/비유:} 주식과 채권의 가격은 매일 변동하며, 거래 수수료도 변동할 수 있습니다. 일부 저축자들은 이러한 가격 위험을 매우 싫어하여 현금을 보유하는 것을 선호할 수 있습니다.
        \item \textbf{기술적 정확한 설명:} 은행의 예금 계약은 매우 편리하고 안전합니다. 정부가 제공하는 예금 보험이 있는 경우, 오늘 은행에 1달러를 예치하면 내일 최소 1달러를 인출할 수 있습니다. 은행 계좌와 관련된 수수료는 있을 수 있지만, 가격 위험은 없습니다.
    \end{itemize}

    \item \textbf{정보 비대칭 (Information Asymmetry):}
    \begin{itemize}
        \item \textbf{한 줄 핵심 요약:} 개별 저축자가 투자 위험과 수익을 평가하기 어려운 문제를 금융 기관이 전문적으로 해결해줍니다.
        \item \textbf{직관적 예시/비유:} Zano라는 아기 드론 프로젝트는 2015년 크라우드펀딩으로 200만 달러 이상을 모았지만, 결국 실패하고 투자자들은 돈을 잃었습니다. 개별 투자자들은 성공적인 사업을 선별하고 손실을 최소화하기 어렵습니다.
        \item \textbf{기술적 정확한 설명:} 정보 비대칭은 저축자가 투자에 대한 위험과 기대 수익을 평가하기 어려울 때 발생합니다. 은행은 차입자를 심사하고 평가하는 전문가이며, 대출 후에도 차입자를 모니터링하고 실적이 좋지 않은 차입자에게는 대출을 중단함으로써 이러한 정보 위험을 해결합니다. 시장 기반 솔루션(예: 주식 분석가, 신용 평가 기관)도 있지만, 정보 위험이 너무 크면 일부 가계는 투자를 꺼릴 수 있습니다.
    \end{itemize}
\end{enumerate}

\begin{tcolorbox}[summarybox]
  \textbf{결론:} 금융 시스템은 저축자로부터 차입자에게 자금을 이동시키며, 이는 두 가지 경로를 통해 이루어집니다. 첫째, 금융 시장을 통한 직접적인 경로로, 브로커의 도움을 받아 투자자들이 서로 증권을 거래합니다. 둘째, 금융 기관을 통한 간접적인 경로로, 금융 기관은 만기 불일치, 가격 위험, 정보 비대칭 문제를 해결하여 가치를 더합니다.
\end{tcolorbox}


\subsubsection*{2.3. 은행 vs. 비은행 금융 기관 (Banks Versus Non-Banks)}
\addcontentsline{toc}{subsubsection}{2.3. 은행 vs. 비은행 금융 기관}

\begin{tcolorbox}[summarybox]
  \textbf{핵심 요약:} 현대 금융 산업은 금융 시장과 금융 기관으로 구성되며, 금융 기관은 다시 '은행'과 '비은행 금융 기관'으로 나뉩니다. 은행은 예금을 받고 대출을 하는 전통적인 역할을 하며, 비은행 기관은 예금을 받지 않으면서 다양한 금융 서비스를 제공합니다.
\end{tcolorbox}

이 세션에서는 미국 금융 산업의 조직 방식과 가장 중요한 은행 및 비은행 금융 기관들을 소개하고, 경제에서 이들의 역할을 설명합니다.

\textbf{금융 산업의 구성}
현대 금융 산업은 크게 금융 시장과 금융 기관으로 구성됩니다.
\begin{itemize}
    \item \textbf{금융 시장:} 저축자와 차입자를 직접 연결하는 중개 서비스를 제공하는 회사들이 있습니다.
    \item \textbf{금융 기관:} 저축자로부터 자금을 조달하여 주식, 대출, 채권, 부동산 등 다양한 금융 상품에 투자하는 회사들입니다. 이들은 대차대조표를 가집니다.
\end{itemize}
우리의 목표는 가장 중요한 금융 기관들이 누구이며, 어떤 고유한 경제적 기능을 수행하고, 어떻게 자금을 조달하며, 어디에 자금을 배치하는지(즉, 대차대조표가 어떻게 생겼는지) 이해하는 것입니다. 또한 시장 점유율 측면에서 어떤 기관이 가장 중요한지도 살펴봅니다.

\textbf{은행과 비은행 금융 기관의 구분}
금융 기관은 크게 '은행(banks)'과 '비은행 금융 기관(non-banks)'으로 나뉩니다.

\textbf{1. 은행 (Banks)}
\begin{itemize}
    \item \textbf{한 줄 핵심 요약:} 우리가 일상적으로 접하는 금융 기관으로, 저축자로부터 예금을 받고 차입자에게 대출을 해주는 전통적인 역할을 합니다.
    \item \textbf{종류:} 상업 은행(commercial banks), 저축 기관(savings institutions), 신용 협동조합(credit unions) 등이 있습니다. 다른 나라에는 다른 유형의 은행도 있지만, 모두 예금을 모으고 대출을 해주는 동일한 기능을 수행합니다.
    \item \textbf{자금 조달 (Funding):} 주로 개인 고객(저축자)으로부터 예금 계좌 형태로 자금을 조달합니다. 일부 대형 은행은 은행 간 시장(interbank market)이나 환매 조건부 채권 시장(repo market)에서 차입하기도 합니다.
    \item \textbf{자금 운용 (Deployment):} 주로 대출에 투자합니다. 사실 은행은 경제에서 대출의 주요 원천입니다. 은행은 어떤 가계가 모기지를 받을지, 대출 규모는 얼마인지, 이자율은 얼마인지를 결정합니다.
    \item \textbf{은행의 특별한 점 (가치 창출):}
    \begin{itemize}
        \item \textbf{예금자에게:} 안전하고 편리하며 지불 시스템에 접근할 수 있는 은행 계좌를 제공합니다.
        \item \textbf{은행 자체:} 소액의 저축을 대량으로 모아 신용도 높은 차입자에게 대출을 해줍니다. 이러한 '대출 원천화(origination)'는 은행의 핵심 기능입니다.
        \item \textbf{고객을 위한 시장 거래:} 고객을 대신하여 금융 시장에서 거래를 수행하여 위험을 분산하고 부를 관리할 수 있도록 돕습니다.
    \end{itemize}
\end{itemize}

\textbf{은행과 증권 시장의 상대적 중요성 (전 세계 데이터)}
세계은행(World Bank)의 최근 데이터에 따르면, 국내 증권 시장(주식 시장 시가총액 + 발행 채권 총액)의 규모와 총 대출 잔액을 비교하여 은행의 상대적 중요성을 파늠할 수 있습니다.

\begin{itemize}
    \item \textbf{주요 시사점 (지난 25년간 미국, 유럽, 아시아 두 경제권 데이터):}
    \begin{enumerate}
        \item \textbf{모두 중요:} 모든 경제권에서 증권 시장과 대출 시장 모두 중요합니다. 이는 은행이 여전히 매우 중요하다는 것을 의미합니다.
        \item \textbf{지역별 차이:} 전 세계적으로 많은 차이가 있습니다. 미국은 금융이 은행에 덜 의존하는 반면, 중국은 그 반대입니다.
    \end{enumerate}
\end{itemize}

\textbf{2. 비은행 금융 기관 (Non-Banks)}
\begin{itemize}
    \item \textbf{한 줄 핵심 요약:} 정부 보증 예금을 받지 않는 금융 기관입니다.
    \item \textbf{설명:} 비은행 금융 기관은 일반 은행과 유사한 많은 경제적 기능을 수행합니다. 부를 저장하는 안전하고 편리한 방법을 제공하고, 소액 저축을 모아 차입자에게 전달하며, 투자와 관련된 정보 비용을 최소화하고, 고객이 위험을 줄일 수 있도록 돕습니다.
    \item \textbf{종류:} 보험 회사(insurance companies), 연기금(pension funds), 증권 회사(securities firms), 할부금융 회사(finance companies), 정부 후원 기업(government-sponsored enterprises, GSEs) 등이 있습니다. 이들은 매우 다양하며, 자금 조달 및 투자 프로필, 그리고 존재 이유가 크게 다릅니다.
\end{itemize}

\begin{tcolorbox}[warningbox]
  \textbf{비은행 금융 기관의 광범위함:} P2P 대출업체부터 전당포, 사모펀드, 그리고 최근의 핀테크 기업까지 매우 광범위한 비은행 금융 기관 생태계가 존재합니다.
\end{tcolorbox}


\subsubsection*{2.4. 비은행 금융 기관 개요 (Overview of Non-Banks)}
\addcontentsline{toc}{subsubsection}{2.4. 비은행 금융 기관 개요}

이 세션에서는 미국에서 자산 관리 규모 기준으로 가장 중요한 비은행 금융 기관들에 초점을 맞춥니다.

\textbf{1. 보험 회사 (Insurance Companies)}
\begin{itemize}
    \item \textbf{경제적 기능:} 위험 완화(risk mitigation)입니다. 개인의 관점에서는 '비상금을 위한 저축'의 한 형태입니다.
    \item \textbf{작동 방식:} 나는 어떤 불리한 사건의 위험을 보험 회사에 이전합니다. 보험료(premium)를 지불하고, 나쁜 상황(예: 자동차 사고)이 발생하면 보험금을 청구하여 보험 회사가 손실을 보전해줍니다. 보험사 입장에서는 더 많은 보험 계약을 체결할수록 보험금 지급이 예측 가능해집니다.
    \item \textbf{주요 유형:}
    \begin{itemize}
        \item \textbf{생명 보험 (Life Insurance):} 장애나 사망으로 인한 소득 손실로부터 보호합니다. (예: 19세기 나폴레옹 전쟁 참전 군인 가족을 위한 Scottish Widows)
        \item \textbf{손해 보험 (Property and Casualty Insurance):} 홍수, 화재와 같은 자연재해나 사고로 인한 가계 및 기업의 손실로부터 보호합니다. (예: 자동차 사고)
    \end{itemize}
    \item \textbf{자금 조달 및 운용:} 보험료는 모아져 투자됩니다. 보험 회사의 자금원은 모아진 보험료이며, 이 자금은 보험금 지급에 사용되지만, 주로 주식, 부동산, 채권과 같은 장기 투자에도 사용됩니다. 보험 회사는 미국 회사채의 가장 큰 투자자 중 하나입니다.
\end{itemize}

\textbf{2. 연기금 (Pension Funds)}
\begin{itemize}
    \item \textbf{경제적 기능:} 미래를 위해 부를 저장하고 성장시키는 데 사용됩니다.
    \item \textbf{작동 방식:} 가계는 젊은 나이부터 은퇴를 위해 저축할 수 있으며, 종종 강제적으로 이루어지기도 합니다 (예: 급여의 일부가 자동 적립). 은퇴 저축 계좌는 일반적으로 소득세 유예와 같은 세금 혜택을 제공하며, 고용주가 기여금을 내기도 합니다. 은퇴 시에는 일시불 또는 종신 연금 형태로 지급받습니다.
    \item \textbf{자금 조달 및 운용:} 연기금은 직원 및 고용주의 기여금을 모아 자금을 조달합니다. 이 자금은 주로 채권, 주식, 부동산과 같은 장기 투자에 사용되며, 동시에 은퇴자에게 지급할 보험금을 관리합니다.
    \item \textbf{보험 회사와의 유사점:} 연기금과 생명 보험 회사는 매우 유사한 사업을 합니다. 연기금은 은퇴할 경우 지급하고, 생명 보험은 사망할 경우 지급합니다. 이러한 유사성 때문에 이들 사업은 동일한 투자를 하는 경향이 있으며, 대형 생명 보험 회사가 연기금을 관리하기도 합니다.
\end{itemize}

\textbf{3. 증권 회사 (Securities Firms)}
주요 유형은 투자 은행과 투자 펀드입니다.

\begin{itemize}
    \item \textbf{투자 은행 (Investment Banks):}
    \begin{itemize}
        \item \textbf{한 줄 핵심 요약:} 상업 은행이 가계의 대출을 결정하는 것과 달리, 투자 은행은 기업이 자본 시장에서 증권을 발행할 수 있는지와 그 조건을 결정합니다.
        \item \textbf{핵심 기능: 인수 (Underwriting):} 기업이 증권을 발행하고자 할 때, 투자 은행은 특정 가격으로 전체 발행 물량을 매입하기로 약속합니다. 그런 다음, 이 증권을 외부 투자자에게 판매하여 수익을 얻습니다.
        \item \textbf{경제적 기능:}
        \begin{enumerate}
            \item \textbf{정보 수집 및 처리:} 잠재적 차입자에 대한 정보를 수집하고 처리하는 사업을 합니다. 자본 시장의 '문지기(gatekeepers)' 역할을 합니다.
            \item \textbf{저축자와 차입자 연결:} 인수 과정을 통해 저축자와 차입자를 연결합니다.
        \end{enumerate}
    \end{itemize}

    \item \textbf{투자 펀드 (Investment Funds):}
    \begin{itemize}
        \item \textbf{한 줄 핵심 요약:} 저축을 모아 대규모의 분산된 자산 포트폴리오에 투자하여 투자자들이 부를 저장하고 성장시킬 수 있도록 돕습니다.
        \item \textbf{종류:}
        \begin{itemize}
            \item \textbf{뮤추얼 펀드 (Mutual Funds) 및 상장지수펀드 (Exchange Traded Funds, ETFs):} 소액 투자자도 최소 투자 요건과 낮은 거래 수수료로 접근할 수 있습니다. 다양한 투자 프로필을 가집니다 (예: 머니마켓 펀드는 매우 안전하고 단기 채권에만 투자, 주식 뮤추얼 펀드는 더 위험). Vanguard와 같은 대형 뮤추얼 펀드 회사는 S\&P 500을 추종하는 펀드를 3,000달러의 최소 투자금으로 제공하여, 개인이 직접 이러한 투자를 복제하기 어려운 '풀링(pooling)'의 이점을 제공합니다.
            \item \textbf{헤지 펀드 (Hedge Funds):} 부유하고 정교한 투자자만을 대상으로 합니다. 100만 달러 이상의 순자산을 가진 99명 이하의 투자자 또는 500만 달러 이상의 순자산을 가진 499명 이하의 투자자를 대상으로 합니다. 일반 소액 투자자를 피함으로써 특정 규제를 우회할 수 있으며, 증권 및 파생상품 시장에서 고위험 고수익 투자 전략을 추구할 수 있습니다.
        \end{itemize}
    \end{itemize}
\end{itemize}

\textbf{4. 할부금융 회사 (Finance Companies)}
\begin{itemize}
    \item \textbf{한 줄 핵심 요약:} 은행과 유사하지만 특정 대출 시장에 특화된 금융 기관입니다.
    \item \textbf{주요 활동:}
    \begin{itemize}
        \item \textbf{소비자 신용 (Consumer Credit):} 가계 구매를 위한 신용을 제공합니다 (예: 주택 모기지, 가전제품 할부 대출, 자동차 대출). Ford Motor Credit Company는 자동차 판매 금융을 담당합니다.
        \item \textbf{사업 대출 (Business Lending):} 항공기 구매 금융부터 농업 장비 리스 제공까지 다양합니다.
    \end{itemize}
    \item \textbf{은행과의 유사점:} 대출 및 리스 제공이 주 활동이며, 차입자를 심사하고 모니터링하여 정보 비용을 최소화하는 등 은행과 유사하게 작동합니다.
    \item \textbf{은행과의 차이점 (자금 조달):} 예금으로 자금을 조달하지 않습니다. 대신 채권 발행이나 은행 대출을 통해 자금을 조달합니다.
\end{itemize}

\textbf{5. 정부 후원 기업 (Government Sponsored Enterprises, GSEs)}
\begin{itemize}
    \item \textbf{한 줄 핵심 요약:} 미국 모기지 시장의 핵심 주체로, 정부가 설립했으나 민간 형태로 운영되며 주택 소유를 지원합니다.
    \item \textbf{주요 기관:}
    \begin{itemize}
        \item \textbf{Fannie Mae (Federal National Mortgage Association):} 1930년대 설립.
        \item \textbf{Freddie Mac (Federal Home Loan Mortgage Corporation):} 1968년 설립.
        \item \textbf{Ginnie Mae (Government National Mortgage Corporation):} 1968년 설립. (정부가 전액 소유)
    \end{itemize}
    \item \textbf{역할:} 저소득층 및 소외 계층을 포함한 미국 내 주택 소유를 지원합니다. 이는 모기지 이자율을 보조함으로써 이루어집니다.
    \item \textbf{보조 방식:}
    \begin{itemize}
        \item \textbf{직접 대출:} 시장 금리보다 낮은 금리로 직접 대출을 제공합니다.
        \item \textbf{상환 보증:} 다른 금융 기관이 발행한 특정 부동산 금융 상품에 대한 상환을 보증합니다. 이러한 보증은 보험과 유사하게 작용하여 보조금 대출의 위험을 줄입니다.
    \end{itemize}
    \item \textbf{자금 조달 및 운용:} GSE는 채권 발행 또는 다른 금융 기관에 대한 대출 보증 판매를 통해 자금을 조달합니다. 이 자금은 주로 모기지 대출을 제공하고, 대출 보증과 관련된 보험금을 지급하는 데 사용됩니다.
\end{itemize}

\textbf{미국 금융 기관의 상대적 중요성 (1970년 vs. 2018년)}
미국 연방준비제도(Federal Reserve)의 자금 흐름(flow of funds) 데이터를 통해 은행과 주요 비은행 금융 기관의 자산 보유 추정치를 비교할 수 있습니다.

\begin{tcolorbox}[summarybox]
  \textbf{주요 시사점:}
  \begin{enumerate}
      \item \textbf{은행의 비중 감소:} 1970년에는 은행이 전체 자산의 절반 이상을 차지했지만, 2018년에는 25\% 미만으로 감소했습니다. 보험 회사도 비중이 감소했습니다.
      \item \textbf{뮤추얼 펀드의 부상:} 뮤추얼 펀드의 자산 비중은 5\% 미만에서 30\% 이상으로 크게 증가했습니다. 이는 Vanguard와 같은 투자 펀드가 제공하는 저비용 인덱스 투자 상품에 대한 접근성이 증가했음을 반영합니다.
  \end{enumerate}
\end{tcolorbox}

\begin{tcolorbox}[summarybox]
  \textbf{결론:} 금융 기관은 저축자로부터 차입자에게 자금을 전달하며, 이 과정에는 많은 은행 및 비은행 금융 기관이 참여합니다. 각 기관은 고유한 경제적 기능과 대차대조표를 가집니다. 전 세계적으로 은행은 여전히 중요한 역할을 하지만, 미국에서는 최근 몇 년간 비은행 금융 기관, 특히 뮤추얼 펀드의 중요성이 커졌습니다. 이러한 자금 재배분이 경제에 미치는 잠재적 영향(긍정적 또는 부정적)은 추후 논의될 주제입니다.
\end{tcolorbox}


\subsection{3. 금융과 실물 경제 (Finance and the Real Economy)}
\addcontentsline{toc}{subsection}{3. 금융과 실물 경제}
이 섹션에서는 금융 시스템의 발전이 실물 경제 활동에 어떤 영향을 미치는지 살펴봅니다.

\subsubsection*{3.1. 금융 발전과 경제 활동 (Financial Development and Economic Activity)}
\addcontentsline{toc}{subsubsection}{3.1. 금융 발전과 경제 활동}

\begin{tcolorbox}[summarybox]
  \textbf{핵심 요약:} 금융 발전은 돈과 금융 시스템의 효율성과 품질 향상을 의미하며, 이는 경제 성장과 위험 분산에 큰 이점을 제공하지만, 동시에 금융 불안정성과 경제 위기를 초래할 수도 있습니다.
\end{tcolorbox}

\textbf{금융 발전의 의미}
금융 발전(Financial Development)은 돈과 금융 시스템이 더 빠르고 더 효율적으로 작동하게 되는 것을 의미합니다. 이는 금융 시스템이 자원을 가장 생산적인 투자 기회로 더 잘 배분할 수 있게 됨을 뜻합니다.

\textbf{금융 발전의 잠재적 이점:}
\begin{itemize}
    \item \textbf{경제 성장 (Economic Growth):} 효율적인 금융 시스템은 자본이 필요한 기업과 개인에게 자금을 효과적으로 공급하여 생산성을 높이고 경제 전반의 성장을 촉진합니다.
    \item \textbf{위험 분산 (Risk-Sharing):} 보험, 뮤추얼 펀드와 같은 금융 상품을 통해 개인과 기업이 위험을 분산하고 관리할 수 있게 합니다. 이는 불확실성을 줄이고 투자를 장려합니다.
\end{itemize}

\textbf{금융 발전의 잠재적 비용 및 위험:}
\begin{itemize}
    \item \textbf{금융 불안정성 (Financial Instability):} 금융 시스템이 너무 빠르게 발전하거나 복잡해지면, 규제가 미치지 못하는 영역이 생기거나 새로운 유형의 위험이 발생할 수 있습니다. 이는 시스템 전체의 불안정성을 초래할 수 있습니다.
    \item \textbf{경제 위기 (Economic Crashes):} 금융 불안정성은 금융 위기로 이어질 수 있으며, 이는 실물 경제에 심각한 타격을 주어 경기 침체나 경제 위기를 유발할 수 있습니다. (예: 2008년 글로벌 금융 위기)
\end{itemize}

\begin{tcolorbox}[summarybox]
  \textbf{모듈 1의 최종 목표:} 이 모듈을 통해 돈과 금융이 무엇인지, 그리고 그것이 가져올 수 있는 잠재적인 이점과 비용에 대한 큰 그림을 이해하는 것입니다. 금융 발전은 양날의 검과 같아서, 그 혜택을 누리면서도 위험을 관리하는 것이 중요합니다.
\end{tcolorbox}


\section{빠르게 훑어보기 (1페이지 요약)}
\addcontentsline{toc}{section}{빠르게 훑어보기 (1페이지 요약)}

\begin{tcolorbox}[summarybox={모듈 1: 돈과 금융 핵심 요약}]
  \textbf{1. 돈의 본질과 기능:}
  \begin{itemize}
    \item 돈은 \textbf{교환의 매개}, \textbf{가치 척도}, \textbf{가치 저장 수단}의 세 가지 핵심 기능을 수행합니다.
    \item \textbf{상품 화폐} (본질적 가치) $\rightarrow$ \textbf{종이 화폐} (금속 보증) $\rightarrow$ \textbf{법정 화폐} (정부 신뢰)로 진화했습니다.
    \item \textbf{돈 $\neq$ 부}: 돈은 부의 한 형태일 뿐, 모든 자산이 돈은 아닙니다.
  \end{itemize}

  \textbf{2. 지불 시스템:}
  \begin{itemize}
    \item \textbf{현금}과 \textbf{비현금 결제} (수표, ACH, 카드, 디지털 지갑)로 나뉩니다.
    \item \textbf{전자 결제}가 증가 추세이며, 지불 시스템은 경제의 '금융 배관' 역할을 합니다.
    \item \textbf{비트코인}은 현재 돈의 핵심 기능(지불 수단, 가치 저장)을 충족하지 못합니다.
  \end{itemize}

  \textbf{3. 통화량과 인플레이션:}
  \begin{itemize}
    \item \textbf{M1} (현금 + 요구불 예금)과 \textbf{M2} (M1 + 저축성 예금)로 통화량을 측정합니다.
    \item \textbf{인플레이션}은 물가 상승률이며, \textbf{CPI}로 측정합니다.
    \item 과거에는 \textbf{통화량 증가}와 \textbf{인플레이션} 간의 강한 상관관계가 있었으나 (1970년대 대인플레이션), 21세기에는 그 관계가 약화되었습니다.
  \end{itemize}

  \textbf{4. 금융 시스템의 역할:}
  \begin{itemize}
    \item \textbf{저축자} (자금 여유)와 \textbf{차입자} (자금 수요)를 연결하여 자원을 효율적으로 배분합니다.
    \item \textbf{금융 상품} (예: 주식, 채권, 예금, 보험)을 통해 자금 교환이 이루어집니다.
  \end{itemize}

  \textbf{5. 금융 시장 vs. 금융 기관:}
  \begin{itemize}
    \item \textbf{금융 시장:} 브로커/딜러를 통해 직접 거래 (예: 주식 시장). 낮은 거래 비용, 효율적 정보, 투자자 보호가 중요합니다.
    \item \textbf{금융 기관:} 은행과 같은 중개자를 통해 간접 거래. \textbf{만기 불일치}, \textbf{가격 위험}, \textbf{정보 비대칭} 문제를 해결하여 가치를 창출합니다.
  \end{itemize}

  \textbf{6. 은행 vs. 비은행 금융 기관:}
  \begin{itemize}
    \item \textbf{은행:} 예금을 받고 대출을 해주는 전통적 기관 (상업 은행, 저축 기관 등).
    \item \textbf{비은행 금융 기관:} 예금을 받지 않으면서 다양한 금융 서비스 제공 (보험사, 연기금, 증권사, 할부금융사, 정부 후원 기업 등).
    \item 미국에서는 \textbf{은행의 비중이 감소}하고 \textbf{뮤추얼 펀드}와 같은 비은행 기관의 중요성이 증가하는 추세입니다.
  \end{itemize}

  \textbf{7. 금융과 실물 경제:}
  \begin{itemize}
    \item \textbf{금융 발전}은 경제 성장과 위험 분산에 기여하지만, 금융 불안정성과 경제 위기를 초래할 수도 있습니다.
  \end{itemize}
\end{tcolorbox}


\begin{tcolorbox}[title={\textbf{개요: 금융 발전과 경제 활동}}]
금융과 화폐 시스템은 경제의 원활한 기능에 필수적입니다. 금융 발전은 사회 전반의 경제적 성과를 향상시키지만, 금융 시스템의 교란은 막대한 비용을 초래할 수 있습니다. 본 문서에서는 금융 발전이 경제 성장을 촉진하는 증거를 제시하고, 금융 위기의 발생 원인과 사회적 비용을 논의합니다. 궁극적으로 금융 효율성과 안정성 사이의 균형을 찾는 것이 중요하며, 이를 위한 정책적 질문들을 탐구합니다.
\end{tcolorbox}

\section{Lesson 1-3.1: 금융 발전과 경제 활동}

\begin{tcolorbox}[title={\textbf{용어 정리}}]
\begin{adjustbox}{width=\textwidth,center}
\begin{tabular}{lllc}
\toprule
\textbf{용어} & \textbf{쉬운 설명} & \textbf{원어} & \textbf{비고} \\
\midrule
금융 발전 & 금융 시장과 기관이 성장하여 자금 흐름이 효율화되는 과정 & Financial Development & 경제 성장의 핵심 동력 \\
금융 시스템 & 저축자의 자금을 차입자에게 연결하는 시장과 기관의 총체 & Financial System & 은행, 비은행 금융기관, 금융 시장 포함 \\
금융 위기 & 금융 시스템 내에서 투자자 신뢰 상실 등으로 발생하는 심각한 교란 & Financial Crisis & 경제에 막대한 피해를 줄 수 있음 \\
증권화 & 위험한 대출을 묶어 새로운 유가증권을 만드는 과정 & Securitization & 2008년 금융 위기의 주요 원인 중 하나 \\
자산유동화증권 & 대출 채권 등 자산을 담보로 발행되는 증권 & Asset-Backed Securities (ABS) & 위험을 분산하거나 변환하는 데 사용 \\
서브프라임 모기지 & 신용 등급이 낮은 사람들에게 제공되는 주택 담보 대출 & Subprime Mortgages & 부실 위험이 높은 대출 \\
머니 마켓 펀드 & 단기 채권에 투자하는 안전 지향 뮤추얼 펀드 & Money Market Funds & 은행 예금과 유사하지만 정부 보증 없음 \\
초인플레이션 & 물가 상승률이 극도로 높아져 화폐 가치가 급락하는 현상 & Hyperinflation & 정부의 과도한 화폐 발행이 원인 \\
디지털 화폐 & 블록체인 기술 기반의 전자 화폐 (예: 비트코인) & Digital Currencies & 정부 통제 밖의 화폐 \\
스테이블코인 & 가치가 법정 화폐에 고정되도록 설계된 디지털 화폐 & Stablecoins & 테더(Tether)가 대표적 \\
상업어음 & 기업이 단기 자금을 조달하기 위해 발행하는 무담보 어음 & Commercial Paper & 단기 금융 시장의 중요한 도구 \\
정부 구제금융 & 금융 위기 시 정부가 금융 기관에 제공하는 재정 지원 & Government Bailout & 시스템 붕괴 방지를 위한 최후의 수단 \\
\bottomrule
\end{tabular}
\end{adjustbox}
\end{tcolorbox}


\section{핵심 개념 및 원리}

\subsection{금융 발전과 경제 활동의 관계}

\begin{tcolorbox}[title={\textbf{핵심 요약}}]
금융 발전은 경제 성장을 촉진하고 사회 전반의 경제적 성과를 향상시키는 핵심 동력입니다. 이는 자금의 효율적인 배분, 시장 가격의 정확성, 그리고 위험 공유의 개선을 통해 이루어집니다.
\end{tcolorbox}

\subsubsection*{금융 시스템의 역할}
경제의 원활한 기능은 화폐와 지급 시스템에 크게 의존합니다. 금융 시스템은 금융 시장과 다양한 은행 및 비은행 금융 기관을 포함하며, 저축자(savers)로부터 차입자(borrowers)에게 자금을 효율적으로 전달하는 중요한 역할을 합니다. 이러한 자금의 흐름은 투자를 촉진하고, 이는 다시 경제 성장의 기반이 됩니다.

\begin{itemize}
    \item \textbf{자금 중개 (Financial Intermediation):} 금융 시스템은 저축된 자금을 필요로 하는 기업이나 개인에게 연결하여 투자를 가능하게 합니다.
    \item \textbf{자원 배분 효율성 (Efficient Resource Allocation):} 금융 발전은 시장 가격을 더욱 정확하게 만들고, 자원이 가장 생산적인 곳으로 배분되도록 돕습니다.
    \item \textbf{위험 공유 (Risk Sharing):} 다양한 금융 상품과 시장을 통해 투자 위험을 여러 주체에게 분산시켜 개별 경제 주체의 부담을 줄여줍니다.
\end{itemize}
이러한 요소들이 결합될 때, 우리는 더 빠른 경제 성장과 사회 전반의 더 나은 경제적 성과를 기대할 수 있습니다.

\subsubsection*{금융 발전의 측정 및 증거}
금융 발전이 경제 성장에 미치는 긍정적인 영향을 확인하기 위해, 우리는 세계은행(World Bank)의 최신 데이터를 활용하여 이 아이디어를 검증할 수 있습니다.

\begin{itemize}
    \item \textbf{측정 지표:} 금융 발전을 측정하기 위해 'GDP 대비 민간 부문 신용 총액(total amount of credit extended to the private sector scaled by GDP)'을 사용합니다. 이는 국가 간 비교를 위해 GDP로 신용 규모를 표준화한 것입니다.
    \item \textbf{경제 발전 지표:} 경제 발전의 척도로는 '1인당 GDP(GDP per capita)'를 사용합니다.
    \item \textbf{분석 결과:} 전 세계 국가들을 대상으로 한 산점도(scatter plot) 분석 결과는 우리의 직관을 뒷받침합니다. 금융 시스템 규모가 큰 국가일수록 시민들의 소득(1인당 GDP)이 더 높은 경향을 보입니다.
\end{itemize}
이는 금융 발전이 긍정적인 경제적 결과를 가져온다는 강력한 증거입니다. 따라서 정책 입안자라면 더 많은 금융 발전, 더 많은 은행, 더 많은 금융 시장, 그리고 이러한 기관에 대한 규제 완화를 우선시해야 한다고 생각할 수 있습니다.

\begin{tcolorbox}[colback=lightgreen, colframe=green!75!black, title={\textbf{정책적 시사점}}] % mygreen -> lightgreen
금융 발전이 경제 성장에 긍정적인 영향을 미친다는 증거는 정책 입안자들이 금융 부문의 성장을 장려해야 함을 시사합니다. 이는 금융 시장의 효율성을 높이고, 자금 조달을 용이하게 하며, 궁극적으로 국가 경제의 번영에 기여할 수 있습니다.
\end{tcolorbox}

\subsection{금융 시스템의 교란과 금융 위기}
금융 발전이 항상 순조로운 것만은 아닙니다. 때로는 금융 시스템 내에서 심각한 교란이 발생할 수 있으며, 이는 사회에 막대한 비용을 초래하는 금융 위기로 이어집니다.

\subsubsection*{금융 위기의 발생 원인}
금융 위기는 주로 두 가지 주요 원인으로 인해 발생합니다.
\begin{enumerate}
    \item \textbf{투자자 신뢰 상실:} 투자자들이 금융 기관에 대한 신뢰를 잃을 때 발생합니다. 이는 은행 예금 인출 사태(bank run)와 같은 형태로 나타날 수 있습니다.
    \item \textbf{빠르고 통제되지 않는 금융 혁신:} 금융 혁신이 너무 빠르게 진행되고 시장이나 규제 당국에 의해 제대로 검증되지 않을 때 발생합니다. 새로운 금융 상품이나 기술의 위험이 충분히 평가되지 않아 시스템 전체의 불안정성을 초래할 수 있습니다.
\end{enumerate}

\subsubsection*{금융 위기의 전개 과정}
금융 위기는 일반적으로 다음과 같은 단계로 전개되며, 이는 대공황(Great Depression)이나 대침체(Great Recession)와 같은 역사적 사건에서 반복적으로 관찰된 패턴입니다.
\begin{enumerate}
    \item \textbf{과잉과 공황 (Excesses and Panic):} 금융 시스템 내의 과도한 위험 추구나 거품이 공황을 유발합니다.
    \item \textbf{자금 인출 (Withdrawals):} 투자자들이 금융 기관에서 저축을 대규모로 인출하기 시작합니다.
    \item \textbf{자산 청산 (Liquidation):} 자금 인출에 대응하기 위해 금융 기관들은 투자 자산을 강제로 청산해야 합니다. 이는 자산 가격 하락을 더욱 부추깁니다.
    \item \textbf{시장 붕괴 (Market Crash):} 주식 시장과 같은 금융 시장이 붕괴됩니다.
    \item \textbf{경제적 부담 (Economic Burden):} 가계는 저축을 잃고, 신용에 접근할 수 없게 됩니다. 기업들은 자금 조달에 어려움을 겪어 파산합니다.
    \item \textbf{장기 불황 (Deep and Long-lasting Recession):} 이 모든 것이 깊고 장기적인 경기 침체로 이어집니다.
\end{enumerate}
이러한 위기 이야기의 초기 버전에서는 은행이 중심적인 역할을 했습니다.

\subsubsection*{역사 속의 금융 공황}
미국에서는 1793년부터 1797년, 1811년, 1813년, 1816년, 그리고 1873년, 1894년, 1890년, 1893년 등 수많은 공황(panics)이 발생했습니다. 이러한 기록을 가진 시스템의 진정한 이름은 '공황 시스템(panic system)'이라고 불릴 정도였습니다.

\begin{itemize}
    \item \textbf{1907년 공황 (Panic of 1907):} 20세기 최초의 세계적인 금융 위기였습니다.
    \item \textbf{1930년대 대공황 (Great Depression):} 1930년대 대공황 기간 동안 은행 공황은 절정에 달했습니다. 1930년에서 1933년 사이에 매년 약 2,000개의 은행이 파산했습니다. 이는 현대 역사상 최악의 글로벌 금융 위기였을 것입니다.
    \item \textbf{FDIC 설립:} 대공황의 경험은 현재의 정부 예금 보험 시스템인 연방예금보험공사(Federal Deposit Insurance Corporation, FDIC)의 설립으로 이어졌습니다. 이는 예금자들을 보호하고 은행 시스템의 안정성을 높이기 위한 조치였습니다.
\end{itemize}
100년이 지난 지금도 금융 위기에는 여전히 공황에 빠진 투자자들이 관련되어 있습니다.


\subsection{현대 금융 위기의 주요 원인}
21세기 금융 위기는 과거와는 다른 새로운 형태의 금융 혁신과 시장 구조에서 비롯되기도 합니다.

\subsubsection*{증권화 (Securitization)}
증권화는 20세기 위대한 금융 혁신 중 하나입니다.
\begin{itemize}
    \item \textbf{개념:} 금융 공학자들이 위험한 대출 풀(pools of risky loans)을 모아 새로운 자산유동화증권(Asset-Backed Securities, ABS)을 만드는 것을 가능하게 합니다.
    \item \textbf{장점 (겉보기):} 일부 ABS는 트리플 A(Triple A) 등급을 받아, 위험을 안전으로 변환하는 것처럼 보였습니다.
    \item \textbf{문제점:} 그러나 증권화 과정 자체는 복잡하고 불투명하여 많은 투자자들이 이해하기 어려웠습니다.
    \item \textbf{2007년 위기 사례:} 2007년부터 가장 위험한 모기지 담보 증권(mortgage-backed securities), 특히 서브프라임 모기지(subprime mortgages)가 부실화되기 시작했습니다. 투자자들은 이것이 불가능하다고 생각했기 때문에 공황에 빠졌습니다. 이 공황은 모기지와 전혀 관련 없는 학자금 대출 및 신용카드 증권화를 포함한 모든 증권화 시장으로 확산되어 신용 시장의 거의 완전한 붕괴를 초래했습니다.
\end{itemize}

\subsubsection*{머니 마켓 펀드 (Money Market Funds)}
머니 마켓 펀드는 21세기 금융 공황의 중심에 있었습니다.
\begin{itemize}
    \item \textbf{개념:} 이 펀드들은 단기 채권 시장에 매우 안전한 투자를 하는 뮤추얼 펀드의 일종입니다. 대규모 기관 투자자들은 이 펀드에 현금을 예치하여 약간의 수익을 얻습니다.
    \item \textbf{특징:} 이 펀드들은 부채 구조를 은행 예금과 유사하게 만들어 위험이 없어 보이고 요구 시 상환이 가능하도록 합니다.
    \item \textbf{위험:} 그러나 머니 마켓 펀드의 저축은 정부에 의해 보증되지 않으며 손실을 입을 수 있습니다.
    \item \textbf{2008년 위기 사례:} 2008년, 리저브 프라이머리 펀드(Reserve Primary Fund)라는 유명한 펀드에서 손실이 발생했습니다. 이는 전혀 예상치 못한 일이었고, 투자자들은 공황에 빠져 전체 머니 마켓 펀드 산업에서 대규모 자금 인출이 발생하여 붕괴 직전까지 몰렸습니다.
\end{itemize}

\subsubsection*{화폐 및 지급 시스템 문제}
화폐와 지급 시스템 자체에서도 문제가 발생할 수 있습니다.
\begin{itemize}
    \item \textbf{초인플레이션 (Hyperinflation):} 1920년대 바이마르 독일, 1946년 헝가리, 2008년 짐바브웨, 2018년 베네수엘라의 초인플레이션은 극단적인 예시입니다.
    \item \textbf{원인:} 각 사례에서 정부는 너무 많은 지폐를 발행했고, 결국 지폐는 가치를 상실했습니다.
    \item \textbf{결과:} 이는 경제와 사회 전체에 참으로 끔찍한 결과를 초래했습니다.
\end{itemize}


\subsection{디지털 화폐와 금융 시스템}
디지털 화폐의 등장은 화폐 시스템에 새로운 도전과 기회를 동시에 제공하고 있습니다.

\subsubsection*{비트코인 (Bitcoin)}
\begin{itemize}
    \item \textbf{특징:} 디지털 화폐는 정부의 지폐 공급 통제를 벗어납니다. 2008년 비트코인 등장 이후, 비트코인은 신뢰할 수 있는 지급 수단이라기보다는 투기적 거품처럼 보였습니다.
    \item \textbf{위험 평가:} 한 가지 관점은 가격 위험이 소매 투자자들 사이에 분산되어 있는 한, 디지털 화폐의 가치가 오르든 내리든 크게 중요하지 않다는 것입니다. 이는 투자자들의 문제일 뿐, 금융 시스템이나 광범위한 경제에 대한 위협은 아니라는 주장입니다.
\end{itemize}

\subsubsection*{스테이블코인 (Stablecoins)}
스테이블코인은 비트코인과는 다릅니다.
\begin{itemize}
    \item \textbf{개념:} 스테이블코인은 가치를 일반 화폐(예: 달러)에 고정시키는 특별한 유형의 디지털 화폐입니다.
    \item \textbf{테더 (Tether) 사례:} 테더는 가장 인기 있는 스테이블코인입니다. 2021년 기준으로 테더는 600억 달러 이상의 디지털 토큰을 발행했으며, 이 토큰들은 개당 1달러로 교환될 수 있습니다. 이 스테이블코인들은 물건을 구매하는 데 사용될 수 있으며, 달러로 교환될 수 있기 때문에 가치를 가집니다.
    \item \textbf{숨겨진 위험:} 테더는 이러한 디지털 토큰을 판매하여 모금한 자금을 투자합니다. 알려진 바에 따르면, 테더는 모은 달러를 상업어음(commercial paper)과 같은 유가증권에 투자합니다. 2021년 기준으로 테더는 약 300억 달러 상당의 상업어음 투자를 보유하고 있었으며, 이는 사실상 세계에서 가장 큰 머니 마켓 펀드 중 하나였습니다.
    \item \textbf{위기 시나리오:} 만약 투자자들이 테더에 대한 신뢰를 잃고 대규모로 디지털 토큰을 상환(redeem)한다면 어떻게 될까요? 일부 투자자들이 공황에 빠져 토큰당 1달러를 받지 못할까 걱정한다면, 그들은 90센트 또는 그 이하의 가격이라도 기꺼이 받아들일 수 있습니다. 이는 테더에 대한 뱅크런(run)으로 이어질 수 있으며, 디지털 화폐의 가치는 즉시 폭락할 수 있습니다. 이러한 공황은 전체 스테이블코인 시장을 교란시킬 수 있습니다.
    \item \textbf{파급 효과:} 더 나쁜 것은, 이러한 교란이 상업어음 시장으로 파급될 수 있다는 점입니다. 이는 정부와 기업이 운전자본을 조달하는 방식에 영향을 미칠 수 있습니다. 일부 디지털 화폐는 다른 화폐보다 더 문제가 될 수 있음을 시사합니다.
\end{itemize}

\subsection{정부 구제금융 (Government Bailouts)}
금융 시스템이 실패할 때마다 정부가 개입하여 지원을 제공하는 것을 우리는 반복적으로 목격했습니다. 이러한 조치를 일반적으로 '정부 구제금융(government bailout)'이라고 합니다.

\begin{itemize}
    \item \textbf{목적:} 금융 시스템의 붕괴를 막고, 그로 인한 경제적 피해를 최소화하는 것이 주된 목적입니다.
    \item \textbf{2008년 미국 사례:} 미국에서는 2008년 금융 위기 동안 및 이후에 다양한 정부 기관이 6천억 달러 이상의 납세자 자금을 지출했습니다. 이는 막대한 금액이지만, 미국은 매우 큰 경제 규모를 가지고 있어 개입이 성공한다면 그 정도의 돈을 지출할 여력이 있습니다. 다행히 2008년에는 대체로 성공적으로 작동했습니다.
    \item \textbf{다른 국가의 한계:} 그러나 다른 정부들은 그 정도 규모의 구제금융을 감당할 수 없습니다.
    \item \textbf{비용의 상대적 규모:} 최근 역사상 가장 큰 은행 구제금융 사례를 GDP 대비 규모로 보면, 1997년 인도네시아는 국가 소득의 절반 이상을 은행 구제금융에 지출했습니다. 이 정도 규모의 구제금융은 정부 재정의 갑작스러운 붕괴로 이어질 수 있습니다.
\end{itemize}


\subsection{금융 정책의 트레이드오프와 규제 질문}
금융 발전과 안정성 사이에는 본질적인 트레이드오프(tradeoff)가 존재합니다.

\begin{tcolorbox}[title={\textbf{핵심 트레이드오프}}]
우리는 효율적이고 경쟁적인 금융 부문을 육성하고, 화폐와 금융의 혁신이 경제를 발전시키도록 장려해야 합니다. 그러나 동시에 금융 부문의 안정성과 위기 발생 위험 사이에서 균형을 찾아야 합니다.
\end{tcolorbox}

위기를 방지하기 위해서는 한계를 설정해야 합니다. 이는 중요한 공공 정책 질문으로 이어집니다.

\begin{enumerate}
    \item \textbf{규제 대상 기관:} 어떤 기관을 규제해야 하는가? 사모펀드(private equity)와 핀테크(FinTech) 기업은 규제 당국이 그대로 두어야 하는가? 이는 여전히 열린 질문입니다.
    \item \textbf{규제 기관의 허용 범위:} 규제 대상 금융 기관은 무엇을 할 수 있도록 허용되어야 하는가?
    \begin{itemize}
        \item 테더는 은행과 동일한 공시 의무를 가져야 하는가?
        \item 은행은 헤지펀드 자회사(hedge funds subsidiaries)를 가질 수 있도록 허용되어야 하는가?
        \item 이러한 헤지펀드 자회사가 비트코인에 투자하는 것을 허용해야 하는가?
    \end{itemize}
    \item \textbf{금융 위기 시 정부 개입:} 금융 위기 동안 어떤 종류의 정부 개입이 적절한가? 우리의 직관은 때때로 실패하는 기관을 구제하는 것이 적절하다는 것입니다. 그러나 그 한계는 어떻게 설정되어야 하는가?
\end{enumerate}
이러한 질문들은 금융 시스템의 미래와 경제 전반의 안정성을 결정하는 데 매우 중요합니다.


\begin{tcolorbox}[title={\textbf{빠르게 훑어보기: Module 1 핵심 요약}}]
\begin{itemize}
    \item \textbf{금융 발전의 중요성:} 금융 시스템의 발전은 자금 중개, 자원 배분 효율성, 위험 공유를 통해 경제 성장을 촉진합니다. (예: GDP 대비 민간 신용 규모가 클수록 1인당 GDP 높음)
    \item \textbf{금융 위기의 그림자:} 금융 시스템은 투자자 신뢰 상실이나 통제되지 않는 혁신으로 교란될 수 있으며, 이는 공황, 자금 인출, 시장 붕괴, 장기 불황으로 이어지는 금융 위기를 초래합니다. (예: 1930년대 대공황, 2008년 대침체)
    \item \textbf{현대 위기의 원인:} 증권화(서브프라임 모기지), 머니 마켓 펀드(2008년 리저브 프라이머리 펀드), 초인플레이션(정부의 과도한 화폐 발행) 등이 현대 금융 위기의 주요 원인입니다.
    \item \textbf{디지털 화폐의 양면성:} 비트코인은 투기적 성격이 강하며, 스테이블코인(테더)은 안정성을 추구하지만, 대규모 인출 시 상업어음 시장 등 광범위한 금융 시스템에 위험을 전이시킬 수 있습니다.
    \item \textbf{정부의 역할과 딜레마:} 금융 위기 시 정부 구제금융은 시스템 붕괴를 막지만, 막대한 비용과 정부 재정 부담을 초래할 수 있습니다. (예: 1997년 인도네시아 GDP의 절반 이상 지출)
    \item \textbf{정책적 트레이드오프:} 금융 효율성과 혁신을 장려하는 동시에 금융 시스템의 안정성을 확보하는 것이 중요합니다. 규제 대상, 규제 범위, 위기 시 개입 한계에 대한 지속적인 논의가 필요합니다.
\end{itemize}
\end{tcolorbox}


%=======================================================================
% Module 2: Module 2: 현대 은행업 (Modern Banking)
%=======================================================================
\chapter{Module 2: 현대 은행업 (Modern Banking)}
\label{ch:module2}

\metainfo{금융기관론}{Lecture 2}{Prof. Rustom Manouchehri Irani}{현대 은행업의 비즈니스 모델, 조직, 기능, 재무제표 분석을 통한 이해}


\section{Module 2: 현대 은행업 (Modern Banking)}
\addcontentsline{toc}{section}{Module 2: 현대 은행업 (Modern Banking)}

\begin{summarybox}[title={개요: Module 2 핵심 요약}]
이 모듈은 현대 은행의 비즈니스 모델을 심층적으로 탐구합니다. 미국 은행업의 간략한 역사부터 시작하여, 현대 은행의 조직과 기능(상업 은행 및 투자 은행)을 상세히 설명합니다. 마지막으로, 현대 은행의 재무제표 분석을 통해 이들 금융기관의 운영 방식을 이해하는 것을 목표로 합니다. 이 모듈을 마치면 현대 은행의 기원, 역할, 그리고 수익 창출 방식을 명확히 이해할 수 있을 것입니다.
\end{summarybox}


\section{용어 정리}
\addcontentsline{toc}{section}{용어 정리}

\begin{adjustbox}{width=\textwidth,center}
\begin{tabular}{lllc}
\toprule
\textbf{용어 (Term)} & \textbf{쉬운 설명 (Easy Explanation)} & \textbf{원어 (Original)} & \textbf{비고 (Notes)} \\
\midrule
텔러 & 은행 창구에서 고객을 응대하는 직원 & Teller & 과거에는 흔했지만 지금은 드물다. \\
유닛 뱅킹 & 은행이 단 하나의 물리적 지점만 가질 수 있는 형태 & Unit Banking & 일리노이주에서 1985년까지 시행. \\
지점 네트워크 & 은행이 여러 물리적 지점을 운영하는 형태 & Branch Network & 캘리포니아주에서 일찍이 허용. \\
은행 공황 & 예금자들이 은행 파산을 우려하여 예금을 대량 인출하는 현상 & Banking Panic & 예금 보험이 없던 시절 흔했다. \\
글래스-스티걸 법 & 1933년 제정된 법으로, 상업 은행과 투자 은행 업무를 분리 & Glass-Steagall Act & 은행 시스템 안정화에 기여했으나 비효율성 비판도 받음. \\
연방 예금 보험 공사 & 은행 파산 시 예금자에게 예금을 보장해주는 기관 & Federal Deposit Insurance Corporation (FDIC) & 예금자 보호를 통해 은행 공황 방지. \\
규제 Q & 예금 계좌에 대한 이자 지급을 금지했던 규제 & Regulation Q & 은행 간 과도한 경쟁을 막기 위함. \\
은행 지주 회사 & 다른 은행이나 비은행 자회사를 소유하는 껍데기 회사 & Bank Holding Company (BHC) & 지점 제한 규제를 우회하는 데 사용. \\
규제 차익 거래 & 법의 허점을 이용하여 규제를 회피하는 행위 & Regulatory Arbitrage & BHC가 지점 제한을 우회한 사례. \\
3-6-3 규칙 & 예금 3\% 이자, 대출 6\% 이자, 오후 3시 골프를 치는 은행업 모델 & 3-6-3 Rule & 규제 보호 하에 은행이 누리던 독점적 이익. \\
머니 마켓 뮤추얼 펀드 & 은행 예금과 유사하지만 규제 Q의 적용을 받지 않는 투자 상품 & Money Market Mutual Fund & 은행 예금의 대체재로 등장. \\
증권화 & 은행 대출을 묶어 유가증권으로 만들어 판매하는 과정 & Securitization & 은행 대출의 새로운 자금 조달 모델. \\
리가-닐 법 & 은행의 주(州) 경계를 넘는 지리적 확장을 허용한 법 & Riegle-Neal Act & 은행 통합 및 효율성 증대. \\
그램-리치-블라일리 법 & 글래스-스티걸 법을 폐지하여 상업/투자 은행 업무 통합을 허용한 법 & Gramm-Leach-Bliley Act & 금융 대형화 및 복합 금융 서비스 등장. \\
금융 지주 회사 & 은행, 증권, 보험 등 다양한 금융 서비스를 제공하는 대형 금융 그룹 & Financial Holding Company & 현대 메가 은행의 조직 형태. \\
소매 금융 & 개인 및 소규모 기업 고객을 대상으로 하는 은행 서비스 & Retail Banking & 예금, 주택 담보 대출, 신용 카드 등. \\
도매 금융 & 대기업 및 기관 투자자를 대상으로 하는 은행 서비스 & Wholesale Banking & 기업 대출, 투자 은행 서비스 등. \\
커뮤니티 은행 & 지역 사회의 개인 및 소규모 기업을 위한 소규모 은행 & Community Bank & 자산 10억 달러 미만. \\
지역/초지역 은행 & 커뮤니티 은행보다 크고 전국적인 지점망을 가진 은행 & Regional/Super-regional Bank & 자산 수백억~수조 달러. \\
머니 센터 은행 & 도매 자금 조달에 크게 의존하는 대형 은행 & Money Center Bank & 뉴욕 등 주요 금융 중심지에 위치. \\
거래 계좌 & 수시 입출금이 가능하고 이자가 없는 예금 계좌 & Transaction Account & 요구불 예금 (Checking Account). \\
비거래 계좌 & 저축성 예금으로, 이자가 지급되지만 인출 제한이 있는 계좌 & Non-transaction Account & 저축 예금, 소액 양도성 예금 증서 (CD). \\
연방 기금 시장 & 은행 간 단기 무담보 대출이 이루어지는 시장 & Federal Funds Market & 은행 간 자금 부족 해소. \\
환매 조건부 채권 & 증권을 담보로 단기 자금을 빌리는 계약 & Repurchase Agreement (Repo) & 담보 대출의 일종. \\
은행 자본 & 은행 소유주의 투자금으로, 주식 및 유보 이익으로 구성 & Bank Capital & 손실 흡수 능력 제공, 가장 안정적인 자금원. \\
소비자 신용 & 개인 고객에게 제공되는 대출 & Consumer Credit & 주택 담보 대출, 자동차 대출, 학자금 대출 등. \\
상업 및 산업 대출 & 기업 고객에게 제공되는 대출 & Commercial and Industrial (C\&I) Lending) & 상업용 부동산 대출, 기간 대출, 사업 신용 한도 등. \\
관계형 대출 & 은행이 고객과의 장기적인 관계를 바탕으로 제공하는 대출 & Relationship Lending & 중견 기업 대출에 주로 적용. \\
대출 개시 & 대출 신청을 심사하고 대출 조건을 협상하여 승인하는 과정 & Origination & 대출 과정의 첫 단계. \\
대출 자금 조달 & 승인된 대출에 필요한 자금을 실제로 제공하는 과정 & Financing & 은행의 자금원 활용. \\
대출 관리 & 대출 실행 후 상환금을 수취하고 차입자의 상태를 모니터링하는 과정 & Servicing & 대출의 지속적인 관리. \\
개시 및 보유 모델 & 은행이 대출을 개시하고 자금을 조달하며 만기까지 보유하는 전통적 모델 & Originate and Hold Model & 모든 과정을 은행이 직접 수행. \\
개시 및 분배 모델 & 은행이 대출을 개시한 후 다른 기관에 판매하거나 증권화하는 모델 & Originate and Distribute Model & 수수료 수입에 중점, 효율성 증대. \\
핀테크 & 금융 기술을 의미하며, 새로운 기술을 활용한 금융 서비스 & Fintech & 온라인 대출 등으로 은행업을 혁신. \\
\bottomrule
\end{tabular}
\end{adjustbox}


\section{Lesson 2-0: 개요 (Overview)}
\addcontentsline{toc}{section}{Lesson 2-0: 개요 (Overview)}

안녕하세요, 금융기관 모듈 2에 오신 것을 환영합니다. 이 모듈은 \textbf{현대 은행업(Modern Banking)}이라고 불립니다.

\begin{summarybox}[title={모듈 2의 학습 목표}]
이 모듈에서는 현대 은행의 비즈니스 모델을 심층적으로 다룰 것입니다.
\begin{itemize}
    \item \textbf{미국 은행업의 간략한 역사:} 현대 은행업을 완전히 이해하기 위해 그 기원을 살펴봅니다.
    \item \textbf{현대 은행의 조직과 기능:} 상업 은행과 투자 은행 사업 모델을 분리하여 각각을 분석합니다.
    \item \textbf{재무제표 분석:} 은행의 대차대조표(Balance Sheet)와 손익계산서(Income Statement)를 면밀히 검토하여 이들 금융기관에서 실제로 어떤 일이 벌어지고 있는지 이해합니다.
\end{itemize}
이 모듈을 마치면 현대 은행이 어디에서 왔고, 무엇을 하며, 어떻게 이익을 창출하는지 명확하게 이해할 수 있을 것입니다.
\end{summarybox}


\section{Lesson 2-1: 미국 은행업의 간략한 역사 (A Brief History of Banking in the United States)}
\addcontentsline{toc}{section}{Lesson 2-1: 미국 은행업의 간략한 역사 (A Brief History of Banking in the United States)}

\subsection{Lesson 2-1.1: 미국 초기 은행업 (Early Banking in the United States)}
\addcontentsline{toc}{subsection}{Lesson 2-1.1: 미국 초기 은행업 (Early Banking in the United States)}

우리가 '은행 시스템'이라고 부르는 것은 지난 100년 동안 극적으로 변화했습니다.

\begin{itemize}
    \item \textbf{과거:} 우리 부모님 세대는 아마도 동네 은행을 방문하여 텔러(Teller)와 직접 대화했을 것입니다. 텔러는 은행에서 일하는 실제 사람을 의미합니다.
    \item \textbf{현재:} 요즘에는 ATM조차 거의 사용하지 않습니다. 대부분의 은행 업무는 온라인이나 스마트폰을 통해 이루어집니다.
\end{itemize}

은행업의 기술적, 경쟁적 환경이 변화함에 따라 은행 산업의 구조도 변화했습니다. 여기서 '구조'는 두 가지를 의미합니다.
\begin{enumerate}
    \item \textbf{은행의 비즈니스 모델:} 은행이 판매하는 상품의 종류를 포함합니다.
    \item \textbf{은행의 수, 규모, 물리적 위치:} 은행의 지리적 분포를 의미합니다.
\end{enumerate}
이번 세션에서는 미국 은행 산업의 역사적 추세에 초점을 맞출 것입니다. 기술과 규제의 변화가 21세기 은행의 모습에 어떻게 영향을 미쳤는지 살펴보겠습니다.

\subsubsection*{전통적인 은행의 비즈니스 모델}
전통적인 은행은 여러 가지 일을 하지만, 그 비즈니스 모델은 크게 다음으로 요약할 수 있습니다.
\begin{itemize}
    \item \textbf{예금 수취:} 예금자로부터 돈(달러)을 받아 은행 계좌를 제공합니다.
    \item \textbf{자금 대출:} 이 자금을 주로 대출 형태로 차입자에게 전달합니다.
\end{itemize}

\begin{examplebox}[title={은행 계좌의 가치}]
은행 계좌는 기본적으로 돈과 같습니다. 우리는 은행에 우리의 자산을 안전하게 보관할 수 있으며, 가격 변동 위험 없이 결제 시스템에 접근할 수 있습니다. 예를 들어, 수표를 발행하거나 직불 카드(Debit Card)를 사용할 수 있습니다.
\end{examplebox}

은행은 또한 신용도 높은 차입자에게만 대출이 이루어지도록 함으로써 가치를 더합니다. 전통적인 은행은 예금을 빌려주고, 대출 이자 수입과 예금 이자 지급의 차이(스프레드)를 통해 수익을 얻습니다. 이들은 노력과 위험 감수에 대한 보상을 받습니다.

\subsubsection*{현대 은행업의 복잡성}
현대 은행업은 매우 복잡합니다. 예를 들어, JP모건 체이스(JPMorgan Chase)의 조직도를 보면 이 회사가 운영하는 다양한 사업 부문을 알 수 있습니다. 여기에는 일반 개인 고객을 위한 \textbf{소비자 금융(Consumer Banking)} 사업뿐만 아니라 훨씬 더 많은 상품을 포함하는 \textbf{도매 금융(Wholesale Banking)} 사업도 포함됩니다. JP모건과 같은 조직은 수만 개의 자회사를 통해 모든 종류의 금융 서비스를 제공합니다.

\begin{center}
    \textit{어떻게 전통적인 은행 모델에서 JP모건 체이스와 같은 '금융 슈퍼마켓'으로 발전하게 되었을까요?}
\end{center}

\subsubsection*{은행의 조직 형태와 주요 논점}
이번 세션에서는 세 가지 주요 조직 형태를 논의할 것입니다.
\begin{enumerate}
    \item \textbf{유닛 뱅크 (Unit Bank):} 단일 지점 은행
    \item \textbf{지점 네트워크 (Branch Network):} 여러 지점을 가진 은행
    \item \textbf{은행 지주 회사 (Bank Holding Company):} 여러 은행 및 비은행 자회사를 소유하는 회사
\end{enumerate}
이 논의를 관통하는 두 가지 핵심 주제가 있습니다.
\begin{enumerate}
    \item \textbf{은행이 법적으로 운영할 수 있는 지역:} 은행이 지리적으로 어디까지 확장할 수 있는가?
    \item \textbf{은행이 판매할 수 있는 금융 상품 및 서비스의 종류:} 은행이 어떤 종류의 사업을 할 수 있는가?
\end{enumerate}
21세기에는 은행이 원하는 곳 어디에서든 운영하고 상상할 수 있는 모든 금융 서비스를 제공할 수 있는 것처럼 보입니다. 그러나 이것이 항상 그랬던 것은 아니라는 점을 인식하는 것이 중요합니다. 역사적으로 은행의 확장이나 사업 변경 허용은 격렬한 정책 논쟁의 원천이었습니다.

\subsubsection*{지리적 확장 제한: 유닛 뱅킹 vs. 지점 네트워크}
은행이 지리적으로 어떻게 분포되었는지 살펴보겠습니다. 이는 은행의 물리적인 지점 수와 위치를 의미합니다. 과거에는 은행이 가질 수 있는 지점 수에 많은 제한이 있었습니다. 주(州)마다 접근 방식이 달랐습니다.

\begin{itemize}
    \item \textbf{엄격한 주 (예: 일리노이):} 일리노이주는 \textbf{유닛 뱅킹(Unit Banking)}만을 허용했습니다. 이는 은행이 단 하나의 물리적 위치만 가질 수 있다는 것을 의미했습니다.
    \begin{examplebox}[title={콘티넨탈 일리노이 은행 (Continental Illinois Bank)}]
    시카고 금융 지구의 라살 스트리트 208번지에 위치한 콘티넨탈 일리노이 은행이 그 예입니다. CEO를 포함한 모든 직원이 한 곳에 있었고, 모든 대출 기록도 이 건물에 보관되었습니다. 고객들은 어디로 가야 할지 명확히 알았습니다. 유닛 뱅킹은 일리노이주에서 1870년부터 1985년까지 법으로 시행되었습니다.
    \end{examplebox}
    \item \textbf{자유로운 주 (예: 캘리포니아):} 캘리포니아주는 항상 은행이 주 경계 내에서 \textbf{지점 네트워크(Branch Network)}를 가질 수 있도록 허용했습니다.
    \begin{examplebox}[title={뱅크 오브 아메리카 (Bank of America) 지점 네트워크 (1929년경)}]
    1929년경 뱅크 오브 아메리카의 지점 네트워크를 보면, 샌프란시스코에 본사를 두고 주 전역에 지점을 운영했습니다.
    \end{examplebox}
\end{itemize}

일리노이주와 다른 주들이 엄격한 지점 규제를 시행한 데에는 여러 가지 이유가 있었습니다.
\begin{itemize}
    \item \textbf{지점 확장에 반대하는 주장:} 은행이 너무 커져 독점력을 행사할 수 있다는 우려가 있었습니다. 특히 소도시의 은행가들은 대도시 은행과의 경쟁으로부터 자신들을 보호하기 위해 지점 제한을 강력히 지지했습니다.
    \item \textbf{지점 확장에 찬성하는 주장:} 다른 주들은 더 크고 지리적으로 분산된 은행이 더 회복력 있고 안전할 것이라고 믿었습니다.
\end{itemize}
유닛 뱅킹과 지점 네트워크 중 어느 것이 더 안전한지에 대한 의견 불일치가 많았지만, 이 시기에는 \textbf{은행 공황(Banking Panics)}이 흔했습니다.

\begin{examplebox}[title={은행 공황에 대한 유명한 인용문}]
\begin{itemize}
    \item "그때는 은행 공황이 얼마나 자주 일어났는지"를 지적하는 첫 번째 인용문.
    \item "결국 전 세계적인 범위로 확대되었다"는 두 번째 인용문.
\end{itemize}
\end{examplebox}

공황은 다음과 같이 작동합니다. 일반적으로 일부 예금자들이 저축을 잃을까 봐 걱정하기 시작합니다. 왜 이런 일이 일어났을까요? 예금자들은 경제가 악화될 수 있고, 자신들의 은행이 대출에서 큰 손실을 입어 결국 파산할 수 있다고 우려하기 시작했습니다. 이때는 \textbf{예금 보험(Deposit Insurance)}이 없던 시절이었습니다. 이는 제때 돈을 인출하지 못한 예금자들만 손실을 입게 된다는 것을 의미했습니다. 따라서 '뱅크 런(Run on the Bank)'이 발생했습니다.

이러한 은행 공황은 1930년대 \textbf{대공황(Great Depression)} 동안 절정에 달했습니다. 이는 아마도 현대 역사상 최악의 글로벌 금융 위기였을 것입니다. 미국 내 모든 은행의 3분의 1 이상이 파산했고, 수백만 가구가 평생 모은 저축을 잃었으며, 경제는 붕괴했습니다.

\subsubsection*{1933년 은행법 (글래스-스티걸 법)}
은행 위기와 그에 따른 대공황에 대한 비상 대응으로, 프랭클린 루즈벨트 대통령은 \textbf{1933년 은행법(Banking Act of 1933)}을 제정했습니다. 이 유명한 법안은 \textbf{글래스-스티걸 법(Glass-Steagall Act)}으로도 알려져 있습니다.

\begin{summarybox}[title={글래스-스티걸 법의 주요 변화}]
\begin{itemize}
    \item \textbf{주(州) 간 은행 지점 설치 금지 (Ban on Interstate Bank Branching):} 은행이 주 경계를 넘어 자회사를 가질 수 없도록 하는 지리적 제한이었습니다.
    \item \textbf{투자 은행 업무 금지 (Ban on Investment Banking):} 상업 은행은 사모 증권의 인수(Underwriting)나 거래, 심지어 보험 판매도 금지되었습니다.
    \item \textbf{규제 Q (Regulation Q) 포함:} 예금 계좌에 대한 이자 지급을 금지했습니다.
    \item \textbf{연방 예금 보험 공사 (Federal Deposit Insurance Corporation, FDIC) 설립:} 은행이 파산할 경우 개별 예금자에게 예금을 보장하여 저축을 잃지 않도록 했습니다.
\end{itemize}
\end{summarybox}

이러한 금지 조치들은 은행 공황의 원인에 대한 당시의 믿음을 반영했습니다.
\begin{itemize}
    \item 주(州) 간 지점 확장은 공황을 들불처럼 퍼뜨린다고 보았습니다.
    \item 투자 은행 업무는 너무 위험하다고 여겨졌습니다.
    \item 예금에 대한 이자 지급은 은행 간 과도한 경쟁을 유발한다고 생각했습니다.
    \item 뱅크 런은 저축을 잃을까 봐 걱정하는 혼란스러운 투자자들 때문에 발생한다고 보았습니다.
\end{itemize}
글래스-스티걸 법은 고도로 규제되고 보호받는 상업 은행 산업을 만들었습니다.

\begin{figure}[h!]
    \centering
    % \includegraphics[width=0.8\textwidth]{example_graph_banks_1935_1956.png}  % Image not found: example_graph_banks_1935_1956.png % 가상의 이미지 파일명
    \caption{1935년에서 1956년 사이 미국 내 보험 가입 상업 은행 수 (단위 은행 및 다중 지점 은행)}
    \label{fig:banks_1935_1956}
    \textit{이 그래프는 1935년에서 1956년 사이의 보험 가입 상업 은행 수를 보여줍니다. 이 기간 동안 은행의 수는 많고 안정적이었으며, 대부분은 지역 독점권을 가진 유닛 뱅크였습니다. 이는 정부가 보장하는 은행 시스템으로, 단순한 상품 제공과 제한된 위험 감수 기회를 가졌습니다.}
\end{figure}

글래스-스티걸 법은 은행 시스템에 대한 신뢰를 회복시켰습니다.

\begin{figure}[h!]
    \centering
    % \includegraphics[width=0.8\textwidth]{example_graph_bank_failures.png}  % Image not found: example_graph_bank_failures.png % 가상의 이미지 파일명
    \caption{미국 내 은행 파산 건수 (1933년 이후)}
    \label{fig:bank_failures}
    \textit{이 차트는 은행 파산 건수가 1933년 약 4,000건에서 효과적으로 0건으로 급감했음을 보여줍니다. 그러나 당시 많은 비평가들은 고도로 규제된 미국 은행 시스템이 너무 파편화되어 있고 비효율적이라고 보았습니다.}
\end{figure}


\subsection{Lesson 2-1.2: 은행 지주 회사 (Bank Holding Companies)}
\addcontentsline{toc}{subsection}{Lesson 2-1.2: 은행 지주 회사 (Bank Holding Companies)}

다음으로 등장한 지배적인 조직 형태는 \textbf{은행 지주 회사(Bank Holding Company, BHC)}였습니다.

\begin{itemize}
    \item \textbf{지주 회사(Holding Company):} 다른 기업이나 자회사 그룹을 소유하는 껍데기 회사(Shell Corporation)입니다.
    \item \textbf{은행 지주 회사(BHC):} 하나 이상의 은행 자회사를 소유합니다.
\end{itemize}

\begin{examplebox}[title={규제 회피를 위한 은행 지주 회사 활용}]
1950년이라고 가정해 봅시다. 제가 샴페인(Champaign)에 매우 성공적인 은행을 소유하고 있고, 지리적으로 다각화하기 위해 세인트 클레어 카운티(Saint Clair County)로 사업을 확장하고 싶습니다. 그러나 일리노이주는 유닛 뱅킹만 허용하므로 이는 불가능합니다.

지점 제한을 회피하는 한 가지 방법은 다음과 같습니다.
\begin{enumerate}
    \item 세인트 클레어에 완전히 별개의 은행을 설립합니다.
    \item 두 은행을 모두 인수하는 지주 회사를 설립합니다.
\end{enumerate}
이러한 은행 자회사들은 독립적인 은행으로 간주될 수 있으므로 법을 준수하는 것으로 보일 수 있습니다.

더 나아가, 제 지주 회사가 일부 자산 관리 서비스를 제공하는 비은행 자회사(Non-bank Subsidiary)를 설립할 수도 있습니다. 이 비은행 자회사는 어떤 은행과도 직접적으로 관련이 없으므로, 아마도 문제가 없을 것입니다.
\end{examplebox}

이러한 유형의 행태를 종종 \textbf{규제 차익 거래(Regulatory Arbitrage)}라고 부르는데, 이는 법의 허점을 이용한 것이기 때문입니다. 1956년과 1970년의 \textbf{은행 지주 회사법(Bank Holding Company Acts)}을 통해 의회는 이러한 허점을 막으려 했지만, 이미 '고양이는 가방 밖으로 나와버린' 상태였습니다. (즉, 이미 널리 퍼져 되돌리기 어려웠다는 의미입니다.)

\begin{figure}[h!]
    \centering
    % \includegraphics[width=0.8\textwidth]{example_graph_banks_1980s.png}  % Image not found: example_graph_banks_1980s.png % 가상의 이미지 파일명
    \caption{1980년대까지의 미국 내 은행 조직 형태 변화}
    \label{fig:banks_1980s}
    \textit{이 차트는 1980년대까지의 은행 조직 형태 변화를 보여줍니다. 1980년대에 이르러 은행 지주 회사를 통한 지점 확장이 지배적인 조직 형태가 되었습니다.}
\end{figure}

1935년부터 1980년까지의 기간은 미국 상업 은행업의 전성기로 여겨지기도 합니다. 은행들은 산업으로서 독점적인 권한을 가졌습니다. 이는 상업 은행업의 유명한 \textbf{3-6-3 규칙(3-6-3 Rule)}을 낳았습니다.

\begin{examplebox}[title={3-6-3 규칙에 대한 인용문}]
"미국 은행가들은 수십 년 동안 3-6-3 규칙에 따라 운영되었습니다. 예금자에게는 3\% 이자를 지급하고, 돈은 6\%에 대출하며, 오후 3시에는 골프를 치러 갔습니다. 연방 및 주 법률이 그들이 운영하는 엄격한 규칙을 정하고 경쟁자로부터 그들을 보호했기 때문에 그들은 그렇게 정확하게 행동할 수 있었습니다."

"그 결과, 은행가들의 권력과 명성은 그들의 금고만큼이나 안전했고, 이익은 꾸준하고 확실했습니다."
\end{examplebox}

그러나 이 모든 것은 1970년대와 80년대부터 변하기 시작했습니다. "갑자기 그 모든 것이 사라졌습니다. 은행가들은 이제 대공황 이후 가장 힘든 생존 시험에 직면해 있습니다."

\begin{figure}[h!]
    \centering
    % \includegraphics[width=0.8\textwidth]{example_graph_commercial_banks_peak.png}  % Image not found: example_graph_commercial_banks_peak.png % 가상의 이미지 파일명
    \caption{미국 내 상업 은행 수의 변화 (1985년 이후)}
    \label{fig:commercial_banks_peak}
    \textit{이 차트는 미국 내 고유 상업 은행의 수가 1985년경 정점을 찍었음을 보여줍니다. 그 후에는 개별 기관의 수가 급격히 감소했습니다. 이는 극적인 통합(Consolidation)의 시기였습니다. 1999년에는 은행 수가 약 절반으로 줄어들어 약 8,000개가 되었습니다.}
\end{figure}

\subsubsection*{은행 통합의 주요 원인}
그렇다면 이러한 통합의 주요 원인은 무엇이었을까요?

\begin{enumerate}
    \item \textbf{기술 (Technology):}
    1970년대부터 우편, 전화, 팩스, 그리고 결국 인터넷이 보편화되면서 은행업의 지역적 특성이 덜 중요해졌습니다. 우리는 더 이상 동네 은행에 직접 가서 대출 담당자와 이야기할 필요가 없어졌습니다. 요즘에는 스마트폰으로 예금, 직불 카드, 주택 담보 대출 상환 등 모든 것을 관리할 수 있습니다. 이러한 변화는 1970년대에 시작되었습니다.

    \item \textbf{채권 시장의 금융 혁신 (Financial Innovations in Debt Markets):}
    이는 예금을 유치하고 대출을 제공하는 전통적인 은행의 능력에 대한 공격이었습니다.
    \begin{itemize}
        \item \textbf{예금 대체재 등장:} 앞서 언급했듯이, 규제 Q(Regulation Q) 하에서는 은행이 예금 계좌에 이자를 지급하는 능력이 제한되었습니다. 1980년대부터 \textbf{머니 마켓 뮤추얼 펀드(Money Market Mutual Funds)}가 등장했습니다. 이 뮤추얼 펀드는 비은행 기관으로, 예금과 유사한 계좌를 제공하지만 규제 Q의 적용을 받지 않아 매우 안전하고 편리했습니다.
        \item \textbf{은행 대출 대체재 등장:} 1980년대에는 공공 채권 시장에서 엄청난 변화가 있었습니다. 기업 금융 측면에서는 단기 기업 어음(Commercial Paper)과 장기 정크 본드(Junk Bond) 시장이 등장했습니다. 그리고 \textbf{증권화(Securitization)}가 은행 대출을 위한 완전히 새로운 자금 조달 모델을 만들어냈습니다. 이는 은행 산업에 더 큰 혼란을 야기했습니다.
    \end{itemize}
\end{enumerate}

이러한 변화는 은행에 대한 주요 경쟁 압력의 원천이었습니다. 지역 은행 독점은 붕괴되었습니다. 이로 인해 은행들은 어려운 선택에 직면했습니다.
\begin{itemize}
    \item 사업을 중단하거나,
    \item 유기적으로 규모를 확장하거나,
    \item 합병 또는 인수를 통해 규모를 확장하는 것이었습니다.
\end{itemize}
규모를 확장함으로써 은행들은 \textbf{규모의 경제(Economies of Scale)}와 \textbf{범위의 경제(Economies of Scope)}를 활용하여 중복 비용을 제거할 수 있었습니다.


\subsection{Lesson 2-1.3: 규제 완화와 현대 은행업 (Deregulation and Modern Banking)}
\addcontentsline{toc}{subsection}{Lesson 2-1.3: 규제 완화와 현대 은행업 (Deregulation and Modern Banking)}

이러한 경쟁의 시기는 은행 규제에 대한 패러다임 전환으로 이어졌습니다. 입법자들과 규제 당국은 과도한 은행 위험 감수나 독점적 은행에 대한 우려를 덜게 되었습니다. 대신, 경쟁에 실패하여 문을 닫는 은행에 대해 더 걱정하게 되었습니다.

상업 은행 산업을 현대화한 두 가지 주요 법안이 있었습니다. 둘 다 빌 클린턴 대통령에 의해 법으로 서명되었습니다. 이러한 규제 완화는 은행이 규모의 경제와 범위의 경제를 활용하여 효율성을 높이는 데 도움이 되도록 설계되었습니다.

\begin{enumerate}
    \item \textbf{리가-닐 법 (Riegle-Neal Act):}
    이 법은 은행이 주(州) 경계를 넘어 운영함으로써 지리적으로 다각화할 수 있도록 허용했습니다. 이제 은행들은 미국 전역으로 사업을 자유롭게 확장할 수 있게 되었습니다. 예를 들어, 다른 주에 있는 은행을 인수하는 방식으로 말입니다. \textbf{유닛 뱅킹(Unit Banking)}은 과거의 일이 되었습니다.

    \item \textbf{그램-리치-블라일리 법 (Gramm-Leach-Bliley Act):}
    이 법은 1933년 글래스-스티걸 법에 있던 투자 은행 업무 금지 조항을 폐지했습니다. 이제 상업 은행업, 투자 은행업, 보험 인수(Insurance Underwriting)가 모두 한 지붕 아래에서 이루어질 수 있게 되었습니다. 이는 은행이 제품 라인을 다각화할 수 있도록 허용했습니다.
\end{enumerate}

\begin{examplebox}[title={제품 라인 다각화의 이점: 정보 비용 절감}]
은행은 정보 사업을 합니다. 여러 제품 라인에 걸쳐 운영하면 정보 비용을 낮출 수 있습니다.
\begin{itemize}
    \item 예를 들어, 고객이 당좌 예금 계좌에 꾸준한 수입 흐름을 가지고 있다는 것을 안다면, 은행은 그 고객에게 생명 보험에 대한 더 좋은 조건을 제시할 의향이 있을 수 있습니다.
    \item 또는 과거에 대출을 해준 기업 고객이라면, 다음 채권 발행을 인수하는 데 더 유리한 위치에 있을 수 있습니다.
\end{itemize}
\end{examplebox}

이러한 규제 완화는 또 다른 통합과 재편성을 촉발했습니다. 여기에는 JP모건 체이스, 뱅크 오브 아메리카, 웰스 파고와 같은 세계 최대 은행들을 탄생시킨 대규모 합병(Mega Mergers)의 물결이 포함되었습니다.

\begin{figure}[h!]
    \centering
    % \includegraphics[width=0.8\textwidth]{example_graph_us_banks_final.png}  % Image not found: example_graph_us_banks_final.png % 가상의 이미지 파일명
    \caption{미국 내 은행 수 변화 (최종 시계열)}
    \label{fig:us_banks_final}
    \textit{이 차트는 미국 은행 수를 마지막으로 보여줍니다. 시계열의 끝 부분에서 지점형 은행(Branched Banks)이 이제 지배적임을 알 수 있습니다.}
\end{figure}

\begin{figure}[h!]
    \centering
    % \includegraphics[width=0.8\textwidth]{example_map_bank_of_america_2010.png}  % Image not found: example_map_bank_of_america_2010.png % 가상의 이미지 파일명
    \caption{뱅크 오브 아메리카의 2010년 지리적 분포}
    \label{fig:bank_of_america_2010}
    \textit{마지막으로, 2010년 뱅크 오브 아메리카의 지리적 분포를 보여줍니다. 각 점은 은행 지점을 나타냅니다. 가장 큰 은행들은 이제 미국 거의 모든 곳에서 운영됩니다.}
\end{figure}

오늘날, 소수의 거대 은행들이 모든 종류의 금융 서비스를 한 지붕 아래에서 제공합니다. 이러한 기관들은 \textbf{금융 지주 회사(Financial Holding Companies)}로 조직되어 있습니다. 이들은 수천 개의 은행 및 비은행 자회사를 포함합니다.

\begin{examplebox}[title={JP모건 체이스의 현대적 조직 구조}]
JP모건 체이스를 다시 살펴보겠습니다. 주요 활동은 다음과 같습니다.
\begin{itemize}
    \item \textbf{소매 및 상업 은행업 (Retail and Commercial Banking):} 일반 개인 및 기업 고객 대상.
    \item \textbf{투자 은행업 (Investment Banking):} 기업의 자금 조달 및 인수합병 자문.
    \item \textbf{증권 회사 기능 (Securities Firm Functions):} 브로커-딜러 활동(Broker-Dealer Activities) 포함.
    \item \textbf{자산 관리 (Asset Management):} 뮤추얼 펀드(Mutual Funds) 및 헤지 펀드(Hedge Funds) 컬렉션을 운영.
\end{itemize}
이러한 새로운 조직 구조를 통해 현대 은행들은 \textbf{수수료 수입(Fee Generation)}에 더 큰 비중을 둡니다. 투자 은행과 증권 회사들은 제공하는 서비스로부터 수수료를 얻습니다. 과거의 3-6-3 은행 모델은 이제 옛말이 되었습니다.
\end{examplebox}

\begin{summarybox}[title={Lesson 2-1 요약}]
이번 세션에서는 20세기 초부터 현재에 이르기까지 은행업의 역사를 논의했습니다.
\begin{itemize}
    \item \textbf{초기 은행:} 단순했으며, 주로 예금 수취와 대출에 집중했고, 종종 한 곳의 물리적 위치에서 운영되었습니다. 규제는 엄격했습니다.
    \item \textbf{변화의 시작:} 시간이 지남에 따라 새로운 기술과 비은행 기관(Nonbanks)과의 경쟁이 전통적인 은행의 시장 지위를 약화시켰습니다.
    \textbf{규제 완화:} 은행의 현대화를 돕기 위해 규제 완화의 물결이 뒤따랐습니다. 은행들은 지리적으로 확장하고 새로운 제품 시장으로 진출하는 것이 허용되었습니다.
    \item \textbf{현대 은행의 탄생:} 이는 오늘날 우리가 보는 금융 지주 회사 구조와 집중된 은행 산업을 낳았습니다.
\end{itemize}
\textbf{질문:} 너무 많은 규제 완화가 있었을까요? 일부 거대 은행들은 '너무 커서 망할 수 없는(Too Big to Fail)' 존재가 되었을까요? 2008년 금융 위기에서 가장 큰 미국 은행들이 핵심적인 역할을 했다는 점은 그 답이 '그렇다'일 수 있음을 시사합니다. 이 복잡한 질문과 규제 당국이 이에 대해 무엇을 했는지는 나중에 다룰 것입니다.
\end{summarybox}


\section{Lesson 2-2: 현대 은행의 조직과 기능 (Organization and Functions of Modern Banks)}
\addcontentsline{toc}{section}{Lesson 2-2: 현대 은행의 조직과 기능 (Organization and Functions of Modern Banks)}

우리는 다른 어떤 금융기관보다 은행과 더 자주 상호작용할 것입니다. 은행은 광범위한 금융 서비스를 제공합니다.
\begin{itemize}
    \item 우리의 돈을 안전하게 보관합니다.
    \item 결제 시스템에 대한 접근을 제공합니다.
    \item 금융 시장에 대한 접근을 제공합니다.
    \item 신용도 높은 차입자에게 대출을 해줍니다.
\end{itemize}
더 큰 은행들은 상업 은행 서비스와 투자 은행 서비스를 모두 한 지붕 아래에서 제공하는 '금융 슈퍼마켓'입니다. 이번 세션에서는 \textbf{상업 은행(Commercial Banks)}의 비즈니스 모델에 초점을 맞출 것입니다. 상업 은행이 어떻게 자금을 조달하고(Borrowing), 무엇을 하는지(Lending) 설명할 것입니다. 그 과정에서 \textbf{소매 금융(Retail Banking)}과 \textbf{도매 금융(Wholesale Banking)}의 차이점을 설명할 것입니다.

\subsection{Lesson 2-2.1: 은행 유형 및 자금 조달 (Bank Types and Borrowing)}
\addcontentsline{toc}{subsection}{Lesson 2-2.1: 은행 유형 및 자금 조달 (Bank Types and Borrowing)}

먼저 상업 은행이 무엇인지 정의해 봅시다.
\begin{summarybox}[title={상업 은행의 정의}]
\textbf{상업 은행(Commercial Bank)}은 예금을 수취하고 소비자 대출, 부동산 대출, 기업 대출을 제공하는 금융기관입니다. 역사적으로 상업 은행은 기업 고객을 대상으로 했기 때문에 '상업(Commercial)'이라는 이름이 붙었지만, 요즘에는 가계와 개인에게도 서비스를 제공합니다.
\end{summarybox}

우리는 일반적으로 상업 은행을 규모와 기능에 따라 분류합니다.

\subsubsection*{상업 은행의 유형}
\begin{enumerate}
    \item \textbf{커뮤니티 은행 (Community Banks):}
    \begin{itemize}
        \item \textbf{특징:} 지역 사회의 소비자와 소규모 기업에 서비스를 제공하는 소규모 은행입니다.
        \item \textbf{운영 방식:} 개인 고객으로부터 예금을 받아 지역 사회의 개인 및 소규모 기업에 대출을 해줍니다. 이를 \textbf{소매 금융(Retail Banking)}이라고 합니다.
        \item \textbf{규모:} 주로 10억 달러 미만의 자산(대부분 대출)을 관리합니다.
    \end{itemize}

    \item \textbf{지역/초지역 은행 (Regional/Super-regional Banks):}
    \begin{itemize}
        \item \textbf{특징:} 커뮤니티 은행보다 규모가 크며, 수백억 달러에서 1조 달러 이상의 자산을 가집니다. 대도시에 본사를 두고 전국적으로 사무실과 지점을 운영할 수 있습니다.
        \item \textbf{조직 형태:} 일부는 여러 은행 및 비은행 자회사를 가진 금융 지주 회사로 조직되어 있습니다.
        \item \textbf{대출 활동:} 상업 은행 부문은 더 큰 규모로 운영됩니다. 소매 대출뿐만 아니라 대기업 고객을 대상으로 하는 상업 및 산업 대출(Commercial and Industrial Lending)과 같은 \textbf{도매 대출(Wholesale Lending)}도 합니다.
        \item \textbf{자금 조달:} 지점 네트워크를 통해 개인 예금자로부터 자금을 조달하는 것 외에도 \textbf{도매 자금 조달 시장(Wholesale Funding Markets)}에 접근할 수 있습니다. (이 빌린 자금에 대해서는 잠시 후에 더 자세히 설명하겠습니다.)
        \item \textbf{예시:} 웰스 파고(Wells Fargo)는 메가 은행이지만, 전국에서 가장 큰 지점 네트워크를 가지고 있어 예금 자금을 쉽게 확보할 수 있습니다. 따라서 웰스 파고는 초지역 은행으로 분류됩니다.
    \end{itemize}

    \item \textbf{머니 센터 은행 (Money Center Banks):}
    \begin{itemize}
        \item \textbf{특징:} 대출 활동을 지원하기 위해 \textbf{도매 자금 조달(Wholesale Funding)}에 크게 의존하는 메가 은행입니다.
        \item \textbf{위치:} 미국 기반의 머니 센터 은행은 주로 뉴욕시에 위치하며, 은행 간 대출(Interbank Loan) 및 통화 교환 시장의 중심에 있습니다.
        \item \textbf{예시:} 씨티그룹(Citigroup)은 머니 센터 은행입니다.
    \end{itemize}
\end{enumerate}

\subsubsection*{은행의 자금 조달 (Borrowing)}
상업 은행은 돈을 빌리고(Borrow), 빌려주며(Lend), 그 차이(Spread)를 벌어들입니다. 돈을 빌리고 그에 대한 비용을 지불한 다음, 돈을 빌려주고 그 투자에서 이자를 벌어들입니다. 이것이 실제로 어떤 의미일까요? 은행은 어떻게 자금을 조달하고, 어떻게 대출을 해줄까요?

먼저 \textbf{자금 조달(Funding)} 측면부터 살펴보겠습니다. 은행은 저축자로부터 자금을 얻습니다. 개인과 기업으로부터 예금을 유치합니다. 예금자들이 은행에 돈을 빌려주도록 설득하기 위해 은행은 안전성과 편리함을 제공합니다. 은행은 또한 금융 시장을 통해 기업 및 다른 금융기관으로부터 자금을 빌릴 수 있습니다. 이는 일반적으로 이자 지급을 필요로 합니다.

\begin{summarybox}[title={은행의 주요 자금 조달원}]
\begin{enumerate}
    \item \textbf{예금 계좌 (Deposit Accounts):}
    \begin{itemize}
        \item \textbf{거래 계좌 (Transaction Accounts):}
        \begin{itemize}
            \item \textbf{요구불 예금 (Demand Deposits):} 우리의 당좌 예금 계좌(Checking Accounts)를 포함합니다. 이자는 지급되지 않지만, 원할 때 언제든지 자금을 인출할 수 있습니다.
        \end{itemize}
        \item \textbf{비거래 계좌 (Non-transaction Deposits):}
        \begin{itemize}
            \item \textbf{저축 계좌 (Savings Accounts) 및 소액 양도성 예금 증서 (Small Certificates of Deposit, CD):} 일반적으로 이자를 제공하는 저축 상품입니다. 대신 이러한 예금은 접근하기가 조금 더 어렵습니다. 한 달에 제한된 횟수만 인출할 수 있거나, CD의 경우 조기 해지 시 상당한 벌금이 부과될 수 있습니다.
        \end{itemize}
    \end{itemize}

    \item \textbf{차입 (Borrowing):}
    특히 대형 은행의 경우, 차입은 두 번째 주요 자금원입니다.
    \begin{itemize}
        \item \textbf{은행 간 시장 (Interbank Market) / 연방 기금 시장 (Federal Funds Market):} 자금 부족을 겪는 한 은행이 다른 은행으로부터 자금을 빌릴 수 있는 시장입니다. 이는 단기 무담보 대출(Short Term Unsecured Loans)입니다.
        \item \textbf{환매 조건부 채권 (Repurchase Agreement, Repo):} \textbf{담보부(Secured Basis)}로 자금을 빌릴 수 있는 금융 상품입니다. 레포는 전당포와 매우 유사하게 작동합니다.
        \begin{examplebox}[title={레포(Repo)의 전당포 비유}]
        제가 다음 일리노이 농구 경기 티켓을 사기 위해 몇백 달러가 정말 필요하다고 가정해 봅시다. 저는 골프채를 전당포에 가져갈 수 있습니다. 우리 모두 이 골프채가 500달러 가치가 있다는 데 동의합니다. 전당포 주인은 골프채를 받고 저에게 400달러를 줄 것입니다. 이것은 골프채를 담보로 한 400달러 대출입니다. 그리고 다음 주에 월급을 받으면 500달러를 주고 골프채를 다시 살 수 있습니다. 만약 제가 다음 주에 나타나지 않으면, 전당포 주인은 제 골프채를 가지게 됩니다. 이것이 기본적으로 환매 조건부 채권(Repo)이 작동하는 방식입니다.
        \end{examplebox}
        은행이 레포를 사용할 때는 연금 기금(Pension Funds)과 같은 현금이 풍부한 기관 투자자로부터 하룻밤 동안 현금을 빌리기 위해 국채(Treasury Bills)와 같은 고품질 증권을 담보로 제공합니다.

        \textbf{레포의 문제점:} 레포와 다른 단기 차입 자금은 투자자들이 걱정하기 시작하면 사라지는 경향이 있습니다. 2007년 공황은 레포 시장의 뱅크 런이었다고 유명하게 관찰되었습니다. 미상환 레포의 규모가 거의 하룻밤 사이에 붕괴되었고, 일부 은행들은 정말 중요한 자금원을 잃었습니다. 리먼 브라더스(Lehman Brothers)가 대표적인 예입니다.
    \end{itemize}

    \item \textbf{은행 자본 (Bank Capital):}
    자금 조달 퍼즐의 마지막 조각은 은행 자본입니다. 이는 은행에 대한 소유주의 투자입니다. 주로 보통주(Common Stock)와 우선주(Preferred Stock)로 구성되며, 다른 사업과 마찬가지로 유보 이익(Retained Earnings), 즉 은행에 재투자된 이익을 포함합니다. 부채와 달리 은행 자본은 만기가 없으므로 가장 안정적인(Sticky) 자금원입니다. 그러나 주식은 가격 위험에 노출되어 있기 때문에 예금보다 조달 비용이 더 비쌉니다.
\end{enumerate}
\end{summarybox}

\begin{table}[h!]
    \centering
    \caption{은행 자금 조달원의 비용 및 안정성 비교}
    \label{tab:funding_comparison}
    \begin{adjustbox}{width=\textwidth,center}
    \begin{tabular}{p{0.25\textwidth}p{0.35\textwidth}p{0.35\textwidth}}
        \toprule
        \textbf{자금 조달원} & \textbf{비용 (Cost)} & \textbf{안정성 (Stickiness)} \\
        \midrule
        \textbf{소매 예금 (Retail Deposits)} &
        \textbullet~대규모 인프라 및 고정 비용 필요 \\
\textbullet~한계 비용은 저렴 (Cheap on the margin) &
        \textbullet~정부 보증 예금 보험 덕분에 공황 시에도 은행 시스템 내에 머무는 경향이 있음 (Sticky) \\
        \addlinespace
        \textbf{도매 자금 (Wholesale Funding)} &
        \textbullet~경쟁력 있는 이자율을 지불해야 함 &
        \textbullet~상황이 나쁠 때 빠르게 사라지는 경향이 있음 (Runs away) \\
        \addlinespace
        \textbf{은행 자본 (Bank Capital)} &
        \textbullet~주식은 가격 위험에 노출되어 있어 예금보다 조달 비용이 비쌈 &
        \textbullet~만기가 없어 가장 안정적 (Never matures, as sticky as it gets) \\
        \bottomrule
    \end{tabular}
    \end{adjustbox}
\end{table}


\subsection{Lesson 2-2.2: 은행 대출 (Bank Lending)}
\addcontentsline{toc}{subsection}{Lesson 2-2.2: 은행 대출 (Bank Lending)}

이제 은행 대출로 넘어가겠습니다. 논의를 세 부분으로 나누겠습니다.
\begin{enumerate}
    \item \textbf{대출의 정의:} 대출이 무엇인지 설명합니다.
    \item \textbf{고객 유형:} 은행이 대출 측면에서 어떤 고객을 다루는지 살펴봅니다.
    \item \textbf{대출 과정의 메커니즘:} 대출이 어떻게 발생하는지 설명합니다.
\end{enumerate}

\subsubsection*{대출의 정의와 조건}
공식적으로 대출이란 무엇일까요? 대출은 \textbf{금융 상품(Financial Instrument)}이며, \textbf{부채(Debt)}입니다. 제가 은행에서 돈을 빌릴 때, 저는 미리 정해진 계약 조건에 따라 상환하겠다는 법적 구속력이 있는 약속을 하는 것입니다.

\begin{itemize}
    \item \textbf{이자율 (Interest Rate):} 변동 금리(Variable Rate)일 수도 있고 고정 금리(Fixed Rate)일 수도 있습니다.
    \item \textbf{담보 (Collateral):} 대출은 어떤 담보로 보증될 수 있습니다.
    \item \textbf{약정 (Covenants):} 계약서에 성과 관련 조항이 명시될 수 있습니다.
    \item \textbf{선순위 채무 (Senior Debt):} 파산 시 사업 자산에 대한 최우선권을 가집니다.
\end{itemize}

은행은 일반적으로 대출 사업을 두 가지 범주로 나눕니다.
\begin{enumerate}
    \item \textbf{소비자 신용 (Consumer Credits):}
    \begin{itemize}
        \item 주거용 부동산 담보 대출: 주택 담보 대출(Mortgages) 및 주택 담보 대출(Home Equity Loans)을 포함합니다.
        \item 신용 카드(Credit Cards), 자동차 대출(Auto Loans), 학자금 대출(Student Loans)도 포함됩니다.
    \end{itemize}
    \item \textbf{상업 및 산업 대출 (Commercial and Industrial Lending):}
    \begin{itemize}
        \item 기업에 대한 대출을 포함합니다.
        \item 상업용 부동산 대출(Commercial Real Estate Loans), 투자 목적의 장기 대출인 기간 대출(Term Loans), 기업의 운전자금(Working Capital)으로 사용될 수 있는 사업 신용 한도(Business Lines of Credit) 등이 있습니다.
    \end{itemize}
\end{enumerate}

\subsubsection*{상업 은행의 고객 기반}
은행은 누구에게 대출을 해줄까요? 상업 은행의 고객 기반을 살펴보겠습니다.

\begin{enumerate}
    \item \textbf{대기업 (Larger Businesses):}
    \begin{itemize}
        \item 상업 은행을 우회하여 금융 시장에서 직접 자금을 조달하는 경향이 있습니다. 채권을 발행하고 주식 자금을 조달합니다.
        \item 이러한 기업에 대한 공개 정보가 많고 정보 비용이 적습니다.
        \item 때로는 대기업이 특별한 상황(예: 인수 거래 또는 부실 채무 구조 조정)에 있을 때 은행을 통해 대출을 주선하기도 합니다.
    \end{itemize}

    \item \textbf{중견 기업 대출 (Middle Market Lending):}
    \begin{itemize}
        \item 상업 은행 사업의 핵심입니다. 때로는 \textbf{관계형 대출(Relationship Lending)}이라고도 불립니다.
        \item 공개 정보가 제한적인 중견 기업들입니다. 아마도 공개 증권 거래소에 상장되어 있지 않을 것이며, 투자자들은 역선택(Adverse Selection)을 우려할 수 있습니다.
        \item 사업 고객을 찾고, 신용 분석을 수행하고, 새로운 대출에 자금을 조달하고, 관계를 유지하는 것이 대부분 은행의 대출 담당자(Loan Officers)와 신용 위원회(Credit Committees)의 핵심 역할입니다.
    \end{itemize}

    \item \textbf{소규모 및 신생 기업 (Small and New Businesses):}
    \begin{itemize}
        \item 일반적으로 어떤 형태의 보증(Guarantee) 없이는 은행 신용에 접근하기 어렵습니다.
        \item 사업주 개인의 보증, 중소기업청(Small Business Administration, SBA)의 부분 보증, 또는 사업주가 개인 자산(주택, 보트, 기타 금융 자산)을 담보로 제공할 수 있습니다.
        \item 소규모 기업은 위험하고 종종 실패합니다. 보증과 담보는 위험을 줄이고 은행이 대출에 참여하도록 돕습니다.
    \end{itemize}

    \item \textbf{개인 대출 (Loans to Individuals):}
    \begin{itemize}
        \item \textbf{주택 담보 대출 (Mortgage Credit):} 커뮤니티 은행과 지역 은행의 핵심 초점입니다. 이러한 대출은 정보 민감성이 높으며 차입자의 개별 상황에 따라 달라질 수 있습니다.
        \item \textbf{신용 카드, 학자금 대출, 자동차 대출:} 건당 대출 규모가 작고 비용이 많이 듭니다. 이 시장에는 큰 규모의 경제(Economies of Scale)가 존재하므로, 이러한 대출을 대규모로 다각화된 풀(Pool)로 만들 수 있는 대형 은행과 비은행 기관이 이 분야를 지배하는 경향이 있습니다.
    \end{itemize}
\end{enumerate}

\begin{figure}[h!]
    \centering
    % \includegraphics[width=0.8\textwidth]{example_jpmorgan_segments.png}  % Image not found: example_jpmorgan_segments.png % 가상의 이미지 파일명
    \caption{JP모건 체이스의 사업 부문별 고객 분류}
    \label{fig:jpmorgan_segments}
    \textit{JP모건 체이스의 사업 부문별 고객 분류를 통해 이러한 다양한 고객들이 어떻게 나타나는지 볼 수 있습니다.}
\end{figure}

\begin{itemize}
    \item \textbf{소비자 및 커뮤니티 금융 (Consumer and Community Banking):} 제가 \textbf{소매 금융(Retail Banking)}이라고 부르는 부분입니다.
    \begin{itemize}
        \item \textbf{소비자 및 사업 금융 (Consumer and Business Banking):} 지점 네트워크를 통해 예금 수취, 소비자 신용 및 소규모 사업 대출을 제공합니다.
        \item \textbf{주택 대출 (Home Lending):} 주거용 주택 담보 대출 및 주택 담보 대출을 담당합니다.
        \item \textbf{카드 및 자동차 (Card and Auto):} 신용 카드와 자동차 대출 및 리스를 처리합니다.
        \item 이 모든 것은 개인 및 소규모 기업을 위한 것입니다.
    \end{itemize}

    \item \textbf{상업 금융 (Commercial Banking):} \textbf{도매 금융(Wholesale Banking)} 사업 내에 있습니다.
    \begin{itemize}
        \item \textbf{상업용 부동산 (Commercial Real Estate):} 사무 공간, 부동산 개발, 산업 공간 등에 대한 자금 조달을 포함합니다.
        \item \textbf{사업 금융 (Business Banking):} JP모건은 고객을 두 그룹으로 나눕니다.
        \begin{itemize}
            \item \textbf{중견 기업 (Middle-Market):} 연간 매출액이 2억 달러에서 5억 달러 사이의 중소기업을 대상으로 합니다. 이보다 작으면 커뮤니티 뱅킹 범주에 속합니다.
            \item \textbf{기업 고객 (Corporate Clients):} 연간 매출액이 5억 달러에서 20억 달러 사이의 대기업을 대상으로 합니다. 이보다 크면 투자 은행 부문에서 처리합니다.
        \end{itemize}
    \end{itemize}
\end{itemize}

\begin{figure}[h!]
    \centering
    % \includegraphics[width=0.8\textwidth]{example_jpmorgan_diagram.png}  % Image not found: example_jpmorgan_diagram.png % 가상의 이미지 파일명
    \caption{JP모건의 비즈니스 모델: 도매 및 소매 금융 매핑}
    \label{fig:jpmorgan_diagram}
    \textit{JP모건의 비즈니스 모델을 이전 다이어그램에 매핑할 수 있습니다. 여기는 중견 기업에 대한 대출을 다루는 도매 사업이고, 여기는 소규모 기업과 개인에 대한 대출을 다루는 소매 사업입니다.}
\end{figure}

\begin{summarybox}[title={Lesson 2-2.2 요약}]
이것이 상업 은행 비즈니스 모델의 개요입니다. 은행은 다양한 출처로부터 자금을 조달하고 광범위한 대출을 제공합니다.
\end{summarybox}


\subsection{Lesson 2-2.3: 대출 과정의 변화 (Changes in the Lending Process)}
\addcontentsline{toc}{subsection}{Lesson 2-2.3: 대출 과정의 변화 (Changes in the Lending Process)}

논의를 마무리하기 전에 두 가지를 더 다루고자 합니다. 첫째, 대출 과정 자체에 대해 좀 더 정확하게 설명하고 싶습니다. 둘째, 지난 50년 동안 이 대출 과정이 어떻게 변했는지 설명하고 싶습니다.

\subsubsection*{대출 과정의 세 가지 필수 단계}
대출 과정에는 세 가지 필수 단계가 있습니다.
\begin{enumerate}
    \item \textbf{개시 (Origination):}
    \begin{itemize}
        \item 대출을 승인하는 단계입니다. 대출 담당자(Loan Officer)가 고객과 만나 신용도를 판단합니다.
        \item 이는 차입자의 재정 상황이나 사업 계획에 대한 신중한 분석을 포함합니다.
        \item 대출 담당자는 대출 금액, 만기, 이자율 및 수수료, 담보 등 조건을 협상합니다.
        \item 우리는 이 단계를 이전 논의에서 \textbf{스크리닝(Screening)}이라고 불렀습니다.
    \end{itemize}

    \item \textbf{자금 조달 (Financing):}
    \begin{itemize}
        \item 대출을 약정한 후, 은행은 차입자에게 실제로 제공할 자금을 마련해야 합니다.
    \end{itemize}

    \item \textbf{관리 (Servicing):}
    \begin{itemize}
        \item 이 단계는 대출이 이루어진 후 은행이 차입자와의 관계를 지속하는 것을 포함합니다.
        \item 은행은 차입자로부터 상환금을 수취합니다.
        \item 은행은 차입자의 진행 상황과 재정 상태를 모니터링하여 은행이 전액 상환받을 가능성을 극대화합니다.
    \end{itemize}
\end{enumerate}

\subsubsection*{전통적인 모델: 개시 및 보유 (Originate and Hold)}
전통적으로 은행은 이 세 단계를 모두 수행했습니다.
\begin{itemize}
    \item 대출을 승인하고,
    \item 예금을 사용하여 대출에 자금을 조달하며,
    \item 대출이 만기될 때까지 장부에 보유하면서 상환금을 수취했습니다.
\end{itemize}
이것이 우리가 \textbf{개시 및 보유 모델(Originate and Hold Model)}이라고 부르는 것입니다.

\subsubsection*{현대적인 모델: 개시 및 분배 (Originate and Distribute)}
이제 재미있는 부분입니다. 현대 시대에는 이 세 단계 각각이 \textbf{개시 및 분배 모델(Originate and Distribute Model)}이라고 불리는 방식으로 분리될 수 있습니다. 이 비즈니스 모델 하에서 상업 은행은 \textbf{수수료 수입(Fee Income)}에 중점을 둡니다.

\begin{itemize}
    \item \textbf{개시 (Origination):}
    \begin{itemize}
        \item 은행은 여전히 차입자와 만나 대출을 개시합니다. 이 부분은 변하지 않았지만, 때로는 브로커(Broker)의 도움을 받기도 합니다.
        \item 주택 담보 대출 시장에서는 온라인 브로커가 수수료를 받고 대출 기관과 고객을 연결해주는 경우를 접했을 수 있습니다.
    \end{itemize}

    \item \textbf{자금 조달 (Financing):}
    \begin{itemize}
        \item 다른 금융기관이 주택 구매 자금을 제공합니다. 이는 보험 회사와 같은 투자자일 수도 있고, \textbf{증권화(Securitization)}라는 추가 단계가 있을 수도 있습니다.
    \end{itemize}

    \item \textbf{관리 (Servicing):}
    \begin{itemize}
        \item 은행이 관리를 수행할 수도 있지만, 이마저도 때로는 아웃소싱됩니다.
        \item 매월 상환금을 은행에 지불하는 대신, 주택 담보 대출 관리자(Mortgage Servicer)에게 지불하고, 관리자는 그 상환금을 투자자에게 전달합니다.
    \end{itemize}
\end{itemize}

\begin{center}
    \textit{왜 이렇게 할까요?}
\end{center}

일부 은행이 정보 비용을 최소화하고 최적의 차입자를 승인하는 데 정말 뛰어나다면, \textbf{비교 우위(Comparative Advantage)} 이론은 그들이 그 활동에 집중해야 한다고 말합니다. 그리고 다른 파트너가 자금을 제공할 수 있습니다. 이것이 바로 \textbf{개시 및 분배 모델(Originate and Distribute Model)}의 핵심입니다.

\subsubsection*{핀테크(Fintech)의 등장과 대출 시장의 변화}
마지막으로 한 가지 더 생각할 점이 있습니다. 이 이론은 21세기에 접어들면서 무너질 수도 있습니다. 요즘에는 로켓 모기지(Rocket Mortgage)와 같은 \textbf{핀테크(Fintech)} 기업들이 주택 담보 대출 시장을 뒤흔들었습니다. 소비자로서 제가 주택 담보 대출을 재융자하거나 특정 품질 기준을 충족하는 우량 차입자라면, 온라인 대출 기관으로부터 빠르고 편리하게 대출을 받을 수 있습니다. 이런 상황에서 은행이 비교 우위를 가질 수 있을까요?

우리는 아마존(Amazon)과 같은 다른 기술 기업들에 의한 기업 대출 시장에서도 유사한 진입을 보고 있습니다. 이들 기업은 고객의 관심과 사업 매출에 대한 방대한 양의 정보를 실시간으로 수집합니다. 아마존이 어떤 차입자에게 대출을 승인해야 할지 결정하는 데 더 나은 위치에 있을 수도 있습니다. 우리는 지켜봐야 할 것입니다.

\begin{summarybox}[title={Lesson 2-2.3 요약}]
이번 세션에서는 상업 은행 사업을 면밀히 살펴보았습니다.
\begin{itemize}
    \item 은행이 어떻게 자금을 조달하고 대출하는지 설명했습니다.
    \item 소매 금융과 도매 금융의 차이점을 설명했습니다.
    \item 마지막으로, 은행이 최적의 차입자를 선택하는 데 비교 우위를 가진다면 합리적인 \textbf{개시 및 분배 모델(Originate and Distribute Model)}을 소개했습니다.
\end{itemize}
\end{summarybox}


\section{개요}
이 문서는 'Banking and Financial Institutions Module 2'의 핵심 내용을 다루며, 투자은행의 조직 구조, 주요 활동, 그리고 현대 은행의 재무제표 분석에 초점을 맞춥니다. 투자은행이 기업과 금융시장을 연결하는 방식, 주요 수익원, 그리고 상업은행과의 통합 과정을 이해합니다. 또한, 상업은행의 대차대조표(자산 및 부채)를 상세히 살펴보고, 대차대조표 외 활동(Off-Balance Sheet Activities)의 중요성과 위험성을 분석하여 금융기관의 복잡한 운영 방식을 종합적으로 파악하는 것을 목표로 합니다.


\section{용어 정리}

\begin{adjustbox}{width=\textwidth,center}
\begin{tabular}{llll}
\toprule
\textbf{용어} & \textbf{쉬운 설명} & \textbf{원어} & \textbf{비고} \\
\midrule
투자은행 & 기업이 주식이나 채권을 발행하여 자금을 조달하도록 돕는 금융기관 & Investment Bank & 상업은행과 달리 대출보다는 자본시장 접근 지원 \\
상업은행 & 예금을 받아 대출을 해주는 전통적인 은행 & Commercial Bank & 일반 개인 및 기업 대상 \\
고액 자산가 & 상당한 금융 자산을 보유한 개인 & High-Net-Worth Individual & 투자은행의 주요 고객 중 하나 \\
증권 인수 & 기업이 발행하는 주식이나 채권을 투자은행이 매입하여 시장에 판매하는 과정 & Underwriting & 투자은행의 핵심 수익원 \\
증권 거래 & 주식, 채권 등 증권을 사고파는 활동 & Securities Trading & 고객 주문 처리 및 자기 계정 거래 포함 \\
M\&A 자문 & 기업의 인수합병(Mergers and Acquisitions) 또는 구조조정에 대한 자문 서비스 & M\&A Advisory & 기업 전략 및 가치 평가 전문성 요구 \\
신디케이트론 & 여러 은행이 공동으로 대규모 대출을 제공하여 위험을 분산하는 방식 & Syndicated Loan & 주로 대규모 기업 대출에 활용 \\
북빌딩 & 투자은행이 신규 증권 발행 시 투자자들의 수요를 파악하고 가격을 결정하는 과정 & Book-building & IPO 과정의 핵심 단계 \\
비드-애스크 스프레드 & 매수호가(Bid Price)와 매도호가(Ask Price)의 차이 & Bid-Ask Spread & 시장 조성자(Market Maker)의 주요 수익원 \\
시장 조성 & 투자자들이 증권을 쉽게 사고팔 수 있도록 매수호가와 매도호가를 제시하는 활동 & Market-making & 유동성 공급 역할 \\
자기 계정 거래 & 투자은행이 고객을 위해서가 아니라 자체 이익을 위해 증권을 사고파는 활동 & Proprietary Trading & 높은 수익 가능성, 높은 위험 수반 \\
대차대조표 & 특정 시점의 기업 자산, 부채, 자본 상태를 보여주는 재무 보고서 & Balance Sheet & 은행의 재무 건전성 파악에 필수적 \\
뱅킹 북 & 전통적인 대출 및 예금 활동과 관련된 은행의 자산 및 부채 & Banking Book & 장기적인 이자 수익 및 비용 발생 \\
트레이딩 북 & 단기적인 증권, 상품, 통화 거래와 관련된 은행의 자산 및 부채 & Trading Book & 단기 매매 차익을 목표로 함 \\
대차대조표 외 활동 & 현재는 자산이나 부채로 기록되지 않지만, 특정 사건 발생 시 재무제표에 영향을 미칠 수 있는 활동 & Off-Balance Sheet (OBS) Activities & 신용카드 약정, 파생상품 등이 대표적 \\
우발 자산/부채 & 특정 미래 사건의 발생 여부에 따라 자산이나 부채가 될 수 있는 항목 & Contingent Asset/Liability & OBS 활동의 핵심 개념 \\
미사용 대출 약정 & 은행이 고객에게 약속했지만 아직 인출되지 않은 대출 한도 & Undrawn Loan Commitment & 신용카드 한도 등이 해당 \\
파생상품 & 기초자산의 가치 변동에 따라 가치가 결정되는 금융상품 & Derivatives & 위험 헤지 또는 투기 목적으로 사용 \\
명목 금액 & 파생상품 계약의 기초자산 총액 & Notional Amount & 실제 시장 가치와는 다름 \\
신용부도스왑 & 채권이나 대출의 부도 위험에 대한 보험 계약 & Credit Default Swap (CDS) & 신용 위험을 전가하는 수단 \\
\bottomrule
\end{tabular}
\end{adjustbox}


\section{투자은행의 조직 구조와 활동}

\subsection{투자은행이란 무엇인가?}
투자은행은 기업이나 다른 기관들이 금융시장에 접근할 수 있도록 돕는 금융기관입니다.
\begin{itemize}
    \item \textbf{상업은행과의 차이:} 일반적인 기업이 대출을 원하면 상업은행(Commercial Bank)에 가지만, 주식이나 채권을 발행하여 자금을 조달하고 싶다면 투자은행(Investment Bank)을 찾습니다.
    \item \textbf{고액 자산가 서비스:} 고액 자산가(High-Net-Worth Individual)가 주식 거래를 실행하고 싶을 때도 투자은행가에게 연락할 수 있습니다.
\end{itemize}
이러한 구분은 단순화된 설명이지만, 투자은행의 핵심 역할을 이해하는 데 도움이 됩니다. 투자은행의 목표는 차입자(Borrower)와 저축자(Saver)를 금융시장에서 직접 연결하는 것입니다.

\begin{summarybox}
\textbf{핵심 요약:} 투자은행은 기업의 자금 조달(주식/채권 발행) 및 M\&A를 돕고, 증권 거래 서비스를 제공하는 금융기관입니다. 상업은행이 대출과 예금에 집중하는 것과 달리, 투자은행은 자본시장을 통한 자금 조달과 투자에 특화되어 있습니다.
\end{summarybox}

\subsection{투자은행의 세 가지 주요 사업 분야}
투자은행은 크게 세 가지 주요 사업 라인을 가지고 있습니다. 이 활동들은 모두 수수료를 발생시키며, 상업은행과 달리 장기 자금 조달이나 광범위한 지점 네트워크가 필요하지 않아 매우 수익성이 높은 사업이 될 수 있습니다.

\begin{enumerate}
    \item \textbf{자금 조달 (Financing):}
    \begin{itemize}
        \item \textbf{정의:} 투자은행은 고객(주로 대기업 및 정부)을 대신하여 증권을 발행(Originate), 인수(Underwrite), 그리고 금융시장에 배치(Place)합니다.
        \item \textbf{예시:} 기업이 새로운 공장을 짓기 위해 자금이 필요할 때, 투자은행은 해당 기업의 주식이나 채권을 발행하여 투자자들에게 판매하는 과정을 총괄합니다.
        \item \textbf{수익원:} 증권 인수 수수료(Underwriting Fees)가 발생합니다.
    \end{itemize}

    \item \textbf{증권 거래 (Securities Trading):}
    \begin{itemize}
        \item \textbf{정의:} 증권 발행 및 판매 사업에 이미 참여하고 있기 때문에, 대부분의 투자은행은 증권 거래에도 관여합니다. 고객의 매수 및 매도 주문을 실행하고, 종종 자체 이익을 위해 거래하기도 합니다.
        \textbf{예시:} 고객이 특정 주식을 사고 싶을 때, 투자은행은 고객을 대신하여 시장에서 주식을 매수하고 수수료를 받습니다. 또한, 은행 자체적으로 주식이나 채권을 사고팔아 이익을 얻기도 합니다 (자기 계정 거래).
        \item \textbf{수익원:} 거래 실행에 대한 수수료(Commissions)가 발생합니다.
    \end{itemize}

    \item \textbf{M\&A 및 기업 구조조정 자문 (M\&A Advisory and Corporate Restructuring):}
    \begin{itemize}
        \item \textbf{정의:} 투자은행은 기업의 인수합병(Mergers and Acquisitions, M\&A) 또는 기타 기업 구조조정 활동을 촉진하는 데 도움을 줍니다.
        \item \textbf{예시:} 한 기업이 다른 기업을 인수하려고 할 때, 투자은행은 인수 대상 기업의 가치를 평가하고, 협상을 돕고, 필요한 자금 조달 방안을 제시합니다.
        \item \textbf{수익원:} 자문 서비스에 대한 수수료(Advisory Fees)가 발생합니다.
    \end{itemize}
\end{enumerate}

\subsection{투자은행의 기업 구조: JP Morgan Chase 사례}
JP Morgan Chase의 사례를 통해 투자은행의 기업 구조를 살펴보겠습니다.
\begin{itemize}
    \item \textbf{도매 사업 부문 (Wholesale Business Arm):} JP Morgan의 기업 및 투자은행 부문은 은행의 도매 사업 부문 내에 위치합니다. 이는 일반 소매 고객(Retail Customers)이 아닌 대규모 고객(Big Clients)을 주로 상대한다는 의미입니다.
    \item \textbf{두 가지 핵심 사업:} 투자은행은 크게 '뱅킹(Banking)'과 '마켓(Markets)' 두 가지 사업으로 구성됩니다.
    \begin{itemize}
        \item \textbf{뱅킹 (Banking):} 기업 전략 및 구조 자문, 자본 조달(주식, 채권 시장), 신디케이트론 시장 자문 등을 포함합니다.
        \item \textbf{마켓 (Markets):} 현금 및 파생상품 거래, 중개 서비스, M\&A 자문, 증권 인수, 증권 거래 등을 포괄합니다.
    \end{itemize}
\end{itemize}
JP Morgan은 2020년에 약 85억 달러의 투자은행 수수료를 벌어들였는데, 이는 연간 수익의 상당 부분을 차지합니다.

\subsection{투자은행 산업의 분류}
투자은행 산업에는 다양한 전문 기업들이 존재하지만, 가장 큰 풀서비스 투자은행은 세 가지 유형으로 분류할 수 있습니다.

\begin{enumerate}
    \item \textbf{상업은행 지주회사의 자회사 (예: JP Morgan):}
    \begin{itemize}
        \item \textbf{특징:} JP Morgan의 투자은행 부문은 상업은행 지주회사(Commercial Bank Holding Company)의 자회사입니다.
        \item \textbf{고객 기반:} 기업, 정부, 금융기관, 고액 자산가 등 광범위한 고객층을 보유한 가장 큰 투자은행들입니다.
    \end{itemize}

    \item \textbf{특정 고객 기반에 특화된 투자은행 (예: Goldman Sachs):}
    \begin{itemize}
        \item \textbf{특징:} Goldman Sachs와 같은 투자은행은 주로 기업 및 증권 거래 활동이 활발한 고객에 집중하는 경향이 있습니다.
        \item \textbf{고객 기반:} 보다 전문화된 고객 기반을 가집니다.
    \end{itemize}

    \item \textbf{독립형 투자은행 (Standalone Investment Banks, 예: Lazard):}
    \begin{itemize}
        \item \textbf{특징:} Lazard와 같은 독립형 투자은행은 모든 범위의 투자은행 서비스를 제공하지만, 지리적 범위가 제한적입니다. 주로 주요 도시에만 운영되며, 대규모 고객에 집중합니다.
        \item \textbf{역사적 변화:} 과거에는 Merrill Lynch, Bear Stearns, Lehman Brothers, Morgan Stanley 등 여러 독립형 투자은행이 있었으나, 규제 변화와 2008년 금융 위기로 인해 대부분 상업은행과 합병되거나 상업은행으로 재편되었습니다.
    \end{itemize}
\end{enumerate}

\begin{infobox}
\textbf{독립형 vs. 계열사:} '독립형(Standalone)'이란 상업은행과 제휴하지 않은 투자은행을 의미합니다. 과거에는 독립형이 많았으나, 1990년대 후반의 규제 완화(상업은행과 투자은행 활동 겸업 허용)와 2008년 금융 위기(대규모 인수합병 및 재편)로 인해 현재는 대부분 상업은행 지주회사 산하에 편입되어 있습니다.
\end{infobox}

\subsection{투자은행 산업의 변화와 통합}
투자은행 산업은 큰 변화를 겪었습니다.
\begin{itemize}
    \item \textbf{규제 변화 (1990년대 후반):} 1990년대 후반의 규제 변화로 상업은행과 투자은행 활동이 같은 회사에서 가능해지면서, 여러 상업은행이 자체 투자은행을 인수하거나 설립했습니다. Gramm-Leach-Bliley Act 통과와 함께 Citigroup이 Citicorp(상업은행)와 Travelers(투자은행을 소유한 보험회사)의 합병으로 탄생한 것이 대표적인 예입니다.
    \item \textbf{금융 위기 (2008년):} 2008년 금융 위기는 대규모 부실 인수(Distressed Acquisitions)의 물결을 일으켰습니다.
    \begin{itemize}
        \item Lehman Brothers는 파산했습니다.
        \item Bear Stearns와 Merrill Lynch는 상업은행에 인수되었습니다 (Bank of America가 Merrill Lynch를 인수한 것이 가장 큰 부실 인수였습니다).
        \item Goldman Sachs와 Morgan Stanley는 생존을 위해 상업은행으로 재편해야 했습니다.
    \end{itemize}
\end{itemize}

\begin{table}[h!]
\centering
\caption{미국 투자은행 산업의 주요 합병 사례}
\label{tab:us_ib_mergers}
\begin{adjustbox}{width=\textwidth,center}
\begin{tabular}{lll}
\toprule
\textbf{합병 주체} & \textbf{피합병 주체} & \textbf{합병 연도} \\
\midrule
Citigroup & Citicorp (상업은행) + Travelers (보험/투자은행) & 1998 \\
Chase & JPMorgan & 직후 \\
Bank of America & Merrill Lynch & 2008 \\
\bottomrule
\end{tabular}
\end{adjustbox}
\end{table}

\subsection{글로벌 투자은행 수익 현황 (2020년)}
2020년 투자은행 수익 분석에 따르면, 전 세계 상위 10개 투자은행 중 Jefferies만이 독립형 회사이며, 나머지는 모두 상업은행과 제휴하여 소매 고객으로부터 예금을 받습니다.
\begin{itemize}
    \item \textbf{시장 집중도:} 이 상위 10개 은행이 전 세계 수수료의 약 50\%를 차지하며, 시장 선두 주자는 약 10\%를 차지합니다.
    \item \textbf{경쟁 환경:} 상당한 통합이 있었음에도 불구하고, 글로벌 규모에서는 여전히 상당히 경쟁적인 시장입니다.
\end{itemize}


\section{투자은행 활동의 세부 사항}

\subsection{증권 인수 (Securities Underwriting)}
증권 인수는 투자은행 수수료의 대부분을 차지하는 핵심 활동입니다. 주로 주식, 채권, 신디케이트론 인수에 초점을 맞춥니다.

\begin{enumerate}
    \item \textbf{과정:}
    \begin{itemize}
        \item \textbf{의뢰 (Mandate):} 기업이 5억 달러와 같은 대규모 자금을 조달하기 위해 금융 상품(주식 또는 채권)을 판매하고자 할 때, 여러 투자은행이 이 인수 과정을 주도하기 위한 의뢰(Mandate)를 따내기 위해 경쟁합니다.
        \item \textbf{주간사 (Lead Underwriter):} 주간사로 선정된 투자은행은 거래 조건을 설정하고 새로 발행된 금융 상품의 판매를 관리하는 역할을 합니다.
    \end{itemize}

    \item \textbf{주식 발행 (Equity Sell):}
    \begin{itemize}
        \item \textbf{IPO (Initial Public Offering):} 기업이 처음으로 주식을 판매하는 경우를 기업공개(IPO)라고 합니다.
        \item \textbf{추가 발행 (Follow-on Deal/Seasoned Offering):} 이미 상장된 기업이 추가로 주식을 판매하는 경우입니다.
        \item \textbf{북빌딩 (Book-building):} 투자은행은 기업 및 잠재 투자자들과 만나 주식 가격과 매수 의사를 협상합니다. 이 과정을 북빌딩이라고 합니다.
        \item \textbf{인수 및 판매:} 투자은행은 전체 발행 물량을 보장된 가격에 매입하여 외부 투자자에게 더 높은 가격에 판매할 수 있습니다. 매입 가격과 판매 가격의 차이(Spread)가 투자은행의 이익이 됩니다.
        \item \textbf{예시:} 투자은행이 1억 주를 주당 5달러에 매입하여 주당 5.25달러에 판매하면, 주당 0.25달러의 스프레드로 2,500만 달러의 이익을 얻습니다.
        \item \textbf{위험:} 그러나 이러한 이익은 보장되지 않습니다. Morgan Stanley는 Facebook IPO 가격을 잘못 책정하여 주가가 몇 주 만에 거의 50\% 하락하는 사태를 겪었습니다. 이는 투자은행의 명성과 향후 사업 전망에 부정적인 영향을 미칠 수 있습니다. 따라서 인수 가격은 재고를 이익을 내고 소진할 수 있으면서도, 기업의 공정한 가치를 반영해야 합니다.
    \end{itemize}

    \item \textbf{채권 인수 (Bond Underwriting):} 주식 인수와 유사한 과정을 따릅니다.

    \item \textbf{신디케이트론 (Syndicated Loan):}
    \begin{itemize}
        \item \textbf{정의:} 일반적으로 상업은행이 기업 대출을 제공하지만, 대규모 대출의 경우 단일 은행이 감당하기에는 위험이 너무 큽니다. 이때 여러 은행과 대출 기관이 모여 공동으로 대출을 제공하여 위험을 분산합니다.
        \item \textbf{주간사 역할:} 투자은행(인수 은행)은 대출 신디케이트를 대표하여 대출 조건을 협상합니다. 대출 과정의 시작(Origination)과 사후 관리(Servicing)를 모두 담당합니다.
        \item \textbf{수익원:} 신디케이트를 주도하는 은행은 많은 수수료를 얻습니다. 다른 신디케이트 구성원들은 자금을 제공하는 '침묵하는 파트너' 역할을 하며, 대출 만기까지 일반적인 이자 수입을 얻습니다.
    \end{itemize}
\end{enumerate}

\subsection{M\&A 자문 (M\&A Advisory)}
투자은행은 기업의 인수합병(M\&A) 거래를 협상하고 실행하는 데 자주 관여합니다.

\begin{itemize}
    \item \textbf{역할:}
    \begin{itemize}
        \item \textbf{인수 대상 탐색:} 기업 X가 기업 Y를 인수하고 싶지만 적절한 가격을 모를 때, 투자은행은 잠재적 인수 대상에 대한 조사를 수행합니다.
        \item \textbf{매수자 탐색:} 기업 Z가 사업 부문을 매각하고 싶을 때, 투자은행은 잠재적 매수자를 찾는 데 도움을 줍니다.
        \item \textbf{가치 평가:} 재무 모델링 기법을 사용하여 상장 및 비상장 기업의 가치를 추정하는 전문가입니다.
        \item \textbf{자금 조달 지원:} M\&A 자문 은행은 거래 자금 조달을 위한 새로운 증권 인수(주식, 채권, 신디케이트론 등 모든 옵션 가능)도 지원할 수 있습니다.
    \end{itemize}
    \item \textbf{보상:} M\&A 거래의 세부 사항을 확정하는 것은 재정적, 법률적 관점에서 복잡하기 때문에, 투자은행은 제공하는 자문 및 서비스에 대해 높은 보상을 받습니다.
\end{itemize}

\subsection{투자은행 상품별 수수료 수익 (2020년)}
2020년 투자은행의 다양한 상품별 수수료 수익을 살펴보면 다음과 같습니다.
\begin{itemize}
    \item \textbf{주식 인수:} IPO(Initial Public Offering)와 추가 발행(Follow-on Offerings)으로 나뉩니다.
    \item \textbf{채권 인수:} 투자 등급(Investment Grade, 신용 등급이 좋은 기업)과 투자 등급 미만(Below Investment Grade, 신용 등급이 낮은 기업)으로 나뉩니다.
    \item \textbf{수익 규모:} 2020년에는 주식 인수, 채권 인수, M\&A 관련 활동에서 발생하는 총 수수료 수익이 모두 비슷한 규모였습니다. 신디케이트론에서 발생하는 수수료는 이들의 약 절반 수준이었습니다.
\end{itemize}


\subsection{증권 거래 (Securities Trading)}
증권 거래는 투자은행의 또 다른 핵심 활동이며, 시장 조성(Market-making)과 자기 계정 거래(Proprietary Trading)로 구성됩니다.

\subsubsection{시장 조성 (Market-making)}
\begin{itemize}
    \item \textbf{정의:} 투자은행이 투자자들이 증권을 거래할 수 있는 시장을 만드는 활동입니다. 많은 증권 회사들이 이 역할을 하지만, 투자은행이 중심적인 역할을 합니다.
    \item \textbf{브로커 역할:} 브로커는 고객을 대신하여 거래를 실행하고 수수료나 커미션을 받습니다.
    \item \textbf{딜러 역할:} 투자은행의 딜러 부문은 자체적으로 재고(Inventory)로 일부 주식을 보유할 수 있습니다. 이 경우 고객은 은행 자체로부터 주식을 매수할 수도 있습니다.
    \item \textbf{비드-애스크 스프레드 (Bid-Ask Spread):} 투자은행이 Facebook 주식을 매도호가(Ask Price) 20달러에 매수하여 고객에게 매수호가(Bid Price) 21달러에 판매한다면, 1달러의 이익을 얻습니다. 이 이익을 비드-애스크 스프레드라고 합니다.
    \item \textbf{시너지 효과:} 증권 딜링과 인수 활동을 결합하면 명확한 시너지 효과가 있습니다. 새로운 주식을 인수한 투자은행은 남은 주식을 재고로 보유하기 쉽습니다.
\end{itemize}

\subsubsection{자기 계정 거래 (Proprietary Trading)}
\begin{itemize}
    \item \textbf{정의:} 은행이 고객을 돕는 것이 아니라, 자체 이익을 위해 금융 상품에 적극적인 포지션을 취하는 활동입니다.
    \item \textbf{포지션 거래 (Position Trade):} 유리한 가격 변동을 예상하여 몇 주 또는 몇 달 동안 증권에 대해 롱(Long) 또는 숏(Short) 포지션을 취하는 것입니다.
    \item \textbf{차익 거래 (Arbitrage Trading):} 금융시장의 작고 일시적인 가격 불일치(Mispricing)를 이용하는 것입니다. 이론적으로는 존재해서는 안 되는 이러한 불일치는 이익을 얻을 기회를 제공하며, 시장의 유동성을 높이고 가격 효율성을 개선하는 데 기여하기도 합니다.
\end{itemize}

\subsection{투자은행 및 거래 수익 (2020년)}
미국 및 유럽의 12개 대형 금융기관의 총 투자은행 및 거래 수익은 2020년에 약 2,000억 달러에 육박했습니다.
\begin{itemize}
    \item \textbf{수익 구성:} 이 중 약 1,500억 달러(약 75\%)가 채권 및 주식 거래에서 발생했습니다.
    \item \textbf{수수료 수익:} 나머지 약 25\%는 투자은행 수수료에서 발생했습니다.
\end{itemize}
이는 거래 활동이 투자은행의 전체 수익에서 매우 큰 비중을 차지함을 보여줍니다.

\begin{summarybox}
\textbf{투자은행의 주요 활동 요약:}
\begin{itemize}
    \item \textbf{증권 인수 (Securities Underwriting):} 기업이 새로운 증권을 발행하고 투자자에게 판매하도록 돕습니다.
    \item \textbf{거래 (Trading):} 이미 발행된 증권의 거래를 돕고, 자체적으로 거래하여 이익을 창출합니다.
    \textbf{기업 전략 자문 (Corporate Strategy):} M\&A 및 기업 구조조정에 대한 자문을 제공합니다.
\end{itemize}
이 세 가지 활동은 모두 투자은행의 중요한 수수료 기반 수입원입니다.
\end{summarybox}


\section{현대 은행의 재무제표 분석}

\subsection{대차대조표 (Balance Sheet)}
은행이 기업으로서 어떻게 조직되고, 은행 산업 및 비즈니스 모델이 어떻게 진화했는지 이해했습니다. 이제 은행이 무엇을 하는지 더 명확히 이해하기 위해 상업은행 지주회사의 대차대조표를 자세히 살펴보겠습니다. 대차대조표는 자산(Assets) 측면에서 투자 목록을, 부채(Liabilities) 측면에서 자금 조달원을 보여줍니다. JP Morgan Chase의 대차대조표와 전체 상업은행 산업의 통합 대차대조표를 사례 연구로 활용할 것입니다.

\begin{summarybox}
\textbf{핵심 요약:} 대차대조표는 특정 시점의 은행 재무 상태를 보여주는 스냅샷입니다. '자산'은 은행이 소유한 것(투자), '부채'는 은행이 빚진 것(자금 조달원), '자본'은 소유주의 몫을 나타냅니다.
\end{summarybox}

\subsubsection{상업은행 대차대조표의 기본 구조}
상업은행 대차대조표의 간단한 구조는 다음과 같습니다.

\begin{itemize}
    \item \textbf{자산 (Asset Side):} 은행이 조달한 자금을 어떻게 사용하는지를 보여줍니다.
    \begin{itemize}
        \item \textbf{현금 (Cash):} 금고에 있는 현금, 다른 은행의 예금 계좌.
        \item \textbf{대출 (Loans):} 주택 담보 대출, 기업 대출, 가계 대출 등.
        \item \textbf{기타 금융 자산 (Other Financial Assets):} 국채, 회사채, 기타 증권 등.
        \item \textbf{유형 자산 (Physical Assets):} 건물, 비품 등 (예: 대출 담당자가 앉는 테이블과 의자).
    \end{itemize}
    \item \textbf{부채 (Liability Side):} 은행이 어떻게 자금을 조달하는지를 보여줍니다.
    \begin{itemize}
        \item \textbf{고객 예금 (Customer Deposits):} 당좌 예금, 저축 예금 등.
        \item \textbf{도매 자금 조달 (Wholesale Funding):} 머니 마켓에서 조달하는 자금.
        \item \textbf{은행 자본 (Bank Capital/Equity):} 은행 소유주가 제공한 자금, 유보 이익(Retained Earnings), 추가 주식 발행 등.
    \end{itemize}
\end{itemize}

\subsubsection{뱅킹 북과 트레이딩 북}
대규모 상업은행 지주회사는 상당한 시장 조성 및 증권 거래 활동을 수행하므로, 대차대조표를 '뱅킹 북(Banking Book)'과 '트레이딩 북(Trading Book)'으로 구분합니다.

\begin{itemize}
    \item \textbf{뱅킹 북 (Banking Book):}
    \begin{itemize}
        \item \textbf{특징:} 전통적인 은행의 장기 차입 및 대출 활동을 나타냅니다.
        \item \textbf{수익원:} 대출에서 발생하는 이자 수입과 수수료를 창출합니다.
        \item \textbf{기록 방식:} 장기적인 관점에서 자산과 부채를 기록합니다.
    \end{itemize}
    \item \textbf{트레이딩 북 (Trading Book):}
    \begin{itemize}
        \item \textbf{특징:} 금융시장에서 빠르게 사고팔 수 있는 증권, 상품, 통화에 대한 투자를 포함합니다.
        \item \textbf{기록 방식:} 롱 포지션(매수 포지션)은 자산 측에 시장 가치로 기록되고, 숏 포지션(차입을 통한 매도 포지션)은 부채로 기록됩니다.
        \item \textbf{수익원:} 거래 포트폴리오에서 발생하는 거래 이익 및 손실은 비이자 수익 및 비용으로 기록됩니다.
    \end{itemize}
\end{itemize}

\subsubsection{JP Morgan Chase의 대차대조표 분석 (2020년)}
JP Morgan Chase는 세계에서 가장 크고 복잡한 은행 중 하나입니다. 2020년 소매 및 도매 사업 부문을 통합한 대차대조표를 살펴보겠습니다. 모든 숫자는 백만 달러 단위입니다.

\begin{itemize}
    \item \textbf{총 자산 (Total Assets):} 약 3.4조 달러 (Trillion).
    \item \textbf{총 부채 (Total Liabilities):} 약 3.1조 달러.
    \item \textbf{총 주주 자본 (Stockholder's Equity):} 약 3,000억 달러 (자산 - 부채).
\end{itemize}

\paragraph{자산 측면의 주요 항목:}
\begin{enumerate}
    \item \textbf{현금 (Cash):} 은행이 보유한 현금과 다른 은행의 예금 계좌에 예치된 현금을 포함합니다. 이는 은행이 보유한 가장 유동적인 자산입니다.
    \item \textbf{연방 기금 판매 및 환매 조건부 계약 (Federal Funds Sold and Reverse Repo Agreements):} JP Morgan이 다른 금융기관(예: 다른 은행)에 단기적으로 자금을 대출해 주는 것을 의미합니다.
    \item \textbf{거래 계정 자산 (Trading Account Assets):} 단기적인 포지션 또는 차익 거래의 일환으로 매도할 의도로 보유하는 증권입니다.
    \item \textbf{투자 증권 (Investment Securities):} 은행의 거래 포트폴리오에 속하지 않는 장기 채무 증권 투자입니다.
    \item \textbf{대출 (Loans):} 대차대조표에서 가장 큰 항목으로, 약 1조 달러에 달합니다.
    \begin{itemize}
        \item \textbf{소매 부문 대출:} 약 5,000억 달러는 개인 및 소기업 대출, 주거용 부동산 담보 대출, 신용카드, 자동차 대출 등 소매 부문에서 발생합니다.
        \item \textbf{도매 부문 대출:} 나머지 5,000억 달러는 대기업 및 정부에 대한 대출 등 도매 부문에서 발생합니다.
    \end{itemize}
\end{enumerate}
JP Morgan은 광범위한 지점 네트워크를 통해 약 1조 달러의 예금을 조달하며, 이는 소매 뱅킹 사업의 대출 규모보다 많습니다. 이는 JP Morgan과 다른 은행들이 상대적으로 저렴하고 안정적인 예금 자금을 사용하여 다른 투자 활동을 지원한다는 것을 의미합니다.

\paragraph{부채 측면의 주요 항목:}
\begin{enumerate}
    \item \textbf{예금 (Deposits):} JP Morgan의 가장 큰 자금 조달원입니다. 당좌 예금(Transaction Accounts)과 비거래 계좌(Non-Transaction Accounts, 주로 단기 채권과 유사한 예금 증서)를 포함합니다.
    \item \textbf{차입금 (Borrowed Money):} 두 번째로 중요한 자금 조달원입니다. 단기 담보 부채(예: 환매 조건부 시장을 통한 차입)부터 장기 무담보 부채(예: 장기 선순위 채권 발행)까지 다양합니다.
    \item \textbf{은행 자본 (Bank Equity Capital):} 은행 소유주로부터 조달된 자금입니다.
    \begin{itemize}
        \item \textbf{납입 자본 (Paid-in Capital):} 보통주 및 우선주 발행을 통해 조달됩니다.
        \item \textbf{유보 이익 (Retained Earnings):} JP Morgan 자본의 대부분은 배당되지 않고 사업에 재투자된 이익에서 나옵니다.
    \end{itemize}
\end{enumerate}
JP Morgan은 투자를 위해 많은 부채를 사용하며, 총 자금 조달에서 자본의 비중은 약 8\%입니다. 이는 은행의 자기자본 승수(Equity Multiplier)가 약 12배라는 것을 의미합니다.

\begin{warningbox}
\textbf{은행의 높은 레버리지:} 은행은 일반 기업에 비해 부채 비율이 매우 높습니다. 이는 예금이라는 저렴하고 안정적인 자금원을 활용하기 때문이지만, 동시에 금융 위기 시 취약성을 높이는 요인이 되기도 합니다.
\end{warningbox}

\subsubsection{미국 상업은행 산업 전체의 통합 대차대조표 (2018년)}
2018년 기준 미국 내 약 5,000개 상업은행의 통합 대차대조표를 살펴보겠습니다. 정부는 은행 시스템의 안전성을 모니터링하기 위해 이러한 통계를 수집합니다.

\paragraph{자산 측면:}
\begin{itemize}
    \item \textbf{총 자산:} 약 16.5조 달러.
    \item \textbf{대출 (Loans):} 약 9조 달러 (총 자산의 약 55\%).
    \begin{itemize}
        \item 상업 및 주거용 부동산 대출, 기업 대출, 개인 대출 순으로 중요합니다.
    \end{itemize}
    \item \textbf{투자 증권 (Investment Securities):} 총 자산의 약 25\%로 두 번째로 중요한 자산입니다.
    \begin{itemize}
        \item 미국 국채, 정부 기관 모기지 담보 증권 등이 가장 큰 하위 구성 요소입니다. 이 자산들은 쉽게 매각할 수 있어 은행의 현금 준비금에 대한 좋은 백업 역할을 합니다.
    \end{itemize}
    \item \textbf{현금 (Cash):} 총 자산의 약 10\%를 차지합니다.
\end{itemize}

\begin{infobox}
\textbf{현금 보유량 증가:} 2008년 금융 위기 이후 많은 은행이 현금 부족을 겪었기 때문에, 은행들은 역사적으로 높은 수준의 현금을 보유하는 경향이 있습니다. 또한, 2008년 이후 낮은 금리는 현금 보유의 기회비용을 낮추는 요인이 되었습니다.
\end{infobox}

\paragraph{가장 작은 지역 은행 (Community Banks)의 대출 포트폴리오:}
총 자산 10억 달러 미만의 지역 소매 은행들은 다음과 같은 특징을 보입니다.
\begin{itemize}
    \item \textbf{신용카드 역할 미미:} 신용카드는 이들 은행 사업에서 거의 역할을 하지 않습니다. 단위당 대출 발생 및 서비스 비용이 너무 비싸기 때문입니다. 대형 은행과 전문 금융 회사가 대규모 신용카드 풀을 다룹니다.
    \item \textbf{기업 대출 비중 낮음:} 기업 대출의 역할도 예상보다 작습니다. 대부분의 상업 대출은 대형 은행이 담당합니다.
    \item \textbf{부동산 대출 집중:} 지역 은행들은 주거용 및 상업용 부동산, 건설 및 개발 대출 등 부동산 대출에 집중합니다. 이들은 지역 시장을 매우 잘 알고 있으며, 이것이 이들의 비교 우위입니다.
\end{itemize}

\paragraph{부채 측면:}
\begin{itemize}
    \item \textbf{주요 자금 조달원:} 예금, 차입금, 자기 자본.
    \item \textbf{예금 (Deposits):} 총 투자의 약 77\%를 조달합니다. 대부분 미국 내 국내 지점에서 보유됩니다.
    \item \textbf{비예금 차입금 (Non-deposit Borrowings):} 자산의 약 11\%를 차지합니다.
    \item \textbf{자기 자본 (Equity Capital):} 총 자산에서 부채를 뺀 나머지 약 11\%를 조달합니다.
\end{itemize}

\begin{figure}[h!]
\centering
\caption{미국 상업은행의 부채 구성 변화 (1970년대~현재)}
\label{fig:us_bank_liabilities_trend}
% 이미지 파일이 없으므로 텍스트로 설명합니다.
\begin{tcolorbox}[title={차트 설명: 미국 상업은행의 부채 구성 변화}]
이 차트는 1970년대부터 현재까지 미국 상업은행의 부채 구성을 보여줍니다.
\begin{itemize}
    \item \textbf{예금의 중요성:} 예금은 총 상업은행 자산의 70\%에서 80\% 사이를 꾸준히 유지하며 가장 중요한 자금 조달원임을 보여줍니다.
    \item \textbf{2008년 금융 위기 이후:} 2008년 금융 위기 이후 은행들은 더욱 신중해졌으며, 안정적이고 신뢰할 수 있는 예금의 역할이 더욱 중요해졌습니다.
\end{itemize}
\end{tcolorbox}
\end{figure}

\begin{summarybox}
\textbf{대차대조표 분석 요약:}
\begin{itemize}
    \item \textbf{자산 측면:} 대출, 증권, 현금이 주요 항목입니다.
    \item \textbf{부채 측면:} 예금, 자본 시장 차입금, 자기 자본이 주요 자금 조달원입니다.
    \item \textbf{최근 동향:} 2008년 금융 위기 이후 은행들은 더 많은 현금을 보유하고, 투자를 위해 더 많은 예금을 활용하는 등 더욱 신중하게 행동하고 있습니다.
\end{itemize}
\end{summarybox}


\section{대차대조표 외 활동 (Off-Balance Sheet Activities)}

\subsection{우발 자산 및 우발 부채의 개념}
대차대조표 외 활동(Off-Balance Sheet, OBS)은 현재 은행의 대차대조표에 자산이나 부채로 기록되지 않지만, 특정 사건이 발생할 경우 은행의 재무제표나 현금 흐름에 중대한 영향을 미칠 수 있는 항목들을 말합니다. 이는 현대 회계 기준에 따라 투자자와 규제 당국에 정보를 제공하기 위해 재무제표에 기록되어야 합니다.

\begin{itemize}
    \item \textbf{우발 자산 (Contingent Asset):} 특정 사건 발생 시 은행의 자산으로 전환되거나 수입에 영향을 미치는 항목입니다.
    \begin{itemize}
        \item \textbf{예시: 신용카드 약정 (Credit Card Commitment):} 은행은 고객에게 신용 한도만큼 대출을 제공하겠다는 약속을 합니다. 고객이 신용카드를 사용할 때(우발 사건), 이 미사용 약정은 대차대조표 외 자산에서 대차대조표 내 자산(소비자 대출)으로 전환됩니다.
    \end{itemize}
    \item \textbf{우발 부채 (Contingent Liability):} 특정 사건 발생 시 은행의 부채로 전환되거나 비용에 영향을 미치는 항목입니다.
    \begin{itemize}
        \item \textbf{예시: 보증 (Guarantees) 또는 보험:} 은행이 보증을 제공하거나 보험을 판매할 때, 이는 대차대조표 외 부채로 기록됩니다. 보증 사건이 발생하거나 보험 청구가 들어오면(우발 사건), 은행은 현금 지출을 해야 하므로 대차대조표에 중대한 영향을 미칩니다.
    \end{itemize}
\end{itemize}
은행이 신용카드를 발행하거나 보험을 판매하는 것을 선호하는 이유는 이러한 활동이 선불 수수료(Upfront Fees)를 발생시키기 때문입니다.

\begin{summarybox}
\textbf{핵심 요약:} 대차대조표 외 활동은 미래의 특정 사건에 따라 은행의 자산이나 부채로 변할 수 있는 잠재적 항목입니다. 이는 은행의 현재 재무 상태를 넘어선 잠재적 위험과 수익 기회를 보여줍니다.
\end{summarybox}

\subsection{주요 대차대조표 외 활동}
은행의 대차대조표 외 활동은 전통적인 은행 대출과 유사한 활동도 있고, 보험과 유사한 활동도 있습니다.

\begin{itemize}
    \item \textbf{자산 측면의 주요 항목:}
    \begin{itemize}
        \item \textbf{미사용 대출 약정 (Undrawn Loan Commitments):} 신용카드 한도와 같이 미래에 인출될 수 있는 대출 약정입니다.
        \item \textbf{롱 파생상품 포지션 (Long Derivatives Positions):} 미래에 이익 또는 손실을 발생시킬 수 있는 파생상품 계약입니다.
    \end{itemize}
    \item \textbf{부채 측면의 주요 항목:}
    \begin{itemize}
        \item \textbf{숏 파생상품 포지션 (Short Derivatives Positions):} 미래에 이익 또는 손실을 발생시킬 수 있는 파생상품 계약입니다.
        \item \textbf{보증 (Guarantees):} 다양한 형태를 취하지만, 대부분 보험과 유사하게 작동합니다. 특정 조건 하에 은행이 보험금 지급을 해야 할 수 있습니다.
    \end{itemize}
\end{itemize}

\subsubsection{미사용 대출 약정 (Loan Commitments)}
대출 약정은 특정 최대 한도 내에서 미리 정해진 이자율로 자금을 제공하겠다는 법적 구속력이 있는 약속입니다. 이는 가계나 기업에 제공될 수 있으며, 무담보 신용카드나 주택 담보 신용 한도 등이 있습니다.

\begin{itemize}
    \item \textbf{은행의 의무:} 은행은 약정 기간 내에 언제든지 자금을 제공할 준비가 되어 있어야 합니다. 고객이 카드를 사용하면 대출이 은행의 대차대조표에 기록됩니다.
    \item \textbf{활용률 (Utilization Rates):} FDIC 자료에 따르면, 주택 담보 대출, 기업 대출, 신용카드 등 다양한 유형의 대출 약정 활용률은 100\% 미만입니다. 이는 대차대조표 외에 미인출 금액이 존재한다는 것을 의미합니다.
    \item \textbf{신용카드 예시:} 신용카드 미인출 잔액은 인출 잔액의 4~5배에 달할 수 있습니다. 은행은 이러한 미인출 잔액에 대해서도 수수료를 받을 수 있습니다.
    \item \textbf{규제 당국의 요구:} 규제 당국은 과도한 인출이 은행에 위험을 초래할 수 있으므로, 은행이 이러한 미인출 잔액을 재무제표에 공개하도록 요구합니다.
\end{itemize}

\subsubsection{파생상품 (Derivatives)}
은행은 파생상품을 사용하여 금리 변동과 같은 불리한 위험 노출을 줄일 수 있습니다. 또한, 고객과의 파생상품 거래에서 상대방이 되어 수수료를 받을 수도 있습니다.

\begin{itemize}
    \item \textbf{현금 흐름 영향:} 파생상품의 매수 및 매도는 미래에 현금 유입 또는 유출로 이어질 수 있으므로, 은행은 이러한 거래를 대차대조표 외에 기록해야 합니다.
    \item \textbf{명목 금액 (Notional Amounts):} 파생상품 보고는 명목 금액으로 이루어집니다. 금리 파생상품의 경우, 150조 달러에 달하는 엄청난 규모를 보입니다. 이는 '보험'으로 생각했을 때, 보험에 가입된 항목의 총 가치를 나타냅니다.
    \item \textbf{시장 가치와의 차이:} 명목 금액은 실제 은행의 잠재적 현금 흐름 노출이나 포지션의 시장 가치와는 다릅니다. (예: 100만 달러 주택 보험의 명목 금액은 100만 달러이지만, 보험료(시장 가치)는 2,000달러일 수 있습니다.)
    \item \textbf{성장률:} 지난 30년간 대출 약정은 은행 대차대조표와 거의 비슷한 5배 성장을 보인 반면, 파생상품은 20배 폭발적인 성장을 보였습니다.
    \item \textbf{집중도:} 이러한 파생상품 활동의 대부분은 JP Morgan Chase, Citi Bank, Bank of America, Goldman Sachs, Morgan Stanley, Wells Fargo, HSBC와 같은 소수의 대형 금융기관에 집중되어 있습니다. 이는 이 사업에서 규모의 경제가 작용함을 반영합니다.
\end{itemize}

\paragraph{JP Morgan Chase의 파생상품 노출 (2020년):}
JP Morgan은 금리 파생상품, 신용 파생상품, 외환, 주식, 상품 계약 등 수천 건의 파생상품 거래에 참여합니다. 대부분의 활동은 고객을 대신하여 이루어집니다.

\begin{itemize}
    \item \textbf{상쇄 (Netting):} JP Morgan이 금리 옵션을 동시에 매수하고 매도하는 것처럼, 많은 거래가 상쇄(Netting)될 수 있습니다. 거래가 많을수록 상쇄 효과가 커집니다.
    \item \textbf{신용부도스왑 (Credit Default Swap, CDS):} CDS는 채권이나 대출의 부도 위험에 대한 보험 계약으로 생각할 수 있습니다.
    \begin{itemize}
        \item JP Morgan은 약 1조 달러 규모의 CDS에 관여하지만, 이는 보험 판매와 보험 매수 사이에 균등하게 분산되어 있습니다.
        \item \textbf{순 노출 (Net Exposure):} 전체 순 노출은 훨씬 작습니다. JP Morgan이 매수하는 보험에 대해 지불하는 금액보다 보험 판매로 징수하는 수수료가 훨씬 적다면, 이 활동은 은행에 너무 많은 위험 노출을 만들지 않으면서 좋은 이익을 가져다줄 수 있습니다.
    \end{itemize}
\end{itemize}

\begin{warningbox}
\textbf{2008년 금융 위기와 OBS 위험:} 대차대조표 외에 숨겨진 위험은 2008년 금융 위기의 핵심 원인이었습니다. 제대로 이해되지 않은 이러한 위험은 세계 최대 금융기관들의 파산, 부실 인수 또는 구제금융으로 이어졌습니다. 이는 수많은 새로운 규제를 낳았습니다.
\end{warningbox}

\begin{summarybox}
\textbf{대차대조표 외 활동 요약:}
\begin{itemize}
    \item \textbf{개념:} 미래의 특정 사건 발생 시 은행의 재무 상태에 영향을 미칠 수 있는 잠재적 자산 및 부채입니다.
    \item \textbf{주요 활동:} 대출 약정(신용카드 한도 등)과 파생상품이 가장 중요하며, 은행에 많은 수수료 수입을 가져다줍니다.
    \item \textbf{위험:} 명목 금액이 크더라도 상쇄 효과(Netting)로 실제 순 노출은 작을 수 있지만, 2008년 금융 위기에서 보았듯이 OBS 활동의 위험은 매우 중요하며 엄격한 규제가 필요합니다.
\end{itemize}
\end{summarybox}


\section{FAQ (자주 묻는 질문)}

\begin{itemize}
    \item \textbf{Q: 투자은행과 상업은행의 가장 큰 차이점은 무엇인가요?}
    \item \textbf{A:} 상업은행은 주로 예금을 받아 대출을 해주는 전통적인 역할을 합니다. 반면 투자은행은 기업의 주식/채권 발행을 통한 자금 조달, 인수합병 자문, 증권 거래 등 자본시장 관련 활동에 특화되어 있습니다. 최근에는 많은 대형 은행들이 이 두 가지 기능을 모두 수행하는 경향이 있습니다.

    \item \textbf{Q: 2008년 금융 위기가 투자은행 산업에 어떤 영향을 미쳤나요?}
    \item \textbf{A:} 2008년 금융 위기는 투자은행 산업에 엄청난 변화를 가져왔습니다. Lehman Brothers는 파산했고, Bear Stearns와 Merrill Lynch는 상업은행에 인수되었습니다. Goldman Sachs와 Morgan Stanley는 생존을 위해 상업은행 지주회사로 전환해야 했습니다. 이로 인해 독립형 투자은행은 거의 사라지고, 대부분 상업은행 지주회사 산하로 편입되었습니다.

    \item \textbf{Q: '뱅킹 북'과 '트레이딩 북'은 왜 구분하나요?}
    \item \textbf{A:} 뱅킹 북은 은행의 전통적인 장기 대출 및 예금 활동을, 트레이딩 북은 단기적인 증권 매매 활동을 기록합니다. 이 둘을 구분하는 이유는 활동의 성격(장기 vs. 단기), 위험 프로필, 수익 인식 방식(이자 수익 vs. 거래 손익)이 다르기 때문입니다. 이는 은행의 다양한 사업 모델을 명확히 보여줍니다.

    \item \textbf{Q: 대차대조표 외 활동(OBS)이 중요한 이유는 무엇인가요?}
    \item \textbf{A:} OBS 활동은 현재 은행의 대차대조표에는 나타나지 않지만, 미래에 특정 사건이 발생하면 은행의 자산, 부채, 현금 흐름에 중대한 영향을 미칠 수 있는 잠재적 위험과 수익 기회를 포함합니다. 신용카드 약정이나 파생상품 등이 대표적이며, 2008년 금융 위기에서 OBS 위험이 은행 시스템에 얼마나 큰 영향을 미칠 수 있는지 드러나면서 그 중요성이 더욱 강조되었습니다.

    \item \textbf{Q: 파생상품의 '명목 금액'이 크다고 해서 은행의 위험이 그만큼 큰 것은 아닌가요?}
    \textbf{A:} 명목 금액은 파생상품 계약의 기초자산 총액을 나타내지만, 실제 은행의 잠재적 손실 노출과는 다릅니다. 은행은 수많은 파생상품 거래를 통해 롱 포지션과 숏 포지션을 동시에 보유하여 위험을 상쇄(Netting)할 수 있습니다. 따라서 순 노출(Net Exposure)은 명목 금액보다 훨씬 작을 수 있습니다. 그러나 명목 금액의 급격한 증가는 잠재적 위험의 증가를 시사하므로 규제 당국이 면밀히 주시합니다.
\end{itemize}


\section{빠르게 훑어보기 (1페이지 요약)}

\begin{tcolorbox}[title={\textbf{투자은행의 역할과 조직}}]
\begin{itemize}
    \item \textbf{역할:} 기업의 자금 조달(주식/채권 발행), M\&A 자문, 증권 거래를 통해 차입자와 저축자를 연결.
    \item \textbf{주요 사업:} 자금 조달(인수), 증권 거래(시장 조성/자기 계정 거래), M\&A 자문.
    \item \textbf{조직:} JP Morgan 사례처럼 상업은행 지주회사 내 '뱅킹'과 '마켓'으로 구성.
    \item \textbf{유형:} 상업은행 계열사(JP Morgan), 전문화된 투자은행(Goldman Sachs), 독립형(Lazard).
    \textbf{산업 변화:} 1990년대 규제 완화와 2008년 금융 위기로 대부분 상업은행과 통합.
\end{itemize}
\end{tcolorbox}

\begin{tcolorbox}[title={\textbf{투자은행 활동의 세부 사항}}]
\begin{itemize}
    \item \textbf{증권 인수:} IPO, 추가 발행, 채권 인수, 신디케이트론 주간사 역할. 북빌딩을 통해 가격 결정 및 판매.
    \item \textbf{M\&A 자문:} 기업 가치 평가, 협상 지원, 자금 조달 지원. 높은 수수료 수익원.
    \item \textbf{증권 거래:}
    \begin{itemize}
        \item \textbf{시장 조성:} 매수/매도 호가 제시로 유동성 공급, 비드-애스크 스프레드 수익.
        \item \textbf{자기 계정 거래:} 은행 자체 이익을 위한 포지션 거래 및 차익 거래.
    \end{itemize}
    \item \textbf{수익 비중:} 거래 수익(약 75\%)이 투자은행 수수료(약 25\%)보다 훨씬 큼.
\end{itemize}
\end{tcolorbox}

\begin{tcolorbox}[title={\textbf{현대 은행의 재무제표 분석}}]
\begin{itemize}
    \item \textbf{대차대조표:} 자산(은행의 투자)과 부채(자금 조달원)를 보여주는 재무 스냅샷.
    \item \textbf{뱅킹 북 vs. 트레이딩 북:}
    \begin{itemize}
        \item \textbf{뱅킹 북:} 전통적 장기 대출/예금 활동.
        \item \textbf{트레이딩 북:} 단기 증권/상품/통화 거래.
    \end{itemize}
    \item \textbf{JP Morgan 자산:} 대출(가장 큼), 투자 증권, 현금, 연방 기금 판매.
    \item \textbf{JP Morgan 부채:} 예금(가장 큼), 차입금, 자기 자본. 높은 레버리지(자본 8\%, 승수 12배).
    \item \textbf{산업 전체:} 대출(55\%), 투자 증권(25\%), 현금(10\%). 예금(77%)이 주요 자금원.
    \item \textbf{지역 은행:} 부동산 대출에 집중, 신용카드/기업 대출 비중 낮음.
\end{itemize}
\end{tcolorbox}

\begin{tcolorbox}[title={\textbf{대차대조표 외 활동 (OBS)}}]
\begin{itemize}
    \item \textbf{개념:} 현재 대차대조표에 없지만, 특정 사건 발생 시 재무에 영향을 미치는 우발 자산/부채.
    \item \textbf{주요 활동:}
    \begin{itemize}
        \item \textbf{미사용 대출 약정:} 신용카드 한도 등. 인출 시 대차대조표 내 자산으로 전환.
        \item \textbf{파생상품:} 위험 헤지 또는 거래 목적. 명목 금액은 크지만 순 노출은 상쇄 효과로 작을 수 있음.
    \end{itemize}
    \item \textbf{중요성:} 선불 수수료 발생. 2008년 금융 위기의 핵심 원인 중 하나로, 잠재적 위험 관리가 중요.
\end{itemize}
\end{tcolorbox}


\section{개요}

\begin{summarybox}[title={핵심 요약}]
이 문서는 은행의 손익계산서(Income Statement)를 분석하여 현대 은행이 어떻게 수익을 창출하는지 설명합니다. 전통적인 대출 사업에서 발생하는 이자수익과 이자비용의 차이인 순이자이익(Net Interest Income)뿐만 아니라, 투자은행 업무, 자산 관리 등 다양한 비즈니스 라인에서 발생하는 비이자수익(Non-Interest Revenue)의 중요성을 강조합니다. JP Morgan Chase 사례 연구를 통해 이러한 수익원들이 어떻게 결합되어 은행의 최종 순이익을 구성하는지 구체적으로 살펴봅니다.
\end{summarybox}


\section{용어 정리}

\begin{table}[h!]
\centering
\caption{주요 용어 설명}
\label{tab:glossary}
\begin{adjustbox}{width=\textwidth,center}
\begin{tabular}{lllc}
\toprule
\textbf{용어 (한글)} & \textbf{쉬운 설명} & \textbf{원어} & \textbf{비고} \\
\midrule
스프레드 & 대출로 벌어들인 이자수익과 예금에 지급한 이자비용의 차이. & Spread & 은행의 전통적인 수익원 \\
순이자이익 & 이자수익에서 이자비용을 뺀 금액. 은행 대출 사업의 핵심 수익성 지표. & Net Interest Income (NII) & 손익계산서의 최상단 지표 \\
비이자수익 & 이자 외의 모든 수익. 수수료, 거래 이익 등이 포함됨. & Non-Interest Revenue & 현대 은행의 중요한 수익원 \\
비이자비용 & 이자 지급과 무관한 모든 운영 비용. 급여, 임대료, 기술 투자 등. & Non-Interest Expense & 은행 운영에 필요한 제반 비용 \\
자기매매 & 은행이 고객을 대신하는 것이 아니라, 자체 자금으로 금융 시장에서 & Proprietary Trading & 고위험 고수익 전략 \\
& 투기적 포지션을 취하여 이익을 얻는 활동. & & \\
인수 업무 & 기업의 주식이나 채권 발행을 돕고 이를 투자자에게 판매하는 업무. & Underwriting & 투자은행의 핵심 기능 \\
시장 조성 & 특정 증권의 매수-매도 호가를 제시하여 시장 유동성을 공급하는 역할. & Market Making & 거래 활성화에 기여 \\
3-6-3 모델 & 과거 은행업의 단순한 수익 모델. 예금에 3\% 지급, 6\%로 대출, 3\% 이익. & 3-6-3 Model & 1980년대 이후 압박받음 \\
\bottomrule
\end{tabular}
\end{adjustbox}
\end{table}


\section{핵심 개념·원리: 은행의 손익계산서}

은행의 손익계산서는 단순히 이자수익과 이자비용의 차이만을 보여주는 것이 아닙니다. 현대 은행은 복잡하고 다각화된 사업 구조를 가지고 있으며, 다양한 방식으로 수익을 창출하고 비용을 지출합니다.

\begin{summarybox}[title={은행의 기본적인 수익 구조: 스프레드}]
\textbf{1. 한 줄 핵심 요약:} 은행은 예금을 받아 대출해주고, 대출 이자수익과 예금 이자비용의 차이인 '스프레드'를 통해 이익을 얻습니다.

\textbf{2. 직관적 예시/비유:} 친구에게 돈을 빌려주고(예금 수취) 다른 친구에게 더 높은 이자로 빌려주는(대출) 중개인과 같습니다. 빌려준 돈에서 받은 이자(대출 이자)와 빌린 돈에 대해 지불한 이자(예금 이자)의 차이가 당신의 이익이 됩니다.

\textbf{3. 기술적 정확한 설명:} 상업은행은 예금을 수취하여 개인, 기업, 부동산 대출에 활용합니다. 이 과정에서 은행은 자금을 차입하고, 대출하며, 이자수익(interest income)과 이자비용(interest expense)의 차이인 스프레드(spread)를 얻습니다.
\end{summarybox}

하지만 현대 은행의 수익 구조는 이보다 훨씬 복잡합니다. 오늘날의 은행은 대규모의 다각화된 사업체이며, 대출 사업에서 발생하는 이자수익 외에도 다양한 사업 라인에서 수수료 수익(fee revenue)을 창출합니다.

\subsection{손익계산서의 세 가지 주요 구성 요소}

은행의 손익계산서는 크게 세 부분으로 나눌 수 있습니다. 이 세 가지 구성 요소를 이해하면 현대 은행이 어떻게 이익을 창출하는지 명확하게 파악할 수 있습니다.

\begin{enumerate}
    \item \textbf{대출 사업의 수익성 (순이자이익):}
    \begin{summarybox}[title={순이자이익 (Net Interest Income, NII)}]
    \textbf{1. 한 줄 핵심 요약:} 은행의 전통적인 대출 및 투자 활동에서 발생하는 이자수익에서 자금 조달에 드는 이자비용을 뺀 금액입니다.

    \textbf{2. 직관적 예시/비유:} 당신이 부동산 임대 사업을 한다고 가정해 봅시다. 임대료 수입(이자수익)에서 대출 이자(이자비용)를 뺀 금액이 순수하게 부동산에서 벌어들인 돈과 같습니다.

    \textbf{3. 기술적 정확한 설명:}
    \begin{itemize}
        \item \textbf{이자수익 (Interest Income):} 은행이 대출을 제공하거나 유가증권에 투자함으로써 얻는 수익입니다. 여기에는 개인 대출, 기업 대출, 부동산 대출, 다른 금융기관에 대한 대출, 그리고 은행이 보유한 유가증권(예: 국채, 회사채) 투자에서 발생하는 이자가 포함됩니다.
        \item \textbf{이자비용 (Interest Expense):} 은행이 예금자에게 지급하거나 금융시장에서 자금을 차입할 때 발생하는 비용입니다. 여기에는 고객 예금에 대한 이자, 단기 및 장기 차입금에 대한 이자, 그리고 은행이 발행한 장기 채권의 쿠폰 지급액 등이 포함됩니다.
    \end{itemize}
    \textbf{순이자이익 (NII) = 이자수익 - 이자비용} 이며, 이는 은행의 대출 사업 수익성을 나타내는 가장 중요한 지표입니다.
    \end{summarybox}

    \item \textbf{비이자수익 (Non-Interest Revenue):}
    \begin{summarybox}[title={비이자수익 (Non-Interest Revenue)}]
    \textbf{1. 한 줄 핵심 요약:} 이자수익 외에 은행이 다양한 서비스 제공을 통해 얻는 모든 수수료 및 기타 수익입니다.

    \textbf{2. 직관적 예시/비유:} 식당을 운영한다고 할 때, 음식 판매 수익(이자수익) 외에 배달 수수료, 예약 수수료, 케이터링 서비스 수수료 등 부가적인 서비스로 버는 돈과 같습니다.

    \textbf{3. 기술적 정확한 설명:} 현대 은행은 대출 외에도 다양한 금융 서비스를 제공하며 수수료를 받습니다. 이는 '수수료 수익(fee income)'이라고도 불리며, 투자은행 업무, 자산 관리, 중개 서비스, 계좌 관리 수수료, 신용카드 수수료 등 매우 광범위합니다. 이는 은행의 수익원을 다각화하고 안정성을 높이는 데 기여합니다.
    \end{summarybox}

    \item \textbf{비이자비용 (Non-Interest Expense):}
    \begin{summarybox}[title={비이자비용 (Non-Interest Expense)}]
    \textbf{1. 한 줄 핵심 요약:} 이자 지급과 직접적인 관련이 없는 은행 운영에 필요한 모든 비용을 포괄합니다.

    \textbf{2. 직관적 예시/비유:} 회사를 운영할 때, 원자재 구매 비용(이자비용) 외에 직원 급여, 사무실 임대료, 전기세, 마케팅 비용 등 회사를 유지하는 데 드는 모든 일반적인 지출과 같습니다.

    \textbf{3. 기술적 정확한 설명:} 비이자비용은 이자 지급과 무관한 모든 운영 비용을 포함하는 포괄적인 항목입니다. 여기에는 직원 급여 및 보너스(가장 큰 비중), 사무실 임대료, 기술 투자 비용, 마케팅 비용, 감가상각비, 그리고 기타 소규모 운영 비용 등이 포함됩니다.
    \end{summarybox}
\end{enumerate}

이 세 가지 요소를 종합하여 은행의 최종 \textbf{순이익 (Net Income)}이 계산됩니다.
\textbf{순이익 = 순이자이익 + 비이자수익 - 비이자비용}


\section{사례 연구: JP Morgan Chase의 손익 구조 (2020년 기준)}

JP Morgan Chase는 세계에서 가장 크고 다각화된 은행 중 하나입니다. 이 은행의 2020년 손익계산서를 통해 위에서 설명한 개념들이 실제 어떻게 적용되는지 구체적으로 살펴보겠습니다.

\subsection{JP Morgan의 사업 부문}

JP Morgan의 사업은 크게 두 가지로 나눌 수 있습니다.

\begin{enumerate}
    \item \textbf{전통적인 은행 업무 (대출):}
    소매 금융(retail banking) 부문과 도매 금융(wholesale business) 내의 상업 대출(commercial lending) 부문이 여기에 해당합니다. 이 부문에서 주로 \textbf{순이자이익}이 발생합니다.

    \item \textbf{비이자수익 창출 업무 (수수료):}
    주로 투자은행(investment banking) 부문에서 발생하며, 인수(underwriting) 업무, 시장 조성(market making), 자산 관리(asset management) 등이 포함됩니다. 여기서 \textbf{비이자수익}이 발생합니다.
\end{enumerate}

\subsection{이자수익과 이자비용}

2020년 JP Morgan의 이자 관련 수익과 비용은 다음과 같습니다.

\begin{itemize}
    \item \textbf{이자수익 (Interest Income):} 약 \textbf{650억 달러}
    \begin{itemize}
        \item 대출 포트폴리오 (loan portfolio)
        \item 은행 계정(banking book) 및 트레이딩 포트폴리오(trading portfolio)에 보유된 투자 유가증권
        \item 다른 금융기관에 대한 대출
    \end{itemize}
    \item \textbf{이자비용 (Interest Expense):}
    \begin{itemize}
        \item 예금에 대한 이자
        \item 기타 차입금 (단기 및 장기 부채, JP Morgan이 발행한 장기 채권의 쿠폰 지급액 포함)
    \end{itemize}
\end{itemize}

이 둘의 차이인 \textbf{순이자이익 (Net Interest Income)}은 2020년에 약 \textbf{550억 달러}였습니다. 이는 JP Morgan의 대출 사업에서 발생한 수익성을 나타냅니다.

\subsection{비이자수익 (Non-Interest Income)}

순이자이익에서 전체 수익성으로 넘어가기 위해서는 수수료 수익을 고려해야 합니다. JP Morgan의 비이자수익은 매우 다양하며, 주로 도매 금융 부문에서 발생합니다. 2020년 JP Morgan의 비이자수익은 약 \textbf{640억 달러}에 달했습니다. 주요 항목은 다음과 같습니다.

\begin{enumerate}
    \item \textbf{투자은행 수수료 (Investment Banking Fees):} 약 \textbf{90억 달러}
    \begin{itemize}
        \item \textbf{M\&A 자문 (M\&A advisory):} 기업 인수합병 관련 자문 서비스.
        \item \textbf{주식 인수 (equity underwriting):} 기업의 주식 발행을 돕고 이를 투자자에게 판매.
        \item \textbf{채권 및 신디케이트론 인수 (bond and syndicated loan underwriting):} 기업의 채권 발행 및 여러 은행이 공동으로 제공하는 대출(신디케이트론) 주선.
    \end{itemize}
    \item \textbf{자기매매 이익 (Proprietary Trading / Principal Transactions):} 약 \textbf{180억 달러}
    \begin{itemize}
        \item JP Morgan의 트레이더들이 고객을 대신하는 것이 아니라, \textbf{은행 자체 계정(on its own account)}으로 금융 시장에서 투기적 포지션을 취하여 이익을 창출하는 활동입니다. 이는 고위험 고수익 사업으로 분류될 수 있습니다.
    \end{itemize}
    \item \textbf{자산 관리 및 중개 서비스 수수료 (Asset Management and Brokerage Services Fees):} 약 \textbf{180억 달러}
    \begin{itemize}
        \item JP Morgan이 고객의 금융 시장 투자를 돕는 서비스입니다.
        \item \textbf{자산 관리 수수료 (asset management fees):} 고객 포트폴리오를 관리해주는 대가로 받는 수수료.
        \item \textbf{중개 수수료 (brokerage commission fees):} 고객의 거래를 중개해주고 받는 수수료.
    \end{itemize}
    \item \textbf{기타 수수료:} 예금 계좌 및 모기지 관련 수수료 (예: 초과 인출 수수료, 연체료), 신용카드 수수료 등.
\end{enumerate}

\begin{warningbox}[title={주목할 점: 비이자수익의 비중}]
JP Morgan과 같이 크고 다각화된 은행의 경우, 2020년 비이자수익(약 640억 달러)이 대출 사업에서 발생한 순이자이익(약 540억 달러)보다 실제로 더 컸습니다. 이는 현대 은행의 수익 구조에서 비이자수익이 얼마나 중요한 부분을 차지하는지 보여줍니다.
\end{warningbox}

\subsection{비이자비용 (Non-Interest Expense)}

JP Morgan의 비이자비용은 다음과 같은 항목들을 포함합니다.

\begin{itemize}
    \item \textbf{급여 및 보너스 (Salaries and bonuses):} 가장 큰 비중을 차지합니다.
    \item 임대료 (Rent)
    \item 기술 투자 (Technology)
    \item 마케팅 (Marketing)
    \item 기타 소규모 항목들
\end{itemize}
은행은 또한 세금(taxes)을 납부해야 합니다.

\subsection{최종 순이익 (Net Income)}

JP Morgan의 2020년 순이익은 약 \textbf{300억 달러}였습니다. 이는 순이자이익과 비이자수익을 합산한 후 비이자비용을 차감하여 계산된 은행의 최종 수익입니다.


\section{현대 은행업의 수익 구조 변화}

JP Morgan의 사례는 현대 은행이 단순히 대출과 예금의 이자 차이로만 수익을 내는 것이 아님을 명확히 보여줍니다.

\subsection{수수료 수익의 중요성}

\begin{summarybox}[title={수수료 수익 증가의 배경}]
\textbf{1. 한 줄 핵심 요약:} 21세기 은행들은 낮은 금리 환경과 대출 시장 경쟁 심화로 인해 이자수익 의존도를 줄이고 수수료 수익을 통해 수익원을 다각화하고 있습니다.

\textbf{2. 직관적 예시/비유:} 스마트폰 제조사가 단순히 기기 판매(이자수익) 외에 앱 스토어 수수료, 클라우드 서비스 구독료(비이자수익) 등으로 수익을 창출하는 것과 유사합니다. 핵심 제품의 수익성이 낮아지면 부가 서비스로 보완하는 것입니다.

\textbf{3. 기술적 정확한 설명:}
\begin{itemize}
    \item \textbf{낮은 금리 환경:} 지속적으로 낮은 금리는 은행의 순이자마진(Net Interest Margin)을 압박하여 대출 사업의 수익성을 저하시킵니다.
    \item \textbf{대출 시장 경쟁:} 핀테크 기업 등 새로운 경쟁자들의 등장으로 대출 시장의 경쟁이 심화되어 전통적인 대출 수익만으로는 성장이 어려워졌습니다.
    \item \textbf{다각화의 필요성:} 수수료 수익은 이러한 환경에서 은행이 잃어버린 이자수익을 보충하고, 수익원을 다각화하여 사업 안정성을 높이는 중요한 수단이 됩니다.
\end{itemize}
\end{summarybox}

\subsection{3-6-3 모델과 수수료 비즈니스로의 전환}

1980년대에 '3-6-3 모델'이라는 전통적인 은행업 모델이 큰 압박을 받기 시작하면서, 은행들은 수수료 기반 비즈니스로의 전환을 가속화했습니다. 3-6-3 모델은 예금에 3\% 이자를 지급하고, 6\% 이자로 대출을 해주며, 3\%의 이익을 남기는 단순한 구조를 의미했습니다. 하지만 규제 완화와 경쟁 심화로 이러한 단순한 모델로는 더 이상 수익을 내기 어려워졌습니다.

\subsection{미국 은행의 평균 수익 구성}

JP Morgan은 미국 최대 은행으로서 다각화된 사업 구조를 가지고 있지만, 수수료 수익에 의존하는 경향은 비단 JP Morgan만의 이야기가 아닙니다. 미국 은행 전체를 평균적으로 보면, 은행 총수익의 약 3분의 1이 수수료에서 발생합니다. 이는 수수료 수익이 현대 은행업에서 보편적으로 중요한 역할을 하고 있음을 보여줍니다.


\section{결론 및 요약}

\begin{summarybox}[title={핵심 정리}]
이 세션에서는 은행의 손익계산서를 심층적으로 분석했습니다. 우리는 현대 은행이 전통적인 대출 사업에서 발생하는 \textbf{순이자이익}뿐만 아니라, 다양한 금융 서비스에서 발생하는 \textbf{비이자수익}을 통해 이익을 창출한다는 것을 확인했습니다. 특히, JP Morgan Chase 사례를 통해 비이자수익이 순이자이익보다 더 큰 비중을 차지할 수 있음을 보았습니다. 21세기 은행들은 낮은 금리와 치열한 경쟁 환경 속에서 수익원을 다각화하고 안정성을 확보하기 위해 수수료 기반 비즈니스의 중요성이 더욱 커지고 있습니다.
\end{summarybox}


%=======================================================================
% Module 3: 개요: 위험과 수익의 균형
%=======================================================================
\chapter{개요: 위험과 수익의 균형}
\label{ch:module3}

\metainfo{Banking and Financial Institutions}{Module 3: Technology Strategy}{Prof. Rustom Manouchehri Irani}{이 문서는 'Banking and Financial Institutions Module 3' 강의 자료를 바탕으로 작성된 독립적인 학습 노트입니다. 현대 은행의 수익원과 그 수익을 둘러싼 다양한 위험 요소를 심층적으로 탐구합니다. 은행이 어떻게 이익을 창출하고, 그 성과를 어떻게 측정하며, 이익이 소유주에게 어떻게 흘러가는지 살펴봅니다. 또한, 은행 사업 모델에 내재된 다양한 위험을 이해하고, 위험 관리자들이 이러한 위험을 어떻게 통제하려 하는지 상세히 설명합니다. 이 모듈을 통해 현대 은행의 '위험-수익 상충 관계(Risk-Return Tradeoff)'를 명확히 이해하는 것이 목표입니다.}


\section{개요: 위험과 수익의 균형}

이 문서는 'Banking and Financial Institutions Module 3' 강의 자료를 바탕으로 작성된 독립적인 학습 노트입니다.
현대 은행의 수익원과 그 수익을 둘러싼 다양한 위험 요소를 심층적으로 탐구합니다.
은행이 어떻게 이익을 창출하고, 그 성과를 어떻게 측정하며, 이익이 소유주에게 어떻게 흘러가는지 살펴봅니다.
또한, 은행 사업 모델에 내재된 다양한 위험을 이해하고, 위험 관리자들이 이러한 위험을 어떻게 통제하려 하는지 상세히 설명합니다.
이 모듈을 통해 현대 은행의 '위험-수익 상충 관계(Risk-Return Tradeoff)'를 명확히 이해하는 것이 목표입니다.

\begin{summarybox}[title=\textbf{학습 목표 요약}]
\begin{itemize}
    \item 은행의 주요 수익원과 이익 창출 방식 이해
    \item 은행 성과 측정 지표 (주식 수익률, ROA, ROE) 분석
    \item 은행 자본 구조와 레버리지의 역할 파악
    \item 현대 은행이 직면하는 다양한 위험 (신용, 금리, 유동성, 시장, 운영 위험) 식별
    \item 각 위험의 측정 및 관리 전략 학습
    \item 위험 조정 수익률의 중요성 이해 및 위험-수익 상충 관계 파악
\end{itemize}
\end{summarybox}


\section{용어 정리}

다음 표는 본 문서에서 다루는 주요 용어들을 비전공자도 쉽게 이해할 수 있도록 설명합니다.

\begin{adjustbox}{width=\textwidth,center}
\begin{tabular}{lllc}
\toprule
\textbf{용어 (Term)} & \textbf{쉬운 설명 (Easy Explanation)} & \textbf{원어 (Original Term)} & \textbf{비고 (Notes)} \\
\midrule
\textbf{자기자본} & 은행의 소유주가 투자한 자금으로, 은행이 손실을 입을 때 가장 먼저 손실을 흡수하는 완충재. & Equity Capital & 손실 흡수 능력 \\
\textbf{잔여 청구권} & 은행이 파산하거나 청산될 때, 모든 부채를 갚고 남은 자산에 대해 소유주가 갖는 권리. & Residual Claim & 부채 상환 후 남은 몫 \\
\textbf{주식 수익률} & 주식 투자에 대한 수익률. 주가 상승과 배당금을 포함. & Stock Returns & 시장 기반 성과 지표 \\
\textbf{총자산 수익률} & 은행이 보유한 총자산 대비 얼마나 효율적으로 이익을 창출했는지 나타내는 지표. & Return on Assets (ROA) & 자산 활용 효율성 \\
\textbf{자기자본 수익률} & 소유주가 투자한 자기자본 대비 얼마나 많은 이익을 창출했는지 나타내는 지표. & Return on Equity (ROE) & 소유주 관점의 수익성 \\
\textbf{레버리지 비율} & 은행이 자산을 조달하기 위해 부채를 얼마나 많이 사용했는지 나타내는 지표. & Leverage Ratio & 부채 활용 정도 \\
\textbf{신용 위험} & 대출을 받은 차입자가 원금이나 이자를 제때 갚지 못할 위험. & Credit Risk & 대출 부도 위험 \\
\textbf{체계적 신용 위험} & 경기 침체와 같이 거시 경제 상황 악화로 인해 많은 차입자가 동시에 부도날 위험. & Systematic Credit Risk & 거시 경제 요인 \\
\textbf{신용 등급} & 차입자가 대출을 갚을 능력이 얼마나 되는지 평가하여 부여하는 점수나 등급. & Credit Rating & 부도 확률 예측 \\
\textbf{약정} & 대출 계약에 포함된 조항으로, 차입자의 특정 행동을 제한하거나 요구하여 위험을 줄임. & Covenants & 대출 조건 \\
\textbf{다각화} & 여러 종류의 자산, 지역, 산업에 분산 투자하여 특정 위험에 대한 노출을 줄이는 전략. & Diversification & 위험 분산 \\
\textbf{신용 파생상품} & 차입자의 부도 위험을 다른 금융기관에 전가하는 금융 상품 (예: CDS). & Credit Derivatives & 위험 전가 수단 \\
\textbf{금리 위험} & 시장 금리 변동으로 인해 은행의 수익성이나 자산 가치가 변동할 위험. & Interest Rate Risk & 금리 변동성 \\
\textbf{만기 불일치} & 은행의 자산(대출) 만기가 부채(예금) 만기보다 길어서 발생하는 구조적 차이. & Maturity Mismatch & 금리 위험의 근원 \\
\textbf{현금 흐름 채널} & 금리 변동이 은행의 이자 수입과 이자 비용에 직접적인 영향을 미쳐 수익성을 변화시키는 경로. & Cash Flow Channel & 단기 수익성 영향 \\
\textbf{평가 채널} & 금리 변동이 은행의 자산과 부채의 시장 가치를 변화시켜 은행의 가치에 영향을 미치는 경로. & Valuation Channel & 장기 가치 영향 \\
\textbf{유동성 위험} & 예금 인출이나 대출 약정 이행 등으로 인해 은행이 갑작스럽게 현금이 부족해질 위험. & Liquidity Risk & 현금 부족 위험 \\
\textbf{뱅크런} & 은행에 대한 신뢰 상실로 인해 많은 예금자가 동시에 예금을 인출하려는 현상. & Bank Run & 대규모 예금 인출 \\
\textbf{헐값 매각} & 유동성 확보를 위해 자산을 시장 가격보다 훨씬 낮은 가격에 급하게 파는 행위. & Fire Sale & 자산 가치 하락 \\
\textbf{시장 위험} & 주식, 채권, 파생상품 등 트레이딩 포트폴리오의 시장 가격 변동으로 인해 손실을 입을 위험. & Market Risk & 트레이딩 손실 위험 \\
\textbf{트레이딩 장부} & 은행이 단기 매매 차익을 목적으로 보유하는 금융 상품 목록. & Trading Book & 단기 투자 자산 \\
\textbf{은행 장부} & 은행이 장기적인 대출 및 투자 목적으로 보유하는 자산 목록. & Bank Book & 장기 투자 자산 \\
\textbf{시가 평가} & 금융 상품의 가치를 매일 시장 가격에 따라 평가하여 손익을 즉시 반영하는 방식. & Mark-to-Market & 시장 가격 반영 \\
\textbf{VaR (Value-at-Risk)} & 특정 신뢰 수준에서 미래 일정 기간 동안 발생할 수 있는 최대 예상 손실 금액. & Value-at-Risk & 최대 손실 예측 \\
\textbf{블랙 스완} & 극히 드물지만 발생하면 엄청난 영향을 미치는 예측 불가능한 사건. & Black Swan & 예측 불가능한 극단적 사건 \\
\textbf{운영 위험} & 내부 프로세스, 인력, 시스템의 부적절하거나 실패, 또는 외부 사건으로 인해 발생하는 손실 위험. & Operational Risk & 내부/외부 실패 위험 \\
\textbf{핀테크} & 금융(Finance)과 기술(Technology)의 합성어로, 정보 기술을 활용한 새로운 금융 서비스. & Fintech & 기술 기반 금융 혁신 \\
\bottomrule
\end{tabular}
\end{adjustbox}


\section{핵심 개념 및 원리}

\subsection{은행 자본과 수익성}

은행의 핵심 목표 중 하나는 이익을 창출하고, 이 이익을 통해 소유주에게 가치를 제공하는 것입니다. 이를 이해하기 위해 은행의 자본 구조와 성과 측정 방법을 살펴봅니다.

\subsubsection{은행 자기자본 (Bank Equity Capital)}

\begin{infobox}[title=\textbf{핵심 요약: 은행 자기자본}]
은행 자기자본은 은행의 소유주가 투자한 자금으로, 은행의 자금 조달원 중 하나이자 손실 발생 시 가장 먼저 손실을 흡수하는 완충재 역할을 합니다.
이는 소유주에게 은행의 이익에 대한 청구권과 의결권을 부여합니다.
\end{infobox}

\textbf{왜 자기자본이 필요한가?}
모든 기업과 마찬가지로 은행도 새로운 프로젝트에 자금을 조달하거나 사업을 확장하기 위해 자금이 필요합니다. 이때 자기자본은 중요한 자금 조달원 중 하나가 됩니다.

\textbf{자기자본의 정의와 역할}
자기자본(Equity Capital)은 은행의 소유권을 나타내는 지분입니다. 은행이 자기자본을 판매한다는 것은 은행의 일부를 외부 투자자에게 판매하는 것을 의미하며, 이 투자자는 은행의 주주(Stockholder 또는 Shareholder)가 됩니다.

주주가 되면 다음과 같은 권리를 갖게 됩니다:
\begin{itemize}
    \item \textbf{이익에 대한 청구권 (Claim on Profits):} 투자에 대한 대가로 주주는 은행 이익의 일부를 받을 자격이 있습니다. 은행이 잘 운영되면 이 이익은 배당금(Dividends) 형태로 지급됩니다.
    \item \textbf{잔여 청구권 (Residual Claim):} 만약 은행이 재정적인 어려움에 처해 매각되거나 폐쇄될 경우, 주주는 모든 부채(Debt)를 상환하고 남은 자산에 대해서만 청구권을 가집니다. 즉, 채권자(Debt holders)가 항상 먼저 변제받으며, 주주는 마지막에 남은 것을 받습니다. 상황이 매우 나쁘면 주주는 투자금을 모두 잃을 수도 있습니다.
    \item \textbf{의결권 (Voting Rights):} 주주는 은행 경영 방식에 영향을 미칠 수 있는 의결권을 가집니다. 이 부분은 중요하지만, 본 논의에서는 주로 재무적 측면에 초점을 맞춥니다.
\end{itemize}

\textbf{실제 사례: JP모건 체이스 (JPMorgan Chase)의 자기자본 (2020년)}
JP모건 체이스는 총 3.4조 달러의 자금을 보유한 대형 은행입니다. 이 중 약 2,800억 달러가 자기자본에서 나옵니다. 자기자본은 주로 보통주(Common Stock)와 우선주(Preferred Stock) 발행을 통해 납입된 자본(Paid-in-capital)으로 구성됩니다. 우선주는 보통주와 유사하지만, 일반적으로 배당금을 먼저 받습니다. JP모건의 자기자본 대부분은 이익잉여금(Retained Earnings)에서 발생합니다. 이익잉여금은 배당으로 지급되지 않고 사업에 재투자된 이익을 의미하며, 궁극적으로 주주에게 돌아가야 할 몫입니다.

\textbf{자기자본의 손실 흡수 역할 (Residual Claim as a Cushion)}
자기자본은 잔여 청구권이기 때문에, 다른 부채 보유자들보다 먼저 손실을 흡수하는 완충재 역할을 합니다.

\begin{examplebox}[title=\textbf{예시: 단순 은행의 청산 및 손실 흡수}]
매우 단순한 은행을 가정해 봅시다.
\begin{itemize}
    \item \textbf{자산:} 장기 투자 (Long-term Investments) = 100달러
    \item \textbf{부채:} 단기 부채 (Short-term Debt, 예금) = 90달러
    \item \textbf{자기자본:} 10달러
\end{itemize}
이 은행은 90달러의 예금과 10달러의 자기자본으로 100달러의 투자를 조달하고 있습니다.

\textbf{정상적인 청산 상황:}
은행 소유주가 은행을 폐쇄하기로 결정했다고 가정합니다.
\begin{itemize}
    \item 자산 100달러를 매각합니다.
    \item 먼저 부채 보유자(예금자)에게 90달러를 상환합니다.
    \item 남은 10달러는 소유주(주주)가 가져갑니다.
\end{itemize}
이 경우, 자기자본 10달러가 그대로 소유주에게 돌아갑니다.

\textbf{경기 침체로 인한 손실 발생 상황:}
경기 침체가 발생하여 일부 차입자가 대출 상환을 하지 못해 대출 가치가 20달러에서 15달러로 감소했다고 가정합니다. 은행의 총 자산 가치는 100달러에서 95달러로 줄어듭니다.
\begin{itemize}
    \item 은행 자산 95달러를 매각합니다.
    \item 부채 보유자(예금자)에게 90달러를 상환합니다. (이들은 여전히 안전합니다.)
    \item 남은 5달러는 소유주(주주)가 가져갑니다.
\end{itemize}
원래 자기자본은 10달러였지만, 손실 5달러를 자기자본이 모두 흡수하여 5달러만 남게 됩니다. 주주는 투자금을 모두 잃지는 않았지만, 손실을 전적으로 부담했습니다.

이것이 바로 자기자본이 '자본 구조의 첫 번째 손실 조각(first loss piece of the capital structure)'이라고 불리는 이유입니다. 자기자본은 부채를 잠재적 손실로부터 보호하는 역할을 합니다. 대부분의 은행 부채는 정부 보증 예금(government-insured deposits)이므로, 규제 당국은 은행이 가능한 한 많은 보호 장치(자기자본)를 갖추기를 원합니다.
\end{examplebox}

\begin{summarybox}[title=\textbf{자기자본 요약}]
은행 자기자본은 은행의 중요한 자금 조달원이며, 대차대조표의 부채 및 자기자본 항목에 부채와 함께 표시됩니다. 주주는 은행의 소유권을 가지며, 이는 의결권과 이익에 대한 청구권을 부여합니다. 자기자본은 잔여 청구권의 성격을 가지므로, 자산에서 손실이 발생할 경우 항상 자기자본이 가장 먼저 손실을 흡수합니다.
\end{summarybox}


\subsubsection{은행 성과 측정 (Measuring Bank Performance)}

\begin{infobox}[title=\textbf{핵심 요약: 은행 성과 측정}]
은행 성과는 주식 수익률, 총자산 수익률(ROA), 자기자본 수익률(ROE)의 세 가지 주요 지표로 측정됩니다. ROA는 자산 활용 효율성을, ROE는 소유주 투자 대비 수익성을 나타내며, 레버리지(부채 사용)는 ROE를 증폭시키는 중요한 요소입니다.
\end{infobox}

은행은 자기자본과 부채를 통해 자금을 조달하고, 이 자금을 대출 및 트레이딩 자산에 투자하여 수익을 창출합니다. 또한 수수료 기반 사업에서도 수익을 얻습니다. 이자 및 비이자 비용을 지불하고 남은 것이 이익(Profits)입니다. 이 이익은 은행 소유주인 주주에게 귀속됩니다.

이 섹션에서는 은행 성과를 측정하는 세 가지 중요한 지표를 다룹니다:
\begin{enumerate}
    \item \textbf{주식 수익률 (Stock Returns)}
    \item \textbf{총자산 수익률 (Return on Assets, ROA)}
    \item \textbf{자기자본 수익률 (Return on Equity, ROE)}
\end{enumerate}
특히, 은행의 차입(Leverage)이 어떻게 이익을 증폭시키는지 논의합니다.

\textbf{1. 주식 수익률 (Stock Returns)}
\begin{itemize}
    \item \textbf{개념:} 주식 수익률은 투자자가 주식에 투자하여 얻는 수익을 나타냅니다. 주가 상승과 배당금 재투자를 포함합니다.
    \item \textbf{장점:} 시장의 평가를 직접적으로 반영하며, 투자자에게 가장 직관적인 성과 지표입니다.
    \item \textbf{단점:}
        \begin{itemize}
            \item 모든 은행이 상장되어 있는 것은 아닙니다. 비상장 은행의 경우 주식 수익률을 측정할 수 없습니다.
            \item 주식 수익률만으로는 은행이 왜 잘하거나 못하는지에 대한 근본적인 이유를 파악하기 어렵습니다.
        \end{itemize}
    \item \textbf{예시: JP모건 체이스 (2015-2020)}
        2015년 12월 31일에 JP모건 보통주에 100달러를 투자하고 배당금을 재투자했다면, 5년 후 225달러에 주식을 팔 수 있었을 것입니다. 이는 JP모건이 KBW 은행 지수(다른 미국 상장 은행 포함)와 더 넓은 금융 및 비금융 기업 지수보다 우수한 성과를 보였음을 의미합니다.
\end{itemize}

\textbf{2. 총자산 수익률 (Return on Assets, ROA)}
\begin{itemize}
    \item \textbf{개념:} ROA는 은행이 자산을 얼마나 효율적으로 사용하여 세후 이익을 창출했는지를 나타내는 지표입니다. 투자된 자산 1달러당 얼마의 이익을 얻었는지를 보여줍니다.
    \item \textbf{계산:}
    \begin{equation*}
        \text{ROA} = \frac{\text{세후 이익 (After-Tax Profits)}}{\text{총자산 (Total Assets)}}
    \end{equation*}
    \item \textbf{영향 요인:}
        \begin{itemize}
            \item \textbf{고수익 자산 투자:} 높은 수익을 내고 손실이 적은 자산에 투자하는 은행.
            \item \textbf{저비용 차입:} 낮은 비용으로 자금을 조달하는 은행.
            \item \textbf{엄격한 운영:} 운영 비용을 효율적으로 통제하는 은행 (예: 오버헤드 비율).
            \item \textbf{세금 최소화:} 세금 부담을 최소화하는 은행.
            \item \textbf{신용 품질:} 대출 포트폴리오의 잠재적 손실 (예: 대손충당금 비율)이 낮을수록 ROA에 긍정적입니다.
        \end{itemize}
    \item \textbf{장점:} 상장 여부와 관계없이 모든 은행에 대해 계산할 수 있습니다. 은행의 내부 효율성을 평가하는 데 유용합니다.
\end{itemize}

\textbf{3. 자기자본 수익률 (Return on Equity, ROE)}
\begin{itemize}
    \item \textbf{개념:} ROE는 은행 소유주(주주)의 투자에 대한 수익률을 나타냅니다. 주주가 투자한 자기자본 1달러당 얼마의 이익을 얻었는지를 보여줍니다.
    \item \textbf{계산:}
    \begin{equation*}
        \text{ROE} = \frac{\text{세후 이익 (After-Tax Profits)}}{\text{자기자본 (Equity)}}
    \end{equation*}
    \item \textbf{레버리지(Leverage)와의 관계:} ROE는 ROA와 레버리지 비율(Leverage Ratio)의 곱으로 다시 쓸 수 있습니다.
    \begin{equation*}
        \text{ROE} = \text{ROA} \times \text{레버리지 비율 (Leverage Ratio)}
    \end{equation*}
    여기서 레버리지 비율은 $\frac{\text{총자산 (Total Assets)}}{\text{자기자본 (Equity)}}$ 입니다.
    \item \textbf{레버리지의 역할:}
        \begin{itemize}
            \item 은행이 자기자본 외에 부채를 사용하여 자산을 조달하는 정도를 나타냅니다.
            \item 레버리지 비율이 높을수록 자기자본 1달러당 더 많은 자산을 운용하게 되므로, ROA가 동일하더라도 ROE는 더 높아집니다.
            \item 은행이 자기자본으로만 자금을 조달한다면 레버리지 비율은 1이 되고, ROE와 ROA는 같아집니다.
            \item \textbf{위험:} 레버리지는 좋은 시기에 ROE를 증폭시키지만, 나쁜 시기에는 손실을 증폭시켜 은행의 취약성을 높입니다.
        \end{itemize}
    \item \textbf{장점:} 주주 관점에서 은행의 수익성을 평가하는 가장 중요한 지표입니다. 상장 여부와 관계없이 모든 은행에 대해 계산할 수 있습니다.
\end{itemize}

\begin{examplebox}[title=\textbf{예시: JP모건 체이스의 ROA와 ROE (2020년)}]
JP모건 체이스의 ROA는 0.91\%로 그리 높아 보이지 않습니다. 하지만 자기자본 수익률(ROE)은 12\%로 훨씬 높습니다. 이 큰 차이는 바로 '레버리지' 때문입니다.
\begin{itemize}
    \item 총자산과 보통주 자기자본을 비교하면, 약 13배의 레버리지 비율을 가집니다.
    \item 이는 보통주 주주가 1달러를 투자할 때, 은행이 12달러의 부채를 빌려 총 13달러의 자산을 운용한다는 의미입니다.
    \item 이 13배의 레버리지가 ROA(0.91\%)를 증폭시켜 ROE(12\%)를 만들어냅니다.
\end{itemize}
ROE는 JP모건의 주요 성과 지표 중 하나입니다.
\end{examplebox}

\textbf{미국 은행의 성과 추이 (지난 25년간)}
미국 연방예금보험공사(FDIC)의 재무제표 데이터를 기반으로 한 지난 25년간 미국 은행의 평균 성과를 살펴보면 다음과 같습니다.
\begin{itemize}
    \item \textbf{평균 ROA:} 2008년 금융 위기 이전에는 약 1.2\%였으나, 2010년에서 2020년 사이에는 약 1\%로 낮아졌습니다.
    \item \textbf{대출 손실 충당금 비율 (Loan Loss Reserves to Total Loans):} 이 비율은 대출 포트폴리오의 잠재적 손실을 나타냅니다. 이 비율이 급증하면 대출 상황이 좋지 않다는 의미이며, 은행 자산의 대부분이 대출로 구성되므로 ROA와 명확한 음의 상관관계를 보입니다. 즉, 대출 손실이 많아지면 ROA는 감소합니다.
    \item \textbf{평균 ROE:} 좋은 시기(ROA가 1.2\% 정도일 때)에는 ROE가 15\%까지 높게 나타납니다. 이는 약 10~11배의 높은 레버리지를 사용했음을 시사합니다.
    \item \textbf{레버리지의 양면성:} 높은 레버리지는 좋은 시기에 ROE를 크게 증폭시키지만, 2008년과 2020년과 같은 위기 시에는 손실을 훨씬 더 극적으로 증폭시켜 ROE를 급락시킵니다.
\end{itemize}

\begin{summarybox}[title=\textbf{은행 성과 측정 요약}]
은행 성과를 측정하는 세 가지 중요한 지표는 주식 수익률, 총자산 수익률(ROA), 자기자본 수익률(ROE)입니다. 시장 기반 지표인 주식 수익률은 상장 은행에만 적용 가능하며, 회계 기반 지표인 ROA와 ROE는 모든 은행에 적용 가능합니다. 이 지표들은 은행이 어떻게 이익을 창출하는지 (훌륭한 투자, 비용 통제, 높은 레버리지 등)를 알려줍니다. 하지만 이 모든 지표는 '위험' 요소를 간과하고 있습니다. 15\%의 ROE가 좋아 보일 수 있지만, 그 수익이 얼마나 취약한지는 별개의 문제입니다. 투자자와 은행 관리자는 위험을 어떻게 고려하고 관리해야 할까요? 이것이 다음 논의의 핵심입니다.
\end{summarybox}


\section{은행 위험 및 위험 관리 개요}

은행 사업 모델은 본질적으로 위험을 수반합니다. 이 섹션에서는 은행이 직면하는 다양한 위험의 종류와 이를 관리하는 기본적인 접근 방식을 소개합니다.

\subsection{위험 개요 (Overview of Risk)}

\begin{infobox}[title=\textbf{핵심 요약: 위험의 정의와 종류}]
위험은 미래의 불확실한 결과로 인해 잠재적 손실이 발생할 가능성을 의미합니다. 은행은 신용, 금리, 유동성, 시장, 운영 위험 등 다양한 위험에 노출되어 있으며, 이러한 위험을 측정하고 관리하는 것이 은행 성과에 매우 중요합니다.
\end{infobox}

\textbf{위험이란 무엇인가?}
금융 의사결정을 할 때, 미래에 발생할 결과에 대한 불확실성(Uncertainty)이 항상 존재합니다. 예를 들어, 오늘 투자한 돈이 내일 얼마의 수익을 가져올지 확실히 알 수 없습니다.
\begin{itemize}
    \item \textbf{잠재적 결과 (Potential Outcomes):} 우리는 무엇이 일어날지 정확히 알 수 없지만, 잠재적인 결과들을 나열할 수는 있습니다. 어떤 결과는 매우 좋을 수 있고, 어떤 결과는 바람직하지 않을 수 있습니다 (예: 대출 포트폴리오에서 손실 발생).
    \item \textbf{확률 할당 (Assigning Probabilities):} 이러한 잠재적 결과들에 확률을 할당하려고 시도할 수 있습니다. 어떤 사건은 발생 확률이 0일 수도 있습니다. 사건의 확률이나 빈도를 추정하는 것은 쉬운 일이 아니지만, 이는 위험 측정의 핵심입니다.
\end{itemize}
은행은 다양한 활동에 참여하며, 이러한 활동 중 일부는 잠재적 손실(Potential Losses)에 노출됩니다. 이러한 손실은 미래에 발생할 수 있으며, 그 규모와 발생 여부는 불확실합니다. 매우 작은 손실은 발생할 가능성이 낮을 수 있지만, 다른 손실은 더 높은 확률을 가질 수 있습니다. 이러한 잠재적 손실을 측정하고 관리하는 것이 은행 성과의 핵심입니다.

\textbf{현대 은행이 직면하는 주요 위험}
은행 관리자들이 직면하는 주요 위험들은 다음과 같습니다. 각 위험에 대해서는 이후 세션에서 더 자세히 다룹니다.

\begin{enumerate}
    \item \textbf{신용 위험 (Credit Risk):}
        \begin{itemize}
            \item \textbf{정의:} 은행이 대출이나 증권과 같은 자산에 투자했을 때, 차입자가 원금이나 이자를 제때 갚지 못할 위험입니다.
            \item \textbf{예시:} 기업 대출이나 주택 담보 대출을 받은 고객이 파산하여 대출금을 상환하지 못하는 경우.
        \end{itemize}
    \item \textbf{금리 위험 (Interest Rate Risk):}
        \begin{itemize}
            \item \textbf{정의:} 은행은 단기 자금을 조달하여 장기 자산에 투자하는 '만기 불일치(Maturity Mismatch)' 구조를 가집니다. 시장 금리가 변동할 때, 이 만기 불일치로 인해 은행의 수익성이나 자산 가치가 변동할 위험입니다.
            \item \textbf{예시:} 단기 예금 금리는 오르는데, 장기 고정 금리 대출의 수익은 그대로여서 은행의 이자 마진이 줄어드는 경우.
        \end{itemize}
    \item \textbf{유동성 위험 (Liquidity Risk):}
        \begin{itemize}
            \item \textbf{정의:} 예금자들이 갑자기 대규모로 예금을 인출하거나, 은행이 약정한 대출을 이행해야 할 때, 은행이 충분한 현금을 확보하지 못할 위험입니다.
            \item \textbf{예시:} 뱅크런(Bank Run)이 발생하여 예금자들이 한꺼번에 돈을 인출하려 할 때, 은행이 장기 자산을 헐값에 팔아야 하는 경우.
        \end{itemize}
    \item \textbf{시장 위험 (Market Risk):}
        \begin{itemize}
            \item \textbf{정의:} 은행의 트레이딩 장부(Trading Book)에 있는 증권, 파생상품 등의 시장 가격이 변동하여 손실을 입을 위험입니다.
            \item \textbf{예시:} 주식 시장이 급락하여 은행이 보유한 주식 포트폴리오의 가치가 크게 하락하는 경우.
        \end{itemize}
    \item \textbf{운영 위험 (Operational Risks):}
        \begin{itemize}
            \item \textbf{정의:} 내부 프로세스, 인력, 시스템의 부적절하거나 실패, 또는 외부 사건(예: 사이버 공격, 자연재해)으로 인해 발생하는 광범위한 손실 위험입니다.
            \item \textbf{예시:} 직원의 사기, 기술 시스템 오류, 사이버 공격으로 인한 데이터 유출, 핀테크 기업의 시장 잠식.
        \end{itemize}
    \item \textbf{지급불능 위험 (Insolvency Risk):}
        \begin{itemize}
            \item \textbf{정의:} 위에서 언급된 하나 이상의 위험으로 인해 손실이 너무 커져 은행의 자기자본이 모두 소진되고, 부채 보유자(예금자 등)까지 손실을 입게 되는 위험입니다.
            \item \textbf{예시:} 손실이 자기자본(X)을 초과하여 은행이 파산하는 경우.
        \end{itemize}
\end{enumerate}

\textbf{위험 통제 및 관리 (Risk Control and Management)}
은행 관리자와 소유주는 지급불능 위험을 피하기 위해 '위험 통제(Risk Control)' 또는 '위험 관리(Risk Management)'라는 광범위한 작업을 수행합니다. 이 과정은 크게 두 단계로 나뉩니다.

\begin{enumerate}
    \item \textbf{위험 노출 측정 (Measure Risk Exposures):}
        \begin{itemize}
            \item 은행이 어떤 위험에 얼마나 노출되어 있는지 정확히 파악하는 것이 첫 번째 단계입니다.
            \item \textbf{예시:} 은행의 이익이 금리 변동에 따라 어떻게 달라지는지 측정합니다. 만약 금리가 오르면 이익이 감소하는 경향이 있다면, 이는 금리 위험에 노출되어 있다는 의미입니다.
        \end{itemize}
    \item \textbf{위험 노출 관리 (Manage Risk Exposures):}
        \begin{itemize}
            \item 측정된 위험 노출을 줄이거나 제거하기 위한 조치를 취합니다.
            \item \textbf{예시:} 금리 변동에 따른 이익 변동성을 줄이기 위해 파생상품을 사용하거나, 자산-부채 만기 구조를 조정합니다. 목표는 이익과 위험원 간의 연결 고리를 끊는 것입니다.
        \end{itemize}
\end{enumerate}

\textbf{위험 관리의 구체적인 방법}
은행은 위험을 관리하기 위해 다음과 같은 방법을 사용합니다.
\begin{itemize}
    \item \textbf{위험 측정 및 가격 책정 (Measure and Price Risks):}
        \begin{itemize}
            \item 은행은 항상 자신이 감수하는 위험을 측정하고 그에 합당한 가격을 책정해야 합니다.
            \item \textbf{예시:} 신뢰할 수 없는 고객에게는 더 높은 이자율의 대출을 제공해야 합니다. 위험이 너무 크다면 아예 대출을 거부하거나, 대출을 실행한 후 다른 기관에 판매(Originate and Sell)하여 대차대조표에서 위험 노출을 직접적으로 줄일 수 있습니다.
        \end{itemize}
    \item \textbf{위험 헤지 (Hedge Risks):}
        \begin{itemize}
            \item 위험을 다른 금융기관으로 이전하는 것을 의미합니다.
            \item \textbf{예시:} 대출 보증(Loan Guarantees), 금리 파생상품(Interest Rate Derivatives) 등을 사용하여 고객의 요구(예: 위험한 장기 대출 제공)를 충족시키면서도 은행이 파산할 위험을 제한할 수 있습니다.
        \end{itemize}
\end{itemize}
궁극적인 목표는 은행의 이익과 특정 위험원(예: 금리) 사이의 연결 고리를 끊는 것입니다.

\begin{summarybox}[title=\textbf{위험 개요 요약}]
위험은 미래의 불확실한 결과로 인한 잠재적 손실을 의미합니다. 은행은 신용, 금리, 유동성, 시장, 운영 위험 등 다양한 특정 위험에 직면합니다. 이러한 위험을 통제하기 위해 은행은 최소한 위험을 측정하고 가격을 책정해야 합니다. 나아가 대차대조표를 직접 조정하거나 헤징을 통해 위험 노출을 제한할 수 있습니다.
\end{summarybox}


\subsection{신용 위험 (Credit Risk)}

\begin{infobox}[title=\textbf{핵심 요약: 신용 위험}]
신용 위험은 차입자가 대출 원리금을 제때 상환하지 못할 위험입니다. 은행은 신용 위험을 측정하기 위해 차입자 데이터를 분석하고 신용 등급을 부여하며, 대출 계약 설계, 다각화, 대출 판매, 신용 파생상품 등을 통해 신용 위험을 관리합니다.
\end{infobox}

은행은 자금을 빌리고, 빌려주고, 그 차이(스프레드)로 이익을 얻습니다. 이 스프레드는 자산에서 발생하는 수익과 자금 조달 비용의 차이에서 나옵니다. 이 과정이 성공하려면 자산에서 발생하는 수익이 실제로 실현되어야 합니다. 즉, 차입자가 대출금을 상환해야 합니다. 이것이 바로 신용 위험의 핵심입니다.

\textbf{신용 위험의 정의}
신용 위험은 은행이 대출을 제공하거나 채권을 구매했을 때, 차입자가 원금이나 이자를 약속된 대로 상환하지 못할 위험을 의미합니다.
\begin{itemize}
    \item \textbf{대출 자산:} 은행의 대출 및 채권은 일반적으로 만기가 길어, 수개월에서 수년에 걸쳐 이자와 원금을 회수할 것으로 예상합니다.
    \item \textbf{부도 (Default):} 대부분의 경우, 계획대로 모든 상환이 이루어지지만, 일부 차입자가 한 번 이상 상환을 놓칠 수 있습니다. 이를 '부도'라고 합니다. 부도는 발생 확률이 낮은 사건이지만, 발생하면 차입자에게는 파산으로 이어질 수 있는 매우 큰 비용을 초래하며, 은행에게도 자산 가치 하락과 수익 감소로 이어지는 해로운 결과를 가져옵니다.
\end{itemize}

\textbf{신용 위험의 영향: 대출 상각률 (Charge-off Rates)}
대출 상각률은 상환되지 않아 은행 대차대조표에서 제거되어야 하는 대출의 비율을 나타냅니다.
\begin{itemize}
    \item \textbf{상품별 차이:}
        \begin{itemize}
            \item \textbf{주택 담보 대출 (Mortgages):} 사람들이 집에 살아야 하므로 일반적으로 부도율이 가장 낮습니다.
            \item \textbf{신용 카드 (Credit Cards):} 무담보이며 상환을 무시하기 쉽기 때문에 평균적으로 부도율이 가장 높습니다.
            \item \textbf{사업 대출 (Business Loans):} 중간 정도의 부도율을 보입니다.
        \end{itemize}
    \item \textbf{경기 순환적 요소 (Business Cycle Component):} 모든 대출 상품의 부도율은 경기 침체기에 급증하는 경향이 있습니다. 이는 경기 침체 시 가계와 기업이 동시에 재정적 어려움을 겪기 때문이며, 이를 '체계적 신용 위험(Systematic Credit Risk)'이라고 합니다.
\end{itemize}

\textbf{실제 사례: JP모건 체이스의 신용 위험 노출 (2019-2020년)}
\begin{itemize}
    \item \textbf{대차대조표상 대출 (On-balance sheet loans):} 일부 차입자는 상환을 놓쳐 '부실 대출(Non-performing bucket)'로 분류될 수 있습니다. 2020년에는 대출 포트폴리오의 약 0.5\%가 상각되었습니다.
    \item \textbf{대차대조표 외 노출 (Off-balance sheet exposures):}
        \begin{itemize}
            \item \textbf{대출 약정 (Loan Commitments):} 신용 카드와 같이 차입자가 언제든지 자금을 인출할 수 있는 잠재적 자산입니다. 신용 카드는 부도 위험이 높으므로, JP모건은 이러한 노출을 공개해야 합니다.
        \end{itemize}
\end{itemize}

\subsubsection{신용 위험 측정 (Measuring Credit Risk)}

은행은 신용 위험을 측정하고 관리합니다. 측정 부분은 종종 '신용 위험 분석(Credit Risk Analysis)' 또는 '차입자 심사(Borrower Screening)'라고 불립니다. 이는 차입자가 대출을 신청할 때 초기 단계에서 이루어집니다. 대출 기관은 잠재적 차입자에 대한 정보를 생성하고 처리하는 데 자원을 투입합니다.

\textbf{신용 위험 측정 과정}
\begin{enumerate}
    \item \textbf{차입자 데이터 입력 (Borrower Data Input):}
        \begin{itemize}
            \item \textbf{신청서:} 대출 담당자는 잠재적 차입자로부터 소득, 자산, 연령, 직업 등의 정보를 담은 서면 또는 온라인 신청서를 받습니다.
            \item \textbf{신용 정보 기관 (Credit Bureau):} 차입자의 기존 부채 및 과거 상환 행동에 대한 데이터를 제공합니다.
            \item \textbf{대안 데이터 (Alternative Data):} 21세기에는 지출 데이터, 인터넷 브라우저 기록, 위치 데이터, 심지어 소셜 네트워크 정보(예: 부도율이 높은 사람들과 페이스북 친구인지 여부)와 같은 다양한 대안 데이터 옵션이 있습니다. 이러한 대안 데이터의 윤리적 측면은 현재 논의 중입니다.
        \end{itemize}
    \item \textbf{신용 위험 모델 (Credit Risk Model):}
        수집된 데이터는 대출 기관의 신용 위험 모델에 입력됩니다. 이 모델은 다양한 형태를 가질 수 있습니다.
        \begin{itemize}
            \item \textbf{정성적 모델 (Qualitative Models):} 차입자의 특정 상황에 대한 일련의 질문(예: 고객이 왜 대출을 받으려 하는가? 무엇이 필요한가? 어떻게 상환할 것인가? 대출 기관은 어떻게 보호받을 수 있는가?)을 따릅니다.
            \item \textbf{통계적 모델 (Statistical Models):} 관찰 가능한 특성을 기반으로 차입자를 위험 범주로 분류합니다. 단순 회귀 분석부터 최첨단 머신러닝(Machine Learning) 방법까지 다양합니다. 이 모델들은 과거 차입자 데이터를 사용하여 현재 신청자의 성과를 예측합니다.
        \end{itemize}
    \item \textbf{신용 등급 출력 (Credit Rating Output):}
        모델의 결과는 '신용 등급(Credit Rating)'입니다. 이 신용 등급은 '부도 확률(Probability of Default)'에 해당합니다.
        \begin{itemize}
            \item \textbf{기업 신용 등급:} 예를 들어, 투자 등급(Investment Grade)의 최하위인 트리플 B(BBB) 등급은 해당 기업의 부도 확률이 0.2\%임을 의미할 수 있습니다.
            \item \textbf{개인 신용 등급 (FICO):} 개인도 신용 등급을 가집니다. FICO(Fair Isaac Corporation) 시스템은 300점에서 850점까지의 점수를 부여하며, 점수가 높을수록 부도 확률이 낮습니다. FICO는 신용 카드, 자동차 대출, 주택 담보 대출, 학자금 대출, 청구서 납부 기록 등의 데이터를 사용하여 독점적인 부도 예측 모델로 점수를 산출합니다.
        \end{itemize}
\end{enumerate}

\subsubsection{신용 위험 관리 (Managing Credit Risk)}

신용 위험을 측정하고 차입자의 신용 등급을 알게 되었다면, 다음 단계는 이 위험을 관리하는 것입니다. 부도는 대출 사업의 피할 수 없는 부분이며, 대출 기관의 목표는 손실을 최소화하고 적절한 보상을 받는 것입니다.

\textbf{1. 대출 계약 설계 (Loan Contract Design)}
\begin{itemize}
    \item \textbf{대출 가격 책정 (Pricing the Loan):} 신용 등급을 사용하여 대출 이자율을 책정합니다. 위험이 높은 차입자는 부도 가능성이 높으므로 더 높은 이자율을 부과하여 잠재적 신용 손실에 대한 보상을 받습니다.
    \item \textbf{약정 (Covenants):} 대출 계약에 성과 관련 조항을 포함하여 부도 확률을 줄이는 특정 행동을 제한하거나 장려합니다.
        \begin{itemize}
            \item \textbf{예시:} 기업 차입자가 발행할 수 있는 신규 부채의 양을 제한하는 약정.
            \item \textbf{위반 시:} 약정 위반 시 차입자는 대출 기관에 보고해야 하며, 문제가 악화되기 전에 대출 기관과 차입자가 대출 조건을 재협상(Workout)할 수 있습니다.
        \end{itemize}
    \item \textbf{실제 사례: JP모건 체이스의 고객 지원 프로그램 (2020년)}
        코로나19 팬데믹으로 인해 많은 고객이 대출 상환에 어려움을 겪었습니다. JP모건은 약 100억 달러 규모의 주택 담보 대출에 대해 고객 지원 프로그램을 제공했습니다. 법적으로는 즉시 상환을 요구할 수 있었지만, 손실 최소화를 위해 단기적인 유예를 제공하는 것이 더 현명했습니다. 많은 경우, JP모건은 미납된 상환금을 대출 만기 시점까지 연기해 주었으며, 그 결과 약 95\%의 차입자가 다시 정상 궤도로 돌아왔습니다.
    \item \textbf{담보 및 우선순위 구조 (Security or Priority Structure):}
        \begin{itemize}
            \item \textbf{담보 대출 (Secured Loan):} 차입자가 부도 시 대출 기관이 청구할 수 있는 자산(담보)으로 뒷받침되는 대출입니다.
            \item \textbf{선순위 대출 (Senior Loan):} 차입자가 부도 시 자산에 대한 첫 번째 청구권을 가지는 대출입니다.
            \item \textbf{효과:} 선순위 및 담보 청구권을 가지면 부도 시 손실을 줄일 수 있습니다. 무디스(Moody's) 데이터에 따르면, 선순위 담보 채권은 후순위 채권보다 더 높은 회수율(Recovery Rate, 부도 시 회수되는 금액)을 보입니다.
        \end{itemize}
\end{itemize}

\textbf{2. 다각화 (Diversification)}
\begin{itemize}
    \item \textbf{개념:} 은행이 다양한 상품, 지역, 산업에 걸쳐 여러 자산에 투자하는 가장 기본적인 보호 수단입니다.
    \item \textbf{효과:} 다각화는 신용 위험을 완전히 제거하지는 못하지만, 자산 포트폴리오에서 매우 나쁜 결과가 발생할 가능성을 제한할 수 있습니다. 주어진 목표 수익률에 대해 위험을 최소화하는 투자 조합을 선택할 수 있습니다.
\end{itemize}

\textbf{3. 대출 판매 (Loan Sales)}
\begin{itemize}
    \item \textbf{개념:} 신용 위험이 너무 높다고 판단될 경우, 은행은 단순히 대출을 다른 기관에 판매하여 신용 위험을 제거할 수 있습니다.
    \item \textbf{시기:} 대출 실행 시점 또는 그 이후에 이루어질 수 있습니다.
    \item \textbf{형태:} 전체 대출을 판매하거나, 대출 신디케이션(Loan Syndication)을 통해 부분적으로 판매할 수 있습니다. 이는 대규모 기업 대출이나 상업용 부동산 대출과 같이 위험이 큰 경우에 흔히 사용됩니다.
    \item \textbf{수익:} 대출을 실행하고 판매하는 경우, 은행은 이자 수입의 일부 또는 전부를 포기하는 대신 선불 수수료(Upfront Fees)를 받습니다.
\end{itemize}

\textbf{4. 신용 파생상품 (Credit Derivatives)}
\begin{itemize}
    \item \textbf{개념:} 차입자의 부도에 대한 보험과 유사하게 작동하는 파생상품입니다.
    \item \textbf{예시: 신용부도스왑 (Credit Default Swaps, CDS):} 가장 일반적인 예시입니다. 은행은 수수료를 지불하고 대출에 대한 CDS를 구매할 수 있습니다. 이 파생상품은 대출이 부도날 때마다 지급금을 제공할 수 있습니다. 은행은 대출에서 얻는 이익의 일부를 포기하지만, 부도 위험에 대해 더 이상 걱정할 필요가 없습니다.
\end{itemize}

\begin{summarybox}[title=\textbf{신용 위험 요약}]
신용 위험은 대출이 제때 완전히 상환되지 않을 위험입니다. 신용 위험 모델은 차입자 데이터를 입력받아 부도 확률에 해당하는 신용 등급을 산출합니다. 대출 기관은 신중하게 설계된 계약(가격 책정, 약정, 담보), 다각화, 그리고 필요하다면 대출 판매나 신용 파생상품 사용을 통해 신용 위험으로부터 자신을 보호할 수 있습니다.
\end{summarybox}


\subsection{금리 위험 (Interest Rate Risk)}

\begin{infobox}[title=\textbf{핵심 요약: 금리 위험}]
금리 위험은 시장 금리 변동으로 인해 은행의 수익성(현금 흐름 채널)이나 가치(평가 채널)가 변동할 위험입니다. 이는 주로 은행의 자산과 부채 간의 '만기 불일치'에서 발생하며, 은행은 만기 조절, 파생상품, 비이자 수익 다각화 등으로 이를 관리합니다.
\end{infobox}

은행은 자금을 조달하고(예금, 자본 시장 차입), 그 자금의 대부분을 대출과 채권에 투자하여 스프레드를 얻습니다. 여기서 중요한 점은 은행이 일반적으로 '단기 차입(Borrow Short)'하여 '장기 대출(Lend Long)'한다는 것입니다. 즉, 은행의 투자는 부채보다 만기가 더 긴 경향이 있습니다. 예를 들어, 주택 담보 대출은 20년 만기일 수 있지만, 예금은 1년 미만의 만기를 가지며, 요구불 예금(Checking Accounts)은 만기가 0일 수 있습니다. 이 섹션에서는 이러한 자산과 부채 간의 '만기 불일치(Maturity Mismatch)'가 어떻게 금리 위험을 초래하는지 논의합니다.

\textbf{금리 변동성 (Interest Rate Volatility)}
금리는 본질적으로 변동성이 큽니다. 미국 연방기금 금리(Federal Funds Rate) 데이터를 보면 1980년대에는 큰 변동이 있었고, 2009년 이후에는 이례적으로 낮고 안정적이었습니다. 이러한 금리 변동의 불확실성은 은행에게 위험을 의미합니다. 우리의 목표는 금리 변동이 은행의 현금 흐름과 가치 평가에 어떻게 영향을 미 미치는지 이해하는 것입니다. 여기서 자산과 부채 간의 만기 불일치가 핵심적인 역할을 합니다.

\textbf{금리 위험이 은행에 영향을 미치는 두 가지 채널}

\textbf{1. 현금 흐름 채널 (Cash Flow Channel)}
이 채널은 은행의 이자 수입(Interest Income)과 이자 비용(Interest Expenses)에 초점을 맞춥니다.
\begin{itemize}
    \item \textbf{이자 수입:} 은행의 대출 및 채권 투자 수익은 금리에 따라 달라집니다. 금리가 높으면 대출에서 더 높은 이자 수입을 얻을 수 있어 은행에 좋은 소식입니다.
        \begin{itemize}
            \item \textbf{예시:} 2000년에 실행된 30년 고정 금리 주택 담보 대출은 2010년이나 2020년에 실행된 대출보다 더 높은 이자 수입을 가져올 것입니다.
        \end{itemize}
    \item \textbf{이자 비용:} 은행은 종종 차입을 통해 투자를 조달합니다. 금리가 높으면 차입 비용(이자 비용)이 증가하여 은행에 나쁜 소식일 수 있습니다.
        \begin{itemize}
            \item \textbf{예시:} 은행이 발행하는 3개월 만기 양도성 예금증서(Certificates of Deposit, CD)의 금리가 오르면, 은행의 이자 비용이 증가합니다.
        \end{itemize}
    \item \textbf{순 이자 수입 (Net Interest Income):} 은행이 금리 상승 시 더 나은지 나쁜지는 은행의 자산과 부채의 만기 구조에 따라 달라집니다.
        \begin{itemize}
            \item \textbf{예시: 단기 차입-장기 대출 은행}
                \begin{itemize}
                    \item \textbf{자산:} 100달러, 2년 만기, 연 5\% 고정 금리.
                    \item \textbf{부채:} 100달러, 1년 만기, 연 4\% 변동 금리 (1년 후 재융자 필요).
                    \item \textbf{1년차 순 이자 수입:} 5달러 (수입) - 4달러 (비용) = 1달러.
                    \item \textbf{2년차 순 이자 수입:} 1년 후 차입 금리가 4\%에서 5\%로 오르면, 순 이자 수입은 5달러 - 5달러 = 0달러가 됩니다. 5\%보다 더 오르면 은행은 손실을 입습니다.
                \end{itemize}
                이것이 '재융자 위험(Refinancing Risk)'입니다.
        \end{itemize}
    \item \textbf{결론:} 현금 흐름 관점에서 은행의 금리 위험 노출은 자산과 부채의 만기 시점, 그리고 고정 금리인지 변동 금리인지에 따라 달라집니다. 변동 금리 대출은 금리 변화에 즉시 반응하지만, 고정 금리 장기 자산은 금리 변화에 둔감합니다.
\end{itemize}

\textbf{2. 평가 채널 (Valuation Channel)}
이 채널은 금리 변동이 은행의 가치 평가에 어떻게 영향을 미치는지 설명합니다.
\begin{itemize}
    \item \textbf{채권 가격 결정 원리:} 채권의 시장 가치는 미래 현금 흐름(쿠폰 지급 및 만기 시 원금)을 현재 금리로 할인하여 결정됩니다.
    \begin{equation*}
        \text{채권 가치} = \sum_{t=1}^{N} \frac{\text{쿠폰 지급}_t}{(1+r)^t} + \frac{\text{원금}}{(1+r)^N}
    \end{equation*}
    여기서 $r$은 이자율입니다. 이자율이 상승하면 할인율이 높아지므로 채권의 가치는 감소합니다. 만기가 긴 채권($N$이 큰)일수록 금리 변화에 더 민감하게 반응하여 가치 하락 폭이 커집니다. 반대로 만기가 짧은 자산은 금리 변화에 덜 민감합니다.
    \item \textbf{은행의 가치 평가:} 은행은 자산 측면에서 채권 포트폴리오로 생각할 수 있습니다. 예금 및 기타 차입금과 같은 부채도 개념적으로 채권과 유사하게 평가할 수 있습니다.
        \begin{itemize}
            \item \textbf{은행 자기자본 가치:} 주주에게 귀속되는 가치는 '자산의 시장 가치 - 부채의 시장 가치'입니다.
            \item \textbf{금리 상승 시 영향:}
                \begin{itemize}
                    \item 채권 가격 결정 원리에 따라 자산 포트폴리오의 가치는 하락합니다.
                    \item 부채의 가치도 하락합니다.
                    \item \textbf{전체 영향:} 은행은 일반적으로 자산의 만기가 부채보다 길기 때문에, 자산이 부채보다 금리 변화에 더 민감합니다. 따라서 금리가 상승하면 자산 가치 하락 폭이 부채 가치 하락 폭보다 커져, 은행의 자기자본 가치는 일반적으로 감소합니다. 즉, 금리 상승은 은행에 나쁜 소식인 경우가 많습니다.
                \end{itemize}
        \end{itemize}
\end{itemize}

\textbf{금리 위험 관리 (Controlling Interest Rate Risk)}
만기 불일치가 금리 위험의 근본 원인이라면, 은행 관리자는 이 격차를 줄이려고 노력할 수 있습니다.
\begin{itemize}
    \item \textbf{만기 일치 (Closing Maturity Mismatches):}
        \begin{itemize}
            \item \textbf{예시:} 은행이 5년 만기 CD를 발행했다면, 그 자금으로 5년 만기 주택 담보 대출에 투자하는 것입니다.
            \item \textbf{문제점:} 하지만 이는 고객의 요구와 상충될 수 있습니다. 고객은 편리한 요구불 예금(0일 만기)을 선호하고, 차입자는 장기 대출을 선호합니다. 만기 불일치는 대출 사업의 근본적인 특성입니다.
        \end{itemize}
    \item \textbf{파생상품 사용 (Use of Derivatives):}
        \begin{itemize}
            \item \textbf{개념:} 금리 파생상품을 사용하여 변동성 있는 금리 위험을 제3자에게 이전할 수 있습니다.
            \item \textbf{예시: 금리 스왑 (Interest Rate Swaps):} 가장 인기 있는 상품으로, 은행이 변동 금리 지급 흐름을 고정 금리 지급 흐름으로 교환할 수 있게 해줍니다. 이를 통해 은행은 단기 예금이나 장기 대출을 원하는 고객의 요구를 충족시키면서도 금리 위험을 제거할 수 있습니다.
        \end{itemize}
    \item \textbf{비이자 수익 다각화 (Diversification into Noninterest Income):}
        \begin{itemize}
            \item \textbf{개념:} 1990년대 후반부터 은행들은 상업 은행 사업 외의 비이자 수익(수수료 수입)을 강조해 왔습니다.
            \item \textbf{예시:} 투자 은행 업무, 보험, 인수 업무, 중개 서비스 등. 이러한 사업 라인에서 발생하는 이익은 금리 변화에 덜 민감할 수 있습니다.
        \end{itemize}
\end{itemize}

\begin{summarybox}[title=\textbf{금리 위험 요약}]
금리는 불확실하며, 금리 변화는 은행의 이자 수입, 이자 비용, 그리고 궁극적으로 은행 이익에 직접적인 영향을 미칩니다. 또한, 금리 변화는 할인율 채널을 통해 은행의 시장 가치에도 영향을 줍니다. 두 경우 모두 자산과 부채 간의 만기 불일치가 핵심적인 역할을 합니다. 위험 관리자는 만기 불일치를 줄이거나, 파생상품을 사용하거나, 비이자 수익으로 다각화하여 금리 위험을 완화할 수 있습니다.
\end{summarybox}


\subsection{유동성 위험 (Liquidity Risk)}

\begin{infobox}[title=\textbf{핵심 요약: 유동성 위험}]
유동성 위험은 은행이 예금 인출이나 대출 약정 이행 등으로 인해 갑작스러운 현금 부족에 직면할 위험입니다. 은행은 자산 관리(현금 보유, 유동성 자산 매각)와 부채 관리(단기 차입)를 통해 이를 관리하며, 예상치 못한 대규모 인출(뱅크런)은 헐값 매각과 지급불능으로 이어질 수 있습니다.
\end{infobox}

모든 금융기관은 부채 보유자들이 현금을 돌려받기를 원할 위험에 직면합니다. 이는 우리가 예금 계좌를 해지하거나 신용 카드를 최대한 사용하는 것과 같은 상황을 포함합니다. 이러한 사건들은 은행이 자금을 마련해야 하는 압력을 가하며, 때로는 매우 짧은 통보 기간 내에 이행되어야 합니다. 이 섹션에서는 은행이 현금 부족(Cash Shortfalls)을 어떻게 관리하는지, 그리고 유동성 위기(Liquidity Crisis)나 뱅크런(Bank Runs)과 같은 대규모 인출 상황에 어떻게 대처하는지 이해하는 것이 목표입니다.

\textbf{유동성 위험의 정의}
유동성 위험은 현금 부족과 그로 인해 발생할 수 있는 비용에 관한 것입니다.
\begin{itemize}
    \item \textbf{현금 흐름:} 은행에는 항상 새로운 예금, 투자 수익, 수수료 수입 등 현금이 유입되고, 예금 인출, 대출 지급 등 현금이 유출됩니다.
    \item \textbf{현금 부족 (Cash Shortfall):} 유출되는 현금이 유입되는 현금보다 많을 때 발생합니다.
    \item \textbf{비용:} 현금 부족을 처리하는 비용은 시장 상황에 따라 달라집니다. 때로는 쉽게 현금을 조달할 수 있지만, 때로는 매우 비쌀 수 있습니다.
\end{itemize}

\textbf{현금 부족의 주요 원인}
\begin{itemize}
    \item \textbf{예금 인출 (Deposits or Withdrawals):} 고객은 요구불 예금 계좌를 매우 짧은 통보 기간 내에 현금으로 전환할 권리가 있습니다.
    \item \textbf{도매 자금 조달 (Wholesale Funding):} 은행은 정기적으로 갱신하거나 롤오버(Rollover)해야 하는 단기 부채인 도매 자금 조달에 의존합니다. 투자자들이 위험 회피적이 되면, 자금 롤오버를 거부하고 현금 상환을 요구할 수 있습니다.
    \item \textbf{신용 약정 (Credit Commitments):} 신용 카드와 같은 신용 약정은 차입자에게 대출을 받을 권리를 부여하며, 이는 은행이 제공해야 하는 현금입니다.
\end{itemize}

\textbf{실제 사례: JP모건 체이스의 유동성 위험 노출 (2020년)}
\begin{itemize}
    \item \textbf{미인출 대출 약정 (Undrawn Loan Commitments):} JP모건의 미인출 대출 약정은 총 대출 잔액의 약 50\%에 달했습니다. 이 대출들이 인출되면 JP모건은 자금을 마련해야 합니다.
    \item \textbf{차입 자금 (Borrowed Funds):} JP모건의 자금 조달에서 차입 자금은 중요한 비중을 차지합니다. 이러한 차입은 단기적으로 이루어질 수 있으며, 이는 '롤오버 위험(Rollover Risk)'에 은행을 노출시킵니다.
\end{itemize}

\textbf{정상적인 유동성 위험 관리 (Normal Times Liquidity Management)}
정상적인 상황에서 예금 인출은 상당히 예측 가능하며, 은행은 두 가지 주요 접근 방식을 사용하여 유동성 위험을 관리합니다.

\begin{enumerate}
    \item \textbf{자산 관리 (Asset Management):}
        \begin{itemize}
            \item \textbf{개념:} 은행은 현금 준비금(Cash Reserves)을 사용하여 현금 부족을 처리합니다.
            \item \textbf{예시: 예금 인출 시}
                \begin{itemize}
                    \item 10달러의 예금이 인출되면, 은행은 단순히 현금 10달러를 내어줍니다. 은행의 자산(현금)과 부채(예금)가 각각 10달러씩 줄어듭니다.
                    \item 또는 은행은 유동성 높은 증권(Liquid Securities)을 매각하여 현금을 확보할 수도 있습니다.
                \end{itemize}
        \end{itemize}
    \item \textbf{부채 관리 (Liability Management):}
        \begin{itemize}
            \item \textbf{개념:} 은행은 추가 자금을 조달하여 현금 부족을 메웁니다.
            \item \textbf{예시: 예금 인출 시}
                \begin{itemize}
                    \item 10달러의 예금이 인출되면, 은행은 다른 예금자나 다른 금융기관으로부터 10달러를 추가로 차입하여 부족분을 충당합니다.
                \end{itemize}
        \end{itemize}
\end{enumerate}
정상적인 상황에서는 이 두 가지 방법이 잘 작동하여 은행이 현금 부족을 관리할 수 있습니다.

\textbf{유동성 위기 및 뱅크런 (Liquidity Crisis and Bank Runs)}
하지만 은행이 예상보다 훨씬 큰 현금 부족에 직면할 때가 있습니다. 부채 보유자들이 한꺼번에 자금을 인출하려 할 수 있는데, 이를 '뱅크런'이라고 합니다. 이는 주로 투자자들이 은행에 대한 신뢰를 잃을 때 발생합니다.
\begin{itemize}
    \item \textbf{관리 도구의 실패:} 뱅크런이 발생하면 일반적인 유동성 관리 도구가 작동하지 않거나 비용이 매우 비싸집니다.
        \begin{itemize}
            \item \textbf{부채 관리 실패:} 은행은 자금 시장에서 차입이 차단될 수 있습니다.
            \item \textbf{자산 관리 실패:} 현금 준비금이 소진되면, 은행은 대출이나 기타 비유동성 자산을 '헐값 매각(Fire Sale)'해야 할 수 있습니다. 이는 시장 가격보다 훨씬 낮은 가격에 자산을 급하게 파는 것을 의미하며, 은행에 막대한 손실을 초래합니다.
        \end{itemize}
    \item \textbf{실제 사례: 2020년 2월의 헐값 매각}
        2020년 2월 말부터 금융 시장에서 거래되는 기업 대출 가격이 급격히 하락했다가 반등했습니다. 만약 은행이 자금 부족을 메우기 위해 이 시점에 기업 대출을 팔아야 했다면, 이는 매우 큰 비용을 초래했을 것입니다.
\end{itemize}

\textbf{뱅크런으로 인한 지급불능 위험 (Insolvency Risk from Bank Runs)}
유동성 위기에 제대로 대비하지 못한 은행은 지급불능(Insolvent) 상태가 되어 자기자본이 소진될 수 있습니다.

\begin{examplebox}[title=\textbf{예시: 단순 은행의 뱅크런}]
\begin{itemize}
    \item \textbf{자산:} 현금 5달러, 증권 20달러, 대출 75달러 (총 100달러)
    \item \textbf{부채:} 예금 70달러, 차입금 20달러, 자기자본 10달러 (총 100달러)
\end{itemize}
뱅크런이 발생하여 예금 20달러와 차입금 20달러, 총 40달러가 인출되었다고 가정합니다. 은행은 차입이 차단되어 부채 관리를 할 수 없습니다.
\begin{itemize}
    \item \textbf{현금 및 증권 사용:} 은행은 현금 5달러와 증권 20달러를 내어줍니다 (총 25달러).
    \item \textbf{대출 매각:} 부족한 15달러(40달러 - 25달러)를 마련하기 위해 대출을 매각해야 합니다. 대출은 정보 민감성이 높아 거래가 어렵고, 잠재적 구매자는 역선택 위험 때문에 상당한 할인을 요구합니다.
    \item \textbf{헐값 매각 손실:} 은행은 필요한 15달러를 마련하기 위해 20달러 상당의 대출을 팔아야 했습니다. 이는 5달러의 손실을 의미하며, 이 손실은 소유주(자기자본)가 부담합니다.
    \item \textbf{결과:} 뱅크런 후 은행은 자산이 75달러(현금 0, 증권 0, 대출 75-20=55)로 줄어들고, 부채는 50달러(예금 70-20=50, 차입금 0)로 줄어듭니다. 자기자본은 10달러에서 5달러 손실을 입어 5달러만 남게 됩니다. 은행은 규모가 작아지고 주주의 자기자본이 줄었지만 생존했습니다.
\end{itemize}
만약 은행이 더 많은 현금을 보유했거나 더 안정적인 자금 조달원을 사용했다면 더 나은 상황이었을 것입니다. 현금 부족은 자기자본 가치에 영향을 미치며, 극단적인 경우 주주를 소멸시키고 은행을 파산시킬 수 있습니다.
\end{examplebox}

\begin{summarybox}[title=\textbf{유동성 위험 요약}]
유동성 위험은 은행이 현금 부족에 직면할 위험입니다. 은행은 자산 관리(현금 보유, 유동성 자산 매각)와 부채 관리(단기 차입)를 통해 이를 처리합니다. 많은 현금을 보유하거나 안정적인 자금 조달원에 접근할 수 있는 은행은 유동성 위기에 더 강합니다. 유동성 위기는 예상치 못한 대규모 인출(뱅크런)로 인해 발생하며, 이는 헐값 매각과 은행의 지급불능으로 이어질 수 있습니다.
\end{summarybox}


\subsection{시장 위험 (Market Risk)}

\begin{infobox}[title=\textbf{핵심 요약: 시장 위험}]
시장 위험은 은행의 트레이딩 포트폴리오에 있는 금융 상품의 시장 가격 변동으로 인해 손실이 발생할 위험입니다. 이는 매일 시가 평가(Mark-to-Market)되며, VaR(Value-at-Risk)은 트레이딩 위험을 평가하는 일반적인 방법이지만, '블랙 스완'과 같은 극단적인 사건에는 한계가 있습니다.
\end{infobox}

많은 은행은 트레이더(Traders)를 고용하여 금융 상품을 매매합니다. 이들은 은행 자기자본의 일부를 투자하여 이익을 창출하려 하지만, 이러한 트레이딩 활동은 본질적으로 위험합니다. 시장 가격은 오르내릴 수 있기 때문입니다. 이 섹션에서는 이러한 '트레이딩 위험(Trading Risk)' 또는 '시장 위험(Market Risk)'을 설명하고, 트레이더의 위험 감수를 평가하는 일반적인 방법인 VaR(Value-at-Risk)을 소개합니다.

\textbf{시장 위험의 발생원: 트레이딩 장부 (Trading Book)}
은행은 크게 두 가지 사업으로 나눌 수 있습니다.
\begin{itemize}
    \item \textbf{대출 사업 (Lending Business):} 은행 장부(Bank Book)에 장기 투자로 보유되는 대출 등으로 구성됩니다.
    \item \textbf{트레이딩 사업 (Trading Business):} 트레이딩 장부(Trading Book)에 보유되는 증권, 파생상품 등 장단기 포지션으로 구성됩니다.
\end{itemize}
트레이딩 장부에 있는 항목들은 매일 '시가 평가(Mark-to-Market)'됩니다. 이는 트레이딩 포지션의 가치가 변동할 때마다 은행이 손실을 즉시 인식해야 한다는 의미입니다. 시장 위험은 이러한 불확실성을 포착합니다. 시장 상황이 변하면 가격이 변하고, 이는 은행의 이익과 자기자본 가치에 영향을 미칠 수 있습니다.

\textbf{실제 사례: 서브프라임 모기지 위기 (2007-2009년)와 리먼 브라더스 (Lehman Brothers)}
\begin{itemize}
    \item \textbf{서브프라임 모기지 위기:} 2006년부터 미국 주택 가격이 하락하기 시작했고, 신용 위험이 가장 높은 서브프라임 차입자들이 부도나기 시작했습니다. 많은 금융기관이 모기지 담보 증권(MBS, Mortgage Backed Securities)을 통해 서브프라임 모기지에 적극적으로 투자하고 있었고, 이 증권들은 막대한 손실을 입었습니다.
        \begin{itemize}
            \item 2009년 중반까지 서브프라임 MBS 손실은 약 1조 달러에 달했습니다. 위기 이전 미국 전체 은행의 총 자기자본이 약 1.3조 달러였음을 감안하면 엄청난 규모였습니다.
            \item 많은 은행이 시장 위험 노출로 인해 지급불능 상태가 되었습니다.
        \end{itemize}
    \item \textbf{리먼 브라더스의 파산:} 리먼 브라더스는 당시 유명한 투자 은행이었으나, 2008년 9월 15일 파산을 선언하며 금융계를 뒤흔들었습니다. 시장 위험이 리먼의 몰락에 핵심적인 역할을 했습니다.
        \begin{itemize}
            \item \textbf{리먼의 트레이딩 포트폴리오 (2008년 5월):} 리먼은 720억 달러 상당의 MBS를 보유하고 있었습니다.
            \item \textbf{가치 하락:} 2007년 7월부터 2008년 7월까지 이 MBS의 가치는 약 50\% 하락했습니다. 이는 대략 360억 달러의 손실을 의미합니다.
            \item \textbf{결과:} 리먼의 자기자본은 260억 달러였습니다. MBS 포트폴리오의 시가 평가 손실(Mark-to-Market Losses)이 은행의 자기자본을 모두 소진시켰고, 이것이 리먼이 파산한 이유입니다.
        \end{itemize}
\end{itemize}

\textbf{시장 위험 관리: VaR (Value-at-Risk)}
트레이딩은 매우 수익성이 높을 수 있지만, 은행은 시장 위험이 지급불능을 초래할 수 있음을 이해합니다. 위험 관리자들은 트레이더의 포트폴리오를 모니터링하고, 시장 위험을 분석하며, 개별 트레이더가 감수할 수 있는 위험의 양에 제한을 둡니다.
\begin{itemize}
    \item \textbf{VaR의 정의:} '일일 VaR(Daily Value-at-Risk)'은 트레이딩 위험을 측정하는 일반적인 방법입니다. 이는 '최악의 시나리오에서 내일 트레이딩 포트폴리오가 얼마나 많은 달러를 잃을 것인가?'를 알려줍니다.
    \item \textbf{VaR의 의미:} "나는 X\%의 신뢰도로 내일 이익이 Y달러를 초과할 것이라고 확신한다." 또는 "1\%의 최악의 시나리오에서 은행의 이익은 Y달러 미만이 될 것이다."
\end{itemize}

\begin{examplebox}[title=\textbf{예시: VaR를 이용한 두 가지 베팅 비교}]
두 가지 베팅이 있다고 가정합니다. 둘 다 1,000달러를 투자해야 합니다.

\textbf{베팅 1:}
\begin{itemize}
    \item 시장 상황이 나쁠 때: 100달러 손실 (확률 50\%)
    \item 시장 상황이 좋을 때: 300달러 이익 (확률 50\%)
    \item 예상 수익: 100달러
\end{itemize}
\textbf{베팅 2:}
\begin{itemize}
    \item 800달러 손실 (확률 10\%)
    \item 100달러 손실 (확률 40\%)
    \item 200달러 이익 (확률 40\%)
    \item 1,000달러 이익 (확률 10\%)
    \item 예상 수익: 100달러
\end{itemize}
VaR는 "최악의 시나리오에서 나의 손실은 얼마인가?"라는 질문에 답합니다.
\begin{itemize}
    \item \textbf{베팅 1의 VaR:} 최대 100달러 손실.
    \item \textbf{베팅 2의 VaR:} 최대 800달러 손실.
\end{itemize}
만약 최악의 결과에 대해 걱정한다면, 베팅 2가 더 위험합니다. 은행은 이와 동일한 개념을 트레이딩 포트폴리오에 적용합니다.
\end{examplebox}

\textbf{VaR의 한계 및 비판}
VaR는 위험 측정 지표로서 종종 비판을 받습니다.
\begin{itemize}
    \item \textbf{비판 이유:} 실제로는 수익 분포가 항상 매끄럽지 않기 때문입니다. 극단적인 시장 상황에서는 예상보다 훨씬 더 나쁜 수익이 발생할 수 있습니다. 이러한 극단적인 사건을 '블랙 스완(Black Swans)'이라고 부르며, 이는 수익 분포의 '두꺼운 꼬리(Fat Tail)'에 해당합니다.
    \item \textbf{문제점:} VaR 통계는 여전히 유효할 수 있습니다 (예: 99\%의 시간 동안 이익이 Y달러를 초과함). 그러나 만약 내일이 '나쁜 날'이 되어 두꺼운 꼬리 영역에 속한다면, 평균 손실은 Y달러보다 훨씬 더 클 수 있습니다.
    \item \textbf{2008년 위기 경험:} 2008년 금융 위기 동안 많은 금융기관의 트레이딩 수익은 분포의 두꺼운 꼬리 영역에 속했습니다. 위험 관리자와 규제 당국은 극단적인 손실이 VaR 예측을 훨씬 초과할 수 있음을 깨달았습니다. 그 이후로 이러한 가능성을 설명하는 새로운 시장 위험 모델이 등장했습니다.
\end{itemize}

\begin{summarybox}[title=\textbf{시장 위험 요약}]
시장 위험은 트레이딩 활동에서 발생하는 잠재적 손실을 의미하며, 시장 상황 변화에 따라 매우 심각할 수 있습니다. VaR는 이러한 꼬리 사건(Tail Events)에 대한 달러 가치를 부여하는 일반적인 측정 방법입니다. 그러나 VaR는 '블랙 스완'과 같은 극단적인 시장 상황에서 발생하는 '두꺼운 꼬리' 위험을 완전히 포착하지 못할 수 있다는 한계가 있습니다.
\end{summarybox}


\subsection{운영 위험 (Operational Risks)}

\begin{infobox}[title=\textbf{핵심 요약: 운영 위험}]
운영 위험은 내부 프로세스, 인력, 시스템의 실패 또는 외부 사건으로 인해 발생하는 광범위한 손실 위험입니다. 이는 기술 오작동, 직원 사기, 사이버 공격, 그리고 핀테크 경쟁으로 인한 비즈니스 모델의 변화 등 다양한 형태로 나타나며, 예측하기 어렵지만 발생 시 막대한 손실을 초래할 수 있습니다.
\end{infobox}

이 섹션에서는 운영 위험에 대한 간략한 설명을 제공하고, 구체적인 사례와 운영 위험으로 인한 손실 규모를 논의합니다.

\textbf{운영 위험의 공식 정의}
운영 위험은 "부적절하거나 실패한 내부 프로세스, 인력, 시스템 또는 외부 사건으로 인해 발생하는 손실 위험"입니다.
\begin{itemize}
    \item \textbf{포괄적인 정의:} 이는 매우 모호하고 포괄적인 정의로, 많은 잠재적 손실을 포함합니다.
    \item \textbf{내부 실패:} 사기(Fraud)와 같은 내부 실패를 포함합니다.
    \item \textbf{외부 사건:} 사이버 공격(Cyber-attack)부터 지진(Earthquake)까지 모든 외부 사건을 포함할 수 있습니다.
\end{itemize}

\textbf{주요 운영 위험원}
가장 많은 주목을 받는 세 가지 잠재적 손실 원인은 다음과 같습니다.

\textbf{1. 운영 실패 (Operational Failures)}
\begin{itemize}
    \item \textbf{기술 오작동 및 데이터 유출 (Technology Malfunctions and Data Breaches):} 기술 시스템의 오작동이나 사이버 공격으로 인한 데이터 유출은 막대한 손실을 초래할 수 있습니다.
    \item \textbf{직원 오류 (Employee Errors):} 잘못된 가격으로 거래를 실수로 실행하는 것과 같은 직원 오류도 손실을 발생시킬 수 있습니다.
\end{itemize}

\textbf{2. 직원 사기 (Employee Fraud)}
\begin{itemize}
    \item \textbf{개념:} 직원의 사기는 운영 위험의 주요 구성 요소입니다.
    \item \textbf{사례:} 수년간 은행들은 미공개 또는 무단 거래로 인해 수십억 달러의 손실을 입었습니다. 이러한 사건들은 대형 은행들 사이에서 놀랍도록 빈번하게 발생했습니다.
    \item \textbf{실제 사례: 은행 벌금 (2008-2015년)}
        \begin{itemize}
            \item \textbf{불법 활동으로 인한 벌금:} 2008년부터 2015년까지 은행들이 지불한 벌금은 주택 담보 대출 관련 사기, 리보(Libor) 금리 조작 스캔들 관련 사기, 런던 고래(London Whale) 사건 관련 사기 등 불법 활동으로 인한 것이었습니다.
            \item \textbf{런던 고래 사건:} JP모건에서 발생한 미공개 트레이딩 사례로, 은행에 약 60억 달러의 손실과 10억 달러의 벌금을 초래했습니다. 이것이 바로 운영 위험의 예시입니다. 예측하기는 어렵지만, 발생하면 그 영향을 명확히 알 수 있습니다.
            \item \textbf{손실 규모:} 뱅크 오브 아메리카(Bank of America)는 이 기간 동안 약 800억 달러의 벌금을 지불했습니다. 2015년 말 뱅크 오브 아메리카의 자기자본이 약 2,000억 달러였음을 감안하면, 이러한 손실은 은행 자기자본의 상당 부분을 차지했습니다.
        \end{itemize}
\end{itemize}

\textbf{3. 핀테크 경쟁 (Fintech Competition)}
\begin{itemize}
    \item \textbf{개념:} 핀테크(Fintech) 기업들은 많은 금융 서비스 기업의 전통적인 비즈니스 모델을 완전히 뒤엎고 있습니다.
    \item \textbf{영향:} 이는 고객 손실과 수익 손실로 이어질 수 있습니다.
    \item \textbf{은행의 대응:} 은행은 자신들의 위치를 방어하기 위해 새로운 기술을 채택하고 기존 IT 시스템을 업그레이드해야 합니다. 물론 이러한 투자와 변화는 위험하며, 성공하지 못할 수도 있습니다.
    \item \textbf{실제 사례: JP모건의 디지털 전략 (2020년)}
        JP모건은 소비자 금융 사업 내에서 온라인 대출 신청, 온라인 수표 입금, 모바일 결제, 디지털 자산 관리 앱 등 디지털 전략에 적극적으로 투자하고 있습니다. 2020년 JP모건의 기술 예산은 약 100억 달러에 달했습니다.
\end{itemize}

\begin{summarybox}[title=\textbf{운영 위험 요약}]
운영 위험은 광범위한 '포괄적(Catch-all)' 위험입니다. 이는 일상적인 대출이나 투자 은행 업무에서 발생하는 나쁜 거래로 인한 손실이 아닙니다. 대신, 실수, 사기, 기술 실패 또는 기술 노후화로 인해 발생할 수 있는 손실을 의미합니다. 이는 예측하기 어렵지만, 발생 시 은행에 막대한 영향을 미칠 수 있습니다.
\end{summarybox}


\section{위험 관리 절차 및 방법}

은행은 다양한 위험에 노출되어 있으며, 이러한 위험을 체계적으로 관리하기 위한 절차와 방법을 갖추고 있습니다.

\subsection{위험 관리의 일반적인 접근 방식}
\begin{enumerate}
    \item \textbf{위험 식별 (Risk Identification):}
        \begin{itemize}
            \item 은행이 직면할 수 있는 모든 잠재적 위험을 파악합니다 (예: 신용, 금리, 유동성, 시장, 운영 위험).
            \item 각 위험이 은행의 이익과 자본에 어떤 영향을 미칠 수 있는지 분석합니다.
        \end{itemize}
    \item \textbf{위험 측정 (Risk Measurement):}
        \begin{itemize}
            \item 식별된 위험의 규모와 발생 확률을 정량적으로 평가합니다.
            \item \textbf{신용 위험:} 신용 등급 모델, 부도 확률(PD), 부도 시 손실률(LGD) 등을 사용합니다.
            \item \textbf{시장 위험:} VaR(Value-at-Risk) 모델을 사용하여 특정 신뢰 수준에서의 최대 잠재 손실을 추정합니다.
            \item \textbf{금리 위험:} 자산-부채 만기 구조 분석, 금리 민감도 분석 등을 수행합니다.
            \item \textbf{유동성 위험:} 현금 흐름 예측, 스트레스 테스트(Stress Test)를 통해 극단적인 상황에서의 현금 부족을 시뮬레이션합니다.
        \end{itemize}
    \item \textbf{위험 통제 및 완화 (Risk Control and Mitigation):}
        측정된 위험을 줄이거나 관리하기 위한 전략을 실행합니다.
        \begin{itemize}
            \item \textbf{대출 계약 설계:} 신용 위험 관리를 위해 대출 이자율 조정, 약정 설정, 담보 및 우선순위 확보.
            \item \textbf{다각화:} 다양한 자산, 지역, 산업에 분산 투자하여 특정 위험에 대한 집중 노출을 줄입니다.
            \textbf{위험 이전 (Risk Transfer):}
                \begin{itemize}
                    \item \textbf{대출 판매:} 신용 위험이 높은 대출을 다른 금융기관에 판매하여 대차대조표에서 위험을 제거합니다.
                    \item \textbf{파생상품 사용:} 신용부도스왑(CDS)으로 신용 위험을, 금리 스왑으로 금리 위험을 헤지합니다.
                \end{itemize}
            \item \textbf{자산-부채 관리 (ALM, Asset-Liability Management):} 금리 위험 및 유동성 위험 관리를 위해 자산과 부채의 만기 구조를 조정합니다.
            \item \textbf{자본 확충:} 손실 흡수 능력을 높이기 위해 자기자본을 늘립니다.
            \item \textbf{내부 통제 강화:} 운영 위험을 줄이기 위해 내부 프로세스, 시스템, 인력에 대한 통제를 강화하고 사기를 방지합니다.
        \end{itemize}
    \item \textbf{위험 모니터링 및 보고 (Risk Monitoring and Reporting):}
        \begin{itemize}
            \item 관리된 위험 노출이 허용 가능한 수준 내에 있는지 지속적으로 감시합니다.
            \item 위험 지표와 관리 성과를 정기적으로 경영진과 규제 당국에 보고합니다.
            \item 시장 상황 변화나 내부 프로세스 변경에 따라 위험 관리 전략을 조정합니다.
        \end{itemize}
\end{enumerate}

\begin{examplebox}[title=\textbf{예시: 금리 위험 관리 절차}]
\begin{enumerate}
    \item \textbf{식별:} 은행의 자산(장기 대출)과 부채(단기 예금) 간의 만기 불일치로 인해 금리 변동에 취약하다는 것을 식별합니다.
    \item \textbf{측정:} 금리가 1\% 상승할 경우 은행의 순 이자 수입이 얼마나 감소하고, 자기자본 가치가 얼마나 하락할지 시뮬레이션(예: 갭 분석, 듀레이션 분석)하여 측정합니다.
    \item \textbf{통제/완화:}
        \begin{itemize}
            \item \textbf{만기 조절:} 단기 대출 비중을 늘리거나 장기 예금 상품을 개발하여 만기 불일치를 줄입니다.
            \item \textbf{파생상품:} 금리 스왑을 매수하여 변동 금리 부채를 고정 금리 부채로 전환함으로써 금리 상승 위험을 헤지합니다.
            \item \textbf{다각화:} 수수료 기반의 비이자 수익 사업을 확장하여 금리 변동에 덜 민감한 수익원을 확보합니다.
        \end{itemize}
    \item \textbf{모니터링:} 매일 또는 매주 금리 변동과 은행의 금리 민감도 지표(예: 순 이자 마진, 자기자본 가치 변화)를 모니터링하고, 필요시 헤지 포지션을 조정합니다.
\end{enumerate}
\end{examplebox}


\section{체크리스트: 모듈 3 학습 점검}

이 체크리스트는 'Banking and Financial Institutions Module 3'의 핵심 내용을 효과적으로 학습하고 이해했는지 스스로 점검하는 데 도움이 됩니다.

\begin{summarybox}[title=\textbf{모듈 3 학습 점검 체크리스트}]
\begin{itemize}
    \item[] \textbf{은행 자본 및 수익성}
    \item[] \quad $\square$ 은행 자기자본의 정의와 역할을 설명할 수 있는가?
    \item[] \quad $\square$ 자기자본이 손실 흡수 완충재 역할을 하는 원리를 이해하는가?
    \item[] \quad $\square$ 주식 수익률, ROA, ROE의 개념과 계산 방법을 설명할 수 있는가?
    \item[] \quad $\square$ 레버리지가 ROE에 미치는 긍정적/부정적 영향을 설명할 수 있는가?
    \item[] \quad $\square$ JP모건 체이스 및 미국 은행 산업의 성과 지표를 해석할 수 있는가?

    \item[] \textbf{은행 위험 개요}
    \item[] \quad $\square$ 위험의 일반적인 정의와 은행 위험의 주요 원천을 설명할 수 있는가?
    \item[] \quad $\square$ 신용, 금리, 유동성, 시장, 운영 위험을 각각 정의할 수 있는가?
    \item[] \quad $\square$ 지급불능 위험이 발생하는 과정을 설명할 수 있는가?
    \item[] \quad $\square$ 위험 관리의 두 가지 주요 단계(측정, 관리)를 설명할 수 있는가?
    \item[] \quad $\square$ 위험 관리의 일반적인 방법(가격 책정, 헤징)을 이해하는가?

    \item[] \textbf{신용 위험}
    \item[] \quad $\square$ 신용 위험의 정의와 부도(Default)의 의미를 설명할 수 있는가?
    \item[] \quad $\square$ 체계적 신용 위험의 개념을 이해하는가?
    \item[] \quad $\square$ 신용 위험 측정 과정(데이터, 모델, 신용 등급)을 설명할 수 있는가?
    \item[] \quad $\square$ FICO와 같은 개인 신용 등급 시스템을 이해하는가?
    \item[] \quad $\square$ 신용 위험 관리 방법(대출 계약 설계, 약정, 담보, 다각화, 대출 판매, 신용 파생상품)을 설명할 수 있는가?

    \item[] \textbf{금리 위험}
    \item[] \quad $\square$ 만기 불일치가 금리 위험을 초래하는 원리를 설명할 수 있는가?
    \item[] \quad $\square$ 금리 위험이 은행 이익에 영향을 미치는 현금 흐름 채널을 설명할 수 있는가?
    \item[] \quad $\square$ 금리 위험이 은행 가치에 영향을 미치는 평가 채널을 설명할 수 있는가?
    \item[] \quad $\square$ 금리 위험 관리 방법(만기 일치, 파생상품, 비이자 수익 다각화)을 설명할 수 있는가?

    \item[] \textbf{유동성 위험}
    \item[] \quad $\square$ 유동성 위험의 정의와 현금 부족의 원인을 설명할 수 있는가?
    \item[] \quad $\square$ 자산 관리와 부채 관리를 통한 유동성 관리 방법을 설명할 수 있는가?
    \item[] \quad $\square$ 뱅크런과 헐값 매각이 유동성 위기에서 발생하는 과정을 설명할 수 있는가?
    \item[] \quad $\square$ 유동성 위기가 은행의 지급불능 위험으로 이어지는 과정을 이해하는가?

    \item[] \textbf{시장 위험}
    \item[] \quad $\square$ 시장 위험의 정의와 트레이딩 장부/은행 장부의 차이를 설명할 수 있는가?
    \item[] \quad $\square$ 시가 평가(Mark-to-Market)의 의미를 이해하는가?
    \item[] \quad $\square$ VaR(Value-at-Risk)의 개념과 계산 원리를 설명할 수 있는가?
    \item[] \quad $\square$ VaR의 한계점(블랙 스완, 두꺼운 꼬리)을 설명할 수 있는가?
    \item[] \quad $\square$ 리먼 브라더스 파산 사례를 통해 시장 위험의 중요성을 설명할 수 있는가?

    \item[] \textbf{운영 위험}
    \item[] \quad $\square$ 운영 위험의 정의와 주요 원천(운영 실패, 직원 사기, 외부 사건)을 설명할 수 있는가?
    \item[] \quad $\square$ 런던 고래 사건과 같은 실제 사례를 통해 운영 위험을 설명할 수 있는가?
    \item[] \quad $\square$ 핀테크 경쟁이 은행의 운영 위험에 미치는 영향을 설명할 수 있는가?
    \item[] \quad $\square$ 은행이 핀테크에 대응하기 위한 기술 투자 사례를 이해하는가?
\end{itemize}
\end{summarybox}


\section{FAQ: 자주 묻는 질문}

초심자들이 혼동할 수 있는 질문들을 Q\&A 형식으로 정리했습니다.

\begin{infobox}[title=\textbf{FAQ: 은행 위험 및 수익성}]
\textbf{Q1: 은행의 자기자본이 왜 중요한가요?}
\textbf{A1:} 자기자본은 은행의 손실을 흡수하는 최후의 보루 역할을 합니다. 은행이 예상치 못한 손실을 입었을 때, 자기자본이 먼저 손실을 부담하여 예금자나 다른 채권자들이 피해를 입는 것을 막아줍니다. 이는 은행의 안정성을 높이고 금융 시스템 전체의 신뢰를 유지하는 데 필수적입니다.

\textbf{Q2: ROA와 ROE는 어떻게 다른가요? 은행 성과를 평가할 때 어떤 지표를 봐야 하나요?}
\textbf{A2:} ROA(총자산 수익률)는 은행이 총자산을 얼마나 효율적으로 운용하여 이익을 창출하는지 보여줍니다. 반면 ROE(자기자본 수익률)는 주주가 투자한 자기자본 대비 얼마나 많은 이익을 얻었는지를 나타냅니다. ROA는 은행의 전반적인 효율성을, ROE는 주주 관점의 수익성을 평가하는 데 중요합니다. 레버리지가 높은 은행은 ROA가 낮아도 ROE가 높을 수 있으므로, 두 지표를 함께 보면서 은행의 효율성과 레버리지 전략을 종합적으로 판단해야 합니다.

\textbf{Q3: 레버리지가 은행에 항상 좋은가요?}
\textbf{A3:} 레버리지는 양날의 검입니다. 좋은 시장 상황에서는 적은 자기자본으로 더 많은 자산을 운용하여 ROE를 크게 증폭시킬 수 있습니다. 하지만 시장 상황이 나빠져 손실이 발생하면, 레버리지는 손실을 더욱 크게 증폭시켜 은행의 자기자본을 빠르게 소진시키고 지급불능 위험을 높일 수 있습니다. 따라서 적절한 레버리지 수준을 유지하는 것이 중요합니다.

\textbf{Q4: 신용 위험과 시장 위험은 어떻게 다른가요?}
\textbf{A4:} 신용 위험은 차입자가 대출 원리금을 갚지 못할 위험, 즉 '부도 위험'에 초점을 맞춥니다. 주로 은행 장부(Bank Book)의 대출 자산에서 발생합니다. 반면 시장 위험은 은행의 트레이딩 장부(Trading Book)에 있는 금융 상품의 시장 가격이 변동하여 손실을 입을 위험입니다. 신용 위험은 차입자의 개별 상황에, 시장 위험은 거시적인 시장 움직임에 더 크게 영향을 받습니다.

\textbf{Q5: 뱅크런은 왜 발생하며, 은행은 어떻게 대응하나요?}
\textbf{A5:} 뱅크런은 예금자들이 은행의 재정 상태에 대한 신뢰를 잃고 한꺼번에 예금을 인출하려 할 때 발생합니다. 은행은 평소에 충분한 현금 준비금을 보유하거나, 유동성 높은 증권을 보유하여 자산 관리로 대응합니다. 또한, 다른 금융기관에서 단기 자금을 차입하는 부채 관리로도 대응합니다. 하지만 대규모 뱅크런 시에는 이러한 방법이 한계에 부딪히고, 은행은 자산을 헐값에 매각해야 하거나 결국 파산할 수도 있습니다.

\textbf{Q6: VaR(Value-at-Risk)이 시장 위험을 측정하는 데 완벽한가요?}
\textbf{A6:} VaR는 특정 신뢰 수준에서 발생할 수 있는 최대 손실을 예측하는 유용한 도구입니다. 하지만 VaR는 '블랙 스완'과 같이 극히 드물지만 발생하면 엄청난 영향을 미치는 극단적인 사건(두꺼운 꼬리 위험)을 충분히 반영하지 못할 수 있다는 한계가 있습니다. VaR가 예측한 손실보다 실제 손실이 훨씬 더 커질 수 있으므로, VaR 외에 스트레스 테스트 등 다른 위험 측정 도구와 함께 사용해야 합니다.

\textbf{Q7: 핀테크가 은행에 어떤 위험을 가져오나요?}
\textbf{A7:} 핀테크 기업들은 혁신적인 기술을 활용하여 전통적인 은행 서비스를 대체하거나 개선하고 있습니다. 이는 은행에게 고객 이탈과 수익 감소라는 '운영 위험'을 가져옵니다. 은행은 이에 대응하기 위해 막대한 기술 투자를 하고 디지털 전략을 강화해야 하지만, 이러한 투자 역시 성공을 보장할 수 없는 새로운 운영 위험을 수반합니다.
\end{infobox}


\section{빠르게 훑어보기: 1페이지 요약}

이 섹션은 모듈 3의 핵심 내용을 한눈에 파악할 수 있도록 요약한 것입니다.

\begin{summarybox}[title=\textbf{모듈 3 핵심 요약: 은행의 위험과 수익}]
\textbf{1. 은행 자본과 수익성}
\begin{itemize}
    \item \textbf{자기자본:} 은행 소유주의 투자금이자 손실 흡수 완충재. 잔여 청구권으로 손실 발생 시 가장 먼저 타격.
    \item \textbf{성과 측정:}
        \begin{itemize}
            \item \textbf{ROA (총자산 수익률):} 자산 효율성.
            \item \textbf{ROE (자기자본 수익률):} 주주 투자 수익률.
            \item \textbf{레버리지:} ROE를 증폭시키지만, 손실도 증폭시켜 위험 증가.
        \end{itemize}
\end{itemize}

\textbf{2. 은행 위험 개요}
\begin{itemize}
    \item \textbf{위험 정의:} 미래 불확실한 결과로 인한 잠재적 손실.
    \item \textbf{주요 위험:}
        \begin{itemize}
            \item \textbf{신용 위험:} 차입자 부도 위험.
            \item \textbf{금리 위험:} 금리 변동으로 인한 수익/가치 변동 (만기 불일치).
            \item \textbf{유동성 위험:} 현금 부족 위험 (뱅크런, 헐값 매각).
            \item \textbf{시장 위험:} 트레이딩 포트폴리오 가격 변동 위험 (VaR, 블랙 스완).
            \item \textbf{운영 위험:} 내부 실패/외부 사건으로 인한 손실 (사기, 기술 오류, 핀테크).
        \end{itemize}
    \item \textbf{위험 관리:} 위험 측정 $\rightarrow$ 위험 관리 (가격 책정, 헤징, 다각화).
\end{itemize}

\textbf{3. 각 위험별 상세 관리 전략}
\begin{itemize}
    \item \textbf{신용 위험:} 대출 계약 설계 (가격, 약정, 담보), 다각화, 대출 판매, 신용 파생상품 (CDS).
    \item \textbf{금리 위험:} 만기 일치, 금리 파생상품 (스왑), 비이자 수익 다각화.
    \item \textbf{유동성 위험:} 자산 관리 (현금, 유동성 증권), 부채 관리 (단기 차입).
    \item \textbf{시장 위험:} VaR를 통한 위험 측정, 포지션 한도 설정, 스트레스 테스트.
    \item \textbf{운영 위험:} 내부 통제 강화, 기술 투자, 디지털 전략.
\end{itemize}

\textbf{결론:} 현대 은행은 수익 창출과 동시에 다양한 위험에 노출되어 있습니다. 성공적인 은행 경영은 이러한 위험을 정확히 측정하고 효과적으로 관리하여 '위험 조정 수익률'을 극대화하는 데 달려 있습니다.
\end{summarybox}


\section{학습 로드맵: FIN-571 모듈 3}

이 로드맵은 'Banking and Financial Institutions Module 3'의 내용을 체계적으로 학습하고 시험에 대비하는 데 도움을 줍니다.

\begin{infobox}[title=\textbf{모듈 3 학습 로드맵}]
\begin{enumerate}
    \item \textbf{기초 다지기: 은행의 기본 구조와 수익성 이해}
    \begin{itemize}
        \item 은행의 자기자본이 무엇이며 왜 중요한지, 손실 흡수 메커니즘을 파악합니다.
        \item ROA, ROE, 주식 수익률 등 은행 성과 지표의 정의와 계산법을 숙지하고, 레버리지의 역할을 이해합니다.
        \item JP모건 체이스와 같은 실제 사례를 통해 이론을 적용해 봅니다.
    \end{itemize}

    \item \textbf{핵심 개념: 은행이 직면하는 주요 위험 파악}
    \begin{itemize}
        \item 신용, 금리, 유동성, 시장, 운영 위험의 정의와 발생 원인을 명확히 구분합니다.
        \item 각 위험이 은행의 이익과 가치에 어떻게 영향을 미치는지, 그리고 궁극적으로 지급불능 위험으로 이어질 수 있는 과정을 이해합니다.
        \item 위험 관리의 두 가지 큰 축인 '측정'과 '관리'의 중요성을 인식합니다.
    \end{itemize}

    \item \textbf{심화 학습: 각 위험별 측정 및 관리 전략 숙달}
    \begin{itemize}
        \item 신용 위험 관리를 위한 대출 계약 설계, 다각화, 대출 판매, 신용 파생상품의 구체적인 방법을 학습합니다.
        \item 금리 위험의 현금 흐름 채널과 평가 채널을 심층적으로 이해하고, 만기 조절, 파생상품, 비이자 수익 다각화 전략을 파악합니다.
        \item 유동성 위험 시 뱅크런과 헐값 매각의 메커니즘을 이해하고, 자산/부채 관리 전략을 익힙니다.
        \item 시장 위험 측정 도구인 VaR의 원리와 한계점(블랙 스완)을 명확히 이해합니다.
        \item 운영 위험의 다양한 형태(사기, 기술 실패, 핀테크 경쟁)와 은행의 대응 전략을 파악합니다.
    \end{itemize}

    \item \textbf{통합 및 적용: 위험-수익 상충 관계 분석}
    \begin{itemize}
        \item 각 위험 관리 전략이 은행의 수익성에 어떤 영향을 미치는지 종합적으로 분석합니다.
        \item '위험 조정 수익률'의 개념을 바탕으로 은행이 어떻게 최적의 위험-수익 균형을 찾아가는지 논의합니다.
        \item 실제 금융 시장의 사례(예: 리먼 브라더스, 런던 고래)를 통해 이론적 지식을 현실에 적용하는 연습을 합니다.
    \end{itemize}
\end{enumerate}
\end{infobox}


%=======================================================================
% Module 4: Lecture 4
%=======================================================================
\chapter{Lecture 4}
\label{ch:module4}

\metainfo{FIN-571: 금융 기관 및 통화 정책}{Module 4: 규제}{Prof. Rustom Manouchehri Irani}{이 문서는 '금융 기관 및 통화 정책' 모듈 4의 강의 내용을 통합하여, 은행 규제의 필요성, 실제 작동 방식, 그리고 21세기 규제 당국이 직면한 주요 과제를 심층적으로 다룹니다. 정부가 은행을 규제하는 이유와 그 목표를 시작으로, 은행 시스템의 취약성과 정부 안전망의 역할을 설명합니다. 실제 규제 및 감독 과정, 특히 자본 요건, 감독 절차, 스트레스 테스트를 상세히 살펴봅니다. 마지막으로, '대마불사' 문제와 그림자 금융이라는 현대적 도전에 대해 논의하며, 규제가 어떻게 진화하고 있는지 이해를 돕습니다. 이 문서를 통해 현대 은행이 왜 규제되는지, 규제 과정이 어떻게 진행되는지, 그리고 새로운 도전에 어떻게 대응하는지 완벽하게 이해할 수 있습니다.}


\section{개요}

\begin{summarybox}[title=모듈 4: 규제 개요]
이 문서는 '금융 기관 및 통화 정책' 모듈 4의 강의 내용을 통합하여, 은행 규제의 필요성, 실제 작동 방식, 그리고 21세기 규제 당국이 직면한 주요 과제를 심층적으로 다룹니다.

정부가 은행을 규제하는 이유와 그 목표를 시작으로, 은행 시스템의 취약성과 정부 안전망의 역할을 설명합니다.

실제 규제 및 감독 과정, 특히 자본 요건, 감독 절차, 스트레스 테스트를 상세히 살펴봅니다.

마지막으로, '대마불사' 문제와 그림자 금융이라는 현대적 도전에 대해 논의하며, 규제가 어떻게 진화하고 있는지 이해를 돕습니다.

이 문서를 통해 현대 은행이 왜 규제되는지, 규제 과정이 어떻게 진행되는지, 그리고 새로운 도전에 어떻게 대응하는지 완벽하게 이해할 수 있습니다.
\end{summarybox}


\section{용어 정리}

다음 표는 본 문서에서 다루는 주요 용어들을 비전공자도 쉽게 이해할 수 있도록 설명합니다.

\begin{adjustbox}{width=\textwidth,center}
\begin{tabular}{lllc}
\toprule
\textbf{용어} & \textbf{쉬운 설명} & \textbf{원어} & \textbf{비고} \\
\midrule
규제 & 정부가 시장 참여자에게 특정 규칙이나 제약을 부과하는 행위 & Regulation & 시장 실패를 교정 \\
시장 실패 & 시장이 자원을 효율적으로 배분하지 못하는 상황 & Market Failure & 규제의 주된 이유 \\
부정적 외부 효과 & 어떤 경제 활동이 관련 없는 제3자에게 해를 끼치는 것 & Negative Externality & 은행 실패의 파급 효과 \\
미시 건전성 규제 & 개별 은행의 위험을 줄여 파산을 막는 규제 & Micro-prudential Regulation & 개별 은행의 안정성 목표 \\
거시 건전성 규제 & 금융 시스템 전체의 위험을 줄여 시스템 위기를 막는 규제 & Macro-prudential Regulation & 시스템 전체의 안정성 목표 \\
정부 안전망 & 은행 시스템의 안정성을 위해 정부가 제공하는 보호 장치 & Government Safety Net & 예금 보험, 최종 대부자 \\
유동성 불일치 & 은행 자산은 장기적이고 비유동적인 반면, 부채는 단기적이고 유동적인 상태 & Liquidity Mismatch & 은행 취약성의 근본 원인 \\
뱅크런 & 예금자들이 은행의 건전성을 의심하여 한꺼번에 예금을 인출하는 현상 & Bank Run & 은행 파산으로 이어질 수 있음 \\
전염 효과 & 한 은행의 문제가 다른 은행으로 확산되는 현상 & Contagion & 금융 위기 심화 요인 \\
예금 보험 & 은행이 파산해도 예금자의 예금을 일정 한도까지 보호해주는 제도 & Deposit Insurance & FDIC (미국) \\
최종 대부자 & 금융 위기 시 일시적인 자금 부족을 겪는 은행에 대출해주는 기관 & Lender of Last Resort & 중앙은행 (미국 연준) \\
도덕적 해이 & 정부의 보호를 믿고 은행이 과도한 위험을 감수하는 경향 & Moral Hazard & 규제의 부작용 \\
바젤 협약 & 국제 은행들이 지켜야 할 최소 자기자본 비율에 대한 국제적 합의 & Basel Accord & 은행의 손실 흡수 능력 강화 \\
자본 요건 & 은행이 보유해야 하는 최소한의 자기자본 비율 & Capital Requirements & 위험 가중 자산 대비 \\
감독 & 규제 준수 여부를 지속적으로 감시하고 위반 시 제재하는 과정 & Supervision & 현장 검사, 스트레스 테스트 \\
CAMELS 평가 & 은행의 건전성을 6가지 기준으로 평가하는 시스템 & CAMELS Rating & 자본, 자산, 경영, 수익, 유동성, 시장 위험 민감도 \\
스트레스 테스트 & 극심한 경제 위기 상황을 가정하여 은행의 손실 흡수 능력을 평가 & Stress Testing & 미래 위험 대비 \\
대마불사 & 너무 커서 망하게 내버려 둘 수 없는 금융 기관 & Too Big to Fail (TBTF) & 정부의 구제금융 가능성 \\
시스템 리스크 & 한 금융 기관의 실패가 전체 금융 시스템에 파급되는 위험 & Systemic Risk & 금융 위기의 핵심 \\
FSOC & 금융 시스템의 안정성을 감독하고 시스템 리스크를 관리하는 미국 기관 & Financial Stability Oversight Council & 시스템적으로 중요한 기관 지정 \\
SIFI / GSIB & 시스템적으로 중요한 금융 기관 / 글로벌 시스템적으로 중요한 은행 & Systemically Important Financial Institution / Global Systemically Important Bank & 강화된 규제 대상 \\
그림자 금융 & 은행과 유사한 금융 서비스를 제공하지만 규제는 덜 받는 기관 & Shadow Banking & 규제 차익거래 가능성 \\
머니마켓펀드 & 단기 금융 상품에 투자하는 펀드로, 은행 예금과 유사한 기능 & Money Market Mutual Fund (MMMF) & 대표적인 그림자 금융 \\
\bottomrule
\end{tabular}
\end{adjustbox}


\section{핵심 개념·원리}

\subsection{금융 규제의 목표 이해}

정부가 민간 시장에 개입하여 금융 기관을 규제하는 것은 단순히 통제하려는 것이 아니라, 시장이 스스로 해결하지 못하는 문제들을 바로잡아 사회 전체에 더 나은 결과를 가져오기 위함입니다. 특히 은행은 그 특성상 규제가 더욱 중요합니다.

\subsubsection{규제의 필요성: 시장 실패와 부정적 외부 효과}

\begin{itemize}
    \item \textbf{한 줄 핵심 요약:} 시장이 자원을 효율적으로 배분하지 못하고, 특정 활동이 의도치 않게 다른 사람들에게 피해를 줄 때 정부 규제가 필요합니다.
    \item \textbf{직관적 예시/비유:} 공장이 오염 물질을 배출하여 강물과 주변 환경을 해치는 경우를 생각해 보세요. 공장 주인은 생산 비용만 고려할 뿐, 오염으로 인한 어업 피해나 주민 건강 악화 같은 '외부 비용'은 신경 쓰지 않습니다. 이처럼 개인에게는 최적의 생산량이라도 사회 전체적으로는 과도한 오염(생산)이 될 수 있습니다. 정부는 환경 규제를 통해 공장이 이러한 외부 비용을 고려하도록 유도합니다.
    \item \textbf{기술적 정확한 설명:} 규제(Regulation)는 정부가 자유 시장에 개입하여 시장 참여자에게 제약을 가하는 행위입니다. 이는 시장 균형에 영향을 미쳐 사회적으로 더 나은 결과를 도출하는 것을 목표로 합니다. 주로 시장 실패(Market Failure)를 교정하기 위해 이루어지며, 특히 부정적 외부 효과(Negative Externality)나 파급 효과(Spillover)가 발생할 때 중요합니다. 시장 참여자들은 자신에게 직접적인 영향을 미치지 않는 외부 효과를 무시하는 경향이 있으므로, 정부는 이들이 의사 결정 시 외부 효과를 고려하도록 유인합니다.
\end{itemize}

\subsubsection{은행의 독특한 외부 효과}

\begin{itemize}
    \item \textbf{한 줄 핵심 요약:} 개별 은행의 실패는 해당 은행을 넘어 다른 금융 기관과 전체 경제에 심각한 피해를 줄 수 있는 '전염성'을 가지고 있습니다.
    \item \textbf{직관적 예시/비유:} 한 은행이 파산하면, 그 은행에 돈을 맡긴 사람들은 저축을 잃거나 돈을 찾지 못하게 됩니다. 또 그 은행에서 대출을 받던 기업들은 자금줄이 끊겨 사업을 접을 수도 있습니다. 더 나아가, 이 은행과 거래하던 다른 은행들도 손실을 입고 불안해져 연쇄적으로 문제가 생길 수 있습니다. 마치 한 건물의 불이 옆 건물로 번져 도시 전체를 위협하는 것과 같습니다.
    \item \textbf{기술적 정확한 설명:} 은행은 사회에 금융 서비스를 제공하며 가치를 창출하지만, 개별 은행 또는 여러 은행의 실패(Bank Failure)는 지역 사회, 기업, 가계, 그리고 납세자에게 큰 손실을 초래할 수 있습니다. 특히 은행들은 서로 고도로 연결되어 있어, 한 은행의 실패가 다른 은행으로 파급되는 연쇄 효과(Knock-on Effects)를 일으킬 수 있습니다. 이러한 파급 효과는 개별 은행이 투자 및 자금 조달 결정을 내릴 때 고려하지 않는 부정적 외부 효과입니다. 이는 사회적 관점에서 볼 때 은행의 위험 감수(Risk-taking)가 과도해질 수 있음을 의미합니다합니다.
\end{itemize}

\begin{examplebox}[title=사례: 리먼 브라더스 (Lehman Brothers)]
2008년 리먼 브라더스의 파산은 한 은행의 과도한 위험 감수가 다른 금융 기관과 광범위한 경제에 어떻게 파급될 수 있는지를 보여주는 명확한 사례입니다. 리먼 브라더스는 파산으로 인한 시스템적 실패(Systemic Failures)의 비용을 스스로 부담하지 않았으며, 이는 정부 개입의 필요성을 부각시켰습니다.
\end{examplebox}

\subsubsection{금융 규제의 두 가지 주요 유형}

\begin{itemize}
    \item \textbf{미시 건전성 규제 (Micro-prudential Regulation):}
    \begin{itemize}
        \item \textbf{목표:} 개별 은행의 파산 위험을 줄이는 데 중점을 둡니다.
        \item \textbf{예시:} 은행의 위험 감수 한도를 설정하여 특정 투자에 대한 노출을 제한하는 것 등이 있습니다. 개별 은행이 건전하면 시스템 전체도 안전하다는 전제에 기반합니다.
    \end{itemize}
    \item \textbf{거시 건전성 규제 (Macro-prudential Regulation):}
    \begin{itemize}
        \item \textbf{목표:} 시스템적 실패(Systemic Failures), 즉 한 은행의 실패가 전체 금융 시스템으로 확산되는 것을 줄이는 데 중점을 둡니다.
        \item \textbf{예시:} 대형 은행에 더 엄격한 규제를 부과하여, 이들 은행의 실패가 사회에 미치는 막대한 비용을 줄이려는 시도입니다. 이는 '오염'이 다른 기관으로 확산되는 것을 막는 것에 비유할 수 있습니다.
    \end{itemize}
\end{itemize}

\subsubsection{규제의 다른 목표: 고객 보호 및 공정성}

\begin{itemize}
    \item \textbf{한 줄 핵심 요약:} 규제는 금융 시스템의 안정성뿐만 아니라, 일반 고객들이 불공정한 대우를 받거나 손실을 입지 않도록 보호하는 역할도 합니다.
    \item \textbf{직관적 예시/비유:} 어린이나 노약자가 사기를 당하지 않도록 법이 보호하는 것과 같습니다. 금융 서비스는 복잡하고 전문적이어서 일반 고객이 은행보다 정보가 부족할 수 있습니다. 따라서 정부는 고객이 부당한 수수료를 내거나, 차별적인 대출을 받지 않도록 보호합니다.
    \item \textbf{기술적 정확한 설명:} 규제 당국은 금융 안정성 외에도 고객 보호(Customer Protection)를 중요한 목표로 삼습니다. 이는 소액 투자자와 차입자를 보호하는 것을 의미합니다. 정보가 부족한 소액 예금자들이 은행 파산 시 평생 저축을 잃지 않도록 보호하고, 은행보다 전문성이 떨어지는 차입자들이 금융 서비스를 이용할 때 부당하게 착취당하지 않도록 합니다.
\end{itemize}

\begin{examplebox}[title=사례: 레드라이닝 (Redlining)]
레드라이닝은 미국 은행 대출에서 소득, 성별, 인종을 기반으로 심각한 차별이 존재했던 악명 높은 사례입니다. 특정 지역 전체를 위험하다고 분류하여 해당 지역 주민들에게 대출을 거부하거나 불리한 조건으로 대출을 제공했습니다. 이는 명백히 불공정한 행위였으며, 이러한 차별을 막기 위해 수많은 정부 규제가 도입되어 은행이 모든 지역 사회 구성원에게 서비스를 제공하도록 보장하고 있습니다.
\end{examplebox}

\subsubsection{규제의 상충 관계: 비용과 부작용}

\begin{itemize}
    \item \textbf{한 줄 핵심 요약:} 규제는 사회에 이점을 제공하지만, 은행의 효율성을 떨어뜨리거나 예상치 못한 부작용, 그리고 도덕적 해이를 유발할 수 있습니다.
    \item \textbf{직관적 예시/비유:} 자동차 제조업체는 파산해도 정부가 구제하지 않습니다. 하지만 은행은 다릅니다. 규제는 은행이 안전하게 운영되도록 돕지만, 동시에 은행이 돈을 버는 데 제약을 가합니다. 예를 들어, 특정 위험한 투자를 금지하면 은행의 수익이 줄어들 수 있습니다. 또한, 정부가 은행을 구제해 줄 것이라는 믿음은 은행이 더 위험한 행동을 하도록 부추길 수 있습니다.
    \item \textbf{기술적 정확한 설명:} 규제는 은행에 제약을 가하여 비효율성(Inefficiency)을 초래할 수 있습니다. 특정 거래를 금지하면 은행의 수익성이 저하될 수 있습니다. 또한, 규제는 거의 항상 의도치 않은 결과(Unintended Consequences)를 낳습니다. 예를 들어, 소외된 지역 사회에 대출을 장려하는 규제가 해당 지역의 임대료를 상승시켜 일부 가구가 감당하기 어렵게 만들 수도 있습니다. 가장 큰 문제는 도덕적 해이(Moral Hazard)입니다. 정부가 과도한 위험 감수를 실패로 처벌하지 않고 구제금융(Bailout)을 제공하면, 은행은 '성공하면 내가 이득, 실패하면 정부가 손실'이라는 생각으로 더 큰 위험을 감수하게 됩니다.
\end{itemize}

\begin{examplebox}[title=사례: 부실자산구제프로그램 (TARP)]
2008년 금융 위기 당시, 미국 정부는 부실자산구제프로그램(Troubled Asset Relief Program, TARP)을 통해 대형 은행들의 부실 자산을 매입하여 이들이 파산하는 것을 막았습니다. 많은 경제학자들은 이러한 구제금융이 은행들이 미래에 더 큰 위험을 감수하도록 부추기는 도덕적 해이를 유발한다고 지적했습니다.
\end{examplebox}

\begin{summarybox}[title=결론: 규제의 균형점]
부정적 외부 효과는 정부 개입의 역할을 만듭니다. 은행 실패로 인한 파급 효과는 은행 규제의 필요성을 야기합니다. 이러한 규제는 사회에 이익을 줄 수 있지만, 공짜가 아닙니다. 정부 개입은 은행을 덜 효율적으로 만들거나 나쁜 유인을 초래할 수 있습니다. 규제가 합리적인지는 궁극적으로 공공 정책 문제입니다. 정책 입안자들은 이점과 비용을 비교하여 결정을 내려야 합니다.
\end{summarybox}


\subsection{정부 안전망}

은행 시스템의 불안정성은 금융 시스템과 경제 전반에 매우 파괴적인 영향을 미칠 수 있습니다. 이러한 불안정성, 특히 뱅크런과 금융 공황으로부터 시스템을 보호하기 위해 전 세계 정부는 두 가지 핵심 제도를 마련했습니다: 예금 보험(Deposit Insurance)과 최종 대부자(Lender of Last Resort). 이 두 제도를 통틀어 '정부 안전망(Government Safety Net)'이라고 부릅니다.

\subsubsection{은행의 취약성: 뱅크런의 원인}

\begin{itemize}
    \item \textbf{한 줄 핵심 요약:} 은행은 단기 예금을 받아 장기 대출에 투자하는 '유동성 불일치' 구조 때문에 본질적으로 취약하며, 이는 뱅크런으로 이어질 수 있습니다.
    \item \textbf{직관적 예시/비유:} 은행은 마치 단기 대여금으로 장기 프로젝트에 투자하는 회사와 같습니다. 고객들은 언제든지 돈을 돌려받을 수 있지만(단기 부채), 은행이 투자한 대출은 당장 현금으로 바꾸기 어렵습니다(장기 비유동 자산). 만약 고객들이 갑자기 모두 돈을 돌려달라고 하면, 은행은 현금이 부족해져 문을 닫을 수밖에 없습니다.
    \item \textbf{기술적 정확한 설명:} 은행은 '단기 차입(borrow short)하여 장기 대출(lend long)'하는 구조를 가지고 있습니다. 대부분의 은행 자산은 정보 민감성이 높고 판매하기 어려운 대출과 같은 비유동 자산(Illiquid Assets)입니다. 반면, 은행은 예금자에게 유동성(Liquidity)을 제공하여 예금자가 요구하면 언제든지 자금을 인출할 수 있도록 합니다. 이러한 자산과 부채 간의 유동성 불일치(Liquidity Mismatch)가 은행을 취약하게 만듭니다. 갑작스러운 현금 부족은 은행을 지급 불능 상태로 몰아넣을 수 있으며, 뱅크런(Bank Run)이 발생하면 예금과 부채가 빠르게 인출됩니다. 은행은 대출 포트폴리오를 청산하려 할 수 있지만, 비유동 자산은 손실을 보고 '헐값 매각(fire sale prices)'될 수 있으며, 이마저도 시간이 걸립니다. 이러한 상황을 예상한 은행은 보통 영업을 중단하게 됩니다.
\end{itemize}

\subsubsection{뱅크런의 사회적 비용과 전염 효과}

\begin{itemize}
    \item \textbf{한 줄 핵심 요약:} 개별 은행의 실패는 예금자, 차입자에게 직접적인 피해를 줄 뿐만 아니라, 금융 시스템 전체로 확산되어 경제 활동을 마비시키는 '금융 공황'을 초래할 수 있습니다.
    \item \textbf{직관적 예시/비유:} 한 은행이 망하면 그 은행 고객만 피해를 보는 것이 아닙니다. 다른 은행 고객들도 '혹시 내 은행도?'라는 불안감에 돈을 인출하기 시작합니다. 이는 마치 독감 환자 한 명이 주변 사람들에게 바이러스를 퍼뜨려 도시 전체가 독감에 걸리는 것과 같습니다. 결국 기업들은 대출을 받지 못해 투자를 줄이고, 가계는 소비를 줄여 경제 전체가 침체됩니다.
    \item \textbf{기술적 정확한 설명:} 은행 실패는 예금자에게 예금 손실이나 지급 시스템 접근 불가라는 피해를 주고, 차입자에게는 금융 서비스 접근 불가라는 피해를 줍니다. 정부의 주된 우려는 이러한 초기 충격이 금융 시스템 전체로 확산되는 전염 효과(Contagion)입니다. 예금자와 채권자들이 금융 시스템 전체의 건전성에 대해 우려하게 되면, 전반적인 인출이 발생하는데 이를 금융 공황(Panic)이라고 합니다.
\end{itemize}

\begin{summarybox}[title=전염 효과의 두 가지 주요 원천]
\begin{enumerate}
    \item \textbf{정보 비대칭 (Information Asymmetry):} 예금자와 은행 경영진 간의 정보 불균형에서 발생합니다. 한 은행에 대한 나쁜 소문(진실 여부와 무관하게)이 퍼지면, 예금자들은 자신의 돈을 인출하기 위해 달려갑니다. 이 과정에서 다른 은행의 예금자들도 불안감을 느끼고 예금을 인출할 수 있습니다. 예금자들이 정보가 부족한 한, 소문은 건전한 은행까지 무너뜨릴 수 있습니다.
    \item \textbf{상호 연결성 (Interconnectedness):} 21세기 은행들은 고도로 상호 연결되어 있습니다. 한 대형 은행이 실패하면, 다른 은행들은 도매 자금 조달(Wholesale Funding)에 접근하지 못할 수 있습니다. 또한, 보험이나 파생상품 계약의 상대방 위험(Counter-party Risk)으로 인해 연쇄적인 손실이 발생할 수 있습니다.
\end{enumerate}
\end{summarybox}

금융 공황은 사회에 막대한 비용을 초래합니다. 건전한 은행들까지 파산할 수 있으며, 차입자들이 신용에 접근하지 못하면 기업 투자, 가계 지출, 전반적인 경제 활동이 악화됩니다. 경제가 나빠지면 더 많은 대출이 부실화되어 상황은 더욱 악화됩니다. 금융과 경제 활동 간의 이러한 부정적인 피드백 루프는 깊고 장기적인 경기 침체로 이어질 수 있습니다. 우리는 부실 은행이 실패하는 것을 원하지만, 금융 공황은 원하지 않습니다. 바로 이 지점에서 정부가 개입하여 안전망을 구축합니다.

\subsubsection{정부 안전망의 두 축: 예금 보험과 최종 대부자}

정부 안전망은 예금 보험(Deposit Insurance)과 최종 대부자(Lender of Last Resort)라는 두 가지 제도로 구성되어 뱅크런에 대한 양면 공격을 펼칩니다.

\begin{enumerate}
    \item \textbf{예금 보험 (Deposit Insurance):}
    \begin{itemize}
        \item \textbf{한 줄 핵심 요약:} 은행이 파산해도 예금자의 돈을 일정 한도까지 보장하여, 예금자들이 뱅크런을 할 유인을 없앱니다.
        \item \textbf{직관적 예시/비유:} 자동차 보험과 같습니다. 사고가 나도 보험사가 수리비를 내주기 때문에 운전자는 작은 사고에 대해 크게 걱정하지 않습니다. 예금 보험도 은행이 망해도 정부가 예금을 돌려주기 때문에 예금자들은 소문에 흔들리지 않고 돈을 인출할 필요가 없습니다.
        \item \textbf{기술적 정확한 설명:} 대부분의 국가에는 예금 보험 프로그램이 있습니다. 미국의 경우 1933년부터 연방예금보험공사(Federal Deposit Insurance Corporation, FDIC)가 운영되고 있습니다. 은행들은 예금 보험료를 지불하며, 은행이 파산하면 소매 예금자(일반 고객)는 돈을 잃지 않습니다. 2020년 기준 보험 한도는 25만 달러로, 일반 고객에게는 충분한 보장입니다. FDIC는 파산한 은행을 정리(Resolve)합니다. 은행은 파산 법원이 아닌 FDIC가 처리하며, 보통 파산 은행을 지역의 건전한 다른 은행에 경매로 넘겨 예금 계좌가 해당 은행으로 이전됩니다. 이는 빠르고 쉬운 과정입니다.
        \item \textbf{작동 원리:} 예금 보험은 뱅크런을 예방합니다. 예금자는 은행에 대한 나쁜 소문을 들어도, 최악의 경우 FDIC로부터 돈을 돌려받을 수 있기 때문에 굳이 은행으로 달려가 돈을 인출할 유인이 없습니다. 이는 '사전 예방(before-the-fact)' 장치입니다.
    \end{itemize}
    \item \textbf{최종 대부자 (Lender of Last Resort):}
    \begin{itemize}
        \item \textbf{한 줄 핵심 요약:} 건전하지만 일시적인 현금 부족을 겪는 은행에 중앙은행이 긴급 대출을 제공하여, 진행 중인 뱅크런을 중단시킵니다.
        \item \textbf{직관적 예시/비유:} 갑자기 집에 불이 났는데 소방차가 출동하여 불을 끄는 것과 같습니다. 은행이 갑자기 현금이 부족해져 뱅크런이 시작되면, 중앙은행이 긴급 자금을 공급하여 뱅크런을 진정시킵니다.
        \item \textbf{기술적 정확한 설명:} 미국에는 100년 이상 된 연방준비제도(Federal Reserve System, Fed)라는 중앙은행이 있습니다. 연준의 주요 기능 중 하나는 최종 대부자 역할을 하는 것입니다. 이는 연준이 지급 능력이 있는(Solvent) 은행, 즉 건전하지만 일시적으로 뱅크런을 겪는 은행에 현금 대출을 제공하는 것을 의미합니다. 은행은 대출을 공개 시장에서 헐값에 매각하는 대신, 연준에 고품질 대출을 담보로 단기 현금 대출을 받을 수 있습니다. 은행은 이 현금을 사용하여 예금자들의 인출 요구를 충족시킬 수 있습니다. 예금자들이 돈을 돌려받을 수 있다는 것을 알게 되면, 뱅크런을 할 유인이 사라집니다. 상황이 진정되면 은행은 연준에 대출금을 상환합니다.
        \item \textbf{작동 원리:} 최종 대부자는 건전한 은행에 현금 대출을 제공하여 진행 중인 뱅크런을 중단시킵니다. 이는 '사후 대응(after-the-fact)' 장치입니다.
    \end{itemize}
\end{enumerate}

\subsubsection{사례 연구: 2007-2009년 금융 위기 개입}

2007-2009년 금융 위기 동안 정부 안전망이 어떻게 작동했는지 살펴보는 것은 매우 중요합니다.

\begin{itemize}
    \item \textbf{위기 전 상황:} 2006년 여름 미국 주택 가격이 정점에 달한 후 급락하기 시작하면서 금융 시스템의 취약점이 드러났습니다.
    \item \textbf{금융 스트레스 지표:} TED 스프레드(TED spread)는 금융 부문 스트레스의 일반적인 척도입니다. 이는 은행 간 대출 금리(interbank rate)와 미국 국채 금리(Treasury bill rate)의 차이입니다. 정상 시기에는 차이가 거의 없지만, 위기 시기에는 스프레드가 확대되어 투자자들이 금융 시스템을 더 위험하게 본다는 것을 나타냅니다.
    \item \textbf{위기 전개:}
    \begin{itemize}
        \item \textbf{2007년 8월:} 프랑스 대형 은행 BNP 파리바(BNP Paribas)가 미국 모기지에 대규모 투자했던 투자 펀드 자회사 세 곳의 인출을 중단했습니다. 이후 투자자들은 어떤 금융 기관이 미국 모기지에 노출되어 있는지 우려하기 시작했고, 금융 스트레스가 증가했습니다.
        \item \textbf{2008년 3월:} 미국 대형 투자 은행 베어 스턴스(Bear Stearns)가 도매 자금 조달(Wholesale Funding)에서 뱅크런을 겪었고, 연준은 JP모건 체이스(JPMorgan Chase)를 통한 긴급 구제를 주선해야 했습니다. 스트레스는 더욱 고조되었습니다.
        \item \textbf{2008년 9월:} 리먼 브라더스(Lehman Brothers)의 파산은 광범위한 공황을 촉발했습니다. 금융 스트레스 지표는 전례 없는 수준으로 치솟았고, 신용 시장은 얼어붙었으며, 주식 시장은 폭락했습니다. 경제 전반으로의 파급 효과가 명백해졌습니다.
    \end{itemize}
    \item \textbf{정부 개입:} 정부는 금융 공황을 막고 경제적 파급 효과를 제한하기 위해 개입했습니다. 주요 개입은 다음과 같습니다.
    \begin{itemize}
        \item \textbf{예금 보험 확대:} 2008년 10월 8일, 리먼 파산 한 달도 채 되지 않아 예금 보험 한도가 10만 달러에서 25만 달러로 인상되었습니다. 이 개입의 목표는 미래의 뱅크런을 예방하는 것이었습니다.
        \item \textbf{최종 대부자 역할:} 연방준비제도는 유동성 지원 시설(Liquidity Facilities)을 통해 임시 현금 대출을 제공했습니다. 2008년 초부터 운영되었지만, 공황 기간 동안 크게 확대되었습니다. 이 시설들은 모두 진행 중인 뱅크런을 중단시키는 동일한 목적을 가졌습니다. 상황이 진정된 후 연준은 대출금을 상환받았습니다.
    \end{itemize}
\end{itemize}

\begin{warningbox}[title=정부 안전망의 함정: 도덕적 해이]
정부 안전망은 은행 경영진에게 실패에 대한 완충 역할을 제공하여 도덕적 해이를 유발할 수 있습니다. 마치 보험에 가입한 운전자가 작은 사고에 대해 덜 걱정하는 것처럼, 은행 경영진은 더 큰 위험을 감수하려 할 수 있습니다. 앞으로 은행 규제와 감독이 이 문제를 어떻게 다루는지 살펴보겠습니다.
\end{warningbox}


\section{규제 및 감독의 실제}

금융 시스템의 심각한 혼란은 사회에 막대한 비용을 초래할 수 있으므로, 은행은 가장 엄격하게 규제받는 기업 중 하나입니다. 이 섹션에서는 정부의 은행 규제가 어떻게 작동하는지, 규제 당국이 은행에 어떤 제약을 가하는지, 그리고 은행이 규칙을 위반했을 때 어떤 일이 발생하는지 설명합니다.

\subsection{은행 규제}

\subsubsection{규제와 감독의 구분}

\begin{itemize}
    \item \textbf{한 줄 핵심 요약:} 규제는 은행이 따라야 할 '규칙'을 만드는 것이고, 감독은 그 규칙이 잘 지켜지는지 '감시하고 집행'하는 것입니다.
    \item \textbf{직관적 예시/비유:} 규제는 도로교통법을 만드는 것과 같습니다. "시속 60km 이상으로 달리지 마시오"와 같은 규칙을 정하는 것이죠. 반면 감독은 경찰이 도로에서 과속하는 차량을 단속하고 벌금을 부과하는 것과 같습니다.
    \item \textbf{기술적 정확한 설명:}
    \begin{itemize}
        \item \textbf{은행 규제 (Bank Regulation):} 정부가 은행이 따라야 할 구체적인 규칙(Specific Rules)을 제정하는 것입니다. 이 규칙들은 보통 정부가 통과시킨 법률(예: 글래스-스티걸 법)의 권한 아래 작성됩니다. 예를 들어, 상업 은행과 투자 은행의 분리를 요구하는 규칙을 만드는 것입니다.
        \item \textbf{은행 감독 (Bank Supervision):} 은행 시스템에 대한 지속적인 감시(Continuous Oversight)를 포함합니다. 은행이 규칙을 위반하면 처벌(Enforcement Actions)이 따릅니다. 감독은 개별 은행의 안전성과 건전성(Safety and Soundness)에 더 좁게 초점을 맞춥니다. 정부는 은행의 재무 보고서를 검사하고, 현장 검사(On-site Examinations)를 수행하며, 연간 스트레스 테스트(Stress Testing)를 실시합니다.
    \end{itemize}
\end{itemize}

\subsubsection{무엇이 규제되는가?}

은행 규제는 크게 경쟁, 활동 및 투자, 자금 조달, 그리고 공시 요건의 네 가지 영역에 걸쳐 이루어집니다.

\begin{enumerate}
    \item \textbf{경쟁 (Competition):}
    \begin{itemize}
        \item \textbf{내용:} 은행을 개설하려면 연방 또는 주 규제 당국의 승인을 받아야 합니다. 은행 합병(Mergers)도 정부 승인이 필요합니다. 최근까지 은행이 지점을 개설할 수 있는 지리적 범위도 제한되었습니다. 산업 퇴출(Industry Exit)의 경우, 은행에는 파산 절차가 없으며 FDIC가 부실 은행을 정리합니다.
        \item \textbf{목표:} 과도한 경쟁은 은행이 생존을 위해 도박을 하게 만들 수 있고, 너무 적은 경쟁은 고객 착취로 이어질 수 있으므로, 이 사이에서 균형을 맞추기 위함입니다.
    \end{itemize}
    \item \textbf{활동 및 투자 (Activities and Investments):}
    \begin{itemize}
        \item \textbf{내용:} 과거에는 예금 수취 상업 은행이 증권 인수(Securities Underwriting)를 할 수 없었습니다. 현재는 자기 계정 거래(Proprietary Trading)에 제한이 있습니다. 은행은 위험성이 높은 주식(Stocks)을 보유할 수 없으며, 정크 본드(Junk Bond) 보유에도 제한이 있습니다.
        \item \textbf{목표:} 이러한 규제는 은행의 위험 감수(Risk-taking)를 직접적으로 통제하여 은행의 건전성을 확보하려는 것입니다.
    \end{itemize}
    \item \textbf{자금 조달 (Funding):}
    \begin{itemize}
        \item \textbf{내용:} 은행은 레버리지(Leverage)에 대한 상한선을 가지고 있습니다. 이는 자기자본 조달(Equity Funding)에 대한 최소 요건으로 표현됩니다. 도매 자금 조달(Wholesale Funding)의 사용도 면밀히 감시됩니다.
        \item \textbf{목표:} 손실에 대한 완충 역할을 할 수 있는 안정적인 자금 조달은 은행의 실패 위험을 줄이는 데 도움이 됩니다.
    \end{itemize}
    \item \textbf{공시 요건 (Disclosure Requirements):}
    \begin{itemize}
        \item \textbf{내용:} 은행은 고객에게 신용 카드 비용을 명확히 알려야 합니다. 또한, 표준화된 형식으로 정기적으로 재무 정보를 공개해야 합니다.
        \textbf{목표:} 이는 규제 당국과 금융 시장 모두가 과도한 위험을 감수하는 은행을 규율할 수 있도록 합니다.
    \end{itemize}
\end{enumerate}

\subsubsection{누가 규제하는가? (미국 중심)}

미국에서는 여러 기관이 은행을 규제하고 감독합니다.

\begin{itemize}
    \item \textbf{연방예금보험공사 (FDIC):} 모든 은행이 예금 보험에 가입되어 있으므로 FDIC가 최우선 규제 기관입니다.
    \item \textbf{연방준비제도 (Federal Reserve System):} 중앙은행으로서 중요한 규제 역할을 합니다.
    \item \textbf{통화감독청 (Office of the Comptroller of the Currency, OCC):} 미국 재무부 산하 기관입니다.
    \item \textbf{주 정부 당국:} 많은 소규모 은행들은 주(State) 차원에서 규제됩니다.
\end{itemize}
미국은 여러 규제 기관이 존재하는 점에서 다소 독특합니다. 다른 국가들은 보통 중앙은행 내에 단일 은행 규제 기관을 두는 경우가 많습니다.

\subsubsection{사례 연구: 바젤 협약과 자본 요건}

\begin{itemize}
    \item \textbf{한 줄 핵심 요약:} 바젤 협약은 전 세계 은행들이 최소한의 자기자본을 보유하도록 요구하여, 은행의 손실 흡수 능력을 강화하고 국제적인 경쟁의 공정성을 유지하려는 국제적 합의입니다.
    \item \textbf{직관적 예시/비유:} 전 세계 자동차 경주에서 모든 참가 차량이 최소한의 안전 장치(에어백, 안전벨트 등)를 갖추도록 의무화하는 것과 같습니다. 어떤 나라의 차량은 안전 장치가 적어 더 가볍고 빠르다면 불공평하겠죠. 바젤 협약은 은행의 '안전 장치'인 자기자본을 국제적으로 표준화하는 것입니다.
    \item \textbf{기술적 정확한 설명:} 바젤 협약(Basel Accord)은 은행 위험 감수에 대한 '글로벌 최소세'로 볼 수 있습니다. 1980년대 은행업이 점점 더 글로벌화되면서, 규제가 느슨한 국가의 은행이 규제가 엄격한 국가의 은행보다 경쟁 우위를 가질 수 있다는 문제가 발생했습니다. 이는 규제 당국 간의 '규제 완화 경쟁(race to the bottom)'을 유발할 수 있었습니다.
    \item \textbf{바젤 협약의 해결책:} 1988년 스위스 바젤에서 전 세계 은행 규제 당국자들이 모여 모든 국가의 은행에 대해 공정한 경쟁 환경(Level Playing Field)을 조성하기 위한 새로운 규칙을 개발했습니다. 바젤 협약은 은행의 손실 흡수 능력(Loss-bearing Capacity)에 초점을 맞췄습니다. 은행이 투자를 자기자본(Equity Capital)으로 조달하도록 장려함으로써, 더 많은 자본은 잠재적 손실에 대한 더 큰 완충 역할을 하여 더 안전한 글로벌 은행 시스템을 구축하는 것을 목표로 했습니다.
    \item \textbf{최소 자본 요건 (Minimum Capital Requirements):}
    \begin{itemize}
        \item 대략적으로 은행은 자산 1달러당 최소 0.08달러의 자기자본으로 자금을 조달해야 합니다.
        \item 이 규칙은 은행 투자의 위험도를 고려합니다. 예를 들어, 매우 안전한 AAA 등급 국채에만 투자하는 은행은 더 적은 자기자본 완충이 필요합니다. 즉, 자본 요건은 위험 기반(Risk-based)입니다.
        \item 시간이 지남에 따라 바젤 1(1988년)에서 바젤 3(2010년)으로 진화했지만, 기본적인 아이디어는 동일합니다: 위험 투자 1달러당 최소 X센트의 자기자본을 보유해야 합니다. 2017년에는 최소 6센트였습니다.
    \end{itemize}
\end{itemize}

\begin{examplebox}[title=미국 대형 은행의 자본 비율 (2017년 데이터)]
JP모건, 뱅크 오브 아메리카 등 미국 대형 은행들은 2017년 기준 최소 자본 요건을 충분히 충족했습니다. 이는 2020년 코로나19 팬데믹이 닥쳤을 때 이들 은행이 잘 대응할 수 있었던 이유 중 하나로 평가됩니다. 여기서 사용되는 자본 비율은 단순히 '자기자본/자산'이 아니라 '자기자본/위험 가중 자산(Risk-weighted Assets)'으로, 위험도를 반영한 복잡한 계산을 거칩니다.
\end{examplebox}

\begin{summarybox}[title=은행 규제 요약]
은행 규제는 은행에 제약을 가하고, 특정 위험 활동을 금지하며, 은행이 간과할 수 있는 행동을 장려합니다. 바젤 협약은 최소 자본 요건을 설정하여 글로벌 표준을 제공하고 공정한 경쟁 환경을 보장하는 중요한 사례입니다.
\end{summarybox}


\subsection{감독 과정}

은행 규제가 정부가 규칙을 제정하는 것이라면, 감독(Supervision)은 그 규칙이 실제로 잘 지켜지고 있는지 확인하는 과정입니다. 규칙을 위반하면 반드시 그에 따른 결과가 있어야 합니다. 감독은 은행의 건전성을 모니터링하기 위해 검사(Exams)와 스트레스 테스트(Stress Tests)를 활용합니다.

\subsubsection{누가 감독하는가?}

규칙을 제정하는 기관과 동일한 기관들이 은행을 감독합니다. 즉, FDIC, 연준(Fed), OCC, 그리고 주(State) 당국이 은행 감독을 담당합니다. 이 기관들은 규칙을 실행에 옮깁니다.

\subsubsection{규칙 위반 시 제재: 집행 조치}

\begin{itemize}
    \item \textbf{한 줄 핵심 요약:} 은행이 규정을 위반하면 규제 당국은 시정 조치를 요구하고, 불이행 시 벌금, 활동 제한, 또는 개인 제재와 같은 처벌을 가합니다.
    \item \textbf{직관적 예시/비유:} 학교에서 교칙을 어긴 학생에게 벌점을 주거나, 반성문을 쓰게 하고, 심하면 정학 처분을 내리는 것과 같습니다. 은행도 규칙을 어기면 '벌점'을 받고, '반성문'에 해당하는 시정 계획을 제출해야 하며, 심하면 '영업 정지'와 같은 제재를 받습니다.
    \item \textbf{기술적 정확한 설명:} 은행이 규칙을 위반할 경우, 규제 당국은 특정 조치(Certain Actions)를 요구하고 불이행 시 제재(Penalize Noncompliance)할 권한을 가집니다. 일반적으로 규제 당국은 은행에 절차상의 결함을 시정하거나 자기자본 확충과 같은 조치를 요구합니다. 은행이 이를 준수하지 않으면, 특정 개인의 은행업 종사 금지, 특정 M\&A 활동 금지, 배당금 지급 금지, 심지어 벌금(Fines)과 같은 처벌이 따를 수 있습니다.
\end{itemize}

\begin{examplebox}[title=사례: 런던 고래 (London Whale) 사건]
2013년 연준은 JP모건 체이스의 '런던 고래' 사건과 관련하여 집행 조치를 발표했습니다. 이 사건에서 한 트레이더가 은행의 부실한 감시 하에 막대한 손실을 입혔습니다. 연준은 은행의 위험 관리(Risk Management) 결함을 지적하고, 이를 시정하기 위한 신속한 조치를 요구했으며, 2억 달러의 벌금을 부과했습니다. 이는 영국 규제 당국의 벌금과 합쳐 약 10억 달러에 달하는 막대한 금액이었습니다. 이러한 처벌은 JP모건이 60억 달러의 손실을 입고 은행을 위험에 빠뜨린 사건에 대한 것이었으며, 적절한 절차가 있었다면 피할 수 있었을 것입니다.
\end{examplebox}

\begin{figure}[h!]
    \centering
    \begin{adjustbox}{width=0.9\textwidth,center}
    \begin{tabular}{ll}
        \toprule
        \textbf{위반 유형} & \textbf{주요 벌금 사례 (2008-2015)} \\
        \midrule
        미판매 미국 모기지 및 관련 증권 & 가장 큰 비중을 차지 \\
        런던 고래 사건 & JP모건 체이스에 2억 달러 벌금 \\
        기타 & 다양한 규제 위반에 대한 벌금 \\
        \bottomrule
    \end{tabular}
    \end{adjustbox}
    \caption{은행별 및 위반 유형별 벌금 (2008-2015년 개념적 요약)}
    \label{tab:bank_fines}
\end{figure}
\vspace{0.5em}
위 표는 2008년부터 2015년까지 은행별 및 위반 유형별 벌금 현황을 개념적으로 보여줍니다. 미판매 미국 모기지 및 관련 증권과 관련된 벌금이 가장 큰 비중을 차지했으며, 런던 고래 사건과 같은 개별 사례도 상당한 벌금을 초래했습니다. 이러한 벌금은 은행이 규칙을 위반할 경우 반드시 결과가 따른다는 것을 보여주며, 은행이 규칙을 준수하고 적절한 위험 관리 관행을 채택하도록 장려하는 역할을 합니다.

\subsubsection{감독 과정의 핵심: 공식 은행 검사}

\begin{itemize}
    \item \textbf{한 줄 핵심 요약:} 감독관은 은행의 재무 보고서를 분석하고 직접 방문하여 은행의 건강 상태를 종합적으로 평가합니다.
    \item \textbf{직관적 예시/비유:} 의사가 환자의 건강 상태를 진단하는 것과 같습니다. 혈액 검사(재무 보고서 분석)를 통해 전반적인 상태를 파악하고, 직접 진찰(현장 검사)을 통해 더 깊이 있는 문제를 찾아냅니다.
    \item \textbf{기술적 정확한 설명:} 공식 은행 검사(Formal Bank Examinations)는 감독 과정의 핵심 부분입니다. 모든 은행은 검사를 받아야 합니다.
    \begin{itemize}
        \item \textbf{비현장 검사 (Offsite Component):} 은행은 규제 당국에 분기별 재무 보고서를 제출해야 합니다. 감독관은 이 보고서를 분석하여 잠재적인 약점을 가진 은행을 찾아냅니다.
        \item \textbf{현장 검사 (Onsite Component):} 검사관이 최소 연 1회 은행을 방문합니다. 때로는 예고 없이 방문하기도 하며, 대형 은행의 경우 검사관이 연중 상주하기도 합니다. 모든 경우에 검사관은 은행의 사업을 심층적으로 조사하여 은행의 상태와 성과를 파악합니다.
    \end{itemize}
\end{itemize}

\subsubsection{CAMELS 평가 시스템}

\begin{itemize}
    \item \textbf{한 줄 핵심 요약:} 은행 검사의 결과는 'CAMELS'라는 6가지 핵심 기준에 따라 은행의 건전성을 평가하는 점수로 나타납니다.
    \item \textbf{직관적 예시/비유:} 학생의 성적표와 같습니다. 국어, 영어, 수학 등 여러 과목(CAMELS 기준)에서 점수를 받고, 이를 종합하여 총점(전반적인 점수)이 나옵니다. 만약 낙제하면 '문제 학생 명단'에 오르겠죠.
    \item \textbf{기술적 정확한 설명:} 검사 과정의 결과는 CAMELS 점수(CAMELS Score)로 나타납니다. 이 점수는 은행이 검사를 통과하는지 여부를 결정합니다. 은행은 다음 6가지 핵심 기준에 따라 평가됩니다.
    \begin{itemize}
        \item \textbf{C}apital adequacy (자본 적정성): 은행의 손실 흡수 능력.
        \item \textbf{A}sset quality (자산 건전성): 자산이 부실화되고 있는지, 아니면 양호한 상태인지.
        \item \textbf{M}anagement (경영): 은행이 법규 및 규정을 잘 준수하며 운영되고 있는지.
        \item \textbf{E}arnings (수익): 은행의 투자가 수익을 창출하고 있는지.
        \item \textbf{L}iquidity (유동성): 은행의 유동성 위험 원천은 무엇이며, 충분한 현금을 보유하고 있는지.
        \item \textbf{S}ensitivity to market risk (시장 위험 민감도): 시장 위험에 대한 은행의 민감도.
    \end{itemize}
    각 기준에 대한 점수가 부여되고, 전반적인 점수가 주어집니다. 은행이 검사에 실패하면 '문제 은행 목록(Problem Bank List)'에 오르게 됩니다. 문제 은행은 공개적으로 명단이 발표되지는 않지만, 규제 당국은 은행에 검사 실패 사실을 비공개로 통보합니다. 이들 은행은 규제 당국의 특정 우려 사항을 해결할 때까지 제한을 받게 됩니다.
\end{itemize}

\begin{figure}[h!]
    \centering
    \begin{adjustbox}{width=0.9\textwidth,center}
    \begin{tabular}{ccc}
        \toprule
        \textbf{연도} & \textbf{문제 은행 수} & \textbf{총 관리 자산 (개념적)} \\
        \midrule
        2008 & (높음) & (높음) \\
        ... & ... & ... \\
        2021 & (낮음) & (낮음) \\
        \bottomrule
    \end{tabular}
    \end{adjustbox}
    \caption{문제 은행 목록 (2008-2021년 개념적 추이)}
    \label{tab:problem_banks}
\end{figure}
\vspace{0.5em}
위 표는 2008년부터 2021년까지 문제 은행의 수와 총 관리 자산의 추이를 개념적으로 보여줍니다. 이 수치는 경기 순환적(Cyclical) 특성을 가집니다. 경기 침체기에는 은행의 결함이 드러나면서 문제 은행의 수가 증가하는 경향이 있습니다.


\subsection{스트레스 테스트}

스트레스 테스트(Stress Testing)는 감독 과정에 비교적 최근에 추가된 요소로, 2007-2009년 위기 이후 인기를 얻었습니다.

\subsubsection{스트레스 테스트의 목표와 작동 방식}

\begin{itemize}
    \item \textbf{한 줄 핵심 요약:} 스트레스 테스트는 은행이 극심한 경제 위기 상황에서도 충분한 자본을 보유하여 생존할 수 있는지 미리 평가하는 제도입니다.
    \item \textbf{직관적 예시/비유:} 건물을 지을 때 지진이나 태풍 같은 최악의 상황을 가정하여 건물이 무너지지 않도록 설계하는 것과 같습니다. 은행도 최악의 경제 상황을 가정하여 얼마나 튼튼한지 미리 시험하는 것입니다.
    \item \textbf{기술적 정확한 설명:} 스트레스 테스트의 목표는 은행이 대규모 잠재적 손실에 대해 생각하도록 강제하는 것입니다. 이는 은행이 심각한 경제 및 금융 스트레스 기간에도 생존할 수 있는 충분한 자본을 보유하고 있는지 확인하는 것입니다. 미래 지향적인 방식으로 은행의 회복 탄력성(Resilience)을 테스트합니다.
\end{itemize}

\begin{summarybox}[title=스트레스 테스트 절차]
\begin{enumerate}
    \item \textbf{경제 시나리오 제공:} 은행은 향후 몇 년간의 경제 시나리오를 받습니다.
    \begin{itemize}
        \item \textbf{기준 시나리오 (Baseline Scenario):} GDP, 실업률, 주택 가격 등이 현재 추세대로 지속되는 상황.
        \item \textbf{불리한 시나리오 (Adverse Scenario):} 경제가 붕괴하는 최악의 시나리오.
        \item \textbf{심각하게 불리한 시나리오 (Severely Adverse Scenario):} 경제가 더욱 심각하게 붕괴하는 최악 중의 최악 시나리오.
    \end{itemize}
    \item \textbf{손실 예측 및 자본 계획 제출:} 은행은 각 시나리오 하에서 예상 손실을 예측하고, 이에 기반하여 자본 계획(Capital Plan)을 제출합니다. 이 계획에는 이익을 은행 내에 유보할지, 배당금으로 지급할지, 폭풍을 견디기 위해 얼마나 추가 주식을 발행할지 등이 포함됩니다.
    \item \textbf{규제 당국의 검토:} 이 계획은 규제 당국(주로 연준)에 제출되어 검토됩니다. 연준은 은행이 제안한 계획을 바탕으로 은행의 성과에 대한 자체 예측을 수행합니다.
    \item \textbf{결과 공개:} 스트레스 테스트 결과는 대중에게 공개됩니다. 은행들은 금융 시장에 놀라움을 주지 않기 위해, 그리고 '공개적인 망신'을 피하기 위해 이 결과를 통과하기를 원합니다.
\end{enumerate}
\end{summarybox}

\subsubsection{사례 연구: SCAP와 CCAR}

\begin{enumerate}
    \item \textbf{감독 자본 평가 프로그램 (Supervisory Capital Assessment Program, SCAP):}
    \begin{itemize}
        \item \textbf{한 줄 핵심 요약:} 2009년 금융 위기 중 실시된 최초의 스트레스 테스트로, 은행 시스템의 부족한 자본을 평가하고 확충을 유도했습니다.
        \item \textbf{기술적 정확한 설명:} SCAP는 2009년 초 금융 위기 중에 은행 시스템의 부족한 자본을 평가하기 위해 실시된 최초의 스트레스 테스트였습니다. 이는 '각 은행이 생존하기 위해 얼마나 더 많은 손실 흡수 능력이 필요한가?'라는 질문에 답했습니다.
        \item \textbf{과정:} 19개 대형 은행을 대상으로 2009년과 2010년의 신용 손실을 기준 시나리오와 불리한 시나리오 하에서 추정하도록 요구했습니다. 불리한 시나리오는 장기적인 경기 침체, 실업률 10\% 이상, 주택 가격 25\% 추가 하락 등을 가정했습니다.
        \item \textbf{결과:} 2009년 5월에 발표된 결과에 따르면, 19개 은행 중 10개 은행이 자본 부족(Capital Deficient) 상태로 판명되었습니다. 미국 정부는 이들 은행에 2009년 11월까지 부족한 자본을 민간에서 조달하거나, 그렇지 않으면 미국 재무부가 필요한 자기자본을 제공하고 향후 배당금을 소각할 것이라는 최후통첩을 보냈습니다. 한 달 이내에 거의 모든 은행이 부족한 자기자본을 조달할 수 있었습니다.
    \end{itemize}
    \item \textbf{종합 자본 적정성 검토 (Comprehensive Capital Adequacy Review, CCAR):}
    \begin{itemize}
        \item \textbf{한 줄 핵심 요약:} SCAP의 성공에 힘입어 대형 은행을 대상으로 매년 실시되는 정기적인 스트레스 테스트입니다.
        \item \textbf{기술적 정확한 설명:} SCAP의 성공을 바탕으로 규제 당국은 대형 은행에 대해 스트레스 테스트를 지속적으로 실시하도록 했습니다. CCAR은 SCAP와 동일하게 '최악의 시나리오에서 얼마나 많은 자본이 부족한가?'라는 문제를 다룹니다.
        \item \textbf{과정:} 매년 실시되며, 세 가지 시나리오(기준, 불리, 심각하게 불리)를 사용하고, 은행은 2년 후의 손실을 예측해야 합니다. 은행은 연준의 검토와 대중 공개를 위해 자본 계획을 계속 제공합니다. 금융 시장이 테스트에 실패한 은행을 규율할 것이라는 믿음이 있습니다.
    \end{itemize}
\end{enumerate}

\begin{examplebox}[title=CCAR 심각하게 불리한 시나리오 (2012년)]
2012년 CCAR의 심각하게 불리한 시나리오는 은행에 제공된 28개의 거시경제 변수를 포함했습니다. GDP 성장률, 실업률, 주식 시장, 주택 가격 등 모든 지표가 이 시나리오에서는 악화되는 것으로 가정되었습니다. 은행들은 이 시나리오에 따라 계획을 제출하고, 연준은 이를 검토한 후 결과를 발표합니다. 2012년 주기에서는 Ally Bank를 제외한 모든 은행이 연준이 설정한 5\% 최소 자본 비율을 통과했습니다.
\end{examplebox}

\begin{figure}[h!]
    \centering
    \begin{adjustbox}{width=0.9\textwidth,center}
    \begin{tabular}{cc}
        \toprule
        \textbf{은행} & \textbf{최악 시나리오 하 자기자본 비율 (2012년 개념적)} \\
        \midrule
        JP모건 & (통과) \\
        뱅크 오브 아메리카 & (통과) \\
        ... & ... \\
        Ally Bank & 5\% 미만 (실패) \\
        \bottomrule
    \end{tabular}
    \end{adjustbox}
    \caption{CCAR 최악 시나리오 하 자기자본 비율 (2012년 개념적)}
    \label{tab:ccar_results}
\end{figure}
\vspace{0.5em}
위 표는 2012년 CCAR에서 최악의 시나리오 하에서 각 은행의 자기자본 비율을 개념적으로 보여줍니다. Ally Bank를 제외한 모든 대형 은행이 연준이 설정한 5\%의 최소 자본 비율을 통과했습니다.

\begin{summarybox}[title=감독 과정 요약]
감독은 은행 규칙과 규제가 실제로 적용되는 곳입니다. 은행은 정기적으로 현장 및 비현장 검사를 받으며, 최근 몇 년간 대형 은행들은 스트레스 테스트를 받았습니다. 은행에 어떤 결함이 발견되면 집행 조치가 뒤따를 수 있습니다. 목표는 허용 가능한 위험 한도 내에서 운영되는 경쟁력 있고 성과가 좋은 은행 시스템을 갖는 것입니다.
\end{summarybox}


\section{21세기 규제 당국의 과제}

정부가 은행 부문을 어떻게 규제하고 감독하는지 살펴보았습니다. 우리는 효율적이면서도 안전한 은행 시스템을 원합니다. 이 섹션에서는 은행 규모의 역할과 규제 당국이 직면한 주요 과제에 대해 논의합니다.

\subsection{대마불사 (Too Big to Fail)}

\subsubsection{미국 은행 시스템의 집중화}

\begin{itemize}
    \item \textbf{한 줄 핵심 요약:} 미국 은행 시스템은 소수의 거대 은행들이 전체 대출 및 예금 활동의 대부분을 차지하며 집중되어 있습니다.
    \item \textbf{직관적 예시/비유:} 한 마을에 수많은 작은 가게들이 있지만, 실제로는 몇 개의 대형 마트가 모든 물건을 팔고 있는 것과 같습니다. 은행도 수천 개가 있지만, 소수의 '메가 은행'이 금융 시장을 지배합니다.
    \textbf{기술적 정확한 설명:} 미국 은행 시스템은 집중화(Concentrated)되어 있습니다. 이는 소수의 은행이 대출 및 예금 활동의 대부분을 통제한다는 의미입니다. JP모건, 뱅크 오브 아메리카, 씨티그룹, 웰스 파고 등 약 10개의 미국 메가 은행이 존재합니다.
\end{itemize}

\begin{figure}[h!]
    \centering
    \begin{adjustbox}{width=0.9\textwidth,center}
    \begin{tabular}{lc}
        \toprule
        \textbf{은행 유형} & \textbf{전체 은행 자산 점유율 (2018년 개념적)} \\
        \midrule
        메가 은행 (각 2,500억 달러 이상) & 약 50\% \\
        커뮤니티 은행 (각 10억 달러 미만) & 약 15\% \\
        \bottomrule
    \end{tabular}
    \end{adjustbox}
    \caption{미국 은행 시스템 자산 집중도 (2018년 FDIC 데이터 기반 개념적)}
    \label{tab:bank_concentration}
\end{figure}
\vspace{0.5em}
위 표는 2018년 FDIC 데이터를 기반으로 미국 은행 시스템의 자산 집중도를 개념적으로 보여줍니다. 수천 개의 커뮤니티 은행이 존재하지만, 소수의 메가 은행이 전체 은행 자산의 약 50\%를 차지하고 있습니다.

\subsubsection{왜 대형 은행이 존재하는가?}

대형 은행이 존재하는 이유에 대한 두 가지 상호 보완적인 관점이 있습니다.

\begin{enumerate}
    \item \textbf{효율성 관점:}
    \begin{itemize}
        \item \textbf{내용:} 은행이 대규모로 성장하는 것이 효율적이라는 주장입니다. 규모의 경제(Economies of Scale)를 활용하기 위해 은행은 규모를 키워야 합니다. 또한, 여러 제품을 판매하여 고정 비용을 공유함으로써 범위의 경제(Economies of Scope)를 활용하기 위해 더 복잡해져야 합니다.
    \end{itemize}
    \item \textbf{정부 보조금 관점 (대마불사):}
    \begin{itemize}
        \item \textbf{내용:} 은행들은 정부가 가장 큰 금융 기관의 실패를 용납하지 않을 것이라는 점을 이해하고 있다는 다소 어두운 관점입니다. 만약 메가 은행이 어려움에 처하면 정부가 구제금융을 제공할 것이라는 기대가 있습니다. 이는 최종 대부자로부터 대출을 받거나, 정부가 신속하게 합병을 주선하거나, 심지어 정부가 자기자본을 대가로 현금을 제공하는 형태일 수 있습니다.
        \item \textbf{결과:} 이러한 '대마불사(Too Big to Fail, TBTF)' 정책은 은행이 이러한 보조금을 포착하기 위해 더 커지도록 장려합니다. 구제금융은 대마불사 정책이 실제로 작동하는 사례입니다.
    \end{itemize}
\end{enumerate}

\begin{examplebox}[title=대마불사 정책의 실제 사례]
\begin{itemize}
    \item \textbf{콘티넨탈 일리노이 은행 (Continental Illinois Bank):} 이 은행이 파산 위기에 처했을 때의 이례적인 처리는 '대마불사'라는 용어를 탄생시켰습니다.
    \item \textbf{2008년 영국 은행:} 로열 뱅크 오브 스코틀랜드(Royal Bank of Scotland)를 포함한 영국 은행들은 파산 대신 부분적으로 국유화되었습니다.
    \item \textbf{미국 은행 구제금융:} 부실자산구제프로그램(TARP)을 통해 수많은 미국 은행들이 자기자본을 대가로 현금을 받았습니다. 골드만삭스(Goldman Sachs)와 모건 스탠리(Morgan Stanley)는 투자 은행에서 상업 은행 지주 회사로 전환하여 이러한 정부 지원에 접근할 수 있었습니다.
    \item \textbf{정부 후원 기업 (GSEs):} 패니 매(Fannie Mae)와 프레디 맥(Freddie Mac)도 정부로부터 막대한 지원을 받았습니다.
    \item \textbf{보험 회사:} 2008년에는 일부 대형 보험 회사들도 정부 지원을 받았습니다.
\end{itemize}
\end{examplebox}

\subsubsection{시스템 리스크와 규제}

\begin{itemize}
    \item \textbf{한 줄 핵심 요약:} 시스템 리스크는 한 금융 기관의 실패가 전체 금융 시스템과 경제에 연쇄적으로 파급되는 위험을 의미하며, 이는 '오염'과 유사하게 규제되어야 합니다.
    \item \textbf{직관적 예시/비유:} 한 공장의 오염 물질이 강물 전체를 오염시키고 생태계를 파괴하는 것과 같습니다. 개별 공장은 자신의 오염이 전체 환경에 미치는 영향을 고려하지 않지만, 사회는 이를 막기 위해 '오염세'를 부과합니다. 시스템 리스크도 마찬가지로, 개별 은행은 자신의 실패가 전체 시스템에 미치는 영향을 고려하지 않으므로 규제가 필요합니다.
    \item \textbf{기술적 정확한 설명:} 금융 기관이 실패할 때 금융 시스템과 광범위한 경제에 파급 효과가 발생할 수 있는데, 시스템 리스크(Systemic Risk)는 이러한 잠재적 파급 효과를 포착하는 개념입니다. 최악의 시나리오에서는 금융 시스템을 통한 지급 및 신용 흐름이 갑자기 멈추고 경제가 마비될 수 있습니다.
    \item \textbf{시스템 리스크를 유발하는 기업:} 규모가 큰 기업은 실패 시 더 큰 파급 효과를 가질 가능성이 높습니다. 또한, 규모는 작지만 더 복잡하고 상호 연결된 금융 기업도 중요할 수 있습니다. 이는 대형 및 상호 연결된 금융 기업이 전통적인 대출 및 차입 활동뿐만 아니라 파생상품과 같은 금융 계약을 통해 기업 및 다른 금융 기업과 더 많은 연결 고리를 가질 가능성이 높기 때문입니다.
    \item \textbf{시스템 리스크 규제:} 과거에는 규제 당국이 개별 은행의 실패 위험을 최소화하려고 노력했습니다. 최근에는 거시 건전성 접근 방식(Macroprudential Approaches)이 대신 금융 기업 간의 파급 효과를 목표로 합니다. 이는 오염에 대한 세금과 유사합니다.
\end{itemize}

\subsubsection{금융 안정성 감독 위원회 (FSOC)}

\begin{itemize}
    \item \textbf{한 줄 핵심 요약:} FSOC는 2010년에 설립된 미국 기관으로, 금융 위기를 예방하기 위해 시스템 리스크를 관리하는 것을 목표로 합니다.
    \item \textbf{기술적 정확한 설명:} 금융 안정성 감독 위원회(Financial Stability Oversight Council, FSOC)는 2010년에 금융 위기를 예방하기 위한 시스템 리스크 관리라는 야심찬 목표를 달성하기 위해 설립되었습니다.
    \item \textbf{구성 및 역할:} FSOC는 재무부, 연방준비제도, SEC, FDIC, 상품선물거래위원회(CFTC) 등 광범위한 금융 규제 당국의 리더들을 한자리에 모읍니다. 이들은 개별 시장 감독에만 집중하는 대신, 금융 시스템 전체에 대한 총체적인 관점(Holistic View)을 가지고 금융 시스템의 안전성에 대해 심도 있게 논의합니다.
    \item \textbf{SIFI 지정 권한:} FSOC는 모든 민간 기업을 시스템적으로 중요한 금융 기관(Systemically Important Financial Institution, SIFI)으로 지정할 권한을 가집니다. SIFI로 지정된 기업은 강화된 감독(Enhanced Oversights)을 받게 됩니다.
\end{itemize}

\subsubsection{글로벌 시스템적으로 중요한 은행 (GSIB)}

\begin{itemize}
    \item \textbf{한 줄 핵심 요약:} GSIB는 국제 은행 규제 당국에 의해 지정된, 전 세계에서 가장 크고 상호 연결된 은행들로, 시스템 리스크에 따라 계층화되어 관리됩니다.
    \item \textbf{기술적 정확한 설명:} SIFI 은행들은 국제 은행 규제 당국에 의해 GSIB(Globally Systemically Important Banks)라고도 불립니다. 이 은행들은 전 세계에서 가장 크고 가장 상호 연결된 은행들입니다. GSIB 은행들은 다양한 시스템 리스크 범주로 계층화됩니다.
\end{itemize}

\begin{figure}[h!]
    \centering
    \begin{adjustbox}{width=0.9\textwidth,center}
    \begin{tabular}{ll}
        \toprule
        \textbf{GSIB 계층} & \textbf{대표 은행 (개념적)} \\
        \midrule
        가장 시스템적으로 중요 & JP모건, HSBC (가장 글로벌) \\
        두 번째 계층 & ICBC (자산 기준 세계 최대, 주로 국내), 기타 메가 은행 \\
        G-SIB 제로 버킷 & 글로벌 시스템적으로 중요한 은행 전체 목록 \\
        \bottomrule
    \end{tabular}
    \end{adjustbox}
    \caption{GSIB 계층 분류 (개념적)}
    \label{tab:gsib_tiers}
\end{figure}
\vspace{0.5em}
위 표는 GSIB 은행들이 시스템 리스크에 따라 어떻게 계층화되는지 개념적으로 보여줍니다. JP모건은 가장 시스템적으로 중요하다고 간주되며, HSBC는 가장 글로벌한 은행으로 평가됩니다. ICBC(중국공상은행)는 자산 기준으로 세계 최대 은행이지만, 주로 국내 범위에 집중되어 있어 다른 계층에 분류될 수 있습니다.

\subsubsection{시스템 리스크 규제 접근 방식}

FSOC가 시스템적으로 중요한 기관을 지정한 후, 미국 정부는 시스템 리스크를 목표로 하는 두 가지 접근 방식을 사용합니다.

\begin{enumerate}
    \item \textbf{스트레스 테스트 (Stress Testing):}
    \begin{itemize}
        \item \textbf{내용:} SIFI 기업들은 연방준비제도의 스트레스 테스트를 받습니다. 이는 이들 기업이 금융 위기 시 잠재적 손실에 대해 미리 생각하고, 이러한 손실을 흡수하고 생존할 수 있는 충분한 자본을 보유하도록 장려합니다.
        \item \textbf{목표:} 시스템적으로 중요한 기업의 경우, 이 연습은 더욱 중요합니다.
    \end{itemize}
    \item \textbf{시스템 리스크 추가 자본금 (Systemic Risk Surcharges):}
    \begin{itemize}
        \item \textbf{내용:} 이러한 기관에 '세금'을 부과하는 것입니다. 오염세가 오염을 유발하는 기업이 깨끗해지도록 유인하는 것처럼, 시스템 리스크 추가 자본금은 금융 기업이 '대마불사'가 되는 것을 억제합니다.
        \item \textbf{예시:} 씨티그룹, HSBC, JP모건과 같은 GSIB 은행들은 지구상에서 가장 시스템적으로 위험한 은행으로 간주되며, 그에 따라 세금이 부과됩니다. 동등한 비시스템적 은행과 비교하여, 이들은 투자 1달러당 추가로 2\%포인트의 자기자본을 보유해야 합니다. 사회적 관점에서 우리는 이러한 시스템적 은행들이 잠재적 손실에 대해 훨씬 더 회복 탄력적(Resilient)이기를 원합니다.
    \end{itemize}
\end{enumerate}

\begin{figure}[h!]
    \centering
    \begin{adjustbox}{width=0.9\textwidth,center}
    \begin{tabular}{lc}
        \toprule
        \textbf{GSIB 은행} & \textbf{시스템 리스크 자본 추가금 (개념적)} \\
        \midrule
        JP모건 & 2\% \\
        HSBC & 2\% \\
        씨티그룹 & 2\% \\
        ... & ... \\
        \bottomrule
    \end{tabular}
    \end{adjustbox}
    \caption{GSIB 은행 시스템 리스크 자본 추가금 (개념적)}
    \label{tab:gsib_surcharges}
\end{figure}
\vspace{0.5em}
위 표는 GSIB 은행에 부과되는 시스템 리스크 자본 추가금을 개념적으로 보여줍니다. 가장 시스템적으로 위험한 은행들은 투자 1달러당 추가로 2\%포인트의 자기자본을 보유해야 합니다.

\begin{summarybox}[title=대마불사 요약]
시스템 리스크는 한 은행의 실패가 시스템의 나머지 부분으로 파급될 수 있는 잠재력을 포착합니다. 더 크고 상호 연결된 은행은 가장 큰 시스템 리스크를 초래하며, 모두가 이들이 '대마불사'라고 믿습니다. 대마불사는 은행 규제 당국에게 큰 도전 과제입니다. 시스템 리스크 규제의 기본 아이디어는 이러한 크고 복잡한 금융 기업을 식별하고 세금을 부과하는 것입니다. 미국에서는 FSOC의 책임입니다. 가장 큰 은행들을 강화된 감독에 노출시킴으로써 시스템이 더 안전해지기를 바랍니다.
\end{summarybox}


\subsection{그림자 금융 정의}

은행 시스템 내의 혼란은 사회에 막대한 비용을 초래할 수 있으므로, 정부는 안전망을 제공하고 은행을 규제합니다. 이러한 규제는 부담스럽지만, 은행 시스템이 원활하고 효율적으로 운영되도록 돕는다는 믿음이 있습니다. 이 섹션에서는 그림자 금융(Shadow Banking)에 대해 살펴봅니다.

\subsubsection{그림자 금융이란 무엇인가?}

\begin{itemize}
    \item \textbf{한 줄 핵심 요약:} 그림자 금융은 은행과 유사한 금융 서비스를 제공하지만, 은행만큼 엄격한 규제를 받지 않는 기관들을 말합니다.
    \item \textbf{직관적 예시/비유:} 정식 식당은 위생 검사를 받고 허가를 받아야 하지만, 길거리에서 음식을 파는 노점상은 상대적으로 규제가 덜한 것과 같습니다. 노점상도 음식을 팔지만, 정식 식당만큼의 규제를 받지 않죠. 그림자 금융도 은행과 비슷한 역할을 하지만, 규제 '그림자'에 숨어있다고 해서 붙여진 이름입니다.
    \textbf{기술적 정확한 설명:} 그림자 은행(Shadow Banks)은 은행과 유사한 금융 서비스를 제공하지만, 규제는 덜 받는(Less Regulation) 기관들입니다. 이는 은행 규제 당국에게 엄청난 도전 과제를 제시하며, 특히 규제 당국이 개입할 권한이 제한적일 때 더욱 그렇습니다.
\end{itemize}

\begin{summarybox}[title=상업 은행의 정의 (1956년 은행 지주 회사법)]
상업 은행(Commercial Bank)은 '요구불 예금(Demand Deposits)을 수취하고 상업 대출(Commercial Loans)을 제공하는 기관'으로 정의됩니다. 즉, 이 두 가지 활동을 같은 지붕 아래에서 모두 수행하는 기관입니다.
\end{summarybox}

\subsubsection{그림자 은행의 특징과 예시}

\begin{itemize}
    \item \textbf{한 줄 핵심 요약:} 그림자 은행은 은행처럼 비유동 자산에 투자하고 예금과 유사한 단기 부채를 발행하지만, 법적으로는 은행으로 분류되지 않아 규제를 피합니다.
    \item \textbf{직관적 예시/비유:} 은행은 고객의 돈을 받아(예금) 기업에 빌려줍니다(대출). 그림자 은행도 고객의 돈을 받아(예금과 유사한 상품) 기업이나 부동산에 투자합니다(대출과 유사한 투자). 하지만 법적으로는 '은행'이 아니라고 주장하며 은행 규제를 받지 않습니다.
    \item \textbf{기술적 정확한 설명:} 수년에 걸쳐 상업 은행 서비스를 제공하지만 기술적으로는 은행으로 간주되지 않는 많은 기관들이 등장했습니다.
    \begin{itemize}
        \item \textbf{비유동 자산에 투자:} 일부 기관은 모기지 담보 증권(MBS)이나 기업 대출 뮤추얼 펀드와 같은 비유동 자산에 자금을 투자합니다.
        \item \textbf{예금과 유사한 부채 발행:} 다른 기관들은 요구불 예금과 유사한 부채를 발행하여 운영 자금을 조달합니다. 예를 들어, 환매 조건부 채권(Repurchase Agreement)은 현금이 풍부한 연기금과 같은 기업이 단기간 현금을 보관하고 단기간 내에 회수할 수 있도록 하여 은행 예금과 유사한 경제적 기능을 수행합니다.
        \item \textbf{두 가지 모두 수행:} 일부 기관은 이 두 가지를 모두 수행합니다.
    \end{itemize}
\end{itemize}

\begin{examplebox}[title=대표적인 그림자 은행]
\begin{itemize}
    \item \textbf{머니마켓펀드 (Money Market Mutual Fund, MMMF):} 고정 가치 주식(Fixed Value Shares)을 판매하여 자금을 조달하는 뮤추얼 펀드입니다. 고객은 1달러를 머니마켓 계좌에 넣고 단기간 내에 1달러(이자 포함)를 인출할 수 있습니다. 이 주식은 고객 관점에서 저축 계좌로 판매되며, 은행 예금과 동일한 기능을 합니다. 자산 측면에서는 다양한 채무 상품에 투자합니다. 이는 기본적으로 은행이 하는 일과 같습니다(예금을 발행하고 비유동 자산에 투자). 이러한 이유로 머니마켓펀드는 종종 그림자 은행이라고 불립니다.
    \item \textbf{헤지펀드 (Hedge Funds):} 단기 부채 조달에 의존하는 일부 헤지펀드도 그림자 은행으로 간주됩니다.
    \item \textbf{그림자 은행이 아닌 경우:} 보험 회사와 같은 비은행 기관은 장기 비유동 자산에 투자하지만, 유동성 부채로 자금을 조달하지 않으므로 그림자 은행으로 간주되지 않습니다.
\end{itemize}
\end{examplebox}

\subsubsection{그림자 은행의 존재 이유: 효율성 vs. 규제 차익거래}

\begin{itemize}
    \item \textbf{한 줄 핵심 요약:} 그림자 은행은 은행보다 효율적일 수 있지만, 주로 은행 규제를 회피하여 '규제 차익거래'를 하려는 목적으로 존재한다는 비판을 받습니다.
    \item \textbf{기술적 정확한 설명:} 그림자 은행이 존재하는 이유에 대해 두 가지 관점이 있습니다.
    \begin{itemize}
        \item \textbf{효율성 관점:} 그림자 은행은 은행보다 더 효율적인 전문 금융 기업일 수 있습니다.
        \item \textbf{규제 차익거래 관점 (더 일반적):} 이들 기관은 은행과 동일한 서비스를 제공하지만, 규제를 회피합니다. 그래서 '그림자'라는 이름이 붙었습니다. 이들 기관은 규제의 그림자 속에 숨어 있습니다.
    \end{itemize}
\end{itemize}

\subsubsection{그림자 은행의 문제점: 유동성 위험과 안전망 부재}

\begin{itemize}
    \item \textbf{한 줄 핵심 요약:} 그림자 은행은 은행과 유사한 유동성 불일치 구조를 가지지만, 정부 안전망의 보호를 받지 못해 뱅크런에 매우 취약합니다.
    \item \textbf{직관적 예시/비유:} 안전벨트나 에어백 없이 고속도로를 달리는 자동차와 같습니다. 사고가 나면 훨씬 더 위험하겠죠. 그림자 은행은 은행처럼 위험한 '유동성 불일치'라는 엔진을 가지고 있지만, '예금 보험'이나 '최종 대부자'라는 안전벨트와 에어백이 없습니다.
    \item \textbf{기술적 정확한 설명:} 그림자 은행은 유동성 위험(Liquidity Risk)에 노출되어 있습니다. 그림자 은행이 예금을 발행하도록 허용한다면, 예금자들의 뱅크런에 취약해질 수밖에 없습니다. 이는 상업 은행의 경우 정부 안전망 도입으로 1세기 전에 해결된 문제였습니다.
    \item \textbf{정부 안전망 부재:} 그림자 은행은 정부 안전망에 접근할 수 없기 때문에 취약합니다. 그림자 예금은 FDIC에 의해 보험에 가입되지 않습니다. 이는 그림자 예금자들이 우려할 경우 그림자 은행에 대해 뱅크런을 할 유인을 가집니다. 또한, 그림자 은행은 금융 위기 시 최종 대부자로부터 차입할 수 없습니다. 이는 뱅크런이 발생할 경우 건전한 그림자 은행이라도 현금이 고갈되어 파산할 수 있음을 의미합니다.
    \item \textbf{규제와 구제금융의 역설:} 이들 기관은 정부 규제가 제한적이므로 위기 시 구제금융을 받아서는 안 됩니다. 그러나 2008년 머니마켓펀드 사례에서 보듯이, 실제로는 구제금융을 받았습니다.
\end{itemize}


\subsection{사례 연구: 머니마켓펀드}

머니마켓펀드(Money Market Mutual Funds, MMMF)는 그림자 금융의 중요한 사례입니다.

\subsubsection{2008년 이전 상황}

\begin{itemize}
    \item \textbf{한 줄 핵심 요약:} 2008년 이전 머니마켓펀드는 예금 보험이나 자본 요건 없이 FDIC 보험 예금과 유사한 저축 계좌를 제공했습니다.
    \item \textbf{기술적 정확한 설명:} 2008년 이전 머니마켓펀드는 FDIC 보험 예금과 유사한 저축 계좌를 제공했습니다. 머니마켓펀드는 기본적으로 예금 보험이나 자본 요건 없이 은행 예금을 제공했습니다.
\end{itemize}

\subsubsection{2008년 공황과 정부 구제금융}

\begin{itemize}
    \item \textbf{한 줄 핵심 요약:} 2008년 리먼 브라더스 파산 후 머니마켓펀드 산업은 대규모 인출 공황을 겪었고, 정부는 '대마불사'로 간주하여 구제금융을 제공했습니다.
    \item \textbf{기술적 정확한 설명:} 규제가 덜한 머니마켓펀드 산업은 2008년 9월 대규모 인출(Massive Withdrawals)을 겪었습니다. 일부 머니마켓펀드는 리먼 브라더스 및 기타 안전하다고 여겨지던 자산에 손실을 입는 투자를 했었습니다. 이러한 손실이 현실화되자 산업 전반에 걸쳐 공황이 촉발되었습니다.
\end{itemize}

\begin{figure}[h!]
    \centering
    \begin{adjustbox}{width=0.9\textwidth,center}
    \begin{tabular}{lc}
        \toprule
        \textbf{날짜} & \textbf{관리 자산 (2008년 9월 10일 기준 100\%)} \\
        \midrule
        2008년 9월 10일 & 100\% \\
        2008년 9월 말 & 약 70\% (30\% 감소) \\
        이후 & 안정화 \\
        \bottomrule
    \end{tabular}
    \end{adjustbox}
    \caption{머니마켓펀드 뱅크런 (2008년 9월 개념적)}
    \label{tab:mmmf_run_2008}
\end{figure}
\vspace{0.5em}
위 표는 2008년 9월 머니마켓펀드 뱅크런 상황을 개념적으로 보여줍니다. 리먼 브라더스 파산 약 일주일 전인 9월 10일을 기준으로 관리 자산을 100\%로 정규화했을 때, 자산이 약 30\% 급감한 후 안정화되었습니다. 이는 명백한 공황 상황이었습니다.

\begin{itemize}
    \item \textbf{공황 진정 이유:} 공황은 정부 구제금융(Government Bailout) 덕분에 진정되었습니다. 머니마켓펀드는 그림자 은행이며, 정부 보험 예금을 발행하지 않고 위기 시 연준으로부터 차입할 수 없어야 합니다. 그러나 2008년 실제 상황에서는 이들이 그렇게 할 수 있도록 허용되었습니다. 전체 산업이 '대마불사'로 간주된 것입니다.
    \item \textbf{구체적인 구제금융 조치:}
    \begin{itemize}
        \item \textbf{2008년 9월 19일:} 미국 정부는 머니마켓펀드의 모든 부채를 보증했습니다. 투자자들은 단 한 푼도 잃을 위험이 없게 되었습니다.
        \item \textbf{연준의 긴급 대출:} 연방준비제도는 비상 권한을 사용하여 머니마켓펀드에 현금 대출을 제공했습니다. 이로써 이들 뮤추얼 펀드는 투자자들의 현금 수요를 쉽게 충족시킬 수 있었습니다.
    \end{itemize}
\end{itemize}

\begin{figure}[h!]
    \centering
    \begin{adjustbox}{width=0.9\textwidth,center}
    \begin{tabular}{lc}
        \toprule
        \textbf{기간} & \textbf{연준의 머니마켓펀드 대출 프로그램 (개념적)} \\
        \midrule
        2008년 9월 말 & 급격한 대출 증가 \\
        이후 & 점진적 감소 및 종료 \\
        \bottomrule
    \end{tabular}
    \end{adjustbox}
    \caption{연준의 머니마켓펀드 대출 프로그램 (2008년 개념적)}
    \label{tab:fed_mmmf_lending}
\end{figure}
\vspace{0.5em}
위 표는 2008년 연준의 머니마켓펀드 대출 프로그램의 개념적인 추이를 보여줍니다. 공황 시기에 대출이 급격히 증가했다가 상황이 진정되면서 점진적으로 감소하고 종료되었습니다. 전반적으로 이들 그림자 은행은 보험에 가입되고 규제되는 상업 은행과 거의 동일한 정부 안전망에 접근했습니다.

\subsubsection{미국 정부의 대응과 새로운 규칙}

\begin{itemize}
    \item \textbf{한 줄 핵심 요약:} 2008년 위기 이후 FSOC는 머니마켓펀드에 대한 새로운 규제를 추진했지만, 2020년 팬데믹 시기에도 또 다른 공황과 구제금융이 발생했습니다.
    \item \textbf{기술적 정확한 설명:} 미국 정부는 이 불행한 사건에 어떻게 대응했을까요? 시스템 리스크 규제 기관인 금융 안정성 감독 위원회(FSOC)가 나섰습니다. FSOC는 재무부, 연준, SEC(뮤추얼 펀드 산업 규제 기관)가 해결책을 찾기 위한 포럼이었습니다.
    \item \textbf{FSOC의 판단:} 개별 머니마켓펀드는 시스템적으로 중요하지 않지만, 전체 산업은 뱅크런에 취약한 것으로 보였습니다. FSOC는 SEC에 금융 안정성을 보호하도록 압력을 가했습니다.
    \item \textbf{새로운 규칙 도입:} SEC는 뱅크런에 취약한 대형 펀드에 대해 '변동 순자산 가치(Floating Net Asset Value)' 규칙을 도입했습니다. 이는 머니마켓펀드 주식을 예금보다는 뮤추얼 펀드 주식처럼 만들었습니다. 또한, 이 새로운 규칙은 머니마켓펀드에 뱅크런을 방지할 도구를 제공했습니다. 뱅크런 발생 시 펀드는 인출에 대해 유동성 수수료(Liquidity Fees)를 부과하거나, 최대 10일 동안 환매(Redemptions)를 동결할 수 있게 되었습니다.
\end{itemize}

\subsubsection{2020년 팬데믹 공황과 규칙의 한계}

\begin{itemize}
    \item \textbf{한 줄 핵심 요약:} 2020년 코로나19 팬데믹 시기에도 머니마켓펀드는 다시 공황을 겪었고, 새로운 규칙은 효과적으로 작동하지 않아 또다시 연준의 구제금융이 필요했습니다.
    \item \textbf{기술적 정확한 설명:} 이 새로운 규칙들이 2020년 뱅크런을 막는 데 도움이 되었을까요? 불행히도 그렇지 않았습니다.
\end{itemize}

\begin{figure}[h!]
    \centering
    \begin{adjustbox}{width=0.9\textwidth,center}
    \begin{tabular}{lc}
        \toprule
        \textbf{날짜} & \textbf{관리 자산 (2020년 3월 9일 기준 100\%)} \\
        \midrule
        2020년 3월 9일 & 100\% \\
        2020년 3월 중순 & 급격한 감소 \\
        이후 & 연준 개입 후 안정화 \\
        \bottomrule
    \end{tabular}
    \end{adjustbox}
    \caption{머니마켓펀드 뱅크런 (2020년 3월 개념적)}
    \label{tab:mmmf_run_2020}
\end{figure}
\vspace{0.5em}
위 표는 2020년 3월 머니마켓펀드 뱅크런 상황을 개념적으로 보여줍니다. 2020년 3월 9일(미국 내 코로나19 충격 시작 시점)을 기준으로 관리 자산을 100\%로 정규화했을 때, 다시 공황이 발생했습니다. 초기 뱅크런 발생 10일 후, 연준은 다시 개입하여 현금 대출로 산업을 구제했습니다.

\begin{itemize}
    \item \textbf{규칙이 작동하지 않은 이유:} 이는 여전히 논쟁의 대상입니다. 한 가지 문제는 인출에 대한 수수료 부과가 투자자들에게 선제적으로 인출할 유인을 제공할 수 있다는 것입니다. '지금 인출하지 않으면 내일 자금이 묶일 수도 있다'는 생각은 오히려 뱅크런을 부추길 수 있습니다. 이는 머니마켓펀드 관리자에게 인출을 차단할 도구를 주는 것이 역효과를 낼 수 있음을 시사합니다. 분명히 미국 정책 입안자들은 또 다른 공황과 구제금융을 막을 수 없었습니다.
\end{itemize}

\begin{summarybox}[title=그림자 금융 요약]
우리는 상업 은행과 그림자 은행의 차이점을 살펴보았습니다. 그림자 은행은 상업 은행과 유사한 금융 서비스를 제공하지만, 동일한 방식으로 규제되지 않습니다. 일반적으로 정부 규제를 덜 받습니다. 그림자 은행은 정부 안전망에 접근할 수 없어야 하지만, 머니마켓펀드 사례에서 보듯이 정부는 이들을 구제할 강력한 유인을 가질 수 있습니다. 그림자 은행을 식별하고 규제하는 것은 금융 규제 당국에게 여전히 엄청난 도전 과제입니다.
\end{summarybox}


\section{체크리스트: 모듈 4 학습 점검}

이 체크리스트를 활용하여 모듈 4의 핵심 내용을 얼마나 잘 이해하고 있는지 스스로 점검해 보세요.

\begin{itemize}
    \item \textbf{규제의 기본 이해:}
    \begin{itemize}
        \item 시장 실패와 부정적 외부 효과가 무엇이며, 왜 규제가 필요한지 설명할 수 있는가?
        \item 미시 건전성 규제와 거시 건전성 규제의 차이점을 아는가?
        \item 고객 보호와 공정성이라는 규제의 다른 목표를 설명할 수 있는가?
        \item 규제의 주요 상충 관계(비효율성, 의도치 않은 결과, 도덕적 해이)를 이해하는가?
    \end{itemize}
    \item \textbf{정부 안전망:}
    \begin{itemize}
        \item 은행의 취약성(유동성 불일치)과 뱅크런의 원인을 설명할 수 있는가?
        \item 전염 효과의 두 가지 주요 원천(정보 비대칭, 상호 연결성)을 아는가?
        \item 예금 보험(FDIC)과 최종 대부자(연준)의 역할과 작동 방식을 설명할 수 있는가?
        \item 2007-2009년 금융 위기 시 정부 개입 사례를 설명할 수 있는가?
    \end{itemize}
    \item \textbf{규제 및 감독 실제:}
    \begin{itemize}
        \item 은행 규제와 은행 감독의 차이점을 명확히 구분할 수 있는가?
        \item 은행 규제의 주요 영역(경쟁, 활동/투자, 자금 조달, 공시)을 설명할 수 있는가?
        \item 바젤 협약의 목적(글로벌 표준, 자본 요건)과 진화 과정을 아는가?
        \item 감독 과정에서 집행 조치(벌금, 제한)가 어떻게 이루어지는지 아는가?
        \item CAMELS 평가 시스템의 각 요소가 무엇을 의미하는지 설명할 수 있는가?
        \item 스트레스 테스트의 목표, 절차, 그리고 SCAP/CCAR 사례를 설명할 수 있는가?
    \end{itemize}
    \item \textbf{21세기 규제 당국의 과제:}
    \begin{itemize}
        \item 미국 은행 시스템의 집중화 현상과 '대마불사' 개념을 이해하는가?
        \item '대마불사'가 은행의 효율성 관점과 정부 보조금 관점에서 어떻게 설명되는지 아는가?
        \item 시스템 리스크의 개념과 FSOC의 역할(SIFI 지정)을 설명할 수 있는가?
        \item GSIB 분류와 시스템 리스크 추가 자본금의 의미를 아는가?
        \item 그림자 금융의 정의와 상업 은행과의 차이점을 설명할 수 있는가?
        \item 머니마켓펀드 사례를 통해 그림자 금융의 유동성 위험과 정부 안전망 부재 문제를 설명할 수 있는가?
        \item 2008년 및 2020년 머니마켓펀드 공황과 정부 개입, 그리고 새로운 규칙의 한계를 아는가?
    \end{itemize}
\end{itemize}


\section{FAQ: 자주 묻는 질문}

초심자들이 금융 규제에 대해 가질 수 있는 몇 가지 질문과 답변입니다.

\begin{examplebox}[title=Q1: 모든 금융 기관이 은행처럼 규제되어야 하나요?]
\textbf{A1:} 반드시 그렇지는 않습니다. 은행은 예금 수취와 대출이라는 독특한 기능 때문에 시스템적으로 중요하며, 뱅크런에 취약합니다. 따라서 은행에 대한 엄격한 규제는 금융 시스템 안정에 필수적입니다. 하지만 모든 금융 기관이 은행과 동일한 위험 프로필을 가지는 것은 아니므로, 각 기관의 특성과 위험에 맞는 맞춤형 규제가 필요합니다. 그림자 금융의 경우, 은행과 유사한 위험을 가지지만 규제가 덜한 것이 문제입니다.
\end{examplebox}

\begin{examplebox}[title=Q2: 규제가 은행의 혁신을 저해하지는 않나요?]
\textbf{A2:} 규제는 은행의 특정 활동에 제약을 가하므로 단기적으로 혁신을 저해할 수 있다는 비판이 있습니다. 그러나 장기적으로 볼 때, 안정적인 금융 시스템은 혁신과 경제 성장의 기반이 됩니다. 과도한 위험 감수로 인한 금융 위기는 훨씬 더 큰 경제적 손실을 초래하므로, 규제는 위험과 혁신 사이의 균형을 찾는 중요한 역할을 합니다.
\end{examplebox}

\begin{examplebox}[title=Q3: '대마불사' 문제는 어떻게 해결할 수 있나요?]
\textbf{A3:} '대마불사'는 규제 당국의 가장 큰 도전 과제 중 하나입니다. 해결책으로는 대형 은행에 더 높은 자본 요건을 부과하거나(시스템 리스크 추가 자본금), 복잡성을 줄여 실패 시 해체하기 쉽게 만드는 방안, 그리고 구제금융 대신 '정리 계획(Resolution Plan)'을 통해 시장에 미치는 충격을 최소화하면서 파산시키는 방안 등이 논의되고 있습니다.
\end{examplebox}

\begin{examplebox}[title=Q4: 스트레스 테스트는 항상 정확한가요?]
\textbf{A4:} 스트레스 테스트는 미래의 불확실한 상황을 예측하는 것이므로 완벽하게 정확할 수는 없습니다. 시나리오 설정의 가정, 은행의 모델링 능력, 그리고 예상치 못한 새로운 유형의 위기 발생 가능성 등 여러 한계가 있습니다. 하지만 과거 위기 경험을 바탕으로 은행의 취약점을 미리 파악하고 대비하는 데 매우 유용한 도구임은 분명합니다.
\end{examplebox}

\begin{examplebox}[title=Q5: 그림자 금융은 왜 규제하기 어려운가요?]
\textbf{A5:} 그림자 금융은 빠르게 진화하고 다양한 형태로 존재하며, 법적 정의가 모호한 경우가 많습니다. 또한, 기존 은행 규제 체계의 사각지대에 놓여 있어 규제 당국의 권한이 제한적입니다. 규제가 너무 엄격하면 금융 활동이 다른 규제 회피 경로로 이동할 수 있어, 규제 당국은 그림자 금융의 위험을 관리하면서도 시장의 효율성을 해치지 않는 균형점을 찾아야 합니다.
\end{examplebox}


\section{빠르게 훑어보기: 모듈 4 핵심 요약 (1페이지)}

\begin{summarybox}[title=모듈 4: 규제 핵심 요약]

\textbf{1. 규제의 목표: 왜 필요한가?}
\begin{itemize}
    \item \textbf{시장 실패 교정:} 은행 실패의 \textbf{부정적 외부 효과} (전염성) 방지.
    \item \textbf{금융 안정성:} \textbf{미시 건전성} (개별 은행) 및 \textbf{거시 건전성} (시스템 전체) 유지.
    \item \textbf{고객 보호:} 소액 예금자 및 차입자 보호, \textbf{레드라이닝} 같은 차별 방지.
    \item \textbf{상충 관계:} \textbf{비효율성}, \textbf{의도치 않은 결과}, \textbf{도덕적 해이} (예: TARP) 발생 가능성.
\end{itemize}

\textbf{2. 정부 안전망: 뱅크런 방지}
\begin{itemize}
    \item \textbf{은행 취약성:} \textbf{유동성 불일치} (단기 차입, 장기 대출)로 인한 \textbf{뱅크런} 발생.
    \item \textbf{전염 효과:} 정보 비대칭, 상호 연결성으로 공황 확산.
    \item \textbf{두 가지 축:}
    \begin{itemize}
        \item \textbf{예금 보험 (FDIC):} 예금자 보호, 뱅크런 \textbf{사전 예방}.
        \item \textbf{최종 대부자 (연준):} 건전 은행에 긴급 대출, 뱅크런 \textbf{사후 중단}.
    \end{itemize}
    \item \textbf{2007-2009 위기:} 예금 보험 확대, 연준 유동성 지원으로 공황 대응.
\end{itemize}

\textbf{3. 규제 및 감독 실제: 어떻게 작동하는가?}
\begin{itemize}
    \item \textbf{규제 vs. 감독:} \textbf{규제}는 규칙 제정, \textbf{감독}은 규칙 준수 감시 및 집행.
    \item \textbf{규제 대상:} 경쟁, 활동/투자, 자금 조달, 공시 요건.
    \item \textbf{바젤 협약:} \textbf{글로벌 최소 자본 요건} (위험 기반) 설정, 공정한 경쟁 환경 조성.
    \item \textbf{감독 과정:}
    \begin{itemize}
        \item \textbf{집행 조치:} 규칙 위반 시 벌금, 활동 제한 (예: 런던 고래).
        \item \textbf{은행 검사:} 비현장 (보고서 분석) 및 현장 (방문) 검사.
        \item \textbf{CAMELS 평가:} 자본, 자산, 경영, 수익, 유동성, 시장 위험 민감도 평가.
    \end{itemize}
    \item \textbf{스트레스 테스트:}
    \begin{itemize}
        \item \textbf{목표:} 극심한 위기 시 은행의 자본 적정성 평가 (미래 지향적).
        \item \textbf{절차:} 경제 시나리오 $\rightarrow$ 손실 예측/자본 계획 $\rightarrow$ 규제 검토 $\rightarrow$ 결과 공개.
        \item \textbf{사례:} SCAP (2009년 위기 대응), CCAR (매년 정기 실시).
    \end{itemize}
\end{itemize}

\textbf{4. 21세기 규제 당국의 과제: 새로운 도전}
\begin{itemize}
    \item \textbf{대마불사 (TBTF):}
    \begin{itemize}
        \item \textbf{문제:} 소수 메가 은행의 집중화, 실패 시 \textbf{시스템 리스크} (전체 경제 마비).
        \item \textbf{원인:} 효율성 (규모/범위의 경제) vs. 정부 보조금 (구제금융 기대).
        \textbf{대응:} \textbf{FSOC} (시스템 리스크 관리), \textbf{SIFI/GSIB} 지정, \textbf{시스템 리스크 추가 자본금} 부과.
    \end{itemize}
    \item \textbf{그림자 금융 (Shadow Banking):}
    \begin{itemize}
        \item \textbf{정의:} 은행과 유사 서비스 제공하나 규제 덜 받는 기관 (예: 머니마켓펀드).
        \item \textbf{문제:} \textbf{유동성 위험}에 취약하나 정부 안전망 부재.
        \item \textbf{사례 (MMMF):} 2008년 및 2020년 공황 시 정부 구제금융 발생, 규제 한계 노출.
    \end{itemize}
\end{itemize}
\end{summarybox}


%=======================================================================
% Module 5: Lecture 5
%=======================================================================
\chapter{Lecture 5}
\label{ch:module5}

\metainfo{금융 시장과 통화 정책}{Module 1: 금리 및 증권화}{Prof. Ralf Meisenzahl}{중앙은행과 통화정책의 핵심 개념, 금리의 종류와 작동 방식, 주요 금융 시장 및 금융 상품에 대한 심층적인 이해를 목표로 함.}


\section{개요}
이 문서는 중앙은행과 통화정책의 핵심 개념을 다루는 모듈 1의 내용을 정리한 것입니다.
통화정책이 어떻게 실행되고 거시 경제에 영향을 미치는지 이해하기 위한 기초 지식을 제공합니다.
특히, 금리의 종류와 작동 방식, 주요 금융 시장 및 금융 상품에 대한 심층적인 이해를 목표로 합니다.
단기 및 장기 금리, 국채, 환매조건부채권(Repo), 머니 마켓 뮤추얼 펀드 등 다양한 시장과 도구를 학습하여,
통화정책의 복잡한 구현 과정을 파악하고 경제 예측에 활용되는 지표들을 분석할 수 있습니다.
이 모듈은 다음 모듈에서 다룰 통화정책의 구체적인 도구와 구현 방식을 이해하기 위한 필수적인 배경 지식을 제공합니다.


\section{용어 정리}
다음 표는 본 문서에서 다루는 주요 용어들을 쉽게 이해할 수 있도록 정리한 것입니다.

\begin{table}[h!]
\centering
\caption{주요 금융 용어 정리}
\label{tab:terms}
\begin{adjustbox}{width=\textwidth,center}
\begin{tabular}{lllc}
\toprule
\textbf{용어 (원어)} & \textbf{쉬운 설명} & \textbf{비고} \\
\midrule
연방기금 시장 (Federal Funds Market) & 은행들이 연방준비제도에 예치된 초과 지급준비금을 서로 빌려주고 빌리는 시장. & 주로 익일물(overnight) 대출. \\
연방기금 금리 (Federal Funds Rate) & 연방기금 시장에서 은행 간에 초과 지급준비금을 거래할 때 적용되는 금리. & 중앙은행의 주요 정책 목표 금리. \\
환매조건부채권 (Repo) & 채권을 담보로 현금을 빌리고, 다음 날 채권을 다시 사들이기로 약속하는 거래. & 담보가 있는 익일물 대출 시장. \\
리보 (LIBOR) & 런던 은행 간 대출 금리. 은행 간 무담보 대출에 적용되는 금리. & 과거 조작 논란으로 SOFR로 대체 중. \\
소퍼 (SOFR) & 담보부 익일물 자금 조달 금리. 국채를 담보로 하는 환매조건부채권 시장 금리. & LIBOR를 대체하는 새로운 무위험 기준 금리. \\
명목 금리 (Nominal Interest Rate) & 인플레이션을 고려하지 않은, 표면적인 금리. & 은행 예금 이자율 등. \\
실질 금리 (Real Interest Rate) & 인플레이션을 고려하여 구매력 변화를 반영한 금리. & 실제 투자 수익률을 나타냄. \\
피셔 방정식 (Fisher Equation) & 명목 금리, 실질 금리, 예상 인플레이션 간의 관계를 설명하는 방정식. & 명목 금리 $\approx$ 실질 금리 + 예상 인플레이션. \\
물가연동국채 (TIPS) & 원금 가치가 인플레이션에 따라 조정되는 국채. & 실질 금리를 지급하며 인플레이션 헤지 기능. \\
국채 (Treasury Securities) & 미국 재무부가 발행하는 채권. 신용 위험이 없어 무위험 자산으로 간주. & T-Bill, T-Note, T-Bond로 구분. \\
무위험 금리 (Risk-Free Rate) & 신용 위험이 없는 투자에 대한 이론적인 수익률. & 보통 미국 국채 금리를 기준으로 삼음. \\
금리 기간 구조 (Term Structure of Interest Rates) & 동일 발행 주체의 만기별 금리 비교. & 수익률 곡선(Yield Curve)으로 시각화. \\
수익률 곡선 (Yield Curve) & 만기가 다른 채권들의 수익률을 연결한 곡선. & 경제 상황 및 미래 금리 기대를 반영. \\
머니 마켓 뮤추얼 펀드 (Money Market Mutual Funds) & 초단기, 초안전 자산에 투자하는 뮤추얼 펀드. & 은행 예금의 대안으로 활용. \\
지급준비금 (Reserves) & 은행이 예금 인출 등에 대비하여 중앙은행에 예치하는 현금. & 법정 지급준비금과 초과 지급준비금으로 나뉠 수 있음. \\
법정 지급준비금 (Required Reserves) & 법으로 정해진 최소한의 지급준비금. & 중앙은행이 결정. \\
초과 지급준비금 (Excess Reserves) & 법정 지급준비금을 초과하여 보유하는 지급준비금. & 은행의 유동성 수요에 대비. \\
헤어컷 (Haircut) & 담보 가치보다 적은 금액을 대출해주는 비율. & 담보 가치 하락 위험을 줄이기 위함. \\
\bottomrule
\end{tabular}
\end{adjustbox}
\end{table}


\section{모듈 1: 금리 (Interest Rates)}

\subsection{개요}
이 모듈은 중앙은행, 통화정책, 그리고 거시 경제에 대한 이해를 돕기 위한 첫걸음입니다. 우리는 통화정책이 어떻게 실행되고, 중앙은행의 결정이 궁극적으로 거시 경제에 어떻게 전달되는지 논의할 것입니다. 이 과정을 통해 통화정책의 목표에 대한 깊이 있는 이해를 얻게 될 것입니다.

\begin{infobox}[title=\textbf{통화정책의 핵심 질문}]
중앙은행이 "금리를 25bp 인상한다"거나 "장기 금리에 하방 압력을 가하기 위해 자산 매입 프로그램을 확대한다"는 발표를 들었을 때, 이것이 실제로 무엇을 의미할까요? 중앙은행은 이러한 정책을 어떻게 실행하며, 어떤 금리가 왜 영향을 받을까요?
\end{infobox}

통화정책의 실행은 여러 금융 시장과 금융 상품을 포함하는 복잡한 주제입니다. 이 모듈은 통화정책이 어떤 금리를 목표로 하는지 정확히 설명하는 데 필요한 배경 지식을 제공합니다.

\begin{warningbox}[title=\textbf{주의: 모듈 1의 특성}]
이 모듈의 영상들은 다음 모듈에서 통화정책을 본격적으로 다루기 때문에 다소 단절된 느낌을 줄 수 있습니다. 각 영상이 전체 그림에 어떻게 들어맞는지 궁금할 수 있습니다. 하지만 이 영상들은 통화정책과 그 실행을 이해하기 위한 배경 지식을 제공하는 것임을 기억하세요. 나중에 특정 통화정책 도구나 금융 위기를 논의할 때 언제든지 이 영상들로 돌아와 개념을 다시 확인할 수 있습니다.
\end{warningbox}

\subsubsection*{모듈 1 영상 개요}
\begin{itemize}
    \item \textbf{단기 자금 시장:} 첫 두 영상은 연방기금 시장(Federal Funds Market)과 환매조건부채권(Repo) 시장이라는 두 가지 단기 자금 시장을 소개합니다. 이 시장들은 금융 기관들이 현금을 빌리고 빌려주는 익일물(overnight) 시장이며, 통화정책의 주요 목표 대상입니다.
    \item \textbf{연방준비제도의 목표:} 연방준비제도(Federal Reserve)는 연방기금 시장의 금리인 연방기금 금리(Federal Funds Rate)를 목표 범위로 설정하여 관리합니다. 환매조건부채권 시장은 연방기금 시장의 대안이므로, 연방준비제도는 모든 익일물 금리가 통화정책 목표와 일치하도록 환매조건부채권 시장의 금리도 목표로 합니다.
    \item \textbf{런던 은행 간 대출 금리 (LIBOR):} 다음 영상에서는 또 다른 익일물 금리인 런던 은행 간 대출 금리(London Interbank Offered Rate, LIBOR)를 환매조건부채권 시장 금리와 비교합니다. 이 금리들은 기업 대출의 벤치마크가 됩니다.
    \item \textbf{주요 금융 상품:} 통화정책의 실행과 전달에 핵심적인 세 가지 금융 상품을 살펴봅니다.
    \begin{enumerate}
        \item \textbf{정부 증권 (Government Securities):} 중앙은행은 통화정책을 실행하기 위해 정부 증권을 사고팔기 때문에, 이 증권이 무엇이고 어떻게 작동하는지 이해하는 것이 중요합니다. 미국 정부 증권은 채무 불이행 위험이 없으므로, 금리의 구성 요소(예: 인플레이션 반영분)를 분석할 수 있습니다.
        \item \textbf{금리 기간 구조 (Term Structure of Interest Rates):} 금리가 인플레이션과 미래에 대한 기대를 반영하므로, 만기가 짧은 정부 증권과 만기가 긴 정부 증권의 금리를 비교하여 시장의 미래 경제 발전 기대를 이해할 수 있습니다. 이는 중앙은행, 투자 은행, 금융 분석가들이 면밀히 주시하는 지표입니다.
        \item \textbf{머니 마켓 뮤추얼 펀드 (Money Market Mutual Funds):} 머니 펀드는 환매조건부채권 시장에서 가장 큰 익일물 현금 대출 기관입니다. 2008년 금융 위기 이후 연방준비제도는 이러한 펀드를 통해 통화정책을 수행하기 시작했습니다.
        \item \textbf{증권화 (Securitization):} 증권화는 신용 공급을 증가시켰습니다. 특히 주택담보부증권(Mortgage-Backed Securities, MBS)과 같은 증권화된 모기지 풀은 매우 중요합니다. 통화정책은 신용 공급을 줄이거나 늘림으로써 경제에 영향을 미치고, 연방준비제도는 특정 MBS를 매입할 수 있으므로, 증권화 시장은 통화정책 고려 사항에서 중요한 요소입니다.
    \end{enumerate}
\end{itemize}
요약하자면, 이 첫 번째 모듈에서 논의되는 모든 시장과 상품은 통화정책 및 그 실행과 직접적으로 관련되어 있습니다. 각 영상은 다음 모듈에서 더 명확해질 큰 그림의 작은 조각들을 제공합니다.


\section{단기 금리 (Short-term Rates)}

\subsection{연방기금 시장 (The Market for Federal Funds)}
이 강의에서는 연방기금 시장의 배경과 작동 방식을 논의합니다.

\subsubsection*{은행이 지급준비금을 보유하는 이유}
\begin{itemize}
    \item \textbf{은행의 사업 모델:} 은행은 단기로 자금을 빌려(예금 유치) 장기로 대출(대출 실행)하는 사업 모델을 가지고 있습니다.
    \item \textbf{현금 수요 (유동성):} 이러한 사업 모델은 은행에 현금, 즉 유동성(liquidity)에 대한 수요를 발생시킵니다. 예를 들어, 예금은 언제든지 인출될 수 있으므로, 은행은 유동성 수요를 충족시키기 위해 일정량의 현금을 보유해야 합니다. 이 현금은 연방준비제도(Federal Reserve)의 지급준비금 계좌에 예치될 수 있습니다.
\end{itemize}

\subsubsection*{지급준비금의 종류: 법정 지급준비금과 초과 지급준비금}
은행은 얼마나 많은 지급준비금을 보유해야 할까요? 이는 어려운 질문입니다.

\begin{itemize}
    \item \textbf{법정 지급준비금 (Required Reserves):}
    \begin{itemize}
        \item 미국에서는 은행이 연방준비제도에 최소한의 지급준비금을 보유하도록 법적으로 요구됩니다. 이를 법정 지급준비금이라고 합니다.
        \item 이 지급준비금 요건은 은행이 일상적인 현금 수요를 충족시킬 수 있도록 설정되며, 연방준비제도가 결정합니다.
        \item \textbf{예시:} 코로나19 팬데믹 이전에는 예금의 3\%였습니다. 팬데믹 동안 은행을 지원하기 위해 이 요건은 0\%로 설정되었습니다.
    \end{itemize}
    \item \textbf{초과 지급준비금 (Excess Reserves):}
    \begin{itemize}
        \item 은행은 일반적으로 법정 지급준비금보다 더 많은 지급준비금을 보유합니다. 이를 초과 지급준비금이라고 합니다.
        \item 은행은 예상치 못한 유동성 수요를 충족시키기 위해 초과 지급준비금을 보유합니다.
        \item \textbf{예시:} 코로나19 팬데믹 초기에 미국 기업들은 신용 한도에서 약 6천억 달러를 인출했습니다. 은행들은 이 예측 불가능한 사건에 대비할 수 없었지만, 단기간에 고객에게 6천억 달러의 현금을 제공해야 했습니다. 이것이 초과 지급준비금의 필요성을 보여줍니다.
    \end{itemize}
\end{itemize}

\subsubsection*{은행의 지급준비금 수준 선택: 유동성과 기회비용}
은행은 지급준비금 수준을 어떻게 선택할까요?

\begin{itemize}
    \item 은행은 현금을 지급준비금 형태로 보유하거나, 이 현금을 투자할 수 있습니다. 예를 들어, 은행에 이자 수입을 안겨주는 대출에 투자할 수 있습니다.
    \item \textbf{상충 관계 (Trade-offs):}
    \begin{itemize}
        \item 지급준비금을 보유하는 것은 일반적으로 대출에 비해 낮은 수입을 창출합니다.
        \item 그러나 예상치 못한 유동성 수요에 직면하면, 은행은 막판에 자금을 빌려야 할 수 있으며, 이는 비용이 많이 들 수 있습니다.
        \item 따라서 은행은 유동성 수요와 이자 지급 자산에 투자했을 때 얻을 수 있었던 수익(기회비용) 사이에서 상충 관계에 직면합니다.
    \end{itemize}
\end{itemize}

\begin{examplebox}[title=\textbf{초과 지급준비금의 기회비용 예시}]
은행이 100억 원의 초과 지급준비금을 보유하고 있다고 가정해 봅시다. 이 돈을 대출해 주면 연 5\%의 이자를 받을 수 있지만, 지급준비금으로 보유하면 연 0.1\%의 이자만 받습니다. 이 경우, 대출을 통해 얻을 수 있었던 4.9\%의 이자 수익이 초과 지급준비금 보유의 기회비용이 됩니다. 은행은 이 기회비용과 예상치 못한 유동성 부족으로 인한 위험 및 비용을 저울질하여 최적의 지급준비금 수준을 결정합니다.
\end{examplebox}

\subsubsection*{미국 상업은행의 지급준비금 보유 현황}
\begin{itemize}
    \item \textbf{역사적으로 낮은 수준:} 1960년대부터 2007년까지 지급준비금은 거의 증가하지 않았습니다.
    \item \textbf{2008년 금융 위기 이후 급증:}
    \begin{itemize}
        \item 2007년 약 400억 달러에서 2014년 거의 3조 달러로, 2021년 3월에는 3.7조 달러로 크게 증가했습니다.
        \item \textbf{급증의 원인:}
        \begin{enumerate}
            \item \textbf{저금리 환경 (2008-2022):} 대출과 같은 이자 지급 자산의 수익률이 감소하여 초과 지급준비금 보유의 기회비용이 낮아졌습니다.
            \item \textbf{지급준비금에 대한 이자 지급:} 연방준비제도가 지급준비금에 이자를 지급하기 시작하여 초과 지급준비금 보유의 기회비용이 더욱 낮아졌습니다.
            \item \textbf{금융 위기 경험:} 금융 위기 동안 은행들은 대규모 현금 수요를 경험했고, 살아남은 은행들은 안전장치로 더 많은 지급준비금을 보유하기로 결정했습니다.
        \end{enumerate}
    \end{itemize}
\end{itemize}

\subsubsection*{연방기금 시장의 작동 방식}
은행이 유동성 수요를 충족시킬 만큼 충분한 지급준비금을 가지고 있지 않을 때 어떻게 될까요?

\begin{itemize}
    \item 은행은 다른 은행으로부터 지급준비금을 빌릴 수 있습니다. 이 거래가 \textbf{연방기금 시장 (Federal Funds Market)}에서 이루어집니다.
    \item \textbf{특징:}
    \begin{itemize}
        \item 대출은 보통 익일물(overnight)이며, 담보를 요구하지 않습니다.
        \item 이는 은행이 채무 불이행을 하지 않을 것이라는 높은 신뢰를 요구합니다.
        \item 이 시장에서 빌리는 비용을 \textbf{연방기금 금리 (Federal Funds Rate)}라고 합니다.
        \item 연방기금 시장은 익일물 현금을 제공하는 소위 \textbf{머니 마켓 (Money Market)}입니다.
    \end{itemize}
\end{itemize}

\subsubsection*{연방기금 시장의 중요성}
연방기금 시장은 두 가지 핵심 기능을 수행합니다.

\begin{enumerate}
    \item \textbf{은행 간 자금 재분배:}
    \begin{itemize}
        \item 일부 은행은 대출 기회가 많지만 현금이 부족합니다. 다른 은행은 지급준비금 형태로 현금이 많지만 대출 기회가 적습니다.
        \item 연방기금 시장을 통해 후자의 은행은 전자의 은행에 현금을 빌려줄 수 있습니다. 이를 통해 전자의 은행은 대출 수요를 충족시킬 수 있습니다.
    \end{itemize}
    \item \textbf{경제 전반의 신용 비용에 영향:}
    \begin{itemize}
        \item 연방기금 금리는 경제 전반의 신용 비용에 영향을 미칩니다.
        \item 은행은 지급준비금을 보유할지, 아니면 모기지와 같은 대출을 해줄지 결정합니다.
        \item 만약 연방기금 금리가 충분히 높다면, 은행은 새로운 모기지에 대해 높은 이자를 받아야 합니다. 그렇지 않으면 은행은 단순히 지급준비금을 보유하고 연방기금 시장에 빌려줄 수 있기 때문입니다.
        \item 따라서 연방기금 금리는 모기지 비용에 영향을 미칠 수 있습니다. 이 논리는 모기지뿐만 아니라 다른 모든 자산에도 적용됩니다. 모든 시장의 금리는 연방기금 금리와 연결되어 있습니다.
        \item 그러나 단기 만기 자산에 더 강한 영향을 미칠 것입니다. 단기 만기 자산은 연방기금 시장에 더 가까운 대체재이기 때문입니다. 이러한 대체 익일물 현금 시장은 연방기금 금리 변화에 즉시 영향을 받습니다.
    \end{itemize}
\end{enumerate}

\begin{infobox}[title=\textbf{중앙은행의 역할}]
미국 중앙은행인 연방준비제도(Federal Reserve)는 정책 목표를 달성하기 위해 연방기금 시장의 지급준비금 공급을 통제하여 연방기금 금리를 조절합니다. 이를 \textbf{공개 시장 조작 (Open Market Operations)}이라고 합니다.
\end{infobox}

\subsubsection*{핵심 요약: 연방기금 시장}
\begin{itemize}
    \item 은행은 법적으로 요구되는 지급준비금과 예상치 못한 유동성 수요를 충족시키기 위한 초과 지급준비금을 보유합니다.
    \item 연방기금 시장은 은행 시스템 전반에 걸쳐 현금을 재분배합니다.
    \item 연방기금 금리는 다른 모든 금리에 영향을 미칩니다.
\end{itemize}


\subsection{환매조건부채권 (Repo) 시장 (The Repo Market)}
이 강의에서는 환매조건부채권(Repurchase Agreement) 시장을 논의합니다.

\subsubsection*{환매조건부채권 시장의 규모}
\begin{itemize}
    \item 환매조건부채권(Repo) 시장은 익일물 차입 및 대출의 핵심 시장입니다.
    \item \textbf{규모:} 2020년 말까지 4조 달러 이상의 환매조건부채권 대출이 미상환 상태였습니다.
    \item \textbf{인기 이유:} 환매조건부채권 시장은 가장 큰 현금 시장이며, 모든 유형의 금융 기관이 참여할 수 있습니다. 이는 은행만 접근할 수 있는 연방기금 시장과의 주요 차이점입니다.
    \item \textbf{주요 참여자:} 머니 마켓 뮤추얼 펀드(Money Market Mutual Funds)는 환매조건부채권 시장에서 가장 큰 대출 기관 중 하나입니다. 2019년 말, 이들은 약 1조 달러의 미상환 환매조건부채권을 보유하고 있었습니다.
\end{itemize}

\subsubsection*{환매조건부채권 시장의 작동 방식}
환매조건부채권 시장은 현금 시장입니다. 차입자는 대출자에게 익일물 현금을 빌립니다.

\begin{itemize}
    \item \textbf{담보 거래 (Secured Transaction):} 환매조건부채권 시장의 거래는 담보가 제공되는 거래입니다.
    \item \textbf{예시:} 오늘 청구서를 지불하기 위해 현금이 필요하고 국채(Treasury bond)를 가지고 있다고 가정해 봅시다. 내일은 현금이 풍부해질 것이므로 오늘 현금 부족을 겪고 있습니다.
    \begin{itemize}
        \item \textbf{대안 1 (비효율적):} 오늘 국채를 팔고, 청구서를 지불한 다음, 내일 새로운 국채를 살 수 있습니다. 그러나 채권을 사고파는 것은 비용이 많이 들 수 있습니다.
        \item \textbf{대안 2 (Repo):} 필요한 현금을 빌릴 수 있습니다. 환매조건부채권 거래에서는 국채를 현금과 교환하여 팔고, 내일 다시 사들이기로 약속합니다. 이를 \textbf{환매조건부채권 매도 (selling a repo)}라고 합니다.
    \end{itemize}
    \item \textbf{역할:}
    \begin{itemize}
        \item \textbf{환매조건부채권 매도자 (Repo Seller):} 다른 당사자로부터 현금을 빌리는 사람입니다.
        \item \textbf{환매조건부채권 매수자 (Repo Buyer):} 현금 대출을 제공하는 사람입니다. 매수자는 거래 기간 동안 국채의 새로운 소유자가 됩니다. 즉, 이는 담보 거래입니다.
        \item \textbf{이자:} 매수자는 보통 약간 더 높은 가격으로 국채를 다시 사들입니다. 이로써 매도자는 현금 대출에 대한 보상을 받습니다.
    \end{itemize}
\end{itemize}

\subsubsection*{환매조건부채권 거래의 위험과 헤어컷 (Haircut)}
\begin{itemize}
    \item \textbf{대출자의 주요 위험:} 차입자가 내일 국채를 다시 사들이지 못할 수 있습니다. 이 경우 대출자는 가치 없는 채권을 보유하게 될 수 있습니다. 이는 대출자가 담보를 보유하는 동안 담보 가치가 하락할 때 발생할 가능성이 더 높습니다.
    \item \textbf{구체적인 예시:}
    \begin{itemize}
        \item 당신이 오늘 200달러 가치의 국채를 팔고 내일 200.02달러에 다시 사들이기로 약속했다고 가정해 봅시다.
        \item 경제 상황이 변하여 시장 가격이 199달러로 떨어지면 어떻게 될까요? 당신은 약속대로 200.02달러에 다시 사들이거나, 약속을 어기고 공개 시장에서 199달러에 다른 국채를 살 수 있습니다.
        \item 대출자는 어떻게 될까요? 그들은 199달러 가치의 국채를 보유하고 있지만 200달러를 빌려주었습니다. 그들은 손실을 볼 가능성이 높습니다.
    \end{itemize}
    \item \textbf{헤어컷 (Haircut):} 이러한 위험을 줄이기 위해 환매조건부채권 거래에는 일반적으로 \textbf{헤어컷}이 적용됩니다. 대출자는 대출 금액에 헤어컷을 적용합니다. 담보의 전체 금액을 빌려주는 대신, 일부만 빌려줍니다.
    \item \textbf{헤어컷의 효과:} 헤어컷은 대부분의 환매조건부채권을 대출자에게 사실상 무위험하게 만듭니다.
    \item \textbf{예시 (헤어컷 적용):} 대출자가 2\%의 헤어컷을 부과한다고 가정해 봅시다. 따라서 차입자는 200달러 가치의 국채에 대해 196달러만 받습니다.
    \begin{itemize}
        \item 시장 가격이 199달러로 떨어져도 차입자는 여전히 채권을 다시 사들이고 싶어 할 것입니다 (199달러에 시장에서 사는 것보다 196달러에 다시 사는 것이 이득이므로).
        \item 채권 가격이 196달러 이상으로 유지된다면 대출자는 안전합니다.
    \end{itemize}
\end{itemize}

\subsubsection*{헤어컷의 크기를 결정하는 요인}
헤어컷의 크기는 담보와 관련된 세 가지 주요 요인에 의해 결정됩니다.

\begin{enumerate}
    \item \textbf{담보의 신용 품질 (Credit Quality):} 신용 위험이 낮을수록 헤어컷이 낮아집니다.
    \item \textbf{담보의 시장 가격 변동성 (Market Price Volatility):} 변동성이 낮을수록 헤어컷이 낮아집니다.
    \item \textbf{전반적인 시장 상황 (Overall Market Condition):} 담보를 쉽게 팔 수 있을수록 헤어컷이 낮아집니다.
\end{enumerate}

\subsubsection*{담보의 종류}
\begin{itemize}
    \item 다양한 자산이 환매조건부채권 거래에 사용될 수 있습니다. 그러나 가장 일반적인 담보는 \textbf{정부 보증 증권 (government backed securities)}입니다.
    \item \textbf{국채 (Treasuries)가 인기 있는 이유:}
    \begin{itemize}
        \item 신용 위험이 없습니다.
        \item 시장 가격이 안정적입니다.
        \item 거래하기 쉽습니다.
        \item 따라서 환매조건부채권 매수자에게 거의 위험을 주지 않으며 헤어컷도 낮습니다.
    \end{itemize}
    \item \textbf{일반 담보 환매조건부채권 (General Collateral Repo):} 정부 보증 증권을 포함하는 환매조건부채권 거래는 위험이 낮습니다. 이 경우 대출자는 특정 담보 유형에 대해 덜 까다롭습니다. 담보가 지정되기 전에 합의되는 환매조건부채권 거래를 일반 담보 환매조건부채권이라고 합니다. 여기서 일반 담보는 적격 정부 보증 증권 목록을 의미합니다. 일반 담보 환매조건부채권은 헤어컷이 가장 낮으며 가장 일반적인 유형의 환매조건부채권 거래입니다.
\end{itemize}

\subsubsection*{핵심 요약: 환매조건부채권 시장}
\begin{itemize}
    \item 환매조건부채권은 단기 담보 대출의 한 형태입니다.
    \item 환매조건부채권 시장은 은행뿐만 아니라 모든 금융 기관에게 중요한 현금 시장입니다.
    \item 헤어컷은 대출자의 손실 위험을 줄입니다.
\end{itemize}


\subsection{단기 금리 비교 (Short-term Interest Rates)}
이 강의에서는 현금 시장에서 중요한 두 가지 단기 금리를 소개합니다. 이 금리들은 모기지 및 신용카드 대출 비용을 통해 가계에, 사업 대출 비용을 통해 기업에 영향을 미칩니다.

\subsubsection*{런던 은행 간 대출 금리 (LIBOR)}
\begin{itemize}
    \item \textbf{정의:} LIBOR는 런던 은행 간 시장의 금리입니다. 이 시장은 은행들이 서로 돈을 빌려주고 빌려주는 시장입니다.
    \item \textbf{통화:} 대부분의 대출은 미국 달러로 이루어지지만, 은행들은 다른 통화도 빌릴 수 있습니다. 유럽에서 이루어지는 미국 달러 대출은 유로달러 대출(Euro dollar loans)이라고도 불립니다.
    \item \textbf{특징:}
    \begin{itemize}
        \item 이 시장의 대출은 \textbf{무담보 (unsecured)}입니다. 즉, 담보가 필요하지 않습니다.
        \item 담보가 필요 없기 때문에, 이 은행 간 대출은 \textbf{신용 위험 (credit risk)}에 노출됩니다. 은행이 대출금을 상환하지 못할 가능성이 있습니다.
        \item \textbf{만기:} 익일물(overnight), 1주, 1개월, 2개월, 3개월, 6개월, 12개월의 7가지 만기가 존재합니다.
    \end{itemize}
    \item \textbf{계산 방식:} LIBOR는 실제 거래보다는 \textbf{설문 조사 (surveys)}를 기반으로 계산됩니다. 매일 가장 크고 거래량이 많은 은행들을 대상으로 설문 조사를 실시하며, LIBOR는 이 설문 조사 결과에 기반한 평균 금리입니다.
\end{itemize}

\begin{examplebox}[title=\textbf{LIBOR와 연방기금 금리의 관계}]
익일물 미국 달러 LIBOR는 최근 몇 년간 미국 은행 간 금리인 유효 연방기금 금리(Effective Federal Funds Rate)와 유사한 패턴을 보입니다. 이는 미국과 해외에서 미국 달러를 빌리는 비용이 매우 밀접하게 움직인다는 것을 의미합니다.

\textbf{왜 그럴까요?} 만약 LIBOR가 더 낮다면, 은행들은 런던에서 달러를 빌려 미국에 빌려줌으로써 이익을 얻을 수 있습니다. 이러한 기본적인 \textbf{차익 거래 (arbitrage)}는 LIBOR와 연방기금 금리가 대략적으로 같도록 보장합니다.
\end{examplebox}

\subsubsection*{LIBOR와 광범위한 경제의 연결}
\begin{itemize}
    \item LIBOR는 은행의 자금 조달 비용을 측정하므로, 기업 대출과 같은 많은 다른 대출의 \textbf{기준 금리 (reference rate)}가 되었습니다.
    \item 이러한 대출은 종종 \textbf{변동 금리 (floating rate)}입니다. 이는 이자율이 변동 기준 금리(예: LIBOR)에 고정된 가산금리(mark up)를 더한 것과 같다는 의미입니다.
    \item \textbf{예시:} 월마트의 2019년 규제 서류에 따르면, 2020년 5월부터 2024년 5월 사이에 만기가 도래하는 약정 신용 한도는 일반적으로 LIBOR + 10bp에서 LIBOR + 75bp 사이의 이자율을 적용받습니다.
    \item 2020년까지 LIBOR는 전 세계적으로 250조 달러 이상의 대출을 뒷받침했습니다.
\end{itemize}

\subsubsection*{LIBOR의 문제점: 조작 논란과 대체}
\begin{itemize}
    \item 2000년부터 2020년까지의 기간을 보면, LIBOR가 여러 차례 급등하여 연방기금 금리를 따르지 않는 경우가 있었습니다.
    \item \textbf{원인 1: 신용 위험:} LIBOR는 무담보 대출이므로, 런던 시장 은행들의 건전성에 의문이 제기되면 은행 간 금리 사이에 격차가 발생할 수 있습니다.
    \item \textbf{원인 2: 조작 논란:} 2008년 금융 위기 이후, LIBOR 조작 의혹이 제기되었습니다. LIBOR는 가장 큰 은행들의 금리 설문 조사를 기반으로 하며, 이 설문 조사는 정직성을 요구합니다. 그러나 이 설문 조사가 실제로 조작되었다는 증거가 나왔습니다. 총 90억 달러 이상의 벌금이 은행들에 부과되었습니다.
    \item \textbf{대체:} 추가 조작을 방지하기 위해 규제 당국은 LIBOR를 대출의 기준 금리에서 제외하도록 요구했습니다. 2022년부터 새로운 금리가 의무화되었습니다.
\end{itemize}

\subsubsection*{담보부 익일물 자금 조달 금리 (SOFR)}
\begin{itemize}
    \item \textbf{정의:} 새로운 금리는 \textbf{담보부 익일물 자금 조달 금리 (Secured Overnight Financing Rate, SOFR)}라고 불립니다. 뉴욕 연방준비은행이 매일 발표합니다.
    \item \textbf{특징:}
    \begin{itemize}
        \item SOFR 금리는 국채를 담보로 하는 익일물 환매조건부채권(repo) 거래의 금리입니다.
        \item SOFR은 설문 조사가 아닌 \textbf{실제 거래 (transactions)}를 기반으로 합니다.
        \item SOFR은 \textbf{담보부 금리 (secured rate)}이므로 \textbf{무위험 (risk free)}으로 간주됩니다.
        \item 반면 LIBOR는 무담보 금리이므로 일부 신용 위험을 포함합니다.
    \end{itemize}
    \item \textbf{규모:} 2020년 기준으로 SOFR은 약 1조 달러 규모의 일일 거래 시장을 기반으로 합니다. LIBOR는 이 규모의 절반 미만에 기반했습니다.
\end{itemize}

\subsubsection*{핵심 요약: 단기 금리}
\begin{itemize}
    \item LIBOR와 SOFR은 단기 금리이며, 각각 무담보 및 담보 대출에 대한 은행의 자금 조달 비용을 나타냅니다.
    \item LIBOR는 신용 위험에 노출되어 있으며 과거에 조작된 사례가 있습니다.
    \item 규제 당국은 무위험 금리인 SOFR이 LIBOR를 기준 금리로 대체하도록 의무화했습니다.
\end{itemize}


\section{장기 금리 (Long-term Rates)}

\subsection{명목 금리와 실질 금리 (Nominal and Real Interest Rates)}
이 강의에서는 명목 금리와 실질 금리를 논의하고, 이 두 금리를 인플레이션율과 연결하는 피셔 방정식(Fisher Equation)을 살펴봅니다. 마지막으로, 실질 금리를 지급하는 몇 안 되는 증권 중 하나인 물가연동국채(Treasury Inflation Protected Securities, TIPS)를 사례 연구로 살펴봅니다.

\subsubsection*{명목 수익률과 실질 수익률의 차이}
투자에 대한 명목 수익률과 실질 수익률의 차이를 예시를 통해 이해할 수 있습니다.

\begin{itemize}
    \item \textbf{예시:} 은행 계좌에 100달러가 있고, 연 10\%의 이자를 지급한다고 가정해 봅시다.
    \begin{itemize}
        \item 1년 후, 10달러를 얻어 총 110달러가 됩니다.
        \item 은행이 지급하는 이 10달러가 100달러 투자에 대한 \textbf{명목 수익률 (nominal return)}입니다. 이 10\%의 이자율을 \textbf{명목 금리 (nominal interest rate)}라고 합니다.
    \end{itemize}
    \item \textbf{실질 수익률 (real return):} 이 이익이 1년 동안 구매하고 소비할 수 있는 실제 상품의 가치로 얼마나 되는가? 이것이 투자에 대한 실질 수익률이 의미하는 바입니다.
\end{itemize}

\begin{examplebox}[title=\textbf{명목 금리와 실질 금리 예시: 주유비}]
\begin{itemize}
    \item \textbf{현재:} 오늘 휘발유 1갤런을 4달러에 살 수 있습니다. 따라서 100달러는 25갤런의 휘발유 가치입니다.
    \item \textbf{인플레이션 가정:} 인플레이션, 즉 경제 전반의 물가 상승률이 5\%라고 가정해 봅시다. 내년에는 휘발유 가격이 4.20달러로 오릅니다. 같은 25갤런의 휘발유를 구매하려면 105달러가 필요합니다.
    \item \textbf{투자 수익률:}
    \begin{itemize}
        \item \textbf{명목 수익률:} 은행은 여전히 저축에 대해 10달러를 지급하므로, 명목 수익률은 10달러입니다.
        \item \textbf{실질 수익률:} 휘발유 가격이 올랐습니다. 25갤런의 휘발유를 구매한 후에는 5달러가 남습니다. 따라서 인플레이션을 고려한 실질 수익률은 5달러입니다.
        \item 명목 수익률 10달러에서 물가 상승분 5달러를 빼면 5달러의 실질 수익률이 나옵니다.
    \end{itemize}
\end{itemize}
이 예시는 명목 금리와 실질 금리가 인플레이션에 의해 연결되어 있음을 보여줍니다.
\end{examplebox}

\subsubsection*{피셔 방정식 (Fisher Equation)}
일반적으로 \textbf{실질 금리 (real interest rate)}는 \textbf{명목 금리 (nominal interest rate)}에서 \textbf{인플레이션 (inflation)}을 뺀 값과 거의 같습니다.

\begin{infobox}[title=\textbf{피셔 방정식}]
명목 금리 $i$는 실질 금리 $r$에 인플레이션 $\pi$를 더한 값과 대략적으로 같습니다.
$$ i \approx r + \pi $$
이 방정식은 미국의 경제학자 어빙 피셔(Irving Fisher)의 이름을 따서 명명되었습니다.
\end{infobox}

\subsubsection*{사례 연구: 물가연동국채 (TIPS)}
물가연동국채(Treasury Inflation Protected Securities, TIPS)는 인플레이션을 명시적으로 고려하는 채권 투자입니다.

\begin{itemize}
    \item \textbf{작동 방식:} TIPS는 채권의 원금 가치가 인플레이션율에 따라 증가합니다. 채권에 지급되는 이자는 원금 가치의 백분율이므로, TIPS에 지급되는 이자도 인플레이션에 따라 증가합니다. 따라서 TIPS의 이자율은 \textbf{실질 금리}입니다.
    \item \textbf{구체적인 예시:}
    \begin{itemize}
        \item 100달러짜리 TIPS 채권을 구매하고 연 5\%의 이자를 지급한다고 가정해 봅시다.
        \item 올해 인플레이션이 2\%라고 가정하면, 내년에 TIPS 채권의 원금은 2\% 증가하여 102달러가 됩니다.
        \item \textbf{채권의 연간 총 수익률:} 5\% 이자 + 2\% 원금 가치 상승 = 총 7\%의 명목 수익률.
        \item \textbf{실질 수익률:} 피셔 방정식에 따르면, 실질 수익률은 명목 수익률에서 인플레이션을 뺀 값입니다. 2\%의 원금 가치 상승은 2\%의 인플레이션을 상쇄합니다. 따라서 실질 수익률은 5\%입니다.
        \item TIPS는 일정한 실질 수익률을 제공하도록 설계되었습니다.
    \end{itemize}
\end{itemize}

\subsubsection*{TIPS 시장의 중요성: 미래 인플레이션 기대}
경제학자들은 TIPS 시장에 왜 주목할까요?

\begin{itemize}
    \item 오늘 투자를 할 때, 우리는 미래 인플레이션이 어떻게 될지 알지 못합니다. 인플레이션은 실질 수익률에 중요합니다.
    \item TIPS 채권의 가격을 사용하여 시장의 \textbf{미래 인플레이션 기대 (market's expectations of future inflation)}를 파악할 수 있습니다.
    \item \textbf{작동 방식:}
    \begin{itemize}
        \item TIPS는 일정한 실질 수익률을 제공하며 인플레이션 보호 기능을 제공합니다.
        \item 예상 인플레이션이 높으면 인플레이션 보호의 가치가 커집니다. TIPS 채권은 인플레이션 보호를 제공하므로, TIPS의 현재 시장 가격은 상승할 것입니다. 가격이 높으면 TIPS의 수익률은 낮아집니다.
        \item 반대로, 인플레이션이 낮을 것으로 예상되면 인플레이션 보호의 가치가 떨어집니다. 이 경우 TIPS 채권 가격은 하락하고, TIPS 채권의 수익률은 상승합니다.
        \item 따라서 TIPS 가격 또는 수익률은 인플레이션 기대를 포착합니다. 실제로 TIPS 수익률과 일반 국채 수익률을 비교하여 예상되는 미래 인플레이션을 측정할 수 있습니다.
    \end{itemize}
\end{itemize}

\begin{examplebox}[title=\textbf{TIPS 수익률 데이터 분석}]
5년 만기 TIPS 채권 수익률과 5년 만기 국채 수익률 데이터를 보면, 두 수익률 사이에 차이가 나는 경우가 있습니다. 수익률 차이가 크다는 것은 투자자들이 인플레이션 보호에 기꺼이 비용을 지불한다는 의미이며, 이는 투자자들이 5년 후 높은 미래 인플레이션을 예상한다는 것을 뜻합니다.

\begin{itemize}
    \item \textbf{2009년:} TIPS 수익률이 높았습니다. 투자자들은 이 시기에 더 높은 인플레이션을 예상했습니다.
    \item \textbf{2013년 초 및 2021년:} 반대의 상황이었습니다. 투자자들은 낮은 미래 인플레이션을 예상했습니다.
\end{itemize}
물론 시장의 예측이 결국 틀릴 수도 있지만, 이는 TIPS 가격에 반영된 기대치였습니다.
\end{examplebox}

\subsubsection*{핵심 요약: 명목 금리와 실질 금리}
\begin{itemize}
    \item 명목 금리는 상품 가격 변화를 고려하지 않습니다.
    \item 명목 금리는 실질 금리 + 인플레이션과 같습니다. 이것이 피셔 방정식입니다.
    \item TIPS 채권은 인플레이션 기대를 포착합니다.
\end{itemize}


\subsection{국채와 무위험 금리 (Treasury Securities and The Risk-Free Rate)}
이 강의에서는 대표적인 무위험 자산인 미국 국채(US Treasury Securities)를 논의합니다.

\subsubsection*{미국 국채 시장의 특징}
\begin{itemize}
    \item \textbf{세계 최대 채권 시장:} 2020년 기준으로 미국 국채 시장은 전 세계에서 가장 큰 채권 시장입니다. 당시 미상환 증권 가치는 24조 달러를 넘었습니다.
    \item \textbf{다양한 투자자:} 거의 모든 유형의 미국 금융 기관(은행, 뮤추얼 펀드, 보험 회사, 연기금, 연방준비제도 등)이 국채를 소유하고 있습니다. 또한 국채의 약 3분의 1은 국제 투자자들이 소유하고 있습니다.
\end{itemize}

\subsubsection*{미국 국채의 종류}
미국 국채는 만기에 따라 세 가지 범주로 나뉩니다. 이들은 평생 고정 이자 지급을 하는 정부 채권입니다.

\begin{enumerate}
    \item \textbf{재무부 단기채 (Treasury Bills, T-Bills):}
    \begin{itemize}
        \item 만기가 가장 짧습니다 (4주에서 52주).
        \item 미국 국채 증권의 약 20\%를 차지합니다.
    \end{itemize}
    \item \textbf{재무부 중기채 (Treasury Notes, T-Notes):}
    \begin{itemize}
        \item 중간 만기 범위의 증권입니다 (2년에서 10년).
        \item 미상환 미국 국채 증권의 대부분을 차지합니다.
    \end{itemize}
    \item \textbf{재무부 장기채 (Treasury Bonds, T-Bonds):}
    \begin{itemize}
        \item 만기가 20년에서 30년인 점을 제외하고는 재무부 중기채와 본질적으로 동일합니다.
        \item '장기 채권 (long bond)'이라고도 불립니다.
    \end{itemize}
\end{enumerate}

\subsubsection*{미국 국채가 투자자들에게 인기 있는 이유}
\begin{itemize}
    \item \textbf{신용 위험 없음 (Credit Risk-Free):} 미국 국채는 신용 위험이 없는 것으로 간주됩니다. 즉, 투자자들은 만기까지 국채를 보유하는 한 이자와 원금 지급을 보장받습니다. 이러한 지급은 미국 정부의 완전한 신뢰와 신용(full faith and credit)으로 보장됩니다.
    \item \textbf{담보 활용:} 미국 국채는 신용 위험이 없기 때문에 차입 약정에서 담보로 자주 사용됩니다. 예를 들어, 수조 달러 규모의 환매조건부채권(repo) 시장에서 담보로 활용됩니다.
    \item \textbf{무위험 금리 (Risk-Free Interest Rate)의 척도:} 국채의 이자율은 무위험 금리의 좋은 척도로 간주됩니다. 무위험 금리는 전 세계 많은 금융 계약에서 벤치마크 이자율로 사용됩니다.
\end{itemize}

\begin{infobox}[title=\textbf{다른 선진국 국채보다 미국 국채가 선호되는 이유}]
독일, 중국, 일본과 같이 채무 불이행 위험이 없는 다른 선진국의 국채보다 미국 국채가 선호되는 세 가지 주요 이유가 있습니다.
\begin{enumerate}
    \item \textbf{시장 규모와 유동성:} 미국 국채 시장은 다른 어떤 시장보다 훨씬 큽니다. 언제든지 다양한 투자자들이 국채 구매에 관심을 가지므로, 거래 비용이 거의 없이 쉽게 매도할 수 있습니다. 시장에서 증권이 적은 비용으로 쉽게 거래될 수 있을 때, 그 시장을 \textbf{유동적 (liquid)}이라고 합니다. 미국 국채 시장은 세계에서 가장 유동적인 증권 시장입니다.
    \item \textbf{달러 자산 보유 수단:} 전 세계 대부분의 금융 거래는 미국 달러로 이루어집니다. 미국 국채를 보유하는 것은 미국 달러를 보유하는 한 가지 방법입니다. 이러한 이유로 미국 국채는 종종 \textbf{현금성 자산 (cash-like assets)}이라고 불립니다. 2020년 말까지 기업과 가계는 미국 국채에만 투자하는 머니 마켓 뮤추얼 펀드에 2조 달러 이상의 현금을 예치했습니다.
    \item \textbf{중앙은행의 달러 준비금:} 전 세계 중앙은행들이 달러 준비금을 보유하는 가장 일반적인 방법이 미국 국채입니다.
\end{enumerate}
\end{infobox}

\subsubsection*{미국 국채 투자에 항상 위험이 없을까?}
만기까지 국채를 보유하면 지급을 받으므로 이 경우에는 위험이 없습니다. 그러나 만기 전에 매도해야 하는 경우 손실을 볼 수 있습니다. 이는 채권의 현재 시장 가치가 하락할 때 발생합니다. 국채 투자자들은 두 가지 주요 위험 원인에 대해 우려합니다.

\begin{enumerate}
    \item \textbf{인플레이션 (Inflation):}
    \begin{itemize}
        \item \textbf{예시:} 액면가 100달러, 명목 이자율 2\%인 52주 만기 T-Bill을 구매했다고 가정해 봅시다. 인플레이션이 2\%일 것으로 예상했다고 가정합니다. 만기 시 100달러의 액면가와 2달러의 이자를 받습니다. 동시에 T-Bill을 구매했을 때 100달러였던 상품은 인플레이션으로 인해 102달러가 됩니다. 벌어들인 이자로 T-Bill을 구매했을 때 살 수 있었던 것과 동일한 상품을 여전히 살 수 있습니다.
        \item \textbf{예상보다 높은 인플레이션:} 올해 인플레이션이 5\%라고 가정해 봅시다. 만기 시 여전히 100달러의 액면가와 2달러의 이자를 받습니다. 그러나 이제 상품은 더 높은 인플레이션으로 인해 105달러가 됩니다. 이는 이전과 동일한 상품을 살 수 없다는 의미입니다.
        \item \textbf{결과:} 명목 수익률은 2\%였지만, 인플레이션을 조정한 실질 수익률은 마이너스 3\%가 됩니다. 인플레이션이 예상보다 높으면 국채의 시장 가치는 하락합니다. 현재 관점에서 국채는 구매력이 적으므로 가치가 떨어집니다.
    \end{itemize}
    \item \textbf{금리 변화 (Changes in the Interest Rate):}
    \begin{itemize}
        \item \textbf{예시:} 100달러에 2\% 이자를 지급하는 2년 만기 국채를 구매했다고 가정해 봅시다. 1년 후 이 투자를 팔고 싶다고 가정합니다. 동시에 이 투자에는 1년 만기가 남아 있습니다. 이 채권으로 얼마를 받을 수 있을까요?
        \item \textbf{결과:} 이는 가장 가까운 대체재인 새로 발행된 52주 만기 T-Bill의 이자율에 따라 달라집니다. 새로 발행된 52주 만기 T-Bill의 이자율이 더 높다면(예: 3\%), 이는 2\%만 지급하는 당신의 채권보다 더 매력적인 투자입니다. 새로운 T-Bill이 100달러라면, 당신의 투자는 100달러 미만이 될 것입니다.
        \item \textbf{결론:} 금리가 상승하면 이전에 발행된 장기 국채의 시장 가치는 하락할 수 있습니다. 장기 국채는 채권의 수명 동안 금리가 변동할 위험이 더 큽니다. 이 위험에 대한 보상을 \textbf{기간 프리미엄 (term premium)}이라고 합니다.
    \end{itemize}
\end{enumerate}

\subsubsection*{장기 국채 명목 금리의 구성 요소}
이러한 위험은 장기 국채의 명목 금리에 반영됩니다. 세 가지 구성 요소가 있습니다.

\begin{enumerate}
    \item \textbf{실질 금리 (Real Interest Rate):} 인플레이션을 고려한 금리입니다.
    \item \textbf{예상 인플레이션율 (Expected Rate of Inflation):} 미래에 예상되는 물가 상승률입니다.
    \item \textbf{기간 프리미엄 (Term Premium):} 금리 변화를 포함한 장기 위험을 포착합니다.
\end{enumerate}

\subsubsection*{핵심 요약: 국채와 무위험 금리}
\begin{itemize}
    \item 미국 국채 시장은 세계에서 가장 크고 유동적인 채권 시장입니다.
    \item 미국 국채는 신용 위험이 없습니다. 이는 그 이자율이 무위험 금리의 좋은 대용치라는 것을 의미합니다.
    \item 미국 국채의 장기 명목 금리는 실질 금리, 예상 인플레이션, 기간 프리미엄으로 분해될 수 있습니다.
\end{itemize}


\subsection{금리 기간 구조 (Term Structure of Interest Rates)}
이 강의에서는 금리 기간 구조가 무엇인지 논의합니다. 금리 기간 구조가 거시 경제에 대한 중요한 정보를 제공하는 이유를 살펴봅니다. 특히, 경기 침체 전에 금리 기간 구조가 어떻게 행동하는지 알아봅니다.

\subsubsection*{금리 기간 구조의 정의}
\begin{itemize}
    \item \textbf{정의:} \textbf{금리 기간 구조 (Term Structure of Interest Rates)}는 특정 시점에 동일한 발행 주체가 발행한 만기가 다른 증권의 금리를 비교하는 것을 의미합니다.
    \item \textbf{수익률 (Yields) 사용:} 금리 기간 구조에는 채권 소유자에게 지급되는 쿠폰 이자율(coupon interest rates)이 아니라, 채권 가격을 고려한 \textbf{수익률 (yields)}을 사용합니다.
    \item \textbf{예시:} 액면가 100달러에 5\% 이자를 지급하는 채권이 있다고 가정해 봅시다.
    \begin{itemize}
        \item 채권을 99달러에 구매하면 여전히 5달러의 이자를 받습니다. 현재 수익률은 5/99, 즉 5.051\%입니다.
        \item 반대로, 같은 5\% 이자 지급 채권을 101달러에 구매하면 수익률은 5/101, 즉 4.95\%가 됩니다.
    \end{itemize}
    \item \textbf{채권 가격과 수익률의 역관계:} 채권 수익률과 채권 가격은 반대 방향으로 움직입니다. 채권 가격이 상승하면 수익률은 하락하고, 채권 가격이 하락하면 수익률은 상승합니다.
\end{itemize}

\subsubsection*{수익률 곡선 (Yield Curve)}
\begin{itemize}
    \item 우리는 매일 수익률의 변화와 그에 따른 금리 기간 구조를 관찰할 수 있으며, 이를 일반적으로 \textbf{수익률 곡선 (Yield Curve)}이라고 합니다.
    \item \textbf{가장 중요한 수익률 곡선:} 미국 국채(US Treasury Securities)를 기반으로 합니다.
    \item \textbf{미국 국채가 좋은 척도가 되는 세 가지 이유:}
    \begin{enumerate}
        \item 시간이 지남에 따라 변할 수 있는 신용 위험이 없습니다.
        \item 다양한 만기로 발행됩니다.
        \item 증권 가격과 수익률을 왜곡할 수 있는 거래 비용이 거의 없습니다.
    \end{enumerate}
\end{itemize}

\begin{examplebox}[title=\textbf{실제 수익률 곡선 예시: 2021년 6월 1일}]
2021년 6월 1일의 미국 국채 수익률 곡선을 보면, 단기 만기 채권의 수익률은 거의 0\%에 가깝습니다. 1년 만기 미국 국채도 연 0.04\%에 불과합니다. 그러나 10년 만기 국채는 연 1.62\%의 수익률을 보입니다.

\textbf{왜 단기 만기와 장기 만기 수익률에 차이가 있을까요?}
원금을 갚는 데 10년이 걸리는 프로젝트에 자금을 조달하는 두 가지 대안을 생각해 봅시다.
\begin{enumerate}
    \item \textbf{대안 1:} 10년 만기 채권을 발행합니다.
    \item \textbf{대안 2:} 1년 만기 채권을 발행하고, 1년 후 새로운 1년 만기 채권으로 부채를 롤오버(roll over)합니다. 이 경우 9번 더 부채를 롤오버해야 합니다.
\end{enumerate}
만약 1년 만기 채권으로 부채를 롤오버할 때 항상 처음과 동일한 수익률을 얻을 수 있다면, 10번의 1년 만기 채권 발행 비용은 10년 만기 채권 발행 비용과 같을 것이고, 따라서 수익률도 같을 것입니다. 만약 같지 않다면, 더 저렴한 대안을 선택하는 것이 유리할 것입니다.

2021년 6월 1일 수익률 곡선에서는 단기 채권의 수익률이 더 낮습니다. 이는 시장 참여자들이 향후 10년 동안 어느 시점에서 금리가 상승할 것으로 예상한다는 의미입니다. 즉, 1년 만기 채권을 롤오버할 때 처음 발행했던 1년 만기 채권보다 더 높은 수익률을 제공해야 할 것으로 예상하는 것입니다.
\end{examplebox}

\subsubsection*{높은 금리 기대를 유발하는 요인}
\begin{enumerate}
    \item \textbf{연방준비제도의 단기 금리 인상 기대:} 시장 참여자들은 연방준비제도가 단기 금리를 인상할 것으로 예상할 수 있습니다. 이는 일반적으로 경제 확장기(economic expansions)에 예상되며, 이때 수익률 곡선은 일반적으로 \textbf{우상향 (upward sloping)}합니다.
    \item \textbf{높은 인플레이션 기대:} 시장 참여자들은 더 높은 인플레이션을 예상할 수 있습니다. 미국 국채의 금리는 명목 금리이므로, 투자자들을 계속 유치하기 위해서는 장기 만기 국채의 수익률이 인플레이션 증가를 상쇄하기 위해 증가해야 합니다.
    \item \textbf{재정 정책으로 인한 정부 부채 증가:} 정부 부채를 증가시키는 재정 정책은 미국 국채의 공급을 늘릴 것입니다. 이 경우, 투자자들이 추가 공급되는 미국 국채를 구매하도록 유인하기 위해 장기 만기 국채의 수익률이 증가해야 합니다.
\end{enumerate}

\subsubsection*{수익률 곡선과 경기 침체 예측}
\begin{itemize}
    \item \textbf{하향 경사 (Downward Sloping) 또는 역전된 (Inverted) 수익률 곡선:} 2019년 9월 4일의 수익률 곡선을 보면, 곡선이 앞쪽에서 하향 경사 또는 역전되어 있습니다. 이는 시장 참여자들이 미래에 금리가 하락할 것으로 예상해야 한다는 의미입니다.
    \item \textbf{언제 금리 하락을 예상할까요?} 바로 \textbf{경기 침체 (recession)} 전입니다.
    \item \textbf{10년물과 2년물 국채 수익률 차이:} 지난 40년간 10년 만기 국채와 2년 만기 국채 수익률의 차이를 보면, 거의 항상 이 차이가 양수였습니다. 즉, 10년물 국채 수익률이 2년물 국채 수익률보다 높았습니다. 대부분의 시간 동안 시장 참여자들은 금리가 상승할 것으로 예상했으며, 이는 경제 확장과 일치합니다.
    \item \textbf{경기 침체 전의 패턴:} 그림의 회색 막대는 경기 침체를 나타냅니다. 경기 침체기에는 수요가 낮으므로 금리가 하락하는 경향이 있습니다. 회색 막대 바로 직전에 수익률 곡선이 어떻게 되는지 주목하세요. 1980년대, 1990년대 초, 2000년 닷컴 버블, 2008년 금융 위기, 그리고 코로나19 팬데믹 이전의 경기 침체 직전에 매번 역전되었습니다.
\end{itemize}

\begin{summarybox}[title=\textbf{수익률 곡선의 중요성}]
요약하자면, 국채 수익률 곡선은 경제의 후속 성과를 강력하게 예측하는 지표이며, 따라서 시장 참여자와 정책 입안자들이 면밀히 주시합니다.
\end{summarybox}

\subsubsection*{핵심 요약: 금리 기간 구조}
\begin{itemize}
    \item 금리 기간 구조 또는 수익률 곡선은 만기가 다른 채권들의 수익률을 비교합니다.
    \item 수익률 곡선은 면밀히 주시되는 경제 지표입니다.
    \item \textbf{우상향 수익률 곡선:} 시장 참여자들이 미래에 금리가 상승할 것으로 예상하며, 이는 경제 확장기에 발생합니다.
    \item \textbf{역전된 하향 경사 수익률 곡선:} 미래 금리 하락을 시사하며, 일반적으로 경제 확장의 끝을 알리는 신호입니다.
\end{itemize}


\section{기타 자금 시장 (Other Funding Markets)}

\subsection{머니 마켓 뮤추얼 펀드 (Money Market Mutual Funds)}
이 강의에서는 머니 마켓 뮤추얼 펀드(Money Market Mutual Funds)를 자세히 논의합니다. 이 펀드들은 경제에서 특히 은행과 비은행 기관에 단기 자금을 제공하는 데 중요한 역할을 합니다.

\subsubsection*{머니 마켓 펀드의 규모와 인기}
\begin{itemize}
    \item \textbf{규모:} 2021년 기준으로 머니 마켓 펀드는 4조 달러 이상의 자산을 운용했습니다. 이는 기업과 가계의 총 은행 예금보다 많은 금액입니다.
    \item \textbf{코로나19 팬데믹의 영향:} 그래프를 보면 코로나19 팬데믹의 영향으로 저축이 급증한 것을 명확히 알 수 있습니다.
    \item \textbf{인기 이유:} 머니 마켓 펀드는 기업과 가계를 위한 저축 수단입니다.
    \item \textbf{탄생 배경:} 1970년대 미국에서 은행 예금에 지급할 수 있는 이자율을 제한하는 은행 규제인 \textbf{규제 Q (Regulation Q)} 때문에 탄생했습니다.
    \begin{itemize}
        \item 규제 Q에 따라 은행이 고객에게 지급할 수 있는 최대 이자율은 5.25\%로 제한되었습니다.
        \item 1970년대와 1980년대에는 인플레이션과 시장 금리가 5\%보다 훨씬 높았기 때문에 은행 저축 계좌는 매력이 없었습니다.
        \item 머니 마켓 펀드는 이 규제를 받지 않아 더 높은 이자율을 지급할 수 있었습니다.
    \end{itemize}
\end{itemize}

\subsubsection*{머니 마켓 펀드의 사업 모델}
\begin{itemize}
    \item \textbf{투자 프로필:} 머니 마켓 펀드는 초안전 단기 채무 증권에 투자합니다. 이러한 투자는 신용 위험이 거의 없습니다.
    \item \textbf{주요 투자처 (2021년 기준):}
    \begin{itemize}
        \item 미국 국채(US Treasury Securities)
        \item 미국 국채를 담보로 하는 환매조건부채권(Repurchase Agreements)
    \end{itemize}
\end{itemize}

\subsubsection*{머니 마켓 펀드의 종류}
모든 머니 마켓 펀드를 포함하는 이전 차트와 달리, 일부 펀드는 특정 투자에 특화되어 있습니다.

\begin{itemize}
    \item \textbf{국채 펀드 (Treasury Funds):} 오직 미국 국채에만 투자하는 머니 마켓 펀드입니다.
    \item \textbf{정부 펀드 (Government Funds):} 다른 정부 증권에도 투자하는 펀드입니다.
    \item \textbf{특징:} 이들은 가장 안전한 투자로 간주되지만, 수익률도 가장 낮습니다.
    \item \textbf{프라임 펀드 (Prime Funds):} 다른 단기 채무 증권에도 투자하는 머니 마켓 펀드입니다.
    \begin{itemize}
        \item 기업이 발행하는 기업어음(commercial paper)
        \item 은행이 발행하는 양도성 예금증서(certificates of deposit)
        \item 헤지 펀드가 발행하는 환매조건부채권(purchase agreements)
    \end{itemize}
    \item \textbf{특징:} 이러한 투자는 약간 더 위험하지만, 수익률도 약간 더 높습니다.
\end{itemize}

\subsubsection*{머니 마켓 펀드와 은행 예금의 비교}
머니 마켓 펀드는 뮤추얼 펀드의 한 종류입니다. 현금을 머니 마켓 계좌에 예치하면 뮤추얼 펀드의 주식을 구매하는 것입니다. 다른 투자 펀드와 달리 머니 마켓 펀드는 초안전 투자에 집중하고 손실을 피하려고 노력합니다.

\begin{itemize}
    \item 투자자는 오늘 1달러에 머니 마켓 펀드 주식을 구매하면 내일 1달러와 이자를 받고 상환할 수 있을 것으로 기대합니다.
    \item 머니 마켓 펀드 주식은 요구불 상환(redeemable on demand)이 가능합니다. 언제든지 돈을 인출할 수 있습니다.
    \item 이러한 특징 때문에 머니 마켓 펀드 투자는 은행 예금과 매우 유사하게 보입니다.
\end{itemize}

그러나 머니 마켓 펀드와 은행 예금 사이에는 두 가지 중요한 차이점이 있습니다.

\begin{enumerate}
    \item \textbf{이자율:} 머니 마켓 펀드는 거의 항상 더 높은 이자율을 지급합니다. 대규모 급여를 관리하는 기업의 경우, 작은 이자율 차이도 큰 금액을 의미할 수 있습니다.
    \item \textbf{예금 보험 (Deposit Insurance):}
    \begin{itemize}
        \item 은행이 파산하면 정부 예금 보험은 예금자들이 모든 돈을 돌려받을 수 있도록 보장합니다.
        \item 머니 마켓 펀드 투자는 정부에 의해 보험에 가입되어 있지 않습니다. 특히 프라임 펀드의 경우 작은 손실 위험이 있습니다.
    \end{itemize}
\end{enumerate}

\subsubsection*{머니 마켓 펀드의 중요성과 문제 발생 시 영향}
머니 마켓 펀드는 금융 시스템의 핵심입니다. 이들은 단기 자금 시장의 초석이며, 광범위한 금융 기관과 기업에 대규모 현금을 대출해 줍니다. 이러한 자금 조달은 단기간에 대체하기 어렵습니다.

\begin{itemize}
    \item 머니 마켓 펀드가 문제가 생기는 경우는 드물지만, 금융 시장과 광범위한 경제에 매우 파괴적일 수 있습니다.
\end{itemize}

\subsubsection*{사례 연구: 2008년 금융 위기 중 리저브 프라이머리 펀드 (Reserve Primary Fund)}
\begin{itemize}
    \item 2008년 기준으로 리저브 프라이머리 펀드는 미국에서 가장 오래되고 큰 머니 마켓 펀드 중 하나였습니다.
    \item 이 펀드는 투자 은행인 리먼 브라더스(Lehman Brothers)의 기업어음에 크게 투자했습니다.
    \item \textbf{파산:} 리먼 브라더스는 2008년 금융 위기 동안 파산했습니다. 이 소식이 전해지자 공황에 빠진 투자자들은 펀드에서 돈을 인출하기 위해 몰려들었습니다.
    \item \textbf{버크 깨짐 (Breaking the Buck):} 결국 펀드는 투자 주당 1달러로 모든 상환을 충족시킬 수 없었습니다. 파산에 들어갔습니다.
    \textbf{결과:} 리저브 프라이머리 펀드는 결국 주당 1달러당 0.97달러만 돌려줄 수 있었습니다. 이 사건을 \textbf{버크 깨짐 (breaking the buck)}이라고 합니다. 머니 마켓 뮤추얼 펀드에는 예금 보험이 없으므로, 펀드가 버크를 깨면 투자자들은 돈을 잃게 됩니다.
    \item \textbf{금융 시장에 미친 연쇄 반응:}
    \begin{itemize}
        \item 리저브 프라이머리 펀드가 실패한 후, 투자자들은 다른 프라임 펀드의 건전성에 대해 우려하게 되었습니다. 안전해 보이던 투자가 갑자기 위험해 보였습니다.
        \item 그들은 이 펀드에서 약 4천억 달러를 인출했습니다. 이는 머니 마켓의 현금 감소로 이어졌습니다.
        \item 머니 마켓 펀드는 단기 자금 시장에서 철수할 수밖에 없었습니다.
        \item 이 시장에서 보통 돈을 빌리던 기업과 금융 기관들은 현금을 확보하거나 자산을 매각하기 위해 분주했습니다. 은행의 경우, 기업과 가계에 대한 신용 공급을 줄였습니다.
    \end{itemize}
\end{itemize}

\begin{summarybox}[title=\textbf{리저브 프라이머리 펀드 사례의 교훈}]
리저브 프라이머리 펀드의 실패는 머니 마켓 펀드가 금융 시스템에 얼마나 중요한지를 보여줍니다. 이 펀드들이 어려움을 겪고 대규모 상환이 발생하면, 이는 금융 시스템의 많은 부분에 매우 문제가 될 수 있습니다. 금융 위기 이후, 금융 시장 규제 당국은 머니 마켓 뮤추얼 펀드의 미래 혼란을 피하기 위한 개혁을 시행했습니다.
\end{summarybox}

\subsubsection*{핵심 요약: 머니 마켓 뮤추얼 펀드}
\begin{itemize}
    \item 머니 마켓 뮤추얼 펀드는 미국 금융 시스템에서 가장 큰 단기 또는 현금 대출 기관입니다. 이 펀드들은 경제 전반에 걸쳐 현금을 효율적으로 배분하는 데 도움을 줍니다.
    \item 투자자들은 머니 마켓 뮤추얼 펀드를 은행 예금의 대체재로 봅니다. 그러나 이러한 투자는 은행 예금처럼 보험에 가입되어 있지 않습니다.
    \item 이 펀드들의 혼란은 단기 신용 감소를 통해 금융 시스템의 많은 부분에 영향을 미칠 수 있습니다.
\end{itemize}

\subsection{증권화 (Securitization)}
\begin{warningbox}[title=\textbf{참고: 증권화 (Securitization) 내용}]
이 문서의 파트 1/2에는 'Lesson 1-3.2. Securitization'이 목차에 포함되어 있으나, 해당 강의의 본문 내용은 제공되지 않았습니다. 따라서 이 섹션에는 추가적인 설명이 포함되어 있지 않습니다.
\end{warningbox}


\section{빠르게 훑어보기 (1페이지 요약)}

\begin{summarybox}[title=\textbf{모듈 1 핵심 요약: 금리 및 금융 시장}]
\begin{itemize}
    \item \textbf{중앙은행의 역할:} 통화정책은 금리 조정을 통해 경제에 영향을 미치며, 그 핵심은 단기 금리 시장에 있습니다.
    \item \textbf{단기 금리 시장:}
    \begin{itemize}
        \item \textbf{연방기금 시장:} 은행 간 초과 지급준비금 거래 시장. 연방기금 금리는 중앙은행의 주요 정책 목표 금리이자 모든 금리에 영향을 미치는 기준점입니다.
        \item \textbf{환매조건부채권 (Repo) 시장:} 담보(주로 국채)를 이용한 익일물 대출 시장. 모든 금융 기관이 참여하며, 헤어컷을 통해 대출자의 위험을 줄입니다.
        \item \textbf{LIBOR vs. SOFR:} LIBOR는 무담보 은행 간 금리로 조작 논란 후 SOFR(담보부 익일물 자금 조달 금리)로 대체되고 있습니다. SOFR은 실제 거래 기반의 무위험 금리입니다.
    \end{itemize}
    \item \textbf{장기 금리 개념:}
    \begin{itemize}
        \item \textbf{명목 금리 vs. 실질 금리:} 명목 금리는 인플레이션을 고려하지 않은 표면적 금리이며, 실질 금리는 인플레이션을 조정한 실제 구매력 변화를 반영합니다 (피셔 방정식: 명목 금리 $\approx$ 실질 금리 + 예상 인플레이션).
        \item \textbf{물가연동국채 (TIPS):} 원금이 인플레이션에 따라 조정되어 실질 금리를 보장하며, 시장의 미래 인플레이션 기대를 측정하는 중요한 지표입니다.
    \end{itemize}
    \item \textbf{국채와 무위험 금리:}
    \begin{itemize}
        \item \textbf{미국 국채:} 세계 최대, 가장 유동적인 채권 시장. 신용 위험이 없어 무위험 금리의 벤치마크로 사용됩니다. T-Bill, T-Note, T-Bond로 구분됩니다.
        \item \textbf{국채 투자 위험:} 인플레이션과 금리 변화 위험이 있으며, 장기 국채의 명목 금리에는 실질 금리, 예상 인플레이션, 기간 프리미엄이 포함됩니다.
    \end{itemize}
    \item \textbf{금리 기간 구조 (수익률 곡선):}
    \begin{itemize}
        \item 만기별 금리를 비교한 것으로, 경제 상황과 미래 금리 기대를 반영합니다.
        \item \textbf{우상향 곡선:} 경제 확장 및 금리 인상 기대.
        \item \textbf{역전된 곡선:} 경기 침체 및 금리 인하 기대의 강력한 예측 지표.
    \end{itemize}
    \item \textbf{기타 자금 시장:}
    \begin{itemize}
        \item \textbf{머니 마켓 뮤추얼 펀드:} 초안전 단기 자산에 투자하는 펀드로, 은행 예금의 대안이자 단기 자금 시장의 핵심입니다. 정부 보험이 없어 '버크 깨짐' 위험이 있습니다 (예: 리저브 프라이머리 펀드 사례).
        \item \textbf{증권화:} (본 문서에서는 내용 미제공)
    \end{itemize}
\end{itemize}
\end{summarybox}


\section{개요: 증권화 (Securitization)}

\begin{tcolorbox}[
    colback=lightblue,
    colframe=darkblue,
    title=\textbf{문서 핵심 요약}
]
  이 문서는 금융 혁신인 \textbf{증권화 (Securitization)}에 대해 다룹니다.
  증권화는 1970년대에 부상하여 미국 경제에 지대한 영향을 미쳤습니다.
  수많은 대출을 한데 모아 \textbf{위험을 분산}하고 \textbf{현금 흐름을 정확하게 예측}함으로써,
  이 대출들을 담보로 \textbf{증권(채권과 지분)}을 발행하여 자금을 조달하는 방식입니다.
  이를 통해 대출 공급이 확대되었지만, \textbf{단기 자금 시장의 변동성}에 민감하다는 단점도 있습니다.
\end{tcolorbox}


\section{용어 정리 (증권화)}

\begin{tcolorbox}[
    colback=lightgreen,
    colframe=darkgreen,
    title=\textbf{핵심 용어 해설}
]
  다음 표는 증권화 개념을 이해하는 데 필수적인 용어들을 정리한 것입니다.
  각 용어의 쉬운 설명과 원어를 함께 제시하여 이해를 돕습니다.

  \begin{adjustbox}{width=\textwidth,center}
  \begin{tabular}{lllc}
  \toprule
  \textbf{용어 (한글)} & \textbf{쉬운 설명} & \textbf{원어} & \textbf{비고} \\
  \midrule
  \textbf{증권화} & 대출과 같은 비유동성 자산을 모아 유동성 있는 증권(채권, 주식)으로 전환하여 투자자에게 판매하는 과정. & Securitization & 금융 혁신의 핵심 \\
  \textbf{모기지 대출} & 부동산을 담보로 제공하고 빌리는 대출. & Mortgage Loan & 주택 구매 시 주로 이용 \\
  \textbf{채무 불이행} & 대출금을 약속대로 상환하지 못하는 것. & Default & 대출의 주요 위험 요소 \\
  \textbf{대출 풀} & 여러 개의 대출을 한데 모아 놓은 집합. & Loan Pool & 위험 분산의 기본 단위 \\
  \textbf{다각화} & 여러 자산에 분산 투자하여 특정 자산의 위험이 전체 포트폴리오에 미치는 영향을 줄이는 전략. & Diversification & '모든 달걀을 한 바구니에 담지 마라' \\
  \textbf{현금 흐름} & 일정 기간 동안 기업이나 자산에서 발생하고 유출되는 현금의 총량. & Cash Flow & 수익성과 안정성 판단 기준 \\
  \textbf{증권} & 주식, 채권 등 재산적 가치를 나타내는 문서 또는 전자 기록. & Securities & 투자 및 자금 조달 수단 \\
  \textbf{채권} & 정부, 기업 등이 자금을 조달하기 위해 발행하는 차용증서. 정해진 이자를 지급하고 만기에 원금을 상환. & Bonds & 비교적 안전한 투자 \\
  \textbf{지분} & 기업의 소유권을 나타내는 부분. 채권자에게 지급하고 남은 이익에 대한 청구권을 가짐. & Equity & 고위험 고수익 가능성 \\
  \textbf{쿠폰 지급} & 채권 투자자에게 정기적으로 지급되는 이자. & Coupon Payments & 채권 수익의 한 형태 \\
  \textbf{우선 청구권} & 특정 자산이나 현금 흐름에 대해 다른 채권자보다 먼저 권리를 행사할 수 있는 권리. & First Claim/Priority & 채권자가 지분보다 우선 \\
  \textbf{기업어음} & 기업이 단기 자금을 조달하기 위해 발행하는 무담보 어음. & Commercial Paper (CP) & 만기가 짧은 단기 채권 \\
  \textbf{머니마켓 뮤추얼 펀드} & 단기 금융 상품에 투자하여 수익을 추구하는 펀드. & Money Market Mutual Funds (MMMF) & 단기 자금 시장의 주요 참여자 \\
  \textbf{정부 후원 기업} & 정부가 설립하거나 후원하여 특정 공공 목적(예: 주택 금융)을 달성하는 데 도움을 주는 민간 기업. & Government Sponsored Enterprises (GSEs) & Fannie Mae, Freddie Mac 등 \\
  \textbf{미국 국채} & 미국 정부가 발행하는 채권. 가장 안전한 투자 자산 중 하나로 간주됨. & US Treasury Securities & 무위험 자산의 대명사 \\
  \bottomrule
  \end{tabular}
  \end{adjustbox}
\end{tcolorbox}


\section{핵심 개념 및 원리: 증권화의 이해}

\subsection{1. 증권화란 무엇인가?}

증권화는 1970년대에 등장하여 금융 시장에 큰 변화를 가져온 중요한 금융 혁신입니다. 이는 단순히 대출을 묶는 것을 넘어, 대출의 위험과 수익을 분리하고 재구성하여 새로운 투자 상품을 만들어내는 복잡한 과정입니다.

\begin{tcolorbox}[
    colback=lightyellow,
    colframe=darkorange,
    title=\textbf{증권화의 정의와 중요성}
]
  \textbf{증권화 (Securitization)}는 은행이나 다른 금융 기관이 보유한 대출(예: 모기지, 자동차 대출, 학자금 대출 등)과 같은 \textbf{비유동성 자산}들을 한데 모아(풀, Pool), 이 자산들에서 발생하는 미래의 현금 흐름(원금 및 이자 상환)을 담보로 \textbf{유동성 있는 증권(Securities)}을 발행하여 투자자들에게 판매하는 과정입니다.

  이러한 증권화는 대출 기관이 대출을 계속해서 제공할 수 있도록 새로운 자금 조달원을 제공하며, 투자자들에게는 다양한 위험-수익 프로필을 가진 새로운 투자 기회를 제공합니다. 특히 미국 경제에서 신용 공급을 확대하는 데 매우 중요한 역할을 해왔습니다.
\end{tcolorbox}

\subsection{2. 증권화의 작동 원리: 불확실성 감소}

증권화의 핵심은 \textbf{대규모의 대출을 모아 위험을 분산}하고, 이를 통해 미래의 현금 흐름을 \textbf{더욱 정확하게 예측}할 수 있게 만드는 것입니다. 이 원리를 이해하기 위해 대출의 규모에 따른 위험 변화를 살펴보겠습니다.

\subsubsection{2.1. 단일 대출의 높은 불확실성}

하나의 고객에게 하나의 모기지 대출을 제공했다고 가정해 봅시다. 고객의 신용도(위험 프로필)를 분석한 결과, 이 고객이 채무 불이행(Default)할 확률이 2\%라고 예상됩니다.

\begin{tcolorbox}[
    colback=boxred,
    colframe=red!70!black,
    title=\textbf{단일 대출의 위험}
]
  단 하나의 대출과 하나의 고객만 있다면, 결과는 둘 중 하나입니다.
  \begin{itemize}
    \item 고객이 대출을 \textbf{완전히 상환}하거나,
    \item 고객이 대출을 \textbf{완전히 불이행}하거나.
  \end{itemize}
  이처럼 단일 대출의 손실은 \textbf{매우 불확실}합니다. '모 아니면 도' (all or nothing)의 상황이기 때문에, 대출 기관은 큰 위험에 직면하게 됩니다. 2\%의 불이행 확률은 낮아 보이지만, 실제로 불이행이 발생하면 전체 대출금을 잃게 됩니다.
\end{tcolorbox}

\subsubsection{2.2. 대출 풀(Pool)을 통한 위험 분산}

이제 이러한 모기지 대출 100개를 모아 '대출 풀(Loan Pool)'을 만들었다고 가정해 봅시다. 각 고객의 채무 불이행 확률은 여전히 2\%입니다. 중요한 것은 이 대출자들이 모두 다른 지역에 살고 있다고 가정하는 것입니다. 이는 한 지역의 경제적 충격이 모든 대출자에게 동시에 영향을 미치지 않도록 하여, \textbf{동시에 채무 불이행할 가능성이 낮다}는 것을 의미합니다.

\begin{tcolorbox}[
    colback=lightblue,
    colframe=darkblue,
    title=\textbf{대출 풀의 위험 감소 효과}
]
  100개의 대출 풀에서 채무 불이행 위험은 어떻게 될까요?
  \begin{itemize}
    \item 평균적으로 100명 중 2명(2\%)이 대출을 상환하지 않을 것으로 예상합니다.
    \item 하지만 실제로는 1명이 불이행할 수도 있고, 3명이 불이행할 수도 있습니다.
    \item 중요한 점은 \textbf{'모 아니면 도'가 아니라는 것}입니다. 100개의 대출 중 일부가 불이행하더라도, 나머지 대출들은 정상적으로 상환됩니다.
  \end{itemize}
  이처럼 100개의 모기지 대출 풀은 단일 대출 예시보다 \textbf{불확실성이 훨씬 줄어듭니다}.

  더 나아가, 10,000개의 모기지 대출 풀을 생각해 봅시다. 각 대출의 불이행 확률은 여전히 2\%이므로, 평균적으로 200명(2\%)이 불이행할 것으로 예상됩니다. 이 경우, 실제 불이행 건수는 190건 또는 210건이 될 수 있으며, 이는 전체의 1.9\% 또는 2.1\%에 해당합니다. 이 수치는 우리가 예상한 2\%에 \textbf{매우 근접}합니다.
\end{tcolorbox}

\subsubsection{2.3. 핵심 원리: 대규모 분산 풀의 예측 가능성}

위의 예시에서 얻을 수 있는 핵심 교훈은 다음과 같습니다.

\begin{tcolorbox}[
    colback=lightgreen,
    colframe=darkgreen,
    title=\textbf{증권화의 핵심: 현금 흐름의 정확한 예측}
]
  우리는 \textbf{크고 잘 다각화된(well-diversified) 모기지 대출 풀}을 원합니다.
  \begin{itemize}
    \item '크다'는 것은 대출의 수가 많다는 의미입니다.
    \item '잘 다각화되었다'는 것은 대출자들이 서로 다른 지역에 살거나, 다른 산업에 종사하는 등, 특정 요인에 의해 동시에 영향을 받을 가능성이 낮다는 의미입니다.
  \end{itemize}
  이러한 대규모의 잘 다각화된 풀을 통해 우리는 대출에서 발생하는 \textbf{현금 흐름(Cash Flow)을 매우 정확하게 추정}할 수 있습니다. 현금 흐름은 대출의 손실(불이행)과 상환액을 모두 포함합니다.

  대출 풀의 현금 흐름을 정확하게 예측할 수 있다는 점이 바로 \textbf{증권화가 작동하는 핵심 원리}입니다. 예측 가능성이 높을수록 투자자들은 더 안심하고 투자할 수 있게 됩니다.
\end{tcolorbox}

\subsection{3. 증권 발행: 채권과 지분}

이제 증권화의 두 번째 중요한 요소인 \textbf{증권 발행}에 대해 알아보겠습니다. 증권화 과정에서 발행 기관(Issuer)은 대출 풀을 담보로 증권을 판매합니다. 이 증권은 주로 \textbf{채권(Bonds)}과 \textbf{지분(Equity)}으로 구성됩니다. 이 증권들을 판매하여 얻은 자금은 모기지 대출에 필요한 자금을 조달하는 데 사용됩니다.

\begin{tcolorbox}[
    colback=lightyellow,
    colframe=darkorange,
    title=\textbf{증권의 구성과 우선순위}
]
  \begin{itemize}
    \item \textbf{채권 (Bonds):}
      \begin{itemize}
        \item 대출 풀의 현금 흐름을 정확하게 추정할 수 있다면, 상당한 양의 채권을 발행할 수 있습니다.
        \item 이 채권에 대한 이자(쿠폰 지급, Coupon Payments)는 대출 풀에서 발생하는 모기지 이자 상환액에서 나옵니다.
        \item 채권은 \textbf{가장 먼저 현금 흐름에 대한 청구권(First Claim)}을 가집니다. 즉, 대출 풀에서 돈이 들어오면 채권 보유자들에게 먼저 지급됩니다.
        \item 이러한 우선순위 때문에 채권은 \textbf{상대적으로 안전한 투자 상품}으로 간주됩니다.
      \end{itemize}
    \item \textbf{지분 (Equity):}
      \begin{itemize}
        \item 지분 증권 또한 모기지 풀에서 현금 흐름을 받습니다.
        \item 하지만 지분 보유자들은 \textbf{두 번째로 지급}받습니다. 모든 채권 보유자들에게 지급이 완료된 후에 남은 현금 흐름을 지분 보유자들이 가져갑니다.
        \item 이 때문에 지분은 채권보다 \textbf{더 높은 위험}을 가지지만, 동시에 \textbf{더 높은 수익}을 얻을 가능성도 있습니다.
      \end{itemize}
  \end{itemize}
  이러한 구조는 다양한 위험 선호도를 가진 투자자들을 유치할 수 있게 합니다. 안전한 투자를 선호하는 투자자는 채권을, 높은 수익을 추구하는 투자자는 지분을 선택할 수 있습니다.
\end{tcolorbox}

\subsection{4. 증권화의 수익성: 구체적인 예시}

증권화가 왜 수익성이 높은지 구체적인 예를 통해 살펴보겠습니다.

\begin{tcolorbox}[
    colback=lightblue,
    colframe=darkblue,
    title=\textbf{증권화 수익성 예시}
]
  \textbf{가정:}
  \begin{itemize}
    \item 100달러 가치의 모기지 대출 풀이 있습니다.
    \item 이 모기지 대출들은 연 5\%의 이자를 지급합니다.
    \item 우리는 이 풀을 담보로 80달러의 채권을 발행하고, 20달러의 지분을 직접 보유합니다.
    \item 발행된 총 100달러의 증권(채권 80달러 + 지분 20달러)이 100달러의 모기지 대출에 자금을 조달합니다.
  \end{itemize}

  채권 보유자들은 우선적으로 지급받기 때문에 상대적으로 안전합니다. 우리는 채권에 연 4\%의 이자를 약속했다고 가정합시다.

  \textbf{시나리오 1: 채무 불이행이 전혀 없는 경우}
  \begin{itemize}
    \item 모기지 풀은 총 5달러의 현금(100달러의 5\%)을 발생시킵니다.
    \item 80달러의 채권은 4\%의 수익률을 받으므로, 채권 보유자들은 3.20달러 (80달러 $\times$ 4\%)를 받습니다.
    \item 지분 보유자들은 남은 1.80달러 (5달러 - 3.20달러)를 받습니다.
  \end{itemize}
  이 경우, 지분 투자액 20달러 대비 1.80달러의 수익은 9\% (1.80/20)로, 채권 수익률 4\%보다 훨씬 높습니다.

  \textbf{시나리오 2: 모기지 중 2\%가 채무 불이행하는 경우}
  \begin{itemize}
    \item 모기지 풀은 총 4.90달러의 현금 (100달러 중 2달러가 불이행하여 98달러만 상환된다고 가정하면, 98달러의 5%는 4.90달러)을 발생시킵니다.
    \item 채권 보유자들은 우선순위가 있으므로 여전히 3.20달러를 받습니다.
    \item 지분 보유자들은 남은 1.70달러 (4.90달러 - 3.20달러)를 받습니다. 이는 시나리오 1보다 0.10달러 적은 금액입니다.
  \end{itemize}

  \textbf{시나리오 3: 예상치 못한 높은 채무 불이행이 발생하는 경우}
  \begin{itemize}
    \item 만약 채무 불이행이 매우 높아서 모기지 풀에서 발생하는 현금 흐름이 3.20달러 이하로 떨어진다면, 채권 보유자들은 약속된 이자를 받지 못할 수도 있습니다.
    \item 하지만 그 전에, 지분 보유자들은 아무것도 받지 못할 수 있습니다. 지분 보유자들은 일반적으로 투자에 대해 더 높은 수익을 기대하지만, 그만큼 더 큰 위험을 감수합니다.
  \end{itemize}
  이 예시는 증권화의 기본 원리를 보여줍니다. 대출 풀의 현금 흐름을 정확하게 추정할 수 있다면, 채권을 사용하여 대출에 자금을 조달할 수 있습니다. 대출 기관은 모든 대출에 필요한 현금을 직접 보유할 필요 없이 시장에서 자금을 빌릴 수 있게 됩니다.
\end{tcolorbox}


\section{증권화의 중요성과 영향}

\subsection{1. 자금 조달 모델로서의 증권화}

증권화는 20세기 가장 중요한 금융 혁신 중 하나로 평가받습니다. 이는 단순한 투자 상품을 넘어, 금융 기관의 자금 조달 방식 자체를 변화시켰습니다.

\begin{tcolorbox}[
    colback=lightgreen,
    colframe=darkgreen,
    title=\textbf{증권화가 신용 공급에 미치는 영향}
]
  \begin{itemize}
    \item \textbf{대출 기관의 부담 경감:} 전통적인 방식에서는 은행이 대출을 제공하려면 그만큼의 예금이나 자체 자본을 보유해야 했습니다. 하지만 증권화를 통해 대출 기관은 대출을 제공한 후 이를 증권화하여 시장에 판매함으로써, 필요한 자금을 투자자들로부터 조달할 수 있게 됩니다.
    \item \textbf{신용 공급 확대:} 대출 기관이 자체 자본에 묶이지 않고 지속적으로 자금을 조달할 수 있게 되면서, 기업과 가계에 대한 대출 공급 능력이 크게 증가했습니다. 1970년대 이후 미국에서 기업과 가계에 대한 신용 공급이 증권화 시장 덕분에 크게 늘어났습니다.
    \item \textbf{자본 효율성 증대:} 은행은 대출을 증권화하여 장부에서 제거함으로써, 대출에 묶여 있던 자본을 다른 대출이나 투자에 활용할 수 있게 되어 자본 효율성을 높일 수 있습니다.
  \end{itemize}
  결론적으로, 증권화는 금융 시스템 내에서 자금이 더욱 효율적으로 순환하고, 더 많은 사람들이 대출을 받을 수 있도록 하는 중요한 메커니즘이 되었습니다.
\end{tcolorbox}

\subsection{2. 미국 내 증권화 현황 데이터}

미국에서는 가계 및 기업 대출의 상당 부분이 증권화되어 있습니다. 이 현상은 1980년대 이후 극적으로 증가했습니다.

\begin{tcolorbox}[
    colback=lightblue,
    colframe=darkblue,
    title=\textbf{미국 증권화 시장의 규모와 특징}
]
  \begin{itemize}
    \item \textbf{시장 규모:} 2008년 금융 위기 이후 잠시 감소했지만, 2021년까지 증권화 관련 채권의 총액은 약 \textbf{12조 달러}에 달했습니다. 이는 엄청난 규모의 금융 시장을 형성하고 있음을 보여줍니다.
    \item \textbf{주요 발행 주체:}
      \begin{itemize}
        \item 이 채권의 대부분, 즉 \textbf{10조 달러 이상}은 \textbf{정부 후원 기업 (Government Sponsored Enterprises, GSEs)}에 의해 발행되었으며, 미국 정부의 보증을 받습니다.
        \item 정부 보증 덕분에 이 채권들은 \textbf{미국 국채 (US Treasury Securities)만큼 안전}하다고 평가됩니다.
        \item 나머지 부분은 은행 및 금융 회사와 같은 \textbf{민간 발행 기관}의 증권화에서 비롯됩니다.
      \end{itemize}
    \item \textbf{주요 증권화 대상:}
      \begin{itemize}
        \item 증권화된 대출의 \textbf{대부분은 모기지}입니다. 여기에는 단독 주택을 위한 주거용 모기지, 다세대 주택 및 아파트 건물을 위한 모기지, 사무실 건물이나 쇼핑몰을 위한 상업용 모기지가 포함됩니다.
        \item 다른 대출 범주에는 \textbf{자동차 대출}과 \textbf{소기업 대출} 등이 있습니다.
      \end{itemize}
  \end{itemize}
  이 데이터는 증권화가 미국 금융 시스템의 핵심적인 부분이며, 특히 주택 금융 시장에서 절대적인 영향력을 가지고 있음을 명확히 보여줍니다.
\end{tcolorbox}

\subsection{3. 단기 자금 시장과의 민감한 관계}

증권화는 매우 효율적인 자금 조달 모델이지만, \textbf{단기 자금 시장 (Short-term Funding Markets)}의 상황에 민감하게 반응할 수 있다는 중요한 단점이 있습니다.

\begin{tcolorbox}[
    colback=boxred,
    colframe=red!70!black,
    title=\textbf{단기 자금 시장과 증권화의 연결고리}
]
  \textbf{왜 증권화가 단기 자금 시장에 민감할까요?}
  \begin{itemize}
    \item \textbf{임시 자금 조달:} 많은 발행 기관들은 대출을 풀로 모아 증권화하기 전에, \textbf{단기적으로 자금을 조달}하여 대출을 제공합니다.
    \item \textbf{기업어음 (Commercial Paper) 활용:} 예를 들어, 이들은 머니마켓 뮤추얼 펀드(Money Market Mutual Funds)에 기업어음(Commercial Paper)을 판매하여 현금을 조달합니다. 이 현금으로 대출을 제공합니다.
    \item \textbf{상환 과정:} 충분한 대출이 모여 풀이 형성되고 증권화되면, 발행 기관은 채권 투자자들로부터 조달한 자금으로 이전에 발행했던 기업어음을 상환합니다.
  \end{itemize}
  이러한 과정은 단기 자금 시장이 원활하게 작동할 때만 가능합니다.

  \textbf{단기 자금 시장의 교란이 경제에 미치는 영향:}
  \begin{itemize}
    \item 만약 단기 자금 시장이 \textbf{경색(tries up)}되면, 즉 기업어음과 같은 단기 자금 조달이 어려워지면, 증권화 모델은 작동을 멈추게 됩니다.
    \item 2008년 금융 위기 때처럼 이러한 상황이 발생하면, 기업과 가계는 대출을 받기 어려워져 \textbf{신용 경색 (Credit Crunch)}이 발생하고 경제 전반에 심각한 영향을 미칠 수 있습니다.
  \end{itemize}
  따라서 증권화는 신용 공급을 확대하는 강력한 도구이지만, 동시에 금융 시스템의 취약성을 증가시킬 수 있는 양날의 검과 같습니다.
\end{tcolorbox}


\section{세션 요약 (증권화)}

이번 세션에서는 금융 시장의 중요한 혁신인 증권화에 대해 깊이 있게 학습했습니다. 핵심 내용은 다음과 같습니다.

\begin{tcolorbox}[
    colback=lightgreen,
    colframe=darkgreen,
    title=\textbf{학습 내용 요약}
]
  \begin{itemize}
    \item \textbf{증권화의 정의와 작동 원리:}
      \begin{itemize}
        \item 증권화는 20세기에 등장한 중요한 자금 조달 모델입니다.
        \item 이는 \textbf{통계 모델}에 기반하며, 많은 수의 대출을 모아 \textbf{불확실성을 줄이는 것}이 핵심입니다.
        \item 대출 풀을 형성하고, 이 풀을 담보로 \textbf{증권(채권과 지분)}을 발행하는 과정을 포함합니다.
      \end{itemize}
    \item \textbf{증권화의 긍정적 영향:}
      \begin{itemize}
        \item 증권화 자금 조달 모델은 미국에서 \textbf{총 신용 공급을 증가}시키는 데 크게 기여했습니다.
        \item 이를 통해 기업과 가계가 더 쉽게 대출을 받을 수 있게 되었습니다.
      \end{itemize}
    \item \textbf{증권화의 잠재적 단점:}
      \begin{itemize}
        \item 증권화는 \textbf{단기 자금 시장의 상황에 민감}할 수 있습니다.
        \item 단기 자금 시장의 교란은 증권화 모델의 붕괴로 이어져, 신용 경색을 유발하고 경제에 부정적인 영향을 미칠 수 있습니다.
      \end{itemize}
  \end{itemize}
  증권화는 금융 시스템의 효율성을 높이는 동시에, 잠재적인 위험 요소를 내포하고 있음을 이해하는 것이 중요합니다.
\end{tcolorbox}


%=======================================================================
% Module 6: 용어 정리
%=======================================================================
\chapter{용어 정리}
\label{ch:module6}

\metainfo{중앙은행과 통화정책}{모듈 2 (파트 1/2)}{Prof. Ralf Meisenzahl}{미국 연방준비제도(Fed)의 역사, 구조, 기능 및 통화정책 도구를 심층적으로 학습하고, 연준의 이중 책무(최대 고용, 물가 안정)를 중심으로 전통적 및 비전통적 통화정책 수단들을 이해한다.}


% 본문 시작
이 문서는 '중앙은행과 통화정책' 강의의 모듈 2 중 첫 번째 파트 내용을 담고 있습니다.
미국 연방준비제도(Federal Reserve System)의 역사, 구조, 기능, 그리고 통화정책 도구들을 심층적으로 다룹니다.
특히, 연준의 이중 책무(Dual Mandate)인 최대 고용과 물가 안정을 중심으로 전통적 및 비전통적 통화정책 수단들을 상세히 살펴봅니다.
화폐 공급의 개념과 측정, 그리고 연준의 대차대조표 변화를 통해 통화정책이 경제에 미치는 영향을 이해하는 것이 목표입니다.
이 문서는 비전공자나 처음 배우는 사람도 완벽히 이해할 수 있도록 쉬운 설명과 풍부한 예시를 제공합니다.

\vspace{1em}

\begin{summarybox}
\textbf{문서 전체 핵심 요약:}
본 문서는 미국 연방준비제도(Fed)의 역사, 구조, 주요 기능 및 통화정책 도구를 다룹니다.
연준의 이중 책무(최대 고용, 물가 안정)를 중심으로 전통적 통화정책(공개시장운영, 재할인 창구)과 비전통적 통화정책(선제 안내, 대규모 자산 매입)을 설명합니다.
화폐 공급의 개념과 측정, 그리고 연준 대차대조표의 변화를 통해 통화정책의 작동 원리를 이해하는 데 중점을 둡니다.
\end{summarybox}

\section{용어 정리}

다음은 본 문서에서 다루는 주요 용어들을 비전공자도 이해하기 쉽게 정리한 표입니다.

\begin{adjustbox}{width=\textwidth,center}
\begin{tabular}{lllc}
\toprule
\textbf{용어 (Term)} & \textbf{쉬운 설명 (Easy Explanation)} & \textbf{원어 (Original Term)} & \textbf{비고 (Notes)} \\
\midrule
중앙은행 & 한 국가의 금융 시스템을 관리하고 통화 정책을 수행하는 기관 & Central Bank & 스웨덴 릭스방크, 미국 연방준비제도 등 \\
연방준비제도 & 미국의 중앙은행 시스템 & Federal Reserve System (Fed) & 줄여서 '연준'이라고도 함 \\
이중 책무 & 연준의 두 가지 주요 목표: 최대 고용과 물가 안정 & Dual Mandate & 통화 정책의 핵심 목표 \\
최대 고용 & 일자리를 찾는 대부분의 사람들이 일하고 있는 상태 & Maximum Employment & 실업률이 낮은 상태를 의미 \\
물가 안정 & 물가 상승률(인플레이션)이 낮고 안정적인 상태 & Stable Prices & 저축과 투자를 장려 \\
통화 정책 & 중앙은행이 경제의 화폐 공급과 신용 조건을 조절하는 활동 & Monetary Policy & 금리 조절, 화폐 공급 조절 등 \\
지급결제 시스템 & 돈이 한 사람/기관에서 다른 사람/기관으로 안전하게 이동하는 방식 & Payment System & 수표, ACH, 대규모 거래 등 \\
화폐 공급 & 경제 내에 유통되는 돈의 총량 & Money Supply & 통화량이라고도 함 \\
통화 승수 & 은행의 대출 활동을 통해 화폐 공급이 얼마나 증가하는지 나타내는 비율 & Money Multiplier & 예금의 일부를 대출하여 새로운 예금을 창출 \\
본원 통화 & 중앙은행이 직접 발행하는 화폐 (현금 + 은행의 중앙은행 예치금) & Monetary Base & 가장 좁은 의미의 화폐 공급 \\
M1 & 본원 통화에 요구불 예금(수표 계좌 잔액)을 더한 화폐 공급 지표 & M1 & 유동성이 높은 자산 포함 \\
M2 & M1에 저축성 예금, 소액 정기 예금 등을 더한 화폐 공급 지표 & M2 & 가장 일반적으로 사용되는 지표 \\
M3 & M2에 대규모 정기 예금, 기관 머니마켓 펀드 등을 더한 화폐 공급 지표 & M3 & 현재 연준에서 발표하지 않음 \\
양적 화폐 이론 & 화폐 공급량과 물가 수준, 경제 성장 간의 관계를 설명하는 이론 & Quantity Theory of Money & 화폐 공급 증가가 인플레이션을 유발할 수 있음 \\
공개시장운영 & 중앙은행이 국채 등을 사고팔아 화폐 공급과 금리를 조절하는 정책 & Open Market Operations (OMO) & 전통적인 통화 정책 수단 \\
연방기금 금리 & 은행들이 서로 하룻밤 동안 돈을 빌려줄 때 적용되는 금리 & Federal Funds Rate & 연준의 주요 목표 금리 \\
재할인 창구 & 중앙은행이 은행에 단기 대출을 제공하는 제도 & Discount Window & 비상시 유동성 공급 수단 \\
선제 안내 & 중앙은행이 미래의 통화 정책 방향에 대해 시장에 미리 알려주는 것 & Forward Guidance & 장기 금리 목표에 사용 \\
대규모 자산 매입 & 중앙은행이 국채나 주택저당증권 등을 대규모로 매입하는 정책 & Large-Scale Asset Purchases (LSAP) & 장기 금리를 직접 낮추는 효과 \\
오버나이트 역환매 조건부 & 중앙은행이 금융기관에 단기 증권을 팔고 다음날 다시 사는 거래 & Overnight Reverse Repurchase Program (ON RRP) & 단기 금리 조절에 사용 \\
대차대조표 & 특정 시점의 자산, 부채, 자본 상태를 보여주는 재무 보고서 & Balance Sheet & 연준의 자산과 부채 현황 \\
자산 & 경제적 가치가 있는 것 (예: 연준이 보유한 국채) & Assets & 미래 경제적 효익을 제공 \\
부채 & 갚아야 할 의무 (예: 시중에 유통되는 화폐, 은행의 중앙은행 예치금) & Liabilities & 미래 경제적 효익을 포기 \\
\bottomrule
\end{tabular}
\end{adjustbox}


\section{모듈 2: 연방준비제도 개요}

안녕하세요, 여러분. 어디에서 참여하시든 잘 지내시길 바랍니다.

이번 모듈에서는 미국 연방준비제도(Federal Reserve System), 즉 연준의 기능과 통화정책 도구들을 심층적으로 학습합니다. 우리는 연준의 역사부터 시작하여, 왜 중앙은행이 필요한지, 어떻게 조직되어 있는지, 그리고 어떤 목적을 수행하는지 살펴볼 것입니다.

\begin{conceptbox}
\textbf{연준의 이중 책무 (Dual Mandate)}
연준의 통화정책 목표는 크게 두 가지입니다: \textbf{최대 고용(Maximum Employment)}과 \textbf{물가 안정(Stable Prices)}. 이 두 가지는 종종 연준의 '이중 책무'라고 불리며, 이 강의의 모든 논의는 이 이중 책무와 연결됩니다.
\end{conceptbox}

우리는 연준의 제도적 특징과 통화정책 의사결정 과정이 어떻게 작동하는지 면밀히 검토할 것입니다. 명확히 하자면, 이 강의는 미국의 연방준비제도 시스템을 사례 연구로 다룹니다. 물론 전 세계에는 많은 다른 중앙은행들이 존재하며, 일부 중앙은행은 다른 책무를 가지거나 약간 다른 핵심 기능을 수행하기도 합니다. 그러나 모든 중앙은행은 통화정책을 수행하고 지급결제 시스템을 제공합니다. 이 강의에서는 주로 중앙은행이 통화정책을 어떻게 수행하는지에 초점을 맞춥니다.

\vspace{1em}

\begin{summarybox}
\textbf{통화정책의 진화: 화폐 공급 vs. 금리}
시간이 지남에 따라 통화정책을 수행하는 최선의 방법에 대한 많은 이견이 있었습니다. 일부 중앙은행은 금리보다는 화폐 공급을 명시적으로 목표로 삼기도 했습니다. 연준도 과거에는 화폐 공급을 목표로 삼았습니다.
\end{summarybox}

\textbf{왜 화폐 공급에 관심을 가져야 할까요?}
다른 상품과 마찬가지로, 돈의 가격은 경제 내의 돈에 대한 수요와 공급에 의해 결정됩니다. 아마도 특이한 점은 돈의 가격이 바로 금리라는 것입니다. 금리를 직접 목표로 삼기보다는 화폐 공급을 통제하는 것이 더 합리적일 수 있습니다. 하지만 이는 화폐 공급을 어떻게 측정할 수 있는가라는 질문을 던집니다. 따라서 우리는 화폐 공급의 개념, 그 측정 방법, 그리고 통화정책에 의해 어떻게 영향을 받을 수 있는지 소개할 것입니다.

그 다음으로, 우리는 연준의 대차대조표와 그것이 시간이 지남에 따라 어떻게 변화해왔는지 살펴볼 것입니다. 통화정책 수행이 화폐 공급을 어떻게 변화시켰는지에 대한 통찰력을 제공할 것입니다.

\vspace{1em}

\begin{conceptbox}
\textbf{전통적 통화정책 (Traditional Monetary Policy)}
화폐 공급과 금리가 정확히 어떻게 관련되어 있는지 이해하기 위해, 우리는 전통적인 통화정책에 초점을 맞출 것입니다.
\begin{itemize}
    \item \textbf{공개시장운영 (Open Market Operations):} 공개시장운영이 연방기금 시장(federal funds market)에서 화폐 공급에 직접적으로 어떻게 영향을 미치는지 논의할 것입니다. 이러한 공개시장운영이 이 시장의 균형 금리인 \textbf{연방기금 금리(Federal funds rate)}를 어떻게 결정하는지 살펴볼 것입니다. 연방기금 금리는 미국 통화정책의 주요 목표 금리입니다.
    \item \textbf{재할인 창구 (Discount Window):} 보완적인 전통적 통화정책 도구는 재할인 창구입니다. 이 정책 도구는 은행에 단기 대출 형태로 단기 유동성을 제공합니다. 이 도구는 주로 비상용으로 설계되었기 때문에 금리는 연방기금 금리보다 높게 설정됩니다. 그러나 재할인 창구 대출은 연방기금의 대안으로 남아 있습니다. 따라서 재할인 창구 대출의 금리는 연준이 설정한 연방기금 금리 목표 범위의 상한선을 강제합니다.
\end{itemize}
\end{conceptbox}

\vspace{1em}

\begin{conceptbox}
\textbf{비전통적 통화정책 (Non-traditional Monetary Policy)}
이 모듈의 마지막 부분에서는 2008년 글로벌 금융 위기 이후 등장한 새로운 통화정책 도구들을 살펴봅니다. 이들은 종종 '비전통적 통화정책 도구'라고 불립니다.
\begin{itemize}
    \item \textbf{선제 안내 (Forward Guidance):} 선제 안내는 시장 참여자들에게 미래의 금리 또는 더 일반적으로 통화정책 조치에 대한 명확한 경로를 제공하며, 장기 금리를 목표로 하는 데 사용될 수 있습니다.
    \item \textbf{대규모 자산 매입 (Large-Scale Asset Purchases):} 단기 금리가 이미 0일 때, 중앙은행은 국채나 주택저당증권 등을 대규모로 매입하여 화폐 공급을 더욱 늘릴 수 있습니다. 이 접근 방식은 중앙은행이 장기 금리를 직접 낮출 수 있도록 합니다.
    \item \textbf{오버나이트 역환매 조건부 프로그램 (Overnight Reverse Repurchase Program, ON RRP):} 이 도구는 환매 시장(repo market)에서 금리를 설정하는 데 사용됩니다. 환매 시장은 연방기금 시장과 유사한 오버나이트 대출 시장입니다. 금융 기관들은 이 시장들 사이를 전환할 수 있으며, 이들은 대체재입니다. 환매 시장의 균형을 관리함으로써 연준은 모든 단기 금리가 통화정책 목표와 일치하도록 보장합니다.
\end{itemize}
\end{conceptbox}

결론적으로, 이 모듈은 통화정책 도구에 대한 심층적인 논의를 제공합니다. 이 비디오들 중 다수는 첫 번째 모듈에서 논의된 금융 시장 및 상품에 대한 이해를 필요로 할 것입니다. 특히, 연방기금 시장, 환매 시장, 국채, 그리고 수익률 곡선에 대한 이해가 중요합니다.

\begin{summarybox}
\textbf{핵심 요약: 통화정책 도구의 두 가지 주요 관점}
모든 통화정책 도구의 기술적인 세부 사항을 모두 이해할 필요는 없습니다. 그러나 다음 두 가지 주요 고수준 포인트를 명심하시기 바랍니다.
\begin{enumerate}
    \item \textbf{다중 시장 목표:} 연방준비제도는 금리 정책을 통해 여러 밀접하게 연결된 시장을 목표로 합니다.
    \item \textbf{제로 금리 하한선 극복:} 단기 금리가 0일 때도 중앙은행은 선택지가 없는 것이 아닙니다. 선제 안내와 대규모 자산 매입을 사용하여 장기 금리를 목표로 할 수 있습니다.
\end{enumerate}
\end{summarybox}


\subsection{레슨 2-1.1. 연방준비제도의 역사와 구조}

안녕하세요, 연방준비제도 시스템의 역사와 구조에 대한 강의에 오신 것을 환영합니다. 이 수업에서는 중앙은행의 개념을 소개할 것입니다. 중앙은행은 전 세계에 존재합니다. 우리는 중앙은행이 왜 설립되었는지, 그들의 목표와 주요 기능은 무엇인지 논의할 것입니다. 그런 다음 미국의 중앙은행인 연방준비제도 시스템에 초점을 맞출 것입니다. 우리는 그 설립을 논의하고, 연방준비제도 시스템의 구조와 의사결정 기구를 검토할 것입니다.

\vspace{1em}

\subsubsection*{세계에서 가장 오래된 중앙은행: 스웨덴 릭스방크 사례}
세계에서 가장 오래된 중앙은행인 스웨덴 릭스방크(Swedish Riksbank)부터 살펴보겠습니다. 스웨덴 스톡홀름에 기반을 둔 이 은행은 국가들이 왜 중앙은행을 설립했는지, 그리고 무엇을 달성하고자 했는지 이해하는 데 도움이 되는 흥미로운 사례 연구입니다.

\begin{examplebox}
\textbf{스웨덴 릭스방크의 설립 배경 (1661년)}
\begin{itemize}
    \item \textbf{초기 화폐:} 1661년 스웨덴에서는 귀금속으로 만든 무거운 동전이 화폐로 사용되었습니다. 이 동전들은 무게가 45파운드(약 20kg)에 달하여 매우 불편했습니다.
    \item \textbf{지폐의 등장:} 1661년, 스웨덴 국왕은 민간 은행들이 종이 은행권을 발행하도록 허용했습니다. 이것이 오늘날 우리가 아는 지폐의 시작입니다. 지폐는 은행과 시민들에게 훨씬 더 편리했습니다.
    \item \textbf{지폐의 신뢰성:} 사람들은 왜 이 은행권을 기꺼이 받아들였을까요? 은행이 예금된 동전과 은행권을 교환해주겠다고 약속했기 때문입니다. 은행권은 요구불 상환(redeemable on demand)이 가능했습니다.
    \item \textbf{은행권 남발과 뱅크런:} 은행권은 큰 성공을 거두었지만, 은행은 점점 더 많은 은행권을 발행했습니다. 이로 인해 사람들은 민간 은행이 모든 은행권을 교환해줄 만큼 충분한 동전을 금고에 보관하고 있는지 의문을 품기 시작했습니다. 이러한 의심은 은행권을 동전으로 상환하려는 사람들이 늘어나게 만들었습니다. 이것이 바로 \textbf{뱅크런(bank run)}의 한 예시입니다. 결국 은행은 모든 은행권을 한 번에 상환할 만큼 충분한 동전을 가지고 있지 않았고, 그 결과 은행은 파산했습니다.
\end{itemize}
\end{examplebox}

\vspace{1em}

이 뱅크런에도 불구하고, 지폐의 아이디어는 여전히 인기가 있었습니다. 그러나 예금된 동전에 비해 너무 많은 지폐가 유통되고 있다는 점이 문제로 인식되었습니다. 즉, 은행 시스템이 보유한 담보(collateral)에 비해 너무 많은 돈이 유통되고 있었던 것입니다.

\begin{conceptbox}
\textbf{스웨덴 릭스방크의 설립 (1668년) 및 중앙은행의 두 가지 주요 역할}
이러한 문제 해결을 위해 1668년 스웨덴 릭스방크가 설립되어 화폐 공급을 통제하게 되었습니다. 법은 이 새로운 중앙은행이 국내 주화의 가치를 올바르고 공정하게 유지하도록 요구했습니다. 또한 법은 릭스방크가 인플레이션을 유발할 정도로 많은 지폐를 발행하는 것을 금지했습니다.

스웨덴 릭스방크의 설립 이야기는 현대 중앙은행이 해결하도록 설계된 두 가지 중요한 문제를 보여줍니다.
\begin{enumerate}
    \item \textbf{물가 안정 보장 (Ensuring Price Stability):} 중앙은행은 일반적으로 화폐 공급에 대한 통제권을 부여받습니다. 이 통제권은 중앙은행이 너무 많은 지폐가 인쇄되는 것을 막을 수 있도록 합니다. 경제에 돈이 너무 많으면 물가 인플레이션을 유발하는 경향이 있습니다.
    \item \textbf{최후의 대부자 역할 (Lender of Last Resort):} 은행은 두려움에 찬 고객들의 뱅크런에 취약합니다. 뱅크런에 대처하기 위해 중앙은행은 '최후의 대부자' 역할을 합니다.
\end{enumerate}
\end{conceptbox}

\textbf{최후의 대부자 역할이란 무엇일까요?}
은행은 단기로 차입하여 장기로 대출합니다. 스웨덴 사례에서 은행은 은행권과 교환하여 금속 동전을 차입했습니다. 그리고 그 동전의 일부를 대출하는 데 사용했습니다. 어느 시점에서든 금고에 있는 동전은 발행된 은행권보다 적었습니다. 모든 사람이 은행권을 동전으로 바꾸고 싶어 할 때, 은행은 동전이 부족하여 파산할 수 있었습니다. 현대 은행업에서는 은행이 예금자로부터 현금을 차입하여 장기로 대출합니다. 이들도 수익성 있는 대출을 좋은 차입자에게 하고 있더라도 때때로 현금 부족에 직면할 수 있습니다.

\begin{examplebox}
\textbf{최후의 대부자 기능의 실제 작동 방식}
중앙은행의 최후의 대부자 기능은 건강한 은행이 현금 부족으로 인해 파산하는 것을 막기 위해 설계되었습니다. 이것이 실제로 어떻게 작동할까요?
\begin{itemize}
    \item \textbf{단기 대출 제공:} 중앙은행은 담보(collateral)를 받고 은행에 단기 대출 형태로 현금을 제공합니다.
    \item \textbf{비상시 사용:} 이는 은행이 스트레스를 받는 비상시에 발생할 수 있습니다.
    \item \textbf{다양한 담보 인정:} 중앙은행은 일반적으로 다양한 유형의 담보를 받습니다. 중요한 것은, 이는 유동성이 낮은 장기 대출, 즉 판매하기 어렵거나 큰 할인 없이는 판매할 수 없는 자산일 수 있다는 점입니다. 예를 들어, 오늘날 연방준비제도 시스템은 부동산 및 소비자 대출을 은행에 대한 현금 대출의 담보로 인정합니다.
\end{itemize}
\end{examplebox}

\subsubsection*{미국 연방준비제도의 설립}
이제 연방준비제도 시스템이 왜 설립되었는지 살펴보겠습니다. 연방준비제도 시스템은 미국 최초의 중앙은행이 아니며, 실제로는 세 번째입니다. 1791년부터 1811년까지 존재했던 제1차 미국 은행(First Bank of America)과 1816년부터 1836년까지 존재했던 제2차 미국 은행(Second Bank of America)이 있었습니다. 그러나 이 두 초기 중앙은행은 오늘날 우리가 중앙은행이라고 생각하는 것과는 매우 달랐습니다. 예를 들어, 이들은 화폐 공급에 대한 통제권이 없었습니다.

미국은 1837년부터 1913년까지 중앙은행이 없었습니다. 스웨덴 사례와 마찬가지로 돈은 민간에서 발행되었습니다. 이 기간 동안 1907년의 공황(panic)을 포함하여 여러 차례의 은행 위기가 발생했습니다. 1907년 공황 동안 주식 시장은 가치의 50\%를 잃었습니다. 전국적으로 은행들은 고객들이 예금을 인출하기 위해 줄을 서는 것을 목격했습니다.

\begin{conceptbox}
\textbf{1907년 공황과 연방준비제도법 (Federal Reserve Act)}
1907년 공황은 최후의 대부자의 필요성을 명확히 보여주었습니다. 그러나 연방준비제도 시스템이 공식적으로 설립되기까지는 시간이 걸렸고 많은 정치적 논쟁이 있었습니다. 1913년 12월 24일, 우드로 윌슨 대통령은 연방준비제도법에 서명했습니다.

\textbf{연방준비제도법이 연준에 허용한 것:}
\begin{enumerate}
    \item \textbf{화폐 공급 통제권 부여:} 연준은 화폐 공급에 대한 통제권을 갖게 되었습니다.
    \item \textbf{은행에 대한 최후의 대부자 역할 수행:} 연준은 은행에 대한 최후의 대부자 역할을 할 수 있게 되었습니다. 이 최후의 대부자 기능은 2008년 금융 위기와 COVID-19 팬데믹을 포함하여 시장 전반의 스트레스 시기에 개입하는 데 사용되었습니다.
\end{enumerate}
\end{conceptbox}

\subsubsection*{연방준비제도의 구조}
연방준비제도법은 새로운 중앙은행의 세 가지 부분을 설립했습니다.

\begin{enumerate}
    \item \textbf{이사회 (Board of Governors):} 워싱턴 D.C.에 위치한 독립적인 연방 기관입니다. 이사회는 7명의 이사로 구성됩니다. 이사회는 전체 연방준비제도 시스템의 정책을 안내합니다.
    \item \textbf{12개 지역 연방준비은행 (Regional Federal Reserve Banks):}
        \begin{itemize}
            \item \textbf{위치:} 보스턴, 뉴욕, 필라델피아, 클리블랜드, 리치먼드, 애틀랜타, 시카고, 세인트루이스, 미니애폴리스, 캔자스시티, 댈러스, 샌프란시스코에 위치합니다.
            \item \textbf{설립 목적:} 지역 연방준비은행들은 중요한 금융 중심지에 배치되어 연준의 권한을 미국 전역에 분산시킵니다.
            \item \textbf{주요 기능:} 이들은 연방준비제도의 정책과 기능을 실행합니다. 예를 들어, 금융 기관을 검사하고 감독하며, 해당 지역 은행에 대한 최후의 대부자 역할을 합니다. 또한 해당 지역 은행에 미국 지급결제 시스템 서비스를 제공합니다.
        \end{itemize}
    \item \textbf{연방공개시장위원회 (Federal Open Market Committee, FOMC):}
        \begin{itemize}
            \item \textbf{구성:} FOMC는 12명의 위원으로 구성됩니다. 항상 이사회 7명의 이사와 뉴욕 연방준비은행 총재가 포함됩니다. 나머지 4명의 위원은 나머지 연방준비은행 총재들의 순환 그룹에서 나옵니다.
            \item \textbf{주요 임무:} FOMC의 주요 임무는 금리 목표를 설정하는 것입니다. 이는 일반적으로 통화정책이라고 불립니다.
        \end{itemize}
\end{enumerate}

\begin{summarybox}
\textbf{레슨 2-1.1 핵심 요약}
\begin{itemize}
    \item 중앙은행이 존재하기 전에는 지폐가 종종 민간에서 발행되었습니다. 이 시스템에서는 화폐 공급이 예측 불가능했고 뱅크런이 자주 발생했습니다.
    \item 화폐 공급을 통제하고 은행 시스템의 안정성을 보장하기 위해 중앙은행이 설립되었습니다.
    \item 중앙은행은 위기 시 은행에 현금을 대출해주는 최후의 대부자 역할도 합니다.
    \item 미국에서는 연방준비제도 시스템이 이사회, 12개 연방준비은행, 연방공개시장위원회라는 세 가지 기관을 통해 이러한 기능을 수행합니다.
\end{itemize}
\end{summarybox}


\subsection{레슨 2-1.2. 연방준비제도의 목적과 기능}

안녕하세요, 연방준비제도 시스템의 다섯 가지 핵심 기능에 대한 강의에 오신 것을 환영합니다. 이 수업에서는 연방준비제도 시스템의 책무에 대해 논의할 것입니다. 법에 따라 연준은 미국 경제의 효과적인 운영을 촉진하기 위해 존재합니다. 이것이 실제로 무엇을 의미하는지, 그리고 연준의 다섯 가지 핵심 기능이 무엇인지 살펴볼 것입니다.

\begin{conceptbox}
\textbf{연준의 다섯 가지 핵심 기능}
\begin{enumerate}
    \item 통화정책 수행 (Conduct monetary policy)
    \item 금융 시스템 안정성 증진 (Promote financial system stability)
    \item 금융 기관 규제 및 감독 (Regulate and supervise financial institutions)
    \item 효율적인 지급결제 시스템 제공 (Provide an efficient payment system)
    \item 소비자 보호 및 지역 사회 개발 증진 (Promote consumer protection and community development)
\end{enumerate}
우리는 이 기능들을 하나씩 논의하고, 이들이 연준의 목표와 어떻게 연결되는지 살펴볼 것입니다.
\end{conceptbox}

\subsubsection*{1. 통화정책 수행 (Conduct Monetary Policy)}
1913년 연방준비제도법은 연방준비제도 시스템을 창설했습니다. 이 법에 따라 연준은 생산적이고 안정적인 경제를 육성하기 위해 통화정책을 수행해야 합니다. 통화정책은 시간이 지남에 따라 진화했으며, 오늘날에는 단기 금리 통제를 포함합니다.

\begin{conceptbox}
\textbf{연준의 법적 통화정책 목표 (이중 책무 + 장기 금리)}
법에 따라 연준은 통화정책을 사용하여 다음 목표를 효과적으로 촉진해야 합니다.
\begin{itemize}
    \item \textbf{최대 고용 (Maximum Employment):} 일반적으로 일자리를 찾는 대부분의 사람들이 유급 고용되어 있는 상황을 의미합니다.
    \item \textbf{물가 안정 (Stable Prices):} 물가가 안정적일 때 저축과 투자가 장려됩니다. 반대로 인플레이션이 높으면 자산 수익률이 낮아질 수 있으며, 이는 가계의 저축 유인과 기업의 투자 유인을 감소시킵니다.
    \item \textbf{온건한 장기 금리 (Moderate Long-Term Interest Rates):} 온건한 장기 금리는 기업의 투자 유인을 증가시킵니다.
\end{itemize}
\textbf{목표 달성:} 최대 고용, 물가 안정, 합리적인 장기 금리 하에서 안정적인 성장과 번성하는 경제가 있을 것이라는 믿음이 있습니다. 이것이 연준이 달성하고자 하는 바입니다.
\end{conceptbox}

\subsubsection*{2. 금융 시스템 안정성 증진 (Promote Financial System Stability)}
금융 안정성이란 무엇일까요? 일반적으로 우리는 금융 시스템이 가계와 기업에 합리적인 비용으로 중단 없이 서비스를 제공할 때 안정적이라고 간주합니다. 금융 시스템은 금융 시장과 은행과 같은 금융 기관으로 구성됩니다. 안정적인 시스템은 저축자로부터 차입자에게 자금을 전달하며, 이는 경제 전반에 걸쳐 자본의 효율적인 배분을 가능하게 합니다.

\begin{conceptbox}
\textbf{금융 시스템 모니터링 접근 방식}
연준은 이 책임을 다하기 위해 금융 시스템을 모니터링합니다. 실제로 두 가지 접근 방식을 사용하여 모니터링합니다.
\begin{enumerate}
    \item \textbf{미시 건전성 접근 방식 (Micro-prudential Approach):} 특정 기관(은행 또는 보험 회사 등)의 위험에 초점을 맞춥니다. 이 미시 건전성 접근 방식은 은행 감독과 관련이 있으며, 이에 대해서는 잠시 후에 논의할 것입니다.
    \item \textbf{거시 건전성 접근 방식 (Macro-prudential Approach):} 대신 전체 금융 시스템에 영향을 미칠 수 있는 위험을 찾습니다. 연방준비제도는 여러 금융 시장 스트레스 측정 지표를 모니터링합니다.
\end{enumerate}
\end{conceptbox}

\begin{examplebox}
\textbf{거시 건전성 모니터링 예시}
\begin{itemize}
    \item \textbf{레버리지 (Leverage):} 금융 시스템 내의 차입 또는 레버리지의 양은 손실을 증폭시키고 채무 불이행을 유발할 수 있습니다.
    \item \textbf{단기 차입 (Short-Term Borrowing):} 금융 기업의 단기 차입량은 신용 흐름을 방해하는 뱅크런을 유발할 수 있습니다.
\end{itemize}
\end{examplebox}

\subsubsection*{3. 금융 기관 규제 및 감독 (Regulate and Supervise Financial Institutions)}
연방준비제도의 세 번째 핵심 기능은 금융 기관을 규제하고 감독하는 것입니다. 규제와 감독은 별개이지만 상호 보완적인 활동입니다.

\begin{conceptbox}
\textbf{규제 (Regulation)와 감독 (Supervision)의 차이}
\begin{itemize}
    \item \textbf{규제 (Regulation):} 규칙을 작성하는 것과 관련이 있습니다. 이러한 규칙은 종종 금융 기관에 제한을 가합니다. 이 과정은 의회의 입법으로 시작되며, 이는 연방준비제도에 규칙을 초안할 권한을 부여합니다. 대중의 의견을 수렴하여 규칙은 최종 확정되고 공표되며, 금융 기관은 이 규칙을 준수해야 합니다.
    \item \textbf{감독 (Supervision):} 반면에 규제 준수와 관련이 있습니다. 이를 위해 연방준비제도는 은행 검사관을 고용하고 훈련시킵니다. 검사관은 규제 대상 금융 기관을 방문하여 규칙 준수 여부를 확인하며, 이는 은행 운영에 대한 심층적인 검사를 포함합니다. 규칙이 위반될 경우, 감독관은 벌금 및 배당금 지급 제한과 같은 다른 제재를 부과할 수 있습니다.
\end{itemize}
\end{conceptbox}

\begin{examplebox}
\textbf{규제 예시: 진실한 대출법 (Truth in Lending Act)}
진실한 대출법은 좋은 예시입니다. 이 규칙에 따라 연방준비제도는 대출 기관이 대출 조건과 비용을 차입자에게 명확하게 공개하도록 요구합니다. 여기에는 신용 카드 이자율이 포함됩니다.

\textbf{감독의 중요성:}
감독은 금융 안정성과 관련이 있습니다. 이는 개별 금융 기관의 안전성과 건전성을 보장하는 미시 건전성 접근 방식입니다.
\end{examplebox}

\subsubsection*{4. 효율적인 지급결제 시스템 제공 (Provide an Efficient Payment System)}
연방준비제도의 네 번째 핵심 기능은 효율적인 지급결제 시스템을 제공하는 것입니다. 현재 연방준비제도는 다양한 지급결제를 지원합니다.

\begin{examplebox}
\textbf{연준이 지원하는 지급결제 시스템}
\begin{itemize}
    \item \textbf{통화 (Currency):} 연방준비제도는 전통적으로 많은 가계 거래에 사용되었던 지폐 또는 현금을 발행합니다.
    \item \textbf{수표 청산 (Check Clearing):} 연방준비제도는 수표 청산을 용이하게 합니다.
    \item \textbf{자동 어음 교환소 (Automated Clearinghouse, ACH):} ACH는 직원에게 지급되는 임금이나 공과금 납부와 같은 소액 은행 이체를 결제합니다.
    \item \textbf{대규모 거래 (Large Value Transactions):} 기업과 은행 간의 대규모 거래 또는 도매 지급결제도 지원합니다.
\end{itemize}
\end{examplebox}

\subsubsection*{5. 소비자 보호 및 지역 사회 개발 증진 (Promote Consumer Protection and Community Development)}
연방준비제도의 다섯 번째 핵심 기능은 소비자 보호 및 지역 사회 개발을 증진하는 것입니다. 저렴한 금융 상품에 대한 접근성을 높이는 것은 가계의 금융 안정성을 제공합니다. 이를 위해 연방준비제도는 다양한 소비자 보호, 공정 대출, 공정 주택 및 지역 사회 재투자 법률을 시행합니다.

\begin{examplebox}
\textbf{소비자 보호 및 지역 사회 개발 관련 법률}
\begin{itemize}
    \item \textbf{진실한 대출법 (Truth in Lending Act):} 앞서 언급했듯이, 이는 투명성과 공정성에 관한 것입니다.
    \item \textbf{지역 사회 재투자법 (Community Reinvestment Act):} 또 다른 중요한 예시입니다. 이 규칙에 따라 연방준비제도는 지역 사회로 유입되는 신용의 양에 대한 데이터를 수집하고, 금융 기관이 서비스하는 전체 지역 사회에 신용 접근성을 제공하도록 장려합니다. 여기에는 저소득 및 중간 소득 지역에 대한 대출이 포함됩니다.
\end{itemize}
\end{examplebox}

\begin{summarybox}
\textbf{레슨 2-1.2 핵심 요약}
이 세션에서는 연방준비제도 시스템의 목적과 기능에 대한 큰 그림을 검토했습니다. 우리는 연준의 다섯 가지 핵심 기능에 대해 논의했습니다.
\begin{itemize}
    \item 통화정책 수행
    \item 금융 시스템 안정성 증진
    \item 금융 기관 규제 및 감독
    \item 효율적인 지급결제 시스템 제공
    \item 소비자 보호 및 지역 사회 개발 증진
\end{itemize}
이러한 기능들은 총체적으로 미국 경제를 지원하고 더 일반적으로 공익에 기여합니다.
\end{summarybox}


\section{레슨 2-2: 전통적 통화정책}

\subsection{레슨 2-2.1. 화폐 공급과 통화정책}

안녕하세요, 화폐 공급에 대한 강의에 오신 것을 환영합니다. 이 수업에서는 화폐 공급이 무엇이며 어떻게 측정되는지 논의할 것입니다. 우리는 화폐 공급에 대한 두 가지 아이디어를 소개할 것입니다. 첫째, 통화 승수(money multiplier), 둘째, 양적 화폐 이론(quantity theory of money)입니다. 이러한 아이디어들은 화폐 공급을 인플레이션을 포함한 중요한 경제 결과와 연결합니다. 중앙은행이 인플레이션을 목표로 한다면, 화폐 공급을 목표로 하는 것이 합리적일 수 있습니다. 우리는 중앙은행이 화폐 공급을 어떻게 통제하는지 논의할 것입니다.

\begin{conceptbox}
\textbf{화폐 공급의 정의}
돈은 우리가 물건을 사고팔 때 사용하는 것입니다. 넓게 말하면, 구매자와 판매자 사이의 신뢰할 수 있는 지불 수단입니다. \textbf{화폐 공급(Money Supply)}은 유통되는 돈의 총량입니다. 화폐 공급에는 달러 지폐와 동전이 포함되지만, 다른 항목들도 포함될 수 있습니다.
\end{conceptbox}

\subsubsection*{화폐 공급의 측정}
화폐 공급을 어떻게 측정하는지 더 정확히 알아봅시다.

\begin{enumerate}
    \item \textbf{본원 통화 (Monetary Base):} 가장 좁은 정의는 본원 통화입니다. 본원 통화는 유통되는 통화량(currency in circulation)을 포함합니다. 또한 은행의 준비금 잔액(bank reserve balances)도 포함합니다. 이는 은행이 연방준비제도 계좌에 예치해 둔 예금입니다.
    \item \textbf{M1:} 화폐 공급의 다소 넓은 정의는 M1입니다. M1은 본원 통화에서 시작하여 거래 예금(transaction deposits)을 추가합니다. 여기에는 은행의 당좌 예금(checking account) 잔액이 포함됩니다. M1의 아이디어는 현금과 은행이 대출할 수 있는 자금(loanable funds)의 양을 포착하는 것입니다.
\end{enumerate}

\textbf{왜 대출 가능 자금이 그렇게 중요하며, 화폐 공급과 어떻게 관련될까요?}
다음 예시를 고려해 봅시다.

\begin{examplebox}
\textbf{통화 승수 (Money Multiplier)의 작동 원리}
\begin{itemize}
    \item 당신이 은행에 100달러를 예금했다고 가정해 봅시다.
    \item 은행은 10\%인 10달러를 연방준비제도에 준비금으로 보관합니다. (이는 의무적이거나 단순히 은행의 신중함 때문일 수 있습니다.)
    \item 은행은 나머지 90달러를 대출해 줍니다.
    \item 이 90달러는 차입자의 당좌 예금에 새로운 예금으로 나타날 것입니다.
    \item 다시 은행은 10\%인 9달러를 연방준비제도에 준비금으로 보관하고, 81달러를 대출해 줍니다.
    \item 이제 81달러는 새로운 예금으로 나타날 것입니다.
\end{itemize}
이 과정은 계속될 수 있으며, 종종 \textbf{통화 승수}라고 불립니다. 자금이 대출될 때마다 은행 예금, 그리고 따라서 화폐 공급은 증가할 것입니다. 사실, 10\%의 준비금(reserve)이 화폐 공급을 제한하는 요인입니다.
\end{examplebox}

\vspace{1em}

\begin{enumerate}
    \setcounter{enumi}{2} % Continue numbering from 2
    \item \textbf{M2:} 화폐 공급의 가장 일반적인 정의는 M2입니다. 이는 M1에 비교적 단기간에 현금으로 전환될 수 있는 다른 항목들을 더한 것입니다. 여기에는 저축성 예금(savings deposits), 10만 달러 미만의 소액 정기 예금(small time deposits), 소매 머니마켓 뮤추얼 펀드(retail money market mutual fund shares)가 포함됩니다.
    \item \textbf{M3:} 가장 넓은 정의는 M3입니다. 이는 M2에 대규모 정기 예금(large time deposits), 기관 머니마켓 펀드(institutional money market funds), 단기 환매 조건부 계약(short-term repurchase agreements), 그리고 더 큰 유동성 자산(larger liquid assets)을 더한 것입니다.
\end{enumerate}

\subsubsection*{화폐 공급과 거시 경제 결과: 양적 화폐 이론}
역사적으로 화폐 공급 측정치는 주요 거시 경제 결과와 밀접하게 관련되어 있었습니다. 과거의 높은 인플레이션 시기는 화폐 공급이 빠르게 증가했던 시기와 일치합니다. 왜 그럴까요? 다른 상품과 마찬가지로, 돈의 가치는 유통되는 돈이 과도할 경우 하락합니다.

\begin{conceptbox}
\textbf{양적 화폐 이론 (Quantity Theory of Money)}
화폐 공급, 인플레이션, 경제 성장 간의 관계를 \textbf{양적 화폐 이론}이라고 합니다. 이는 경제학에서 가장 오래된 아이디어 중 하나이며, 시카고 대학의 밀턴 프리드먼(Milton Friedman)에 의해 대중화되었습니다.

\textbf{왜 이 관계가 중요할까요?}
중앙은행은 물가 안정과 최대 고용을 목표로 합니다. 이것이 그들의 책무입니다. 역사적으로 화폐 공급 측정치는 인플레이션 및 경제 성장과 밀접하게 관련되어 있었습니다. 이러한 이유로 중앙은행은 통화정책을 결정할 때 화폐 공급 측정치를 사용해 왔습니다.
\end{conceptbox}

\textbf{데이터로 살펴본 화폐 공급과 인플레이션/GDP}
미국의 1960년대 이후 데이터를 살펴보겠습니다. 여기에 M3 화폐 공급이 있습니다. 이제 인플레이션을 측정하기 위해 소비자 물가 지수(consumer price index)를 추가합니다. 1980년대 경기 침체 이후 약간의 차이가 있지만, M3 화폐 공급이 인플레이션을 상당히 잘 추적하는 것을 볼 수 있습니다. 마찬가지로, M3는 2000년까지 명목 국내총생산(nominal cross domestic product)을 잘 추적합니다.

\subsubsection*{중앙은행의 화폐 공급 통제 방법}
원칙적으로 연방준비제도는 화폐 공급을 통제하여 그 책무를 달성할 수 있습니다. 실제로 그들은 몇 가지 방법으로 은행을 목표로 삼아 이를 수행할 수 있습니다.

\begin{enumerate}
    \item \textbf{지급준비율 변경 (Change the Reserve Requirement):} 연방준비제도는 은행의 지급준비율을 변경할 수 있습니다. 통화 승수를 상기해 봅시다. 지급준비율이 높을수록 은행은 더 많은 준비금을 보유해야 하고, 더 적은 대출을 할 수 있습니다. 다시 말해, 연방준비제도는 통화 승수를 통제할 수 있습니다.
    \item \textbf{공개시장운영 수행 (Conduct Open Market Operations):} 연방준비제도는 공개시장운영을 수행할 수 있습니다. 공개시장운영에서 연준은 은행과 증권(예: 국채)을 거래합니다. 연방준비제도가 은행에 증권을 판매하면, 은행은 대출 가능 자금(loanable funds)을 사용하여 연준에 지불해야 합니다. 결과적으로 은행은 대출할 자금이 줄어들고, 이는 통화 승수와 화폐 공급을 감소시킬 것입니다.
\end{enumerate}

\textbf{하지만 문제가 있습니다.}
불행히도 시간이 지남에 따라 화폐 공급, 인플레이션, GDP 간의 관계는 데이터에서 덜 신뢰할 수 있게 되었습니다.
\begin{itemize}
    \item \textbf{측정의 어려움:} 첫째, 이제 새로운 지급결제 기술과 가계가 현금을 보관하는 방법이 생겨 화폐 공급을 측정하기가 더 어려워졌습니다. 연방준비제도는 M3 발표를 중단했습니다.
    \item \textbf{관계의 불안정성:} 둘째, M3와 인플레이션 또는 GDP 간의 관계는 2000년 이후 안정적이지 않았습니다. 2008년 금융 위기 이후 M3와 GDP 간의 관계는 무너졌습니다. 여기에서 명확한 차이를 볼 수 있습니다. COVID-19 팬데믹 동안에도 다시 발생했습니다.
\end{itemize}
대부분의 중앙은행은 이제 통화정책 결정을 내릴 때 화폐 공급 측정치를 덜 중요하게 여깁니다.

\begin{summarybox}
\textbf{레슨 2-2.1 핵심 요약}
이 세션에서는 화폐 공급의 개념과 통화정책과의 연결성에 대해 논의했습니다.
\begin{itemize}
    \item \textbf{화폐 공급 측정:} 우리는 본원 통화(Monetary Base), M1, M2, M3 등 여러 가지 방법으로 화폐 공급을 측정할 수 있음을 논의했습니다.
    \item \textbf{통화 승수:} 통화 승수는 은행 준비금과 화폐 공급을 연결합니다.
    \item \textbf{양적 화폐 이론:} 우리는 화폐 공급, 인플레이션, GDP 간의 관계를 살펴보았습니다.
    \item \textbf{통화정책의 영향:} 통화정책이 화폐 공급을 증가시키거나 감소시킬 수 있음을 보았습니다.
\end{itemize}
\end{summarybox}


\subsection{레슨 2-2.2. 연방준비제도의 대차대조표}

안녕하세요, 연방준비제도 대차대조표에 대한 강의에 오신 것을 환영합니다. 연방준비제도 시스템은 미국의 중앙은행입니다. 다른 금융 기관과 마찬가지로 연준도 대차대조표를 가지고 있습니다. 모든 대차대조표와 마찬가지로 자산과 부채가 있을 것입니다. 다른 점은 연준이 대차대조표를 사용하여 통화정책 목표를 달성한다는 것입니다. 예를 들어, 연준이 자산을 매입하고 대차대조표를 확장하기로 결정하는 것은 통화정책 조치입니다.

우리는 연방준비제도 대차대조표를 면밀히 살펴볼 것입니다. 우리는 그 규모를 살펴볼 것입니다. 우리는 그 구성 요소를 검토할 것입니다. 주요 자산은 무엇이며, 이러한 자산은 어떻게 조달될까요? 그런 다음 최근 몇 년 동안 연준 대차대조표의 구성이 어떻게 변했는지 살펴볼 것입니다. 위기에 대응한 극적인 통화정책 조치가 이러한 극적인 변화를 어떻게 이끌었는지 볼 것입니다. 이것은 이 수업의 중요한 주제가 될 것입니다.

\textbf{연준 대차대조표의 규모 변화}
매주 연방준비제도는 상세한 대차대조표를 발표합니다. 이 정보는 온라인에서 제공되며 시장 참여자들은 면밀히 주시합니다.

시간에 따른 연방준비제도 대차대조표의 규모를 살펴보겠습니다. 규모는 연방준비제도의 총 자산을 의미합니다.

\begin{examplebox}
\textbf{연준 대차대조표 규모의 폭발적인 성장}
\begin{itemize}
    \item \textbf{2000년대 초:} 연방준비제도는 약 8천억 달러 상당의 자산을 보유했습니다.
    \item \textbf{2008년 금융 위기:} 대차대조표는 두 배 이상 증가했습니다. 이때 연방준비제도는 몇 가지 주요 통화정책 개입을 수행했습니다. 이러한 개입은 극적인 영향을 미쳤습니다. 2009년까지 연방준비제도는 2조 달러가 조금 넘는 자산을 보유했습니다.
    \item \textbf{2009년-2014년:} 대차대조표는 다시 두 배 이상 증가하여 4조 달러를 넘어섰습니다.
    \item \textbf{COVID-19 팬데믹:} 대차대조표 규모는 다시 두 배 증가하여 2021년 여름까지 8조 달러에 도달했습니다.
\end{itemize}
\textbf{핵심 요약:} 2008년 이후 연방준비제도 대차대조표는 폭발적으로 성장했습니다. 이러한 성장은 통화정책 조치를 반영합니다.
\end{examplebox}

\subsubsection*{연준 대차대조표의 구성 요소}

\textbf{1. 2008년 이전의 대차대조표 (정상 시기)}
가장 중요한 자산의 내역부터 살펴보겠습니다. 2008년 이전의 대차대조표는 정상 시기의 연준 운영에 대해 알려줄 것입니다.

\begin{examplebox}
\textbf{2000년 1월 12일 연준 대차대조표 (자산 측면)}
\begin{itemize}
    \item \textbf{총 자산:} 약 6,300억 달러.
    \item \textbf{정부 증권 (Government securities):} 4,850억 달러 상당으로, 전체 자산의 4분의 3을 차지하는 가장 큰 자산 포지션입니다.
    \item \textbf{환매 거래 (Report transactions):} 두 번째로 큰 포지션입니다. 이는 연방준비제도가 단기적으로 현금 대출을 하는 경우입니다. 이는 일반적으로 공개시장운영의 일부입니다.
    \textbf{통화 (Currency), 금 (Gold):} 연방준비제도는 동전과 지폐를 포함한 통화와 금도 보유합니다.
\end{itemize}
\end{examplebox}

\begin{examplebox}
\textbf{2000년 1월 12일 연준 대차대조표 (부채 측면)}
\begin{itemize}
    \item \textbf{유통 통화 (Currency in circulation):} 5,900억 달러 상당으로, 가장 큰 포지션입니다. 이는 여러분의 주머니에 있는 달러 지폐와 동전입니다.
    \item \textbf{은행의 준비금 잔액 (Reserve balances of banks):} 연방준비제도에 예치된 은행의 준비금으로, 의무 준비금과 초과 준비금 모두를 포함합니다.
    \item \textbf{본원 통화 (Monetary Base):} 유통 통화와 준비금을 합친 것을 본원 통화라고 합니다. 본원 통화는 연준의 부채입니다. 이는 연준에 의해 발행되며 연방준비제도 시스템의 자산에 의해 뒷받침됩니다.
    \item \textbf{재무부 일반 계정 (Treasury General Account):} 미국 정부의 당좌 예금 계좌입니다. 모든 미국 정부 지불은 이 계좌에서 이루어집니다. 세금 수입도 이 계좌에 예치됩니다.
\end{itemize}
2000년 1월 12일의 대차대조표는 20년 후와 비교하면 작았습니다. 실제로 2000년대 초에는 COVID-19 팬데믹 기간의 대차대조표 규모의 10분의 1도 되지 않았습니다. 이러한 대차대조표 확장은 어떻게 일어났을까요?
\end{examplebox}

\textbf{2. 2008년 이후의 대차대조표 (확장 시기)}
먼저 자산 측면을 살펴보겠습니다. 다음은 2021년 1월 7일 데이터입니다. 2021년 데이터에서 주목할 점이 몇 가지 있습니다.

\begin{examplebox}
\textbf{2021년 1월 7일 연준 대차대조표 (자산 측면)}
\begin{itemize}
    \item \textbf{총 자산:} 7.3조 달러. (이전 6,300억 달러에 비해 훨씬 커졌습니다.)
    \item \textbf{항목 변화:} 주택저당증권(mortgage-backed securities)과 위기 시설(crisis facilities)이 나타납니다.
    \item \textbf{미국 국채 (US treasury securities):} 4.6조 달러 이상으로 여전히 가장 큰 항목입니다.
    \item \textbf{주택저당증권 (Mortgage-backed securities, MBS):} 정부 보증 기업(government sponsored enterprises)의 주택저당증권으로 뒷받침되는 채권으로, 2조 달러 이상을 보유했습니다. 이 채권들은 미국 정부에 의해 보증됩니다.
    \item \textbf{자본 이득 (Capital gains):} 아직 판매되지 않은 증권의 이익입니다. 이는 미실현 자본 이득(unrealized capital gains)이라고도 불립니다.
    \item \textbf{위기 시설 (Crisis facilities):} 다양한 대출 프로그램입니다. 주로 은행에 대한 것이며, 시장 전반의 스트레스 시기에 자금이 제공되었습니다. 은행은 이 대출금을 상환해야 합니다.
\end{itemize}
\textbf{자산 성장의 주요 동인:} 분명히 연방준비제도가 보유한 정부 및 정부 보증 증권입니다.
\end{examplebox}

\textbf{연준은 이 대규모 대차대조표 확장을 어떻게 조달했을까요?}
일반적인 기업이나 가계는 확장을 위해 차입해야 할 수도 있습니다. 그러나 연준은 화폐 공급을 통제하므로 다른 방식으로 이를 수행할 수 있었습니다.

\textbf{연방준비제도는 더 많은 지폐를 인쇄했을까요?}
같은 기간 동안 유통되는 총 통화량을 살펴보겠습니다. 보시다시피, 유통 통화량은 많이 증가했습니다. 20년 동안 거의 4배 증가했습니다. 그러나 이 증가가 전체 이야기는 아닙니다.

\begin{examplebox}
\textbf{2021년 1월 7일 연준 대차대조표 (부채 측면)}
\begin{itemize}
    \item \textbf{지불 방식:} 연방준비제도는 미국 재무부로부터 직접 정부 증권을 매입하는 것이 아니라 은행으로부터 매입합니다. 연준은 이 은행들에게 어떻게 지불했을까요? 지폐로 지불하지 않았습니다. 대신 디지털 화폐로 지불했습니다. 이는 연준에 있는 은행의 준비금 계좌에 대한 신용(credit) 형태로 이루어집니다.
    \item \textbf{은행 준비금 증가:} 연준은 시스템 내 은행 준비금의 양을 증가시켰습니다. 보시다시피, 연방준비제도에 있는 은행 준비금은 2021년까지 부채의 약 40\%로 확장되었습니다.
    \item \textbf{재무부 일반 계정:} 여기에 빠진 것은 무엇일까요? 나머지 부채의 대부분은 미국 정부 일반 계정(US government general account)에 속하며, 최고점에서는 약 20\%를 차지했습니다. 왜 그럴까요? 2008년 이후, 특히 2020년에 미국 정부의 역할이 커졌습니다. 이로 인해 연방준비제도에 더 많은 돈이 예치되었습니다. 정부는 이러한 예치금을 다양한 재정 프로그램에 지불하는 데 사용했습니다.
\end{itemize}
\end{examplebox}

\begin{summarybox}
\textbf{레슨 2-2.2 핵심 요약}
이 강의에서는 연방준비제도 대차대조표와 그 구성 요소를 논의했습니다.
\begin{itemize}
    \item \textbf{주요 자산:} 연방준비제도가 보유한 주요 자산은 정부 증권입니다.
    \item \textbf{주요 부채:} 주요 부채는 본원 통화입니다. 이는 유통 통화와 은행 준비금으로 구성됩니다.
    \item \textbf{규모 변화:} 대차대조표 규모는 2008년과 2020년 사이에 10배 증가했습니다. 이때 연준은 경제를 지원하기 위해 통화정책 도구를 사용했습니다. 연준은 본원 통화, 특히 금융 시스템 내 은행 준비금의 양을 늘림으로써 이러한 지원에 대한 비용을 지불했습니다.
\end{itemize}
\end{summarybox}


\section{중앙은행의 통화정책 도구와 비전통적 접근}

\begin{summarybox}
  \textbf{문서 전체 핵심 요약}
  이 문서는 중앙은행의 핵심 통화정책 도구인 공개시장운영과 재할인창구를 심층적으로 다룹니다.
  특히 2008년 금융위기 이후 등장한 비전통적 통화정책인 포워드 가이던스, 대규모 자산 매입(양적 완화)을 설명합니다.
  또한, 풍부한 은행 준비금 환경에서 금리 통제를 위한 새로운 도구인 지급준비금 이자(IORB)와 역환매조건부채권(ON RRP)의 역할과 작동 방식을 상세히 분석합니다.
  각 정책이 경제에 미치는 영향과 중앙은행의 목표 달성 과정을 이해하는 것이 이 문서의 주요 목표입니다.
\end{summarybox}


\section{용어 정리}

\begin{table}[h!]
\centering
\caption{주요 통화정책 용어 및 설명}
\label{tab:glossary}
\begin{adjustbox}{width=\textwidth,center}
\begin{tabular}{lllc}
\toprule
\textbf{용어} & \textbf{쉬운 설명} & \textbf{원어} & \textbf{비고} \\
\midrule
공개시장운영 & 중앙은행이 국채 등을 사고팔아 시중의 돈(준비금)의 양을 조절하는 것. & Open Market Operations (OMO) & 가장 기본적인 통화정책 수단 \\
연방기금금리 & 은행들이 서로 하룻밤 동안 돈을 빌려주고 빌릴 때 적용되는 금리. & Federal Funds Rate & 미국 통화정책의 주요 목표 금리 \\
재할인창구 & 중앙은행이 은행들에게 단기 대출을 해주는 창구. & Discount Window & 은행의 단기 유동성 지원 및 금리 상한선 역할 \\
재할인율 & 재할인창구를 통해 은행이 중앙은행에서 돈을 빌릴 때 적용되는 금리. & Discount Rate & 연방기금금리의 상한선 역할을 함 \\
포워드 가이던스 & 중앙은행이 미래 통화정책 방향에 대해 시장에 미리 알려주는 소통 방식. & Forward Guidance & 불확실성 감소 및 장기 금리 영향 목적 \\
대규모 자산 매입 & 중앙은행이 장기 국채나 주택저당증권 등을 대량으로 매입하는 정책. & Large-Scale Asset Purchases (LSAPs) & 양적 완화(QE)라고도 불리며, 장기 금리 인하 목적 \\
양적 완화 & 대규모 자산 매입과 동일한 의미로 사용되며, 경제 활성화를 위한 비전통적 정책. & Quantitative Easing (QE) & 단기 금리가 0에 가까울 때 사용 \\
지급준비금 이자 & 중앙은행이 은행이 예치한 준비금에 대해 지급하는 이자. & Interest on Reserve Balances (IORB) & 풍부한 준비금 환경에서 금리 하한선 역할 \\
역환매조건부채권 & 중앙은행이 비은행 금융기관으로부터 증권을 담보로 돈을 빌리고 이자를 지급하는 것. & Overnight Reverse Repurchase Agreement (ON RRP) & 비은행 기관의 단기 금리 하한선 역할 \\
주요 딜러 & 중앙은행의 공개시장운영에 참여하는 대형 금융기관. & Primary Dealers & 주로 대형 은행들이 포함됨 \\
시스템 공개시장계정 & 뉴욕 연방준비은행이 연준의 공개시장운영을 관리하는 계정. & System Open Market Account (SOMA) & 연준의 자산 보유 현황을 보여줌 \\
\bottomrule
\end{tabular}
\end{adjustbox}
\end{table}


\section{Lesson 2-2.3. 공개시장운영 (Open Market Operations)}

\begin{mybox}
  \textbf{핵심 요약}
  공개시장운영은 중앙은행이 국채 등의 증권을 사고팔아 시중의 은행 준비금 양을 조절하고, 이를 통해 연방기금금리(Federal Funds Rate)를 목표 수준으로 유도하는 핵심 통화정책 수단입니다.
\end{mybox}

공개시장운영(Open Market Operations, OMO)은 중앙은행의 통화정책 도구 중 가장 핵심적인 역할을 합니다. 만약 연방준비제도(Federal Reserve, Fed)가 금리를 인상하기로 결정한다면, 이는 공개시장운영을 통해 달성될 수 있습니다.

\subsection{1. 공개시장운영이란 무엇인가?}
공개시장운영은 중앙은행이 증권을 매각하거나 매입하는 행위를 말합니다. 이러한 행위가 실제로 어떻게 작동하는지, 그리고 연방기금시장(Federal Funds Market)에서 준비금(reserves)의 양에 직접적으로 어떤 영향을 미치는지 살펴보겠습니다. 궁극적으로 이러한 공개시장운영이 시장의 균형 이자율을 어떻게 결정하는지 이해하는 것이 중요합니다. 연방기금금리는 미국의 통화정책에서 주요 목표 이자율이라는 점을 기억해야 합니다.

\subsection{2. 통화정책의 목표 이자율: 연방기금금리}
통화정책은 목표 이자율을 선택하는 것을 포함합니다. 이 목표 이자율은 경제에 영향을 미쳐 연방준비제도가 목표를 달성할 수 있도록 돕습니다. 연방준비제도의 목표 이자율은 연방기금금리(Federal Funds Rate)라고 불리며, 이는 단기 이자율입니다. 사실상 이는 은행 간 시장(interbank market)의 이자율입니다. 미국 은행 간 시장은 연방기금시장이라고 불립니다. 모든 은행은 중앙은행에 준비금 계좌를 가지고 있으며, 이 준비금을 연방기금(federal funds)이라고 부릅니다. 공개시장운영은 이 연방기금금리에 영향을 미치도록 설계되었습니다.

\subsection{3. 공개시장운영의 작동 방식 (개념적 이해)}
연방준비제도는 이자율을 직접적으로 변경하지 않습니다. 대신, 공개시장운영을 통해 준비금의 공급을 변경합니다. 준비금의 공급과 수요의 교차점이 이 준비금의 가격을 결정하며, 이 가격이 바로 시장의 이자율, 즉 연방기금금리입니다.

\begin{examplebox}
  \textbf{예시: 금리 인하 목표 시}
  만약 연방준비제도가 이자율을 낮추고 싶다고 가정해 봅시다. 이를 달성하기 위해 연방준비제도는 준비금의 공급 곡선을 오른쪽으로 이동시킬 수 있습니다. 새로운 균형에서는 더 많은 연방기금이 더 낮은 이자율로 거래될 것입니다.
\end{examplebox}

\subsection{4. 준비금 조작을 통한 금리 통제}
공개시장운영은 연방준비제도가 시스템 내 준비금의 양을 조작할 수 있도록 합니다. 연준은 어떻게 이를 수행할까요?

\begin{examplebox}
  \textbf{예시: 은행 준비금 증가 목표 시}
  연준이 은행 준비금을 늘리고 싶다고 가정해 봅시다. 연준은 은행으로부터 증권을 매입하고 그 대가로 은행에 더 많은 준비금을 지급할 수 있습니다. 은행은 이미 이 증권을 소유하고 있었을 것이며, 연준은 공정한 가격을 지불해야 합니다.
\end{examplebox}

다시 한번 정리하면, 연방준비제도가 연방기금금리를 낮추고 싶다면, 은행으로부터 증권을 매입할 수 있습니다. 이러한 매입은 준비금으로 지급됩니다. 따라서 연방기금시장에서 전체 가용 자금은 증가할 것입니다. 이는 공급 곡선을 오른쪽으로 이동시킵니다. 새로운 균형에서는 더 많은 연방기금과 더 낮은 이자율이 형성됩니다.

\subsection{5. 공개시장운영의 실제 절차}
실제로 공개시장운영이 어떻게 작동하는지 몇 단계를 통해 살펴보겠습니다.

\begin{enumerate}
    \item \textbf{1단계: 목표 이자율 결정} \\
    연방준비제도 공개시장위원회(Federal Open Market Committee, FOMC)는 목표 이자율을 결정합니다.
    \item \textbf{2단계: 운영 위임} \\
    실제 증권 매입 및 매각은 뉴욕 연방준비은행(Federal Reserve Bank of New York)에 위임됩니다. 이들은 연방준비제도 시스템 공개시장계정(System Open Market Account, SOMA)을 관리합니다.
    \item \textbf{3단계: 증권 매입} \\
    뉴욕 연방준비은행의 트레이딩 데스크는 소위 주요 딜러(primary dealers)로부터 증권을 매입합니다. 이 딜러들은 연방준비제도에 준비금 계좌를 가지고 있으며, 가장 큰 은행들이 포함됩니다. 주로 미국 국채(US treasuries)가 매입됩니다.
    \item \textbf{4단계: 연준 대차대조표 확장} \\
    트레이딩 데스크가 매입을 완료하면 연방준비제도의 대차대조표가 확장됩니다. 자산 측면에서는 매입한 증권이 추가되고, 부채 측면에서는 준비금이 증가합니다. 이 준비금은 연방준비제도가 새로 창출하는 것입니다. 연준은 통화 기반(monetary base)을 통제하므로 이러한 준비금을 자유롭게 창출할 수 있습니다.
\end{enumerate}

\begin{examplebox}
  \textbf{예시: 연방준비제도의 공개시장 매입}
  2021년 7월 19일, 연준은 액면가 14억 달러의 미국 국채를 매입했습니다. 이 증권들은 15년에서 22년 사이의 잔여 만기를 가지고 있었습니다.
\end{examplebox}

\begin{figure}[h!]
  \centering
  % 실제 이미지가 없으므로 텍스트로 설명합니다.
  \begin{tcolorbox}[mybox, title={\textbf{시스템 공개시장계정 포트폴리오 (2021년 7월 19일)}}]
    이 시점의 연방준비제도 시스템의 국내 증권 보유 현황을 보여주는 스냅샷입니다.
    모든 증권은 미국 정부 부채 또는 미국 정부가 보증하는 부채였습니다.
    포트폴리오는 이전 주 대비 약 1,000억 달러 증가했습니다.
  \end{tcolorbox}
  \caption{시스템 공개시장계정 포트폴리오 (2021년 7월 19일) 설명}
  \label{fig:soma_portfolio}
\end{figure}

\subsection{6. 2008년 금융위기 이후의 변화}
2008년 금융위기는 공개시장운영 방식에 전환점을 가져왔습니다. 연준은 2007년 중반 5.25\%였던 연방기금금리 목표를 2008년 12월 0\%에서 0.25\% 범위로 낮췄습니다. 단기 이자율이 거의 0에 가까워지자, 단순히 더 많은 증권을 매입하여 준비금을 늘리는 방식은 더 이상 효과적이지 않았습니다. 어려움을 겪는 경제를 자극하기 위한 새로운 접근 방식이 필요했습니다.

금융위기 이후, 은행 준비금의 총 공급량은 그 어느 때보다 많아졌습니다. 이는 연방준비제도의 대규모 개입의 직접적인 결과였습니다. 이는 전통적인 공개시장운영이 단기 이자율을 통제하는 데 덜 효과적일 것임을 의미했습니다. 연준은 이러한 문제에 대처하기 위한 새로운 통화정책 도구가 필요했습니다. 이 문제는 별도의 강의에서 논의될 것입니다.

\begin{summarybox}
  \textbf{핵심 정리: 공개시장운영}
  \begin{itemize}
    \item 공개시장운영은 준비금 공급에 영향을 미치며, 이는 연방준비제도의 목표 이자율인 연방기금금리를 결정합니다.
    \item 2008년 금융위기 이전에는 공개시장운영이 주요 통화정책 도구였습니다.
    \item 2008년 이후, 준비금 공급이 증가하면서 전통적인 공개시장운영의 중요성은 감소했습니다.
  \end{itemize}
\end{summarybox}


\section{Lesson 2-2.4. 재할인창구 (The Discount Window)}

\begin{mybox}
  \textbf{핵심 요약}
  재할인창구는 중앙은행이 은행들에게 단기 현금 대출을 제공하는 전통적인 통화정책 도구입니다. 이는 은행의 유동성 관리를 돕고, 재할인율을 통해 연방기금금리의 상한선을 설정하여 통화정책에 영향을 미칩니다.
\end{mybox}

이 강의에서는 연방준비제도의 재할인창구(Discount Window)에 대해 논의합니다. 연방준비제도가 재할인창구를 사용하여 은행에 다양한 단기 현금 대출을 제공하는 방법을 살펴보겠습니다. 주요 신용(primary credit), 보조 신용(secondary credit), 계절 신용(seasonal credit)을 검토하고, 이러한 대출의 특징과 존재 이유를 이해할 것입니다. 그런 다음 재할인창구 대출과 통화정책 간의 연관성을 파악하고, 재할인율(discount rate)이 연방준비제도의 목표 이자율에 어떻게 영향을 미치는지 이해할 것입니다.

\subsection{1. 재할인 대출의 기본 개념}
재할인 대출(Discount lending)은 은행업의 오랜 전통입니다. 기본 아이디어는 대출 기관이 먼저 필요한 이자와 수수료 금액을 계산한다는 것입니다. 그런 다음 대출 기관은 이 금액을 대출의 액면가에서 할인합니다. 차용인은 오늘 할인된 대출 금액을 받습니다. 이 대출 금액은 이자와 수수료를 미리 반영하여 이미 줄어든 상태이지만, 차용인은 만기 시 전체 액면가를 상환해야 합니다.

\begin{examplebox}
  \textbf{예시: 재할인 대출 계산}
  액면가 100달러의 재할인 대출을 받고 싶다고 가정해 봅시다. 대출 만기는 1년입니다. 대출 기관이 10\%의 재할인율을 계산했다고 가정해 봅시다. 이는 대출 이자와 수수료를 포함합니다. 오늘 당신이 받는 대출 금액은 90달러입니다. 이는 액면가 100달러에서 이자와 수수료 10달러를 뺀 금액입니다. 하지만 1년 후에는 전체 100달러를 상환해야 합니다.
\end{examplebox}

\subsection{2. 재할인 대출의 유용성}
이러한 유형의 대출은 왜 유용할까요? 대출 비용이 선불로 처리되기 때문입니다. 이는 일련의 월별 이자 지급 대신 한 번의 상환만 필요하다는 것을 의미합니다. 이러한 이유로 재할인 대출은 종종 단기 자금 조달 필요성을 해결하는 데 도움이 됩니다.

\subsection{3. 연방준비제도의 재할인창구}
1913년부터 연방준비제도는 은행에 재할인 대출을 제공해 왔습니다. 연준은 왜 이런 일을 할까요? 이러한 대출은 은행에 단기간에 현금을 제공하여 은행 시스템에 충분한 현금이 있도록 보장합니다. 극단적인 경우, 이는 뱅크런(bank runs)과 은행 공황(banking panics)을 방지할 수 있습니다. 재할인 대출은 각 지역 연방준비은행의 특별 창구인 재할인창구에서 제공되었으며, 그래서 재할인창구 대출(discount window lending)이라는 이름이 붙었습니다. 오늘날 연준은 적격 담보로 담보된 단기 현금 대출을 제공하며, 이를 담보 재할인 대출(secured discount loans)이라고 합니다. 그러나 재할인창구라는 용어는 여전히 연방준비제도가 은행에 직접 대출할 때 사용됩니다.

\subsection{4. 재할인창구 대출의 종류: 주요 신용 (Primary Credit)}
연방준비제도에서 은행은 현재 여러 유형의 재할인 대출에 접근할 수 있습니다. 우리는 가장 중요한 범주인 주요 신용(primary credit)에 초점을 맞출 것입니다. 주요 신용의 주된 목적은 다른 자금 조달이 불가능한 경우 현금을 제공하는 것입니다. 주요 신용 재할인 대출은 자본이 충분하고 건전한 은행에 제공됩니다. 이들은 최대 만기 90일의 단기 대출입니다.

\begin{mybox}
  \textbf{주요 신용의 특징}
  \begin{itemize}
    \item \textbf{이자율:} 주요 신용 대출의 이자율은 목표 연방기금금리에 상대적으로 설정됩니다. 이는 일반적으로 다른 형태의 단기 신용 이자율보다 높습니다.
    \item \textbf{목적:} 이는 은행이 민간 현금 시장에서 차입하도록 유도하기 위함입니다. 연준은 최후의 수단(last resort)이 되어야 합니다.
    \item \textbf{담보:} 이들은 담보 대출입니다. 연방준비제도는 상업 및 산업 대출, 소비자 대출, 주거 및 상업용 부동산 대출, 미국 국채를 포함한 광범위한 담보를 허용합니다.
    \item \textbf{마진:} 일부 담보는 다른 담보보다 위험합니다. 위험을 반영하기 위해 연준은 은행이 각 유형의 담보에 대해 얼마나 빌릴 수 있는지를 나타내는 마진(margins)을 적용합니다.
  \end{itemize}
\end{mybox}

\begin{figure}[h!]
  \centering
  % 실제 이미지가 없으므로 텍스트로 설명합니다.
  \begin{tcolorbox}[mybox, title={\textbf{마진 테이블 예시}}]
    이 마진은 시장 가치의 백분율로 표시됩니다.
    \begin{itemize}
      \item 잔여 만기 1년의 미국 국채는 매우 안전합니다. 은행은 1달러당 99센트를 빌릴 수 있습니다.
      \item 대조적으로, 잔여 만기 1년의 상업용 부동산 대출은 여전히 상당히 위험합니다. 은행은 대출 가치의 35\%에서 95\% 사이만 빌릴 수 있습니다.
    \end{itemize}
  \end{tcolorbox}
  \caption{담보 유형별 마진 예시 설명}
  \label{fig:margin_table}
\end{figure}

\subsection{5. 재할인창구와 통화정책의 연관성}
재할인창구는 연방준비제도가 경제의 단기 이자율을 통제하는 데 어떻게 도움이 될까요? 두 가지를 염두에 두어야 합니다.
\begin{enumerate}
    \item \textbf{대체재 관계:} 연준으로부터의 재할인 대출과 다른 은행으로부터의 은행 간 대출은 밀접한 대체재입니다. 둘 다 단기 대출입니다. 모든 건전한 은행은 은행 간 시장이나 재할인창구에 접근할 수 있습니다.
    \item \textbf{재할인율의 상한선 역할:} 재할인율은 연방기금 목표 범위의 상한선 근처에 설정된다는 점을 기억하십시오.
\end{enumerate}

그렇다면 연방기금의 실현 이자율이 재할인율을 초과하면 어떻게 될까요? 은행이 재할인창구에서 빌리는 것이 더 저렴할 것입니다. 은행 간 대출에서 멀어지는 이러한 수요 이동은 연방기금금리에 하향 압력을 가할 것입니다. 따라서 재할인율은 중앙은행이 설정한 연방기금 목표의 상한선을 강제합니다. 이는 실현된 연방기금금리가 너무 높아지지 않도록 보장합니다.

\begin{figure}[h!]
  \centering
  % 실제 이미지가 없으므로 텍스트로 설명합니다.
  \begin{tcolorbox}[mybox, title={\textbf{주요 신용 금리와 연방기금금리 비교}}]
    데이터에서 이를 확인할 수 있습니다. 2003년부터 주요 신용 금리는 항상 연방기금금리보다 높게 유지되었습니다.
  \end{tcolorbox}
  \caption{주요 신용 금리와 연방기금금리 추이 설명}
  \label{fig:primary_credit_vs_fed_funds}
\end{figure}

\begin{summarybox}
  \textbf{핵심 정리: 재할인창구}
  \begin{itemize}
    \item 재할인창구 대출은 은행에 대한 단기 담보 대출입니다. 이는 은행이 단기 현금 필요성을 관리하는 데 도움이 됩니다.
    \item 재할인 대출은 은행 간 시장에서 연방기금을 차입하는 것의 대체재입니다.
    \item 결과적으로 재할인율은 연방기금금리의 상한선을 제공합니다.
  \end{itemize}
\end{summarybox}


\section{Lesson 2-3. 비전통적 통화정책 (Non-traditional Monetary Policy)}

\subsection{Lesson 2-3.1. 포워드 가이던스 (Forward Guidance)}

\begin{mybox}
  \textbf{핵심 요약}
  포워드 가이던스는 중앙은행이 미래 통화정책 의도를 시장에 소통하여 시장의 기대를 형성하고 불확실성을 줄이는 비전통적 통화정책 도구입니다. 이는 단기 금리가 0에 가까워졌을 때 장기 금리에 영향을 미치기 위해 도입되었습니다.
\end{mybox}

이 강의에서는 새로운 통화정책 도구인 포워드 가이던스(Forward Guidance)를 살펴보겠습니다. 포워드 가이던스는 연방준비제도가 미래 통화정책 의도를 소통할 수 있도록 합니다. 이는 오늘 취해진 조치보다는 미래의 조치에 대한 소통입니다. 이 소통은 미래 인플레이션이나 경제 성장에 관한 것일 수 있으며, 종종 연준이 미래에 이자율을 어떻게 조정할 계획인지 포함합니다. 목표는 시장 기대를 형성하고 불확실성을 줄이는 것입니다. 포워드 가이던스의 여러 예시를 논의하고, 연준이 왜 포워드 가이던스를 사용하는지, 그리고 그것이 금융 상황에 어떻게 영향을 미치는지 이해할 것입니다.

\subsection{1. 포워드 가이던스의 등장 배경}
포워드 가이던스는 비전통적 통화정책 도구입니다. 전통적인 통화정책은 단기 이자율에 작용하며, 실제로는 공개시장운영과 재할인창구 대출을 의미합니다. 포워드 가이던스는 비교적 새로운 개념으로, 2009년 연방기금금리가 0에 가까워졌을 때 등장했습니다. 이 시점에서는 단기 이자율을 목표로 하는 것이 더 이상 선택 사항이 아니었습니다. 이 문제를 극복하기 위해 연방공개시장위원회(FOMC)는 미래에 통화정책을 어떻게 조정할 의도인지 소통하기로 결정했습니다. 연준은 장기 이자율을 목표로 삼으려 했습니다.

\subsection{2. 포워드 가이던스의 형태와 목표}
포워드 가이던스는 FOMC의 공식적인 소통입니다. 이는 보도 자료, 연설, 의회 증언 등의 형태로 나타날 수 있습니다. 이들은 일반적으로 대중에게 매우 신중하게 작성된 성명서로, 혼란을 피하는 것이 목표입니다. 포워드 가이던스는 대중에게 통화정책 의도를 알립니다. 이는 향후 몇 달 또는 몇 년 동안의 조치일 수 있으며, 경제가 어떻게 진화하는지에 따라 달라질 수 있습니다. 목표는 장기 이자율을 움직이는 것일 수 있습니다. 포워드 가이던스는 또한 불확실성을 줄일 수 있습니다.

\subsection{3. 초기 FOMC 성명서의 특징 (2009년 이전)}
연준이 어떻게 소통했는지에 대한 초기 예시를 먼저 살펴보겠습니다. 1994년부터 FOMC는 가장 최근의 통화정책 회의 결정 사항을 발표하는 성명서를 발행했습니다. 2009년 이전에는 이러한 성명서가 간결했습니다.

\begin{examplebox}
  \textbf{예시: 2004년 3월 16일 성명서}
  이 성명서는 이자율에 대한 한 문장, 현재 경제 상황을 요약한 한 단락, 그리고 잠재적 위험에 대한 한 단락으로 구성되어 있습니다. 또한 어떤 FOMC 위원들이 조치에 찬성했고, 반대한 위원이 있다면 누구인지도 명시합니다. 마지막으로, 미래 통화정책 조치에 대한 어떠한 언급도 없다는 점에 주목하십시오.
\end{examplebox}

\begin{tcolorbox}[mybox, title={\textbf{성명서 비교를 통한 시장 분석 (2004년 1월 vs 3월)}}]
  2004년 1월 28일 성명서와 비교해 봅시다. 두 번째 단락에서 경제에 대한 설명이 약간 변경되었습니다. 1월에는 확장이 위험했지만, 3월에는 "견고한 속도의 확장"에 불과했습니다. 이는 성장률이 여전히 긍정적이지만 감소했음을 나타냅니다. 마찬가지로 노동 시장에 대한 설명도 덜 우호적입니다. 마지막 단락은 "that"이라는 단어의 삽입을 제외하고는 동일합니다. 이러한 언어의 변화는 대중에게 면밀히 관찰됩니다.
  
  \textbf{시사점:} 이 두 성명서를 비교하여 무엇을 알 수 있을까요? 6주 동안 연방준비제도는 경제 상황에 대한 의견을 변경했습니다. 1월 성명서는 상당히 낙관적이었지만, 3월 성명서는 덜 낙관적입니다. 이것이 통화정책에 무엇을 의미할까요? 아마도 연방준비제도는 다음 회의에서 이자율을 인상하지 않을 것입니다. 이는 사실입니다. 비록 미래 통화정책에 대한 명확한 언급은 포함되어 있지 않지만 말입니다.
\end{tcolorbox}

\subsection{4. 명시적 포워드 가이던스의 등장 (2009년 이후)}
이제 더 최근의 FOMC 성명서를 살펴보겠습니다. 여기에는 명시적인 포워드 가이던스가 포함되어 있습니다.

\begin{examplebox}
  \textbf{예시: 2009년 3월 18일 성명서}
  이 시점에서 연방기금금리는 0에 가까웠습니다. FOMC는 무엇이라고 말했을까요?
  "위원회는 연방기금금리 목표 범위를 0에서 0.25\%로 유지할 것이며, 경제 상황이 \textbf{상당히 오랜 기간 동안(for an extended period)} 연방기금금리의 예외적으로 낮은 수준을 정당화할 가능성이 있다고 예상합니다."
  
  이것이 바로 포워드 가이던스입니다. 이 성명서는 연방기금금리가 "상당히 오랜 기간 동안" 예외적으로 낮은 수준, 즉 0으로 유지될 것임을 시사합니다. 시간 프레임은 정확히 명확하지 않습니다. 대중은 이를 어떻게 해석해야 할까요? 금융 시장 분석가들은 아마도 연방기금금리가 최소 1년, 그리고 훨씬 더 오랫동안 0으로 유지될 것으로 예상했을 것입니다.
\end{examplebox}

\begin{examplebox}
  \textbf{예시: 2012년 12월 12일 성명서 (조건부 포워드 가이던스)}
  이제 2012년 12월 12일 성명서를 고려해 봅시다. 이 시점에서 연방기금금리는 여전히 0에 가까웠지만, 금리 인상이 가능성이 있었습니다. 이 성명서는 훨씬 더 길다는 것을 알 수 있습니다. FOMC는 경제에 대해 더 자세히 설명합니다. 그들은 가계 및 기업 투자와 같은 특정 부문을 언급합니다.
  
  그러나 여기에서 FOMC는 새로운 것을 합니다. 5번째 단락에서 다음과 같이 명시합니다.
  "위원회는 연방기금금리 목표 범위를 0에서 0.25\%로 유지하기로 결정했으며, 현재 이 예외적으로 낮은 연방기금금리 범위가 적절할 것으로 예상합니다. \textbf{적어도 실업률이 6.5\%를 초과하는 한, 1~2년 후 인플레이션이 위원회의 2\% 장기 목표보다 0.5\% 포인트 이상 높지 않을 것으로 예상되는 한, 그리고 장기 인플레이션 기대치가 잘 고정되어 있는 한 말입니다.} 이 문장은 상당히 길고 복잡합니다."
\end{examplebox}

\begin{tcolorbox}[mybox, title={\textbf{조건부 포워드 가이던스 해설}}]
  이것을 풀어봅시다. FOMC는 다음과 같은 정확한 조건 하에서 목표 연방기금금리를 0으로 유지할 것이라고 말합니다.
  \begin{enumerate}
    \item 실업률이 6.5\%를 초과하는 한.
    \item 1~2년 후 예상 인플레이션이 2.5\%를 초과하지 않는 한.
    \item 인플레이션 기대치가 잘 고정되어 있는 한 (즉, 모든 사람이 장기적으로 2\% 인플레이션을 여전히 예상하는 한).
  \end{enumerate}
  
  다시 말해, FOMC는 이러한 정확한 조건 하에서 이자율을 낮게 유지할 것입니다. 이러한 유형의 포워드 가이던스는 계획의 확실성을 제공합니다. 이는 "상당히 오랜 기간"이 무엇을 의미하는지 명확히 합니다.
\end{tcolorbox}

\subsection{5. FOMC 회의록 및 경제 전망 요약 (SEP)}
FOMC는 또한 금리 설정 회의록을 공개합니다. 이들은 회의에서 논의된 내용을 그대로 기록한 것입니다. 회의록은 각 회의 후 3주 후에 공개됩니다. 회의록은 논의된 정책 관련 주제, 참가자들이 표명한 견해, 전망에 대한 위험과 불확실성, 위원회 결정의 이유에 대한 추가 통찰력을 제공합니다. 물론 금융 시장 분석가들은 미래 통화정책에 대한 통찰력을 얻기 위해 이 회의록을 면밀히 검토합니다.

2007년부터 FOMC는 경제 전망 요약(Summary of Economic Projections, SEP)도 발행하고 있습니다. 이는 매년 4회 발행됩니다. 이 전망은 FOMC 참가자들의 미래 경제 결과에 대한 평가를 제공합니다. 여기에는 GDP, 실업률, 인플레이션, 이자율이 포함됩니다. 이러한 결과는 중기 및 장기적으로 고려됩니다. 각 FOMC 참가자가 평가를 제공하므로 대중은 다양한 예측 범위를 볼 수 있습니다.

\begin{figure}[h!]
  \centering
  % 실제 이미지가 없으므로 텍스트로 설명합니다.
  \begin{tcolorbox}[mybox, title={\textbf{2021년 3월 경제 전망 요약 (실질 GDP 변화)}}]
    2021년 3월 17일 경제 전망 요약을 살펴보겠습니다. 당시 경제는 COVID-19 팬데믹으로 인해 어려움을 겪고 있었습니다.
    
    \textbf{실질 GDP 변화 차트:} 인플레이션을 고려한 GDP 성장률입니다. 2021년 예측은 5\%에서 7.3\% 사이로 다양하다는 것을 알 수 있습니다. 단기 예측의 차이는 회복에 대한 불확실성을 반영합니다. 장기적으로는 참가자들의 추정치가 서로 매우 가깝습니다. 장기적으로는 의견 불일치가 거의 없습니다.
  \end{tcolorbox}
  \caption{2021년 3월 실질 GDP 변화 전망 설명}
  \label{fig:sep_gdp}
\end{figure}

\begin{figure}[h!]
  \centering
  % 실제 이미지가 없으므로 텍스트로 설명합니다.
  \begin{tcolorbox}[mybox, title={\textbf{2021년 3월 경제 전망 요약 (닷 플롯)}}]
    2021년 3월 경제 전망 요약에서 소위 닷 플롯(dot plots)도 있습니다. 이 차트는 FOMC가 연방기금금리가 어떻게 될 것이라고 생각하는지를 요약합니다. 금융 분석가들은 이 포워드 가이던스에 세심한 주의를 기울입니다.
    
    \textbf{닷 플롯 분석:} 당시 이자율은 0에 가까웠습니다. 질문은 금리가 얼마나 빨리 인상될 것인가였습니다.
    \begin{itemize}
      \item 2021년에는 어떤 참가자도 0에서 인상을 예상하지 않았습니다.
      \item 2022년에는 4명의 참가자가 경제가 연방기금금리를 최소 한 번 인상할 만큼 충분히 회복될 것이라고 생각했습니다.
      \item 2023년에는 거의 절반의 참가자가 연방기금금리가 인상될 것으로 보았습니다.
      \item 장기적으로 참가자들은 연방기금금리가 2\%에서 3\% 사이가 될 것으로 예상합니다.
    \end{itemize}
    물론 이러한 예측은 수정될 수 있습니다. 그러나 이들은 FOMC 참가자들의 사고에 대한 귀중한 통찰력을 제공합니다. 이러한 종류의 포워드 가이던스는 금융 시장 참가자들에게 더 많은 계획의 확실성을 제공합니다.
  \end{tcolorbox}
  \caption{2021년 3월 닷 플롯 설명}
  \label{fig:sep_dot_plot}
\end{figure}

\begin{summarybox}
  \textbf{핵심 정리: 포워드 가이던스}
  \begin{itemize}
    \item 포워드 가이던스는 중앙은행이 경제 및 통화정책이 어떻게 진화할 것으로 예상하는지 신호를 보낼 수 있는 통화정책 도구입니다.
    \item 대중, 특히 금융 시장 참가자들은 이에 세심한 주의를 기울입니다.
    \item 일상적인 포워드 가이던스는 단기 이자율이 매우 낮았던 2009년에 등장했습니다.
    \item 동시에 연방준비제도는 장기 이자율에 압력을 가하기 위해 포워드 가이던스를 도입했습니다.
  \end{itemize}
\end{summarybox}


\subsection{Lesson 2-3.2. 대규모 자산 매입 (Large-Scale Asset Purchases)}

\begin{mybox}
  \textbf{핵심 요약}
  대규모 자산 매입(LSAPs), 또는 양적 완화(QE)는 단기 금리가 0에 가까울 때 장기 금리에 직접적으로 영향을 미쳐 경제를 자극하기 위해 사용되는 비전통적 통화정책입니다. 중앙은행이 장기 채권을 대량으로 매입하여 채권 가격을 올리고 수익률을 낮춥니다.
\end{mybox}

이 강의에서는 대규모 자산 매입(Large-Scale Asset Purchases, LSAPs)에 대해 다룹니다. LSAPs는 종종 양적 완화(Quantitative Easing, QE)라고 불립니다. 이들은 장기 이자율에 직접적으로 영향을 미치도록 설계된 비전통적 통화정책의 한 형태입니다. 연방준비제도는 2008년 금융위기 동안 이 도구를 처음 사용했습니다. 당시 연준은 단기 이자율이 0이었음에도 불구하고 경제를 자극하려고 노력했습니다. 이 세션에서는 LSAPs가 어떻게 작동하는지, 이러한 프로그램이 어떻게 설계되었는지, 장기 이자율에 어떻게 영향을 미쳤는지, 그리고 경제 전반에 미친 광범위한 영향을 살펴보겠습니다.

\subsection{1. 대규모 자산 매입의 필요성}
전통적인 통화정책은 단기 이자율, 일반적으로 연방기금금리 또는 재할인율에만 작용합니다. 연준이 단기 이자율을 인하하면 장기 이자율도 낮아질 수 있습니다. 그러나 장기 이자율은 덜 움직입니다. 2008년 후반, 미국 경제는 심각한 경기 침체에 진입하고 있었습니다. 연방준비제도는 경제를 자극하고 싶었지만, 연방기금금리는 이미 0으로 인하된 상태였습니다. 통화정책에 대한 새로운 접근 방식이 필요했습니다.

\subsection{2. 중앙은행의 대응: QE 프로그램}
중앙은행은 이 문제에 어떻게 대응했을까요? 연방준비제도는 대량의 장기 채권을 매입하기 시작했습니다. 이들이 소위 QE 프로그램입니다. QE1, QE2, QE3, 그리고 오퍼레이션 트위스트(Operation Twist)의 네 가지 프로그램이 있었습니다.

\begin{mybox}
  \textbf{QE 프로그램의 작동 원리}
  \begin{itemize}
    \item \textbf{목표:} 이러한 매입은 장기 채권에 대한 대규모 수요 충격처럼 작용했습니다.
    \item \textbf{가격 상승:} 초과 수요는 가격을 상승시킵니다.
    \item \textbf{수익률 하락:} 가격이 높다는 것은 채권 수익률이 낮다는 것을 의미합니다. 채권 가격과 채권 수익률 사이에는 역관계가 있다는 것을 기억하십시오.
    \item \textbf{경제 자극:} 따라서 이러한 매입은 장기 채권 수익률을 직접적으로 낮추는 것을 목표로 합니다. 장기 채권 수익률이 낮아지면 장기 투자 기간에 대한 차입 비용이 줄어듭니다. 이는 투자와 경제 활동을 활성화해야 합니다.
  \end{itemize}
\end{mybox}

\subsection{3. QE1 프로그램의 규모와 효과}
QE1 프로그램의 맥락에서 이러한 프로그램을 살펴보고, 그 규모와 효과를 이해해 봅시다. QE1은 2008년 11월에 시작되어 2009년 3월에 확대되었습니다. 2010년 3월에 종료될 때까지 연준은 1조 2,500억 달러의 주택저당증권(mortgage-backed securities, MBS), 1,750억 달러의 GSE(Fannie Mae 및 Freddie Mac) 발행 부채, 그리고 3,000억 달러의 미국 국채를 매입했습니다. 이러한 매입은 당시 시장 가치의 20\% 이상을 차지했습니다. 이는 중앙은행의 대규모 개입입니다.

첫 번째 양적 완화는 성공적이었을까요? 장기 수익률을 살펴보겠습니다.

\begin{figure}[h!]
  \centering
  % 실제 이미지가 없으므로 텍스트로 설명합니다.
  \begin{tcolorbox}[mybox, title={\textbf{QE1의 장기 국채 수익률 영향}}]
    \textbf{장기 국채:} 2년 만기 국채 수익률 차트를 보면, 프로그램이 시행된 후 2008년 10월 후반 1.6\%에서 2009년 4월 1\%로 떨어졌습니다.
  \end{tcolorbox}
  \caption{QE1 시행 후 2년 만기 국채 수익률 변화 설명}
  \label{fig:qe1_treasury_yield}
\end{figure}

\begin{figure}[h!]
  \centering
  % 실제 이미지가 없으므로 텍스트로 설명합니다.
  \begin{tcolorbox}[mybox, title={\textbf{QE1의 회사채 수익률 영향}}]
    \textbf{회사채:} BAA 등급 회사채 수익률을 살펴보겠습니다. 이는 자본 접근성이 제한적인 기업의 차입 비용입니다. 차트에서 볼 수 있듯이, QE1이 시작되면서 수익률은 최고 9.5\%에서 약 8\%로 떨어졌습니다. 이는 기업의 차입 비용을 낮췄습니다.
  \end{tcolorbox}
  \caption{QE1 시행 후 BAA 등급 회사채 수익률 변화 설명}
  \label{fig:qe1_corporate_bond_yield}
\end{figure}

\begin{figure}[h!]
  \centering
  % 실제 이미지가 없으므로 텍스트로 설명합니다.
  \begin{tcolorbox}[mybox, title={\textbf{QE1의 주택담보대출 시장 영향}}]
    \textbf{주택담보대출 시장:} 연준은 주택저당증권(MBS)을 대량으로 매입했습니다. MBS 수익률과 주택담보대출 금리 사이에는 밀접한 관계가 있습니다. MBS 수익률이 하락하면 일반적으로 주택담보대출 금리도 하락합니다. 차트에서 QE1 프로그램이 시작되었을 때 30년 고정 주택담보대출 금리가 2008년 10월 후반 거의 6.5\%에서 2009년 4월 초 4.85\%로 떨어지기 시작한 것을 볼 수 있습니다. 이는 상당한 효과를 가져왔습니다. 주택담보대출 금리 인하는 주택 구매를 장려했습니다. 이러한 개입은 주택 시장을 지원하는 데 도움이 되었을 수 있습니다.
  \end{tcolorbox}
  \caption{QE1 시행 후 30년 고정 주택담보대출 금리 변화 설명}
  \label{fig:qe1_mortgage_rate}
\end{figure}

\subsection{4. 다른 QE 프로그램과 연준 대차대조표의 변화}
다른 양적 완화 프로그램들도 비슷한 성격을 가졌습니다. QE2는 2010년 10월에 발표되었고, QE3는 2012년 9월에 발표되었습니다. 오퍼레이션 트위스트(Operation Twist)는 연방준비제도가 단기 증권을 매각하고 그 수익금으로 장기 증권을 매입하는 것을 포함했습니다. 각 경우에 의도는 장기 이자율을 낮추는 것이었습니다. 이는 가계와 기업의 차입 비용을 줄입니다. 이러한 프로그램은 2014년 10월에 중단되었습니다. 이 시점에는 노동 시장이 회복되었습니다.

\begin{figure}[h!]
  \centering
  % 실제 이미지가 없으므로 텍스트로 설명합니다.
  \begin{tcolorbox}[mybox, title={\textbf{QE 프로그램이 연준 대차대조표에 미친 영향}}]
    이 차트는 QE 프로그램이 연방준비제도의 대차대조표에 미친 영향을 보여줍니다. 이러한 대규모 자산 매입은 결국 연준의 자산으로 귀결되었습니다.
    
    2014년까지 연방준비제도는 거의 4조 달러에 달하는 증권을 보유했습니다. 이들은 주로 미국 국채와 MBS였습니다. 이는 2008년 금융위기 이전 기간에 비해 대차대조표의 극적인 확장입니다.
  \end{tcolorbox}
  \caption{QE 프로그램에 따른 연준 대차대조표 변화 설명}
  \label{fig:fed_balance_sheet_qe}
\end{figure}

\subsection{5. COVID-19 팬데믹 대응으로서의 LSAPs}
대규모 자산 매입은 COVID-19 팬데믹에 대한 연방준비제도의 대응의 일부이기도 했습니다. 2020년 3월, FOMC는 연방준비제도가 "원활한 시장 기능과 통화정책의 광범위한 금융 상황으로의 효과적인 전달을 지원하는 데 필요한 만큼" 국채와 기관 주택저당증권을 계속 매입할 것이라고 발표했습니다. 이 프로그램은 경제를 지원하기 위한 또 다른 개방형 약속이었습니다.

\begin{summarybox}
  \textbf{핵심 정리: 대규모 자산 매입 (LSAPs/QE)}
  \begin{itemize}
    \item 단기 이자율이 0\%일 때, 연방준비제도는 대규모 자산 매입(QE)을 사용하여 통화정책을 수행할 수 있습니다.
    \item 대규모 자산 매입은 장기 이자율을 목표로 할 수 있습니다.
    \item 대규모 자산 매입은 장기 이자율을 낮춤으로써 소비자와 기업의 신용 비용을 줄입니다.
  \end{itemize}
\end{summarybox}


\subsection{Lesson 2-3.3. 풍부한 은행 준비금 환경에서의 통화정책 (Monetary Policy With Abundant Bank Reserves)}

\begin{mybox}
  \textbf{핵심 요약}
  2008년 금융위기 이후 대규모 자산 매입으로 인해 은행 준비금이 풍부해지면서, 전통적인 공개시장운영의 금리 통제 효과가 감소했습니다. 이에 연준은 지급준비금 이자(IORB)와 역환매조건부채권(ON RRP)이라는 새로운 도구를 도입하여 풍부한 준비금 환경에서도 단기 금리를 효과적으로 통제할 수 있게 되었습니다.
\end{mybox}

이 강의에서는 은행 준비금이 풍부할 때 통화정책을 어떻게 구현하는지 다룹니다. 2008년 이후, 그리고 대규모 자산 매입 프로그램 이후, 은행 시스템은 준비금으로 넘쳐났습니다. 이는 전통적인 공개시장운영이 이자율을 관리하는 데 덜 효과적이게 만들었습니다. 이 강의에서는 이러한 환경을 위해 만들어진 두 가지 새로운 통화정책 도구를 검토합니다. 첫째, 지급준비금 이자(Interest on Reserve Balances, IORB)와 둘째, 역환매조건부채권(Overnight Reverse Repurchase Agreement, ON RRP) 시설입니다. 이러한 도구는 은행 준비금 수준이 높을 때 연방준비제도가 단기 이자율을 통제하는 데 도움이 됩니다.

\subsection{1. 전통적인 금리 목표 방식의 한계}
연방준비제도는 어떻게 이자율을 목표로 삼을까요? 역사적으로 연방준비제도는 공개시장운영을 수행했습니다. 이는 연방기금시장에서 준비금 공급 수준을 변경하기 위해 증권을 사고파는 것을 포함합니다.

\begin{figure}[h!]
  \centering
  % 실제 이미지가 없으므로 텍스트로 설명합니다.
  \begin{tcolorbox}[mybox, title={\textbf{전통적인 연방기금시장 모델}}]
    이 그림은 연방기금시장을 보여줍니다. 은행 준비금의 공급을 조정함으로써 연방준비제도는 연방기금금리를 통제합니다.
    \textbf{그래프 설명:} 수직축은 연방기금금리, 수평축은 준비금의 양을 나타냅니다. 준비금 공급 곡선은 수직선이며, 수요 곡선은 우하향합니다. 공급 곡선을 이동시켜 균형 금리를 조절합니다.
  \end{tcolorbox}
  \caption{전통적인 연방기금시장 모델 설명}
  \label{fig:traditional_fed_funds_market}
\end{figure}

그러다 2008년 금융위기가 닥쳤습니다. 연방기금금리는 0으로 인하되었고, 대규모 자산 매입이 이루어졌습니다. 이러한 매입은 은행 준비금의 급격한 증가로 자금이 조달되었습니다. 준비금은 2007년 140억 달러에서 2014년 2조 6천억 달러 이상으로 증가했습니다.

\begin{figure}[h!]
  \centering
  % 실제 이미지가 없으므로 텍스트로 설명합니다.
  \begin{tcolorbox}[mybox, title={\textbf{풍부한 준비금 환경에서의 연방기금시장}}]
    우리는 다른 강의에서 이 차트를 검토했습니다. 연방기금시장의 균형에 미치는 영향을 볼 수 있습니다. 준비금 공급이 훨씬 더 많아졌고, 연방기금금리는 0입니다.
    \textbf{그래프 설명:} 준비금 공급 곡선이 극도로 오른쪽으로 이동하여 수요 곡선의 거의 평평한 부분과 교차합니다. 이로 인해 준비금의 양이 크게 증가했음에도 불구하고 연방기금금리는 0에 가깝게 유지됩니다.
  \end{tcolorbox}
  \caption{풍부한 준비금 환경에서의 연방기금시장 설명}
  \label{fig:abundant_reserves_fed_funds_market}
\end{figure}

\subsection{2. 풍부한 준비금 환경의 문제점}
이것이 왜 문제일까요? 결국 연방준비제도는 이자율을 인상하고 싶을 것입니다. 그러나 연방기금금리는 더 이상 은행 준비금의 작은 변화에 반응하지 않을 것입니다. 전통적인 공개시장운영만으로는 충분하지 않을 것입니다. 연방기금금리가 반응하려면 준비금 공급에 큰 변화가 필요했습니다. 이론적으로 연방준비제도는 대규모 자산 매각을 통해 준비금을 줄일 수 있었습니다. 매입이 아닌 자산 매각 말입니다. 그러나 이는 너무 비용이 많이 들 것이라는 우려가 있었습니다. 1조 달러의 증권을 매각하는 것은 금융 시장과 차용인에게 큰 의도치 않은 결과를 초래할 수 있었습니다. 이는 너무 위험하다고 판단되었습니다.

\subsection{3. 새로운 통화정책 도구의 도입}
대신 두 가지 대안이 채택되었습니다.
\begin{enumerate}
    \item \textbf{지급준비금 이자 (IORB):} 연방준비제도는 이제 지급준비금 잔액에 대해 이자를 지급합니다. 이 도구는 은행을 대상으로 합니다.
    \item \textbf{역환매조건부채권 시설 (ON RRP):} 연방준비제도는 역환매조건부채권 시설(Overnight Reverse Repurchase Agreement, ON RRP)을 만들었습니다. ON RRP는 머니마켓 뮤추얼 펀드와 같은 비은행 기관을 대상으로 합니다. 이러한 비은행 기관도 현금 시장의 활발한 참여자이므로 무시할 수 없습니다.
\end{enumerate}
이러한 도구들은 함께 연방준비제도가 준비금이 풍부할 때에도 단기 이자율을 통제할 수 있도록 합니다.

\subsection{4. 지급준비금 이자 (IORB)의 작동 방식}
어떻게 작동하는지 살펴보겠습니다. 지급준비금 잔액에 대한 이자부터 시작하겠습니다. 연방준비제도는 이제 지급준비금 잔액에 대해 이자를 지급합니다. 이들은 은행이 연준 계좌에 하룻밤 동안 보유하는 자금입니다. 이 자금은 은행 간 시장에서 다른 은행에 대출되지 않습니다. 이것이 통화정책 구현에 어떻게 도움이 될까요?

\begin{mybox}
  \textbf{IORB의 금리 하한선 역할}
  연방준비제도가 지급준비금 잔액에 대한 이자를 목표 연방기금금리와 동일하게 설정한다고 가정해 봅시다. 이제 은행들은 연준에 하룻밤 동안 투자하는 것과 은행 간 대출을 하는 것 사이에서 무관심해집니다. 어느 쪽이든 목표 연방기금금리를 벌게 될 것입니다. 이는 균형 상태에서 반드시 그래야 합니다. 만약 연방기금금리가 IORB보다 높다면, 은행들은 준비금을 은행 간 시장에 대출할 것입니다. 이러한 차익 거래 과정은 두 이자율이 같아질 때까지 계속되어야 합니다. 실제로 지급준비금 이자는 연방기금금리의 하한선을 제공합니다.
\end{mybox}

\begin{figure}[h!]
  \centering
  % 실제 이미지가 없으므로 텍스트로 설명합니다.
  \begin{tcolorbox}[mybox, title={\textbf{IORB가 연방기금금리에 미치는 영향}}]
    여기 연방기금시장의 공급과 수요가 있습니다. 지급준비금 잔액에 이자를 지급함으로써 이는 연방기금금리에 하한선을 설정할 것입니다. 균형 상태에서 어떤 은행도 은행 간 시장에 대출하여 IORB보다 적게 받으려 하지 않을 것입니다.
    \textbf{그래프 설명:} 준비금 수요 곡선이 IORB 수준에서 평평해지는 구간을 보여줍니다. 이는 IORB가 연방기금금리의 효과적인 하한선으로 작용함을 나타냅니다.
  \end{tcolorbox}
  \caption{IORB가 연방기금금리에 미치는 영향 설명}
  \label{fig:iorb_fed_funds_market}
\end{figure}

이는 연방준비제도가 준비금 공급을 변경하지 않고도 연방기금금리를 인상할 수 있도록 합니다.

\subsection{5. 역환매조건부채권 시설 (ON RRP)의 작동 방식}
ON RRP 시설은 왜 필요할까요? 어떻게 작동할까요? 미국에서는 비은행 기관이 현금 시장에서 중요한 역할을 합니다. 예를 들어, 머니마켓 펀드(money market funds)가 있습니다. 이러한 대출 기관은 IORB에 접근할 수 없기 때문에 때로는 더 낮은 금리로 대출할 의향이 있을 수 있습니다. 이는 유휴 현금에 대해 아무것도 받지 않는 것보다 낫습니다. 이는 단기 이자율이 IORB 아래로 떨어질 수 있음을 의미합니다. 이러한 효과에 대처하기 위해 연방준비제도는 ON RRP 시설을 만들었습니다. 이는 이러한 비은행 기관을 위한 IORB와 동일합니다. 그들은 환매조건부채권(repurchase agreement)을 통해 연준에 하룻밤 동안 투자할 수 있습니다. 이 시설은 연준이 단기 이자율에 하한선을 설정할 수 있도록 보장합니다. 이러한 비은행 기관은 ON RRP 금리보다 적게 대출하려 하지 않을 것입니다.

\begin{figure}[h!]
  \centering
  % 실제 이미지가 없으므로 텍스트로 설명합니다.
  \begin{tcolorbox}[mybox, title={\textbf{ON RRP 금리와 연방기금금리 추이}}]
    이 차트는 ON RRP 금리와 연방기금금리를 함께 보여줍니다. 보시다시피 두 금리는 잘 추적됩니다. 이는 ON RRP가 연방준비제도가 연방기금금리를 인상하는 데 도움이 될 수 있음을 의미합니다. 기억하십시오, 이 시기는 시스템에 막대한 양의 은행 준비금이 있었던 때입니다. 이러한 조건에서는 전통적인 공개시장운영이 작동하지 않았을 것입니다.
  \end{tcolorbox}
  \caption{ON RRP 금리와 연방기금금리 추이 설명}
  \label{fig:on_rrp_fed_funds_rate}
\end{figure}

\begin{summarybox}
  \textbf{핵심 정리: 풍부한 준비금 환경에서의 통화정책}
  \begin{itemize}
    \item 준비금이 많을 때, 전통적인 통화정책 도구는 이자율에 영향을 미치지 못할 수 있습니다.
    \item 연방준비제도는 연방기금금리에 하한선을 제공하기 위해 IORB와 ON RRP라는 새로운 도구를 만들었습니다.
  \end{itemize}
\end{summarybox}


\section{빠르게 훑어보기 (1페이지 요약)}

\begin{tcolorbox}[summarybox, title={\textbf{중앙은행 통화정책 도구 요약}}]
  \textbf{1. 공개시장운영 (OMO)}
  \begin{itemize}
    \item \textbf{개념:} 중앙은행이 증권(주로 국채)을 사고팔아 은행 준비금의 양을 조절.
    \item \textbf{목표:} 연방기금금리(Fed Funds Rate)를 목표 수준으로 유도.
    \textbf{작동:} 증권 매입 $\rightarrow$ 준비금 증가 $\rightarrow$ 연방기금금리 하락 (공급 곡선 우측 이동).
    \item \textbf{변화:} 2008년 금융위기 이후 준비금 풍부로 전통적 OMO 효과 감소.
  \end{itemize}

  \textbf{2. 재할인창구 (Discount Window)}
  \begin{itemize}
    \item \textbf{개념:} 중앙은행이 은행에 단기 현금 대출 제공 (최후의 대출자 역할).
    \item \textbf{목표:} 은행 유동성 지원, 뱅크런 방지, 연방기금금리 상한선 설정.
    \textbf{작동:} 재할인율(Discount Rate)을 연방기금금리 목표 상단에 설정. 연방기금금리가 재할인율 초과 시 은행은 재할인창구 이용 $\rightarrow$ 연방기금금리 하향 압력.
    \item \textbf{종류:} 주요 신용(Primary Credit)이 가장 중요 (건전한 은행 대상, 단기, 담보 대출).
  \end{itemize}

  \textbf{3. 비전통적 통화정책 (Unconventional Monetary Policy)}
  \begin{itemize}
    \item \textbf{등장 배경:} 2008년 금융위기 이후 단기 금리가 0에 가까워지면서 전통적 정책의 한계 발생.
    \item \textbf{포워드 가이던스 (Forward Guidance):}
      \begin{itemize}
        \item \textbf{개념:} 중앙은행이 미래 통화정책 의도(금리, 경제 전망 등)를 시장에 소통.
        \item \textbf{목표:} 시장 기대 형성, 불확실성 감소, 장기 금리 영향.
        \textbf{예시:} "상당히 오랜 기간" (extended period), 특정 경제 조건(실업률, 인플레이션)에 따른 금리 유지 약속 (조건부 가이던스), 닷 플롯(dot plots)을 통한 금리 전망 공개.
      \end{itemize}
    \item \textbf{대규모 자산 매입 (LSAPs / QE):}
      \begin{itemize}
        \item \textbf{개념:} 중앙은행이 장기 채권(국채, MBS 등)을 대량으로 매입.
        \item \textbf{목표:} 장기 채권 가격 상승 $\rightarrow$ 장기 수익률 하락 $\rightarrow$ 가계/기업 차입 비용 감소 $\rightarrow$ 투자/경제 활동 자극.
        \textbf{예시:} QE1, QE2, QE3, 오퍼레이션 트위스트, COVID-19 팬데믹 대응. 연준 대차대조표의 극적인 확장 초래.
      \end{itemize}
  \end{itemize}

  \textbf{4. 풍부한 준비금 환경에서의 통화정책}
  \begin{itemize}
    \item \textbf{문제점:} QE로 인해 은행 준비금이 과도하게 많아지면서, 전통적 OMO로는 금리 조절이 어려워짐.
    \item \textbf{지급준비금 이자 (IORB):}
      \begin{itemize}
        \item \textbf{개념:} 중앙은행이 은행이 보유한 준비금에 대해 이자를 지급.
        \item \textbf{목표:} 연방기금금리의 하한선 설정 (은행은 IORB보다 낮은 금리로 대출하지 않음).
      \end{itemize}
    \item \textbf{역환매조건부채권 시설 (ON RRP):}
      \begin{itemize}
        \item \textbf{개념:} 중앙은행이 비은행 금융기관(머니마켓 펀드 등)으로부터 증권을 담보로 돈을 빌리고 이자를 지급.
        \item \textbf{목표:} 비은행 기관의 단기 금리 하한선 설정 (ON RRP 금리보다 낮은 금리로 대출하지 않음).
      \end{itemize}
    \item \textbf{결과:} IORB와 ON RRP는 풍부한 준비금 환경에서도 연방기금금리를 효과적으로 통제할 수 있는 새로운 도구로 기능.
  \end{itemize}
\end{tcolorbox}


%=======================================================================
% Module 7: 개요: 중앙은행과 통화정책 모듈 3
%=======================================================================
\chapter{개요: 중앙은행과 통화정책 모듈 3}
\label{ch:module7}

\metainfo{중앙은행과 통화정책}{Lecture 3}{Prof. Ralf Meisenzahl}{통화정책이 경제에 미치는 영향을 심층적으로 분석하고, 연방준비제도(Fed)의 이중 책무인 최대 고용과 물가 안정 달성을 위한 통화정책의 작동 방식, 인플레이션과 실업률 간의 상충 관계, 통화정책 전달 메커니즘, 그리고 통화정책의 한계점을 다룹니다.}


\section{개요: 중앙은행과 통화정책 모듈 3}

이 문서는 '중앙은행과 통화정책' 모듈 3의 핵심 내용을 담고 있으며, 통화정책이 경제에 미치는 영향을 심층적으로 분석합니다. 연방준비제도(Fed)의 이중 책무인 최대 고용과 물가 안정 달성을 위한 통화정책의 작동 방식, 인플레이션과 실업률 간의 상충 관계, 통화정책 전달 메커니즘, 그리고 통화정책의 한계점을 다룹니다. 특히 비전공자도 쉽게 이해할 수 있도록 각 개념에 대한 상세한 설명과 현실적인 예시를 제공하여 시험 대비 및 심화 학습에 도움을 줍니다.


\section{용어 정리}

다음 표는 본 문서에서 사용되는 주요 경제 및 통화정책 관련 용어들을 정리한 것입니다. 각 용어의 쉬운 설명과 원어를 함께 제시하여 이해를 돕습니다.

\begin{adjustbox}{width=\textwidth,center}
\begin{tabular}{lllc}
\toprule
\textbf{용어} & \textbf{쉬운 설명} & \textbf{원어} & \textbf{비고} \\
\midrule
\textbf{이중 책무} & 중앙은행의 두 가지 주요 목표: 물가 안정과 최대 고용 & Dual Mandate & Fed의 핵심 목표 \\
\textbf{GDP} & 한 국가에서 생산된 모든 재화와 서비스의 총 가치 & Gross Domestic Product & 경제 규모 측정 \\
\textbf{실질 GDP} & 인플레이션을 조정한 GDP & Real GDP & 실제 경제 성장률 반영 \\
\textbf{실업률} & 경제활동인구 중 실업자의 비율 & Unemployment Rate & 고용 상태 측정 \\
\textbf{자연 실업률} & 경제가 완전 고용 상태일 때도 존재하는 실업률 & Natural Rate of Unemployment & 마찰적/구조적 실업 포함 \\
\textbf{오쿤의 법칙} & 실업률과 GDP 성장률 간의 음의 상관관계 & Okun's Law & 통화정책의 지침 \\
\textbf{인플레이션} & 물가 수준의 지속적인 상승 & Inflation & 화폐 가치 하락 \\
\textbf{CPI} & 소비자가 구매하는 상품 및 서비스 가격의 변화 & Consumer Price Index & 소비자 물가 변동 측정 \\
\textbf{PPI} & 생산자가 판매하는 상품 및 서비스 가격의 변화 & Producer Price Index & 생산자 물가 변동 측정 \\
\textbf{필립스 곡선} & 인플레이션과 실업률 간의 상충 관계 & Phillips Curve & 단기적 정책 딜레마 \\
\textbf{인플레이션 기대} & 미래 인플레이션에 대한 경제 주체들의 예상 & Inflation Expectations & 필립스 곡선 이동 요인 \\
\textbf{생산량 갭} & 실제 GDP와 잠재 GDP의 차이 & Output Gap / GDP Gap & 경제 과열/침체 지표 \\
\textbf{인플레이션 갭} & 실제 인플레이션과 중앙은행 목표 인플레이션의 차이 & Inflation Gap & 물가 안정 목표 달성 여부 \\
\textbf{테일러 준칙} & 중앙은행이 금리를 설정하는 규칙 & Taylor Rule & 인플레이션 및 생산량 갭 반영 \\
\textbf{연방 기금 금리} & 은행 간 초단기 대출 금리 & Federal Funds Rate & Fed의 주요 정책 수단 \\
\textbf{순이자 마진} & 은행의 대출 이자 수입과 예금 이자 비용의 차이 & Net Interest Margin & 은행 수익성 지표 \\
\textbf{공급 충격} & 생산 비용에 갑작스러운 변화를 주는 사건 & Supply Shocks & 스태그플레이션 유발 가능 \\
\textbf{스태그플레이션} & 경기 침체와 높은 인플레이션이 동시에 발생하는 현상 & Stagflation & 통화정책의 큰 도전 \\
\bottomrule
\end{tabular}
\end{adjustbox}


\section{모듈 3: 통화정책의 경제적 영향}

\subsection{레슨 3-0.1: 개요}

안녕하세요, 여러분. 어디에서 참여하시든 환영합니다. 연방준비제도(Federal Reserve, Fed)의 목적은 \textbf{최대 고용(maximum employment)}과 \textbf{물가 안정(stable prices)}을 달성하는 것입니다. 이를 \textbf{Fed의 이중 책무(dual mandate)}라고 부릅니다. 이상적인 세상에서는 경제가 최소한의 인플레이션을 겪고 모든 사람이 직업을 가질 수 있겠지만, 현실은 항상 그렇지 않습니다.

그렇다면 Fed의 통화정책 조치가 경제에 어떻게 영향을 미칠까요? 이것이 이 모듈의 목표입니다. 우리는 통화정책 조치와 가장 중요한 거시경제적 결과물들 사이의 연관성을 살펴볼 것입니다. 또한 이러한 변화가 어떻게 발생하는지도 논의할 것입니다.

\begin{conceptbox}
\textbf{Fed의 이중 책무 (Dual Mandate)}
\begin{itemize}
    \item \textbf{최대 고용 (Maximum Employment)}: 가능한 한 많은 사람들이 일자리를 가질 수 있도록 하는 것입니다. 단순히 실업률을 0으로 만드는 것이 아니라, 경제가 건강하게 작동할 때 자연적으로 발생하는 실업률 수준을 의미합니다.
    \item \textbf{물가 안정 (Stable Prices)}: 물가가 급격하게 오르거나 내리지 않고 안정적인 수준을 유지하는 것입니다. 이는 인플레이션이 너무 높거나 디플레이션이 발생하지 않도록 관리하는 것을 포함합니다.
\end{itemize}
이 두 가지 목표는 때때로 상충될 수 있어, Fed는 이 둘 사이에서 균형을 찾아야 합니다.
\end{conceptbox}

우리는 최대 고용이라는 아이디어부터 시작할 것입니다. 이중 책무는 일자리 측면에서 명시되어 있지만, 연방준비제도를 포함한 많은 중앙은행들은 \textbf{경제 성장(economic growth)}에 대해 많이 이야기합니다. 그 이유를 설명하기 위해 먼저 \textbf{국내총생산(Gross Domestic Product, GDP)}과 실업률 간의 관계를 살펴볼 것입니다.

그 다음, Fed의 목표 금리 인하가 어떻게 GDP를 증가시키고 결과적으로 실업률을 감소시킬 수 있는지 알아볼 것입니다. 이 논의의 일환으로, 미국 맥락에서 최대 고용의 목표가 무엇을 의미하는지 살펴볼 것입니다.

다음으로, 실업률, 임금, 인플레이션 간의 연관성을 연구할 것입니다. 실업률과 인플레이션 사이에는 \textbf{상충 관계(trade-off)}가 있다는 것을 알게 될 것입니다. 통화정책을 완화함으로써 중앙은행은 적어도 단기적으로는 실업률을 낮출 수 있지만, 이는 더 높은 인플레이션이라는 대가를 치를 수 있습니다. 실업률과 인플레이션 간의 이러한 상충 관계는 연방준비제도의 의사 결정에 핵심적인 요소입니다. 우리는 이 상충 관계가 유명한 \textbf{테일러 준칙(Taylor Rule)}과 같은 통화정책 규칙에 어떻게 내재되어 있는지 살펴볼 것입니다. 이러한 규칙은 경제 상황에 따라 금리가 정확히 어떻게 설정되어야 하는지를 지시합니다.

인플레이션과 실업률 간의 상충 관계를 염두에 두고, 통화정책이 어떻게 작동하는지, 즉 금리 변화가 경제 성장, 그리고 결과적으로 인플레이션과 실업률에 어떤 경로를 통해 영향을 미치는지 연구할 것입니다. 우리는 통화정책이 경제 전체에 미치는 영향을 보여주는 개괄적인 개요부터 시작할 것입니다. 그 다음, 통화정책이 은행의 신용 공급에, 그리고 결과적으로 경제의 통화 공급에 어떻게 영향을 미치는지 자세히 살펴볼 것입니다.

다음 레슨에서는 통화정책이 기업과 가계에 어떻게 영향을 미치고, 이것이 다시 투자 및 소비 결정에 어떤 영향을 미치는지 논의할 것입니다. 소비와 투자는 모두 GDP의 일부이며, 이러한 결정은 실업률에 영향을 미칠 것입니다. 통화정책은 경기 확장기에는 금리를 인상하고 경기 수축기에는 금리를 인하함으로써 경제 주기의 해로운 호황과 불황을 완화할 수 있습니다. 이것이 중앙은행이 달성하고자 하는 목표입니다.

마지막 두 레슨은 통화정책의 한계를 보여주는 사례 연구입니다. 우리는 통화정책이 유가 급등과 같은 공급 측면 충격(supply-side shocks)을 상쇄하는 데 거의 할 수 없다는 것을 알게 될 것입니다. 또한 구조적 경제 환경 변화로 인해 영구적으로 쇠퇴하는 산업의 고용을 지원하는 데도 많은 것을 할 수 없습니다.

요약하자면, 이전 모듈에서는 중앙은행이 사용할 수 있는 통화정책 도구와 그것이 금리에 어떻게 영향을 미치는지 논의했습니다. 이 모듈에서는 금리 변화가 경제에 어떻게 영향을 미치고, 이것이 연방준비제도가 물가 안정과 최대 고용이라는 이중 책무를 선택하는 데 어떻게 도움이 되는지 살펴봅니다.


\section{레슨 3-1: 인플레이션과 실업률 간의 연관성}

\subsection{레슨 3-1.1: 경제 성장과 실업률 이해하기}

안녕하세요, 경제 성장과 실업률에 대한 강의에 오신 것을 환영합니다. 연방준비제도는 물가 안정과 최대 고용을 달성하기 위해 통화정책을 사용한다는 것을 기억하십시오. 하지만 최대 고용이란 정확히 무엇을 의미할까요? 이 수업에서는 실업률을 살펴볼 것입니다. 우리는 실업률에 대한 몇 가지 새로운 정의를 소개할 것입니다: \textbf{자연 실업률(natural rate of unemployment)}, \textbf{완전 고용(full employment)}, \textbf{최대 고용(maximum employment)}. 우리는 이 용어들이 무엇을 의미하고 왜 필요한지 이해할 것입니다. 또한 \textbf{국내총생산(Gross Domestic Product, GDP)}을 정의할 것입니다. 이것은 한 경제의 총 생산량을 측정하는 지표입니다. 우리는 데이터에서 실업률과 GDP 성장률 간의 관계를 살펴볼 것입니다. 마지막으로, 연방준비제도가 실업률을 0으로 줄이려고 하지 않는 이유를 이해할 것입니다.

\subsubsection*{실업률 데이터}

실업률 데이터부터 시작해 봅시다. 개인은 직업이 없고 구직 활동을 하고 있다면 실업 상태입니다. \textbf{실업률(unemployment rate)}은 경제활동인구 중 현재 실업 상태인 사람들의 비율입니다.

\begin{examplebox}
\textbf{실업률의 의미}
만약 한 국가의 경제활동인구가 100명이고, 이 중 5명이 직업이 없지만 적극적으로 일자리를 찾고 있다면, 실업률은 5\%가 됩니다. 나머지 95명은 고용 상태이거나, 아예 구직 활동을 하지 않는 비경제활동인구(예: 학생, 은퇴자)입니다.
\end{examplebox}

그래프를 보면 현대 시대 미국의 실업률을 알 수 있습니다. 그래프의 음영 처리된 부분은 \textbf{경기 침체(recessions)}를 나타냅니다. 경기 침체기에는 실업률이 급격히 증가하는 것을 볼 수 있습니다. 예를 들어, 1973년 10월부터 1975년 5월까지 실업률은 4.6\%에서 9\%로 거의 두 배가 되었습니다. 동시에, 경기 침체가 끝난 후에는 실업률이 빠르게 하락합니다. 1976년 3월에는 실업률이 이미 7.6\%로 떨어졌습니다.

\subsubsection*{경제 성장 측정: GDP}

이제 경제 성장의 측정 지표인 \textbf{국내총생산(Gross Domestic Product, GDP)}으로 넘어가겠습니다. GDP는 한 경제가 생산하는 모든 재화와 서비스의 총 달러 가치입니다. \textbf{실질 GDP(Real GDP)}는 이 측정치를 인플레이션에 대해 조정합니다. 미국의 실질 GDP는 시간이 지남에 따라 성장해 왔습니다. 그러나 경기 침체기에는 감소하며, 이는 경제가 수축한다는 것을 의미합니다. 데이터를 통해 이를 살펴보겠습니다.

그래프는 실질 GDP 성장률의 분기별 퍼센트 변화를 보여줍니다. 음영 처리된 경기 침체기에는 실질 GDP 성장률이 음수입니다. 1974년 각 분기 동안 GDP가 1\%만큼 감소한 것을 볼 수 있습니다.

\subsubsection*{실업률과 실질 GDP 성장률의 관계: 오쿤의 법칙}

실업률과 실질 GDP 성장률 간의 관계는 무엇일까요? 경기 침체기에는 실업률이 상승하고 GDP는 하락합니다. 경기 확장기에는 그 반대가 일어납니다. 이는 \textbf{음의 상관관계(negative correlation)}를 시사합니다.

이것이 어떻게 작동하는지 보기 위해 이전 두 차트를 결합해 봅시다. 각 점은 한 분기를 나타냅니다. X축에는 이전 분기 대비 실업률 증가율이 있고, Y축에는 해당 분기의 GDP 성장률이 있습니다. 데이터에서 음의 상관관계를 볼 수 있으며, 이는 아래로 기울어진 빨간색 선으로 표시됩니다. 역사적인 미국 데이터에 따르면, 이 선의 기울기는 약 -2입니다. 이 경험적 사실은 경제학자 아서 오쿤(Arthur Okun)에 의해 처음 언급되었으며, \textbf{오쿤의 법칙(Okun's Law)}으로 알려져 있습니다.

\begin{conceptbox}
\textbf{오쿤의 법칙 (Okun's Law)}
실업률과 실질 GDP 성장률 사이에는 음의 상관관계가 있다는 경제학적 원리입니다. 즉, 실업률이 1\% 감소하면 실질 GDP 성장률은 약 2\% 증가하는 경향이 있습니다 (미국 데이터 기준).

\textbf{왜 중요한가요?}
통화정책 입안자들은 오쿤의 법칙을 유용하게 활용합니다. 경제를 부양하기 위해 금리를 인하할 때 실업률이 얼마나 감소할지에 대한 지침을 제공하기 때문입니다. 금리 인하 후 실질 GDP가 얼마나 증가할지에 대한 추정치가 있다면, 실업률이 얼마나 감소할지도 알 수 있습니다. 이러한 거시경제 변수들 간의 관계를 이해하면 Fed가 통화정책을 더 잘 목표로 삼을 수 있습니다.
\end{conceptbox}

이제 다른 질문을 해봅시다.

\subsubsection*{자연 실업률: 실업률을 0으로 만들 수 없는 이유}

여러분은 중앙은행이 실업률을 0으로 줄일 수 없는 이유가 무엇인지 궁금할 수 있습니다. 답은 \textbf{자연 실업률(natural rate of unemployment)}이 존재하기 때문입니다. 이는 경제가 최고 성과를 내고 있을 때조차도 일부 근로자는 여전히 실업 상태에 있을 것이라는 의미입니다.

\begin{conceptbox}
\textbf{자연 실업률 (Natural Rate of Unemployment)}
경제가 완전 고용 상태에 있을 때도 존재하는 실업률입니다. 이는 경기 변동과 무관하게 발생하는 실업으로, 주로 다음과 같은 이유로 발생합니다.
\begin{itemize}
    \item \textbf{마찰적 실업 (Frictional Unemployment)}: 사람들이 새로운 직업을 찾거나, 교육을 마치고 첫 직장을 구하거나, 이사 등으로 인해 직업을 전환하는 데 시간이 걸리기 때문에 발생합니다.
    \item \textbf{구조적 실업 (Structural Unemployment)}: 기술 변화, 산업 구조 변화 등으로 인해 특정 기술을 가진 근로자들이 일자리를 잃고 새로운 기술을 배우거나 다른 산업으로 이동해야 할 때 발생합니다. (예: 자율 주행 기술이 트럭 운전사에게 미칠 영향)
    \item \textbf{제도적 요인}: 높은 실업 수당은 실업자들이 더 나은 일자리를 찾는 데 더 많은 시간을 할애하도록 유도할 수 있습니다.
\end{itemize}
따라서 중앙은행은 실업률을 0으로 만드는 것을 목표로 하지 않습니다. 대신 자연 실업률 수준에 도달하는 것을 '최대 고용'으로 간주합니다.
\end{conceptbox}

다시 미국의 실업률 시계열을 살펴보겠습니다. 대략적으로 말해, 우리는 좋은 경제 시기의 실업률을 기준으로 자연 실업률을 추정할 수 있습니다.

\begin{table}[h!]
\centering
\caption{미국 자연 실업률 추정치 (역사적 데이터)}
\label{tab:natural_unemployment_rates}
\begin{tabular}{lc}
\toprule
\textbf{기간} & \textbf{자연 실업률 추정치} \\
\midrule
1950년대 및 1960년대 & 4\% \\
1970년 ~ 1983년 & 6\% \\
1984년 ~ 2000년 & 5.5\% \\
\bottomrule
\end{tabular}
\end{table}

차트는 또한 실업률의 큰 주기적 변동을 보여줍니다. 특히 경기 침체기에는 실업률이 자연 실업률을 초과합니다. 바로 이 지점에서 통화정책이 효과적일 수 있습니다.

\subsubsection*{오늘날의 자연 실업률}

마지막 질문입니다. 오늘날의 자연 실업률은 얼마일까요? 새천년의 실업률을 살펴보겠습니다. 2001년 경기 침체 이후 실업률은 20세기 후반의 경기 침체에 비해 느리게 회복되었습니다. 마찬가지로, 2008년 금융 위기 이후 실업률은 예상만큼 빠르게 회복되지 않았습니다. 이는 자연 실업률이 증가했을 것이라는 추측을 불러일으켰습니다. 그러나 미국에서 가장 긴 경제 확장기가 끝나갈 무렵인 2019년 12월에는 실업률이 3.6\%였습니다. 마찬가지로, 코로나19 팬데믹 이후에도 실업률은 역사적으로 낮았습니다. 이는 자연 실업률이 감소했을 수 있음을 시사합니다. 전반적으로, 새로운 기술이나 노동 시장의 변화가 역할을 할 수 있지만, 우리가 결정하기까지는 더 많은 시간이 필요합니다.

\begin{summarybox}
\textbf{레슨 3-1.1 요약}
\begin{itemize}
    \item \textbf{오쿤의 법칙 (Okun's Law)}: 실업률 변화와 실질 GDP 변화는 데이터에서 음의 상관관계를 보입니다. 이 관계는 통화정책 결정에 중요한 지침이 됩니다.
    \item \textbf{자연 실업률 (Natural Rate of Unemployment)}: 경제가 완전한 생산 능력을 발휘할 때조차도 실업률은 0보다 클 것입니다. 이는 마찰적 실업과 구조적 실업 등 경제의 자연스러운 부분입니다.
\end{itemize}
\end{summarybox}


\subsection{레슨 3-1.2: 실업률과 인플레이션 연결하기}

안녕하세요, 임금, 실업률, 인플레이션 간의 관계에 대한 강의에 오신 것을 환영합니다. \textbf{인플레이션 목표제(inflation targeting)}는 통화정책의 핵심입니다. 이 수업에서는 \textbf{노동 비용(labor costs)} 또는 \textbf{임금 인플레이션(wage inflation)}이 \textbf{소비자 물가 인플레이션(consumer price inflation)}으로 어떻게 이어지는지 논의할 것입니다. 우리는 임금과 인플레이션이 어떻게 함께 움직이는지 살펴볼 것입니다. 그 다음, 인플레이션과 실업률 간의 관계를 검토할 것입니다. 이를 위해 \textbf{필립스 곡선(Phillips Curve)}이라는 아이디어를 소개할 것입니다. 이것은 경제학에서 가장 유명한 차트 중 하나입니다. 인플레이션과 실업률 간의 상충 관계를 어떻게 설명하는지 알게 될 것입니다. 일반적으로 인플레이션이 낮아지면 실업률이 높아지는 대가를 치르게 됩니다. 마지막으로, \textbf{인플레이션 기대(inflation expectations)}의 변화가 필립스 곡선을 어떻게 이동시킬 수 있는지 이해할 것입니다.

\subsubsection*{소비자 물가 인플레이션의 원인}

경제학자, 정책 입안자, 언론은 항상 인플레이션에 대해 이야기합니다. 그들은 보통 \textbf{소비자 물가 인플레이션(consumer price inflation)} 또는 \textbf{CPI(Consumer Price Index)}를 언급합니다. 이는 소비자 상품 묶음의 가격 상승을 의미합니다.

여러분은 소비자 물가가 왜 오르는지 궁금할 수 있습니다. 한 가지 답은 \textbf{투입 가격(input prices)}입니다. 원유 가격이 높으면 자동차 연료비가 더 비싸다는 것을 우리 모두 알고 있습니다. 더 일반적으로, 투입 가격이 상승하면 상품 생산 및 서비스 제공 비용도 증가합니다. 다르게 말하면, \textbf{생산자 물가 인플레이션(producer price inflation)}이 높아지면 CPI가 증가할 수 있습니다. 이에 대한 데이터를 살펴보겠습니다.

그래프는 1980년대 이후 소비자 물가 인플레이션과 생산자 물가 인플레이션을 보여줍니다. 생산자 물가 지수(Producer Price Index, PPI)가 더 변동성이 크지만, 두 시리즈가 서로를 추적하는 것을 볼 수 있습니다. 생산자 물가가 높아지면 소비자 상품 가격도 높아집니다.

대부분 산업의 가장 큰 비용 구성 요소는 무엇일까요? 바로 \textbf{노동 비용(labor expenses)}입니다. 이 그래프에 \textbf{고용 비용 지수(employment cost index)}를 추가해 봅시다. 이 고용 비용 지수에는 임금과 복리후생이 모두 포함됩니다. 이 시리즈는 2000년에 시작되므로 시간 프레임을 제한해야 합니다. 우리는 다시 고용 비용과 소비자 물가 인플레이션이 서로를 추적하는 것을 봅니다. 이는 임금 상승이 소비자 물가 인플레이션으로 이어진다는 것을 시사합니다.

\subsubsection*{임금 인플레이션과 실업률: 필립스 곡선}

이러한 추론을 한 단계 더 나아가 봅시다. 임금은 언제 가장 많이 오를까요? 무엇이 임금 인플레이션을 유발할까요? 임금은 보통 직원들이 가장 많은 \ bargaining power를 가질 때 오릅니다. 이는 실업률이 낮고 일자리 경쟁이 적을 때입니다. 실업률이 높을 때는 고용주가 구직자를 유치하기 위해 더 높은 임금을 지불할 필요가 없습니다.

이 관계는 아서 필립스(Arthur Phillips)에 의해 처음 언급되었습니다. 그는 영국에서 실업률과 임금 상승률을 그래프로 그렸습니다. 그는 두 결과물 사이에 \textbf{음의 상관관계(negative correlation)}를 발견했습니다.

\begin{conceptbox}
\textbf{필립스 곡선 (Phillips Curve)}
원래는 실업률과 임금 상승률 간의 음의 상관관계를 나타냈지만, 오늘날에는 주로 \textbf{실업률과 소비자 물가 인플레이션 간의 상충 관계(trade-off)}를 설명하는 데 사용됩니다.
\begin{itemize}
    \item \textbf{실업률이 낮을 때}: 노동 시장이 타이트해져 기업들이 인력을 확보하기 위해 더 높은 임금을 지불해야 합니다. 이는 생산 비용을 증가시키고, 결국 소비자 물가 상승(인플레이션)으로 이어집니다.
    \item \textbf{실업률이 높을 때}: 노동 시장에 여유가 있어 기업들이 임금을 크게 올릴 필요가 없습니다. 이는 인플레이션 압력을 낮춥니다.
\end{itemize}
\end{conceptbox}

실업률이 낮을 때는 임금이 오르는 경향이 있습니다. 실업률이 높을 때는 임금이 줄어드는 경향이 있습니다. 이 관계를 종종 \textbf{필립스 곡선(Phillips Curve)}이라고 부릅니다. 오늘날에는 필립스 곡선이라는 용어가 실업률과 소비자 물가 인플레이션 간의 관계를 설명하는 데 사용됩니다. 우리는 앞서 CPI와 임금 상승률이 서로를 추적한다는 것을 보았으므로 이는 이해할 수 있습니다. 이제 최근 몇 년간 미국의 필립스 곡선을 살펴보겠습니다.

여기 1960년대 데이터가 있습니다. 이전과 마찬가지로 X축에는 실업률이 있지만, 이제 Y축에는 소비자 물가 인플레이션이 있습니다. 인플레이션과 실업률 간의 음의 관계가 유지되는 것을 볼 수 있습니다.

\subsubsection*{정책 입안자에게 미치는 영향}

왜 우리는 신경 써야 할까요? 정책 입안자에게 미치는 영향은 무엇일까요? 이 그래프는 실업률과 인플레이션 간의 명확한 상충 관계를 시사합니다. 이는 실업률을 줄이는 모든 정부 정책이 인플레이션을 증가시킬 수 있음을 의미합니다. 이것은 통화정책에 중요합니다. 물가 안정과 최대 고용이 연방준비제도의 책무라는 것을 기억하십시오. 이 목표들은 서로 상충됩니다. 두 목표를 동시에 달성하는 것은 균형을 잡는 행위입니다. 실제로 Fed의 주요 의사 결정자들이 필립스 곡선에 대해 생각하는 것을 볼 수 있습니다.

예를 들어, 2008년 당시 연방준비제도 이사회 부의장이었던 도널드 콘(Donald Cohn)은 "인플레이션의 경제 모델은 통화정책 심의에 필수적인 투입물이다. 필립스 곡선 전통의 모델은 대부분의 학술 연구자들과 정책 입안자들, 나를 포함하여, 인플레이션 변동에 대해 생각하는 방식의 핵심으로 남아 있다. 실제로 대안적인 프레임워크는 견고한 경제적 기반과 경험적 지지를 결여하고 있는 것으로 보인다"고 말했습니다.

\subsubsection*{필립스 곡선의 이동: 인플레이션 기대의 역할}

우리는 끝났을까요? 글쎄요, 아직 아닙니다. 1960년대와 1970년대 초에는 실업률과 인플레이션 간의 관계가 안정적이라고 생각되었습니다. 그러나 다음 두 십 년 동안 다른 패턴이 나타났습니다. 여기 1970년대와 1980년대 데이터가 있습니다. 불행히도 음의 상관관계는 사라진 것처럼 보입니다.

그러나 조금 자세히 보면, 일반적인 패턴이 나타나는 것 같습니다. 하나의 관계가 아니라 세 가지 관계가 있는 것처럼 보입니다. 1970년대 초, 1970년대 후반, 그리고 1980년대 초에 각각 하나씩 말입니다. 세 개의 필립스 곡선이 있는 것처럼 보입니다. 이상하죠? 이 곡선들은 시간이 지남에 따라 바깥쪽으로 이동하는 것처럼 보입니다. 바깥쪽으로 이동하는 것은 좋은 징조가 아닙니다. 필립스 곡선이 바깥쪽으로 더 많이 이동할수록 인플레이션을 억제하기 위해 더 많은 실업률이 필요합니다. 예를 들어, 인플레이션을 6\%로 유지하고 싶다고 가정해 봅시다. 1973년에는 5\%의 실업률이 필요했을 것입니다. 그러나 1982년에는 동일한 인플레이션율을 달성하기 위해 실업률이 10\%에 도달했을 것입니다.

1970년대와 1980년대는 인플레이션과 실업률 간의 전반적인 관계에 의문을 제기했지만, 데이터는 적어도 단기적으로는 필립스 곡선이 작동한다는 것을 보여줍니다. 이제 1970년대와 1980년대의 이러한 단기 필립스 곡선을 결정하는 요인을 살펴보겠습니다.

1960년에서 1995년 사이 미국의 소비자 물가 인플레이션을 살펴보겠습니다. 분명히 1970년대 초, 1970년대 후반, 1980년대 초에는 소비자 물가 인플레이션이 1960년대에 비해 더 높고 증가했습니다. 따라서 필립스 곡선의 바깥쪽 이동에 대한 그럴듯한 설명은 1970년대 초부터 1980년대 초까지 \textbf{인플레이션 기대(inflation expectations)}의 상향 이동일 수 있습니다.

\begin{conceptbox}
\textbf{인플레이션 기대 (Inflation Expectations)}
미래의 인플레이션에 대한 경제 주체들(가계, 기업, 투자자)의 예상입니다. 이 기대는 실제 인플레이션에 큰 영향을 미칩니다.
\begin{itemize}
    \item \textbf{왜 중요한가요?} 만약 사람들이 미래에 물가가 많이 오를 것이라고 예상한다면, 그들은 현재 더 높은 임금을 요구하거나, 상품 가격을 미리 올리려 할 것입니다. 이는 실제 인플레이션을 더욱 부추기는 결과를 낳습니다.
    \item \textbf{필립스 곡선에 미치는 영향}: 인플레이션 기대가 상승하면, 동일한 실업률 수준에서도 더 높은 인플레이션이 발생하게 되어 필립스 곡선이 바깥쪽(오른쪽)으로 이동합니다. 이는 정책 입안자들이 인플레이션을 억제하기 위해 더 높은 실업률을 감수해야 함을 의미합니다.
\end{itemize}
\end{conceptbox}

미래 인플레이션에 대한 기대가 왜 중요할까요? 이해하기 어렵지 않습니다. 고용주가 5\% 임금 인상을 발표했다고 가정해 봅시다. 이는 큰 인상처럼 들립니다. 그러나 소비자 물가 인플레이션이 현재 10\%로 진행 중이고 내년에도 10\%로 유지될 것으로 예상된다면 어떨까요? 이는 인플레이션을 조정한 소득, 즉 \textbf{실질 임금(real wage)}이 5\% 감소한다는 것을 의미합니다. 좋은 소식이 아닙니다. 실질 임금 인상만이 생활 수준을 높입니다. 이는 근로자들이 더 높은 인플레이션을 예상할 때 더 높은 임금을 요구하게 된다는 것을 시사합니다. 이는 임금과 실업률 간의 관계를 변화시키고, 결과적으로 실업률과 소비자 물가 인플레이션 간의 관계에 연쇄적인 영향을 미칩니다.

인플레이션 기대의 증가는 1970년대와 1980년대에 우리가 보았던 단기 필립스 곡선의 바깥쪽 이동을 설명합니다. 또한 정책 입안자들이 인플레이션 기대를 면밀히 주시하는 이유를 설명합니다. 그들은 인플레이션 기대를 \textbf{잘 고정(well anchored)}시키려고 노력합니다.

\begin{summarybox}
\textbf{레슨 3-1.2 요약}
\begin{itemize}
    \item \textbf{임금과 인플레이션의 연동}: 임금 상승은 생산 비용 증가를 통해 소비자 물가 인플레이션으로 이어지는 경향이 있습니다.
    \item \textbf{필립스 곡선}: 인플레이션과 실업률 사이에는 단기적인 상충 관계가 존재합니다. 중앙은행은 이 두 목표 사이에서 균형을 찾아야 합니다.
    \item \textbf{인플레이션 기대의 중요성}: 인플레이션 기대의 변화는 필립스 곡선을 이동시킬 수 있습니다. 기대 인플레이션이 높아지면, 동일한 실업률에서도 더 높은 인플레이션이 발생하여 필립스 곡선이 바깥쪽으로 이동합니다. 중앙은행은 인플레이션 기대를 안정적으로 유지하는 것을 중요하게 생각합니다.
\end{itemize}
\end{summarybox}


\subsection{레슨 3-1.3: 통화정책 규칙}

안녕하세요, 통화정책 규칙 설계에 대한 수업에 오신 것을 환영합니다. 오늘날 중앙은행은 목표 달성을 위해 금리를 설정합니다. 연방준비제도의 경우, \textbf{연방 기금 금리(federal funds rate)}를 조정합니다. 그러나 Fed는 물가 안정과 최대 고용이라는 목표 사이에서 상충 관계에 직면합니다. 필립스 곡선을 기억하시나요? 이 세션에서는 이 상충 관계가 실제로 어떻게 접근되는지 논의할 것입니다. 우리는 \textbf{생산량 갭(output gap)}과 \textbf{인플레이션 갭(inflation gap)}이라는 아이디어를 소개할 것입니다. 그 다음, 중앙은행이 이러한 갭에 대해 무엇을 해야 하는지 질문할 것입니다. 우리는 통화정책이 경제 상황에 반응해서는 안 된다고 말하는 밀턴 프리드먼(Milton Friedman)의 \textbf{K-규칙(K-rule)}을 살펴볼 것입니다. 그리고 나서 금리를 인플레이션 및 고용의 갭 또는 변동과 연결하는 유명한 \textbf{테일러 준칙(Taylor Rule)}을 검토할 것입니다.

\subsubsection*{최대 고용과 물가 안정 목표 측정}

연방준비제도의 이중 책무, 즉 최대 고용과 물가 안정을 기억하십시오. 먼저 최대 고용에 초점을 맞추고 이 목표가 어떻게 달성될 수 있는지 설명해 봅시다. 실업률과 GDP는 밀접하게 연결되어 있습니다. 이 관계를 \textbf{오쿤의 법칙(Okun's Law)}이라고 합니다. 결과적으로 우리는 종종 경제 성장과 실업률을 서로 바꿔가며 이야기합니다.

경제가 자원을 완전히 활용하고 있을 때, 우리는 경제가 \textbf{잠재 GDP(potential GDP)}를 생산하고 있다고 말합니다. 경제가 잠재 GDP에 있을 때, 이는 \textbf{완전 고용(full employment)} 또는 \textbf{최대 고용(maximum employment)}을 의미합니다. 이것이 우리가 바랄 수 있는 최선입니다. 최대 고용을 달성하기 위해 중앙은행은 경제가 잠재 GDP에 가까운지 알아야 합니다.

\begin{conceptbox}
\textbf{잠재 GDP (Potential GDP)}
한 경제가 모든 생산 요소를 (노동, 자본, 토지 등) 정상적으로 활용하여 달성할 수 있는 최대 생산량 수준을 의미합니다. 이는 지속 가능한 성장을 나타내며, 이 수준에서는 인플레이션 압력이 크게 발생하지 않고 실업률은 자연 실업률 수준에 머무릅니다.

\textbf{생산량 갭 (Output Gap 또는 GDP Gap)}
실제 GDP와 잠재 GDP 간의 차이를 백분율로 나타낸 것입니다.
\begin{itemize}
    \item \textbf{양의 생산량 갭 (Positive Output Gap)}: 실제 GDP가 잠재 GDP보다 높을 때 발생합니다. 경제가 과열되어 자원을 과도하게 사용하고 있음을 의미하며, 인플레이션 압력이 높아질 수 있습니다. (예: 근로자들이 초과 근무를 하고 공장이 최대 생산 능력을 초과하여 가동될 때)
    \item \textbf{음의 생산량 갭 (Negative Output Gap)}: 실제 GDP가 잠재 GDP보다 낮을 때 발생합니다. 경제가 침체되어 자원이 충분히 활용되지 못하고 있음을 의미하며, 실업률이 높아질 수 있습니다.
\end{itemize}
중앙은행은 이 생산량 갭을 통화정책 결정의 중요한 지표로 활용하여 경제를 잠재 GDP 수준으로 되돌리려 합니다.
\end{conceptbox}

데이터에서 이것은 어떻게 보일까요? 차트에서 2030년까지 미국의 잠재 실질 GDP 추정치를 볼 수 있습니다. 미국 정부 기관인 의회 예산국(Congressional Budget Office)이 잠재 GDP를 추정합니다. 이것은 순전히 이론적인 구성물입니다. 기술, 교육 및 인구 통계학적 추세의 예상 개선으로 인해 시간이 지남에 따라 성장합니다.

여기 실제 실질 GDP가 있습니다. 두 가지가 밀접하게 추적되지만, 눈에 띄는 편차가 있습니다.

이 차트는 이러한 편차를 보여줍니다. 이는 실제 실질 GDP와 잠재 GDP 간의 퍼센트 포인트 차이로 계산됩니다. 이를 \textbf{생산량 갭(output gap)} 또는 \textbf{GDP 갭(GDP gap)}이라고 합니다. 생산량 갭이 양수인 시기가 있습니다. 이들은 경제가 과열되는 호황기입니다. 자원이 장기적으로 사용될 것보다 더 많이 사용됩니다. 근로자들은 초과 근무를 하고 추가 교대 근무를 하는 등입니다. 양의 갭은 덜 흔하고 규모도 작습니다.

대조적으로, 음의 생산량 갭은 크고 더 오래 지속될 수 있습니다.

이제 실업률을 겹쳐서 봅시다. 생산량 갭이 음수일 때 실업률이 평소보다 높다는 것을 알 수 있습니다. 2008년에 시작된 십 년은 명확한 예시입니다. 여기서 핵심 통찰력은 무엇일까요? 최대 고용이 목표 중 하나일 때 생산량 갭을 통화정책의 투입물로 사용하는 것은 매우 합리적입니다.

\subsubsection*{물가 안정 목표 측정: 인플레이션 갭}

다음으로, 연방준비제도 책무의 두 번째 부분은 물가 안정입니다. 물가 안정은 낮은 인플레이션율을 의미합니다. 하지만 얼마나 낮아야 할까요? 존 테일러(John Taylor)와 같은 학자들은 중앙은행이 2\% 목표를 사용해야 한다고 제안했습니다. 연방준비제도도 동의하는 것 같습니다. 2012년부터 연방공개시장위원회(FOMC)는 연간 2\%의 소비자 물가 인플레이션율이 연방준비제도의 법정 책무와 장기적으로 가장 일치한다고 발표했습니다. 마찬가지로 유럽중앙은행(ECB)은 2\% 미만이지만 2\%에 가까운 인플레이션을 목표로 합니다.

생산량 갭과 마찬가지로 \textbf{인플레이션 갭(inflation gap)}을 계산할 수 있습니다. 이는 인플레이션이 중앙은행의 2\% 목표율과 얼마나 다른지를 측정합니다.

\begin{conceptbox}
\textbf{인플레이션 갭 (Inflation Gap)}
실제 인플레이션율과 중앙은행이 목표로 하는 인플레이션율(예: 2\%) 간의 차이입니다.
\begin{itemize}
    \item \textbf{양의 인플레이션 갭}: 실제 인플레이션이 목표보다 높을 때 발생합니다. 중앙은행은 인플레이션을 낮추기 위해 긴축 통화정책을 고려할 수 있습니다.
    \item \textbf{음의 인플레이션 갭}: 실제 인플레이션이 목표보다 낮을 때 발생합니다. 중앙은행은 경제를 부양하고 인플레이션을 목표 수준으로 끌어올리기 위해 완화 통화정책을 고려할 수 있습니다.
\end{itemize}
\end{conceptbox}

이 그래프는 개인 소비 지출(Personal Consumption Expenditure, PCE)의 연간 성장률을 보여줍니다. 이는 소비자 물가 인플레이션의 측정치입니다. 1970년대와 1980년대에는 인플레이션이 2\% 목표보다 높았다는 것을 알 수 있습니다. 반면에 1990년대 중반부터 2020년까지는 인플레이션이 지속적으로 2\% 미만이었습니다.

\subsubsection*{통화정책 규칙: K-규칙과 테일러 준칙}

중앙은행은 실업률과 인플레이션의 변동에 어떻게 대응해야 할까요? 전혀 대응하지 않아야 할까요? 노벨상 수상 경제학자 밀턴 프리드먼(Milton Friedman)은 중앙은행이 큰 차이를 만들 수 있다는 것에 회의적이었습니다.

대신 프리드먼은 오늘날 \textbf{K-규칙(K-rule)}으로 알려진 것을 제안했습니다. 중앙은행은 단순히 매년 일정한 K 퍼센트, 예를 들어 3\%만큼 통화 공급을 증가시켜야 합니다. 분명히 이 규칙은 경제 상황을 무시합니다. 중앙은행은 재량권이 없고 미리 프로그래밍된 알고리즘을 따를 뿐입니다. 왜 이것이 좋은 규칙일까요? 통화 공급의 대규모 증가가 인플레이션을 유발한다면, 적당한 통화 성장은 물가를 안정적으로 유지해야 합니다. 이 규칙의 또 다른 장점은 금융 시장과 기업이 모든 통화정책 조치를 완벽하게 예측하고 이에 따라 결정을 내릴 수 있다는 것입니다. 통화 공급을 통제하고 예측 가능하게 유지하는 것이 경제 및 금융 안정을 촉진할 것이라는 믿음입니다.

다른 종류의 통화정책 규칙은 경제 상황에 반응하도록 설계되었습니다. 경기 순환이 비용이 많이 든다면, 이를 완화하는 것이 사회에 이로울 수 있습니다. 가장 유명한 규칙은 소위 \textbf{테일러 준칙(Taylor Rule)}으로, 금리 결정을 생산량과 인플레이션의 변동과 연결합니다.

\begin{conceptbox}
\textbf{밀턴 프리드먼의 K-규칙 (Friedman's K-rule)}
\begin{itemize}
    \item \textbf{핵심}: 중앙은행은 매년 통화 공급을 일정한 비율(K\%)로만 증가시켜야 합니다. 경제 상황에 관계없이 기계적으로 적용됩니다.
    \textbf{장점}: 통화 공급의 예측 가능성을 높여 인플레이션을 안정시키고, 시장의 불확실성을 줄입니다. 중앙은행의 재량적 판단으로 인한 오류를 방지합니다.
    \item \textbf{단점}: 경제 상황 변화(예: 경기 침체, 공급 충격)에 유연하게 대응할 수 없어, 경기 안정화에 실패할 수 있습니다.
\end{itemize}
\end{conceptbox}

\subsubsection*{테일러 준칙 (Taylor Rule)}

테일러 준칙을 살펴보겠습니다. 이는 1990년대에 학자 존 테일러(John Taylor)에 의해 처음 제안되었습니다. 먼저, \textbf{명목 금리(nominal interest rate)}는 \textbf{실질 금리(real interest rate, R)}와 \textbf{인플레이션($\pi$)}의 합과 같다는 것을 기억하십시오. 인플레이션이 2\% 목표에 있다고 가정해 봅시다. 또한 균형 실질 금리가 0.5\%라고 가정해 봅시다. 인플레이션이 목표에 머물러 있다고 가정하면, 명목 금리는 2.5\%가 될 것입니다. 이것이 중앙은행의 개입 없이 시장이 제공할 금리입니다. 그러나 중앙은행은 원한다면 다른 명목 금리를 선택할 수 있습니다.

사실, 중앙은행이 인플레이션 갭과 생산량 갭을 가중치 알파($\alpha$)와 베타($\beta$)로 명시적으로 고려하여 목표 금리를 선택한다고 가정해 봅시다. 이것이 여러분의 간단한 테일러 준칙입니다.

\begin{equation*}
\text{목표 금리} = \text{실질 금리} + \text{목표 인플레이션} + \alpha \times (\text{인플레이션 갭}) + \beta \times (\text{생산량 갭})
\end{equation*}

이 공식은 목표 금리가 시장이 원하는 금리에 경제 상황에 대한 조정을 더한 것이어야 한다고 말합니다.

더 정확히 말하면, 인플레이션 갭이 양수이거나 생산량 갭이 양수일 때 중앙은행은 금리를 인상해야 한다고 말합니다. 이는 경제를 식힐 것입니다.

그리고 인플레이션 갭이 음수이거나 생산량 갭이 음수일 때는 그 반대로 해야 한다고 말합니다. 이는 경제를 부양할 것입니다. 이 간단한 규칙은 연방준비제도의 의사 결정 요소를 포착합니다. 물가 안정과 최대 고용 모두 공식에 나타납니다. 물론 인플레이션 가중치 알파($\alpha$)와 생산량 갭 가중치 베타($\beta$)에 대한 값을 찾아야 합니다. 테일러의 원래 권고는 동일한 가중치를 부여하여 둘 다 0.5로 설정하는 것이었습니다. 경제 연구자들은 이후 이러한 가중치가 무엇이어야 하는지에 대한 수많은 연구를 수행했습니다. 오늘날 인플레이션 가중치 알파($\alpha$)에 대한 합의는 1.5인 반면, 생산량 갭 가중치 베타($\beta$)에 대한 추정치는 0.5에서 1 사이로 다양합니다.

\begin{examplebox}
\textbf{테일러 준칙 예시}
\begin{itemize}
    \item \textbf{기본 설정}: 실질 금리 = 0.5\%, 목표 인플레이션 = 2\%
    \item \textbf{가중치}: $\alpha = 1.5$, $\beta = 0.5$ (테일러의 원래 권고)
\end{itemize}
\textbf{시나리오 1: 경제 과열 (양의 갭)}
\begin{itemize}
    \item 실제 인플레이션 = 3\% (인플레이션 갭 = +1\%)
    \item 실제 GDP = 잠재 GDP + 2\% (생산량 갭 = +2\%)
    \item \textbf{목표 금리} = 0.5\% + 2\% + (1.5 $\times$ 1\%) + (0.5 $\times$ 2\%) \\
    = 0.5\% + 2\% + 1.5\% + 1\% = \textbf{5\%}
\end{itemize}
$\rightarrow$ 경제가 과열되어 인플레이션과 생산량이 목표보다 높으므로, 중앙은행은 금리를 5\%로 인상하여 경제를 식히려고 합니다.

\textbf{시나리오 2: 경제 침체 (음의 갭)}
\begin{itemize}
    \item 실제 인플레이션 = 1\% (인플레이션 갭 = -1\%)
    \item 실제 GDP = 잠재 GDP - 2\% (생산량 갭 = -2\%)
    \item \textbf{목표 금리} = 0.5\% + 2\% + (1.5 $\times$ -1\%) + (0.5 $\times$ -2\%) \\
    = 0.5\% + 2\% - 1.5\% - 1\% = \textbf{0\%}
\end{itemize}
$\rightarrow$ 경제가 침체되어 인플레이션과 생산량이 목표보다 낮으므로, 중앙은행은 금리를 0\%로 인하하여 경제를 부양하려고 합니다.
\end{examplebox}

통화정책을 사용하여 경기 순환을 안정화하는 것은 현대 시대 많은 중앙은행의 표준이었습니다. 하지만 이 접근 방식은 성공적이었을까요? 1980년대 이후 미국의 경험을 바탕으로 조심스럽게 "예"라고 결론 내릴 수 있습니다.

여기 1980년 이후 미국의 인플레이션, 실업률, 금리가 있습니다. 1990년 경기 침체가 시작되었을 때, 연방준비제도는 금리를 인하했습니다. 실업률은 4년 이내에 경기 침체 이전 수준 이하로 떨어졌고, 인플레이션은 1990년대 내내 낮게 유지되었습니다. 이 패턴은 2000년대 초 닷컴 버블 붕괴 이후에도 반복됩니다. 더 일반적으로, 1990년대 이후 거시경제 변수의 낮은 변동성은 중앙은행과 그들의 금리 결정에 기인합니다.

\begin{summarybox}
\textbf{레슨 3-1.3 요약}
\begin{itemize}
    \item \textbf{통화정책 목표 측정}: 최대 고용은 생산량 갭(실제 GDP와 잠재 GDP의 차이)으로, 물가 안정은 인플레이션 갭(실제 인플레이션과 목표 인플레이션의 차이)으로 측정됩니다.
    \item \textbf{통화정책 규칙}:
    \begin{itemize}
        \item \textbf{프리드먼의 K-규칙}: 경제 상황과 무관하게 통화 공급을 일정한 비율로 증가시키는 규칙입니다. 예측 가능성을 높이지만 유연성이 부족합니다.
        \item \textbf{테일러 준칙}: 금리 결정을 인플레이션 갭과 생산량 갭에 연동시키는 규칙입니다. 경제 상황에 따라 금리를 조정하여 경기 안정화를 목표로 합니다.
    \end{itemize}
    \item \textbf{성공적인 개입}: 1980년대와 2000년대 초 미국의 경험은 중앙은행의 금리 개입이 거시경제 변동성을 낮추는 데 성공적이었음을 시사합니다.
\end{itemize}
\end{summarybox}


\section{레슨 3-2: 통화정책의 전달}

\subsection{레슨 3-2.1: 통화정책 변화가 경제에 어떻게 영향을 미치는가?}

안녕하세요, 통화정책의 경제 전달에 대한 강의에 오신 것을 환영합니다. 이 세션에서는 통화정책 결정이 경제 전반에 어떻게 확산되는지 요약할 것입니다. 미국에서 연방준비제도는 변화하는 경제 상황에 대응하여 통화정책을 조정합니다. 목표는 최대 고용과 물가 안정이라는 것을 기억하십시오. 가장 일반적인 대응은 단기 금리인 \textbf{연방 기금 금리(federal funds rate)}를 인상하거나 인하하는 것입니다.

연방 기금 금리의 변화는 다른 금리와 더 넓은 금융 시장 상황의 변화로 이어질 것입니다. 차례로, 이러한 변화는 개별 가계와 기업의 지출 결정에 영향을 미칠 것입니다. 마지막으로, 이러한 개별 효과는 거시경제적 규모로 집계될 것입니다. 이는 경제 전반의 성장, 고용 및 인플레이션으로 이어져야 합니다. 우리의 목표는 이러한 일련의 사건이 어떻게 전개되는지 이해하는 것입니다. 이는 통화정책이 어떻게 작동하는지에 대한 우리의 이해를 향상시킬 것입니다. 시작해 봅시다.

\subsubsection*{통화정책이 시장 금리에 미치는 영향}

2013년 당시 연방준비제도 이사였던 제레미 스타인(Jeremy Stein)이 연설에서 말했듯이, "통화정책은 한 가지 중요한 장점을 가지고 있는데, 그것은 모든 틈새에 스며든다는 것입니다. 상업 은행, 브로커-딜러, 역외 헤지 펀드, 특수 목적 자산 담보부 기업 어음(ABCP) 특수 목적 법인(SPV)이 공통적으로 직면하는 한 가지는 모두 동일한 시장 금리 세트를 마주한다는 것입니다. 시장 금리가 위험 선호도나 만기 변환에 참여하려는 인센티브에 영향을 미치는 한, 금리 변화는 시장의 모든 구석에 도달할 수 있습니다."

먼저 통화정책이 시장 금리에 미치는 영향을 살펴보겠습니다. 연방 기금 금리는 \textbf{초단기 금리(overnight interest rate)}입니다. 은행들은 원하면 이 금리를 벌 수 있습니다. 다른 단기 금리들은 대체 투자입니다. 이는 다른 단기 금리들이 연방 기금 금리의 변화를 추적해야 한다는 것을 의미합니다.

\begin{examplebox}
\textbf{연방 기금 금리와 단기 금리}
연방 기금 금리는 은행들이 서로에게 하루짜리 돈을 빌려줄 때 적용되는 금리입니다. 만약 Fed가 이 금리를 올리면, 은행들은 더 높은 비용으로 돈을 빌리게 되므로, 다른 단기 대출(예: 3개월 만기 국채)에 대한 금리도 따라서 오르게 됩니다. 이는 은행들이 더 높은 수익률을 요구하거나, 더 높은 비용을 대출자에게 전가하기 때문입니다.
\end{examplebox}

3개월 만기 국채(Treasury Bills) 수익률, 즉 수익률과 연방 기금 금리를 비교해 봅시다. 그래프는 둘 사이의 차이가 0에 가깝다는 것을 보여줍니다. 이는 3개월 만기 국채 수익률이 연방 기금 금리와 거의 1대1로 움직인다는 것을 의미합니다.

장기 금리는 어떨까요? 그래프는 이제 10년 만기 국채(Treasury note) 수익률과 연방 기금 금리 간의 차이를 보여줍니다. 1대1 관계가 없다는 것을 알 수 있습니다.

더 일반적으로, 연방 기금 금리 변화가 단기 금리에 미치는 영향은 클 것입니다. 그러나 장기 금리에 미치는 영향은 완화될 것입니다. 장기 금리는 더 긴 기간에 대한 기대를 포함합니다. 이러한 금리는 사람들이 미래에 연방 기금 금리가 어떻게 변할 것으로 예상하는지에 반응하며, 단지 오늘날의 변화에만 반응하는 것이 아닙니다. 예를 들어, 시장이 통화정책이 향후 몇 년 동안 상당히 긴축될 것이라고 믿는다고 가정해 봅시다. 이 경우 오늘날의 장기 금리는 단기 금리보다 훨씬 높을 수 있습니다. 이것이 미래 연방 기금 금리 변화에 대한 \textbf{선제적 발언(forward-looking statements)}이 금융 시장 참가자들에 의해 면밀히 조사되는 한 가지 이유입니다.

\subsubsection*{가계 및 기업에 미치는 영향}

다음으로, 금리 변화가 가계와 기업에 어떻게 영향을 미칠까요? 통화정책 완화로 인해 장기 금리가 하락하고 있다고 가정해 봅시다. 이는 기업이 대출을 받아 투자하는 비용을 더 저렴하게 만들 것입니다. 마찬가지로, 소비자 대출과 모기지(mortgages)도 더 저렴해집니다. 이는 가전제품과 같은 내구재 및 주택 소유에 대한 지출을 자극합니다.

장기 금리 하락은 또한 미국 달러 자산의 수익률을 국제 투자자들에게 덜 매력적으로 만들 것입니다.

이러한 투자자들은 달러 표시 자산에 대한 수요를 줄여 달러 가치를 낮출 수 있습니다. \textbf{약한 달러(weaker dollar)}는 미국 상품이 해외에서 더 저렴해지므로 미국 수출을 증가시키는 경향이 있습니다. 동시에 수입품은 더 비싸져 미국 소비자와 기업이 국내 생산품을 구매하도록 유도합니다.

이 모든 조각들을 합쳐 봅시다. 기업과 가계에 대한 영향이 있습니다. 장기 금리 변화는 주가에도 영향을 미칩니다. 일반적으로 우리는 기업의 시장 가치를 모든 미래 배당금의 할인된 가치로 생각합니다. 장기 금리가 하락하면 미래 배당금에 대한 할인율이 더 유리해집니다. 이는 기업의 시장 가치를 증가시키고 결과적으로 주가를 높입니다. 이러한 주가 상승은 개인의 부를 증가시키고, 이는 다시 소비를 더욱 증가시킵니다.

\begin{examplebox}
\textbf{금리 인하의 경제적 파급 효과}
Fed가 금리를 인하하면:
\begin{enumerate}
    \item \textbf{대출 비용 감소}: 기업은 더 저렴하게 자금을 조달하여 투자를 늘리고, 가계는 모기지나 자동차 대출을 더 낮은 이자로 받을 수 있어 소비를 늘립니다.
    \item \textbf{주가 상승}: 미래 기업 이익의 현재 가치가 높아져 주가가 상승하고, 이는 가계의 자산 증가로 이어져 소비를 더욱 촉진합니다.
    \item \textbf{환율 하락}: 달러 자산의 매력이 감소하여 달러 가치가 하락합니다. 이는 수출을 늘리고 수입을 줄여 국내 생산을 활성화합니다.
\end{enumerate}
이러한 효과들이 종합적으로 작용하여 경제 성장, 고용 증가, 그리고 필립스 곡선에 따라 인플레이션 상승으로 이어질 수 있습니다.
\end{examplebox}

낮은 금리로 경제를 부양하는 것이 실업률을 낮출 것이라는 믿음입니다. 이는 또한 필립스 곡선에 따라 물가를 높일 것입니다. 이것이 이론입니다. 이제 이러한 효과가 데이터에서 어떻게 나타나는지 살펴보겠습니다. 우리는 연방 기금 금리의 예상치 못한 인상을 검토할 것입니다. 산업 생산, 통화 공급, 실업률, 인플레이션에 어떤 일이 일어나는지 알고 싶습니다. 우리가 보게 될 차트들은 월별 데이터를 기반으로 추정된 시계열 통계 분석에 기반합니다. 우리는 2년의 시간 범위에 걸쳐 효과를 측정할 것입니다.

\subsubsection*{데이터로 본 통화정책의 영향}

첫 번째 그래프는 연방 기금 금리가 50bp(베이시스 포인트) 예상치 못하게 한 번 인상된 후 어떻게 되는지 보여줍니다. 단기적으로는 금리가 더 인상될 가능성이 있다는 점에 유의하십시오. 그러나 시간이 지남에 따라 금리 인상의 효과는 사라지고 우리는 시작점으로 돌아갑니다.

여기 M2 통화 공급이 있습니다. 연방준비제도는 공개 시장 조작(open market operations)을 통해 연방 기금 금리 인상을 시행합니다. 공개 시장 조작은 Fed가 증권을 매각하고 그 대가로 은행 준비금 형태로 현금을 받는다는 것을 기억하십시오. 준비금은 M2로 측정되는 통화 공급의 일부입니다. 연방 기금 금리를 인상하면 통화 공급이 줄어듭니다. 이 효과가 가라앉는 데 약 2년이 걸립니다.

여기 산업 생산이 있습니다. 그래프는 산업 생산이 감소하는 데 약 6개월이 걸린다는 것을 보여줍니다. 이것은 중요한 점입니다. 통화정책은 시차를 두고 경제에 영향을 미칩니다. 이는 적어도 두 가지 이유로 발생합니다. 첫째, 기업은 이미 투자 계획을 위한 자금을 받았을 수 있으며 여전히 이를 실행하고 있을 수 있습니다. 둘째, 금리 변화가 소비자 지출에 영향을 미치는 데 시간이 걸릴 수 있습니다.

실업률은 산업 생산의 거울 이미지입니다. 몇 달이 걸리지만, 산업 생산이 감소함에 따라 실업률은 상승하기 시작합니다. 또한 연방 기금 금리가 인상되면 인플레이션이 하락한다는 증거가 있습니다. 중앙은행이 경제 과열을 우려하여 연방 기금 금리를 인상하면, 중기적으로 실업률은 증가하고 인플레이션은 하락할 것입니다.

\begin{summarybox}
\textbf{레슨 3-2.1 요약}
\begin{itemize}
    \item \textbf{통화정책 전달 경로}: 연방 기금 금리 변화 $\rightarrow$ 다른 금리 및 금융 시장 상황 변화 $\rightarrow$ 가계 및 기업의 지출 결정 변화 $\rightarrow$ 거시경제적 성장, 고용, 인플레이션 변화.
    \item \textbf{시장 금리 영향}: 연방 기금 금리는 단기 금리에 큰 영향을 미치지만, 장기 금리에는 미래 기대가 반영되어 영향이 완화됩니다.
    \item \textbf{경제적 영향}: 금리 인하(완화 정책)는 기업 투자 및 가계 소비를 자극하고, 달러 약세를 유발하여 수출을 증대시키며, 주가를 상승시켜 부를 증가시킵니다.
    \item \textbf{데이터로 본 효과}: 연방 기금 금리 인상 시, 통화 공급은 즉시 감소하고, 산업 생산은 약 6개월 후 감소하며, 실업률은 상승하고, 인플레이션은 하락합니다. 이러한 효과는 \textbf{시차(time lag)}를 두고 나타납니다.
\end{itemize}
\end{summarybox}


\subsection{레슨 3-2.2: 은행을 통한 통화정책 전달}

안녕하세요, 통화정책의 은행 대출 채널에 대한 강의에 오신 것을 환영합니다. 이 수업에서는 통화정책이 은행의 신용 공급에 어떻게 영향을 미치는지 살펴볼 것입니다. 우리는 세 가지 메커니즘을 논의할 것입니다: \textbf{순이자 마진(net interest rate margin)}, \textbf{지급 준비금(reserves)}, 그리고 \textbf{예금(deposits)}.

\subsubsection*{은행의 사업 모델과 순이자 마진}

은행의 사업 모델을 생각해 봅시다. 은행은 예금자와 머니 마켓 펀드로부터 단기적으로 자금을 빌려, 이 자금을 사용하여 장기 대출을 실행합니다. 단기 금리가 보통 장기 금리보다 낮기 때문에, 은행은 단기 금리와 장기 금리 간의 차이를 이용하여 돈을 법니다.

\begin{conceptbox}
\textbf{순이자 마진 (Net Interest Margin, NIM)}
은행의 핵심 수익성 지표입니다. 은행이 대출 및 모기지에서 얻는 이자 수입과 예금자에게 지급하는 이자 비용 간의 차이를 나타냅니다.
\begin{itemize}
    \item \textbf{계산}: (대출 이자 수입 - 예금 이자 비용) / 총 자산
    \item \textbf{의미}: NIM이 높을수록 은행의 수익성이 좋다는 것을 의미하며, 이는 새로운 대출을 실행할 수 있는 자본이 더 많다는 것을 뜻합니다.
\end{itemize}
\end{conceptbox}

은행의 대출 및 모기지에서 발생하는 순이자 수입과 은행이 예금자에게 지급하는 이자 비용 간의 차이를 \textbf{순이자 마진(net interest margin)}이라고 합니다. 순이자 마진은 은행의 수익성을 측정합니다. 순이자 마진이 클수록 이익이 높습니다.

이 그래프는 1980년대 이후 미국 은행의 순이자 마진을 보여줍니다. 보시다시피, 은행의 순이자 마진은 시간이 지남에 따라 감소했습니다.

동시에, 2010년경 연방준비제도가 연방 기금 금리를 처음으로 0으로 인하했을 때 순이자 마진이 크게 증가하는 것을 볼 수 있습니다.

\begin{examplebox}
\textbf{통화정책이 순이자 마진에 미치는 영향}
\begin{itemize}
    \item \textbf{시나리오 1: 금리 인상 (긴축 정책)}
    \begin{itemize}
        \item 은행이 100달러를 2\% 이자로 예금받아, 5\% 이자로 100달러 고정 금리 모기지를 대출했다고 가정합니다. 순이자 마진은 3달러(3\%)입니다.
        \item Fed가 연방 기금 금리를 2\%에서 3\%로 인상하면, 은행은 예금 이자를 3\%로 올려야 예금을 유지할 수 있습니다.
        \item 하지만 고정 금리 모기지 이자는 5\%로 변하지 않습니다.
        \item \textbf{결과}: 순이자 마진은 5달러 - 3달러 = 2달러(2\%)로 감소합니다.
    \end{itemize}
    \item \textbf{시나리오 2: 금리 인하 (완화 정책)}
    \begin{itemize}
        \item 위와 동일한 초기 조건에서 Fed가 금리를 인하하면, 은행은 예금 이자를 낮출 수 있습니다.
        \item 고정 금리 대출의 이자는 변하지 않으므로, 순이자 마진은 증가하여 은행의 수익성이 개선됩니다.
    \end{itemize}
\end{itemize}
\textbf{결론}: 금리 인상은 은행의 순이자 마진을 압박하여 수익을 감소시키고, 이는 새로운 대출을 위한 자본을 줄여 신용 공급을 위축시킵니다. 반대로 금리 인하는 순이자 마진을 개선하여 신용 공급을 확대할 수 있습니다.
\end{examplebox}

통화정책이 순이자 마진에 어떻게 영향을 미치는지 간단한 예시로 살펴보겠습니다. 100달러의 예금으로 전액 자금을 조달하는 은행을 생각해 봅시다. 예금 이자율은 2\%이므로 총 이자 비용은 2달러입니다.

은행은 5\% 이자를 지급하는 100달러 고정 금리 모기지 하나를 발행했습니다.

이자 수입은 5달러이고, 2달러의 이자 비용을 제외하면 순이자 마진은 3달러 또는 3\%입니다.

이제 연방준비제도가 연방 기금 금리를 2\%에서 3\%로 인상한다고 가정해 봅시다. 연방 기금 시장은 초단기 대출 시장이므로, 다른 모든 초단기 및 단기 금리도 인상될 것입니다. 예금은 요구불 예금이며, 따라서 초단기 대출과 매우 유사합니다. 은행이 예금을 잃고 싶지 않다면, 그에 따라 예금 금리를 인상해야 할 것입니다. 은행은 예금에 대해 3\%를 지불해야 하지만, 고정 금리 모기지의 이자율은 변하지 않았습니다. 따라서 순이자 마진은 3\%에서 2\%로 하락합니다.

순이자 마진 하락은 이익 감소를 의미합니다. 이익 감소는 다시 새로운 대출을 위한 자본 감소를 의미합니다. 다시 말해, 통화정책은 은행의 이익에 직접적으로 영향을 미치고, 따라서 은행의 신용 공급에 영향을 미칩니다. 물론 통화정책이 이익에 얼마나 강하게 영향을 미치는지는 은행이 대차대조표에 보유하고 있는 고정 금리 대출의 비중에 따라 달라집니다. 기업 대출과 같은 일부 대출은 일반적으로 변동 금리 대출이므로, 이러한 대출의 순이자 마진은 통화정책의 영향을 받지 않습니다.

\subsubsection*{지급 준비금 채널}

통화정책이 은행에 영향을 미치는 두 번째 방법은 은행이 연방준비제도 계좌에 보유하고 있는 \textbf{지급 준비금(reserves)}을 통해서입니다. 이를 \textbf{연방 기금(federal funds)}이라고 합니다. 연방준비제도가 연방 기금 금리를 인상할 때, 일반적으로 공개 시장 조작을 통해 그렇게 합니다. 즉, 연방준비제도는 은행에 증권을 매각합니다. 은행은 준비금으로 증권 구매 대금을 지불합니다. 연방 기금 시장에서 사용할 수 있는 준비금이 적어지면 단기 자금 시장의 신용 공급이 감소합니다. 은행은 단기 유동성 필요, 예를 들어 대출을 발행할 때 연방 기금 시장을 사용합니다. 물론 은행은 다른 시장에서도 자금을 빌릴 수 있지만, 종종 그 자금은 더 비싸거나 환매 조건부 채권(Repo) 시장과 같이 담보를 요구합니다. 은행에 사용할 수 있는 자금이 적거나, 남아 있는 자금이 더 비쌀 때, 은행은 신용 공급을 줄입니다. 따라서 통화정책은 준비금을 줄임으로써 은행의 신용 공급을 줄입니다.

\subsubsection*{예금 채널}

통화정책이 은행에 영향을 미치는 또 다른 방법은 \textbf{예금(deposits)}을 통해서입니다. 순이자 마진에 대해 제가 들었던 간단한 예시를 기억하십시오. 그 예시에서 은행은 연방 기금 금리 인상에 맞춰 예금 금리를 인상했습니다. 그러나 여러분 대부분이 당좌 예금 및 저축 계좌에서 알고 있듯이, 은행은 일반적으로 당좌 예금 및 저축 계좌의 금리 인상에 빠르게 대응하지 않습니다.

대부분의 예금자들은 투자 대안을 찾지 않는 경향이 있습니다. 시간이 걸리기 때문입니다. 모든 은행 정보를 변경해야 하므로 은행을 바꾸는 데는 훨씬 더 많은 시간이 걸립니다. 그럼에도 불구하고 일부 예금자, 특히 상당한 예금을 가진 예금자들은 은행을 바꿉니다. 그들은 예금 금리와 대체 투자 금리 간의 차이가 증가할 때 은행에서 예금을 인출하는 경향이 있습니다. 예금의 더 가까운 대체재는 \textbf{머니 마켓 뮤추얼 펀드(money market mutual funds)}입니다.

그래프는 머니 마켓 뮤추얼 펀드의 총 자산과 연방 기금 금리를 보여줍니다. 보시다시피, 연방준비제도가 연방 기금 금리를 인상할 때, 예를 들어 2004년에서 2006년 사이 또는 2016년에서 2019년 사이에 머니 마켓 뮤추얼 펀드의 자산이 훨씬 더 빠르게 증가했습니다.

예금 감소는 또한 은행이 자금을 덜 보유하고 있으므로 신용 공급을 줄여야 한다는 것을 의미합니다. 여러분은 왜 이러한 행동이 은행에 최적인지 궁금할 수 있습니다. 은행은 일부 예금을 잃고 일부 대출 기회를 포기하는 것과 모든 예금의 이자율을 1\% 포인트 인상하고 현재 대출 속도를 유지하는 것 사이에서 선택에 직면합니다. 다르게 말하면, 예금의 5\%를 잃고 대출을 5\% 줄이는 것과 모든 예금의 이자율을 1\% 포인트 인상하고 현재 대출 속도를 유지하는 것 중 어느 것을 선호하시겠습니까? 최근 연구에 따르면 은행은 예금 금리를 인상하는 시점에 대해 전략적입니다. 미국 내에서는 은행 경쟁이 더 심한 곳에서 예금 금리가 더 빠르게 인상되는 경향이 있습니다. 이는 은행을 바꾸기 쉬울수록 은행이 예금을 유지하기 위해 금리 변화에 더 많이 반응해야 한다는 것을 시사합니다. 이 발견의 반대 측면은 은행 시스템이 고도로 집중된 지역이나 국가에서는 이러한 통화정책의 예금 채널이 덜 중요하다는 것입니다.

\begin{summarybox}
\textbf{레슨 3-2.2 요약: 은행을 통한 통화정책 전달 3가지 채널}
\begin{itemize}
    \item \textbf{순이자 마진 (Net Interest Margin) 채널}: 연방 기금 금리 인상은 은행의 순이자 마진을 압박하여 수익을 감소시키고, 이는 새로운 대출을 위한 자본을 줄여 신용 공급을 위축시킵니다.
    \item \textbf{지급 준비금 (Reserves) 채널}: 연방 기금 금리 인상은 공개 시장 조작을 통해 은행 시스템 내의 지급 준비금을 감소시킵니다. 이는 은행의 단기 자금 조달 비용을 높여 신용 공급을 줄입니다.
    \item \textbf{예금 (Deposits) 채널}: 연방 기금 금리 인상은 예금자들이 더 높은 수익률을 찾아 머니 마켓 뮤추얼 펀드와 같은 대체 투자로 이동하게 만듭니다. 이는 은행의 예금 기반 자금을 감소시켜 신용 공급을 줄입니다.
\end{itemize}
이 세 가지 채널 모두 연방 기금 금리 인상이 은행의 전반적인 신용 공급을 감소시킨다는 것을 의미합니다.
\end{summarybox}


\subsection{레슨 3-2.3: 기업과 가계를 통한 통화정책 전달}

안녕하세요, 기업과 가계를 통한 통화정책 전달에 대한 강의에 오신 것을 환영합니다. 이 수업의 목표는 금리 변화가 차입 비용으로 어떻게 전달되는지 더 잘 이해하는 것입니다. 우리는 이것이 기업, 특히 변동 금리 은행 대출에 의존하는 소규모 및 재정적으로 제약된 기업에 어떻게 중요한지 살펴볼 것입니다. 이는 경제의 투자 및 일자리에 연쇄적인 영향을 미 미칠 수 있습니다. 그 다음, 동일한 채널이 가계에 어떻게 영향을 미치는지 검토할 것입니다. 여기서는 주택 시장과 모기지 대출에 초점을 맞출 것입니다. 우리는 금리 변화가 모기지를 재융자할 옵션이 있는 가계에 어떻게 중요한지 살펴볼 것입니다. 그리고 나서 경제의 총수요에 대한 함의를 논의할 것입니다. 시작해 봅시다.

\subsubsection*{기업에 미치는 영향: 변동 금리 대출}

먼저 기업을 살펴보겠습니다. 거의 모든 기업은 어떤 형태의 부채를 안고 있습니다. 가장 큰 기업들은 고정 금리 부채인 채권 시장에 접근할 수 있습니다. 소규모 기업들은 일반적으로 은행 대출에만 접근할 수 있습니다. 이러한 사업 대출은 종종 \textbf{변동 금리(floating or variable interest rate)}를 가집니다. 이는 대출 이자율이 일부 변동 벤치마크에 고정 마크업을 더한 것과 같다는 것을 의미합니다. 변동 벤치마크는 과거에는 런던 은행 간 금리(LIBOR), 담보부 초단기 자금 조달 금리(SOFR) 또는 최우수 고객에게 제공되는 금리인 프라임 금리(prime rate) 중 하나일 수 있습니다. 일반적으로 이는 기업이 지불하는 이자율이 단기 금리와 함께 움직인다는 것을 의미합니다.

\begin{conceptbox}
\textbf{변동 금리 대출 (Variable Rate Loans)}
대출 이자율이 시장 금리(벤치마크 금리)에 따라 변동하는 대출입니다.
\begin{itemize}
    \item \textbf{구성}: 벤치마크 금리 (예: SOFR, Prime Rate) + 고정 마크업 (은행의 이윤)
    \item \textbf{특징}: 중앙은행의 정책 금리 변화에 직접적으로 영향을 받습니다. 정책 금리가 오르면 대출 이자도 오르고, 내리면 대출 이자도 내립니다.
    \item \textbf{주요 대상}: 주로 소규모 기업 대출이나 일부 가계 대출(예: 변동 금리 모기지)에 적용됩니다. 대기업은 고정 금리 채권 발행을 통해 금리 변동 위험을 회피하는 경향이 있습니다.
\end{itemize}
\end{conceptbox}

이것이 통화정책과 어떻게 연결되는지 예시를 통해 살펴보겠습니다. 한 기업이 100만 달러의 사업 대출을 가지고 있다고 가정해 봅시다. 이 대출의 이자율은 SOFR + 2\%라고 가정해 봅시다. 연방 기금 금리와 마찬가지로 SOFR은 초단기 금리이며 둘은 밀접하게 움직입니다. 연방 기금 금리와 SOFR이 현재 4\%라고 가정해 봅시다. 이는 대출 이자율이 6\%이고 총 연간 이자 지급액이 60,000달러라는 것을 의미합니다.

이제 FOMC가 연방 기금 금리를 1\%로 인하하기로 결정했다고 가정해 봅시다. SOFR도 이 하락 움직임을 따를 것입니다. 이는 기업 대출의 이자율이 6\%에서 3\%로 재설정되어 이자 지급액이 30,000달러로 감소한다는 것을 의미합니다. 따라서 연방 기금 금리 인하는 기업이 투자 및 고용에 사용할 수 있는 30,000달러를 확보해 줍니다.

더 일반적으로, 연방 기금 금리 인하는 레버리지가 있는 기업의 이익을 증가시킬 것입니다. 더 높은 이익의 또 다른 함의는 기업의 대차대조표와 주식 가치가 개선될 것이라는 점입니다. 더 높은 이익으로 기업은 더 나은 재정 상태에 있습니다. 이는 차입 비용을 더욱 줄일 수 있습니다.

데이터가 이러한 아이디어를 지지할까요? 글쎄요, 그렇기도 하고 아니기도 합니다. 연구는 기업이 연방 기금 금리 인상에 어떻게 반응하는지, 그리고 그것이 경제 전체에 중요한지 여부를 조사했습니다. 연구에 따르면 대기업은 금리가 인상될 때 대부분 영향을 받지 않습니다. 이자 비용이 증가하고 수익이 감소하더라도 생산 및 고용 수준을 유지합니다. 이러한 대기업은 기업 어음 시장 및 기타 단기 신용원에 접근하는 데 문제가 없습니다. 이는 이익이 감소하더라도 단기 차입을 늘릴 수 있도록 합니다. 동시에 대기업의 재고는 증가합니다.

이와는 대조적으로, 통화정책이 소규모 기업의 차입 비용에 미치는 영향은 훨씬 더 중요합니다. 소규모의 재정적으로 제약된 기업은 이자 비용 증가에 재고를 줄이고 근무 시간과 생산을 줄임으로써 반응합니다. 경제 전체에 중요한 것은, 그 효과가 여전히 상당하다는 것입니다. 이는 소규모 기업이 여전히 미국 전체 근로자의 절반 이상을 고용하고 있기 때문입니다.

\subsubsection*{가계에 미치는 영향: 주택 시장과 재융자}

다음으로, 통화정책에 가계가 어떻게 반응하는지 살펴보겠습니다. 모든 가계가 FOMC의 성명을 주의 깊게 읽는다고 말하는 것은 아닙니다. 그러나 그들의 의사 결정은 변화하는 시장 금리를 포함한 지배적인 경제 상황에 반응할 것입니다. 가계 대차대조표에서 가장 큰 항목은 보통 \textbf{모기지(mortgage)}입니다.

이 그래프는 30년 고정 금리 모기지 이자율과 연방 기금 금리를 보여줍니다. 보시다시피, 두 금리는 함께 움직이므로 모기지 금리는 통화정책 변화에 반응할 것입니다.

모기지 금리 변화에서 얼마나 큰 효과를 기대할 수 있을까요? 답은 가계가 얼마나 부채를 지고 있는지에 달려 있습니다. 가계가 부채가 전혀 없다면, 차입 비용 변화는 무관할 것입니다. 지도는 2007년, 2008년 금융 위기 직전 미국 모든 카운티의 가계 중위 부채 대비 소득 비율을 보여줍니다. 보시다시피, 많은 카운티에는 연간 소득의 세 배가 넘는 부채를 가진 가계가 있습니다. 연간 소득 대비 이렇게 높은 차입으로 인해 금리 변화는 가계 재정에 상당한 영향을 미칠 수 있습니다.

미국의 일부 모기지는 변동 금리이지만, 대부분의 모기지는 고정 금리입니다. FOMC가 연방 기금 금리를 인하하기로 결정하면 모기지 금리도 하락하는 경향이 있습니다. 이에 대응하여 가계는 현재 모기지를 \textbf{재융자(refinance)}하기로 결정할 수 있습니다. 이는 가계가 현재 모기지를 더 낮은 이자율 모기지로 대체하는 것을 의미합니다. 또한 이는 선불 비용을 수반할 수 있습니다. 재융자로 인한 이자율 절감은 상당할 수 있습니다. 모기지 이자율 그래프로 돌아가서 21세기에 초점을 맞춰 봅시다.

2006년 모기지 이자율은 약 6.5\%였고, 2011년 12월에는 4\% 미만으로 떨어졌습니다. 원금을 전혀 갚지 않은 가계는 모기지를 재융자함으로써 모기지 이자 비용을 거의 40\% 절감할 수 있었고, 많은 사람들이 그렇게 했습니다. 최근 연구에 따르면 재융자를 한 평균 가계는 월 모기지 상환액에서 거의 1,000달러를 절약했습니다.

가계는 이 추가 현금으로 무엇을 했을까요? 모기지 상환액을 절약한 가계는 더 많이 소비하는 경향이 있습니다. 예를 들어, 사람들은 외식을 하거나 집을 개조하기 시작할 수 있습니다. 내구재와 자동차에 대한 지출이 증가합니다. 결과적으로 상품 및 서비스에 대한 수요가 증가하여 기업에 이익이 되고 경제를 활성화합니다.

\begin{summarybox}
\textbf{레슨 3-2.3 요약: 기업과 가계를 통한 통화정책 전달}
\begin{itemize}
    \item \textbf{기업에 미치는 영향}:
    \begin{itemize}
        \item \textbf{변동 금리 대출}: 소규모 기업은 주로 변동 금리 은행 대출에 의존하므로, 정책 금리 인하 시 이자 비용이 감소하여 이익이 증가하고 투자 및 고용을 늘릴 수 있습니다.
        \item \textbf{대기업과의 차이}: 대기업은 채권 시장 접근성 등으로 인해 금리 변화에 덜 민감합니다.
    \end{itemize}
    \item \textbf{가계에 미치는 영향}:
    \begin{itemize}
        \item \textbf{모기지 금리}: 정책 금리 변화는 모기지 금리에 영향을 미칩니다.
        \item \textbf{재융자}: 금리 인하 시 가계는 모기지를 재융자하여 월 상환액을 줄일 수 있으며, 이로 인해 가처분 소득이 증가하여 소비(내구재, 자동차 등)를 늘립니다.
    \end{itemize}
    \item \textbf{총수요 증가}: 금리 인하로 인한 기업의 투자 증가와 가계의 소비 증가는 경제의 전반적인 상품 및 서비스 수요를 증가시켜 경제를 활성화합니다.
\end{itemize}
\end{summarybox}


\section{레슨 3-3: 통화정책의 한계}

\subsection{레슨 3-3.1: 통화정책의 한계: 공급 충격}

안녕하세요, 스태그플레이션에 대한 강의에 오신 것을 환영합니다. 이 수업에서는 낮은 성장, 높은 인플레이션, 높은 실업률로 특징지어졌던 1970년대 기간을 살펴볼 것입니다. 이 사례 연구는 통화정책의 상충 관계와 한계를 더 잘 이해하는 데 도움이 될 것입니다. \textbf{스태그플레이션(stagflation)}이라는 용어는 \textbf{경기 침체(stagnation)}와 \textbf{인플레이션(inflation)}의 조합입니다. 스태그플레이션은 중앙은행에게 특히 어려운 문제입니다. 경제를 부양하기 위해 중앙은행은 금리를 인하하지만, 낮은 금리는 또한 인플레이션을 증가시키는 경향이 있습니다. 성장이 낮고 인플레이션이 높은 상황에서 중앙은행은 선택의 폭이 제한적입니다.

\begin{conceptbox}
\textbf{스태그플레이션 (Stagflation)}
경제 성장(stagnation)이 둔화되거나 침체되는 동시에 물가 상승(inflation)이 지속되는 현상을 말합니다. 일반적으로 경제학 이론에서는 경기 침체 시 물가가 하락하거나 안정되는 경향이 있다고 보지만, 스태그플레이션은 이러한 일반적인 관계를 벗어나는 예외적인 상황입니다.
\begin{itemize}
    \item \textbf{중앙은행의 딜레마}:
    \begin{itemize}
        \item \textbf{경기 부양 (금리 인하)}: 실업률을 낮추고 성장을 촉진할 수 있지만, 이미 높은 인플레이션을 더욱 악화시킬 수 있습니다.
        \item \textbf{인플레이션 억제 (금리 인상)}: 인플레이션을 낮출 수 있지만, 이미 낮은 성장을 더욱 둔화시키고 실업률을 높일 수 있습니다.
    \end{itemize}
    이러한 딜레마 때문에 스태그플레이션은 중앙은행에게 매우 다루기 어려운 문제입니다.
\end{itemize}
\end{conceptbox}

그래프는 1960년부터 1990년까지의 소비자 물가 인플레이션을 보여줍니다. 1960년대 후반에 인플레이션이 가속화되기 시작한 것을 볼 수 있습니다. 가장 큰 급등은 1974년과 1980년에 발생했으며, 이때 인플레이션율은 각각 11\%와 13.5\%에 달했습니다. 이러한 급등의 원인을 살펴보겠습니다.

\subsubsection*{공급 충격: 1970년대 오일 쇼크}

주요 원인은 두 번의 큰 \textbf{오일 쇼크(oil shocks)}입니다. 첫 번째 오일 쇼크는 욤 키푸르 전쟁(Yom Kippur War) 이후 OPEC이 미국과 다른 서방 국가들에 대한 석유 금수 조치를 취하면서 발생했습니다. 몇 달 만에 유가는 배럴당 3달러에서 약 11달러로 거의 300\% 상승했습니다.

\begin{conceptbox}
\textbf{공급 충격 (Supply Shocks)}
생산 비용에 갑작스럽고 예상치 못한 변화를 가져오는 사건을 말합니다. 이는 주로 생산 요소의 가격이나 공급량에 영향을 미칩니다.
\begin{itemize}
    \item \textbf{긍정적 공급 충격}: 생산 비용을 낮추고 생산량을 늘리는 충격 (예: 기술 혁신, 풍작).
    \item \textbf{부정적 공급 충격}: 생산 비용을 높이고 생산량을 줄이는 충격 (예: 유가 급등, 자연재해, 전쟁).
\end{itemize}
\textbf{스태그플레이션과의 관계}: 부정적 공급 충격은 생산 비용을 높여 물가를 상승시키고(인플레이션), 동시에 생산량을 감소시켜 경제 성장을 둔화시키거나 침체시킬 수 있습니다(경기 침체). 이는 스태그플레이션을 유발하는 주요 원인 중 하나입니다.
\end{conceptbox}

이는 경제에 엄청난 비용 충격이었습니다. 거의 모든 경제 부문이 휘발유 형태나 발전용, 화학 산업 등에서 석유를 수입으로 사용했기 때문입니다. 생산자 물가는 급격히 상승했습니다. 동시에 소비자들은 휘발유와 난방에 더 많은 돈을 지불해야 했고, 이는 다른 소비재에 대한 지출을 줄였습니다.

불과 몇 년 후, 두 번째 오일 쇼크가 경제를 강타했습니다. 1979년 이란 혁명과 1980년 이란-이라크 전쟁 발발 이후 석유 생산량이 감소했습니다. 석유 공급 감소와 미래 석유 공급에 대한 불확실성으로 인해 유가는 폭등했습니다. 1980년에는 유가가 배럴당 39달러 50센트에 달했습니다. 이는 1973년 초 이후 10배 이상 증가한 수치입니다. 미국에서는 자동차 제조업체들이 더 연료 효율적인 자동차를 제공하기 위해 고군분투하면서 특히 큰 타격을 입었습니다. 이러한 대규모 충격은 경제 둔화로 이어졌습니다.

그래프에서 실업률을 볼 수 있습니다. 1975년에는 9\%로 제2차 세계 대전 이후 최고치를 기록했으며, 1981년 두 번째 오일 쇼크와 함께 다시 감소하여 7.8\%에 도달했습니다. 연방준비제도의 책무는 물가 안정과 최대 고용이라는 것을 기억하십시오. 그러나 오일 쇼크는 그 반대 효과, 즉 높은 인플레이션과 높은 실업률을 초래했습니다. 오일 쇼크는 \textbf{공급 충격(supply shocks)}의 예시입니다.

그래프는 공급 충격이 실업률과 인플레이션 간의 상충 관계, 즉 필립스 곡선을 어떻게 바깥쪽으로 이동시키는지 보여줍니다. 우리는 별도의 강의에서 필립스 곡선을 자세히 논의했습니다.

\subsubsection*{통화정책의 한계와 대응}

공급 충격에 직면했을 때 통화정책의 선택지는 제한적입니다. 낮은 금리는 공급 제약에 영향을 미치지 않기 때문입니다. 실제로 오일 쇼크에 대한 많은 대응은 통화정책이 아니라 단기 및 장기적으로 석유 수요에 영향을 미쳤습니다. 1973년 오일 쇼크에 대한 가장 극적인 단기 대응은 미국에서 휘발유 배급제를 시행한 것이었습니다. 예를 들어, 홀수 번호판을 가진 차량은 홀수 날에만 휘발유를 살 수 있었고, 다른 차량은 짝수 날에만 살 수 있었습니다.

석유 수요에 상당한 장기적 영향을 미친 대응은 결국 유가를 낮추는 데 도움이 되었습니다. 몇 가지 예시가 있습니다. 휘발유 소비를 줄이기 위해 1974년에는 전국 최고 속도 제한이 시속 55마일로 부과되었습니다. 1년 후인 1975년에는 기업 평균 연비(CAFE) 기준이 만들어졌습니다. 이 기준은 미국에서 자동차와 경트럭의 연비 개선을 요구했습니다. 동시에 미국은 \textbf{전략 비축유(Strategic Petroleum Reserve)}를 개발했습니다.

이러한 조치와 생산량 증가는 결국 1980년대에 소위 \textbf{오일 과잉(oil glut)}으로 이어졌습니다. 유가는 상당히 하락했고 인플레이션을 낮추는 데 기여했습니다.

1980년대의 이러한 공급 충격은 인플레이션과 실업률 간의 상충 관계를 다시 안쪽으로 이동시키는 데 도움이 되었습니다. 그래프에서 볼 수 있듯이, 1980년대 후반에는 1980년대 초에 비해 실업률을 줄이기 위해 훨씬 적은 인플레이션이 필요했습니다.

그럼에도 불구하고 통화정책은 공급 충격에 대응하는 데 거의 할 수 없습니다. 이 기간은 인플레이션 기대의 변화를 야기했습니다. 즉, 합의는 낮은 인플레이션 체제에서 인플레이션이 높게 유지될 체제로 바뀌었습니다.

높은 인플레이션 기대는 높은 임금 요구로 이어졌고, 이는 다시 인플레이션을 증가시켰습니다.

폴 볼커(Paul Volcker)가 1979년 연방준비제도 의장이 되었을 때, 그는 인플레이션과 인플레이션 기대를 낮추기를 원했습니다.

그는 연방 기금 금리를 인상함으로써 그렇게 했으며, 1981년 12월에는 19\%에 달했습니다. 금리의 급격한 인상은 유가 상승과 함께 1980년과 1982년 경기 침체를 촉발했습니다.

실업률은 10\% 이상으로 상승하여 제2차 세계 대전 종전부터 코로나19 팬데믹까지 최고 수준을 기록했습니다. 다시 말해, 이 기간 동안 연방준비제도는 물가 안정과 최대 고용이라는 이중 책무에도 불구하고 실업률보다는 인플레이션에 초점을 맞추기로 선택했습니다. 연방 기금 금리 인상은 인플레이션의 큰 하락으로 이어졌습니다.

더 중요하게는, 연방준비제도는 인플레이션을 낮추기 위해 일시적으로 더 높은 실업률이라는 대가를 받아들였습니다. 이는 미래의 낮은 인플레이션에 대한 신뢰할 수 있는 약속으로 여겨졌습니다. 이 약속은 다시 인플레이션 기대를 낮추는 데 도움이 되었고, 이는 낮은 인플레이션 체제를 확립하는 데 기여했습니다.

\begin{summarybox}
\textbf{레슨 3-3.1 요약: 통화정책의 한계와 공급 충격}
\begin{itemize}
    \item \textbf{스태그플레이션}: 1970년대는 낮은 성장, 높은 인플레이션, 높은 실업률이 동시에 나타나는 스태그플레이션으로 특징지어졌습니다. 이는 중앙은행에게 큰 딜레마를 안겨줍니다.
    \item \textbf{공급 충격의 영향}: 유가 급등과 같은 공급 충격은 생산 비용을 증가시켜 물가를 상승시키고(인플레이션), 동시에 경제 활동을 위축시켜 실업률을 높입니다. 이는 필립스 곡선을 바깥쪽으로 이동시킵니다.
    \item \textbf{통화정책의 한계}: 통화정책은 금리 조정을 통해 총수요에 영향을 미치지만, 공급 측면의 제약(예: 원유 생산량 감소)을 직접적으로 해결할 수 없습니다. 따라서 공급 충격에 대응하는 데는 한계가 있습니다.
    \item \textbf{인플레이션 기대 관리}: 공급 충격이 인플레이션 기대를 영구적으로 변화시킬 경우, 중앙은행은 인플레이션과 인플레이션 기대를 낮추기 위해 일시적으로 높은 실업률을 감수하는 강력한 긴축 정책을 펼칠 수 있습니다 (예: 폴 볼커 의장 시기). 이는 미래의 물가 안정에 대한 신뢰를 구축하는 데 중요합니다.
\end{itemize}
\end{summarybox}


\subsection{레슨 3-3.2: 통화정책의 한계: 고용 없는 회복과 위험 감수 인센티브}

(이 레슨의 본문 내용은 제공된 텍스트 파트 1/2에 포함되어 있지 않습니다.)


\section{체크리스트: 학습 및 복습 가이드}

이 문서를 통해 학습한 내용을 점검하고 복습하는 데 도움이 되는 체크리스트입니다.

\begin{itemize}
    \item \textbf{개요 및 기본 개념 이해}
    \begin{itemize}
        \item Fed의 이중 책무(최대 고용, 물가 안정)를 정확히 설명할 수 있는가?
        \item GDP, 실질 GDP, 실업률의 정의와 중요성을 이해하는가?
    \end{itemize}
    \item \textbf{인플레이션과 실업률의 관계}
    \begin{itemize}
        \item 오쿤의 법칙(실업률과 GDP 성장률의 음의 상관관계)을 설명할 수 있는가?
        \item 자연 실업률의 개념과 중앙은행이 실업률을 0으로 만들지 않는 이유를 아는가?
        \item 소비자 물가 인플레이션과 생산자 물가 인플레이션, 임금 인플레이션의 관계를 이해하는가?
        \item 필립스 곡선(인플레이션과 실업률의 상충 관계)을 설명할 수 있는가?
        \item 인플레이션 기대가 필립스 곡선에 미치는 영향(곡선 이동)을 이해하는가?
    \end{itemize}
    \item \textbf{통화정책 규칙}
    \begin{itemize}
        \item 잠재 GDP와 생산량 갭(output gap)의 개념을 아는가?
        \item 인플레이션 갭의 개념을 아는가?
        \item 밀턴 프리드먼의 K-규칙과 테일러 준칙의 차이점을 설명할 수 있는가?
        \item 테일러 준칙의 구성 요소(실질 금리, 목표 인플레이션, 인플레이션 갭, 생산량 갭)를 이해하는가?
        \item 테일러 준칙에 따른 금리 결정 예시를 설명할 수 있는가?
    \end{itemize}
    \item \textbf{통화정책 전달 메커니즘}
    \begin{itemize}
        \item 연방 기금 금리 변화가 단기 및 장기 시장 금리에 미치는 영향을 설명할 수 있는가?
        \item 금리 변화가 환율, 주가, 소비, 투자에 미치는 영향을 설명할 수 있는가?
        \item 통화정책이 경제에 시차를 두고 영향을 미치는 이유를 아는가?
        \item 은행을 통한 통화정책 전달의 세 가지 채널(순이자 마진, 지급 준비금, 예금)을 설명할 수 있는가?
        \item 기업(특히 소규모 기업의 변동 금리 대출)과 가계(모기지 재융자)를 통한 통화정책 전달을 설명할 수 있는가?
    \end{itemize}
    \item \textbf{통화정책의 한계}
    \begin{itemize}
        \item 스태그플레이션의 개념과 중앙은행의 딜레마를 설명할 수 있는가?
        \item 공급 충격(예: 오일 쇼크)이 스태그플레이션을 유발하는 과정을 설명할 수 있는가?
        \item 통화정책이 공급 충격에 효과적으로 대응하기 어려운 이유를 아는가?
        \item 인플레이션 기대를 안정화하기 위한 중앙은행의 역할과 그 중요성을 이해하는가?
    \end{itemize}
\end{itemize}


\section{FAQ: 자주 묻는 질문}

초심자들이 혼동할 수 있는 질문들을 Q\&A 형식으로 정리했습니다.

\begin{itemize}
    \item \textbf{Q: Fed의 이중 책무인 '최대 고용'과 '물가 안정'은 항상 함께 달성될 수 있나요?}
    \textbf{A:} 이상적으로는 그렇지만, 현실에서는 종종 상충 관계에 놓입니다. 예를 들어, 경제를 과열시켜 고용을 늘리려 하면 인플레이션이 높아질 수 있고, 반대로 인플레이션을 잡으려 하면 고용이 위축될 수 있습니다. Fed는 이 두 목표 사이에서 균형을 찾아야 합니다.

    \item \textbf{Q: 오쿤의 법칙에서 실업률 1\% 감소가 GDP 2\% 증가로 이어진다는 것은 항상 정확한가요?}
    \textbf{A:} 오쿤의 법칙은 경험적 관계이며, '약 -2'라는 기울기는 역사적 데이터에서 관찰된 평균적인 경향입니다. 특정 시기나 국가에 따라 이 관계는 달라질 수 있으며, 정확히 2배가 아닌 근사치로 이해하는 것이 중요합니다.

    \item \textbf{Q: 필립스 곡선이 '사라졌다'는 말은 무슨 의미인가요?}
    \textbf{A:} 1970년대와 1980년대에 높은 인플레이션과 높은 실업률이 동시에 나타나는 스태그플레이션이 발생하면서, 기존의 안정적인 필립스 곡선 관계가 깨진 것처럼 보였습니다. 이는 인플레이션 기대의 변화로 인해 필립스 곡선 자체가 이동했기 때문이며, 단기적인 상충 관계는 여전히 유효하다고 해석됩니다.

    \item \textbf{Q: 테일러 준칙은 중앙은행이 반드시 따라야 하는 법적 규칙인가요?}
    \textbf{A:} 아닙니다. 테일러 준칙은 학자들이 제안한 통화정책의 '지침' 또는 '규칙' 모델입니다. 중앙은행이 금리 결정을 할 때 참고할 수 있는 유용한 프레임워크를 제공하지만, 법적 구속력은 없으며, 실제 정책 결정은 다양한 요소를 종합적으로 고려하여 이루어집니다.

    \item \textbf{Q: 통화정책이 경제에 영향을 미치는 데 왜 '시차'가 발생하나요?}
    \textbf{A:} 금리 변화가 기업의 투자 결정이나 가계의 소비 결정에 영향을 미치기까지는 시간이 걸립니다. 기업은 이미 투자 계획을 세웠을 수 있고, 가계는 대출 금리 변화를 인지하고 행동을 바꾸기까지 시간이 필요하기 때문입니다. 이러한 시차 때문에 중앙은행은 미래 경제 상황을 예측하여 선제적으로 정책을 결정해야 합니다.

    \item \textbf{Q: 공급 충격이 발생했을 때 통화정책은 왜 효과적이지 않나요?}
    \textbf{A:} 통화정책은 주로 총수요를 조절하여 경제에 영향을 미칩니다. 하지만 공급 충격(예: 유가 급등)은 생산 비용을 직접적으로 증가시켜 물가를 올리고 생산량을 줄입니다. 금리 인하로 수요를 늘려도 생산 비용 문제는 해결되지 않으며, 오히려 인플레이션을 악화시킬 수 있습니다. 따라서 공급 충격에는 통화정책보다 공급 측면의 구조적 정책이 더 효과적일 수 있습니다.
\end{itemize}


\section{빠르게 훑어보기: 1페이지 요약}

\begin{tcolorbox}[colback=myblue!10!white, colframe=myblue!50!black, title={\textbf{모듈 3 핵심 요약: 통화정책의 경제적 영향}}]
\textbf{1. Fed의 이중 책무:}
\begin{itemize}
    \item \textbf{최대 고용}: 자연 실업률 수준의 고용 달성.
    \item \textbf{물가 안정}: 낮은 인플레이션 유지 (목표 2\%).
\end{itemize}

\textbf{2. 인플레이션과 실업률의 관계:}
\begin{itemize}
    \item \textbf{오쿤의 법칙}: 실업률 $\uparrow$ $\leftrightarrow$ GDP 성장률 $\downarrow$ (음의 상관관계).
    \item \textbf{필립스 곡선}: 인플레이션 $\leftrightarrow$ 실업률 (단기적 상충 관계).
    \item \textbf{인플레이션 기대}: 필립스 곡선을 이동시키는 핵심 요인. 기대 $\uparrow$ $\rightarrow$ 곡선 바깥 이동.
\end{itemize}

\textbf{3. 통화정책 규칙:}
\begin{itemize}
    \item \textbf{생산량 갭}: 실제 GDP와 잠재 GDP의 차이.
    \item \textbf{인플레이션 갭}: 실제 인플레이션과 목표 인플레이션의 차이.
    \item \textbf{테일러 준칙}: 금리 = 실질금리 + 목표인플 + $\alpha$(인플 갭) + $\beta$(생산량 갭). 경제 상황에 따라 금리 조정.
\end{itemize}

\textbf{4. 통화정책 전달 메커니즘:}
\begin{itemize}
    \item \textbf{금리 채널}: 정책 금리 $\rightarrow$ 시장 금리 $\rightarrow$ 투자/소비/환율/주가.
    \item \textbf{은행 대출 채널}: 정책 금리 $\rightarrow$ 순이자 마진, 지급 준비금, 예금 $\rightarrow$ 은행 신용 공급.
    \item \textbf{기업/가계 채널}: 정책 금리 $\rightarrow$ 변동 금리 대출 비용, 모기지 재융자 $\rightarrow$ 기업 투자, 가계 소비.
    \item \textbf{시차}: 정책 효과는 경제에 시차를 두고 나타남.
\end{itemize}

\textbf{5. 통화정책의 한계:}
\begin{itemize}
    \item \textbf{공급 충격 (예: 오일 쇼크)}: 스태그플레이션 유발. 통화정책은 수요 조절에 강하고 공급 충격 대응에 약함.
    \item \textbf{인플레이션 기대 관리}: 공급 충격 시, 인플레이션 기대를 안정화하기 위해 일시적 고실업 감수 필요 (볼커 의장 사례).
\end{itemize}
\end{tcolorbox}


\section{개요: 통화 정책의 한계}

\begin{summarybox}[title={문서 전체 핵심 요약}]
이 문서는 통화 정책이 효과적으로 작동하지 않는 두 가지 주요 상황을 탐구합니다.
첫째, \textbf{고용 없는 회복(Jobless Recovery)} 현상으로, 경기 침체 후 GDP는 성장하지만 노동 시장은 느리게 회복되는 경우입니다.
이는 높은 불확실성이나 경제의 구조적 변화로 인해 발생할 수 있으며, 통화 정책만으로는 해결하기 어렵습니다.
둘째, \textbf{완화적 통화 정책의 부작용}으로, 장기간 낮은 금리가 과도한 위험 감수를 유발하여 금융 불안정이나 자산 거품을 초래할 수 있습니다.
특히 2000년대 초반의 주택 시장 거품 사례를 통해 이러한 위험을 상세히 분석합니다.
\end{summarybox}


\section{통화 정책의 한계: 고용 없는 회복과 위험 감수 유인}

\subsection{개요}
이번 강의에서는 통화 정책이 기대만큼 효과를 발휘하지 못했던 시기들을 살펴봅니다. 특히 두 가지 주요 관련 문제에 집중할 것입니다. 첫째는 '고용 없는 회복(Jobless Recovery)'이라는 개념입니다. 이 용어는 1991년 경기 침체 이후의 회복을 설명하기 위해 처음 사용되었습니다. 연방준비제도(Federal Reserve)의 상당한 금리 인하에도 불구하고, 노동 시장은 예상보다 느리게 회복되었습니다. 우리는 이러한 현상이 왜 발생하는지 이해하려고 노력할 것입니다.

둘째는 완화적 통화 정책의 의도치 않은 결과인 '과도한 위험 감수(Excessive Risk-Taking)'입니다. 금리 인하는 경제를 활성화할 수 있지만, 만약 금리가 너무 오랫동안 너무 낮게 유지된다면 어떻게 될까요? 일부에서는 이것이 금융 불안정을 초래할 수 있다고 주장합니다. 우리는 이 주장을 자세히 검토할 것입니다。

\subsection{고용 없는 회복 (Jobless Recoveries)}

\subsubsection{개념 및 역사적 배경}
역사적으로 경기 침체가 끝난 후에는 실업률이 급격히 하락하는 경향이 있었습니다. 예를 들어, 1982년 경기 침체 이후의 상황을 보면 이를 명확히 알 수 있습니다. 그러나 1991년 경기 침체 이후에는 대조적으로 실업률이 점진적으로만 하락했습니다. 이는 연방준비제도가 3년 동안 금리를 낮게 유지하고 1992년에서 1995년 사이에 GDP가 5\% 이상 성장했음에도 불구하고 발생한 현상입니다. 이러한 노동 시장의 느린 회복은 경제학자들을 당혹스럽게 했습니다.

\begin{examplebox}[title={1991년 경기 침체 후 고용 없는 회복}]
1991년 경기 침체 이후, 연방준비제도는 경제 활성화를 위해 금리를 크게 인하했습니다. 일반적으로 금리 인하는 기업의 투자와 고용을 촉진하여 실업률을 빠르게 낮추는 효과가 있습니다. 실제로 1992년부터 1995년까지 미국의 GDP는 5\% 이상 성장하며 경제는 회복세를 보였습니다. 하지만 이러한 경제 성장에도 불구하고, 실업률은 과거의 경기 회복기처럼 급격히 떨어지지 않고 매우 점진적으로 하락했습니다. 이는 경제가 성장해도 일자리가 충분히 늘어나지 않는 '고용 없는 회복'의 전형적인 예시로 기록되었습니다.
\end{examplebox}

2001년 경기 침체 이후에도 동일한 패턴이 나타났습니다. 연방준비제도는 경기 침체에 대응하여 연방기금금리(federal funds rate)를 6.5\%에서 1\%로 대폭 인하했으며, 2004년 여름까지 연방기금금리를 인상하지 않았습니다. 그러나 실업률은 2003년 여름이 되어서야 하락하기 시작했고, 그마저도 하락세는 완만했습니다. 이처럼 낮은 금리 수준에서의 통화 정책 완화는 전례 없는 일이었습니다. 이러한 이유로 경제학자들은 2000년대 초반을 '고용 없는 회복' 시기로 지칭합니다.

\subsubsection{2001년 고용 없는 회복의 특징}
2001년 고용 없는 회복 시기의 노동 시장 데이터를 통해 어떤 특징이 나타났는지 살펴보겠습니다. 노동 시장 지표들은 근로자들이 일자리를 찾는 데 상당한 어려움을 겪었음을 시사했습니다.

\begin{itemize}
    \item \textbf{노동 시장 참여율 감소:} 특히 젊은 근로자들 사이에서 노동 시장 참여율이 감소했습니다. 이는 일자리를 찾으려는 의지를 가진 사람 자체가 줄어들었음을 의미할 수 있습니다.
    \item \textbf{실업 기간 증가:} 실업 상태가 지속되는 기간이 상당히 길어졌습니다. 이는 새로운 일자리를 찾는 데 더 많은 시간이 소요되었음을 보여줍니다.
    \item \textbf{구인 광고 지수 하락:} 일자리 공고의 척도인 구인 광고 지수가 경기 침체 이전 수준보다 낮게 유지되었습니다. 이는 기업들이 예상만큼 적극적으로 채용하지 않았다는 증거입니다.
\end{itemize}
이러한 증거들은 기업들이 예상만큼 많이 고용하지 않았음을 시사합니다.

또한, 2003년 9월 말 현재의 고용 통계 조사(Current Employment Statistics survey)에 따르면, 2001년 3월 경기 침체 시작 이후 비농업 부문 고용이 약 280만 개 감소했습니다. 이 중 제조업 부문이 가장 큰 타격을 입어, 전체 감소한 280만 개 일자리 중 약 240만 개가 제조업에서 발생했습니다.

\subsubsection{느린 노동 시장 회복의 원인}
그렇다면 왜 노동 시장의 반응은 그렇게 느렸을까요? 그리고 왜 제조업 부문이 가장 큰 타격을 입었을까요? 여기에는 두 가지 주요 설명이 있습니다.

\begin{enumerate}
    \item \textbf{경제에 대한 높은 불확실성 (High Uncertainty):}
    기업들은 경제 상황에 대한 불확실성이 높을 때 투자나 고용을 꺼리는 경향이 있습니다. 2000년대 초반에는 9.11 테러 공격과 그 여파(아프가니스탄 및 이라크 전쟁 포함)가 가장 명백한 불확실성의 원천이었습니다. 이러한 외부 충격은 기업들이 미래를 예측하기 어렵게 만들었고, 이는 고용 결정에 부정적인 영향을 미쳤습니다.

    \item \textbf{미국 경제의 구조적 변화 (Structural Changes):}
    제조업 고용 손실은 미국 경제의 구조적 변화 때문이었습니다. 구조적 변화는 특정 산업의 영구적인 쇠퇴를 의미합니다. 제조업의 쇠퇴는 새로운 현상이 아니었습니다. 1980년대 초반부터 제조업 부문 고용 근로자 비중은 급격히 감소하기 시작했습니다. 이는 부분적으로 국제 경쟁 심화에 의해 주도되었습니다.

    \begin{examplebox}[title={구조적 변화의 예시: 제조업 쇠퇴}]
    1980년대 오일 위기 이후, 일본 및 서유럽 자동차 제조업체들이 미국 시장으로 확장하면서 자동차 시장의 경쟁이 심화된 것이 가장 두드러진 사례입니다. 이러한 국제 경쟁은 미국 제조업의 일자리 감소에 기여했습니다. 구조적 변화는 일반적으로 느리게 진행되는 과정이므로, 정책 입안자들이 실시간으로 그 결과를 측정하기 어렵게 만듭니다.

    2000년대 초반의 논의를 되돌아보면, 한 가지 요인이 눈에 띄게 빠져 있습니다. 2001년 12월 11일, 중국이 세계무역기구(WTO)에 가입했습니다. 그 결과, 미국의 무역 정책은 중국산 수입품에 대한 잠재적인 관세 인상을 철회했습니다. 최근 연구에 따르면, 제조업 일자리 손실은 중국과의 수입 경쟁에 직면한 미국 기업들에 집중되었습니다. 이는 중국의 WTO 가입이 미국 제조업 일자리 감소에 상당한 영향을 미쳤음을 시사합니다.
    \end{examplebox}
\end{enumerate}

\subsubsection{통화 정책의 한계점}
불확실성이 높거나 구조적 변화가 진행 중일 때 통화 정책은 효과적일 수 있을까요?

\begin{itemize}
    \item \textbf{불확실성 감소:} 일반적으로 통화 정책은 현재와 미래에 완화적인 통화 정책을 약속함으로써 미래 경제 발전에 대한 불확실성을 줄일 수 있습니다. 이는 '선제적 안내(forward guidance)'의 사용을 포함합니다. 그러나 2000년대 초반의 경우, 불확실성은 테러 공격과 그에 따른 전쟁으로 인해 발생했습니다. 이러한 유형의 불확실성은 통화 정책으로 줄이기 어려웠을 것입니다.

    \item \textbf{구조적 변화 대응:} 마찬가지로 통화 정책은 구조적 변화로 인한 고용 효과를 줄이는 데 거의 할 수 있는 것이 없습니다. 쇠퇴하는 산업에서 사라진 일자리는 결국 성장하는 산업(예: 서비스 또는 기술 산업)의 일자리로 대체될 것으로 예상됩니다. 그러나 이 과정은 시간이 걸릴 수 있습니다. 새로운 일자리가 창출되어야 하고, 근로자들은 재배치되고 재교육을 받아야 합니다.

    \item \textbf{최적의 작동 조건:} 통화 정책은 경기 순환적 변동(cyclical fluctuation)에 대응할 때 가장 잘 작동합니다. 구조적 변화가 실업률을 주도할 때는 할 수 있는 것이 거의 없습니다.
\end{itemize}

\begin{cautionbox}{통화 정책의 한계}
통화 정책은 경제의 단기적인 경기 변동(예: 경기 침체와 회복)을 완화하는 데 효과적입니다. 하지만 9.11 테러와 같은 외부 충격으로 인한 불확실성이나, 제조업 쇠퇴와 같은 장기적인 산업 구조 변화로 인한 실업 문제에는 직접적인 해결책을 제공하기 어렵습니다. 이러한 문제들은 통화 정책보다는 재정 정책, 교육 및 훈련 프로그램, 산업 정책 등 다른 정책 수단이 더 효과적일 수 있습니다.
\end{cautionbox}


\subsection{완화적 통화 정책의 위험: 과도한 위험 감수}
이제 이 세션의 나머지 부분에서는 완화적 통화 정책의 위험을 살펴보겠습니다. 2001년, 투자 운용사 PIMCO의 경제학자 폴 맥컬리(Paul McCully)는 연방준비제도가 설정한 극도로 낮은 연방기금금리가 주택 가격 상승을 유발할 수 있다고 예측했습니다.

\subsubsection{주택 시장 거품 형성 과정}
맥컬리의 예측은 현실이 되었습니다. 다음은 주택 시장 거품이 어떻게 형성되었는지에 대한 설명입니다.

\begin{enumerate}
    \item \textbf{모기지 금리 하락:}
    2000년 8\% 이상이던 모기지 금리는 2003년 여름 5\%로 떨어졌습니다.

    \item \textbf{주택 수요 및 가격 상승:}
    모기지 금리가 낮아지면 사람들이 새 집이나 더 큰 집을 구매할 여력이 생깁니다. 이는 주택 수요를 증가시키고 주택 가격을 상승시킵니다. 실제로 2000년에서 2006년 사이에 주택 가격은 60\% 이상 상승했습니다.

    \item \textbf{건설 경기 호황 및 부의 효과:}
    주택 가격 상승의 결과로 건설 경기가 호황을 누렸습니다. 또한, '부의 효과(wealth effect)'로 인해 주택 소유자들은 모기지 재융자(mortgage refinancing)와 주택 담보 대출(home equity lines of credit, HELOC)을 이용하여 주택을 담보로 더 많은 돈을 빌렸습니다. 이렇게 얻은 자금은 자동차나 가전제품과 같은 소비재에 지출되었습니다.

    \item \textbf{신용 공급의 극적인 증가:}
    주택 담보 대출(HELOC)은 2000년 약 1,500억 달러에서 2005년 말까지 6,500억 달러로 극적으로 증가했습니다. 이러한 신용 공급의 증가는 주택 건설, 소비자 지출 및 일자리 창출을 촉진했습니다.
\end{enumerate}

\begin{examplebox}[title={모기지 금리 하락과 주택 시장의 변화}]
\textbf{2000년:} 모기지 금리 8\% 이상. 주택 구매 부담이 컸음.
\textbf{2003년 여름:} 모기지 금리 5\%로 하락.
\begin{itemize}
    \item \textbf{구매력 증가:} 같은 월 상환액으로 더 비싼 집을 살 수 있게 됨.
    \item \textbf{수요 증가:} 더 많은 사람이 주택 시장에 진입하거나 더 큰 집으로 옮기려 함.
    \item \textbf{가격 상승:} 수요가 공급을 초과하면서 주택 가격이 2000년 대비 2006년까지 60\% 이상 폭등.
\end{itemize}
\textbf{부의 효과:} 주택 가격이 오르자, 주택 소유자들은 자신이 부자가 되었다고 느끼고, 주택 담보 대출을 통해 현금을 인출하여 소비를 늘렸습니다. 이는 경제 활성화에 기여했지만, 동시에 과도한 부채와 자산 거품의 위험을 키웠습니다.
\end{examplebox}

\subsubsection{금융 불안정 경고}
2005년, 당시 IMF 수석 경제학자 라구람 라잔(Raghuram Rajan)은 주택 거품에 대한 우려를 표명했습니다. 그는 낮은 금리와 금융 혁신이 금융 시장 참여자들이 점점 더 많은 위험을 감수하는 상황을 만들었다고 주장했습니다. 그는 이러한 위험이 금융 시스템의 안정성을 위협할 수 있다고 결론지었습니다. 돌이켜보면, 그의 분석은 정확했습니다.

\subsubsection{일반적인 원칙 및 중앙은행의 역할}
기억해야 할 더 일반적인 요점은 낮은 금리가 위험 축적(buildup of risk)으로 이어질 수 있다는 것입니다. 낮은 금리는 과도한 차입을 장려할 수 있으며, 이는 주택 시장 및 금융 시장에서 자산 가격 거품을 부채질할 수 있습니다. 낮은 금리와 위험 감수 사이의 연관성은 미국뿐만 아니라 모든 주요 경제권에서 나타났습니다.

위험 감수를 장려하는 것은 경제를 활성화하는 한 가지 방법입니다. 그러나 중앙은행은 완화적 통화 정책이 금융 시스템의 전반적인 건전성에 미치는 영향을 염두에 두어야 합니다. 2008년 금융 위기 이후, 금융 안정성(financial stability)은 연방준비제도의 임무에 추가되었습니다. 연방준비제도는 이제 경제 내의 위험 축적을 모니터링합니다.

\begin{cautionbox}{과도한 위험 감수의 위험}
낮은 금리는 기업과 가계가 돈을 빌리는 비용을 낮춰 투자를 늘리고 소비를 촉진하는 긍정적인 효과가 있습니다. 하지만 금리가 너무 오랫동안 너무 낮게 유지되면, 사람들은 더 높은 수익을 추구하기 위해 과도한 위험을 감수하게 됩니다. 이는 주식, 부동산 등 자산 시장에 거품을 형성하고, 부채 수준을 위험하게 높여 금융 시스템 전체의 안정성을 해칠 수 있습니다. 2008년 글로벌 금융 위기는 이러한 위험이 현실화된 대표적인 사례입니다.
\end{cautionbox}


\subsection{결론 및 요약}
우리는 통화 정책의 두 가지 잠재적 한계를 논의했습니다.

\begin{enumerate}
    \item \textbf{구조적 변화 또는 외부 사건으로 인한 불확실성:}
    통화 정책은 구조적 변화나 외부 사건에 대한 불확실성이 실업률을 증가시킬 때 할 수 있는 것이 거의 없습니다. 이러한 문제들은 통화 정책의 범위를 넘어섭니다.

    \item \textbf{완화적 통화 정책의 부작용:}
    2000년대 초반의 고용 없는 회복 이후의 완화적 통화 정책은 미국에서 주택 가격 거품을 부채질했을 수 있습니다. 이 경험은 더 일반적인 문제를 강조합니다. 낮은 금리는 위험 감수를 장려하고 경제 내에 위험을 축적시킬 수 있습니다.
\end{enumerate}


\section{용어 정리}

\begin{table}[h!]
\centering
\caption{주요 용어 정리}
\label{tab:glossary}
\begin{adjustbox}{width=\textwidth,center}
\begin{tabular}{lllc}
\toprule
\textbf{용어} & \textbf{쉬운 설명} & \textbf{원어} & \textbf{비고} \\
\midrule
고용 없는 회복 & 경제는 성장하지만 일자리는 느리게 늘어나는 현상 & Jobless Recovery & 1991년, 2001년 침체 후 관찰 \\
연방준비제도 & 미국의 중앙은행 & Federal Reserve & 통화 정책을 결정 \\
연방기금금리 & 은행 간 초단기 대출 금리 목표치 & Federal Funds Rate & 중앙은행의 주요 정책 금리 \\
완화적 통화 정책 & 금리를 낮춰 경제 활동을 촉진하는 정책 & Loose/Accommodative Monetary Policy & 경기 부양 목적 \\
과도한 위험 감수 & 낮은 금리로 인해 투자자들이 지나치게 위험한 투자를 하는 경향 & Excessive Risk-Taking & 금융 불안정의 원인 \\
구조적 변화 & 특정 산업의 영구적인 쇠퇴나 산업 구조의 변화 & Structural Change & 제조업 쇠퇴 등이 해당 \\
선제적 안내 & 중앙은행이 미래 통화 정책 방향을 미리 알려주는 것 & Forward Guidance & 시장의 불확실성 감소 목적 \\
부의 효과 & 자산 가치 상승으로 소비가 증가하는 현상 & Wealth Effect & 주택 가격 상승 시 발생 \\
모기지 재융자 & 기존 주택 담보 대출을 새로운 조건으로 다시 받는 것 & Mortgage Refinancing & 낮은 금리 활용 \\
주택 담보 대출 & 주택을 담보로 돈을 빌리는 신용 한도 & Home Equity Line of Credit (HELOC) & 주택 가치 상승 시 활용 증가 \\
금융 안정성 & 금융 시스템이 충격에 강하고 원활하게 작동하는 상태 & Financial Stability & 중앙은행의 중요한 목표 \\
\bottomrule
\end{tabular}
\end{adjustbox}
\end{table}


\section{빠르게 훑어보기: 핵심 요약}

\begin{summarybox}[title={통화 정책의 한계: 1페이지 요약}]
\textbf{1. 고용 없는 회복 (Jobless Recovery)}
\begin{itemize}
    \item \textbf{정의:} GDP는 성장하지만 실업률은 느리게 하락하는 현상.
    \item \textbf{사례:} 1991년 및 2001년 미국 경기 침체 후 관찰.
    \textbf{원인:}
    \begin{itemize}
        \item \textbf{높은 불확실성:} 9.11 테러, 전쟁 등 외부 충격.
        \item \textbf{구조적 변화:} 제조업 쇠퇴, 국제 경쟁(중국 WTO 가입).
    \end{itemize}
    \item \textbf{통화 정책의 한계:} 경기 순환적 문제에는 효과적이나, 외부 충격으로 인한 불확실성이나 구조적 실업에는 대응하기 어려움.
\end{itemize}

\textbf{2. 완화적 통화 정책의 위험: 과도한 위험 감수}
\begin{itemize}
    \item \textbf{문제점:} 금리가 너무 낮게, 너무 오랫동안 유지될 때 발생.
    \item \textbf{메커니즘:}
    \begin{itemize}
        \item 낮은 모기지 금리 $\rightarrow$ 주택 수요 증가 $\rightarrow$ 주택 가격 급등 (2000-2006년 60\% 이상 상승).
        \item 부의 효과 $\rightarrow$ 주택 담보 대출(HELOC) 증가 (2000년 1,500억 달러 $\rightarrow$ 2005년 6,500억 달러).
        \item 과도한 차입 및 자산 거품 형성.
    \end{itemize}
    \item \textbf{경고:} 라구람 라잔(2005년)은 낮은 금리와 금융 혁신이 금융 불안정을 초래할 수 있다고 경고.
    \item \textbf{중앙은행의 역할:} 2008년 금융 위기 이후, 금융 안정성 감시 및 위험 축적 모니터링이 연방준비제도의 주요 임무에 추가됨.
\end{itemize}
\end{summarybox}


%=======================================================================
% Module 8: 용어 정리
%=======================================================================
\chapter{용어 정리}
\label{ch:module8}

\metainfo{FIN-571: Financial Institutions \& Monetary Policy}{모듈 4: 금융 위기 (파트 1/2)}{Prof. Ralf Meisenzahl}{이 문서는 '중앙은행과 통화 정책' 모듈 4의 첫 번째 부분으로, 금융 위기의 본질과 중앙은행의 역할에 초점을 맞춥니다. 우리는 1907년 공황, 대공황, 2008년 금융 위기 중 자동차 대출 시장 사례를 통해 금융 위기가 발생하는 원인과 그 파급 효과를 심층적으로 탐구합니다. 특히 중앙은행이 '최종 대부자(Lender of Last Resort)'로서 어떻게 금융 시장을 안정화하고 위기를 완화하는지, 그리고 그 과정에서 발생하는 정책적 한계와 비용은 무엇인지 살펴봅니다. 이 학습 자료는 비전공자도 금융 위기의 복잡한 메커니즘과 정책 대응을 명확하게 이해할 수 있도록 구성되었습니다.}


\section{용어 정리}
금융 위기 관련 핵심 용어들을 쉽게 이해할 수 있도록 정리했습니다.

\begin{adjustbox}{width=\textwidth,center}
\begin{tabular}{lllc}
\toprule
\textbf{용어} & \textbf{쉬운 설명} & \textbf{원어} & \textbf{비고} \\
\midrule
뱅크런 & 많은 예금자가 동시에 은행에서 돈을 인출하려는 현상 & Bank Run & 은행 파산의 주요 원인 \\
만기 불일치 & 은행이 단기로 돈을 빌려(예금) 장기로 빌려주는(대출) 구조 & Maturity Mismatch & 은행 사업 모델의 핵심이자 취약점 \\
유동성 부족 & 은행에 당장 현금이 부족하지만, 자산 가치는 충분한 상태 & Illiquid & 일시적인 현금 부족 \\
지급 불능 & 은행의 자산 가치가 부채보다 적어 파산 상태인 경우 & Insolvent & 자본 잠식 상태 \\
숏 스퀴즈 & 주식을 빌려 판(공매도) 투자자들이 주가 상승으로 손실을 줄이기 위해 & Short Squeeze & 주가 급등을 유발하는 현상 \\
& 다시 주식을 사들이면서 주가가 더 오르는 현상 & & \\
최종 대부자 & 금융 시스템 위기 시, 다른 곳에서 돈을 빌릴 수 없을 때 & Lender of Last Resort & 중앙은행의 핵심 기능 \\
& 마지막으로 돈을 빌려주는 기관 (주로 중앙은행) & & \\
금본위제 & 통화 가치를 금에 고정하여, 화폐를 금으로 교환할 수 있도록 하는 제도 & Gold Standard & 통화량 조절에 제약이 있음 \\
자산유동화기업어음 & 자동차 대출 같은 자산을 담보로 발행하는 단기 어음 & Asset-Backed Commercial Paper (ABCP) & 비은행 금융기관의 주요 자금 조달원 \\
캡티브 금융 회사 & 자동차 제조사 등이 소유한 금융 회사로, 자사 제품 구매자에게 & Captive Finance Company & 자동차 대출 시장의 주요 플레이어 \\
& 대출을 제공함 & & \\
할인 창구 & 중앙은행이 시중 은행에 단기 자금을 대출해주는 제도 & Discount Window & 전통적인 중앙은행의 유동성 공급 수단 \\
프라이머리 딜러 & 중앙은행의 공개 시장 운영 거래 상대방이 되는 주요 투자은행 & Primary Dealer & 금융 시장의 핵심 중개자 \\
재정 정책 & 정부가 세금, 지출 등을 통해 경제에 영향을 미치는 정책 & Fiscal Policy & 정부 주도의 경제 안정화 정책 \\
통화 정책 & 중앙은행이 금리, 통화량 등을 조절하여 경제에 영향을 미치는 정책 & Monetary Policy & 중앙은행 주도의 경제 안정화 정책 \\
\bottomrule
\end{tabular}
\end{adjustbox}


\section{모듈 4: 금융 위기}

\subsection{Lesson 4-0.1: 개요}
안녕하세요, 여러분. 이번 모듈에서는 중앙은행이 '최종 대부자(Lender of Last Resort)'로서 수행하는 역할에 집중합니다. 이를 위해 우리는 여러 금융 위기 사례를 깊이 있게 연구하고, 위기가 발생한 근본적인 원인을 논의할 것입니다. 이어서 중앙은행이 이러한 위기에 어떻게 정책적으로 대응했는지 살펴봅니다. 중앙은행의 노력이 성공을 거두기도 하지만, 단독으로 모든 문제를 해결하는 데는 한계가 있다는 점을 알게 될 것입니다.

\begin{tcolorbox}[title={\textbf{모듈 4 학습 목표}}]
\begin{itemize}
    \item 다양한 금융 위기 사례를 통해 위기 발생 원인을 이해합니다.
    \item 중앙은행의 '최종 대부자' 역할과 그 중요성을 파악합니다.
    \item 금융 위기에 대한 중앙은행의 정책 대응 방안을 학습합니다.
    \item 중앙은행 단독 정책의 한계와 재정 정책의 필요성을 인식합니다.
\end{itemize}
\end{tcolorbox}

우리가 처음 살펴볼 위기는 1907년 공황(Panic of 1907)입니다. 이 위기는 20세기 최초의 주요 금융 위기였으며, 연방준비제도(Federal Reserve system)가 설립되기 전에 발생한 마지막 미국 은행 공황이기도 합니다. 이 사례를 통해 최종 대부자로서의 중앙은행이 왜 필요한지 명확하게 이해할 수 있습니다.

다음으로 대공황(Great Depression)을 연구합니다. 대공황은 오늘날 존재하는 사회 안전망(social safety net)의 대부분이 당시에는 없었다는 점에서 또 다른 유용한 사례입니다. 이를 통해 정부 개입이 없을 경우 경제가 어떻게 될 수 있는지 배울 수 있습니다.

그다음에는 금융 위기에 대한 국제적인 관점을 취하고, 위기의 원인을 파악하는 데 도움이 되는 공통된 주제들을 식별합니다. 이 거시적인 관점은 단기 자금이나 취약한 자금 조달(phone funding)로 조달되고 장기 프로젝트나 지속 불가능한 정부 정책에 투자된 '신용 붐(credit booms)'이 금융 위기를 예고했음을 강조합니다.

2008년 금융 위기 당시 미국 자동차 대출 시장을 연구함으로써 단기 자금 조달에 의존하는 위험을 다시 한번 강조합니다. 이 시기에 비은행 대출 기관들은 단기 자금 시장에 크게 의존하여 자동차 대출을 제공했습니다. 단기 자금 시장이 고갈되자 자동차 판매가 급감했습니다.

이 모든 사례를 살펴본 후, 여러분은 신용 붐이 항상 금융 위기와 경기 침체로 이어지는지 궁금해할 수 있습니다. 짧은 답변은 '아니오'입니다. 다음 강의에서 논의하겠지만, 이는 신용 붐이 어떤 부문에서 발생하는지에 따라 달라집니다. 예를 들어, 제조업 부문의 신용 붐은 적어도 역사적으로는 종종 지속 가능했습니다.

금융 위기의 원인과 결과를 검토한 후, 우리는 정책 대응으로 넘어갑니다. 구체적으로, 중앙은행이 뱅크런과 신용 공급의 급격한 감소에 어떻게 대응할 수 있는지 살펴봅니다. 이를 위해 2008년 금융 위기와 COVID-19 팬데믹에 대한 연방준비제도의 대응을 연구하여 '최종 대부자' 개념을 설명합니다. 즉, 연방준비제도는 다른 누구도 대출하지 않을 때 대출합니다.

통화 정책만이 금융 위기에 대응하는 유일한 방법은 아닙니다. 정부 이전(government transfers)과 지출을 포함하는 재정 정책(Fiscal Policy)도 경제를 도울 수 있습니다. 따라서 마지막 강의에서는 금융 위기에 대한 재정적 대응을 논의할 것입니다.


\section{Lesson 4-1: 금융 위기 이해하기}

\subsection{Lesson 4-1.1: 뱅크런과 1907년 공황}
뱅크런(Bank Run)과 1907년 공황(Panic of 1907)에 대한 수업에 오신 것을 환영합니다. 우리는 두 단계로 진행할 것입니다. 첫째, 뱅크런이 무엇이며 은행이 왜 뱅크런에 취약한지 이해할 것입니다. 둘째, 우리의 첫 번째 주요 사례 연구인 1907년 공황을 살펴볼 것입니다. 이것은 미국 역사상 가장 두드러진 금융 위기 중 하나였습니다. 대규모 은행 공황이었으며, 현대 연방준비제도 시스템의 형성을 이끌었습니다. 우리는 공황이 왜 발생했으며 어떻게 해결되었는지 이해할 것입니다. 이러한 역사적 사건을 이해하는 것은 즐겁고 매우 유익할 수 있습니다. 정확한 역사적 세부 사항에 대해서는 너무 걱정하지 마십시오. 많은 세부 사항을 언급하겠지만, 이는 주로 재미를 위한 것이며, 서술과 우리가 배울 수 있는 교훈에 집중하십시오.

\subsubsection*{은행은 왜 뱅크런에 취약한가?}
은행은 단기적으로 예금을 받아 돈을 빌리고, 장기적으로 대출을 해줍니다. 이것은 자산(대출)과 부채(예금) 사이에 '만기 불일치(Maturity Mismatch)'를 만듭니다. 이러한 만기 변환(maturity transformation)은 은행 사업 모델의 핵심 특징입니다.

\begin{tcolorbox}[title={\textbf{핵심 개념: 만기 불일치 (Maturity Mismatch)}}]
\textbf{한 줄 요약:} 은행이 단기 자금(예금)을 조달하여 장기 자산(대출)에 투자하는 구조입니다.

\textbf{직관적 비유:} 친구에게 "내일 갚을게" 하고 돈을 빌려, 그 돈으로 1년 뒤에나 수익이 나는 사업에 투자하는 것과 같습니다. 친구가 갑자기 "지금 당장 돈 갚아!"라고 하면 곤란해지겠죠?

\textbf{기술적 설명:} 예금자들은 언제든지 자신의 자금을 인출할 수 있지만, 은행이 대출해준 돈은 장기간 묶여 있습니다. 이 구조는 은행이 정상적으로 운영될 때는 효율적이지만, 예금자들이 동시에 돈을 인출하려 할 때 은행이 현금 부족에 직면하게 만듭니다.
\end{tcolorbox}

그렇다면 뱅크런은 왜 발생할까요? 예금자들은 언제든지 자신의 자금을 인출할 수 있다는 점을 기억하십시오. 어떤 충격이나 나쁜 소식이 발생하면, 이는 예금자들에게 위험을 초래합니다. 그들은 안전하게 돈을 인출하기로 결정할 수 있습니다. 아마도 은행 앞에 사람들이 줄 서 있는 것을 보았기 때문일 수도 있습니다. 뱅크런에 동참하는 것은 보통 합리적인 행동입니다. 만약 예금자가 줄에 합류하면, 돈을 돌려받을 기회가 있습니다. 줄의 맨 앞에 있는 사람들은 보통 돈을 돌려받을 것입니다. 만약 그들이 뱅크런에 동참하지 않으면, 돈을 잃을 수도 있습니다. 줄의 맨 뒤에 있는 사람들은 아무것도 받지 못할 수도 있습니다.

\subsubsection*{뱅크런의 결과}
많은 예금자가 동시에 도착하면, 은행은 충분한 현금을 가지고 있지 않을 수 있습니다. 대부분의 예금은 장기 투자에 묶여 있습니다. 이것이 만기 변환이 작동하는 방식입니다. 이 경우 은행은 '유동성 부족(illiquid)' 상태가 됩니다. 현금이 바닥난 것입니다.

\begin{tcolorbox}[title={\textbf{핵심 개념: 유동성 부족(Illiquid) vs. 지급 불능(Insolvent)}}]
\textbf{한 줄 요약:} 유동성 부족은 현금이 일시적으로 없는 상태이고, 지급 불능은 자산보다 빚이 더 많은 파산 상태입니다.

\textbf{직관적 비유:}
\begin{itemize}
    \item \textbf{유동성 부족:} 지갑에 현금은 없지만, 비싼 집이나 차 같은 자산은 충분히 가지고 있는 상태입니다. 당장 현금이 필요할 때 자산을 팔아야 하지만, 팔 시간이 없거나 제값을 받기 어렵습니다.
    \item \textbf{지급 불능:} 지갑에 현금도 없고, 가지고 있는 모든 자산을 팔아도 빚을 다 갚을 수 없는 상태입니다. 완전히 파산한 것입니다.
\end{itemize}

\textbf{기술적 설명:}
\begin{itemize}
    \item \textbf{유동성 부족 (Illiquid):} 은행은 현금이 부족하지만, 그 투자는 여전히 수익성이 있을 수 있습니다. 유동성 부족 은행은 여전히 사업을 영위할 수 있으며, 지급 능력이 있을 수 있습니다.
    \item \textbf{지급 불능 (Insolvent):} 은행은 자본(equity)이 바닥나 파산 상태입니다.
\end{itemize}
\end{tcolorbox}

가장 최악의 시나리오는 건전한 은행이 뱅크런 때문에 파산하는 경우입니다. 이는 은행이 예금 인출 요구를 충족하기 위해 자산을 대폭 할인하여 강제로 매각해야 할 때 발생할 수 있습니다. 우리의 사례 연구에서 이러한 역학 관계가 실제로 어떻게 작용하는지 보게 될 것입니다. 그리고 이러한 최악의 시나리오를 방지하기 위해 다양한 금융 규제와 정부 기관이 존재하는 이유를 알게 될 것입니다.

\subsubsection*{사례 연구: 1907년 공황}
이제 우리의 사례 연구인 1907년 공황을 살펴보겠습니다. 1907년 공황은 뱅크런에 대해 배우는 데 유용한 사건입니다. 당시 현대 금융 시스템의 두 가지 특징이 없었습니다. 첫째, 중앙은행이 없었습니다. 이는 은행들이 정부 지원을 받는 비상 대출에 접근할 수 없었다는 것을 의미합니다. 둘째, 예금 보험이 없었습니다. 예금 보험이 없었기 때문에 예금자들은 훨씬 더 공황 상태에 빠졌습니다. 은행이 파산하면 평생 모은 돈을 잃을 수도 있었기 때문입니다. 이 역사적 사건은 뱅크런을 방지하기 위해 고안된 현대 정부 개입이 없을 때 어떤 일이 일어날 수 있는지 이해하는 데 도움이 됩니다. 이 금융 위기가 어떻게 전개되었는지 살펴보겠습니다.

대부분의 위기와 마찬가지로, 단일 원인이 아니라 나쁜 소식이 쌓여 뱅크런으로 이어졌습니다. 이야기는 1906년 샌프란시스코 지진으로 시작됩니다. 이 지진은 GDP의 약 1.5\%에 해당하는 피해를 입혔습니다. 이러한 대규모 충격은 즉시 금융 시스템 전체에 파급되었습니다. 보험사들은 보험금 청구를 충족하기 위해 자산을 매각해야 했습니다. 예금자들은 은행에서 돈을 인출했고, 뉴욕과 런던의 주식 시장은 하락했습니다. 이러한 추세는 계속되었습니다. 1907년 3월 며칠 동안 주식은 10\% 이상 하락했습니다. 단기 자금 시장(money markets)은 스트레스를 받았고, 이자율은 며칠 동안 2.5\%에서 10\% 이상으로 치솟았습니다. 이 사건은 단 하루 만에 일어난 것이 아니라 며칠에 걸쳐 발생했기 때문에 "조용한 폭락(the silent crash)"이라고 불리기도 합니다.

두 번째 충격은 1907년 작황이 전년도만큼 풍성하지 않아 수출이 감소한 것이었습니다. 미국 재무부는 만기가 도래하는 채권을 재융자할 수 없게 되자 은행에서 3천만 달러를 인출했습니다. 정부와 예금자들이 은행에서 더 많은 돈을 인출하면서 통화 공급이 위축되고 신용이 부족해지며 이자율이 더욱 상승했습니다. 요컨대, 1907년 여름 동안 금융 상황은 상당히 악화되었습니다.

이러한 배경 속에서 한 가지 더 중요한 요인이 있었습니다. 단기 은행 대출로 자금을 조달하여 주식 시장에서 위험한 투기를 하던 투기꾼들이었습니다. 이 위험한 투기가 실패하자 은행 대출이 부도 처리되었고, 이것이 은행 공황의 시작이었습니다. 이것이 짧은 버전의 이야기입니다. 여기 여러분이 흥미롭게 여길 만한 더 긴 버전이 있습니다.

하인츠 형제(Heinze process)는 매우 부유한 구리 광산 소유주였습니다. 1907년, 그들은 뉴욕시로 이주하여 찰스 W. 모스(Charles W. Morse)와 합류하여 구리 주식 시장에서 투기했습니다. 그들은 유나이티드 코퍼 컴퍼니(United Copper Company)의 많은 주식을 매입했습니다. 그들은 이 주식들을 담보로 삼아 20개에 달하는 은행들과 거래했습니다. 하인츠 형제와 그들의 동료들은 단기 은행 대출로 자금을 조달하여 주식 시장에서 위험한 투기를 하던 투기꾼들이었습니다.

1907년 10월 9일, 오토 하인츠(Otto Heinze)는 유나이티드 코퍼 주식의 거래 행태를 분석하기로 결정했습니다. 그는 이 주식 약 45만 주가 거래되었음을 발견했습니다. 그러나 회사에 따르면 발행 주식은 42만 5천 주에 불과했습니다. 이는 일부 브로커들이 주식을 빌려 공매도(shorted the stock)했다는 것을 의미했습니다. 즉, 이 다른 투자자들은 주식을 빌려 나중에 더 싼 가격에 다시 사들여 이익을 얻으려는 희망으로 주식을 팔았습니다. 하인츠 형제는 브로커들이 자신들이 소유하고 담보로 제공했던 주식을 빌려주었다고 믿었습니다. 그들은 주식을 회수함으로써 공매도자들이 자신들이 소유한 주식을 매우 높은 가격에 강제로 사들이게 되는 상황을 만들 수 있다고 판단했습니다. 이것을 '숏 스퀴즈(short squeeze)'라고 합니다.

\begin{tcolorbox}[title={\textbf{핵심 개념: 숏 스퀴즈 (Short Squeeze)}}]
\textbf{한 줄 요약:} 공매도한 주식의 가격이 예상과 달리 급등하여, 공매도 투자자들이 손실을 줄이기 위해 주식을 다시 사들이면서 주가가 더욱 폭등하는 현상입니다.

\textbf{직관적 비유:} 친구에게 "나중에 싸게 사서 돌려줄게" 하고 연필을 빌려 팔았는데, 갑자기 연필 가격이 폭등한 상황입니다. 친구가 "지금 당장 연필 돌려줘!"라고 하면, 비싸게라도 연필을 사서 돌려줘야 하므로 손해가 커집니다. 이때 많은 공매도자들이 동시에 연필을 사들이면 가격은 더 오르게 됩니다.

\textbf{기술적 설명:} 하인츠 형제는 자신들이 담보로 제공한 주식을 회수하여 시장에 유통되는 주식 수를 줄임으로써, 공매도자들이 주식을 되사기 어렵게 만들고 가격을 올리려 했습니다.
\end{tcolorbox}

숏 스퀴즈를 유도하기 위해 하인츠 형제는 주식을 담보로 한 대출금을 갚기 위해 150만 달러를 빌려야 했습니다. 유나이티드 코퍼 주식이 45.50달러에서 37.75달러로 다시 하락하자, 하인츠 형제는 대출금 상환을 요청하며 주식을 회수하기 시작했습니다. 몇 시간 만에 주가는 23달러 상승했지만, 이후 약 52달러로 하락했습니다. 그러나 하인츠 형제는 오판했습니다. 주식을 담보로 대출을 연장했던 20개 은행 모두가 모든 주식을 인도했습니다. 결과적으로 하인츠 형제는 필요한 모든 대출금을 상환할 수 없었고, 주식 인수를 거부했습니다. 대출 기관들은 즉시 거부된 담보를 매각했고 주가는 폭락했습니다. 3일 후, 유나이티드 코퍼 주식은 가치의 50\%를 잃었고 결국 주당 10달러에 거래되었습니다. 숏 스퀴즈 계획은 처참하게 실패했습니다. 하인츠 형제는 파산했습니다.

10월 16일, 그들의 증권 중개 회사가 문을 닫았습니다. 하인츠 형제 중 한 명이 소유한 작은 은행도 마찬가지였습니다. 그러나 그 작은 은행은 뉴욕에 기반을 둔 머캔타일 은행(Mercantile Bank)을 결제 서비스에 이용했으며, 약 100만 달러의 예금을 가지고 있었습니다. 머캔타일 은행 또한 가치가 폭락한 유나이티드 코퍼 주식에 직접적으로 노출되어 있었습니다. 이러한 사건들은 너무나 충격적이어서 사람들은 하인츠 형제와 관련된 모든 은행에서 돈을 인출하기 시작했습니다. 공황이 시작된 것입니다.

모든 채권의 질서 있는 청산을 조정하기 위해 뉴욕 청산소(New York Clearing House), 즉 뉴욕 은행들의 컨소시엄이 10월 20일에 개입했습니다. 이것은 은행 공황을 진정시키기 위해 고안된 민간 부문의 개입이었습니다. 그러나 청산소 회의 중에 뉴욕에서 세 번째로 큰 은행인 니커보커 트러스트(Knickerbocker Trust)가 하인츠 형제와의 연관성 때문에 금융 부정에 연루되었다는 사실이 밝혀졌습니다. 이는 니커보커 트러스트가 심각한 문제에 처해 있음을 시사했습니다. 니커보커 트러스트가 재정적으로 건전했는지는 확실하지 않지만, 은행은 신뢰에 의존합니다. 예금자들이 돈을 돌려받을 수 있을지 의문을 품는다면, 그들은 단순히 돈을 인출할 것입니다.

그래서 1907년 10월 22일 아침, 니커보커 트러스트 사무실 문이 열리기 한 시간 전부터 줄이 형성되기 시작했습니다. 줄은 분 단위로 늘어나는 것 같았고, 모든 예금자들이 돈을 돌려달라고 요구했습니다. 줄 형성을 돕기 위해 경찰이 불려왔습니다. 정오까지 니커보커는 약 6,900만 달러의 자산과 6,400만 달러의 부채를 가지고 있으며 모든 사람에게 돈을 지급할 수 있다고 발표했습니다. 아무도 줄에서 이탈하지 않았습니다. 니커보커는 그날 800만 달러의 예금을 잃었습니다. 이것은 전형적인 뱅크런의 예입니다. 니커보커는 재정적으로 건전했지만, 걱정하는 예금자들이 어쨌든 돈을 돌려달라고 요구했습니다.

다음 날, 니커보커는 영업을 중단했습니다. 그러나 뉴욕의 다른 많은 건전한 은행들도 공황에 휩쓸렸습니다. 1907년 10월 24일까지 여러 은행이 파산했습니다. 트러스트 컴퍼니 오브 아메리카(Trust Company of America)에 대한 뱅크런이 시작되자, 뉴욕의 가장 유명한 은행가인 J.P. 모건(J.P. Morgan)과 다른 사람들이 개입했습니다. 그들은 예금자들의 요구를 충족시키기 위해 현금 대출을 제공했습니다. 현금 부족은 뉴욕 증권 거래소를 붕괴 직전으로 몰고 갔습니다. 다시 한번 J.P. 모건이 개입하여 증권 거래소를 유지하는 데 필요한 자본을 조달했습니다.

더 많은 은행들이 붕괴 직전에 처하자, J.P. 모건은 11월 초 뉴욕에 있는 자신의 집으로 약 120명의 뉴욕 은행장들을 소환했습니다. 그는 그들에게 위기 해결을 도울 것을 요구했습니다. 해결책을 강요하기 위해 그는 은행가들을 자신의 서재에 가두었다는 유명한 일화가 있습니다. 결국 은행가들은 동의했고, J.P. 모건은 위기 해결을 위해 2,500만 달러를 추가로 기부했습니다. 이 좋은 소식은 은행 공황을 멈추었지만, 이미 많은 피해가 발생했습니다. 은행 파산을 막기에는 너무 늦었고, 심각한 경기 침체를 막기에도 너무 늦었습니다.

\subsubsection*{1907년 공황의 교훈}
1907년 공황은 미국 역사상 가장 심각한 금융 위기 중 하나로 간주됩니다. 이 사건에서 무엇을 배울 수 있을까요?

\begin{tcolorbox}[title={\textbf{1907년 공황의 주요 교훈}}]
\begin{enumerate}
    \item \textbf{신뢰 부족과 현금 부족:} 예금자들은 저축을 잃을까 봐 걱정했기 때문에 공황 상태에 빠졌습니다. 신뢰와 현금의 부족은 예금자들의 공황을 초래할 수 있습니다. 이것이 예금 보험이 뱅크런을 방지하는 이유를 설명합니다. 보험에 가입된 예금자들은 잠재적인 손실로부터 보호받을 것입니다. 이것이 오늘날 예금자들이 줄을 서는 고전적인 뱅크런이 드문 이유일 것입니다.
    \item \textbf{현금 대출의 중요성:} 두 번째 교훈은 현금 대출이 진행 중인 뱅크런을 예금자들의 요구를 충족시킴으로써 막는 데 도움이 될 수 있다는 것입니다. 1907년 공황에서 J.P. 모건은 '최종 대부자(Lender of Last Resort)' 역할을 했습니다.
    \item \textbf{정부 지원의 한계:} 1907년에는 J.P. 모건이 현금이 바닥나거나 자신의 은행에 문제가 생기지 않았다는 점에서 운이 좋았습니다. 정부의 지원 없이는 민간 최종 대부자의 선택지는 제한적입니다. 이것이 무제한적인 자원을 가진 연방준비제도가 최종 대부자로서 더 효과적일 수 있는 이유를 설명합니다.
\end{enumerate}
\end{tcolorbox}


\subsection{Lesson 4-1.2: 대공황}
대공황(The Great Depression)에 대한 수업에 오신 것을 환영합니다. 현대 시대에 대공황은 미국과 서구 세계 대부분에서 가장 심각한 금융 위기이자 경제 침체였습니다. 우리는 대공황 동안 금융 상황이 어떻게 전개되었는지 살펴볼 것입니다. 우리는 엄청난 수의 은행 파산을 보게 될 것입니다. 우리는 연방준비제도 시스템이 은행 공황을 막는 데 왜 비효율적이었는지 이해할 것입니다. 마지막으로, 경제를 재활성화하고 은행 부문에 대한 신뢰를 회복시킨 대규모 정부 대응을 살펴볼 것입니다. 여기에는 정부 지출 프로그램이 포함되었습니다. 또한 예금 보험과 같은 새로운 기관의 설립도 포함되었습니다. 이러한 역사적 사건을 이해하는 것은 즐겁고 매우 유익할 수 있습니다. 정확한 역사적 세부 사항에 대해서는 너무 걱정하지 마십시오. 많은 세부 사항을 언급하겠지만, 이는 주로 재미를 위한 것이며, 서술과 우리가 배울 수 있는 교훈에 집중하십시오. 대공황은 예금 보험이나 실업 수당과 같이 금융 위기의 영향을 완화할 대부분의 정책이 미국에 존재하지 않았다는 점에서 흥미롭습니다. 이는 정부가 금융 위기 동안 개입하지 않았을 때 경제에서 어떤 일이 일어날지 이해하는 데 도움이 됩니다.

\subsubsection*{대공황의 시작}
대공황의 시작을 간략하게 요약해 보겠습니다. 1929년 10월 24일, 이른바 '검은 목요일(Black Thursday)'에 개장과 동시에 시장은 11\% 하락했습니다. 이 폭락은 대규모 투자자들에 대한 사기 혐의로 인한 런던 증권 거래소의 하락에 뒤이어 발생했습니다. 1929년 10월 28일, 주가 하락으로 인해 점점 더 많은 투자자들이 신용으로 매수한 주식에 대한 마진콜(margin calls)에 직면했습니다. 결과적으로 더 많은 투자자들이 주식을 팔아 손실을 줄였습니다. 다우 지수는 그날 거의 13\% 하락했습니다. 다음 날, 투자자들은 기록적인 수의 주식을 거래했습니다. 하루 만에 1,600만 주가 거래되었습니다. 다음 날, 주식 시장은 유동성 부족의 징후를 보였습니다. 일부 주식은 어떤 가격에도 팔리지 않았습니다. 다우 지수는 또 다른 12\% 하락했습니다.

\begin{tcolorbox}[title={\textbf{주식 시장 붕괴 (1929년 9월 ~ 1932년 7월)}}]
\textbf{핵심 내용:}
\begin{itemize}
    \item 1929년 10월의 사건들은 주가 장기 하락의 서막에 불과했습니다.
    \item 1929년 9월부터 1932년 7월까지 주가는 80\% 이상 폭락했습니다.
\end{itemize}
\textbf{시사점:} 주식 시장의 붕괴는 단순한 일시적 현상이 아니라, 장기적인 경제 침체의 시작을 알리는 신호였습니다.
\end{tcolorbox}

위기가 시작될 당시, 8,000개 이상의 상업 은행이 연방준비제도 시스템에 속해 있었지만, 거의 16,000개는 그렇지 않았습니다. 당시 은행들은 종종 주식을 직접 소유하고 있었고 현금 준비금이 제한적이어서, 연방준비제도 시스템에 접근하여 현금을 확보할 수 없다는 것은 문제가 되었습니다. 연방준비제도에 접근할 수 없는 은행들은 종종 다른 대형 은행에 현금을 예치했는데, 이는 소규모 은행이 대형 은행에서 예금을 인출하거나 대형 은행 자체가 어려움에 처할 경우 전염(contagion) 가능성을 만들었습니다. 이러한 환경에서 예금자들은 당연히 은행의 건전성에 대해 의문을 품었습니다. 대공황의 어쩌면 놀라운 특징은 하나의 전국적인 은행 공황이 아니라 수많은 지역 은행 뱅크런이 있었다는 것입니다. 총 1929년에서 1933년 사이에 약 9,000개의 은행이 문을 닫았습니다.

\begin{tcolorbox}[title={\textbf{은행 파산의 파급 효과}}]
\textbf{핵심 내용:}
\begin{itemize}
    \item 대규모 은행 폐쇄는 예금을 급격히 감소시켰습니다.
    \item 은행들은 현금 준비금을 확보하기 위해 대출을 줄였습니다.
    \item 결과적으로 통화 공급이 위축되었습니다.
\end{itemize}
\textbf{시사점:} 은행 파산은 단순한 개별 은행의 문제가 아니라, 전체 금융 시스템의 신용 경색과 통화 공급 감소로 이어져 경제 전반에 악영향을 미쳤습니다.
\end{tcolorbox}

여러분은 연방준비제도가 왜 통화 공급을 늘리지 않았는지 궁금할 수 있습니다. 당시 미국은 다른 대부분의 국가와 마찬가지로 금본위제(gold standard)를 채택하고 있었습니다. 이는 화폐가 고정된 환율로 금에 연동되어 있었고, 이는 중앙은행의 선택지를 심각하게 제한했습니다.

\subsubsection*{디플레이션의 심화}
낮은 통화 공급은 물가에 하방 압력을 가했습니다. 이 차트에서 볼 수 있듯이, 대공황 동안 소비자 물가는 약 25\% 디플레이션되었습니다. 식품의 경우 하락폭이 훨씬 더 컸습니다. 이러한 디플레이션은 농부들과 기업들이 실질적으로, 즉 생산된 상품의 측면에서 더 높은 부채에 직면했음을 의미했습니다. 결과적으로 많은 농부들과 기업들이 파산했습니다. 이러한 파산은 은행 시스템에 추가적인 스트레스를 유발했습니다.

\begin{tcolorbox}[title={\textbf{핵심 개념: 디플레이션 (Deflation)}}]
\textbf{한 줄 요약:} 물가가 지속적으로 하락하는 현상입니다.

\textbf{직관적 비유:} 오늘 100원짜리 빵이 내일 90원이 되고 모레 80원이 된다고 상상해 보세요. 사람들은 "내일 사면 더 싸지겠네?" 하고 소비를 미루게 됩니다. 기업은 물건이 안 팔리니 생산을 줄이고, 직원도 해고하게 됩니다.

\textbf{기술적 설명:} 물가 하락, 특히 미래의 물가 하락 예상은 투자를 할 유인을 감소시킵니다. 만약 현금을 보유하면 미래에 더 많은 상품을 살 수 있다면, 왜 미래에 낮은 가격만 청구할 수 있는 위험한 사업에 현금을 투자하겠습니까? 이는 경제 활동을 더욱 위축시킵니다.
\end{tcolorbox}

디플레이션이 시작되고 신용 가용성이 줄어들면서 기업 투자는 급락했습니다. 투자 감소는 대규모 파산과 함께 산업 생산의 장기적인 감소로 이어졌습니다. 이 차트에서 볼 수 있듯이, 대공황 동안 산업 생산량은 50\% 감소했습니다. 이는 전례 없는 경제 활동의 감소였습니다.

\begin{tcolorbox}[title={\textbf{경제 활동 위축 (산업 생산 및 실업률)}}]
\textbf{핵심 내용:}
\begin{itemize}
    \item 산업 생산량은 대공황 동안 50\% 급감했습니다.
    \item 실업률은 대공황 초기의 거의 0\%에서 끝 무렵에는 25\% 이상으로 증가했습니다.
\end{itemize}
\textbf{시사점:} 경제가 위축되면서 일자리가 급격히 사라졌고, 이는 빈곤과 사회 불안으로 이어졌습니다. 실업 보험이 없었기 때문에 실업자들은 더욱 큰 고통을 겪었습니다.
\end{tcolorbox}

경제 위축과 함께 실업률이 빠르게 상승했습니다. 이 차트에서 볼 수 있듯이, 실업률은 대공황 초기의 거의 0\%에서 대공황 말기에는 25\% 이상으로 증가했습니다. 실업 보험이 존재하지 않았기 때문에 빈곤과 사회 불안이 뒤따랐습니다.

\subsubsection*{사회 불안과 정부 대응}
사회 불안의 가장 두드러진 예 중 하나는 1차 세계대전 참전 용사들의 행진이었습니다. 1932년, 1차 세계대전 참전 용사들은 '보너스 원정군(Bonus Expeditionary Forces)'이라는 단체를 조직하여 워싱턴 D.C.로 행진했습니다. 약 2만 명의 참전 용사들이 수도에 모여 워싱턴 D.C.에 야영하며 약속되었지만 의회에서 지연되고 있던 보너스 지급을 요구했습니다. 시위는 결국 강제로 해산되었지만, 굶주린 참전 용사들을 해산시키는 탱크의 사진은 후버 대통령의 재선 도전을 약화시켰습니다. 11월, 프랭클린 델라노 루스벨트(Franklin Delano Roosevelt)는 경제 포퓰리즘을 내세워 선거에서 승리했고, 참전 용사 보너스는 결국 1936년에 지급되었습니다.

루스벨트가 취임한 후 첫 번째 조치 중 하나는 1933년 3월 6일에 '뱅크 홀리데이(bank holiday)', 더 정확히는 '뱅크 홀리데이 주간'을 선언한 것이었습니다. 일주일 내내 아무도 은행이나 은행 서비스에 접근할 수 없었습니다. 이는 돈을 인출하거나 이체할 수 없다는 것을 의미했으며, 그해 초 다시 시작되었던 뱅크런을 중단시켰습니다. 뱅크 홀리데이의 목적은 은행 시스템에 대한 신뢰를 회복하고 미래의 뱅크런을 방지할 법안을 제정할 시간을 버는 것이었습니다.

1933년 긴급 은행법(Emergency Banking Act of 1933)은 이 목표를 달성했습니다. 요컨대, 이 법은 건전한 은행만이 재개하고 부실한 은행은 폐쇄되도록 보장했습니다. 또한 이 법은 연방준비제도에 더 많은 권한을 부여하고 은행에 대출하기 쉽게 만들었습니다. 연방준비제도는 사실상 재개된 은행의 예금 보증인이 되었습니다.

\begin{tcolorbox}[title={\textbf{긴급 은행법의 즉각적인 효과}}]
\textbf{핵심 내용:}
\begin{itemize}
    \item 법안 통과 후 다우존스 산업 주가 지수는 하루 만에 15\% 이상 상승했습니다.
    \item 뱅크런이 중단되었습니다.
\end{itemize}
\textbf{시사점:} 정부의 신속하고 단호한 조치가 시장의 신뢰를 회복하고 금융 공황을 진정시키는 데 결정적인 역할을 했음을 보여줍니다.
\end{tcolorbox}

1933년 은행법(Banking Act of 1933), 일명 글래스-스티걸 법(Glass-Steagall Act)은 1933년 6월에 예금자 보호를 강화했습니다. 상업 은행(예금 수취)과 투자 은행을 분리함으로써, 예금자들은 더 이상 주식이나 일반적으로 증권 시장에서의 은행의 위험 감수에 노출되지 않게 되었습니다. 이 법은 또한 연방예금보험공사(Federal Deposit Insurance Corporation, FDIC)를 설립하여 예금에 대한 보험을 제공함으로써 추가적인 뱅크런을 방지하는 데 도움이 되었습니다. 은행법은 또한 은행 감독을 강화하고 연방준비제도 시스템의 일부 특징을 변경했습니다. 은행 공황은 끝났지만, 실물 경제는 여전히 심각한 상황에 처해 있었습니다.

경제 위기에서 정부의 역할을 근본적으로 재편한 '뉴딜(New Deal)' 정책의 일환으로, 루스벨트의 첫 번째 목표는 실업자들의 고통을 완화하는 것이었습니다. 정부는 공공사업진흥국(Works Progress Administration, WPA)과 민간산림보존단(Civilian Conservation Corps, CCC)을 통해 일자리를 창출했습니다. 공공사업진흥국은 주로 건설 분야에서 약 850만 명에게 일자리를 제공했습니다. 65만 마일 이상의 도로, 12만 5천 개의 공공 건물, 7만 5천 개의 다리, 8천 개의 공원이 건설되었습니다. 민간산림보존단은 나무 심기, 홍수 방지벽 건설, 산불 진압, 산림 도로 유지 보수 등의 작업을 제공했습니다. 이러한 프로젝트들은 단기간에 시작될 수 있었고 즉각적인 효과를 가져왔다는 점이 중요합니다.

\begin{tcolorbox}[title={\textbf{뉴딜 정책의 고용 효과}}]
\textbf{핵심 내용:}
\begin{itemize}
    \item 뉴딜 정책 시행 후 1년 이내에 실업률이 9\% 포인트 감소했습니다.
\end{itemize}
\textbf{시사점:} 정부의 대규모 공공 사업과 일자리 창출 프로그램이 실업 문제 해결에 직접적이고 효과적인 영향을 미쳤음을 보여줍니다.
\end{itemize}
\end{tcolorbox}

사람들의 경제적 어려움은 또한 사회 안전망의 필요성을 보여주었습니다. 사회보장국(Social Security Administration)을 통한 연금 시스템 구축, 실업 수당, 장애 수당은 '뉴딜'의 일부였습니다. '뉴딜'은 또한 전국노동관계위원회(National Labor Relations Board)를 설립하여 노동 관계를 재편했습니다. 경제는 다시 정상 궤도에 오르는 것처럼 보였습니다. 그러나 '뉴딜'은 자금 조달이 필요했고, 세금이 인상되었습니다. 더욱이 경제가 개선되면서 재정 부양책이 철회되었습니다. 동시에 연방준비제도는 은행의 지급준비율(reserve requirements)을 인상하여 신용 가용성을 감소시켰습니다.

\begin{tcolorbox}[title={\textbf{1937년 경기 침체: 정책 오류의 경고}}]
\textbf{핵심 내용:}
\begin{itemize}
    \item 1937년, 긴축 통화 정책(연준의 지급준비율 인상)과 긴축 재정 정책(재정 부양책 철회)이 동시에 시행되었습니다.
    \item 결과적으로 산업 생산이 다시 붕괴되었습니다.
\item 경기 침체는 연방준비제도가 지급준비율을 인하하면서 끝났습니다.
\end{itemize}
\textbf{시사점:} 1937년 사건은 경제 지원을 너무 일찍 철회하는 것에 대한 경고를 제공합니다. 위기에서 벗어나기 위한 정책적 노력이 충분히 지속되지 않으면, 경제는 다시 침체에 빠질 수 있습니다.
\end{tcolorbox}

이 차트에서 볼 수 있듯이, 1937년 경기 침체 이후 실업률은 빠르게 감소했습니다. 그러나 또 다른 부양책이 경제를 빠르게 성장시키는 데 도움이 되었다는 점에 유의하는 것이 중요합니다. 유럽의 국방력 증강과 그에 따른 정부 지출 증가였습니다. 1940년 중반부터 실업률 감소가 가속화되는 것을 볼 수 있습니다.

\subsubsection*{대공황에서 얻은 교훈}
이 강의에서 무엇을 배웠습니까?

\begin{tcolorbox}[title={\textbf{대공황의 주요 교훈}}]
\begin{enumerate}
    \item \textbf{금융 위기의 파괴적 영향:} 대공황을 금융 위기의 결과에 대한 사례 연구로 사용하면서, 우리는 그 영향이 특히 사회 안전망이나 정부 개입이 없을 때 얼마나 파괴적일 수 있는지 보았습니다.
    \item \textbf{금융 시스템의 대규모 변화 필요성:} 금융 시스템에 대한 신뢰를 회복하려면 대규모 변화가 필요합니다.
    \item \textbf{정부 지원의 조기 철회 위험:} 1937년의 경기 침체는 경제에 대한 정부 지원을 너무 적게 제공하거나 너무 일찍 철회하는 것이 또 다른 경제 침체로 이어질 수 있음을 보여줍니다.
\end{enumerate}
\end{tcolorbox}


\subsection{Lesson 4-1.3: 금융 위기의 공통 특징}
금융 위기의 공통 특징에 대한 수업에 오신 것을 환영합니다. 우리는 지난 150년 동안의 주요 글로벌 금융 위기에 대한 광범위한 개요를 제공할 것입니다. 우리의 목표는 위기가 어떻게 전개되었는지, 왜 위기가 전 세계적이었는지, 그리고 이러한 위기들 사이에 어떤 유사점이 있는지 이해하는 것입니다. 우리는 금융 위기의 이러한 공통된 주제들을 밝혀낼 것입니다. 과도한 부채가 항상 역할을 한다는 것을 알게 될 것입니다. 은행 예금과 같은 단기 부채에 대한 런이 항상 있습니다. 지속 불가능한 장기 투자나 정부 정책도 종종 역할을 합니다. 신용의 급격한 확장이 항상 금융 위기로 끝나는가요? 강의가 끝날 무렵에 이 질문에도 답할 것입니다. 시작해 봅시다.

\subsubsection*{주요 글로벌 금융 위기 사례}
지난 150년 동안 가장 심각했던 5가지 글로벌 금융 위기는 1873년 공황(Panic of 1873), 1890년 베어링스 위기(Barings Crisis of 1890), 1907년 공황(Panic of 1907), 1930년대 대공황(Great Depression of the 1930s), 그리고 2008년 금융 위기(2008 Financial Crisis)였습니다. 각각을 차례로 살펴보겠습니다. 이것은 간략한 설명이며, 우리의 목표는 계속해서 나타나는 특징들을 찾는 것임을 기억하십시오. 모든 작은 세부 사항보다는 큰 그림에 집중하십시오.

\begin{tcolorbox}[title={\textbf{금융 위기의 공통 특징을 찾는 여정}}]
우리는 다음의 주요 글로벌 금융 위기들을 통해 공통된 패턴을 찾아낼 것입니다.
\begin{itemize}
    \item \textbf{1873년 공황:} 철도 투자 버블과 국제적 확산
    \item \textbf{1890년 베어링스 위기:} 아르헨티나 정부 부채와 국제적 자금 조달 문제
    \item \textbf{1907년 공황:} 구리 투기와 은행 시스템의 취약성
    \item \textbf{1930년대 대공황:} 금본위제와 고정 환율을 통한 전 세계적 전파
    \item \textbf{2008년 금융 위기:} 주택 시장 붕괴와 글로벌 금융 시스템의 상호 연결성
\end{itemize}
\end{tcolorbox}

\paragraph{1873년 공황}
1873년 공황은 최초의 글로벌 은행 위기였습니다. 철도는 1800년대 중반의 뜨거운 투자 대상이었습니다. 1869년 미국 동부와 서부 해안을 연결하는 대륙횡단철도(Transcontinental railroad) 완성과 같은 이정표는 철도 투자에 대한 관심을 부추겼습니다. 이는 미국뿐만 아니라 유럽에서도 마찬가지였습니다. 예를 들어, 독일 전역에서 부채로 자금을 조달한 철도 회사들이 생겨났습니다. 철도 주가는 전 세계적으로 급등했습니다. 많은 은행들이 철도 채권을 투자로 보유했습니다. 한 국가의 은행 시스템의 주요 취약점은 철도 채권 투기였습니다. 실제로 철도는 많은 선행 투자가 필요했고 초기 몇 년 동안은 수익이 거의 또는 전혀 발생하지 않아 투자자들에게 위험을 안겨주었습니다.

철도 투기 거품 붕괴로 이어진 여러 요인들이 있었습니다. 예를 들어, 위기 직전 보스턴과 시카고의 화재는 은행 준비금을 감소시켰습니다. 독일과 미국은 은본위제(silver currencies)를 포기했는데, 이는 광부들에게 부담을 주고 철도 네트워크의 필요성을 감소시켰습니다. 그리고 프랑스-독일 전쟁으로 인한 혼란이 여전히 존재했습니다. 위기의 첫 징후는 5월 9일 비엔나에서 주식 시장이 폭락하고 일련의 파산이 뒤따르면서 나타났습니다. 독일에서는 대규모 철도 회사가 붕괴되어 공황을 가중시켰습니다. 미국에서는 주요 은행인 J. 코흐 앤 컴퍼니(J Koch and Company)가 철도 채권을 매각할 수 없게 되어 9월 중순에 붕괴되었고, 이는 미국 전역에서 뱅크런과 파산을 촉발했습니다.

미국 은행 위기의 저명한 희생자 중 하나는 프리드먼 저축 은행(Friedman Savings Bank)이었습니다. 이 은행은 당시 자유를 얻은 전 노예들을 위한 최초이자 유일한 은행이었습니다. 미국과 유럽은 실업률이 빠르게 상승하는 깊은 경기 침체에 빠졌습니다. 이 위기는 오스만 제국과 라틴 아메리카에도 영향을 미쳤습니다. 종합적으로 볼 때, 이 위기는 투기적 투자, 즉 철도 채권에서 시작되어 대규모 손실과 은행 시스템에 대한 신뢰 부족으로 이어졌습니다.

\paragraph{1890년 베어링스 위기}
정부 부채라는 다른 유형의 투자가 이른바 1890년 베어링스 위기를 촉발했습니다. 1890년까지 아르헨티나는 10년간 대규모 자본 유입을 경험했습니다. 자본은 주로 런던에서 베어링스 은행(Barings bank)을 포함한 유럽 은행들에 의해 조달되었습니다. 이 돈은 철도 및 운송 네트워크 건설, 농지 개선을 포함한 장기 인프라 투자 프로젝트에 자금을 조달하는 데 사용되었습니다. 그러나 1880년대 후반까지 해외 차입금의 40\%가 부채 상환에, 수입의 60\%가 소비재에 사용되었습니다. 아르헨티나 정부 또한 상당한 재정 적자를 겪고 있었습니다. 철도 순이익은 1888년에 정점을 찍었고, 결국 국제 투자자들은 아르헨티나 정부의 부채 부담에 대해 경계심을 갖게 되었습니다. 결국 베어링스는 아르헨티나의 채권 발행을 위한 투자자를 찾을 수 없었습니다. 아르헨티나는 뱅크런을 경험했고, 새로운 자금을 조달할 수 없게 되자 정부 채무 불이행(default)을 선언했습니다. 이 채무 불이행은 라틴 아메리카에 심각한 경기 침체를 초래했을 뿐만 아니라 유럽과 미국의 주식 시장에도 급격한 하락을 가져왔습니다. 1873년 위기와 마찬가지로, 투기적인 장기 투자에 대한 신용 붐은 결국 나쁜 결과를 초래했습니다. 그러나 국제적 연계를 주목하십시오. 투기 활동은 국제 투자자들에 의해 부추겨졌습니다.

\paragraph{1907년 공황}
다음 위기인 1907년 공황은 구리 투기에서 시작되었습니다. 1907년 공황은 별도의 강의에서 논의했으므로 간략하게 설명하겠습니다. 샌프란시스코 지진은 보험 회사, 철도, 은행에 큰 손실을 입혔습니다. 이러한 배경 속에서 구리 주식 투기가 잘못되어 은행 시스템에 더 많은 손실을 초래했습니다. 이는 은행 공황, 은행 폐쇄, 그리고 미국에서 유럽으로 확산된 경기 침체로 이어졌습니다. 미국에서는 실업률이 16\% 이상으로 상승했습니다. 글로벌 시장도 큰 타격을 입었습니다. 예를 들어, 영국의 실업률은 1906년과 1908년 사이에 3.6\%에서 7.8\%로 두 배가 되었습니다. 독일, 일본, 이탈리아도 뱅크런을 경험했습니다. 1907년 공황은 금융 시장 투기가 잘못된 예입니다. 그러나 그 영향은 다른 위기와 유사합니다. 은행 손실은 뱅크런을 유발했습니다. 그리고 은행 시스템의 혼란은 신용과 실물 경제를 위축시켰습니다. 여기에도 국제적 연계가 있었습니다. 여기에는 샌프란시스코의 재해 위험을 인수하는 해외 보험 회사와 외국인 투자자들의 위험한 주식 시장 투자가 포함되었습니다. 이는 위기가 미국에서 전 세계로 확산되도록 했습니다.

\paragraph{대공황}
대공황은 국제적 연계의 중요성을 가장 명확하게 보여주는 예입니다. 대공황은 별도로 논의했으므로 국제적 요소에만 집중하겠습니다. 1929년 미국에서 주식 시장 붕괴가 대공황을 촉발했을 때, 대부분의 국가들은 금본위제(gold standard)를 채택하고 있었습니다. 금본위제는 모든 지폐가 금으로 전환될 수 있었고, 따라서 금 준비금이 통화 공급을 결정한다는 것을 의미했습니다. 이는 또한 각 통화가 금에 고정되어 있었기 때문에 대부분의 국가가 고정 환율 시스템으로 운영되었다는 것을 의미했습니다.

당시 미국은 글로벌 시장에서 금의 순 대출국이었습니다. 그러나 은행 위기가 전개되자 미국은 금을 대출하기보다는 비축했습니다. 이는 다른 국가들이 각자의 통화 공급 감소를 피하기 위해 금 준비금을 유지하기 위해 고군분투하게 만들었습니다. 낮은 통화 공급은 은행에 부담을 주고 전 세계적으로 디플레이션을 초래했습니다. 예를 들어, 독일 라이히스방크(Reichsbank)는 단기 채권자들이 독일에 대한 대출을 중단했을 때 사실상 금이 바닥났습니다. 결국, 국가들이 자체 통화 공급을 통제하고 뱅크런과 디플레이션을 피할 수 있도록 금본위제는 폐지되었습니다. 요컨대, 대공황은 고정 환율 때문에 글로벌 금융 시스템을 통해 전파되었습니다. 따라서 이전 금융 위기와 달리, 여기서는 나쁜 장기 투자에 대한 공통된 노출이 문제가 아니었습니다. 이번에는 전 세계의 모든 통화와 모든 통화 공급이 연결되어 있었습니다. 변동 환율(floating exchange rates)에서는 국제적 충격에 직면했을 때 통화 가치가 조정됩니다. 고정 환율에서는 통화 공급이 조정되어야 하며, 이는 큰 부정적인 결과를 초래합니다.

\paragraph{2008년 금융 위기}
이제 2008년 금융 위기로 넘어갑니다. 이 사건이 과거 위기들과 얼마나 유사했는지 이제 알 수 있을 것입니다. 위기는 미국 주택 시장에서 시작되었습니다. 스페인과 아일랜드와 같은 다른 많은 국가들도 주택 붐과 붕괴를 경험했습니다. 주택 붐에 자금을 조달하는 많은 장기 모기지 투자는 단기 신용으로 조달되었습니다. 이 단기 신용이 고갈되자 모기지 시장과 주택 가격이 붕괴되었습니다. 이전 금융 위기와 마찬가지로, 국제 투자자들은 이러한 투자에 노출되어 있었습니다. 이러한 노출의 정도는 2007년 6월 중견 독일 은행인 ECA Bay Deutsche Industriebank가 미국 서브프라임 모기지 손실로 파산했을 때 명확해졌습니다. 또한 국내 플레이어들, 가장 악명 높은 투자 은행인 리먼 브라더스(Lehman Brothers)가 모기지로 담보된 단기 자금 조달에 의존하고 있다는 것이 분명해졌습니다.

리먼 브라더스는 2008년 가을에 붕괴되었습니다. 가장 큰 보험 회사인 아메리칸 인터내셔널 그룹(American International Group)도 마찬가지였습니다. 이는 전 세계적으로 주식 시장 폭락과 자금 시장 공황을 촉발했습니다.

2008년 위기는 자본이 국경을 자유롭게 넘나드는 글로벌 통합 금융 시스템 때문에 전 세계적으로 전개되었습니다. 글로벌 통합 금융 시스템은 충격이 작을 때 쉽게 흡수될 수 있도록 돕습니다. 그러나 가장 큰 경제에 대한 충격은 모든 사람에게 영향을 미칩니다.

\subsubsection*{신용 붐은 항상 금융 위기로 이어지는가?}
마무리하기 전에 마지막 질문을 살펴보겠습니다. 신용 붐이 항상 금융 위기로 이어지는가요? 짧은 답변은 '아니오'입니다. 최근 연구에 따르면 많은 신용 붐이 금융 위기로 이어지지 않고, 대신 지속적인 경제 성장을 가져올 수 있습니다. 이 차트는 1910년부터 2014년까지 189개국 표본에 대한 신용 붐 정점 전후의 생산량 역학을 보여줍니다. 신용 붐의 정점은 이벤트 시간 0, 즉 수직선으로 표시됩니다. 두 가지 다른 유형의 신용 붐을 나타내는 두 개의 선이 있습니다. 대략적으로 말하면, 하나는 건설 및 부동산 부문이고 다른 하나는 제조업 부문입니다.

\begin{tcolorbox}[title={\textbf{신용 붐과 경제 성장: 부동산 vs. 제조업}}]
\textbf{차트 분석 (텍스트 기반 설명):}
\begin{itemize}
    \item \textbf{신용 붐 기간:} 평균적으로 GDP는 붐 기간 동안 강하게 성장합니다.
    \item \textbf{부동산 신용 붐 이후:} 정점 이후에만 생산량이 감소하는 경향을 보입니다.
    \item \textbf{제조업 신용 붐 이후:} 생산량이 지속적으로 성장하거나 안정적인 경향을 보입니다.
\end{itemize}
\textbf{왜 이런 현상이 발생할까요?}
\begin{itemize}
    \item \textbf{부동산 신용 붐:} 일반적으로 주택 가격 상승, 신규 주택 건설, 고용 증가 등을 유발합니다. 그러나 결국 붐을 지속할 만큼 충분한 지역 구매자가 없게 됩니다. 이는 주택 가격에 하방 압력을 가하고, 건설 부문의 붕괴로 이어질 수 있습니다.
    \item \textbf{제조업 신용 붐:} 글로벌 수요나 신제품 수요에 의해 주도될 때 더 지속 가능할 수 있습니다. 예를 들어, 1960년대부터 1980년대까지 한국 제조업의 빠른 성장을 생각해 보십시오. 대부분의 제품은 수출되었습니다. 제조업 부문에 대한 신용은 수익성 있는 투자에 자금을 조달할 수 있습니다. 이러한 더 높은 이익은 신용 붐 동안 발생한 부채를 상환하는 데 도움이 됩니다.
\end{itemize}
\textbf{결론:} 모든 신용 붐이 금융 위기로 이어지는 것은 아니며, 어떤 부문에서 신용 붐이 발생하는지가 중요합니다.
\end{tcolorbox}

\subsubsection*{금융 위기의 공통 특징 요약}
이 세션에서 우리는 금융 위기의 공통 특징을 설명했습니다.

\begin{tcolorbox}[title={\textbf{금융 위기의 공통 특징}}]
\begin{itemize}
    \item \textbf{과도한 부채:} 특히 단기 부채가 많습니다.
    \item \textbf{실패한 장기 투자:} 결국 나쁜 결과를 초래하는 장기 투자가 있습니다.
    \textbf{지속 불가능한 정부 정책:} 과도한 지출이나 고정 환율과 같은 지속 불가능한 정부 정책이 위기를 심화시킵니다.
    \item \textbf{국제적 전파:} 일부 금융 위기는 국제적 연계를 통해 전 세계적으로 확산됩니다.
\end{itemize}
\textbf{추가 교훈:} 모든 신용 붐이 은행 위기나 GDP 하락으로 이어지는 것은 아닙니다. 역사적으로 제조업 부문의 신용 붐은 더 지속 가능했습니다.
\end{itemize}
\end{tcolorbox}


\subsection{Lesson 4-1.4: 2008년 금융 위기 중 자동차 대출}
2008년 금융 위기 중 자동차 대출에 대한 강의에 오신 것을 환영합니다. 이 사례 연구에서는 단기 자금 조달이 고갈되었을 때 비은행 대출 기관이 큰 역할을 하는 자동차 대출 시장에서 어떤 일이 일어났는지 살펴볼 것입니다. 은행뿐만 아니라 다른 금융 기관도 런에 취약하며, 신용 공급의 감소가 자동차 판매 감소로 이어진다는 것을 알게 될 것입니다.

\subsubsection*{캡티브 금융 회사와 자금 조달}
2008년 금융 위기 이전에 전체 자동차 대출의 약 절반은 '캡티브 금융 회사(captive finance companies)'에 의해 시작되었습니다. 캡티브 금융 회사는 자동차 제조업체가 소유한 대출 기관으로, 예를 들어 포드 모터 크레딧(Ford Motor Credit), GM 파이낸셜(GM Financial), 토요타 파이낸셜 서비스(Toyota Financial Services) 등이 있습니다.

\begin{tcolorbox}[title={\textbf{캡티브 금융 회사의 시장 점유율}}]
\textbf{핵심 내용:}
\begin{itemize}
    \item 캡티브 금융 회사의 시장 점유율은 자동차 제조업체가 위치한 미시간주, 그리고 역사적으로 북부보다 은행 접근성이 낮았던 남부와 서부에서 특히 높았습니다.
\end{itemize}
\textbf{시사점:} 캡티브 금융 회사는 은행 신용 접근성이 낮은 지역에서 특히 강력한 역할을 수행했습니다.
\end{tcolorbox}

은행과 달리 캡티브 금융 회사는 예금을 받지 않습니다. 대신 캡티브 금융 회사는 '자산유동화기업어음(Asset-Backed Commercial Paper, ABCP)'과 '자산유동화증권(Asset-Backed Securities, ABS)'을 발행하여 자금을 조달했습니다.

\paragraph{자산유동화기업어음 (ABCP)의 작동 방식}
이러한 맥락에서 자산유동화기업어음이 어떻게 작동했는지 먼저 논의해 봅시다. 캡티브 금융 회사는 신탁(trust)을 설립하고 자동차 대출을 그 신탁으로 이전합니다. 그러면 신탁은 자동차 대출과 이 대출에서 발생하는 현금 흐름을 담보로 하는 기업어음(commercial paper)을 발행합니다. 신탁은 자산유동화기업어음 발행 수익금을 사용하여 캡티브 금융 회사에 대출금을 지급합니다.

\begin{tcolorbox}[title={\textbf{ABCP의 회전식(Revolving) 구조}}]
\textbf{핵심 내용:}
\begin{itemize}
    \item 일반적으로 이러한 신탁은 '회전식(revolving)'입니다. 즉, 오래된 자동차 대출이 상환되면 캡티브 금융 회사는 새로운 자동차 대출을 신탁에 넣습니다.
    \item 이를 통해 캡티브 금융 회사는 지속적으로 단기 자금 조달에 접근할 수 있으며, 이는 다시 새로운 대출을 발행할 수 있게 합니다.
    \item 그러나 기업어음은 만기가 최대 270일이며 지속적으로 롤오버(roll over)되어야 합니다.
\end{itemize}
\textbf{시사점:} 캡티브 금융 회사는 은행과 마찬가지로 단기적으로 차입하고 장기적으로 대출합니다. 즉, '만기 변환(maturity transformation)'에 참여합니다. 이는 뱅크런과 유사한 위험에 노출될 수 있음을 의미합니다.
\end{tcolorbox}

\paragraph{자산유동화증권 (ABS)의 작동 방식}
자산유동화증권은 약간 다릅니다. 여기서는 캡티브 금융 회사가 신탁을 설정하고 자동차 대출을 그 신탁으로 이전합니다. 신탁은 더 장기적인 채권(notes)을 발행하여 자금을 조달합니다.

\begin{tcolorbox}[title={\textbf{ABS의 특징 (예: Ford Motor Credit Owner Trust 2018-B)}}]
\textbf{핵심 내용:}
\begin{itemize}
    \item 포드 모터 크레딧 오너 트러스트 2018-B의 예시에서, 거의 10억 달러에 달하는 채권을 발행합니다.
    \item 이 채권은 포드 모터 크레딧이 구매한 자동차, 경트럭, 유틸리티 차량 대출 풀을 담보로 합니다.
    \item 최대 만기는 7년으로, 훨씬 더 긴 자금 조달 기간입니다. 이 거래에서는 만기 변환이 많지 않습니다.
    \item 만기가 길수록 채권의 이자율이 높아집니다.
    \item ABCP와 달리, 이것은 일회성 거래이며 대출 풀이 보충되지 않고 대출 상환금이 채권 상환에 사용됩니다.
\end{itemize}
\textbf{시사점:} 2008년 금융 위기 이전에는 ABCP가 ABS보다 저렴하고 회전식이었기 때문에 캡티브 금융 회사의 주요 자금 조달원이었습니다.
\end{tcolorbox}

\subsubsection*{2008년 금융 위기 중 ABCP 런}
그러나 2006년과 2007년에 서브프라임 위기가 모기지 시장을 강타하기 시작하자, 투자자들은 자산유동화기업어음을 담보하는 대출 풀의 품질에 대해 의심을 품기 시작했습니다. 제기된 질문 중 하나는 해당 대출 풀에 얼마나 많은 서브프라임 자동차 대출이 포함되어 있는가였습니다. 투자자들이 제너럴 모터스 어셉턴스 코퍼레이션(General Motors Acceptance Corporation, GMAC)의 풀이 낮은 품질이라는 것을 발견하자, 이러한 우려는 자동차 대출 시장으로 확산되었습니다.

투자자들은 자동차 대출을 담보로 하는 자산유동화기업어음을 롤오버하려는 의지가 점점 줄어들었고, 이는 캡티브 금융 회사에 가용 자금을 감소시켰습니다. GMAC의 대응은 극적이었습니다. 2008년 10월, GMAC는 신용 기준을 상당히 강화하고 신용 점수가 700점 이상인 소비자에게만 대출할 것이라고 발표했습니다. 당시 중간 신용 점수는 720점이었습니다. 이는 잠재적인 자동차 구매자 중 상당수가 신용을 거부당할 것이라는 의미였습니다. 다른 캡티브 금융 회사들도 이미 신용 기준을 강화했습니다. 그럼에도 불구하고 투자자들은 자산유동화기업어음 롤오버를 거부했고, 이 자금 조달은 완전히 고갈되었습니다.

신탁의 붕괴는 캡티브 금융 회사들이 자산유동화기업어음을 상환해야 한다는 것을 의미했습니다. GMAC와 크라이슬러 파이낸셜 인코퍼레이티드(Chrysler Financial Incorporated)는 의무를 이행할 수 없었고 붕괴되었습니다. 투자자들이 자산유동화기업어음 롤오버를 거부한 것은 본질적으로 뱅크런과 동일합니다. 모든 투자자들이 동시에 단기 자금 조달을 철회하여 신탁의 붕괴와 청산으로 이어지는 것입니다. 자산유동화기업어음 신탁에 대한 이러한 런과 그로 인한 신용 경색이 자동차 대출에 어떤 의미를 가졌는지 살펴보겠습니다.

\begin{tcolorbox}[title={\textbf{자동차 대출 및 판매에 미친 영향}}]
\textbf{차트 분석 (텍스트 기반 설명):}
\begin{itemize}
    \item \textbf{자동차 대출 성장률:} 2006년 초 모기지 시장에서 서브프라임 위기가 전개되기 시작하면서 자동차 대출의 전년 대비 성장률은 마이너스로 전환되었습니다.
    \item \textbf{ABCP 런 이후:} 2008년 중반 투자자들이 ABCP 롤오버를 중단하자, 대출 발생은 전년 대비 30\% 이상 급감했습니다.
    \item \textbf{자동차 판매:} 금융 회사에 대한 자금 조달이 고갈되기 전에는 미국 가구가 연간 약 1,600만 대의 신차를 구매했습니다. 금융 위기가 닥치자 2009년 1분기에는 연간 790만 대로 소매 자동차 판매가 붕괴되었습니다.
\end{itemize}
\textbf{수요 감소 vs. 신용 공급 감소:}
\begin{itemize}
    \item 타이밍으로 볼 때 자동차 판매 감소는 신용 공급 감소에 의해 주도되었음을 시사합니다.
    \item 하지만 수요 감소 때문이 아니었을까요? 금융 위기 동안 많은 노동자들이 일자리를 잃었고, 암울한 경제 전망은 소비자들이 신차 구매를 주저하게 만들었을 수 있습니다.
    \item 만약 붕괴가 수요에 의해 주도되었다면, 캡티브 금융 회사들이 대출을 할 수 없을 때 다른 대출 기관들도 대출을 하지 않았을 것입니다.
    \item 그러나 2008년에는 더 많은 자동차 구매자들이 현금으로 지불하거나 은행 또는 신용 조합에서 신용을 얻었습니다. 이는 ABCP 런이 자동차 판매에 상당한 영향을 미 미쳤음을 시사합니다.
\end{itemize}
\textbf{결론:} 캡티브 금융 회사의 사례는 뱅크런이 은행 부문 밖에서도 발생할 수 있음을 보여줍니다. 실제로 만기 변환에 참여하는 모든 금융 기관은 런의 대상이 될 수 있습니다.
\end{tcolorbox}

\subsubsection*{비은행 금융 기관의 규제 공백}
그러나 2008년 금융 위기 이전에 규제 당국은 비은행 금융 기관의 만기 변환 및 신용 제공을 대체로 간과했습니다. 이러한 비은행 금융 기관은 단기 자금 시장의 런을 막는 데 도움이 될 최종 대부자로서의 연방준비제도에 접근할 수 없었습니다. 이러한 런은 신용 공급에 미치는 영향을 통해 뒤따른 경기 침체의 깊이에 기여했습니다.

\subsubsection*{2008년 금융 위기 중 자동차 대출 사례의 교훈}
이 강의에서 무엇을 배웠습니까?

\begin{tcolorbox}[title={\textbf{자동차 대출 사례의 주요 교훈}}]
\begin{enumerate}
    \item \textbf{비은행 금융 기관의 만기 변환:} 캡티브 금융 회사를 예로 들어, 비은행 금융 기관도 만기 변환에 참여한다는 것을 보여주었습니다.
    \item \textbf{비은행 금융 기관의 런:} 비은행 금융 기관은 2008년 금융 위기 동안 자금 시장에서 런을 경험했습니다.
    \item \textbf{신용 공급 감소의 영향:} 런은 자동차 신용을 감소시켰고, 따라서 자동차 판매를 감소시켰습니다.
\end{enumerate}
\end{tcolorbox}


\section{Lesson 4-2: 금융 위기에 대한 정책 대응}

\subsection{Lesson 4-2.1: 최종 대부자로서의 연준}
최종 대부자(Lender of Last Resort) 기능에 대한 수업에 오신 것을 환영합니다. 1907년 공황, 대공황, 2008년 글로벌 금융 위기 등 은행 공황과 금융 위기는 금융 부문과 궁극적으로 실물 경제에 파괴적인 영향을 미칠 수 있습니다. 이 세션에서는 중앙은행이 금융 위기 동안 단기 현금 대출을 제공함으로써 금융 시장을 안정화하고 최악의 시나리오를 방지하는 데 어떻게 도움이 될 수 있는지 살펴볼 것입니다. 이러한 개입이 원칙적으로 어떻게 설계되어야 하는지, 그리고 실제로 어떻게 작동하는지 보게 될 것입니다. 또한 이러한 개입의 잠재적 비용도 이해할 것입니다. 우리는 2008년 금융 위기와 2020년 글로벌 팬데믹의 맥락에서 이러한 문제들을 검토할 것입니다.

\subsubsection*{월터 배젓의 원칙}
정부 지원을 받는 중앙은행이 금융 위기에 개입하여 위기를 막아야 한다는 생각은 19세기 영국 언론인 월터 배젓(Walter Bagehot)에게서 시작되었습니다.

1873년, 배젓은 금융 위기 시 중앙은행은 '건전한 예금 기관에 자유롭게 대출해야 하지만, 건전한 담보에 대해서만, 그리고 진정으로 필요한 차입자가 아닌 이들을 단념시킬 만큼 충분히 높은 이자율로 대출해야 한다'고 유명하게 말했습니다. 이 진술은 중앙은행이 금융 위기에 어떻게 대응해야 하는지에 대한 시금석이 되었습니다. 이것이 무엇을 의미하는지 빠르게 풀어봅시다.

\paragraph{건전한 예금 기관에 자유롭게 대출}
"중앙은행은 건전한 예금 기관에 자유롭게 대출해야 한다"는 부분부터 시작해 봅시다. 때때로 건전한 은행들도 은행 공황에 휩쓸립니다. 모든 은행은 단기적으로 차입하고 장기적으로 대출하기 때문에, 이러한 건전한 은행 중 일부는 현금이 바닥날 수 있습니다. 중앙은행이 지원해야 하는 대상은 바로 이러한 '건전한(solvent)' 은행들입니다.

\paragraph{건전한 담보와 충분히 높은 이자율}
두 번째 부분은 어떻습니까? "건전한 담보에 대해서만, 그리고 진정으로 필요한 차입자가 아닌 이들을 단념시킬 만큼 충분히 높은 이자율로 대출해야 한다." 이것은 무엇을 달성할까요?

\begin{tcolorbox}[title={\textbf{월터 배젓 원칙의 두 가지 목표}}]
\begin{enumerate}
    \item \textbf{납세자 손실 방지:} 첫째, 중앙은행은 개입할 때 손실을 입어서는 안 됩니다. 이러한 프로그램은 납세자에게 손실을 초래해서는 안 됩니다.
    \item \textbf{일시적 지원:} 둘째, 차입자들은 다른 선택지가 없을 때만 중앙은행의 도움을 받아야 합니다. 이는 일시적인 지원이어야 하며 장기적인 해결책이 되어서는 안 됩니다. 차입자들은 가능한 한 빨리 중앙은행 대출을 상환하기를 원해야 합니다.
\end{enumerate}
\textbf{결론:} 배젓은 정부가 개입해야 하지만 신중하게 해야 한다고 생각했습니다. 궁극적인 목표는 납세자에게 손실을 주지 않고 은행 공황을 막는 것입니다.
\end{tcolorbox}

이것이 이론입니다. 이제 이것이 실제로 어떻게 작동하는지 살펴보겠습니다. 최근 두 가지 위기, 즉 2008년 금융 위기와 2020년 글로벌 팬데믹에서 중앙은행의 개입을 고려해 봅시다.

\subsubsection*{2008년 금융 위기 대응}
2008년 사건부터 시작하겠습니다. 이 위기는 2006년에 시작된 미국 주택 시장의 위축으로 촉발되었습니다. 주택 가격 하락은 모기지 연체율 상승을 동반했습니다. 손실이 증가하면서 신용 공급, 특히 단기 자금 시장에서 급격한 감소가 있었습니다. 이는 유명한 투자 은행인 리먼 브라더스(Lehman Brothers)의 붕괴로 절정에 달했으며, 이는 글로벌 금융 시장에 충격을 주었습니다. 2008년 가을까지 미국은 깊은 경기 침체 직전에 있었습니다.

\paragraph{연방준비제도의 대응}
연방준비제도는 어떻게 대응했을까요? 첫 번째로 한 일은 2007년 중반에 금리를 5.25\%에서 사실상 0\%로 인하한 것이었습니다. 또한 연방준비제도는 금융 기관과 신용 시장을 안정화하기 위해 다양한 비상 유동성 프로그램(emergency liquidity programs)을 설정했습니다. 연준은 리먼 브라더스 파산 이후의 공황을 막고 싶었습니다. 이것이 연준이 최종 대부자로서 행동한 것입니다.

전통적으로 연방준비제도는 '할인 창구(discount window)'를 통해 은행에만 비상 현금 대출을 제공했을 것입니다. 그러나 2008년에는 이것이 작동하지 않았습니다. 은행만이 유일한 플레이어가 아니었습니다. 비은행 금융 기관(non-banks)이 많은 자금 시장에서 중요한 역할을 했고, 이들 비은행 기관 중 다수가 공황에 휩쓸렸습니다. 예를 들어, 리먼 브라더스는 투자 은행이었습니다. 할인 창구에 접근할 수 없었습니다.

자금 시장을 안정화하기 위해 연방준비제도는 세 가지 차원에서 할인 창구의 범위를 확장함으로써 단호하게 행동했습니다.

\begin{tcolorbox}[title={\textbf{2008년 연준의 할인 창구 확장 프로그램}}]
\begin{enumerate}
    \item \textbf{TAF (Term Auction Facility):} 은행의 장기 자금 조달이 고갈되자 연방준비제도는 TAF를 설립했습니다. 이는 은행에 장기 유동성을 경매 방식으로 제공하여 자금 조달 압박을 완화했습니다.
    \item \textbf{TALF (Term Asset-backed-securities Loan Facility):} 리먼 브라더스 파산 이후 증권화 시장이 붕괴되자 연방준비제도는 이 시장을 해동하기 위해 TALF를 설립했습니다. TALF 참여자들은 연준으로부터 직접 장기 대출을 받아 증권화된 대출을 구매할 수 있었습니다.
    \item \textbf{PDCF (Primary Dealer Credit Facility) \& TSLF (Term Securities Lending Facility):} 자금 시장의 핵심인 많은 비은행 기관들이 절실히 현금을 필요로 했습니다. 가장 중요하고 평판 좋은 투자 은행들은 '프라이머리 딜러(primary dealers)'라고도 불립니다. 프라이머리 딜러는 연방준비제도의 공개 시장 운영의 거래 상대방 역할을 합니다. 이들은 통화 정책에 매우 중요합니다. 또한 잘 기능하는 국채 시장에도 필수적입니다. PDCF와 TSLF는 이들의 필요를 충족시키기 위해 설립되었습니다.
    \item \textbf{AMLF (Asset-Backed Commercial Paper Money Market Mutual Fund Liquidity Facility):} 머니 마켓 펀드(money market funds)가 런을 경험하고 있었기 때문에 연준은 AMLF를 설립했습니다. 이는 머니 마켓 펀드에 유동성을 제공하여 투자자들의 대규모 인출을 막았습니다.
    \item \textbf{CPFF (Commercial Paper Funding Facility):} 마지막으로 CPFF는 미국 상업어음 발행자들에게 지원을 제공했습니다. 이것은 2008년의 유일한 직접 구매 프로그램의 예입니다. 이는 연방준비제도가 상업어음을 직접 구매했다는 것을 의미합니다. 이것은 주로 금융 기관이 발행한 것이었지만, 연준은 할리 데이비슨(Harley Davidson)이나 맥도날드(McDonald's)와 같은 비금융 기관이 발행한 어음도 구매했습니다. 역사적 기준으로 볼 때, 이것은 매우 이례적인 조치였습니다.
\end{enumerate}
\end{tcolorbox}

\paragraph{2008년 개입의 성공 여부와 비용}
그것은 크고 복잡한 개입이었습니다. 성공적이었을까요? 월터 배젓은 만족했을까요?

\begin{tcolorbox}[title={\textbf{2008년 연준 개입의 평가}}]
\textbf{긍정적 측면:}
\begin{itemize}
    \item \textbf{금융 공황 중단:} 무엇보다도 연방준비제도는 금융 공황을 막았습니다.
    \item \textbf{손실 없음:} 각 프로그램에서 비상 대출은 고품질 담보로 뒷받침되었고 시장 금리 이상의 이자율이 부과되었습니다. 대출은 신속하게 상환되었고, 프로그램은 손실 없이 종료되었습니다.
\end{itemize}
\textbf{잠재적 비용 및 우려:}
\begin{itemize}
    \item \textbf{정부 안전망의 확장:} 다른 한편으로, 이것은 정부 안전망의 대규모적이고 전례 없는 확장(unprecedented expansion)이었습니다. 지원을 받은 금융 기관 중 다수는 예금 기관이 아니었습니다. 그들은 연방준비제도나 다른 은행 규제 기관의 감독이나 규제를 받지 않았습니다.
    \item \textbf{도덕적 해이(Moral Hazard) 우려:} 만약 시장 참여자들이 상황이 나빠질 때 연방준비제도가 계속해서 지원을 제공할 것이라고 기대한다면, 이것이 위험 감수 행동에 어떤 영향을 미칠까요? 이 질문에 대한 명확한 답은 없습니다.
\end{itemize}
\end{tcolorbox}

\subsubsection*{2020년 글로벌 팬데믹 대응}
2020년으로 빠르게 넘어가 봅시다. 글로벌 팬데믹은 2020년 1분기에 미국에 도착했습니다. 2020년 3월까지 국채 시장에 심각한 변동성이 있었고, 이는 단기 자금 시장에 혼란을 야기할 가능성이 있었습니다. 이는 전 세계적으로 경고음을 울렸습니다. 2008년 금융 위기가 얼마나 심각해질 수 있는지를 보여주었기 때문에, 정책 입안자들은 국채 시장에 대한 첫 번째 지원을 제공하고 더 광범위하게 지원하기 위해 서둘렀습니다. 이전과 마찬가지로, 여기에는 최종 대부자로서 행동하는 연방준비제도가 포함되었습니다.

충격은 매우 달랐지만, 중앙은행의 대응은 유사했습니다. 2008년 금융 위기 동안 사용되었던 여러 유동성 시설이 다시 가동되었습니다. 여기에는 TALF(Term-Asset-Backed Securities Loan Facility), CPFF(Commercial Paper Funding Facility), MMLF(Money Market Mutual Fund Liquidity Facility)가 포함되었습니다.

\begin{tcolorbox}[title={\textbf{2020년 연준의 주요 유동성 프로그램}}]
\begin{enumerate}
    \item \textbf{PPPLF (Paycheck Protection Program Liquidity Facility):} 아마도 가장 큰 프로그램은 PPPLF였을 것입니다. 급여 보호 프로그램(Paycheck Protection Program)은 중소기업에 6천억 달러의 상환 가능한 대출을 제공하는 정부 프로그램이었습니다. PPPLF에 따라 대출 기관은 중앙은행으로부터 이러한 대출에 대한 자금을 받을 수 있었습니다.
    \item \textbf{PSMCCF (Primary and Secondary Market Corporate Credit Facilities):} 직접 구매 프로그램은 이번에는 훨씬 더 큰 역할을 했습니다. PSMCCF는 회사채 시장을 통해 대기업에 직접적인 지원을 제공했습니다. 연준은 대규모로 회사채를 직접 구매했습니다.
    \item \textbf{Main Street Lending Program:} 2020년 4월까지 연방준비제도는 중소기업에 대한 직접 대출에도 관여했습니다. 메인 스트리트 대출 프로그램(Main Street Lending Program)은 연준이 은행과 협력하여 중소기업 대출을 시작할 수 있도록 했습니다.
\end{enumerate}
\end{tcolorbox}

\paragraph{2020년 개입의 교훈}
이것은 우리에게 무엇을 남길까요? 2020년 사건에서 최종 대부자에 대해 무엇을 더 배울 수 있을까요?

\begin{tcolorbox}[title={\textbf{2020년 연준 개입의 평가 및 교훈}}]
\textbf{긍정적 측면:}
\begin{itemize}
    \item \textbf{경제 붕괴 방지:} 이러한 개입은 경제의 완전한 붕괴를 피하는 데 성공한 것으로 보입니다.
    \item \textbf{재정-통화 정책 협력:} 위기 개입은 재정 정책과 통화 정책이 협력할 때 가장 잘 작동하는 것으로 보입니다. 아마도 가장 성공적인 정부 개입은 급여 보호 프로그램이었을 것입니다. 연방준비제도가 이러한 대출에 6천억 달러의 자금을 제공한 것이 위기 시기에 얼마나 중요했는지 쉽게 알 수 있습니다. 이 시점에서 은행들은 추가 대출을 하기보다는 현금을 비축했을 것입니다.
\end{itemize}
\textbf{월터 배젓의 관점과 우려:}
\begin{itemize}
    \item 월터 배젓은 이 모든 것에 대해 어떻게 생각했을까요? 그는 중앙은행이 예금 기관만 지원해야 한다고 생각했습니다. 그것은 과거의 일이 된 것 같습니다.
    \item 그는 또한 납세자에게 잠재적인 손실이 발생할까 봐 우려했습니다. 불행히도 2008년에서 2020년으로 가면서 연방준비제도가 중소기업 및 대기업에 대한 직접 대출에 훨씬 더 많이 관여하는 것을 보았습니다. 이는 물론 기업 부문에서 채무 불이행이 발생할 경우 연방준비제도가 직접적인 손실에 노출된다는 것을 의미합니다. 이는 분명히 배젓을 불안하게 만들 것입니다.
\end{itemize}
\end{tcolorbox}

\subsubsection*{최종 대부자 역할의 진화}
이 세션에서 무엇을 배웠습니까?

\begin{tcolorbox}[title={\textbf{최종 대부자 역할의 주요 교훈}}]
\begin{enumerate}
    \item \textbf{위기 중단 능력:} 최종 대부자는 비상 대출 프로그램을 통해 금융 위기를 막을 수 있습니다.
    \item \textbf{역사적 원칙:} 역사적으로 이러한 자금은 은행으로 흘러갔고, 대출은 고품질 담보와 시장 금리 이상의 이자율을 요구했습니다.
    \item \textbf{최근 개입 사례:} 우리는 연방준비제도가 2008년과 2020년에 시장을 진정시키기 위해 어떻게 개입했는지 보았습니다.
    \item \textbf{프로그램 범위 확장:} 마지막으로, 이러한 비상 대출 프로그램의 범위가 최근 몇 년 동안 비은행 기관뿐만 아니라 직접 대출 프로그램까지 포함하도록 어떻게 확장되었는지 논의했습니다.
\end{enumerate}
\end{tcolorbox}


\subsection{Lesson 4-2.2: 금융 위기에 대한 재정적 대응}
이 강의는 금융 위기에 대한 재정적 대응에 대해 다룹니다. (이 파트 1/2 텍스트에는 해당 강의의 내용이 포함되어 있지 않습니다. 전체 모듈 4의 파트 2/2에서 다루어질 예정입니다.)


\section{개요: 재정 위기에 대한 재정 대응}

이 문서는 금융 위기와 그에 따른 경기 침체에 정부가 어떻게 재정 정책(Fiscal Policy)으로 대응하는지 설명합니다. 통화 정책(Monetary Policy)만으로는 경제를 완전히 회복시키기 어렵기 때문에, 정부는 다양한 개입을 통해 경제 회복을 지원합니다. 우리는 세 가지 주요 재정 대응 방안을 살펴보고, 각 방안의 효과를 2008년 금융 위기와 2020년 팬데믹 사례를 통해 분석합니다.

\begin{itemize}
    \item \textbf{주요 목표:} 금융 위기로 인한 경제적 피해를 줄이고 경기 침체에서 벗어나기 위한 재정 정책의 역할 이해.
    \item \textbf{핵심 개념:} 정부의 소비자 직접 이전(세금 감면, 경기 부양 수표), 기업 지원, 직접 지출(인프라 프로젝트).
    \item \textbf{평가 포인트:} 각 재정 대응 방안의 효과성, 케인즈주의적 관점과 반대 관점, 통화 정책의 보조적 역할.
\end{itemize}


\section{용어 정리}

다음은 본 문서에서 사용되는 주요 경제 용어와 그에 대한 쉬운 설명입니다.

\begin{adjustbox}{width=\textwidth,center}
\begin{tabular}{lllc}
    \toprule
    \textbf{용어} & \textbf{쉬운 설명} & \textbf{원어} & \textbf{비고} \\
    \midrule
    재정 정책 & 정부가 세금과 지출을 조절하여 경제를 안정시키려는 정책 & Fiscal Policy & 통화 정책과 대비 \\
    통화 정책 & 중앙은행이 이자율과 통화량을 조절하여 경제를 안정시키려는 정책 & Monetary Policy & 재정 정책과 대비 \\
    금융 위기 & 금융 시스템이 불안정해져 경제 전반에 큰 충격을 주는 상황 & Financial Crisis & 2008년, 2020년 사례 \\
    경기 침체 & 경제 활동이 전반적으로 위축되고 생산과 고용이 감소하는 시기 & Recession & GDP 감소, 실업률 증가 \\
    GDP & 한 나라에서 일정 기간 생산된 모든 최종 재화와 서비스의 시장 가치 합계 & Gross Domestic Product & 경제 규모를 나타내는 지표 \\
    소비 (C) & 가계가 재화와 서비스를 구매하는 지출 & Consumption & GDP의 가장 큰 구성 요소 \\
    투자 (I) & 기업이 생산 설비 등에 지출하거나 가계가 주택을 구매하는 지출 & Investment & 미래 생산 능력 증대 \\
    정부 지출 (G) & 정부가 재화와 서비스를 구매하는 지출 & Government Spending & 인프라 건설, 공무원 급여 등 \\
    순수출 (NX) & 수출에서 수입을 뺀 값 & Net Exports & (수출 - 수입) \\
    경기 부양 수표 & 정부가 경제 활성화를 위해 가계에 직접 지급하는 현금 & Stimulus Checks & 소비 진작 목적 \\
    실업 급여 & 실직자에게 정부가 지급하는 소득 보전금 & Unemployment Benefits & 소비 감소 방지 \\
    임금 보조금 & 정부가 기업의 임금 일부를 지원하여 해고를 막는 정책 & Wage Bill Subsidies & 고용 유지 목적 \\
    급여 보호 프로그램 & 팬데믹 시기 미국에서 기업의 고용 유지를 위해 제공된 대출 프로그램 & Paycheck Protection Program (PPP) & 대출 상환 면제 조건 \\
    정부 지출 승수 & 정부 지출 1달러가 GDP를 얼마나 증가시키는지 나타내는 값 & Government Spending Multiplier & 케인즈주의 이론의 핵심 \\
    \bottomrule
\end{tabular}
\end{adjustbox}


\section{핵심 개념 및 원리: 재정 위기 대응 방안}

금융 위기와 경기 침체 시, 경제를 회복시키기 위해 정부는 재정 정책을 활용합니다. 통화 정책만으로는 한계가 있기 때문에, 정부는 다양한 개입을 통해 경제 회복을 지원합니다.

\subsection{1. GDP의 구성 요소와 소비의 중요성}

\begin{tcolorbox}[mybox={GDP의 구성 요소}]
    \textbf{한 줄 핵심 요약:} GDP는 소비, 투자, 정부 지출, 순수출의 합으로, 이 중 소비가 가장 큰 비중을 차지합니다.
\end{tcolorbox}

경제의 총 생산량, 즉 국내총생산(GDP)은 항상 다음 네 가지 요소로 구성됩니다:
\begin{itemize}
    \item \textbf{민간 소비 (C, Private Consumption):} 가계가 재화와 서비스를 구매하는 지출입니다.
    \item \textbf{기업 투자 (I, Business Investment):} 기업이 공장, 기계 등 생산 설비에 지출하거나 가계가 주택을 구매하는 지출입니다.
    \textbf{정부 지출 (G, Government Spending):} 정부가 재화와 서비스를 구매하는 지출입니다.
    \item \textbf{순수출 (NX, Net Exports):} 수출에서 수입을 뺀 값입니다. (수출 - 수입)
\end{itemize}

\begin{examplebox}[title={예시: 2019년 미국 GDP 구성}]
    2019년 미국의 경우, GDP에서 소비가 약 68\%를 차지했습니다. 기업 투자는 17\%, 정부 지출은 19\%, 그리고 순수출은 -4\%였습니다. 순수출이 마이너스라는 것은 미국이 수출보다 수입을 더 많이 했다는 의미입니다. 이처럼 소비는 GDP에서 가장 큰 비중을 차지하는 요소입니다.
\end{examplebox}

\begin{tcolorbox}[mybox={경기 침체 시 소비의 변화}]
    \textbf{한 줄 핵심 요약:} 경기 침체 시 소비는 급격히 감소하며, 이는 GDP에 막대한 영향을 미칩니다.
\end{tcolorbox}

경기가 좋을 때는 소비가 꾸준히 증가하지만, 금융 위기나 팬데믹과 같은 어려운 시기에는 소비가 급격히 감소합니다. 예를 들어, 2008년 금융 위기와 그 여파, 그리고 2020년 글로벌 팬데믹 시기에는 소비가 크게 줄었습니다. 이 두 경우 모두 실업률이 급증했고, 이는 가계의 소득 감소로 이어져 소비를 더욱 위축시켰습니다. 소비가 GDP의 3분의 2를 차지하기 때문에, 소비의 급락은 경제 전반에 엄청난 충격을 줍니다. 따라서 재정 위기에 대한 정부의 대응은 주로 소비를 지탱하는 데 초점을 맞춥니다.

\subsection{2. 재정 대응 방안 1: 소비자에게 직접 이전 (Government Transfers to Consumers)}

정부가 소비를 지탱하기 위해 가장 먼저 고려하는 방법 중 하나는 가계에 직접 돈을 이전하는 것입니다.

\subsubsection*{2.1. 세금 감면 vs. 경기 부양 수표}

\begin{tcolorbox}[mybox={세금 감면의 한계}]
    \textbf{한 줄 핵심 요약:} 실업자에게는 소득이 없으므로, 소득세 감면은 경기 침체 시 소비 진작에 효과적이지 않습니다.
\end{tcolorbox}

새롭게 실업자가 된 사람들은 소득이 없기 때문에, 소득세를 감면해주는 정책은 이들에게 아무런 도움이 되지 않습니다. 이들은 소비를 가장 많이 줄이는 계층이므로, 이들을 직접 지원하는 다른 방법이 필요합니다.

\begin{tcolorbox}[mybox={경기 부양 수표의 도입}]
    \textbf{한 줄 핵심 요약:} 정부는 모든 가구에 경기 부양 수표를 보내 소비를 직접적으로 늘리려 합니다.
\end{tcolorbox}

세금 감면의 한계를 보완하기 위해 정부는 모든 가구에 경기 부양 수표(Stimulus Checks)를 보냅니다. 이는 가계에 직접 현금을 지급하여 소비를 유도하는 방식입니다.

\begin{examplebox}[title={예시: 미국 경기 부양 수표 지급 사례}]
    \begin{itemize}
        \item \textbf{2008년:} 개인 납세자는 최대 600달러, 부부는 최대 1,200달러를 받았습니다.
        \item \textbf{2020년:} 개인 납세자는 최대 1,200달러, 부부는 최대 2,400달러를 받았습니다.
    \end{itemize}
\end{examplebox}

\begin{tcolorbox}[mybox={경기 부양 수표의 효과성}]
    \textbf{한 줄 핵심 요약:} 경기 부양 수표는 소비 진작에 부분적으로 효과적이지만, 모든 금액이 소비로 이어지지는 않습니다.
\end{tcolorbox}

경기 부양 수표가 과연 효과적일까요? 연구 결과에 따르면, 2020년 지급된 경기 부양 수표의 경우 약 42\%만이 소비에 사용되었고, 27\%는 저축되었으며, 31\%는 부채 상환에 사용되었습니다.

\begin{cautionbox}{왜 모든 돈이 소비되지 않았을까요?}
    \begin{itemize}
        \item \textbf{소득 수준에 따른 차이:} 저소득층은 수표의 대부분을 소비했지만, 고소득층은 주로 저축했습니다. 이는 고소득층이 이미 충분한 소비 여력이 있거나, 미래의 불확실성에 대비해 저축을 선호하기 때문입니다.
        \item \textbf{소비 기회의 제한:} 2020년에는 미국 경제가 봉쇄(Lockdown)되어 소비할 기회가 제한적이었습니다. 상점들이 문을 닫고 외식이나 여행이 어려워지면서, 돈을 쓰고 싶어도 쓸 곳이 마땅치 않았던 것입니다.
    \end{itemize}
\end{cautionbox}

\subsubsection*{2.2. 목표 지향적 직접 지급: 실업 급여 증액 및 고용 유지 프로그램}

\begin{tcolorbox}[mybox={실업 급여의 효과성}]
    \textbf{한 줄 핵심 요약:} 실업 급여 증액은 실직 가구의 소비 감소를 직접적으로 막아주어 매우 효과적인 개입입니다.
\end{tcolorbox}

경기 부양 수표보다 더 목표 지향적인(Targeted) 직접 지급 방식은 실업 급여(Unemployment Benefits)를 늘리는 것입니다.
\begin{itemize}
    \item \textbf{2009년:} 실업 급여 지급 기간이 연장되었습니다.
    \item \textbf{코로나19 팬데믹 시기:} 실업 급여가 증액되고 지급 기간이 다시 연장되었습니다.
\end{itemize}
이러한 개입은 최근 일자리를 잃은 사람들에게만 지급되므로, 소비를 가장 많이 줄이는 가구에 직접적인 도움을 줍니다. 이는 경기 부양 수표보다 효율적일 수 있습니다.

\begin{tcolorbox}[mybox={고용 유지 프로그램의 장점}]
    \textbf{한 줄 핵심 요약:} 기업이 직원을 해고하지 않도록 임금의 일부를 지원하는 프로그램은 고용 관계를 유지하고, 직원의 기술 손실을 막으며, 경제 회복 시 재고용 비용을 줄여줍니다.
\end{tcolorbox}

또 다른 접근 방식은 기업에 돈을 지급하여 직원을 해고하지 않도록 하는 것입니다. 정부가 기업의 임금 청구서(Wage Bill)의 일부를 부담하는 방식입니다.

\begin{examplebox}[title={예시: 독일과 미국의 고용 유지 프로그램}]
    \begin{itemize}
        \item \textbf{독일 (2008년 금융 위기):} 이 접근 방식은 2008년 금융 위기 당시 독일에서 처음 시도되었습니다.
        \item \textbf{미국 (코로나19 팬데믹):} 미국의 급여 보호 프로그램(Paycheck Protection Program, PPP)은 이와 유사한 목적을 가졌습니다. 이 프로그램은 중소기업에 대출을 제공했는데, 만약 기업이 직원을 해고하지 않으면 이 대출을 탕감해 주었습니다.
    \end{itemize}
\end{examplebox}

이러한 프로그램의 주요 장점은 다음과 같습니다:
\begin{itemize}
    \item \textbf{고용 관계 유지:} 고용주와 직원 간의 관계를 그대로 유지합니다.
    \textbf{기술 손실 방지:} 직원의 기술이 감가상각(Depreciate)되지 않습니다. 즉, 일자리를 잃고 쉬는 동안 기술이 녹슬지 않습니다.
    \item \textbf{재고용 비용 절감:} 수요가 다시 증가했을 때 기업이 직원을 다시 고용할 필요가 없어 재고용에 드는 시간과 비용을 절약할 수 있습니다.
    \item \textbf{소비 유지:} 직원들이 소득을 유지하므로 소비를 줄일 필요가 없어집니다.
\end{itemize}

\begin{summarybox}[title={소비자 직접 이전의 결론}]
    가계에 대한 직접 이전은 경기 침체 시 소비와 GDP를 지탱하는 데 도움이 됩니다. 특히, 새로 실직한 사람들을 대상으로 하거나 고용을 유지하는 프로그램이 가장 효과적입니다.
\end{summarybox}


\subsection{3. 재정 대응 방안 2: 기업 지원 (투자 촉진)}

GDP의 두 번째 주요 구성 요소는 투자(Investment)입니다.

\begin{tcolorbox}[mybox={경기 침체 시 투자의 변화}]
    \textbf{한 줄 핵심 요약:} 경기 침체 시 투자는 급격히 감소하며, 정부의 개입 옵션은 제한적입니다.
\end{tcolorbox}

투자는 2008년 금융 위기 때 급락했고, 2020년에도 다시 감소했습니다. 이러한 투자 감소에 대응하기 위해 정부는 다음과 같은 제한적인 옵션을 가지고 있습니다.

\begin{itemize}
    \item \textbf{세금 감면 (Tax Breaks):} 기업의 투자 비용을 줄여주기 위해 세금을 감면해 줄 수 있습니다.
    \item \textbf{보조금 (Subsidies):} 특정 투자 활동에 대해 직접적인 보조금을 지급할 수 있습니다.
\end{itemize}

\begin{cautionbox}{투자의 근본적인 동인}
    결국 투자는 경제의 전반적인 상태에 의해 좌우됩니다. 예를 들어, 기업 투자는 상품에 대한 수요에 의해 결정됩니다. 수요가 없으면 아무리 세금 감면이나 보조금을 줘도 기업은 투자를 늘리지 않을 것입니다. 따라서 정부의 직접적인 투자 촉진 정책은 한계가 있을 수 있습니다.
\end{cautionbox}

\subsection{4. 재정 대응 방안 3: 정부의 직접 지출 (Direct Government Spending)}

\begin{tcolorbox}[mybox={정부 지출의 역할}]
    \textbf{한 줄 핵심 요약:} 정부는 인프라 프로젝트와 같은 직접 지출을 늘려 일자리를 창출하고 경제 활동을 촉진할 수 있습니다.
\end{tcolorbox}

정부가 경기 침체에 대응하는 마지막 유형의 재정 정책은 정부 자체의 지출을 늘리는 것입니다.

\begin{examplebox}[title={예시: 미국 정부의 직접 지출 프로그램}]
    \begin{itemize}
        \item \textbf{2009년 미국 경기 회복 및 재투자법 (American Recovery and Reinvestment Act of 2009):} 이 법안은 2,750억 달러를 연방 계약, 보조금, 대출에 할당하여 일자리를 창출했습니다.
        \item \textbf{2021년 인프라 투자 및 일자리 법 (Infrastructure Investment and Jobs Act):} 글로벌 팬데믹 이후 통과된 이 법안은 전국적으로 많은 공공 건설 프로젝트를 포함했습니다.
    \end{itemize}
\end{examplebox}

이러한 정부 지출 법안들은 주로 도로, 교량, 공항 등 공공 인프라 건설 프로젝트를 포함합니다. 하지만 정부 지출을 늘리는 것이 경기 침체에 맞서는 효과적인 도구일까요?

\subsubsection*{4.1. 정부 지출 승수 (Government Spending Multiplier)}

\begin{tcolorbox}[mybox={케인즈주의적 관점}]
    \textbf{한 줄 핵심 요약:} 케인즈는 경기 침체 시 정부 지출이 GDP를 1달러 이상 증가시키는 승수 효과를 가진다고 주장했습니다.
\end{tcolorbox}

영국의 경제학자 존 메이너드 케인즈(John Maynard Keynes)는 경기 침체 시 소비가 감소하면 정부가 직접 상품을 구매하여 이 손실된 수요를 보충해야 한다고 주장했습니다. 케인즈는 이것이 GDP를 안정화시킬 것이라고 보았습니다. 여기서 핵심 질문은 "정부가 1달러를 지출했을 때 GDP가 1달러보다 더 많이 증가하는가?"입니다. 이것을 \textbf{정부 지출 승수(Government Spending Multiplier)}라고 부릅니다.

케인즈주의적 관점에서는 경기 침체 시 정부 지출 승수가 1보다 크다고 봅니다. 그 이유는 정부가 손실된 수요를 보충하면 일자리가 창출되고, 이는 다시 소비를 증가시키기 때문입니다. 즉, 정부가 1달러를 지출하면 그 돈이 누군가의 소득이 되고, 그 소득의 일부가 다시 소비되어 또 다른 사람의 소득이 되는 연쇄 효과가 발생한다는 것입니다.

\begin{tcolorbox}[mybox={케인즈주의 반대 관점}]
    \textbf{한 줄 핵심 요약:} 반대론자들은 정부 지출이 민간 수요를 단순히 대체할 뿐이며, 비효율적인 자원 재배치를 초래하여 승수 효과가 1보다 작다고 주장합니다.
\end{tcolorbox}

케인즈주의에 반대하는 사람들은 정부 지출 승수가 1보다 작다고 주장합니다. 이들은 정부가 단순히 민간 소비를 대체할 수 없다고 봅니다. 정부는 민간이 원하지 않는 다른 종류의 상품에 대한 수요를 창출할 수 있다는 것입니다.

\begin{examplebox}[title={예시: 비효율적인 자원 재배치}]
    주택 부문에서 발생한 충격에 대응하여 정부가 군사 지출을 늘린다고 가정해 봅시다. 건설 부문의 노동자들은 국방 부문에서 일자리를 찾아야 할 것입니다. 하지만 건설 노동자들이 국방 산업에 필요한 기술을 가지고 있지 않을 수 있습니다. 노동자들의 기술을 재교육하고 산업 간에 재배치하는 것은 비용이 많이 들고 비효율적일 수 있습니다. 따라서 정부 지출 승수는 1보다 작다는 주장이 나옵니다.
\end{examplebox}

\begin{summarybox}[title={정부 지출 승수에 대한 증거}]
    대부분의 연구는 경기 침체 시 정부 지출 승수가 1보다 크다는 것을 시사하지만, 증거는 혼합되어 있습니다. 예상대로, 그 효과는 돈이 어떻게 사용되는지에 따라 달라집니다.
\end{summarybox}

\subsubsection*{4.2. 정부 지출의 실제적 어려움}

\begin{cautionbox}{정부 지출의 시간 지연 및 측정의 어려움}
    \textbf{한 줄 핵심 요약:} 경기 침체 시 정부 지출은 즉각적으로 이루어져야 하지만, 대규모 프로젝트는 계획, 승인, 실행에 시간이 오래 걸리며, 그 효과를 측정하기도 어렵습니다.
\end{cautionbox}

경기 침체 시 정부는 즉시 지출해야 하지만, 실제로는 프로젝트에 시간이 걸리는 경우가 많습니다.
\begin{itemize}
    \item \textbf{시간 지연:} 대규모 인프라 프로젝트를 생각해 보세요. 신중한 계획, 이해관계자들의 동의, 정부 구매 절차를 거쳐야 합니다. 이 모든 과정은 프로그램의 실행을 지연시킬 가능성이 큽니다.
    \item \textbf{효과 측정의 어려움:} 이러한 지출의 효과를 측정하는 것도 어렵습니다. 예를 들어, 수리된 다리의 이점을 어떻게 측정할 수 있을까요? 교통량 증가, 안전 개선, 경제 활동 촉진 등 다양한 측면을 고려해야 하지만, 이를 정량화하기는 쉽지 않습니다.
\end{itemize}

\begin{summarybox}[title={정부 지출의 결론}]
    결론적으로, 특정 조건 하에서는 정부 지출을 늘리는 것이 경기 침체 시 경제를 지원하는 효과적인 방법이 될 수 있습니다.
\end{summarybox}


\section{통화 정책의 보조적 역할}

\begin{tcolorbox}[mybox={재정 정책을 지원하는 통화 정책}]
    \textbf{한 줄 핵심 요약:} 중앙은행은 이자율을 낮추고 금융 기관에 자금을 제공하여 정부의 재정 부양책을 지원하고 그 효과를 증폭시킬 수 있습니다.
\end{tcolorbox}

마지막으로, 통화 정책은 이러한 재정 대응을 지원하는 데 중요한 역할을 할 수 있습니다. 정부의 세 가지 유형의 개입(소비자 이전, 기업 지원, 직접 지출)은 모두 비용이 듭니다. 중앙은행은 이자율을 통제함으로써 정부의 경기 부양 프로그램을 두 가지 방식으로 지원할 수 있습니다.

\begin{enumerate}
    \item \textbf{정부의 차입 비용 직접 인하:} 중앙은행이 이자율을 낮추면 정부가 돈을 빌리는 비용이 직접적으로 줄어듭니다. 이는 정부가 더 많은 재정 지출을 할 수 있도록 재정적 여유를 제공합니다.
    \item \textbf{민간 부문 활동 증폭:} 이자율이 낮아지면 기업이 투자하고 가계가 소비를 위해 대출을 받는 것이 더 저렴해집니다. 이는 정부의 경기 부양책 효과를 증폭시킬 수 있습니다. 예를 들어, 정부가 인프라 프로젝트를 시작하면 기업은 낮은 이자율로 자금을 조달하여 프로젝트에 참여하고, 이는 다시 고용과 소비를 늘리는 선순환을 만듭니다.
\end{enumerate}

중앙은행은 또한 금융 기관에 자금을 제공하여 경기 부양 프로그램을 지원할 수 있습니다.

\begin{examplebox}[title={예시: 급여 보호 프로그램 유동성 지원 시설 (2020년)}]
    2020년의 급여 보호 프로그램 유동성 지원 시설(Paycheck Protection Program Liquidity Facility)이 대표적인 예입니다. 이 프로그램 하에서 대출 기관들은 중앙은행으로부터 자금을 받아 정부의 급여 보호 프로그램(PPP)에 따라 중소기업에 제공되는 보조금을 지원했습니다. 중앙은행의 이러한 지원이 없었다면, 자체적인 제약에 직면한 대출 기관들은 정부 프로그램에 협조하기를 꺼렸을 수도 있습니다.
\end{examplebox}


\section{빠르게 훑어보기: 재정 위기 대응 요약}

\begin{summarybox}[title={재정 위기 대응의 핵심}]
    \textbf{목표:} 금융 위기로 인한 최악의 시나리오를 피하고 경제를 지원하는 것.
\end{summarybox}

\begin{tcolorbox}[mybox={세 가지 주요 재정 대응 방안}]
    \begin{itemize}
        \item \textbf{수요 진작:}
        \begin{itemize}
            \item 가계에 직접 돈을 이전 (경기 부양 수표, 실업 급여 증액).
            \item 정부 지출 증가 (인프라 프로젝트).
        \end{itemize}
        \item \textbf{목표 지향적 개입의 효과성:}
        \begin{itemize}
            \item 임금 보조금 지급이나 실업 급여 증액과 같이 특정 계층이나 기업을 대상으로 하는 개입이 가장 효과적입니다. 이는 소비 감소를 직접적으로 막고 고용을 유지하는 데 기여합니다.
        \end{itemize}
        \item \textbf{통화 정책의 지원:}
        \begin{itemize}
            \item 중앙은행의 지원적인 통화 정책(이자율 인하, 금융 기관 유동성 공급)은 정부의 재정 부양 프로그램을 촉진하고 그 효과를 증폭시킬 수 있습니다.
        \end{itemize}
    \end{itemize}
\end{tcolorbox}


\section{FAQ: 자주 묻는 질문}

\begin{tcolorbox}[mybox={Q\&A: 재정 위기 대응}]
    \textbf{Q1: 통화 정책만으로는 경기 침체를 극복할 수 없나요?}
    \textbf{A1:} 통화 정책은 이자율 조절을 통해 경제를 자극할 수 있지만, 금융 위기나 심각한 경기 침체 시에는 가계와 기업의 신뢰가 낮아져 이자율 인하만으로는 투자와 소비를 충분히 늘리기 어렵습니다. 이때 정부의 직접적인 재정 개입이 필요합니다.

    \textbf{Q2: 경기 부양 수표는 왜 모든 사람에게 똑같이 효과적이지 않나요?}
    \textbf{A2:} 경기 부양 수표는 저소득층에게는 즉각적인 소비로 이어지는 경향이 강하지만, 고소득층은 이미 소비 여력이 있거나 미래 불확실성에 대비해 저축하거나 부채를 상환하는 데 사용하는 경향이 있습니다. 또한, 경제 봉쇄와 같은 상황에서는 소비할 기회 자체가 부족할 수 있습니다.

    \textbf{Q3: 정부 지출 승수가 1보다 크다는 것이 정확히 무엇을 의미하나요?}
    \textbf{A3:} 정부 지출 승수가 1보다 크다는 것은 정부가 1달러를 지출했을 때, 그 1달러가 직접적인 GDP 증가 외에 추가적인 경제 활동(예: 일자리 창출, 소득 증가, 추가 소비)을 유발하여 최종적으로 GDP가 1달러보다 더 많이 증가한다는 의미입니다. 이는 연쇄적인 경제 효과를 나타냅니다.

    \textbf{Q4: 정부 지출이 항상 좋은가요? 단점은 없나요?}
    \textbf{A4:} 정부 지출은 경기 침체 시 수요를 보충하고 일자리를 창출하는 긍정적인 효과가 있지만, 단점도 있습니다. 대규모 프로젝트는 계획 및 실행에 시간이 오래 걸려 적절한 시기를 놓칠 수 있고, 비효율적인 자원 배분으로 이어질 수 있으며, 정부 부채를 증가시키는 요인이 될 수 있습니다.

    \textbf{Q5: 중앙은행은 재정 정책을 어떻게 지원하나요?}
    \textbf{A5:} 중앙은행은 주로 이자율을 낮춰 정부의 차입 비용을 줄이고, 기업과 가계의 투자 및 소비 대출을 저렴하게 만들어 재정 부양책의 효과를 증폭시킵니다. 또한, 금융 기관에 유동성을 공급하여 정부 프로그램(예: PPP)이 원활하게 작동하도록 돕습니다.
\end{tcolorbox}


\end{document}
